\input{preamble}

% OK, start here.
%
\begin{document}

\title{Homologie de Chow et classes de Chern}

\maketitle

\phantomsection
\label{section-phantom}


\tableofcontents


\section{Introduction}
\label{section-introduction}

\noindent
Dans ce chapitre, nous étudions les groupes d'homologie de Chow et la construction
des classes de Chern des fibrés vectoriels comme éléments des groupes
de cohomologie de Chow opérationnelle (le tout à coefficients dans $\mathbf{Z}$).

\medskip\noindent
Nous commençons ce chapitre en donnant la démonstration algébrique la plus courte
possible du lemme clé \ref{lemma-milnor-gersten-low-degree}.
Nous définissons d'abord le quotient de Herbrand
(section \ref{section-periodic-complexes})
et le calculons dans certains cas
(section \ref{section-calculation}).
Ensuite, nous démontrons quelques lemmes algébriques simples sur
l'existence de factorisations convenables après des modifications
(section \ref{section-preparation-tame-symbol}).
À l'aide de ceux-ci, nous construisons et définissons le symbole modéré à la
section \ref{section-tame-symbol}.
Seules les propriétés les plus élémentaires du symbole modéré
sont nécessaires pour démontrer le lemme clé, ce que nous faisons
à la section \ref{section-key-lemma}.

\medskip\noindent
Ensuite, nous introduisons le cadre de base dans lequel nous travaillons dans la suite de ce
chapitre à la section \ref{section-setup}. Afin de rendre la matière
un peu plus stimulante, nous avons décidé de traiter un cas un peu plus général
que celui habituellement considéré. Plus précisément, nous supposons que nos schémas $X$ sont localement
de type fini sur un schéma de base fixé, localement noethérien, qui est universellement
caténaire et muni d'une fonction de dimension. Ces hypothèses suffisent
pour pouvoir définir les groupes d'homologie de Chow $\CH_*(X)$ et l'action sur
ceux-ci du produit cap avec les classes de Chern. Cela indique que nous devrions
également pouvoir les définir pour les champs algébriques localement de type fini
sur une telle base.

\medskip\noindent
Ensuite, nous suivons les premiers chapitres de \cite{F} afin de définir
les cycles, l'image réciproque plate, l'image directe propre et l'équivalence rationnelle,
si ce n'est que nous avons été moins précis quant aux supports des cycles
considérés.

\medskip\noindent
Nous nous écartons de la présentation donnée dans \cite{F} en utilisant le
lemme clé mentionné ci-dessus pour démontrer une relation de commutativité fondamentale à la
section \ref{section-key}. Nous en déduisons que l'opération
d'intersection avec un faisceau inversible passe au quotient par l'équivalence
rationnelle et est commutative, voir la section \ref{section-commutativity}.
Une nouvelle application du lemme
clé démontre que l'application de Gysin d'un diviseur de Cartier effectif
passe au quotient par l'équivalence rationnelle, voir la section \ref{section-gysin}.
Cela étant démontré, il est immédiat de définir les classes de Chern
des fibrés vectoriels, d'établir l'additivité, de démontrer le principe de scindage,
d'introduire les caractères de Chern et les classes de Todd, puis d'énoncer le
théorème de Grothendieck--Riemann--Roch.

\medskip\noindent
Il y a deux appendices. Dans l'appendice A (section \ref{section-appendix-A}),
nous étudions une autre construction, plus longue, du
symbole modéré et la démonstration correspondante du lemme clé.
Enfin, dans l'appendice B (section \ref{section-appendix-chow}),
nous examinons brièvement le rapport avec la $K$-théorie des faisceaux
cohérents et étudions quelques lemmes sur les éclatements.
Nous suggérons au lecteur de consulter leurs introductions pour
de plus amples informations.

\medskip\noindent
Nous reviendrons, au chapitre suivant, sur les groupes de Chow $\CH_*(X)$ des variétés projectives lisses
sur des corps algébriquement clos. À l'aide d'un lemme de déplacement
comme dans \cite{Samuel}, \cite{ChevalleyI} et \cite{ChevalleyII},
ainsi que de la formule de Serre pour les Tor
(voir \cite{Serre_local_algebra} ou \cite{Serre_algebre_locale}),
nous définirons une structure d'anneau sur $\CH_*(X)$. Voir
Théorie de l'intersection, section \ref{intersection-section-introduction} et suivantes.








\section{Complexes périodiques et quotients de Herbrand}
\label{section-periodic-complexes}

\noindent
Il existe bien entendu une notion très générale de complexe périodique.
On peut imposer la périodicité des applications ou celle des objets.
Nous ajouterons ici ces conditions selon les besoins. Pour le moment, seuls les
cas suivants nous sont nécessaires.

\begin{definition}
\label{definition-periodic-complex}
Soit $R$ un anneau.
\begin{enumerate}
\item Un {\it complexe $2$-périodique} sur $R$ est donné
par un quadruplet $(M, N, \varphi, \psi)$ constitué de
$R$-modules $M$, $N$ et d'applications de $R$-modules $\varphi : M \to N$,
$\psi : N \to M$ telles que
$$
\xymatrix{
\ldots \ar[r] &
M \ar[r]^\varphi &
N \ar[r]^\psi &
M \ar[r]^\varphi &
N \ar[r] & \ldots
}
$$
soit un complexe. Dans ce cadre, nous définissons les {\it modules de cohomologie}
du complexe comme étant les $R$-modules
$$
H^0(M, N, \varphi, \psi) = \Ker(\varphi)/\Im(\psi)
\quad\text{et}\quad
H^1(M, N, \varphi, \psi) = \Ker(\psi)/\Im(\varphi).
$$
Nous disons que le complexe $2$-périodique est {\it exact} si ses modules de cohomologie
sont nuls.
\item Un {\it complexe $(2, 1)$-périodique} sur $R$ est donné
par un triplet $(M, \varphi, \psi)$ constitué d'un $R$-module $M$ et
d'applications de $R$-modules $\varphi : M \to M$, $\psi : M \to M$
telles que
$$
\xymatrix{
\ldots \ar[r] &
M \ar[r]^\varphi &
M \ar[r]^\psi &
M \ar[r]^\varphi &
M \ar[r] & \ldots
}
$$
soit un complexe. Comme il s'agit d'un cas particulier d'un complexe $2$-périodique,
nous disposons de ses {\it modules de cohomologie} $H^0(M, \varphi, \psi)$,
$H^1(M, \varphi, \psi)$ et d'une notion d'exactitude.
\end{enumerate}
\end{definition}

\noindent
Dans la suite, nous utiliserons tout résultat démontré pour les complexes $2$-périodiques
sans autre mention dans le cas des complexes $(2, 1)$-périodiques.
Il est clair que la collection des complexes $2$-périodiques forme une
catégorie, dont les morphismes
$(f, g) : (M, N, \varphi, \psi) \to (M', N', \varphi', \psi')$
sont les couples d'applications $f : M \to M'$ et $g : N \to N'$ tels
que $\varphi' \circ f = g \circ \varphi$ et $\psi' \circ g = f \circ \psi$.
On obtient une catégorie abélienne, avec noyaux et conoyaux comme dans
Homologie, Lemme \ref{homology-lemma-cat-chain-abelian}.

\begin{definition}
\label{definition-periodic-length}
Soit $(M, N, \varphi, \psi)$ un complexe $2$-périodique
sur un anneau $R$ dont les modules de cohomologie sont de longueur finie.
Dans ce cas, nous définissons la {\it multiplicité} de $(M, N, \varphi, \psi)$
comme étant l'entier
$$
e_R(M, N, \varphi, \psi) =
\text{longueur}_R(H^0(M, N, \varphi, \psi))
-
\text{longueur}_R(H^1(M, N, \varphi, \psi))
$$
Dans le cas d'un complexe $(2, 1)$-périodique $(M, \varphi, \psi)$,
nous la notons $e_R(M, \varphi, \psi)$ et l'appellerons parfois
le {\it quotient de Herbrand (additif)}.
\end{definition}

\noindent
Si les groupes de cohomologie de $(M, \varphi, \psi)$
sont des groupes abéliens finis, il est d'usage d'appeler
{\it quotient de Herbrand (multiplicatif)}
$$
q(M, \varphi, \psi) =
\frac{\# H^0(M, \varphi, \psi)}{\# H^1(M, \varphi, \psi)}
$$
Autrement dit, le quotient de Herbrand multiplicatif est le nombre d'éléments de
$H^0$ divisé par le nombre d'éléments de $H^1$. Si $R$ est local et si
le corps résiduel de $R$ est fini à $q$ éléments, alors
$$
q(M, \varphi, \psi) = q^{e_R(M, \varphi, \psi)}
$$

\medskip\noindent
Un exemple de complexe $(2, 1)$-périodique sur un anneau $R$ est tout triplet
de la forme $(M, 0, \psi)$, où $M$ est un $R$-module et $\psi$ une application
$R$-linéaire. Si le noyau et le conoyau de $\psi$ sont de longueur finie,
alors on obtient
\begin{equation}
\label{equation-multiplicity-coker-ker}
e_R(M, 0, \psi) = \text{longueur}_R(\Coker(\psi)) - \text{longueur}_R(\Ker(\psi))
\end{equation}
Énonçons et démontrons les lemmes usuels relatifs à ces notations.

\begin{lemma}
\label{lemma-additivity-periodic-length}
Soit $R$ un anneau. Supposons donnée une suite exacte courte de
complexes $2$-périodiques
$$
0 \to (M_1, N_1, \varphi_1, \psi_1)
\to (M_2, N_2, \varphi_2, \psi_2)
\to (M_3, N_3, \varphi_3, \psi_3)
\to 0
$$
Si deux d'entre eux ont des modules de cohomologie de longueur finie, il en va de même
du troisième, et l'on a
$$
e_R(M_2, N_2, \varphi_2, \psi_2) =
e_R(M_1, N_1, \varphi_1, \psi_1) +
e_R(M_3, N_3, \varphi_3, \psi_3).
$$
\end{lemma}

\begin{proof}
Abrégeons $A = (M_1, N_1, \varphi_1, \psi_1)$,
$B = (M_2, N_2, \varphi_2, \psi_2)$ et $C = (M_3, N_3, \varphi_3, \psi_3)$.
Nous avons une suite exacte longue de cohomologie
$$
\ldots
\to H^1(C)
\to H^0(A)
\to H^0(B)
\to H^0(C)
\to H^1(A)
\to H^1(B)
\to H^1(C)
\to \ldots
$$
Elle donne une suite exacte finie
$$
0 \to I
\to H^0(A)
\to H^0(B)
\to H^0(C)
\to H^1(A)
\to H^1(B)
\to K \to 0
$$
avec $0 \to K \to H^1(C) \to I \to 0$ comme filtration. Par additivité de
la fonction longueur (Algèbre, Lemme \ref{algebra-lemma-length-additive}),
on obtient le résultat.
\end{proof}

\begin{lemma}
\label{lemma-finite-periodic-length}
Soit $R$ un anneau. Si $(M, N, \varphi, \psi)$ est un complexe $2$-périodique
tel que $M$ et $N$ soient de longueur finie, alors
$e_R(M, N, \varphi, \psi) = \text{longueur}_R(M) - \text{longueur}_R(N)$.
En particulier, si $(M, \varphi, \psi)$ est un complexe $(2, 1)$-périodique
tel que $M$ soit de longueur finie, alors
$e_R(M, \varphi, \psi) = 0$.
\end{lemma}

\begin{proof}
Observons que, dans la catégorie des complexes $2$-périodiques où $M$ et $N$
sont de longueur finie, la quantité « $\text{longueur}_R(M) - \text{longueur}_R(N)$ »
est additive pour les suites exactes courtes (l'énoncé précis est laissé au lecteur).
Considérons la suite exacte courte
$$
0 \to (M, \Im(\varphi), \varphi, 0) \to
(M, N, \varphi, \psi) \to (0, N/\Im(\varphi), 0, 0) \to 0
$$
La remarque initiale, combinée à l'additivité du
Lemme \ref{lemma-additivity-periodic-length},
nous ramène aux cas (a) $M = 0$ et (b) $\varphi$ est surjective.
Nous les laissons au lecteur.
\end{proof}

\begin{lemma}
\label{lemma-compare-periodic-lengths}
Soit $R$ un anneau. Soit $f : (M, \varphi, \psi) \to (M', \varphi', \psi')$
une application de complexes $(2, 1)$-périodiques dont les modules de cohomologie
sont de longueur finie. Si $\Ker(f)$ et $\Coker(f)$ sont de longueur finie,
alors $e_R(M, \varphi, \psi) = e_R(M', \varphi', \psi')$.
\end{lemma}

\begin{proof}
Appliquons l'additivité du Lemme \ref{lemma-additivity-periodic-length}
et observons que $(\Ker(f), \varphi, \psi)$ et
$(\Coker(f), \varphi', \psi')$ ont une multiplicité nulle d'après le
Lemme \ref{lemma-finite-periodic-length}.
\end{proof}




\section{Calcul de quelques multiplicités}
\label{section-calculation}

\noindent
Pour démontrer plus loin l'égalité de certains cycles, nous devons
calculer quelques multiplicités. Notre principal outil, outre les
lemmes élémentaires sur les multiplicités donnés à la section précédente,
sera Algèbre, Lemme \ref{algebra-lemma-order-vanishing-determinant}.

\begin{lemma}
\label{lemma-length-multiplication}
Soit $R$ un anneau local noethérien.
Soit $M$ un $R$-module fini. Soit $x \in R$. Supposons que
\begin{enumerate}
\item $\dim(\text{Supp}(M)) \leq 1$, et
\item $\dim(\text{Supp}(M/xM)) \leq 0$.
\end{enumerate}
Écrivons
$\text{Supp}(M) = \{\mathfrak m, \mathfrak q_1, \ldots, \mathfrak q_t\}$.
Alors
$$
e_R(M, 0, x) =
\sum\nolimits_{i = 1, \ldots, t}
\text{ord}_{R/\mathfrak q_i}(x)
\text{longueur}_{R_{\mathfrak q_i}}(M_{\mathfrak q_i}).
$$
\end{lemma}

\begin{proof}
Commençons par quelques remarques préparatoires.
Le résultat du lemme vaut si $M$ est de longueur finie, c'est-à-dire si $t = 0$,
car le membre de gauche et celui de droite sont alors tous deux nuls,
voir le Lemme \ref{lemma-finite-periodic-length}.
De même, si l'on a une suite exacte courte $0 \to M \to M' \to M'' \to 0$
de modules satisfaisant (1) et (2), le lemme pour deux de ces trois
modules implique le lemme pour le troisième, par
additivité des longueurs (Algèbre, Lemme \ref{algebra-lemma-length-additive}) et
additivité des multiplicités (Lemme \ref{lemma-additivity-periodic-length}).

\medskip\noindent
Notons $M_i$ l'image de $M$ dans $M_{\mathfrak q_i}$, de sorte que
$\text{Supp}(M_i) = \{\mathfrak m, \mathfrak q_i\}$.
Le noyau et le conoyau de l'application $M \to \bigoplus M_i$
ont pour support $\{\mathfrak m\}$ et sont donc de longueur finie.
Il résulte de nos remarques préparatoires qu'il suffit de
démontrer le lemme pour chacun des $M_i$. Nous pouvons donc supposer que
$\text{Supp}(M) = \{\mathfrak m, \mathfrak q\}$.
Dans ce cas, nous avons une filtration finie
$M \supset \mathfrak qM \supset \mathfrak q^2M \supset \ldots \supset
\mathfrak q^nM = 0$ d'après Algèbre, Lemme
\ref{algebra-lemma-Noetherian-power-ideal-kills-module}.
Une nouvelle application de l'additivité montre qu'il suffit de démontrer le lemme
lorsque $M$ est annulé par $\mathfrak q$.
Nous pouvons alors considérer $M$ comme un $R/\mathfrak q$-module,
c'est-à-dire supposer que $R$ est un anneau local noethérien intègre
de dimension $1$, de corps des fractions $K$.
En quotientant par le sous-module de torsion, c'est-à-dire par le
noyau de $M \to M \otimes_R K = V$ (la torsion étant de
longueur finie, nos remarques préliminaires s'y appliquent),
nous pouvons supposer que $M \subset V$ est un réseau
(Algèbre, Définition \ref{algebra-definition-lattice}).
Alors $x : M \to M$ est injective et
$\text{longueur}_R(M/xM) = d(M, xM)$
(Algèbre, Définition \ref{algebra-definition-distance}). Comme
$\text{longueur}_K(V) = \dim_K(V)$,
on voit que $\det(x : V \to V) = x^{\dim_K(V)}$ et
$\text{ord}_R(\det(x : V \to V)) = \dim_K(V) \text{ord}_R(x)$.
L'égalité cherchée résulte donc de
Algèbre, Lemme \ref{algebra-lemma-order-vanishing-determinant}
dans ce cas.
\end{proof}

\begin{lemma}
\label{lemma-additivity-divisors-restricted}
Soit $R$ un anneau local noethérien.
Soit $x \in R$. Si $M$ est un module de Cohen--Macaulay fini sur $R$
avec $\dim(\text{Supp}(M)) = 1$ et $\dim(\text{Supp}(M/xM)) = 0$, alors
$$
\text{longueur}_R(M/xM)
=
\sum\nolimits_i \text{longueur}_R(R/(x, \mathfrak q_i))
\text{longueur}_{R_{\mathfrak q_i}}(M_{\mathfrak q_i}).
$$
où $\mathfrak q_1, \ldots, \mathfrak q_t$ sont les idéaux premiers
minimaux du support de $M$. Si $I \subset R$ est un idéal
tel que $x$ soit un non-diviseur de zéro dans $R/I$ et $\dim(R/I) = 1$, alors
$$
\text{longueur}_R(R/(x, I))
=
\sum\nolimits_i \text{longueur}_R(R/(x, \mathfrak q_i))
\text{longueur}_{R_{\mathfrak q_i}}((R/I)_{\mathfrak q_i})
$$
où $\mathfrak q_1, \ldots, \mathfrak q_n$ sont les idéaux premiers minimaux
contenant $I$.
\end{lemma}

\begin{proof}
Ce sont des cas particuliers du Lemme \ref{lemma-length-multiplication}.
\end{proof}

\noindent
Voici un autre cas où l'on peut déterminer la valeur d'une multiplicité.

\begin{lemma}
\label{lemma-powers-period-length-zero}
Soit $R$ un anneau. Soit $M$ un $R$-module.
Soit $\varphi : M \to M$ un endomorphisme et soit $n > 0$
tel que $\varphi^n = 0$ et que $\Ker(\varphi)/\Im(\varphi^{n - 1})$
soit de longueur finie comme $R$-module.
Alors
$$
e_R(M, \varphi^i, \varphi^{n - i}) = 0
$$
pour $i = 0, \ldots, n$.
\end{lemma}

\begin{proof}
Les cas $i = 0, n$ sont triviaux puisque, par convention, $\varphi^0 = \text{id}_M$.
Considérons $M$ comme un $R[t]$-module où la multiplication par $t$
est donnée par $\varphi$. Écrivons
$K_i = \Ker(t^i : M \to M)$ et
$$
a_i = \text{longueur}_R(K_i/t^{n - i}M),\quad
b_i = \text{longueur}_R(K_i/tK_{i + 1}),\quad
c_i = \text{longueur}_R(K_1/t^iK_{i + 1})
$$
Les valeurs au bord sont $a_0 = a_n = b_0 = c_0 = 0$.
Les $c_i$ sont des entiers pour $i < n$, car $K_1/t^iK_{i + 1}$
est un quotient de $K_1/t^{n - 1}M$, qui est supposé de longueur finie.
Nous utiliserons souvent l'égalité $K_i \cap t^jM = t^jK_{i + j}$.
Pour $0 < i < n - 1$, nous avons une suite exacte
$$
0 \to
K_1/t^{n - i - 1}K_{n - i} \to
K_{i + 1}/t^{n - i - 1}M \xrightarrow{t} K_i/t^{n - i}M
\to K_i/tK_{i + 1} \to 0
$$
Par récurrence sur $i$, on en déduit que $a_i$ et $b_i$ sont
des entiers pour $i < n$ et que
$$
c_{n - i - 1} - a_{i + 1} + a_i - b_i = 0
$$
Pour $0 < i < n - 1$, il existe une suite exacte
$$
0 \to
K_i/tK_{i + 1} \to
K_{i + 1}/tK_{i + 2} \xrightarrow{t^i}
K_1/t^{i + 1}K_{i + 2} \to
K_1/t^iK_{i + 1} \to 0
$$
qui donne
$$
b_i - b_{i + 1} + c_{i + 1} - c_i = 0
$$
Comme $b_0 = c_0$, on en conclut que $b_i = c_i$ pour $i < n$.
On voit alors que
$$
a_2 = a_1 + b_{n - 2} - b_1,\quad
a_3 = a_2 + b_{n - 3} - b_2,\quad \ldots
$$
On vérifie immédiatement que cela implique $a_i = a_{n - i}$, comme voulu.
\end{proof}

\begin{lemma}
\label{lemma-multiply-period-length}
Soit $(R, \mathfrak m)$ un anneau local noethérien. Soit
$(M, \varphi, \psi)$ un complexe $(2, 1)$-périodique sur $R$,
avec $M$ fini et dont les groupes de cohomologie sont de longueur finie sur $R$.
Soit $x \in R$ tel que $\dim(\text{Supp}(M/xM)) \leq 0$. Alors
$$
e_R(M, x\varphi, \psi) = e_R(M, \varphi, \psi) - e_R(\Im(\varphi), 0, x)
$$
et
$$
e_R(M, \varphi, x\psi) = e_R(M, \varphi, \psi) + e_R(\Im(\psi), 0, x)
$$
\end{lemma}

\begin{proof}
Nous ne démontrerons que la première formule, la seconde se démontrant
exactement de la même manière.
Soit $M' = M[x^\infty]$ le sous-module de $x$-torsion de $M$.
Considérons la suite exacte courte $0 \to M' \to M \to M'' \to 0$.
Alors $M''$ est sans $x$-torsion (Compléments sur l'algèbre, Lemme
\ref{more-algebra-lemma-divide-by-torsion}).
Comme $\varphi$ et $\psi$ envoient $M'$ dans $M'$,
nous obtenons une suite exacte courte
$$
0 \to (M', \varphi', \psi') \to (M, \varphi, \psi) \to
(M'', \varphi'', \psi'') \to 0
$$
de complexes $(2, 1)$-périodiques. De plus, on obtient une suite exacte courte
$0 \to M' \cap \Im(\varphi) \to \Im(\varphi) \to \Im(\varphi'') \to 0$.
Nous avons
$e_R(M', \varphi, \psi) = e_R(M', x\varphi, \psi) =
e_R(M' \cap \Im(\varphi), 0, x) = 0$
d'après le Lemme \ref{lemma-compare-periodic-lengths}.
Par additivité (Lemme \ref{lemma-additivity-periodic-length}),
on voit qu'il suffit de démontrer le lemme pour $(M'', \varphi'', \psi'')$.
Cela nous ramène au cas étudié dans le paragraphe suivant.

\medskip\noindent
Supposons que $x : M \to M$ soit injective.
Dans ce cas, $\Ker(x\varphi) = \Ker(\varphi)$.
D'autre part, nous avons une suite exacte courte
$$
0 \to \Im(\varphi)/x\Im(\varphi) \to
\Ker(\psi)/\Im(x\varphi) \to \Ker(\psi)/\Im(\varphi) \to 0
$$
Ceci, joint à (\ref{equation-multiplicity-coker-ker}), démontre la formule.
\end{proof}







\section{Préparation aux symboles modérés}
\label{section-preparation-tame-symbol}

\noindent
Dans cette section, nous donnons quelques lemmes qui nous aideront à définir le
symbole modéré à la section suivante.

\begin{lemma}
\label{lemma-glue-at-max}
Soit $A$ un anneau noethérien. Soient $\mathfrak m_1, \ldots, \mathfrak m_r$
des idéaux maximaux deux à deux distincts de $A$. Pour $i = 1, \ldots, r$, soit
$\varphi_i : A_{\mathfrak m_i} \to B_i$ un morphisme d'anneaux dont le
noyau et le conoyau sont annulés par une puissance
de $\mathfrak m_i$. Alors il existe un morphisme d'anneaux $\varphi : A \to B$ tel
que
\begin{enumerate}
\item le localisé de $\varphi$ en $\mathfrak m_i$ soit
isomorphe à $\varphi_i$, et
\item $\Ker(\varphi)$ et $\Coker(\varphi)$ soient annulés
par une puissance de $\mathfrak m_1 \cap \ldots \cap \mathfrak m_r$.
\end{enumerate}
De plus, si chaque $\varphi_i$ est fini, injectif ou
surjectif, il en va de même de $\varphi$.
\end{lemma}

\begin{proof}
Posons $I = \mathfrak m_1 \cap \ldots \cap \mathfrak m_r$. Posons
$A_i = A_{\mathfrak m_i}$ et $A' = \prod A_i$.
Alors $IA' = \prod \mathfrak m_i A_i$ et $A \to A'$
est un morphisme d'anneaux plat tel que $A/I \cong A'/IA'$.
Nous pouvons donc appliquer Compléments sur l'algèbre, Lemme
\ref{more-algebra-lemma-application-formal-glueing}
pour voir qu'il existe une application de $A$-modules $\varphi : A \to B$
tel que $\varphi_i$ soit isomorphe à la localisation de $\varphi$
en $\mathfrak m_i$. On peut ensuite utiliser la discussion de
Compléments sur l'algèbre, Remarque \ref{more-algebra-remark-formal-glueing-algebras}
pour munir $B$ d'une structure de $A$-algèbre
qui induit la structure de $A$-algèbre donnée sur $B_i$.
La dernière assertion du lemme résulte aisément du
fait que $\Ker(\varphi)_{\mathfrak m_i} \cong \Ker(\varphi_i)$
et $\Coker(\varphi)_{\mathfrak m_i} \cong \Coker(\varphi_i)$.
\end{proof}

\noindent
Le lemme suivant est très semblable à
Algèbre, Lemme \ref{algebra-lemma-nonregular-dimension-one}.

\begin{lemma}
\label{lemma-Noetherian-domain-dim-1-two-elements}
Soit $(R, \mathfrak m)$ un anneau local noethérien de dimension $1$.
Soient $a, b \in R$ des non-diviseurs de zéro.
Il existe une extension finie d'anneaux $R \subset R'$
telle que $R'/R$ soit annulé par une puissance de $\mathfrak m$,
et des non-diviseurs de zéro $t, a', b' \in R'$ tels que
$a = ta'$, $b = tb'$ et $R' = a'R' + b'R'$.
\end{lemma}

\begin{proof}
Si $a$ ou $b$ est une unité, le lemme est vrai avec $R = R'$.
Nous pouvons donc supposer $a, b \in \mathfrak m$.
Posons $I = (a, b)$. L'idée consiste à éclater $R$ le long de $I$.
Au lieu d'un raisonnement algébrique, raisonnons géométriquement.
Soit $X = \text{Proj}(\bigoplus_{d \geq 0} I^d)$.
D'après Diviseurs, Lemme
\ref{divisors-lemma-blowing-up-gives-effective-Cartier-divisor},
le morphisme $X \to \Spec(R)$ est un isomorphisme au-dessus du
spectre épointé $U = \Spec(R) \setminus \{\mathfrak m\}$.
Nous pouvons et allons donc considérer $U$ comme un sous-schéma ouvert de $X$.
Le morphisme $X \to \Spec(R)$ est projectif d'après
Diviseurs, Lemme \ref{divisors-lemma-blowing-up-projective}.
De plus, tout point générique de $X$ appartient à $U$, par exemple
d'après Diviseurs, Lemme \ref{divisors-lemma-blow-up-and-irreducible-components}.
Il résulte de Variétés, Lemme \ref{varieties-lemma-finite-in-codim-1}
que $X \to \Spec(R)$ est fini. Ainsi $X = \Spec(R')$ est
affine et $R \to R'$ est fini. Nous avons $R_a \cong R'_a$ puisque $U = D(a)$.
Par conséquent, une puissance de $a$ annule le $R$-module fini $R'/R$.
Comme $\mathfrak m = \sqrt{(a)}$, le module $R'/R$ est annulé
par une puissance de $\mathfrak m$. D'après
Diviseurs, Lemme \ref{divisors-lemma-blowing-up-gives-effective-Cartier-divisor},
on voit que $IR'$ est un idéal localement principal.
Comme $R'$ est semi-local, $IR'$ est principal,
voir Algèbre, Lemme \ref{algebra-lemma-locally-free-semi-local-free},
disons $IR' = (t)$. Alors $a = a't$ et $b = b't$, et tout est
clair.
\end{proof}

\begin{lemma}
\label{lemma-not-infinitely-divisible}
Soit $(R, \mathfrak m)$ un anneau local noethérien de dimension $1$.
Soient $a, b \in R$ des non-diviseurs de zéro avec $a \in \mathfrak m$.
Il existe un entier $n = n(R, a, b)$ tel que, pour toute extension finie
d'anneaux $R \subset R'$, si $b = a^m c$ pour un certain $c \in R'$, alors $m \leq n$.
\end{lemma}

\begin{proof}
Choisissons un idéal premier minimal $\mathfrak q \subset R$. Observons que
$\dim(R/\mathfrak q) = 1$ ; en particulier, $R/\mathfrak q$ n'est pas un corps.
On peut choisir un anneau de valuation discrète $A$ dominant $R/\mathfrak q$
et ayant le même corps des fractions, voir
Algèbre, Lemme \ref{algebra-lemma-dominate-by-dimension-1}. Observons que
$a$ et $b$ s'envoient sur des éléments non nuls de $A$, puisque les non-diviseurs de zéro de $R$
n'appartiennent pas à $\mathfrak q$. Soit $v$ la valuation discrète sur $A$.
Alors $v(a) > 0$ puisque $a \in \mathfrak m$.
Nous affirmons que $n = v(b)/v(a)$ convient.

\medskip\noindent
Soit donnée une extension $R \subset R'$. Posons $A' = A \otimes_R R'$.
Comme $\Spec(R') \to \Spec(R)$ est surjectif
(Algèbre, Lemme \ref{algebra-lemma-integral-overring-surjective}),
$\Spec(A') \to \Spec(A)$ est aussi surjectif
(Algèbre, Lemme \ref{algebra-lemma-surjective-spec-radical-ideal}).
Choisissons un idéal premier $\mathfrak q' \subset A'$ au-dessus de $(0) \subset A$.
Alors $A \subset A'' = A'/\mathfrak q'$ est une extension finie d'anneaux
(induisant de nouveau une surjection sur les spectres).
Choisissons un idéal maximal $\mathfrak m'' \subset A''$
au-dessus de l'idéal maximal de $A$, ainsi qu'un anneau de valuation discrète
$A'''$ dominant $A''_{\mathfrak m''}$ (voir le lemme cité ci-dessus).
Alors $A \to A'''$ est une extension d'anneaux de valuation discrète
et nous avons $b = a^m c$ dans $A'''$. Ainsi $v'''(b) \geq mv'''(a)$.
Comme $v''' = ev$, où $e$ est l'indice de ramification
de $A'''/A$, on obtient $m \leq n$, comme voulu.
\end{proof}

\begin{lemma}
\label{lemma-prepare-tame-symbol}
Soit $(A, \mathfrak m)$ un anneau local noethérien de dimension $1$.
Soit $r \geq 2$ et soient $a_1, \ldots, a_r \in A$ des non-diviseurs de zéro
qui ne sont pas tous des unités.
Alors il existe
\begin{enumerate}
\item une extension finie d'anneaux $A \subset B$ telle que
$B/A$ soit annulé par une puissance de $\mathfrak m$,
\item pour tout idéal maximal $\mathfrak m_j \subset B$,
un non-diviseur de zéro $\pi_j \in B_j = B_{\mathfrak m_j}$, et
\item des factorisations $a_i = u_{i, j} \pi_j^{e_{i, j}}$ dans $B_j$,
où $u_{i, j} \in B_j$ est une unité et $e_{i, j} \geq 0$.
\end{enumerate}
\end{lemma}

\begin{proof}
Puisqu'au moins un $a_i$ n'est pas une unité, $\mathfrak m$
n'est pas un idéal premier associé de $A$. De plus, pour toute extension $A \subset B$
comme dans l'énoncé, $\mathfrak m$ n'est pas un idéal premier associé de $B$
et $\mathfrak m_j$ n'est pas un idéal premier associé de $B_j$.
Garder cela à l'esprit aidera à vérifier les raisonnements ci-dessous.

\medskip\noindent
Tout d'abord, nous affirmons qu'il suffit de démontrer le lemme pour $r = 2$.
Nous allons le montrer par récurrence sur $r$ ; nous suggérons au lecteur
de sauter la démonstration. Supposons donnés $A \subset B$ et des $\pi_j$ dans
$B_j = B_{\mathfrak m_j}$, ainsi que des factorisations
$a_i = u_{i, j} \pi_j^{e_{i, j}}$ pour $i = 1, \ldots, r - 1$ dans $B_j$,
où $u_{i, j} \in B_j$ est une unité et $e_{i, j} \geq 0$.
En appliquant le cas $r = 2$ à $\pi_j$ et $a_r$ dans $B_j$,
on peut trouver des extensions $B_j \subset C_j$ et, pour tout idéal maximal
$\mathfrak m_{j, k}$ de $C_j$, un non-diviseur de zéro
$\pi_{j, k} \in C_{j, k} = (C_j)_{\mathfrak m_{j, k}}$
et des factorisations
$$
\pi_j = v_{j, k} \pi_{j, k}^{f_{j, k}}
\quad\text{et}\quad
a_r = w_{j, k} \pi_{j, k}^{g_{j, k}}
$$
comme dans le lemme. Il existe une unique extension finie $B \subset C$
telle que $C/B$ soit annulé par une puissance de $\mathfrak m$ et
que $C_j \cong C_{\mathfrak m_j}$ pour tout $j$, voir
le Lemme \ref{lemma-glue-at-max}.
Les idéaux maximaux de $C$ correspondent $1$ à $1$
aux idéaux maximaux $\mathfrak m_{j, k}$ dans les localisés,
et dans ces localisés nous avons
$$
a_i = u_{i, j} \pi_j^{e_{i, j}} =
u_{i, j} v_{j, k}^{e_{i, j}} \pi_{j, k}^{e_{i, j}f_{j, k}}
$$
pour $i \leq r - 1$. Comme $a_r$ se factorise lui aussi correctement, la
démonstration de l'étape de récurrence est achevée.

\medskip\noindent
Démonstration du cas $r = 2$. Nous raisonnons par récurrence sur
$$
\ell = \min(\text{longueur}_A(A/a_1A),\ \text{longueur}_A(A/a_2A)).
$$
Si $\ell = 0$, alors $a_1$ ou $a_2$ est une unité et
le lemme vaut avec $A = B$. Nous pouvons donc supposer $\ell > 0$.

\medskip\noindent
Supposons donnée une extension finie d'anneaux $A \subset A'$ telle que
$A'/A$ soit annulé par une puissance de $\mathfrak m$ et que
$\mathfrak m$ ne soit pas un idéal premier associé de $A'$.
Soient $\mathfrak m_1, \ldots, \mathfrak m_r \subset A'$
les idéaux maximaux, et posons $A'_i = A'_{\mathfrak m_i}$.
Si l'on peut résoudre le problème pour $a_1, a_2$ dans chaque $A'_i$,
alors on peut appliquer le Lemme \ref{lemma-glue-at-max}
pour produire une solution pour $a_1, a_2$ dans $A$.
Choisissons $x \in \{a_1, a_2\}$ tel que $\ell = \text{longueur}_A(A/xA)$.
D'après le Lemme \ref{lemma-compare-periodic-lengths} 
et (\ref{equation-multiplicity-coker-ker}),
nous avons $\text{longueur}_A(A/xA) = \text{longueur}_A(A'/xA')$.
D'autre part,
$$
\text{longueur}_A(A'/xA') =
\sum [\kappa(\mathfrak m_i) : \kappa(\mathfrak m)]
\text{longueur}_{A'_i}(A'_i/xA'_i)
$$
d'après Algèbre, Lemme \ref{algebra-lemma-pushdown-module}.
Comme $x \in \mathfrak m$, chaque terme du membre de droite
est positif. On en conclut que l'hypothèse de récurrence s'applique
à $a_1, a_2$ dans chaque $A'_i$ si $r > 1$, ou si $r = 1$ et
$[\kappa(\mathfrak m_1) : \kappa(\mathfrak m)] > 1$.
On en conclut que l'on peut supposer que tout $A'$ comme ci-dessus est local et a
le même corps résiduel que $A$.

\medskip\noindent
En appliquant la discussion du paragraphe précédent,
nous pouvons remplacer $A$ par l'anneau construit dans le
Lemme \ref{lemma-Noetherian-domain-dim-1-two-elements}
pour $a_1, a_2 \in A$. Puisque $A$ est local, nous obtenons alors,
quitte à échanger $a_1$ et $a_2$, que $a_2 \in (a_1)$.
Écrivons $a_2 = a_1^m c$, avec $m > 0$ maximal. En fait, d'après le
Lemme \ref{lemma-not-infinitely-divisible},
nous pouvons supposer que $m$ demeure maximal même après avoir remplacé $A$
par toute extension finie $A \subset A'$ comme au paragraphe précédent.
Si $c$ est une unité, nous avons terminé. Sinon, nous remplaçons
$A$ par l'anneau construit dans le
Lemme \ref{lemma-Noetherian-domain-dim-1-two-elements}
pour $a_1, c \in A$. Alors ou bien (1) $c = a_1 c'$, ou bien
(2) $a_1 = c a'_1$. Le premier cas est impossible, puisqu'il
donnerait $a_2 = a_1^{m + 1} c'$, ce qui contredirait la
maximalité de $m$. Dans le second cas, on obtient
$a_1 = c a'_1$ et $a_2 = c^{m + 1} (a'_1)^m$.
Il suffit alors de démontrer le lemme pour $A$ et $c, a'_1$.
Si $a'_1$ est une unité, nous avons terminé ; sinon,
$\text{longueur}_A(A/cA) < \ell$, car $cA$ est un idéal strictement
plus grand que $a_1A$. Nous concluons donc par l'hypothèse de récurrence.
\end{proof}






\section{Symboles modérés}
\label{section-tame-symbol}

\noindent
Considérons un anneau local noethérien $(A, \mathfrak m)$ de dimension $1$.
Notons $Q(A)$ l'anneau total des fractions de $A$, voir
Algèbre, Exemple \ref{algebra-example-localize-at-prime}.
Le {\it symbole modéré} sera une application
$$
\partial_A(-, -) : Q(A)^* \times Q(A)^* \longrightarrow \kappa(\mathfrak m)^*
$$
qui satisfait les propriétés suivantes :
\begin{enumerate}
\item $\partial_A(f, gh) = \partial_A(f, g) \partial_A(f, h)$
\label{item-bilinear}
pour $f, g, h \in Q(A)^*$,
\item $\partial_A(f, g) \partial_A(g, f) = 1$
\label{item-skew}
pour $f, g \in Q(A)^*$,
\item $\partial_A(f, 1 - f) = 1$
\label{item-1-x}
pour $f \in Q(A)^*$ tel que $1 - f \in Q(A)^*$,
\item $\partial_A(aa', b) = \partial_A(a, b)\partial_A(a', b)$
\label{item-bilinear-better}
et $\partial_A(a, bb') = \partial_A(a, b)\partial_A(a, b')$
pour des non-diviseurs de zéro $a, a', b, b' \in A$,
\item $\partial_A(b, b) = (-1)^m$
\label{item-skew-better}
avec $m = \text{longueur}_A(A/bA)$,
pour tout non-diviseur de zéro $b \in A$,
\item $\partial_A(u, b) = u^m \bmod \mathfrak m$
\label{item-normalization}
avec $m = \text{longueur}_A(A/bA)$, pour toute unité $u \in A$ et tout
non-diviseur de zéro $b \in A$, et
\item
\label{item-1-x-better}
$\partial_A(a, b - a)\partial_A(b, b) = \partial_A(b, b - a)\partial_A(a, b)$
pour $a, b \in A$ tels que $a, b, b - a$ soient des non-diviseurs de zéro.
\end{enumerate}
Comme il est plus facile de travailler avec des éléments de $A$, nous
considérerons souvent $\partial_A$ comme une application définie sur les couples de
non-diviseurs de zéro de $A$ qui satisfait (\ref{item-bilinear-better}),
(\ref{item-skew-better}), (\ref{item-normalization}),
(\ref{item-1-x-better}). C'est un exercice de vérifier qu'en
posant
$$
\partial_A(\frac{a}{b}, \frac{c}{d}) =
\partial_A(a, c) \partial_A(a, d)^{-1} \partial_A(b, c)^{-1} \partial_A(b, d)
$$
on obtient une application bien définie $Q(A)^* \times Q(A)^* \to \kappa(\mathfrak m)^*$
qui satisfait (\ref{item-bilinear}), (\ref{item-skew}), (\ref{item-1-x}),
ainsi que les autres propriétés.

\medskip\noindent
Nous ne prétendons pas qu'il existe une unique application ayant ces propriétés.
Nous allons plutôt donner une recette pour construire une telle application.
Plus précisément, étant donnés des non-diviseurs de zéro $a_1, a_2 \in A$, nous choisissons
une extension d'anneaux $A \subset B$ et des factorisations locales
comme dans le Lemme \ref{lemma-prepare-tame-symbol}.
Puis nous définissons
\begin{equation}
\label{equation-tame-symbol}
\partial_A(a_1, a_2) = \prod\nolimits_j
\text{Norme}_{\kappa(\mathfrak m_j)/\kappa(\mathfrak m)}
((-1)^{e_{1, j}e_{2, j}}u_{1, j}^{e_{2, j}}u_{2, j}^{-e_{1, j}}
\bmod \mathfrak m_j)^{m_j}
\end{equation}
où $m_j = \text{longueur}_{B_j}(B_j/\pi_j B_j)$ et le produit
est pris sur les idéaux maximaux $\mathfrak m_1, \ldots, \mathfrak m_r$ de $B$.

\begin{lemma}
\label{lemma-well-defined-tame-symbol}
La formule (\ref{equation-tame-symbol}) détermine un
élément bien défini de $\kappa(\mathfrak m)^*$. Autrement dit, le
membre de droite ne dépend ni du choix des
factorisations locales, ni du choix de $B$.
\end{lemma}

\begin{proof}
Indépendance du choix des factorisations. Supposons donnés
un anneau local noethérien de dimension $1$, noté $B$, des éléments $a_1, a_2 \in B$,
et des non-diviseurs de zéro $\pi, \theta$ tels que l'on puisse écrire
$$
a_1 = u_1 \pi^{e_1} = v_1 \theta^{f_1},\quad
a_2 = u_2 \pi^{e_2} = v_2 \theta^{f_2}
$$
avec des entiers $e_i, f_i \geq 0$ et des unités $u_i, v_i$ de $B$.
Observons que cela implique
$$
a_1^{e_2} = u_1^{e_2}u_2^{-e_1}a_2^{e_1},\quad
a_1^{f_2} = v_1^{f_2}v_2^{-f_1}a_2^{f_1}
$$
D'autre part, en posant
$m = \text{longueur}_B(B/\pi B)$ et $k = \text{longueur}_B(B/\theta B)$,
on obtient $e_2 m = \text{longueur}_B(B/a_2 B) = f_2 k$.
En développant $a_1^{e_2m} = a_1^{f_2 k}$ à l'aide des égalités ci-dessus, on trouve
$$
(u_1^{e_2}u_2^{-e_1})^m =  (v_1^{f_2}v_2^{-f_1})^k
$$
Cela démontre l'égalité voulue au signe près. Pour vérifier que les signes
conviennent, il faut montrer que $me_1e_2$ est pair si et seulement si
$kf_1f_2$ est pair. Cela résulte des deux égalités $me_2 = kf_2$ et
$me_1 = kf_1$ (par le même raisonnement que ci-dessus).

\medskip\noindent
Indépendance du choix de $B$. Supposons données deux extensions
$A \subset B$ et $A \subset B'$ comme dans le Lemme \ref{lemma-prepare-tame-symbol}.
Alors
$$
C = (B \otimes_A B')/(\mathfrak m\text{-torsion})
$$
en fournira une troisième. Nous pouvons donc supposer que nous avons
$A \subset B \subset C$, ainsi que des factorisations sur les
anneaux locaux de $B$, et il faut montrer que l'emploi
des mêmes factorisations sur les anneaux locaux de $C$
donne le même élément de $\kappa(\mathfrak m)$.
Par transitivité des normes
(Corps, Lemme \ref{fields-lemma-trace-and-norm-tower}),
cela se ramène au problème suivant :
si $B$ est un anneau local noethérien de dimension $1$
et $\pi \in B$ un non-diviseur de zéro, alors
$$
\lambda^m = \prod \text{Norme}_{\kappa_k/\kappa}(\lambda)^{m_k}
$$
Nous avons employé ici les notations suivantes :
(1) $\kappa$ est le corps résiduel de $B$,
(2) $\lambda$ est un élément de $\kappa$,
(3) les $\mathfrak m_k \subset C$ sont les idéaux maximaux de $C$,
(4) $\kappa_k = \kappa(\mathfrak m_k)$ est le corps résiduel de
$C_k = C_{\mathfrak m_k}$,
(5) $m = \text{longueur}_B(B/\pi B)$, et
(6) $m_k = \text{longueur}_{C_k}(C_k/\pi C_k)$.
L'égalité affichée vaut parce que
$\text{Norme}_{\kappa_k/\kappa}(\lambda) = \lambda^{[\kappa_k : \kappa]}$,
puisque $\lambda \in \kappa$, et parce que $m = \sum m_k[\kappa_k:\kappa]$.
Premièrement, nous avons $m = \text{longueur}_B(B/xB) = \text{longueur}_B(C/\pi C)$
d'après le Lemme \ref{lemma-compare-periodic-lengths} 
et (\ref{equation-multiplicity-coker-ker}).
Enfin, $\text{longueur}_B(C/\pi C) = \sum m_k[\kappa_k:\kappa]$
d'après Algèbre, Lemme \ref{algebra-lemma-pushdown-module}.
\end{proof}

\begin{lemma}
\label{lemma-tame-symbol}
Le symbole modéré (\ref{equation-tame-symbol}) satisfait
(\ref{item-bilinear-better}), (\ref{item-skew-better}),
(\ref{item-normalization}), (\ref{item-1-x-better}) et fournit donc
une application $\partial_A : Q(A)^* \times Q(A)^* \to \kappa(\mathfrak m)^*$
qui satisfait (\ref{item-bilinear}), (\ref{item-skew}), (\ref{item-1-x}).
\end{lemma}

\begin{proof}
Démontrons (\ref{item-bilinear-better}).
Soient $a_1, a_2, a_3 \in A$ des non-diviseurs de zéro.
Choisissons $A \subset B$ comme dans le Lemme \ref{lemma-prepare-tame-symbol}
pour $a_1, a_2, a_3$. Alors l'égalité
$$
\partial_A(a_1a_2, a_3) = \partial_A(a_1, a_3) \partial_A(a_2, a_3)
$$
résulte de l'égalité
$$
(-1)^{(e_{1, j} + e_{2, j})e_{3, j}}
(u_{1, j}u_{2, j})^{e_{3, j}}u_{3, j}^{-e_{1, j} - e_{2, j}} =
(-1)^{e_{1, j}e_{3, j}}
u_{1, j}^{e_{3, j}}u_{3, j}^{-e_{1, j}}
(-1)^{e_{2, j}e_{3, j}}
u_{2, j}^{e_{3, j}}u_{3, j}^{-e_{2, j}}
$$
dans $B_j$. Les propriétés (\ref{item-skew-better}) et
(\ref{item-normalization}) sont tout aussi immédiates.

\medskip\noindent
Démontrons (\ref{item-1-x-better}). Soient $a_1, a_2, a_1 - a_2 \in A$
des non-diviseurs de zéro et posons $a_3 = a_1 - a_2$.
Choisissons $A \subset B$ comme dans le Lemme \ref{lemma-prepare-tame-symbol}
pour $a_1, a_2, a_3$. Il suffit alors de montrer que
$$
(-1)^{e_{1, j}e_{2, j} + e_{1, j}e_{3, j} + e_{2, j}e_{3, j} + e_{2, j}}
u_{1, j}^{e_{2, j} - e_{3, j}}
u_{2, j}^{e_{3, j} - e_{1, j}}
u_{3, j}^{e_{1, j} - e_{2, j}} \bmod \mathfrak m_j = 1
$$
C'est clair si $e_{1, j} = e_{2, j} = e_{3, j}$.
Supposons $e_{1, j} > e_{2, j}$. Alors $e_{3, j} = e_{2, j}$,
car $a_3 = a_1 - a_2$, et $u_{3, j}$
a la même classe résiduelle que $-u_{2, j}$. Par conséquent,
la formule est vraie ; les signes conviennent eux aussi,
et cette vérification explique le choix des signes
dans (\ref{equation-tame-symbol}).
Les autres cas se traitent exactement de la même manière.
\end{proof}

\begin{lemma}
\label{lemma-norm-down-tame-symbol}
Soit $(A, \mathfrak m)$ un anneau local noethérien de dimension $1$.
Soit $A \subset B$ une extension finie d'anneaux telle que $B/A$
soit annulé par une puissance de $\mathfrak m$ et que $\mathfrak m$ ne soit pas
un idéal premier associé de $B$.
Pour des non-diviseurs de zéro $a, b \in A$, nous avons
$$
\partial_A(a, b) = \prod
\text{Norme}_{\kappa(\mathfrak m_j)/\kappa(\mathfrak m)}(\partial_{B_j}(a, b))
$$
où le produit est pris sur les idéaux maximaux $\mathfrak m_j$ de $B$
et $B_j = B_{\mathfrak m_j}$.
\end{lemma}

\begin{proof}
Choisissons $B_j \subset C_j$ comme dans le
Lemme \ref{lemma-prepare-tame-symbol} pour $a, b$.
D'après le Lemme \ref{lemma-glue-at-max}, on peut choisir une extension finie
d'anneaux $B \subset C$ telle que $C_j \cong C_{\mathfrak m_j}$ pour tout $j$.
Soient $\mathfrak m_{j, k} \subset C$ les idéaux maximaux de $C$
au-dessus de $\mathfrak m_j$. Soient
$$
a = u_{j, k}\pi_{j, k}^{f_{j, k}},\quad
b = v_{j, k}\pi_{j, k}^{g_{j, k}}
$$
les factorisations locales dont l'existence résulte de notre choix de
$C_j \cong C_{\mathfrak m_j}$. Par définition, nous avons
$$
\partial_A(a, b) = 
\prod\nolimits_{j, k}
\text{Norme}_{\kappa(\mathfrak m_{j, k})/\kappa(\mathfrak m)}
((-1)^{f_{j, k}g_{j, k}}u_{j, k}^{g_{j, k}}v_{j, k}^{-f_{j, k}}
\bmod \mathfrak m_{j, k})^{m_{j, k}}
$$
et
$$
\partial_{B_j}(a, b) = 
\prod\nolimits_k
\text{Norme}_{\kappa(\mathfrak m_{j, k})/\kappa(\mathfrak m_j)}
((-1)^{f_{j, k}g_{j, k}}u_{j, k}^{g_{j, k}}v_{j, k}^{-f_{j, k}}
\bmod \mathfrak m_{j, k})^{m_{j, k}}
$$
Le résultat découle de la transitivité des normes
pour $\kappa(\mathfrak m_{j, k})/\kappa(\mathfrak m_j)/\kappa(\mathfrak m)$, voir
Corps, Lemme \ref{fields-lemma-trace-and-norm-tower}.
\end{proof}

\begin{lemma}
\label{lemma-tame-symbol-formally-smooth}
Soit $(A, \mathfrak m, \kappa) \to (A', \mathfrak m', \kappa')$
un homomorphisme local d'anneaux locaux noethériens. Supposons $A \to A'$
plat et $\dim(A) = \dim(A') = 1$. Posons
$m = \text{longueur}_{A'}(A'/\mathfrak mA')$.
Pour des non-diviseurs de zéro $a_1, a_2 \in A$,
$\partial_A(a_1, a_2)^m$ s'envoie sur $\partial_{A'}(a_1, a_2)$
par $\kappa \to \kappa'$.
\end{lemma}

\begin{proof}
Si $a_1, a_2$ sont tous deux des unités, alors $\partial_A(a_1, a_2) = 1$,
$\partial_{A'}(a_1, a_2) = 1$, et le résultat est vrai.
Sinon, on peut choisir une extension d'anneaux $A \subset B$ et
des factorisations locales comme dans le Lemme \ref{lemma-prepare-tame-symbol}.
Notons $\mathfrak m_1, \ldots, \mathfrak m_m$
les idéaux maximaux de $B$. Soient $\mathfrak m_1, \ldots, \mathfrak m_m$
les idéaux maximaux de $B$, de corps résiduels $\kappa_1, \ldots, \kappa_m$.
Pour chaque $j \in \{1, \ldots, m\}$, notons $\pi_j \in B_j = B_{\mathfrak m_j}$
un non-diviseur de zéro tel que nous ayons des factorisations
$a_i = u_{i, j}\pi_j^{e_{i, j}}$ comme dans le lemme.
Par définition, nous avons
$$
\partial_A(a_1, a_2) = \prod\nolimits_j
\text{Norme}_{\kappa_j/\kappa}
((-1)^{e_{1, j}e_{2, j}}u_{1, j}^{e_{2, j}}u_{2, j}^{-e_{1, j}}
\bmod \mathfrak m_j)^{m_j}
$$
où $m_j = \text{longueur}_{B_j}(B_j/\pi_j B_j)$.

\medskip\noindent
Posons $B' = A' \otimes_A B$. Comme $A'$ est plat sur $A$, on voit
que $A' \subset B'$ est une extension d'anneaux telle que $B'/A'$ soit annulé
par une puissance de $\mathfrak m'$. Soient
$$
\mathfrak m'_{j, l},\quad l = 1, \ldots, n_j
$$
les idéaux maximaux de $B'$ au-dessus de $\mathfrak m_j$. Notons
$\kappa'_{j, l}$ le corps résiduel de $\mathfrak m'_{j, l}$. Notons
$B'_{j, l}$ le localisé de $B'$ en $\mathfrak m'_{j, l}$.
Comme factorisations de $a_1$ et $a_2$ dans $B'_{j, l}$,
nous utilisons l'image des factorisations
$a_i = u_{i, j} \pi_j^{e_{i, j}}$ qui nous sont données dans $B_j$.
Par définition, nous avons
$$
\partial_{A'}(a_1, a_2) = \prod\nolimits_{j, l}
\text{Norme}_{\kappa'_{j, l}/\kappa'}
((-1)^{e_{1, j}e_{2, j}}u_{1, j}^{e_{2, j}}u_{2, j}^{-e_{1, j}}
\bmod \mathfrak m'_{j, l})^{m'_{j, l}}
$$
où $m'_{j, l} = \text{longueur}_{B'_{j, l}}(B'_{j, l}/\pi_j B'_{j, l})$.

\medskip\noindent
En comparant les formules ci-dessus, on voit qu'il suffit de montrer que
pour tout $j$ et toute unité $u \in B_j$, on a
\begin{equation}
\label{equation-to-prove}
\left(\text{Norme}_{\kappa_j/\kappa}(u \bmod \mathfrak m_j)^{m_j}\right)^m
=
\prod\nolimits_l
\text{Norme}_{\kappa'_{j, l}/\kappa'}(u \bmod \mathfrak m'_{j, l})^{m'_{j, l}}
\end{equation}
dans $\kappa'$. Pour le démontrer, nous allons utiliser la construction des déterminants
d'endomorphismes de modules de longueur finie donnée dans
Compléments sur l'algèbre, section \ref{more-algebra-section-determinants-finite-length}.
Posons $M = B_j/\pi_j B_j$. D'après
Compléments sur l'algèbre, Lemme \ref{more-algebra-lemma-multiplication}, nous avons
$$
\text{Norme}_{\kappa_j/\kappa}(u \bmod \mathfrak m_j)^{m_j} =
\det\nolimits_\kappa(u : M \to M)
$$
Ainsi, d'après
Compléments sur l'algèbre, Lemme \ref{more-algebra-lemma-flat-base-change-det},
le membre de gauche de (\ref{equation-to-prove}) est égal à
$\det_{\kappa'}(u : M \otimes_A A' \to M \otimes_A A')$.
Nous avons un isomorphisme
$$
M \otimes_A A' = (B_j/\pi_j B_j) \otimes_A A' =
\bigoplus\nolimits_l B'_{j, l}/\pi_j B'_{j, l}
$$
de $A'$-modules. En posant $M'_l = B'_{j, l}/\pi_j B'_{j, l}$, on voit que
$\text{Norme}_{\kappa'_{j, l}/\kappa'}(u \bmod \mathfrak m'_{j, l})^{m'_{j, l}}
= \det_{\kappa'}(u_j : M'_l \to M'_l)$ d'après, de nouveau,
Compléments sur l'algèbre, Lemme \ref{more-algebra-lemma-multiplication}.
Ainsi (\ref{equation-to-prove}) résulte de la multiplicativité de la construction
du déterminant, voir Compléments sur l'algèbre, Lemme \ref{more-algebra-lemma-ses}.
\end{proof}






\section{Un lemme clé}
\label{section-key-lemma}

\noindent
Dans cette section, nous appliquons les résultats précédents pour démontrer le
Lemme \ref{lemma-milnor-gersten-low-degree}.
Ce lemme est le cas en bas degré de l'assertion selon laquelle
il existe, pour la K-théorie de Milnor, un complexe analogue
au complexe de Gersten--Quillen en K-théorie de Quillen.
Voir la Remarque \ref{remark-gersten-complex-milnor}.

\begin{lemma}
\label{lemma-perpare-key}
Soit $(A, \mathfrak m)$ un anneau local noethérien de dimension $2$.
Soit $t \in \mathfrak m$ un non-diviseur de zéro. Écrivons
$V(t) = \{\mathfrak m, \mathfrak q_1, \ldots, \mathfrak q_r\}$.
Soit $A_{\mathfrak q_i} \subset B_i$ une extension finie d'anneaux
telle que $B_i/A_{\mathfrak q_i}$ soit annulé par une puissance de
$t$. Alors il existe une extension finie $A \subset B$
d'anneaux locaux, identifiant les corps résiduels,
telle que $B_i \cong B_{\mathfrak q_i}$ et que $B/A$ soit annulé
par une puissance de $t$.
\end{lemma}

\begin{proof}
Choisissons $n > 0$ tel que $B_i \subset t^{-n}A_{\mathfrak q_i}$.
Soient $M \subset t^{-n}A$, respectivement $M' \subset t^{-2n}A$, les
sous-$A$-modules formés des éléments dont les images appartiennent à $B_i$
dans $t^{-n}A_{\mathfrak q_i}$ et $t^{-2n}A_{\mathfrak q_i}$, respectivement.
L'inclusion $M \subset M'$ est alors une inclusion de $A$-modules finis, puisque $A$ est noethérien,
et $M_{\mathfrak q_i} = M'_{\mathfrak q_i} = B_i$, car la localisation
est exacte. Ainsi $M'/M$ est annulé par $\mathfrak m^c$ pour un certain
$c > 0$. Observons que $M \cdot M \subset M'$ pour la multiplication
$t^{-n}A \times t^{-n}A \to t^{-2n}A$. Par conséquent,
$B = A + \mathfrak m^{c + 1}M$ est une $A$-algèbre finie ayant les bons
localisés. Nous omettons de vérifier que $B$ est local, d'idéal maximal
$\mathfrak m + \mathfrak m^{c + 1}M$.
\end{proof}

\begin{lemma}
\label{lemma-key-nonzerodivisors}
Soit $(A, \mathfrak m)$ un anneau local noethérien de dimension $2$.
Soient $a, b \in A$ des non-diviseurs de zéro.
Alors
$$
\sum
\text{ord}_{A/\mathfrak q}(\partial_{A_{\mathfrak q}}(a, b))
=
0
$$
où la somme porte sur les idéaux premiers de hauteur $1$, notés $\mathfrak q$, de $A$.
\end{lemma}

\begin{proof}
Si $\mathfrak q$ est un idéal premier de hauteur $1$ de $A$ tel que $a, b$
s'envoient sur une unité de $A_\mathfrak q$, alors $\partial_{A_\mathfrak q}(a, b) = 1$.
La somme est donc finie. En fait, si
$V(ab) = \{\mathfrak m, \mathfrak q_1, \ldots, \mathfrak q_r\}$,
la somme porte sur $i = 1, \ldots, r$.
Pour chaque $i$, choisissons une extension $A_{\mathfrak q_i} \subset B_i$
comme dans le Lemme \ref{lemma-prepare-tame-symbol} pour $a, b$.
D'après le Lemme \ref{lemma-perpare-key}, appliqué avec $t = ab$ et la liste d'idéaux premiers donnée,
nous pouvons supposer donnée une extension locale finie $A \subset B$
telle que $B/A$ soit annulé par une puissance de $ab$ et que,
pour chaque $i$, on ait $B_{\mathfrak q_i} \cong B_i$.
Observons que, si les $\mathfrak q_{i, j}$ sont les idéaux premiers de $B$
au-dessus de $\mathfrak q_i$, alors
$$
\text{ord}_{A/\mathfrak q_i}(\partial_{A_{\mathfrak q_i}}(a, b))
=
\sum\nolimits_j
\text{ord}_{B/\mathfrak q_{i, j}}(\partial_{B_{\mathfrak q_{i, j}}}(a, b))
$$
d'après le Lemme \ref{lemma-norm-down-tame-symbol} et
Algèbre, Lemme \ref{algebra-lemma-finite-extension-dim-1}.
Nous pouvons donc remplacer $A$ par $B$ et
nous ramener au cas traité dans le paragraphe suivant.

\medskip\noindent
Supposons que, pour chaque $i$, il existe un non-diviseur de zéro
$\pi_i \in A_{\mathfrak q_i}$ et des unités $u_i, v_i \in A_{\mathfrak q_i}$
tels que, pour certains entiers $e_i, f_i \geq 0$, nous ayons
$$
a = u_i \pi_i^{e_i},\quad b = v_i \pi_i^{f_i}
$$
dans $A_{\mathfrak q_i}$. En posant
$m_i = \text{longueur}_{A_{\mathfrak q_i}}(A_{\mathfrak q_i}/\pi_i)$,
nous avons
$\partial_{A_{\mathfrak q_i}}(a, b) =
((-1)^{e_if_i}u_i^{f_i}v_i^{-e_i})^{m_i}$ par définition.
Comme $a, b$ sont des non-diviseurs de zéro, le
complexe $(2, 1)$-périodique $(A/(ab), a, b)$ a une cohomologie nulle.
Notons $M_i$ l'image de $A/(ab)$ dans $A_{\mathfrak q_i}/(ab)$.
Nous avons alors une application
$$
A/(ab) \longrightarrow \bigoplus M_i
$$
dont le noyau et le conoyau ont leur support dans $\{\mathfrak m\}$
et sont donc de longueur finie. On voit ainsi que
$$
\sum e_A(M_i, a, b) = 0
$$
d'après le Lemme \ref{lemma-compare-periodic-lengths}. Il suffit donc de montrer
$e_A(M_i, a, b) =
- \text{ord}_{A/\mathfrak q_i}(\partial_{A_{\mathfrak q_i}}(a, b))$.

\medskip\noindent
Démontrons d'abord ceci dans le cas où $\pi_i, u_i, v_i$ sont les images
d'éléments $\pi_i, u_i, v_i \in A$ (l'emploi des mêmes symboles
ne devrait prêter à aucune confusion). Dans ce cas, on obtient
\begin{align*}
e_A(M_i, a, b) & =
e_A(M_i, u_i\pi_i^{e_i}, v_i\pi_i^{f_i}) \\
& =
e_A(M_i, \pi_i^{e_i}, \pi_i^{f_i}) -
e_A(\pi_i^{e_i}M_i, 0, u_i) +
e_A(\pi_i^{f_i}M_i, 0, v_i) \\
& =
0 -
f_im_i\text{ord}_{A/\mathfrak q_i}(u_i) +
e_im_i\text{ord}_{A/\mathfrak q_i}(v_i) \\
& =
-m_i\text{ord}_{A/\mathfrak q_i}(u_i^{f_i}v_i^{-e_i}) =
-\text{ord}_{A/\mathfrak q_i}(\partial_{A_{\mathfrak q_i}}(a, b))
\end{align*}
La deuxième égalité résulte du Lemme \ref{lemma-multiply-period-length}.
Observons que
$M_i \subset (M_i)_{\mathfrak q_i} = A_{\mathfrak q_i}/(\pi_i^{e_i + f_i})$
et que
$(\pi_i^{e_i}M_i)_{\mathfrak q_i} \cong A_{\mathfrak q_i}/\pi_i^{f_i}$ et
$(\pi_i^{f_i}M_i)_{\mathfrak q_i} \cong A_{\mathfrak q_i}/\pi_i^{e_i}$.
Le $0$ de la troisième égalité provient du
Lemme \ref{lemma-powers-period-length-zero},
et les deux autres termes proviennent du
Lemme \ref{lemma-length-multiplication}.
Les deux dernières égalités résultent de la multiplicativité de
la fonction ordre et de la définition de notre symbole modéré.

\medskip\noindent
En général, on peut d'abord choisir $c \in A$, $c \not \in \mathfrak q_i$
tel que $c\pi_i \in A$. En remplaçant $\pi_i$ par $c\pi_i$,
$u_i$ par $c^{-e_i}u_i$ et $v_i$ par $c^{-f_i}v_i$,
nous pouvons et allons supposer que $\pi_i$ appartient à $A$.
Ensuite, choisissons un $c \in A$, $c \not \in \mathfrak q_i$,
tel que $cu_i, cv_i \in A$. On observe alors que
$$
e_A(M_i, ca, cb) = e_A(M_i, a, b) - e_A(aM_i, 0, c) + e_A(bM_i, 0, c)
$$
d'après le Lemme \ref{lemma-length-multiplication}.
D'autre part, nous avons
$$
\partial_{A_{\mathfrak q_i}}(ca, cb) =
c^{m_i(f_i - e_i)}\partial_{A_{\mathfrak q_i}}(a, b)
$$
dans $\kappa(\mathfrak q_i)^*$, puisque $c$ est une unité de $A_{\mathfrak q_i}$.
Les raisonnements du paragraphe précédent montrent que
$e_A(M_i, ca, cb) = -
\text{ord}_{A/\mathfrak q_i}(\partial_{A_{\mathfrak q_i}}(ca, cb))$.
Il suffit donc de démontrer
$$
e_A(aM_i, 0, c) = \text{ord}_{A/\mathfrak q_i}(c^{m_if_i})
\quad\text{et}\quad
e_A(bM_i, 0, c) = \text{ord}_{A/\mathfrak q_i}(c^{m_ie_i})
$$
ce qui résulte du Lemme \ref{lemma-length-multiplication}
par la description donnée ci-dessus de ce qui se produit
après localisation en $\mathfrak q_i$.
\end{proof}

\begin{lemma}[Lemme clé]
\label{lemma-milnor-gersten-low-degree}
\begin{reference}
Lorsque $A$ est un anneau excellent, il s'agit de \cite[Proposition 1]{Kato-Milnor-K}.
\end{reference}
Soit $A$ un anneau local noethérien intègre de dimension $2$, de corps des fractions $K$.
Soient $f, g \in K^*$.
Soient $\mathfrak q_1, \ldots, \mathfrak q_t$ les idéaux premiers de hauteur
$1$, notés $\mathfrak q$, de $A$ tels que $f$ ou $g$ n'appartienne pas à
$A^*_{\mathfrak q}$.
Alors
$$
\sum\nolimits_{i = 1, \ldots, t}
\text{ord}_{A/\mathfrak q_i}(\partial_{A_{\mathfrak q_i}}(f, g))
=
0
$$
On peut aussi l'écrire
$$
\sum\nolimits_{\text{hauteur}(\mathfrak q) = 1}
\text{ord}_{A/\mathfrak q}(\partial_{A_{\mathfrak q}}(f, g))
=
0
$$
car, en tout idéal premier de hauteur $1$, noté $\mathfrak q$,
de $A$ tel que $f, g \in A^*_{\mathfrak q}$,
nous avons $\partial_{A_{\mathfrak q}}(f, g) = 1$.
\end{lemma}

\begin{proof}
Comme les symboles modérés $\partial_{A_{\mathfrak q}}(f, g)$ sont
bilinéaires et les fonctions ordre $\text{ord}_{A/\mathfrak q}$
additives, il suffit de démontrer la formule lorsque
$f$ et $g$ sont des éléments de $A$. Ce cas est démontré au
Lemme \ref{lemma-key-nonzerodivisors}.
\end{proof}

\begin{remark}[K-théorie de Milnor]
\label{remark-gersten-complex-milnor}
Pour un corps $k$, notons $K^M_*(k)$ le quotient de
l'algèbre tensorielle sur $k^*$ par l'idéal bilatère
engendré par les éléments $x \otimes 1 - x$ pour $x \in k \setminus \{0, 1\}$.
Ainsi $K^M_0(k) = \mathbf{Z}$, $K_1^M(k) = k^*$, et
$$
K^M_2(k) = k^* \otimes_\mathbf{Z} k^* / \langle x \otimes 1 - x \rangle
$$
Si $A$ est un anneau de valuation discrète de corps des fractions $F = \text{Frac}(A)$
et de corps résiduel $\kappa$, il existe un symbole modéré
$$
\partial_A : K_{i + 1}^M(F) \to K_i^M(\kappa)
$$
défini comme à la section \ref{section-tame-symbol} ; voir \cite{Kato-Milnor-K}.
Plus généralement, cette application peut être étendue au cas où $A$ est un
anneau local intègre excellent de dimension $1$, à l'aide de la normalisation et des applications
normes sur $K_i^M$, voir \cite{Kato-Milnor-K} ; vraisemblablement, la méthode de la
section \ref{section-tame-symbol} permet d'étendre la construction
du symbole modéré $\partial_A$ à tout anneau local noethérien intègre $A$
de dimension $1$. Ensuite, soit $X$ un schéma noethérien muni d'une
fonction de dimension $\delta$. On peut alors utiliser ces symboles modérés pour obtenir
les flèches suivantes :
$$
\bigoplus\nolimits_{\delta(x) = j + 1} K^M_{i + 1}(\kappa(x))
\longrightarrow
\bigoplus\nolimits_{\delta(x) = j} K^M_i(\kappa(x))
\longrightarrow
\bigoplus\nolimits_{\delta(x) = j - 1} K^M_{i - 1}(\kappa(x))
$$
Cependant, il n'est pas clair que la composée soit nulle, c'est-à-dire que
l'on obtienne un complexe de groupes abéliens.
Pour $X$ excellent, cela est démontré dans \cite{Kato-Milnor-K}.
Lorsque $i = 1$ et $j$ est quelconque, cela résulte du
Lemme \ref{lemma-milnor-gersten-low-degree}.
\end{remark}





















\section{Cadre}
\label{section-setup}

\noindent
Nous travaillerons partout sur une base localement noethérienne et universellement
caténaire $S$, munie d'une fonction de dimension $\delta$.
Bien qu'il soit probablement possible de généraliser certaines parties de la
discussion de ce chapitre, cela semble constituer une bonne première
approximation. C'est exactement le degré de généralité étudié dans \cite{Thorup}.
Nous ne supposons généralement pas nos schémas
séparés ni quasi-compacts. De nombreux champs algébriques intéressants
ne sont ni séparés ni quasi-compacts, et c'est un bon
cas d'étude pour voir comment élaborer aussi pour eux une théorie raisonnable.
Afin de pouvoir faire référence à ces hypothèses, nous les numérotons.

\begin{situation}
\label{situation-setup}
Ici, $S$ est un schéma localement noethérien et universellement caténaire.
De plus, nous supposons $S$ muni d'une fonction de dimension
$\delta : S \longrightarrow \mathbf{Z}$.
\end{situation}

\noindent
Voir Morphismes, Définition \ref{morphisms-definition-universally-catenary},
pour la notion de schéma universellement caténaire, et
Topologie, Définition \ref{topology-definition-dimension-function},
pour la notion de fonction de dimension. Rappelons que tout schéma localement
noethérien et caténaire possède localement une fonction de dimension, voir
Propriétés, Lemme \ref{properties-lemma-catenary-dimension-function}.
De plus, beaucoup de schémas sont universellement caténaires,
voir Morphismes, Lemme \ref{morphisms-lemma-ubiquity-uc}.

\medskip\noindent
Soit $(S, \delta)$ comme dans la Situation \ref{situation-setup}.
Tout schéma $X$ localement de type fini sur $S$ est localement noethérien
et caténaire. En fait, $X$ possède une fonction de dimension canonique
$$
\delta = \delta_{X/S} : X \longrightarrow \mathbf{Z}
$$
associée à $(f : X \to S, \delta)$ et donnée par la règle
$\delta_{X/S}(x) = \delta(f(x)) + \text{trdeg}_{\kappa(f(x))}\kappa(x)$.
Voir Morphismes, Lemme \ref{morphisms-lemma-dimension-function-propagates}.
De plus, si $h : X \to Y$ est un morphisme de schémas localement de type
fini sur $S$, et si $x \in X$, $y = h(x)$,
alors, de manière évidente,
$\delta_{X/S}(x) = \delta_{Y/S}(y) + \text{trdeg}_{\kappa(y)}\kappa(x)$.
Nous utiliserons librement cette fonction et ses propriétés dans la suite.

\medskip\noindent
Voici les exemples fondamentaux de cadres comme ci-dessus.
En fait, le principal intérêt réside dans le cas où la base
est le spectre d'un corps, ou dans celui où la base
est le spectre d'un anneau de Dedekind (par exemple $\mathbf{Z}$
ou un anneau de valuation discrète).

\begin{example}
\label{example-field}
Ici, $S = \Spec(k)$ et $k$ est un corps.
On pose $\delta(pt) = 0$, où $pt$ désigne l'unique point de $S$.
Le couple $(S, \delta)$ est un exemple de situation comme dans la
Situation \ref{situation-setup}, d'après
Morphismes, Lemme \ref{morphisms-lemma-ubiquity-uc}.
\end{example}

\begin{example}
\label{example-domain-dimension-1}
Ici, $S = \Spec(A)$, où $A$ est un anneau noethérien intègre
de dimension $1$.
On peut par exemple considérer $A = \mathbf{Z}$.
On pose $\delta(\mathfrak p) = 0$ si
$\mathfrak p$ est un idéal maximal, et $\delta(\mathfrak p) = 1$
si $\mathfrak p = (0)$ correspond au point générique.
C'est un exemple de Situation \ref{situation-setup}, d'après
Morphismes, Lemme \ref{morphisms-lemma-ubiquity-uc}.
\end{example}

\begin{example}
\label{example-CM-irreducible}
Ici, $S$ est un schéma de Cohen--Macaulay. Alors $S$ est universellement caténaire d'après
Morphismes, Lemme \ref{morphisms-lemma-ubiquity-uc}.
On pose $\delta(s) = -\dim(\mathcal{O}_{S, s})$.
Si $s' \leadsto s$ est une spécialisation non triviale de points de $S$,
alors $\mathcal{O}_{S, s'}$ est le localisé de $\mathcal{O}_{S, s}$
en un idéal premier non maximal $\mathfrak p \subset \mathcal{O}_{S, s}$, voir
Schémas, Lemme \ref{schemes-lemma-specialize-points}.
Ainsi $\dim(\mathcal{O}_{S, s}) =
\dim(\mathcal{O}_{S, s'}) + \dim(\mathcal{O}_{S, s}/\mathfrak p) >
\dim(\mathcal{O}_{S, s'})$ d'après
Algèbre, Lemme \ref{algebra-lemma-CM-dim-formula}.
Par conséquent, $\delta(s') > \delta(s)$. Si $s' \leadsto s$ est
une spécialisation immédiate, alors il n'existe aucun idéal
premier strictement compris entre $\mathfrak p$ et $\mathfrak m_s$,
et l'on obtient $\delta(s') = \delta(s) + 1$. Ainsi $\delta$
est une fonction de dimension. Autrement dit, le couple $(S, \delta)$
est un exemple de Situation \ref{situation-setup}.
\end{example}

\noindent
Si $S$ est jacobsonien et si $\delta$ envoie les points fermés sur zéro, alors $\delta$
est la fonction qui associe à un point la dimension de son adhérence.

\begin{lemma}
\label{lemma-delta-is-dimension}
Soit $(S, \delta)$ comme dans la Situation \ref{situation-setup}.
Supposons en outre que $S$ soit un schéma jacobsonien et que $\delta(s) = 0$ pour tout
point fermé $s$ de $S$. Soit $X$ localement de type fini sur $S$.
Soit $Z \subset X$ un sous-schéma fermé intègre et soit
$\xi \in Z$ son point générique. Les entiers suivants sont égaux :
\begin{enumerate}
\item $\delta_{X/S}(\xi)$,
\item $\dim(Z)$, et
\item $\dim(\mathcal{O}_{Z, z})$, où $z$ est un point fermé de $Z$.
\end{enumerate}
\end{lemma}

\begin{proof}
Soient $X \to S$ et $\xi \in Z \subset X$ comme dans le lemme.
Comme $X$ est localement de type fini sur $S$, on voit que
$X$ est jacobsonien, voir
Morphismes, Lemme \ref{morphisms-lemma-Jacobson-universally-Jacobson}.
Par conséquent, les points fermés de $X$ sont denses dans tout fermé de $Z$
et s'envoient sur des points fermés de $S$. Ainsi, toute chaîne
de fermés irréductibles de $Z$ peut être terminée par un point fermé de $Z$.
Il s'ensuit que $\dim(Z) = \sup_z(\dim(\mathcal{O}_{Z, z})$
(voir Propriétés, Lemme \ref{properties-lemma-codimension-local-ring}),
où $z \in Z$ parcourt les points fermés de $Z$.
Notons que $\dim(\mathcal{O}_{Z, z}) = \delta(\xi) - \delta(z)$
d'après les propriétés d'une fonction de dimension.
Pour tout point fermé $z \in Z$, l'extension de corps
$\kappa(z)/\kappa(f(z))$ est finie, voir Morphismes,
Lemme \ref{morphisms-lemma-jacobson-finite-type-points}.
Ainsi $\delta_{X/S}(z) = \delta(f(z)) = 0$ pour $z \in Z$ fermé.
Il s'ensuit que les trois entiers sont égaux.
\end{proof}

\noindent
Dans la situation du lemme précédent, la
valeur de $\delta$ au point générique d'un fermé irréductible
est la dimension de ce fermé irréductible.
Cependant, on ne peut s'attendre à ce que cette égalité vaille en général.
Par exemple, si $S = \Spec(\mathbf{C}[[t]])$ et
$X = \Spec(\mathbf{C}((t)))$, on obtiendrait alors
$\delta(x) = 1$ pour l'unique point de $X$, mais $\dim(X) = 0$.
Nous voulons néanmoins considérer $\delta_{X/S}$ comme donnant la
dimension des sous-schémas fermés irréductibles. Nous introduisons donc
la terminologie suivante.

\begin{definition}
\label{definition-delta-dimension}
Soit $(S, \delta)$ comme dans la Situation \ref{situation-setup}.
Pour tout schéma $X$ localement de type fini sur $S$
et tout fermé irréductible $Z \subset X$, nous définissons
$$
\dim_\delta(Z) = \delta(\xi)
$$
où $\xi \in Z$ est le point générique de $Z$.
Nous appellerons cet entier la {\it $\delta$-dimension de $Z$}.
Si $Z$ est un sous-schéma fermé de $X$, nous définissons
$\dim_\delta(Z)$ comme la borne supérieure des $\delta$-dimensions
de ses composantes irréductibles.
\end{definition}







\section{Cycles}
\label{section-cycles}

\noindent
Comme nous ne supposons pas nos schémas quasi-compacts, il faut
être un peu prudent dans la définition des cycles. Nous devons autoriser
les sommes infinies, car une fonction rationnelle peut, par exemple, avoir une infinité de
pôles. Quoi qu'il en soit, si $X$ est quasi-compact, un
cycle est une somme finie, comme d'habitude.

\begin{definition}
\label{definition-cycles}
Soit $(S, \delta)$ comme dans la Situation \ref{situation-setup}.
Soit $X$ localement de type fini sur $S$.
Soit $k \in \mathbf{Z}$.
\begin{enumerate}
\item Un {\it cycle sur $X$} est une somme formelle
$$
\alpha = \sum n_Z [Z]
$$
où la somme porte sur les sous-schémas fermés intègres $Z \subset X$,
où chaque $n_Z \in \mathbf{Z}$, et où la collection
$\{Z; n_Z \not = 0\}$ est localement finie
(Topologie, Définition \ref{topology-definition-locally-finite}).
\item Un {\it $k$-cycle} sur $X$ est un cycle
$$
\alpha = \sum n_Z [Z]
$$
tel que $n_Z \not = 0 \Rightarrow \dim_\delta(Z) = k$.
\item Le groupe abélien de tous les $k$-cycles sur $X$ est noté $Z_k(X)$.
\end{enumerate}
\end{definition}

\noindent
Autrement dit, un $k$-cycle sur $X$
est une combinaison $\mathbf{Z}$-linéaire formelle et localement finie
de sous-schémas fermés intègres de $\delta$-dimension $k$.
L'addition des $k$-cycles $\alpha = \sum n_Z[Z]$ et
$\beta = \sum m_Z[Z]$ est donnée par
$$
\alpha + \beta = \sum (n_Z + m_Z)[Z],
$$
c'est-à-dire par addition des coefficients.

\begin{remark}
\label{remark-cycles-pointwise}
Soit $(S, \delta)$ comme dans la Situation \ref{situation-setup}.
Soit $X$ localement de type fini sur $S$. Soit $k \in \mathbf{Z}$.
Alors on peut écrire
$$
Z_k(X) = \bigoplus\nolimits_{\delta(x) = k}' K_0^M(\kappa(x))
\quad\subset\quad
\bigoplus\nolimits_{\delta(x) = k} K_0^M(\kappa(x))
$$
avec les notations et conventions suivantes :
\begin{enumerate}
\item $K_0^M(\kappa(x)) = \mathbf{Z}$ est la partie de degré $0$ de
la K-théorie de Milnor du corps résiduel $\kappa(x)$ du point
$x \in X$ (voir la Remarque \ref{remark-gersten-complex-milnor}), et
\item la somme directe de droite porte sur tous les points $x \in X$
tels que $\delta(x) = k$,
\item la notation $\bigoplus'_x$ signifie que nous considérons le
sous-groupe constitué des éléments localement finis, à savoir des éléments
$\sum_x n_x$ tels que, pour tout ouvert quasi-compact $U \subset X$,
l'ensemble des $x \in U$ tels que $n_x \not = 0$ soit fini.
\end{enumerate}
\end{remark}

\begin{definition}
\label{definition-support-cycle}
Soit $(S, \delta)$ comme dans la Situation \ref{situation-setup}.
Soit $X$ localement de type fini sur $S$.
Le {\it support} d'un cycle $\alpha = \sum n_Z [Z]$ sur $X$
est
$$
\text{Supp}(\alpha) = \bigcup\nolimits_{n_Z \not = 0} Z \subset X
$$
\end{definition}

\noindent
Comme la collection $\{Z; n_Z \not = 0\}$ est localement finie,
on voit que $\text{Supp}(\alpha)$ est un fermé de $X$.
Si $\alpha$ est un $k$-cycle, toute composante irréductible $Z$
de $\text{Supp}(\alpha)$ est de $\delta$-dimension $k$.

\begin{definition}
\label{definition-effective-cycle}
Soit $(S, \delta)$ comme dans la Situation \ref{situation-setup}.
Soit $X$ localement de type fini sur $S$.
Un cycle $\alpha$ sur $X$ est {\it effectif} s'il
peut s'écrire $\alpha =\sum n_Z [Z]$ avec $n_Z \geq 0$ pour tout $Z$.
\end{definition}

\noindent
L'ensemble de tous les cycles effectifs est un monoïde, car la somme de
deux cycles effectifs est effective, mais ce n'est pas un groupe
(sauf si $X = \emptyset$).




\section{Cycle associé à un sous-schéma fermé}
\label{section-cycle-of-closed-subscheme}

\begin{lemma}
\label{lemma-multiplicity-finite}
Soit $(S, \delta)$ comme dans la Situation \ref{situation-setup}.
Soit $X$ localement de type fini sur $S$.
Soit $Z \subset X$ un sous-schéma fermé.
\begin{enumerate}
\item Soit $Z' \subset Z$ une composante irréductible et
soit $\xi \in Z'$ son point générique.
Alors
$$
\text{longueur}_{\mathcal{O}_{X, \xi}} \mathcal{O}_{Z, \xi} < \infty
$$
\item Si $\dim_\delta(Z) \leq k$ et si $\xi \in Z$ vérifie
$\delta(\xi) = k$, alors $\xi$ est le point générique d'une
composante irréductible de $Z$.
\end{enumerate}
\end{lemma}

\begin{proof}
Soient $Z' \subset Z$ et $\xi \in Z'$ comme dans (1).
Alors $\dim(\mathcal{O}_{Z, \xi}) = 0$ (par exemple d'après
Propriétés, Lemme \ref{properties-lemma-codimension-local-ring}).
Ainsi $\mathcal{O}_{Z, \xi}$ est un anneau noethérien
local de dimension zéro, donc de longueur finie sur
lui-même (voir
Algèbre, Proposition \ref{algebra-proposition-dimension-zero-ring}).
Il est donc aussi de longueur finie sur $\mathcal{O}_{X, \xi}$, voir
Algèbre, Lemme \ref{algebra-lemma-length-independent}.

\medskip\noindent
Supposons $\xi \in Z$ et $\delta(\xi) = k$.
Considérons l'adhérence $Z' = \overline{\{\xi\}}$. C'est un sous-schéma
fermé irréductible tel que $\dim_\delta(Z') = k$ par définition.
Comme $\dim_\delta(Z) = k$, ce doit être une composante irréductible
de $Z$. Cela démontre (2).
\end{proof}

\begin{definition}
\label{definition-cycle-associated-to-closed-subscheme}
Soit $(S, \delta)$ comme dans la Situation \ref{situation-setup}.
Soit $X$ localement de type fini sur $S$.
Soit $Z \subset X$ un sous-schéma fermé.
\begin{enumerate}
\item Pour toute composante irréductible $Z' \subset Z$ de point générique $\xi$,
l'entier
$m_{Z', Z} = \text{longueur}_{\mathcal{O}_{X, \xi}} \mathcal{O}_{Z, \xi}$
(Lemme \ref{lemma-multiplicity-finite})
est appelé la {\it multiplicité de $Z'$ dans $Z$}.
\item Supposons $\dim_\delta(Z) \leq k$.
Le {\it $k$-cycle associé à $Z$} est
$$
[Z]_k
=
\sum m_{Z', Z}[Z']
$$
où la somme porte sur les composantes irréductibles de $Z$
de $\delta$-dimension $k$. (C'est un $k$-cycle d'après
Diviseurs, Lemme \ref{divisors-lemma-components-locally-finite}.)
\end{enumerate}
\end{definition}

\noindent
Il importe de noter que nous ne définissons $[Z]_k$ que si la $\delta$-dimension
de $Z$ ne dépasse pas $k$. Autrement dit, par convention, écrire
$[Z]_k$ implique que $\dim_\delta(Z) \leq k$.



\section{Cycle associé à un faisceau cohérent}
\label{section-cycle-of-coherent-sheaf}



\begin{lemma}
\label{lemma-length-finite}
Soit $(S, \delta)$ comme dans la Situation \ref{situation-setup}.
Soit $X$ localement de type fini sur $S$.
Soit $\mathcal{F}$ un $\mathcal{O}_X$-module cohérent.
\begin{enumerate}
\item La collection des composantes irréductibles du support de
$\mathcal{F}$ est localement finie.
\item Soit $Z' \subset \text{Supp}(\mathcal{F})$
une composante irréductible, et
soit $\xi \in Z'$ son point générique.
Alors
$$
\text{longueur}_{\mathcal{O}_{X, \xi}} \mathcal{F}_\xi < \infty
$$
\item Si $\dim_\delta(\text{Supp}(\mathcal{F})) \leq k$
et si $\xi \in \text{Supp}(\mathcal{F})$ vérifie $\delta(\xi) = k$, alors $\xi$ est un
point générique d'une composante irréductible de $\text{Supp}(\mathcal{F})$.
\end{enumerate}
\end{lemma}

\begin{proof}
D'après Cohomologie des schémas, Lemme \ref{coherent-lemma-coherent-support-closed},
le support $Z$ de $\mathcal{F}$ est un fermé de $X$.
Nous pouvons considérer $Z$ comme un sous-schéma fermé réduit de $X$
(Schémas, Lemme \ref{schemes-lemma-reduced-closed-subscheme}).
Ainsi (1) résulte de
Diviseurs, Lemme \ref{divisors-lemma-components-locally-finite}, appliqué à $Z$,
et (3) résulte du
Lemme \ref{lemma-multiplicity-finite}, appliqué à $Z$.

\medskip\noindent
Soit $\xi \in Z'$ comme dans (2). Dans ce cas, pour toute spécialisation
$\xi' \leadsto \xi$ dans $X$, nous avons $\mathcal{F}_{\xi'} = 0$.
Rappelons que les idéaux premiers non maximaux de $\mathcal{O}_{X, \xi}$ correspondent
aux points de $X$ qui se spécialisent en $\xi$
(Schémas, Lemme \ref{schemes-lemma-specialize-points}).
Ainsi $\mathcal{F}_\xi$ est un $\mathcal{O}_{X, \xi}$-module fini
de support $\{\mathfrak m_\xi\}$. Il est donc de longueur finie
d'après Algèbre, Lemme \ref{algebra-lemma-support-point}.
\end{proof}

\begin{definition}
\label{definition-cycle-associated-to-coherent-sheaf}
Soit $(S, \delta)$ comme dans la Situation \ref{situation-setup}.
Soit $X$ localement de type fini sur $S$.
Soit $\mathcal{F}$ un $\mathcal{O}_X$-module cohérent.
\begin{enumerate}
\item Pour toute composante irréductible $Z' \subset \text{Supp}(\mathcal{F})$
de point générique $\xi$, l'entier
$m_{Z', \mathcal{F}} = \text{longueur}_{\mathcal{O}_{X, \xi}} \mathcal{F}_\xi$
(Lemme \ref{lemma-length-finite})
est appelé la {\it multiplicité de $Z'$ dans $\mathcal{F}$}.
\item Supposons $\dim_\delta(\text{Supp}(\mathcal{F})) \leq k$.
Le {\it $k$-cycle associé à $\mathcal{F}$} est
$$
[\mathcal{F}]_k
=
\sum m_{Z', \mathcal{F}}[Z']
$$
où la somme porte sur les composantes irréductibles de
$\text{Supp}(\mathcal{F})$ de $\delta$-dimension $k$.
(C'est un $k$-cycle d'après le Lemme \ref{lemma-length-finite}.)
\end{enumerate}
\end{definition}

\noindent
Il importe de noter que nous ne définissons $[\mathcal{F}]_k$
que si $\mathcal{F}$ est cohérent et si la $\delta$-dimension
de $\text{Supp}(\mathcal{F})$ ne dépasse pas $k$. Autrement dit,
par convention, écrire $[\mathcal{F}]_k$ implique que
$\mathcal{F}$ est cohérent sur $X$ et que
$\dim_\delta(\text{Supp}(\mathcal{F})) \leq k$.

\begin{lemma}
\label{lemma-cycle-closed-coherent}
Soit $(S, \delta)$ comme dans la Situation \ref{situation-setup}.
Soit $X$ localement de type fini sur $S$.
Soit $Z \subset X$ un sous-schéma fermé.
Si $\dim_\delta(Z) \leq k$, alors $[Z]_k = [{\mathcal O}_Z]_k$.
\end{lemma}

\begin{proof}
Cela vient de ce que, dans ce cas, les multiplicités $m_{Z', Z}$ et
$m_{Z', \mathcal{O}_Z}$ coïncident par définition.
\end{proof}

\begin{lemma}
\label{lemma-additivity-sheaf-cycle}
Soit $(S, \delta)$ comme dans la Situation \ref{situation-setup}.
Soit $X$ localement de type fini sur $S$.
Soit $0 \to \mathcal{F} \to \mathcal{G} \to \mathcal{H} \to 0$
une suite exacte courte de faisceaux cohérents sur $X$.
Supposons que la $\delta$-dimension des supports
de $\mathcal{F}$, $\mathcal{G}$ et $\mathcal{H}$ soit $\leq k$.
Alors $[\mathcal{G}]_k = [\mathcal{F}]_k + [\mathcal{H}]_k$.
\end{lemma}

\begin{proof}
Cela résulte immédiatement de l'additivité des longueurs, voir
Algèbre, Lemme \ref{algebra-lemma-length-additive}.
\end{proof}









\section{Préparation à l'image directe propre}
\label{section-preparation-pushforward}

\begin{lemma}
\label{lemma-equal-dimension}
Soit $(S, \delta)$ comme dans la Situation \ref{situation-setup}.
Soient $X$ et $Y$ localement de type fini sur $S$.
Soit $f : X \to Y$ un morphisme.
Supposons $X$ et $Y$ intègres et $\dim_\delta(X) = \dim_\delta(Y)$.
Alors ou bien $f(X)$ est contenu dans un sous-schéma fermé strict
de $Y$, ou bien $f$ est dominant et l'extension de corps de fonctions
$R(X)/R(Y)$ est finie.
\end{lemma}

\begin{proof}
L'adhérence $\overline{f(X)} \subset Y$ est irréductible, puisque $X$
est irréductible (Topologie, Lemmes
\ref{topology-lemma-image-irreducible-space} et
\ref{topology-lemma-irreducible}).
Si $\overline{f(X)} \not = Y$, nous avons terminé.
Si $\overline{f(X)} = Y$, alors $f$ est dominant et, d'après
Morphismes,
Lemme \ref{morphisms-lemma-dominant-finite-number-irreducible-components},
le point générique $\eta_Y$ de $Y$ appartient à l'image de $f$.
Cela implique bien sûr que $f(\eta_X) = \eta_Y$, où $\eta_X \in X$
est le point générique de $X$. Comme $\delta(\eta_X) = \delta(\eta_Y)$,
on voit que $R(Y) = \kappa(\eta_Y) \subset \kappa(\eta_X) = R(X)$
est une extension de degré de transcendance $0$.
Par conséquent, $R(Y) \subset R(X)$ est une extension finie d'après
Morphismes, Lemme \ref{morphisms-lemma-finite-degree}
(qui s'applique en vertu de
Morphismes, Lemme \ref{morphisms-lemma-permanence-finite-type}).
\end{proof}

\begin{lemma}
\label{lemma-quasi-compact-locally-finite}
Soit $(S, \delta)$ comme dans la Situation \ref{situation-setup}.
Soient $X$ et $Y$ localement de type fini sur $S$.
Soit $f : X \to Y$ un morphisme.
Supposons $f$ quasi-compact et $\{Z_i\}_{i \in I}$ une collection localement
finie de fermés de $X$.
Alors $\{\overline{f(Z_i)}\}_{i \in I}$ est une collection localement finie
de fermés de $Y$.
\end{lemma}

\begin{proof}
Soit $V \subset Y$ un ouvert quasi-compact.
Puisque $f$ est quasi-compact, l'ouvert $f^{-1}(V)$ est
quasi-compact. Ainsi l'ensemble
$\{i \in I \mid Z_i \cap f^{-1}(V) \not = \emptyset \}$
est fini par un argument topologique simple que nous omettons.
Comme il s'agit du même ensemble que
$$
\{i \in I \mid f(Z_i) \cap V \not = \emptyset \} =
\{i \in I \mid \overline{f(Z_i)} \cap V \not = \emptyset \}
$$
le lemme est démontré.
\end{proof}









\section{Image directe propre}
\label{section-proper-pushforward}

\begin{definition}
\label{definition-proper-pushforward}
Soit $(S, \delta)$ comme dans la Situation \ref{situation-setup}.
Soient $X$ et $Y$ localement de type fini sur $S$.
Soit $f : X \to Y$ un morphisme.
Supposons $f$ propre.
\begin{enumerate}
\item Soit $Z \subset X$ un sous-schéma fermé intègre
tel que $\dim_\delta(Z) = k$. Nous définissons
$$
f_*[Z] =
\left\{
\begin{matrix}
0 & \text{si} & \dim_\delta(f(Z))< k, \\
\deg(Z/f(Z)) [f(Z)] & \text{si} & \dim_\delta(f(Z)) = k.
\end{matrix}
\right.
$$
Ici, nous considérons $f(Z) \subset Y$ comme un sous-schéma fermé intègre.
Le degré de $Z$ sur $f(Z)$ est fini si
$\dim_\delta(f(Z)) = \dim_\delta(Z)$,
d'après le Lemme \ref{lemma-equal-dimension}.
\item Soit $\alpha = \sum n_Z [Z]$ un $k$-cycle sur $X$. Nous appelons
{\it image directe} de $\alpha$ la somme
$$
f_* \alpha = \sum n_Z f_*[Z]
$$
où chaque $f_*[Z]$ est défini comme ci-dessus. La somme est localement finie
d'après le Lemme \ref{lemma-quasi-compact-locally-finite} ci-dessus.
\end{enumerate}
\end{definition}

\noindent
Par définition, l'image directe propre des cycles
$$
f_* : Z_k(X) \longrightarrow Z_k(Y)
$$
est un homomorphisme de groupes abéliens. Elle fait de $X \mapsto Z_k(X)$
un foncteur covariant sur la catégorie des schémas localement
de type fini sur $S$, dont les morphismes sont les morphismes propres.

\begin{lemma}
\label{lemma-compose-pushforward}
Soit $(S, \delta)$ comme dans la Situation \ref{situation-setup}.
Soient $X$, $Y$ et $Z$ localement de type fini sur $S$.
Soient $f : X \to Y$ et $g : Y \to Z$ des morphismes propres.
Alors $g_* \circ f_* = (g \circ f)_*$ comme applications $Z_k(X) \to Z_k(Z)$.
\end{lemma}

\begin{proof}
Soit $W \subset X$ un sous-schéma fermé intègre de dimension $k$.
Considérons $W' = f(W) \subset Y$ et $W'' = g(f(W)) \subset Z$.
Comme $f$ et $g$ sont propres, $W'$ (respectivement $W''$) est
un sous-schéma fermé intègre de $Y$ (respectivement de $Z$).
Il faut montrer que $g_*(f_*[W]) = (g \circ f)_*[W]$.
Si $\dim_\delta(W'') < k$, les deux membres sont nuls.
Si $\dim_\delta(W'') = k$, les morphismes induits
$$
W \longrightarrow
W' \longrightarrow
W''
$$
satisfont tous deux les hypothèses du Lemme \ref{lemma-equal-dimension}. Ainsi
$$
g_*(f_*[W]) = \deg(W/W')\deg(W'/W'')[W''],
\quad
(g \circ f)_*[W] = \deg(W/W'')[W''].
$$
On peut alors appliquer
Morphismes, Lemme \ref{morphisms-lemma-degree-composition},
pour conclure.
\end{proof}

\noindent
Une immersion fermée est propre. Si $i : Z \to X$ est une immersion fermée,
alors les applications
$$
i_* : Z_k(Z) \longrightarrow Z_k(X)
$$
sont toutes {\it injectives}.

\begin{lemma}
\label{lemma-exact-sequence-closed}
Soit $(S, \delta)$ comme dans la Situation \ref{situation-setup}.
Soit $X$ localement de type fini sur $S$. Soient $X_1, X_2 \subset X$
des sous-schémas fermés tels que $X = X_1 \cup X_2$ ensemblistement.
Pour tout $k \in \mathbf{Z}$, la suite de groupes abéliens
$$
\xymatrix{
Z_k(X_1 \cap X_2) \ar[r] &
Z_k(X_1) \oplus Z_k(X_2) \ar[r] &
Z_k(X) \ar[r] &
0
}
$$
est exacte. Ici, $X_1 \cap X_2$ est l'intersection schématique et
les applications sont les applications d'image directe, l'une étant multipliée par $-1$.
\end{lemma}

\begin{proof}
Supposons d'abord $X$ quasi-compact. Alors $Z_k(X)$ est un $\mathbf{Z}$-module libre
ayant pour base les éléments $[Z]$, où $Z \subset X$ est fermé intègre
de $\delta$-dimension $k$. Les groupes
$Z_k(X_1)$, $Z_k(X_2)$ et $Z_k(X_1 \cap X_2)$ sont libres sur le sous-ensemble des
$Z$ tels que $Z \subset X_1$, $Z \subset X_2$, $Z \subset X_1 \cap X_2$.
Cela démontre immédiatement le lemme dans ce cas. Le cas général est analogue
et la démonstration est omise.
\end{proof}

\begin{lemma}
\label{lemma-cycle-push-sheaf}
Soit $(S, \delta)$ comme dans la Situation \ref{situation-setup}.
Soit $f : X \to Y$ un morphisme propre de schémas qui sont
localement de type fini sur $S$.
\begin{enumerate}
\item Soit $Z \subset X$ un sous-schéma fermé tel que $\dim_\delta(Z) \leq k$.
Alors
$$
f_*[Z]_k = [f_*{\mathcal O}_Z]_k.
$$
\item Soit $\mathcal{F}$ un faisceau cohérent sur $X$ tel que
$\dim_\delta(\text{Supp}(\mathcal{F})) \leq k$. Alors
$$
f_*[\mathcal{F}]_k = [f_*{\mathcal F}]_k.
$$
\end{enumerate}
Notons que l'énoncé a un sens puisque $f_*\mathcal{F}$ et
$f_*\mathcal{O}_Z$ sont des $\mathcal{O}_Y$-modules cohérents d'après
Cohomologie des schémas, Proposition
\ref{coherent-proposition-proper-pushforward-coherent}.
\end{lemma}

\begin{proof}
La partie (1) résulte de (2) et du Lemme \ref{lemma-cycle-closed-coherent}.
Soit $\mathcal{F}$ un faisceau cohérent sur $X$.
Supposons que $\dim_\delta(\text{Supp}(\mathcal{F})) \leq k$.
D'après Cohomologie des schémas, Lemme \ref{coherent-lemma-coherent-support-closed},
il existe un sous-schéma fermé $i : Z \to X$ et un
$\mathcal{O}_Z$-module cohérent $\mathcal{G}$ tels que
$i_*\mathcal{G} \cong \mathcal{F}$ et que le support
de $\mathcal{F}$ soit $Z$. Soit $Z' \subset Y$ l'image schématique
de $f|_Z : Z \to Y$. Considérons le diagramme commutatif de schémas
$$
\xymatrix{
Z \ar[r]_i \ar[d]_{f|_Z} &
X \ar[d]^f \\
Z' \ar[r]^{i'} & Y
}
$$
Nous avons $f_*\mathcal{F} = f_*i_*\mathcal{G} = i'_*(f|_Z)_*\mathcal{G}$
en parcourant le diagramme des deux façons. Supposons le résultat établi
pour les immersions fermées et pour $f|_Z$. Alors
$$
f_*[\mathcal{F}]_k = f_*i_*[\mathcal{G}]_k
= (i')_*(f|_Z)_*[\mathcal{G}]_k =
(i')_*[(f|_Z)_*\mathcal{G}]_k =
[(i')_*(f|_Z)_*\mathcal{G}]_k = [f_*\mathcal{F}]_k
$$
comme voulu. Le cas d'une immersion fermée est immédiat (et omis).
Notons que $f|_Z : Z \to Z'$ est un morphisme dominant (voir
Morphismes, Lemme \ref{morphisms-lemma-quasi-compact-scheme-theoretic-image}).
Nous nous sommes donc ramenés au cas où
$\dim_\delta(X) \leq k$ et $f : X \to Y$ est propre et dominant.

\medskip\noindent
Supposons $\dim_\delta(X) \leq k$ et $f : X \to Y$ propre et dominant.
Puisque $f$ est dominant, pour toute composante irréductible $Z \subset Y$
de point générique $\eta$, il existe un point $\xi \in X$ tel
que $f(\xi) = \eta$. Ainsi $\delta(\eta) \leq \delta(\xi) \leq k$.
On voit donc que, dans les expressions
$$
f_*[\mathcal{F}]_k = \sum n_Z[Z],
\quad
\text{et}
\quad
[f_*\mathcal{F}]_k = \sum m_Z[Z].
$$
dès que $n_Z \not = 0$ ou que $m_Z \not = 0$, le sous-schéma fermé
intègre $Z$ est en fait une composante irréductible de $Y$ de
$\delta$-dimension $k$. Choisissons un tel sous-schéma fermé intègre
$Z \subset Y$ et notons $\eta$ son point générique. Notons que, pour
tout $\xi \in X$ tel que $f(\xi) = \eta$, nous avons $\delta(\xi) \geq k$,
et que $\xi$ est donc lui aussi le point générique d'une composante irréductible
de $X$ de $\delta$-dimension $k$
(voir le Lemme \ref{lemma-multiplicity-finite}). Comme $f$ est quasi-compact
et $X$ localement noethérien, il ne peut y en avoir qu'un nombre fini,
et $f^{-1}(\{\eta\})$ est donc fini.
D'après Morphismes, Lemme \ref{morphisms-lemma-generically-finite}, il existe
un voisinage ouvert $\eta \in V \subset Y$ tel que $f^{-1}(V) \to V$
soit fini. En remplaçant $Y$ par $V$ et $X$ par $f^{-1}(V)$, on se ramène au
cas où $Y$ est affine et où $f$ est fini.

\medskip\noindent
Écrivons $Y = \Spec(R)$ et $X = \Spec(A)$ (ce qui est possible car
un morphisme fini est affine).
Alors $R$ et $A$ sont des anneaux noethériens, et $A$ est fini sur $R$.
De plus, $\mathcal{F} = \widetilde{M}$ pour un certain $A$-module fini
$M$. Notons que $f_*\mathcal{F}$ correspond à $M$ considéré comme $R$-module.
Soit $\mathfrak p \subset R$ l'idéal premier minimal correspondant
à $\eta \in Y$. Le coefficient de $Z$ dans $[f_*\mathcal{F}]_k$
est manifestement $\text{longueur}_{R_{\mathfrak p}}(M_{\mathfrak p})$.
Soient $\mathfrak q_i$, $i = 1, \ldots, t$, les idéaux premiers de $A$
au-dessus de $\mathfrak p$. Alors $A_{\mathfrak p} = \prod A_{\mathfrak q_i}$,
puisque $A_{\mathfrak p}$ est un anneau artinien, étant fini sur l'anneau
local noethérien de dimension zéro $R_{\mathfrak p}$.
Le coefficient de $Z$ dans $f_*[\mathcal{F}]_k$ est manifestement
$$
\sum\nolimits_{i = 1, \ldots, t}
[\kappa(\mathfrak q_i) : \kappa(\mathfrak p)]
\text{longueur}_{A_{\mathfrak q_i}}(M_{\mathfrak q_i})
$$
L'égalité voulue résulte donc de
Algèbre, Lemme \ref{algebra-lemma-pushdown-module}.
\end{proof}




















\section{Préparation à l'image réciproque plate}
\label{section-preparation-flat-pullback}


\noindent
Rappelons qu'un morphisme $f : X \to Y$ qui est localement de type fini
est dit de dimension relative $r$ si toute fibre non vide
est équidimensionnelle de dimension $r$. Voir
Morphismes, Définition \ref{morphisms-definition-relative-dimension-d}.

\begin{lemma}
\label{lemma-flat-inverse-image-dimension}
Soit $(S, \delta)$ comme dans la Situation \ref{situation-setup}.
Soient $X$ et $Y$ localement de type fini sur $S$.
Soit $f : X \to Y$ un morphisme.
Supposons $f$ plat de dimension relative $r$.
Pour tout fermé $Z \subset Y$, nous avons
$$
\dim_\delta(f^{-1}(Z)) = \dim_\delta(Z) + r.
$$
pourvu que $f^{-1}(Z)$ soit non vide.
Si $Z$ est irréductible et si $Z' \subset f^{-1}(Z)$ est une composante irréductible,
alors $Z'$ domine $Z$ et
$\dim_\delta(Z') = \dim_\delta(Z) + r$.
\end{lemma}

\begin{proof}
Il suffit de démontrer la dernière assertion.
Nous pouvons remplacer $Y$ par le sous-schéma fermé intègre $Z$ et
$X$ par l'image réciproque schématique $f^{-1}(Z) = Z \times_Y X$.
Nous pouvons donc supposer $Z = Y$ intègre et $f$ un morphisme plat
de dimension relative $r$. Comme $Y$ est localement noethérien, le
morphisme $f$, qui est localement de type fini,
est en fait localement de présentation finie. Par conséquent,
Morphismes, Lemme \ref{morphisms-lemma-fppf-open},
s'applique et montre que $f$ est ouvert.
Soit $\xi \in X$ le point générique d'une composante irréductible
de $X$. Comme $f$ est ouvert, $f(\xi)$ est le
point générique $\eta$ de $Z = Y$. Notons que $\dim_\xi(X_\eta) = r$,
car $f$ est, par hypothèse, de dimension relative $r$. D'autre
part, comme $\xi$ est un point générique de $X$, on voit que
$\mathcal{O}_{X, \xi} = \mathcal{O}_{X_\eta, \xi}$ ne possède qu'un
idéal premier et est donc de dimension $0$. Ainsi, d'après
Morphismes, Lemme \ref{morphisms-lemma-dimension-fibre-at-a-point},
on conclut que le degré de transcendance
de $\kappa(\xi)$ sur $\kappa(\eta)$ est $r$.
Autrement dit, $\delta(\xi) = \delta(\eta) + r$, comme voulu.
\end{proof}

\noindent
Voici le lemme que nous utiliserons pour démontrer que l'image réciproque plate
d'une collection localement finie de sous-schémas fermés est localement finie.

\begin{lemma}
\label{lemma-inverse-image-locally-finite}
Soit $(S, \delta)$ comme dans la Situation \ref{situation-setup}.
Soient $X$ et $Y$ localement de type fini sur $S$.
Soit $f : X \to Y$ un morphisme.
Supposons que $\{Z_i\}_{i \in I}$ soit une collection localement
finie de fermés de $Y$.
Alors $\{f^{-1}(Z_i)\}_{i \in I}$ est une collection localement finie
de fermés de $X$.
\end{lemma}

\begin{proof}
Soit $U \subset X$ un ouvert quasi-compact.
Comme l'image $f(U) \subset Y$ est un sous-ensemble quasi-compact,
il existe un ouvert quasi-compact $V \subset Y$ tel que
$f(U) \subset V$. Notons que
$$
\{i \in I \mid f^{-1}(Z_i) \cap U \not = \emptyset \}
\subset
\{i \in I \mid Z_i \cap V \not = \emptyset \}.
$$
Le membre de droite étant fini par hypothèse, nous avons terminé.
\end{proof}



\section{Image réciproque plate}
\label{section-flat-pullback}

\noindent
Dans la suite, nous appelons $f^{-1}(Z)$
l'{\it image réciproque schématique} d'un sous-schéma fermé
$Z \subset Y$ par un morphisme de schémas $f : X \to Y$.
Rappelons que l'image réciproque schématique est le produit fibré
$$
\xymatrix{
f^{-1}(Z) \ar[r] \ar[d] & X \ar[d] \\
Z \ar[r] & Y
}
$$
et qu'elle est aussi le sous-schéma fermé de $X$ défini par le
faisceau quasi-cohérent d'idéaux $f^{-1}(\mathcal{I})\mathcal{O}_X$, si
$\mathcal{I} \subset \mathcal{O}_Y$ est le faisceau quasi-cohérent d'idéaux
qui correspond à $Z$ dans $Y$.
(Ceci est étudié dans
Schémas, section \ref{schemes-section-closed-immersion} et
Lemme \ref{schemes-lemma-fibre-product-immersion},
ainsi que dans la Définition \ref{schemes-definition-inverse-image-closed-subscheme}.)

\begin{definition}
\label{definition-flat-pullback}
Soit $(S, \delta)$ comme dans la Situation \ref{situation-setup}.
Soient $X$ et $Y$ localement de type fini sur $S$.
Soit $f : X \to Y$ un morphisme.
Supposons $f$ plat de dimension relative $r$.
\begin{enumerate}
\item Soit $Z \subset Y$ un sous-schéma fermé intègre de
$\delta$-dimension $k$. Nous définissons $f^*[Z]$ comme le
$(k+r)$-cycle sur $X$ associé à l'image réciproque schématique
$$
f^*[Z] = [f^{-1}(Z)]_{k+r}.
$$
Cela a un sens puisque $\dim_\delta(f^{-1}(Z)) = k + r$
d'après le Lemme \ref{lemma-flat-inverse-image-dimension}.
\item Soit $\alpha = \sum n_i [Z_i]$ un
$k$-cycle sur $Y$. L'{\it image réciproque plate de $\alpha$ par $f$}
est la somme
$$
f^* \alpha = \sum n_i f^*[Z_i]
$$
où chaque $f^*[Z_i]$ est défini comme ci-dessus.
La somme est localement finie d'après le Lemme \ref{lemma-inverse-image-locally-finite}.
\item Nous notons $f^* : Z_k(Y) \to Z_{k + r}(X)$ l'application de groupes
abéliens ainsi obtenue.
\end{enumerate}
\end{definition}

\noindent
Une immersion ouverte est plate. C'est un cas particulier important, quoique trivial,
d'un morphisme plat. Si $U \subset X$ est ouvert, on appelle parfois
l'image réciproque par $j : U \to X$ d'un cycle la {\it restriction} du
cycle à $U$. Notons que, dans ce cas, les applications
$$
j^* : Z_k(X) \longrightarrow Z_k(U)
$$
sont toutes {\it surjectives}. En effet, étant donné tout sous-schéma fermé
intègre $Z' \subset U$, on peut prendre l'adhérence $Z$ de $Z'$ dans $X$
et la considérer comme un sous-schéma fermé réduit de $X$ (voir
Schémas, Lemme \ref{schemes-lemma-reduced-closed-subscheme}).
Et, manifestement, $Z \cap U = Z'$ ; autrement dit,
$j^*[Z] = [Z']$, d'où la surjectivité. En fait, un peu plus
est vrai.

\begin{lemma}
\label{lemma-exact-sequence-open}
Soit $(S, \delta)$ comme dans la Situation \ref{situation-setup}.
Soit $X$ localement de type fini sur $S$.
Soit $U \subset X$ un sous-schéma ouvert et notons
$i : Y = X \setminus U \to X$ le sous-schéma fermé réduit de $X$.
Pour tout $k \in \mathbf{Z}$, la suite
$$
\xymatrix{
Z_k(Y) \ar[r]^{i_*} & Z_k(X) \ar[r]^{j^*} & Z_k(U) \ar[r] & 0
}
$$
est un complexe exact de groupes abéliens.
\end{lemma}

\begin{proof}
Supposons d'abord $X$ quasi-compact. Alors $Z_k(X)$ est un $\mathbf{Z}$-module libre
ayant pour base les éléments $[Z]$, où $Z \subset X$ est fermé intègre
de $\delta$-dimension $k$. Un tel élément de base s'envoie
soit sur l'élément de base $[Z \cap U]$, soit sur zéro si $Z \subset Y$.
Le lemme est donc clair dans ce cas. Le cas général est analogue
et la démonstration est omise.
\end{proof}

\begin{lemma}
\label{lemma-compose-flat-pullback}
Soit $(S, \delta)$ comme dans la Situation \ref{situation-setup}.
Soient $X, Y, Z$ localement de type fini sur $S$.
Soient $f : X \to Y$ et $g : Y \to Z$ des morphismes plats de dimensions relatives
$r$ et $s$. Alors $g \circ f$ est plat de dimension relative
$r + s$ et
$$
f^* \circ g^* = (g \circ f)^*
$$
comme applications $Z_k(Z) \to Z_{k + r + s}(X)$.
\end{lemma}

\begin{proof}
La composée est plate de dimension relative $r + s$ d'après
Morphismes, Lemme \ref{morphisms-lemma-composition-relative-dimension-d}.
Supposons que
\begin{enumerate}
\item $W \subset Z$ soit un sous-schéma fermé intègre de $\delta$-dimension $k$,
\item $W' \subset Y$ soit un sous-schéma fermé intègre de $\delta$-dimension
$k + s$ tel que $W' \subset g^{-1}(W)$, et
\item $W'' \subset Y$ soit un sous-schéma fermé intègre de $\delta$-dimension
$k + s + r$ tel que $W'' \subset f^{-1}(W')$.
\end{enumerate}
Il faut montrer que le coefficient $n$ de $[W'']$ dans
$(g \circ f)^*[W]$ coïncide avec le coefficient $m$ de
$[W'']$ dans $f^*(g^*[W])$. Le fait qu'il suffise de vérifier le lemme dans ces
cas résulte du Lemme \ref{lemma-flat-inverse-image-dimension}.
Soient $\xi'' \in W''$, $\xi' \in W'$
et $\xi \in W$ les points génériques. Considérons les anneaux locaux
$A = \mathcal{O}_{Z, \xi}$, $B = \mathcal{O}_{Y, \xi'}$
et $C = \mathcal{O}_{X, \xi''}$. Nous avons alors des morphismes locaux plats d'anneaux
$A \to B$, $B \to C$ et, de plus,
$$
n = \text{longueur}_C(C/\mathfrak m_AC),
\quad
\text{et}
\quad
m = \text{longueur}_C(C/\mathfrak m_BC) \text{longueur}_B(B/\mathfrak m_AB)
$$
L'égalité résulte donc de
Algèbre, Lemme \ref{algebra-lemma-pullback-transitive}.
\end{proof}

\begin{lemma}
\label{lemma-pullback-coherent}
Soit $(S, \delta)$ comme dans la Situation \ref{situation-setup}.
Soient $X, Y$ localement de type fini sur $S$.
Soit $f : X \to Y$ un morphisme plat de dimension relative $r$.
\begin{enumerate}
\item Soit $Z \subset Y$ un sous-schéma fermé tel que
$\dim_\delta(Z) \leq k$. Alors
$\dim_\delta(f^{-1}(Z)) \leq k + r$
et $[f^{-1}(Z)]_{k + r} = f^*[Z]_k$ dans $Z_{k + r}(X)$.
\item Soit $\mathcal{F}$ un faisceau cohérent sur $Y$ tel que
$\dim_\delta(\text{Supp}(\mathcal{F})) \leq k$.
Alors $\dim_\delta(\text{Supp}(f^*\mathcal{F})) \leq k + r$
et
$$
f^*[{\mathcal F}]_k = [f^*{\mathcal F}]_{k+r}
$$
dans $Z_{k + r}(X)$.
\end{enumerate}
\end{lemma}

\begin{proof}
Les assertions sur les dimensions résultent immédiatement du
Lemme \ref{lemma-flat-inverse-image-dimension}.
La partie (1) résulte de la partie (2), du Lemme \ref{lemma-cycle-closed-coherent}
et du fait que $f^*\mathcal{O}_Z = \mathcal{O}_{f^{-1}(Z)}$.

\medskip\noindent
Démonstration de (2).
Comme $X$ et $Y$ sont localement noethériens, nous pouvons appliquer
Cohomologie des schémas, Lemme \ref{coherent-lemma-coherent-Noetherian}, pour voir
que $\mathcal{F}$ est de type fini, donc que $f^*\mathcal{F}$ est
de type fini (Modules, Lemme \ref{modules-lemma-pullback-finite-type}),
et donc que $f^*\mathcal{F}$ est cohérent
(de nouveau Cohomologie des schémas, Lemme \ref{coherent-lemma-coherent-Noetherian}).
Le lemme a donc un sens. Soit $W \subset Y$ un sous-schéma fermé intègre
de $\delta$-dimension $k$, et soit $W' \subset X$ un sous-schéma fermé
intègre de dimension $k + r$ qui s'envoie dans $W$
par $f$. Il faut montrer que le coefficient $n$ de
$[W']$ dans $f^*[{\mathcal F}]_k$ coïncide avec le coefficient
$m$ de $[W']$ dans $[f^*{\mathcal F}]_{k+r}$. Soient $\xi \in W$ et
$\xi' \in W'$ soient les points génériques. Posons
$A = \mathcal{O}_{Y, \xi}$, $B = \mathcal{O}_{X, \xi'}$
et considérons $M = \mathcal{F}_\xi$ comme un $A$-module. (Notons que
$M$ est de longueur finie d'après nos hypothèses sur les dimensions, mais nous
n'avons en fait pas besoin de le vérifier. Voir le
Lemme \ref{lemma-length-finite}.)
Nous avons $f^*\mathcal{F}_{\xi'} = B \otimes_A M$.
Nous voyons ainsi que
$$
n = \text{longueur}_B(B \otimes_A M)
\quad
\text{et}
\quad
m = \text{longueur}_A(M) \text{longueur}_B(B/\mathfrak m_AB)
$$
L'égalité résulte donc du
Lemme \ref{algebra-lemma-pullback-module} d'Algèbre.
\end{proof}



\section{Images directes et images réciproques}
\label{section-push-pull}

\noindent
Dans cette section, nous vérifions que l'image directe propre et l'image réciproque plate
sont compatibles lorsque cela a un sens. D'après le travail effectué ci-dessus, ceci
est une conséquence de la cohomologie et du changement de base.

\begin{lemma}
\label{lemma-flat-pullback-proper-pushforward}
Soit $(S, \delta)$ comme dans la Situation \ref{situation-setup}.
Soit
$$
\xymatrix{
X' \ar[r]_{g'} \ar[d]_{f'} & X \ar[d]^f \\
Y' \ar[r]^g & Y
}
$$
un diagramme cartésien de schémas localement de type fini sur $S$.
Supposons $f : X \to Y$ propre et $g : Y' \to Y$ plat de dimension relative $r$.
Alors $f'$ est aussi propre et $g'$ est aussi plat de dimension relative $r$.
Pour tout $k$-cycle $\alpha$ sur $X$, nous avons
$$
g^*f_*\alpha = f'_*(g')^*\alpha
$$
dans $Z_{k + r}(Y')$.
\end{lemma}

\begin{proof}
Le fait que $f'$ soit propre résulte du
Lemme \ref{morphisms-lemma-base-change-proper} de Morphismes.
Le fait que $g'$ soit plat de dimension relative $r$ résulte des
Lemmes \ref{morphisms-lemma-base-change-relative-dimension-d}
et \ref{morphisms-lemma-base-change-flat} de Morphismes.
Il suffit de démontrer l'égalité de cycles lorsque $\alpha = [W]$
pour un sous-schéma fermé intègre $W \subset X$ de $\delta$-dimension $k$.
Notons que, dans ce cas, nous avons $\alpha = [\mathcal{O}_W]_k$ ; voir le
Lemme \ref{lemma-cycle-closed-coherent}.
D'après les Lemmes \ref{lemma-cycle-push-sheaf} et
\ref{lemma-pullback-coherent}, il suffit donc
de montrer que $f'_*(g')^*\mathcal{O}_W$ est isomorphe à
$g^*f_*\mathcal{O}_W$. Ceci résulte de la cohomologie et du
changement de base ; voir le
Lemme \ref{coherent-lemma-flat-base-change-cohomology} de Cohomologie des schémas.
\end{proof}

\begin{lemma}
\label{lemma-finite-flat}
Soit $(S, \delta)$ comme dans la Situation \ref{situation-setup}.
Soient $X$, $Y$ localement de type fini sur $S$.
Soit $f : X \to Y$ un morphisme fini localement libre
de degré $d$ (voir la
Définition \ref{morphisms-definition-finite-locally-free} de Morphismes).
Alors $f$ est à la fois propre et plat de dimension relative $0$, et
$$
f_*f^*\alpha = d\alpha
$$
pour tout $\alpha \in Z_k(Y)$.
\end{lemma}

\begin{proof}
Un morphisme fini localement libre est plat et fini d'après le
Lemme \ref{morphisms-lemma-finite-flat} de Morphismes,
et un morphisme fini est propre
d'après le Lemme \ref{morphisms-lemma-finite-proper} de Morphismes.
Nous omettons de montrer qu'un morphisme fini
est de dimension relative $0$. La formule a donc un sens.
Pour la démontrer, soit $Z \subset Y$ un sous-schéma fermé intègre
de $\delta$-dimension $k$. Il suffit de démontrer la formule
pour $\alpha = [Z]$. Puisque le changement de base d'un morphisme fini localement libre
est fini localement libre
(Lemme \ref{morphisms-lemma-base-change-finite-locally-free} de Morphismes),
nous voyons que $f_*f^*\mathcal{O}_Z$ est un faisceau fini localement libre de
rang $d$ sur $Z$. Ainsi
$$
f_*f^*[Z] = f_*f^*[\mathcal{O}_Z]_k =
[f_*f^*\mathcal{O}_Z]_k = d[Z]
$$
où nous avons utilisé les Lemmes \ref{lemma-pullback-coherent} et
\ref{lemma-cycle-push-sheaf}.
\end{proof}








\section{Préparation aux diviseurs principaux}
\label{section-preparation-principal-divisors}

\noindent
Une partie du contenu de cette section recoupe partiellement la
discussion de Diviseurs, section \ref{divisors-section-Weil-divisors}.

\begin{lemma}
\label{lemma-divisor-delta-dimension}
Soit $(S, \delta)$ comme dans la Situation \ref{situation-setup}.
Soit $X$ localement de type fini sur $S$. Supposons $X$
intègre.
\begin{enumerate}
\item Si $Z \subset X$ est un sous-schéma fermé intègre, alors
les conditions suivantes sont équivalentes :
\begin{enumerate}
\item $Z$ est un diviseur premier,
\item $Z$ est de codimension $1$ dans $X$, et
\item $\dim_\delta(Z) = \dim_\delta(X) - 1$.
\end{enumerate}
\item Si $Z$ est une composante irréductible d'un diviseur de Cartier
effectif sur $X$, alors $\dim_\delta(Z) = \dim_\delta(X) - 1$.
\end{enumerate}
\end{lemma}

\begin{proof}
La partie (1) résulte de la définition d'un diviseur premier
(Définition \ref{divisors-definition-Weil-divisor} de Diviseurs)
et de la définition d'une fonction de dimension
(Définition \ref{topology-definition-dimension-function} de Topologie).
Soit $\xi \in Z$ le point générique d'une composante irréductible $Z$ d'un
diviseur de Cartier effectif $D \subset X$.
Alors $\dim(\mathcal{O}_{D, \xi}) = 0$ et
$\mathcal{O}_{D, \xi} = \mathcal{O}_{X, \xi}/(f)$ pour un certain
élément non diviseur de zéro $f \in \mathcal{O}_{X, \xi}$ (Lemme
\ref{divisors-lemma-effective-Cartier-in-points} de Diviseurs).
Alors $\dim(\mathcal{O}_{X, \xi}) = 1$ d'après le
Lemme \ref{algebra-lemma-one-equation} d'Algèbre. Ainsi $Z$ est comme dans (1) d'après le
Lemme \ref{properties-lemma-codimension-local-ring} de Propriétés,
ce qui achève la démonstration.
\end{proof}

\begin{lemma}
\label{lemma-finite-in-codimension-one}
Soit $f : X \to Y$ un morphisme de schémas.
Soit $\xi \in Y$ un point.
Supposons que
\begin{enumerate}
\item $X$, $Y$ sont intègres,
\item $Y$ est localement noethérien,
\item $f$ est propre, dominant et l'extension $R(Y) \subset R(X)$ est finie, et
\item $\dim(\mathcal{O}_{Y, \xi}) = 1$.
\end{enumerate}
Alors il existe un voisinage ouvert $V \subset Y$ de $\xi$
tel que $f|_{f^{-1}(V)} : f^{-1}(V) \to V$ soit fini.
\end{lemma}

\begin{proof}
Ce lemme est un cas particulier du
Lemme \ref{varieties-lemma-finite-in-codim-1} de Variétés.
Voici un argument direct dans ce cas.
D'après le Lemme
\ref{coherent-lemma-proper-finite-fibre-finite-in-neighbourhood} de Cohomologie des schémas,
il suffit de démontrer que $f^{-1}(\{\xi\})$ est fini.
Nous remplaçons $Y$ par un ouvert affine, disons $Y = \Spec(R)$.
Notons que $R$ est noethérien, puisque $Y$ est supposé localement noethérien.
Puisque $f$ est propre, il est quasi-compact. Nous pouvons donc trouver un recouvrement
ouvert affine fini $X = U_1 \cup \ldots \cup U_n$, avec
chaque $U_i = \Spec(A_i)$. Notons que $R \to A_i$ est un
homomorphisme injectif de type fini de domaines tel que
l'extension induite des corps des fractions soit finie.
Le lemme résulte donc
du Lemme \ref{algebra-lemma-finite-in-codim-1} d'Algèbre.
\end{proof}


\section{Diviseurs principaux}
\label{section-principal-divisors}

\noindent
La définition suivante est l'analogue de la
Définition \ref{divisors-definition-principal-divisor} de Diviseurs
dans notre cadre actuel.

\begin{definition}
\label{definition-principal-divisor}
Soit $(S, \delta)$ comme dans la Situation \ref{situation-setup}.
Soit $X$ localement de type fini sur $S$. Supposons $X$
intègre avec $\dim_\delta(X) = n$.
Soit $f \in R(X)^*$. Le {\it diviseur principal
associé à $f$} est le $(n - 1)$-cycle
$$
\text{div}(f) = \text{div}_X(f) = \sum \text{ord}_Z(f) [Z]
$$
défini dans Diviseurs, Définition \ref{divisors-definition-principal-divisor}.
Ceci a un sens parce que les diviseurs premiers sont de $\delta$-dimension $n - 1$ d'après le
Lemme \ref{lemma-divisor-delta-dimension}.
\end{definition}

\noindent
Dans la situation de la définition, pour $f, g \in R(X)^*$, nous avons
$$
\text{div}_X(fg) = \text{div}_X(f) + \text{div}_X(g)
$$
dans $Z_{n - 1}(X)$. Voir Diviseurs, Lemme \ref{divisors-lemma-div-additive}.
Le lemme suivant sera supplanté par le Lemme plus général
\ref{lemma-flat-pullback-rational-equivalence}.

\begin{lemma}
\label{lemma-flat-pullback-principal-divisor}
Soit $(S, \delta)$ comme dans la Situation \ref{situation-setup}.
Soient $X$, $Y$ localement de type fini sur $S$. Supposons $X$, $Y$
intègres et $n = \dim_\delta(Y)$.
Soit $f : X \to Y$ un morphisme plat de dimension relative $r$.
Soit $g \in R(Y)^*$. Alors
$$
f^*(\text{div}_Y(g)) = \text{div}_X(g)
$$
dans $Z_{n + r - 1}(X)$.
\end{lemma}

\begin{proof}
Notons que, puisque $f$ est plat, il est dominant, de sorte que
$f$ induit un plongement $R(Y) \subset R(X)$ ; nous pouvons donc
considérer $g$ comme un élément de $R(X)^*$.
Soit $Z \subset X$ un sous-schéma fermé intègre de
$\delta$-dimension $n + r - 1$. Soit $\xi \in Z$
son point générique. Si $\dim_\delta(f(Z)) > n - 1$,
alors nous voyons que le coefficient de $[Z]$ dans les membres de gauche et
de droite de l'équation est nul.
Nous pouvons donc supposer que $Z' = \overline{f(Z)}$ est un
sous-schéma fermé intègre de $Y$ de $\delta$-dimension $n - 1$.
Soit $\xi' = f(\xi)$. C'est le point générique de $Z'$.
Posons $A = \mathcal{O}_{Y, \xi'}$, $B = \mathcal{O}_{X, \xi}$.
L'homomorphisme d'anneaux $A \to B$ est un homomorphisme local plat de
domaines locaux noethériens de dimension $1$.
Nous avons $g$ dans le corps des fractions de $A$. Ce que nous devons montrer est que
$$
\text{ord}_A(g) \text{longueur}_B(B/\mathfrak m_AB)
=
\text{ord}_B(g).
$$
Ceci résulte du Lemme \ref{algebra-lemma-pullback-module} d'Algèbre
(détails omis).
\end{proof}











\section{Diviseurs principaux et image directe}
\label{section-two-fun}

\noindent
Le premier lemme implique que l'image directe d'un diviseur
principal par un morphisme génériquement fini est un diviseur principal.

\begin{lemma}
\label{lemma-proper-pushforward-alteration}
Soit $(S, \delta)$ comme dans la Situation \ref{situation-setup}.
Soient $X$, $Y$ localement de type fini sur $S$. Supposons $X$, $Y$
intègres et $n = \dim_\delta(X) = \dim_\delta(Y)$.
Soit $p : X \to Y$ un morphisme propre dominant.
Soit $f \in R(X)^*$. Posons
$$
g = \text{Nm}_{R(X)/R(Y)}(f).
$$
Alors nous avons
$p_*\text{div}(f) = \text{div}(g)$.
\end{lemma}

\begin{proof}
Soit $Z \subset Y$ un sous-schéma fermé intègre de $\delta$-dimension
$n - 1$. Nous voulons montrer que les coefficients de $[Z]$ dans
$p_*\text{div}(f)$ et $\text{div}(g)$ sont égaux. Nous pouvons appliquer le
Lemme \ref{lemma-finite-in-codimension-one}
au morphisme $p : X \to Y$ et au point générique $\xi \in Z$.
Nous pouvons donc remplacer $Y$ par un
voisinage ouvert affine de $\xi$ et supposer que $p : X \to Y$ est fini.
Écrivons $Y = \Spec(R)$ et $X = \Spec(A)$, où $p$ est induit
par un homomorphisme fini $R \to A$ de domaines noethériens qui induit
une extension finie $L/K$ des corps des fractions.
Nous avons alors $f \in L$, $g = \text{Nm}(f) \in K$,
et un idéal premier $\mathfrak p \subset R$ avec $\dim(R_{\mathfrak p}) = 1$.
Le coefficient de $[Z]$ dans $\text{div}_Y(g)$ est
$\text{ord}_{R_\mathfrak p}(g)$.
Le coefficient de $[Z]$ dans $p_*\text{div}_X(f)$ est
$$
\sum\nolimits_{\mathfrak q\text{ au-dessus de }\mathfrak p}
[\kappa(\mathfrak q) : \kappa(\mathfrak p)]
\text{ord}_{A_{\mathfrak q}}(f)
$$
L'égalité souhaitée résulte donc du
Lemme \ref{algebra-lemma-finite-extension-dim-1} d'Algèbre.
\end{proof}

\noindent
Dans l'étude des diviseurs principaux, un rôle important
est joué par le diviseur principal « universel » $[0] - [\infty]$
sur $\mathbf{P}^1_S$. Pour préciser cela, notons
\begin{equation}
\label{equation-zero-infty}
D_0, D_\infty \subset
\mathbf{P}^1_S = \underline{\text{Proj}}_S(\mathcal{O}_S[T_0, T_1])
\end{equation}
les sous-schémas fermés définis par la section $T_1$, resp.\ $T_0$
de $\mathcal{O}(1)$. Ce sont des diviseurs de Cartier effectifs ; voir la
Définition \ref{divisors-definition-effective-Cartier-divisor}
et le Lemme \ref{divisors-lemma-characterize-OD} de Diviseurs.
Le lemme suivant dit que, de façon informelle, nous avons
« $\text{div}(T_1/T_0) = [D_0] - [D_1]$ » et que c'est le
diviseur principal universel.

\begin{lemma}
\label{lemma-rational-function}
Soit $(S, \delta)$ comme dans la Situation \ref{situation-setup}.
Soit $X$ localement de type fini sur $S$. Supposons $X$
intègre et $n = \dim_\delta(X)$. Soit $f \in R(X)^*$.
Soit $U \subset X$ un ouvert non vide tel que $f$
corresponde à une section $f \in \Gamma(U, \mathcal{O}_X^*)$.
Soit $Y \subset X \times_S \mathbf{P}^1_S$ l'adhérence du
graphe de $f : U \to \mathbf{P}^1_S$.
Alors
\begin{enumerate}
\item le morphisme de projection $p : Y \to X$ est propre,
\item $p|_{p^{-1}(U)} : p^{-1}(U) \to U$ est un isomorphisme,
\item les images réciproques $Y_0 = q^{-1}D_0$ et $Y_\infty = q^{-1}D_\infty$
par le morphisme $q : Y \to \mathbf{P}^1_S$ sont définies
(Diviseurs, Définition
\ref{divisors-definition-pullback-effective-Cartier-divisor}),
\item nous avons
$$
\text{div}_Y(f) = [Y_0]_{n - 1} - [Y_\infty]_{n - 1}
$$
\item nous avons
$$
\text{div}_X(f) = p_*\text{div}_Y(f)
$$
\item si nous considérons $Y_0$ et $Y_\infty$ comme des sous-schémas fermés de $X$
par le morphisme $p$, alors nous avons
$$
\text{div}_X(f) = [Y_0]_{n - 1} - [Y_\infty]_{n - 1}
$$
\end{enumerate}
\end{lemma}

\begin{proof}
Puisque $X$ est intègre, nous voyons que $U$ est intègre.
Ainsi $Y$ est intègre, et $(1, f)(U) \subset Y$ est un sous-schéma ouvert dense.
Notons aussi que le sous-schéma fermé $Y \subset X \times_S \mathbf{P}^1_S$
ne dépend pas du choix de l'ouvert $U$, puisqu'il s'agit en définitive de
l'adhérence de l'ensemble à un point $\{\eta'\} = \{(1, f)(\eta)\}$,
où $\eta \in X$ est le point générique. Cela étant dit, démontrons les
assertions du lemme.

\medskip\noindent
Pour (1), notons que $p$ est la composée de l'immersion fermée
$Y \to X \times_S \mathbf{P}^1_S = \mathbf{P}^1_X$ et du morphisme propre
$\mathbf{P}^1_X \to X$. Comme une composée de morphismes propres
est propre (Morphismes, Lemme \ref{morphisms-lemma-composition-proper}),
nous concluons.

\medskip\noindent
Il est clair que $Y \cap U \times_S \mathbf{P}^1_S = (1, f)(U)$.
Ainsi (2) en résulte. Il en résulte aussi que $\dim_\delta(Y) = n$.

\medskip\noindent
Notons que $q(\eta') = f(\eta)$ n'est contenu ni dans $D_0$ ni dans $D_\infty$,
puisque $f \in R(X)^*$. D'où (3) par le
Lemme \ref{divisors-lemma-pullback-effective-Cartier-defined} de Diviseurs.
Nous obtenons $\dim_\delta(Y_0) = n - 1$
et $\dim_\delta(Y_\infty) = n - 1$ par le
Lemme \ref{lemma-divisor-delta-dimension}.

\medskip\noindent
Considérons le diviseur de Cartier effectif $Y_0$.
En tout point $\xi \in Y_0$, nous avons $f \in \mathcal{O}_{Y, \xi}$ et
une équation locale de $Y_0$ est donnée par $f$.
En particulier, si $\delta(\xi) = n - 1$, de sorte que $\xi$ est le point générique
d'un sous-schéma fermé intègre $Z$ de $\delta$-dimension $n - 1$,
alors nous voyons que le coefficient de $[Z]$ dans $\text{div}_Y(f)$ est
$$
\text{ord}_Z(f) =
\text{longueur}_{\mathcal{O}_{Y, \xi}}
(\mathcal{O}_{Y, \xi}/f\mathcal{O}_{Y, \xi}) =
\text{longueur}_{\mathcal{O}_{Y, \xi}}
(\mathcal{O}_{Y_0, \xi})
$$
qui est le coefficient de $[Z]$ dans $[Y_0]_{n - 1}$. Un argument semblable
utilisant la fonction rationnelle $1/f$ montre que
$-[Y_\infty]$ coïncide avec les termes de coefficients négatifs dans
l'expression de $\text{div}_Y(f)$. D'où (4).

\medskip\noindent
Notons que $D_0 \to S$ est un isomorphisme. Nous voyons donc que
$X \times_S D_0 \to X$ est également un isomorphisme. Il est clair que
nous avons $Y_0 = Y \cap X \times_S D_0$ (intersection schématique)
dans $X \times_S \mathbf{P}^1_S$. Il s'agit donc bien du cas où
$Y_0 \to X$ est une immersion fermée. Il en résulte que
$$
p_*\mathcal{O}_{Y_0} = \mathcal{O}_{Y'_0}
$$
où $Y'_0 \subset X$ est l'image de $Y_0 \to X$.
D'après le Lemme \ref{lemma-cycle-push-sheaf}, nous
avons $p_*[Y_0]_{n - 1} = [Y'_0]_{n - 1}$. Il en va de même
pour $D_\infty$ et $Y_\infty$. Ainsi (6) est une conséquence de (5).
Enfin, (5) résulte immédiatement du
Lemme \ref{lemma-proper-pushforward-alteration}.
\end{proof}

\noindent
Le lemme suivant dit que le degré d'un diviseur principal sur
une courbe propre est nul.

\begin{lemma}
\label{lemma-curve-principal-divisor}
Soit $K$ un corps quelconque. Soit $X$ un schéma intègre de dimension $1$
muni d'un morphisme propre $c : X \to \Spec(K)$.
Soit $f \in K(X)^*$ une fonction rationnelle inversible.
Alors
$$
\sum\nolimits_{x \in X \text{ fermé}}
[\kappa(x) : K] \text{ord}_{\mathcal{O}_{X, x}}(f)
=
0
$$
où $\text{ord}$ est comme dans la
Définition \ref{algebra-definition-ord} d'Algèbre.
Autrement dit, $c_*\text{div}(f) = 0$.
\end{lemma}

\begin{proof}
Considérons le diagramme
$$
\xymatrix{
Y \ar[r]_p \ar[d]_q & X \ar[d]^c \\
\mathbf{P}^1_K \ar[r]^-{c'} & \Spec(K)
}
$$
que nous avons construit dans le Lemme \ref{lemma-rational-function}
à partir de $X$ et de la fonction rationnelle $f$ sur $S = \Spec(K)$.
Nous utiliserons sans autre mention tous les résultats de ce lemme.
Nous devons montrer que $c_*\text{div}_X(f) = c_*p_*\text{div}_Y(f) = 0$.
Cela revient à démontrer que $c'_*q_*\text{div}_Y(f) = 0$.
Si $q(Y)$ est un point fermé de $\mathbf{P}^1_K$, alors nous
voyons que $\text{div}_X(f) = 0$ et le lemme est démontré.
Nous pouvons donc supposer que $q$ est dominant.
Supposons que nous puissions montrer que $q : Y \to \mathbf{P}^1_K$ est fini
localement libre de degré $d$ (voir la
Définition \ref{morphisms-definition-finite-locally-free} de Morphismes).
Puisque $\text{div}_Y(f) = [q^{-1}D_0]_0 - [q^{-1}D_\infty]_0$,
nous voyons (par définition de l'image réciproque plate) que
$\text{div}_Y(f) = q^*([D_0]_0 - [D_\infty]_0)$.
Alors le Lemme \ref{lemma-finite-flat} donne
$q_*\text{div}_Y(f) = d([D_0]_0 - [D_\infty]_0)$.
Comme il est clair que $c'_*[D_0]_0 = c'_*[D_\infty]_0$, nous avons le résultat.

\medskip\noindent
Il reste à montrer que $q$ est fini localement libre.
(Il aura automatiquement un certain degré donné puisque $\mathbf{P}^1_K$
est connexe.)
Comme $\dim(\mathbf{P}^1_K) = 1$, nous voyons que $q$ est fini, par exemple
d'après le Lemme \ref{lemma-finite-in-codimension-one}.
Tous les anneaux locaux de $\mathbf{P}^1_K$ aux
points fermés sont des anneaux locaux réguliers de dimension $1$
(autrement dit, des anneaux de valuation discrète), puisqu'ils sont des
localisés de $K[T]$ (voir le
Lemme \ref{algebra-lemma-dim-affine-space} d'Algèbre).
Ainsi, pour $y\in Y$ fermé, l'anneau local $\mathcal{O}_{Y, y}$
sera plat sur $\mathcal{O}_{\mathbf{P}^1_K, q(y)}$ dès qu'il est
sans torsion (Plus sur l'algèbre, Lemme
\ref{more-algebra-lemma-dedekind-torsion-free-flat}).
C'est manifestement le cas puisque
$\mathcal{O}_{Y, y}$ est un domaine et $q$ est dominant.
Ainsi $q$ est plat. Par conséquent, $q$ est fini localement libre d'après le
Lemme \ref{morphisms-lemma-finite-flat} de Morphismes.
\end{proof}





\section{Équivalence rationnelle}
\label{section-rational-equivalence}

\noindent
Dans cette section, nous définissons l'{\it équivalence rationnelle} sur les $k$-cycles.
Nous autoriserons les sommes localement finies d'images de
diviseurs principaux (par des immersions fermées). Cela conduit à des
phénomènes assez étranges ; voir l'Exemple \ref{example-weird}.
Cependant, si nous ne les autorisons pas, nous ne savons pas démontrer que
le produit avec les classes de Chern de fibrés inversibles se factorise par l'équivalence
rationnelle.

\begin{definition}
\label{definition-rational-equivalence}
Soit $(S, \delta)$ comme dans la Situation \ref{situation-setup}.
Soit $X$ un schéma localement de type fini sur $S$.
Soit $k \in \mathbf{Z}$.
\begin{enumerate}
\item Étant donnée toute famille localement finie $\{W_j \subset X\}$
de sous-schémas fermés intègres avec $\dim_\delta(W_j) = k + 1$,
et des $f_j \in R(W_j)^*$ quelconques, nous pouvons considérer
$$
\sum (i_j)_*\text{div}(f_j) \in Z_k(X)
$$
où $i_j : W_j \to X$ est le morphisme d'inclusion.
Ceci a un sens car le morphisme
$\coprod i_j : \coprod W_j \to X$ est propre.
\item Nous disons que $\alpha \in Z_k(X)$ est {\it rationnellement équivalent à zéro}
si $\alpha$ est un cycle de la forme affichée ci-dessus.
\item Nous disons que $\alpha, \beta \in Z_k(X)$ sont
{\it rationnellement équivalents} et nous écrivons $\alpha \sim_{rat} \beta$
si $\alpha - \beta$ est rationnellement équivalent à zéro.
\item Nous définissons
$$
\CH_k(X) = Z_k(X) / \sim_{rat}
$$
comme étant le {\it groupe de Chow des $k$-cycles sur $X$}. On l'appelle parfois
le {\it groupe de Chow des $k$-cycles modulo équivalence rationnelle sur $X$}.
\end{enumerate}
\end{definition}

\noindent
Il existe beaucoup d'autres relations d'équivalence (adéquates) intéressantes.
L'équivalence rationnelle est la moins fine de toutes.

\begin{remark}
\label{remark-chow-group-pointwise}
Soit $(S, \delta)$ comme dans la Situation \ref{situation-setup}.
Soit $X$ localement de type fini sur $S$. Soit $k \in \mathbf{Z}$.
Montrons que nous avons une présentation
$$
\bigoplus\nolimits_{\delta(x) = k + 1}' K_1^M(\kappa(x))
\xrightarrow{\partial}
\bigoplus\nolimits_{\delta(x) = k}' K_0^M(\kappa(x)) \to
\CH_k(X) \to 0
$$
Nous utilisons ici les notations et conventions introduites dans la
Remarque \ref{remark-cycles-pointwise} et, de plus,
\begin{enumerate}
\item $K_1^M(\kappa(x)) = \kappa(x)^*$ est la partie de degré $1$ de
la K-théorie de Milnor du corps résiduel $\kappa(x)$ du point
$x \in X$ (voir la Remarque \ref{remark-gersten-complex-milnor}), et
\item la différentielle $\partial$ est définie comme suit :
étant donné un élément $\xi = \sum_x f_x$, nous notons $W_x = \overline{x}$
le sous-schéma fermé intègre de $X$ de point générique $x$ et nous posons
$$
\partial(\xi) = \sum (W_x \to X)_*\text{div}(f_x)
$$
dans $Z_k(X)$, ce qui a un sens puisque nous avons vu que le deuxième
terme du complexe est égal à $Z_k(X)$ d'après la
Remarque \ref{remark-cycles-pointwise}.
\end{enumerate}
Le fait que nous obtenions une présentation de $\CH_k(X)$ résulte
immédiatement de la comparaison avec la Définition \ref{definition-rational-equivalence}.
\end{remark}

\noindent
Voici un lemme très simple mais important.

\begin{lemma}
\label{lemma-restrict-to-open}
Soit $(S, \delta)$ comme dans la Situation \ref{situation-setup}.
Soit $X$ un schéma localement de type fini sur $S$.
Soit $U \subset X$ un sous-schéma ouvert, et notons
$i : Y = X \setminus U \to X$ le sous-schéma fermé réduit de $X$.
Soit $k \in \mathbf{Z}$.
Supposons $\alpha, \beta \in Z_k(X)$.
Si $\alpha|_U \sim_{rat} \beta|_U$, alors il existe un cycle
$\gamma \in Z_k(Y)$ tel que
$$
\alpha \sim_{rat} \beta + i_*\gamma.
$$
Autrement dit, la suite
$$
\xymatrix{
\CH_k(Y) \ar[r]^{i_*} & \CH_k(X) \ar[r]^{j^*} & \CH_k(U) \ar[r] & 0
}
$$
est un complexe exact de groupes abéliens.
\end{lemma}

\begin{proof}
Soit $\{W_j\}_{j \in J}$ une famille localement finie de sous-schémas fermés
intègres de $U$, de $\delta$-dimension $k + 1$, et soient $f_j \in R(W_j)^*$
des éléments tels que $(\alpha - \beta)|_U = \sum (i_j)_*\text{div}(f_j)$
comme dans la définition. Posons $W_j' \subset X$ égal
à l'adhérence de $W_j$. Supposons que $V \subset X$ soit un ouvert
quasi-compact. Alors $V \cap U$ est aussi un ouvert quasi-compact de $U$, car
$V$ est noethérien. Par conséquent, l'ensemble
$\{j \in J \mid W_j \cap V \not = \emptyset\}
= \{j \in J \mid W'_j \cap V \not = \emptyset\}$
est fini puisque $\{W_j\}$ est localement finie. Autrement dit, nous voyons que
$\{W'_j\}$ est également localement finie. Comme $R(W_j) = R(W'_j)$, nous voyons
que
$$
\alpha - \beta - \sum (i'_j)_*\text{div}(f_j)
$$
est un cycle supporté sur $Y$, et le lemme en résulte (voir le
Lemme \ref{lemma-exact-sequence-open}).
\end{proof}

\begin{lemma}
\label{lemma-exact-sequence-closed-chow}
Soit $(S, \delta)$ comme dans la Situation \ref{situation-setup}.
Soit $X$ localement de type fini sur $S$. Soient $X_1, X_2 \subset X$
des sous-schémas fermés tels que $X = X_1 \cup X_2$ ensemblistement.
Pour tout $k \in \mathbf{Z}$, la suite de groupes abéliens
$$
\xymatrix{
\CH_k(X_1 \cap X_2) \ar[r] &
\CH_k(X_1) \oplus \CH_k(X_2) \ar[r] &
\CH_k(X) \ar[r] &
0
}
$$
est exacte. Ici $X_1 \cap X_2$ est l'intersection schématique, et les
applications sont les applications d'image directe, l'une étant multipliée par $-1$.
\end{lemma}

\begin{proof}
D'après le Lemme \ref{lemma-exact-sequence-closed}, la flèche
$\CH_k(X_1) \oplus \CH_k(X_2) \to \CH_k(X)$ est surjective.
Supposons que $(\alpha_1, \alpha_2)$ soit envoyé sur zéro par cette application.
Écrivons $\alpha_1 = \sum n_{1, i}[W_{1, i}]$ et
$\alpha_2 = \sum n_{2, i}[W_{2, i}]$. Nous obtenons alors une famille
localement finie $\{W_j\}_{j \in J}$ de sous-schémas fermés intègres
de $X$ de $\delta$-dimension $k + 1$ et des $f_j \in R(W_j)^*$
tels que
$$
\sum n_{1, i}[W_{1, i}] + \sum n_{2, i}[W_{2, i}] = \sum (i_j)_*\text{div}(f_j)
$$
comme cycles sur $X$, où $i_j : W_j \to X$ est le morphisme d'inclusion.
Choisissons une décomposition en union disjointe $J = J_1 \amalg J_2$ telle que
$W_j \subset X_1$ si $j \in J_1$ et $W_j \subset X_2$ si $j \in J_2$.
(C'est possible parce que les $W_j$ sont intègres.) Nous pouvons alors écrire
l'équation ci-dessus sous la forme
$$
\sum n_{1, i}[W_{1, i}] - \sum\nolimits_{j \in J_1} (i_j)_*\text{div}(f_j) =
- \sum n_{2, i}[W_{2, i}] + \sum\nolimits_{j \in J_2} (i_j)_*\text{div}(f_j)
$$
Cette expression est donc un cycle (!) sur $X_1 \cap X_2$. Autrement dit,
l'élément $(\alpha_1, \alpha_2)$ appartient à l'image de la première flèche,
ce qui achève la démonstration.
\end{proof}

\begin{example}
\label{example-weird}
Voici un exemple « étrange ».
Supposons que $S$ soit le spectre d'un corps $k$,
avec $\delta$ comme dans l'Exemple \ref{example-field}.
Supposons que $X = C_1 \cup C_2 \cup \ldots$ soit une union infinie
de courbes $C_j \cong \mathbf{P}^1_k$ recollées
de la manière suivante : le point $\infty \in C_j$ est recollé
transversalement au point $0 \in C_{j + 1}$ pour $j = 1, 2, 3, \ldots$.
Prenons le point $0 \in C_1$. Cela donne un zéro-cycle
$[0] \in Z_0(X)$. L'« étrangeté » de cette situation est
qu'en fait $[0] \sim_{rat} 0$ ! En effet, nous pouvons choisir
la fonction rationnelle $f_j \in R(C_j)$ comme étant la fonction
qui a un zéro simple en $0$, un pôle simple en $\infty$
et aucun autre zéro ni pôle. Nous voyons alors que la somme
$\sum (i_j)_*\text{div}(f_j)$ est exactement le $0$-cycle
$[0]$. En fait, il s'avère que $\CH_0(X) = 0$ dans cet exemple.
Si vous trouvez cela trop bizarre, il vous suffit de
veiller à ce que vos espaces soient toujours quasi-compacts
(ainsi $X$ n'existe même pas pour vous).
\end{example}

\begin{remark}
\label{remark-infinite-sums-rational-equivalences}
Soit $(S, \delta)$ comme dans la Situation \ref{situation-setup}.
Soit $X$ un schéma localement de type fini sur $S$.
Supposons données des familles infinies $\alpha_i, \beta_i \in Z_k(X)$,
$i \in I$, de $k$-cycles sur $X$. Supposons que les supports
des $\alpha_i$ et des $\beta_i$ forment des familles localement finies
de fermés de $X$, de sorte que $\sum \alpha_i$
et $\sum \beta_i$ soient définies comme cycles. Supposons de plus que
$\alpha_i \sim_{rat} \beta_i$ pour tout $i$. Il n'est alors pas
évident que $\sum \alpha_i \sim_{rat} \sum \beta_i$. En effet,
le problème est que les équivalences rationnelles peuvent être
données par des
familles localement finies $\{W_{i, j}, f_{i, j} \in R(W_{i, j})^*\}_{j \in J_i}$,
mais que l'union $\{W_{i, j}\}_{i \in I, j\in J_i}$ peut ne pas
être localement finie.

\medskip\noindent
Dans de nombreux cas pratiques, on dispose d'une famille localement finie de fermés
$\{T_i\}_{i \in I}$ telle que $\alpha_i, \beta_i$
soient supportés sur $T_i$ et telle que $\alpha_i = \beta_i$
dans $\CH_k(T_i)$ ; autrement dit, les familles
$\{W_{i, j}, f_{i, j} \in R(W_{i, j})^*\}_{j \in J_i}$
sont constituées de sous-schémas $W_{i, j} \subset T_i$. Dans ce cas, il est vrai que
$\sum \alpha_i \sim_{rat} \sum \beta_i$ sur $X$, simplement parce que
la famille $\{W_{i, j}\}_{i \in I, j\in J_i}$ est alors automatiquement
localement finie.
\end{remark}






\section{Équivalence rationnelle et images directe et réciproque}
\label{section-properties-rational-equivalence}

\noindent
Dans cette section, nous montrons que l'image réciproque plate et l'image directe propre
commutent avec l'équivalence rationnelle.

\begin{lemma}
\label{lemma-prepare-flat-pullback-rational-equivalence}
Soit $(S, \delta)$ comme dans la Situation \ref{situation-setup}.
Soient $X$, $Y$ des schémas localement de type fini sur $S$.
Supposons $Y$ intègre avec $\dim_\delta(Y) = k$.
Soit $f : X \to Y$ un morphisme plat de
dimension relative $r$. Alors, pour $g \in R(Y)^*$, nous avons
$$
f^*\text{div}_Y(g) =
\sum n_j i_{j, *}\text{div}_{X_j}(g \circ f|_{X_j})
$$
comme $(k + r - 1)$-cycles sur $X$, où la somme porte sur les composantes
irréductibles $X_j$ de $X$ et $n_j$ est la multiplicité de $X_j$ dans $X$.
\end{lemma}

\begin{proof}
Soit $Z \subset X$ un sous-schéma fermé intègre de $\delta$-dimension
$k + r - 1$. Nous devons montrer que le coefficient $n$ de $[Z]$ dans
$f^*\text{div}(g)$ est égal au coefficient
$m$ de $[Z]$ dans $\sum i_{j, *} \text{div}(g \circ f|_{X_j})$.
Soit $Z'$ l'adhérence de $f(Z)$, qui est un sous-schéma fermé
intègre de $Y$. D'après le Lemme \ref{lemma-flat-inverse-image-dimension},
nous avons $\dim_\delta(Z') \geq k - 1$. Ainsi, ou bien $Z' = Y$,
ou bien $Z'$ est un diviseur premier sur $Y$. Si $Z' = Y$, alors les coefficients
$n$ et $m$ sont tous deux nuls : c'est clair pour $n$ par définition
de $f^*$, et cela résulte pour $m$ du fait que $g \circ f|_{X_j}$ est
une unité en tout point de $X_j$ envoyé sur le point générique de $Y$.
Nous pouvons donc supposer que $Z' \subset Y$ est un diviseur premier.

\medskip\noindent
Nous allons traduire en algèbre l'égalité de $n$ et $m$.
En effet, soient $\xi' \in Z'$ et $\xi \in Z$ les points génériques.
Posons $A = \mathcal{O}_{Y, \xi'}$ et $B = \mathcal{O}_{X, \xi}$.
Notons que $A$, $B$ sont noethériens, que $A \to B$ est plat et local,
que $A$ est un domaine et que $\mathfrak m_AB$ est un idéal de définition
de l'anneau local $B$. La fonction rationnelle $g$ est un élément
du corps des fractions $Q(A)$ de $A$.
Par construction, les sous-schémas fermés $X_j$
qui rencontrent $\xi$ correspondent $1$-à-$1$ aux idéaux premiers minimaux
$$
\mathfrak q_1, \ldots, \mathfrak q_s \subset B
$$
Les entiers $n_j$ sont les longueurs correspondantes
$$
n_i = \text{longueur}_{B_{\mathfrak q_i}}(B_{\mathfrak q_i})
$$
Les fonctions rationnelles $g \circ f|_{X_j}$ correspondent à l'image
$g_i \in \kappa(\mathfrak q_i)^*$ de $g \in Q(A)$.
En rassemblant tout cela, nous voyons que
$$
n = \text{ord}_A(g) \text{longueur}_B(B/\mathfrak m_AB)
$$
et que
$$
m = \sum \text{ord}_{B/\mathfrak q_i}(g_i)
\text{longueur}_{B_{\mathfrak q_i}}(B_{\mathfrak q_i})
$$
En écrivant $g = x/y$ pour des éléments non nuls $x, y \in A$, nous voyons qu'il suffit
de démontrer
$$
\text{longueur}_A(A/(x)) \text{longueur}_B(B/\mathfrak m_AB) =
\text{longueur}_B(B/xB)
$$
(l'égalité utilise Algèbre, Lemme \ref{algebra-lemma-pullback-module})
égale à
$$
\sum\nolimits_{i = 1, \ldots, s}
\text{longueur}_{B/\mathfrak q_i}(B/(x, \mathfrak q_i))
\text{longueur}_{B_{\mathfrak q_i}}(B_{\mathfrak q_i})
$$
et de même pour $y$. Comme $A \to B$ est plat, il s'ensuit que $x$
est un non-diviseur de zéro dans $B$. L'égalité souhaitée résulte alors du
Lemme \ref{lemma-additivity-divisors-restricted}.
\end{proof}

\begin{lemma}
\label{lemma-flat-pullback-rational-equivalence}
Soit $(S, \delta)$ comme dans la Situation \ref{situation-setup}.
Soient $X$, $Y$ des schémas localement de type fini sur $S$.
Soit $f : X \to Y$ un morphisme plat de dimension relative $r$.
Soient $\alpha \sim_{rat} \beta$ des $k$-cycles rationnellement équivalents sur $Y$.
Alors $f^*\alpha \sim_{rat} f^*\beta$ comme $(k + r)$-cycles sur $X$.
\end{lemma}

\begin{proof}
Que devons-nous montrer ? Supposons que l'on nous donne une famille
$$
i_j : W_j \longrightarrow Y
$$
d'immersions fermées, avec chaque $W_j$ intègre de $\delta$-dimension $k + 1$,
et des fonctions rationnelles $g_j \in R(W_j)^*$. Supposons de plus que
la famille $\{i_j(W_j)\}_{j \in J}$ soit localement finie sur $Y$.
Nous devons alors montrer que
$$
f^*(\sum i_{j, *}\text{div}(g_j)) = \sum f^*i_{j, *}\text{div}(g_j)
$$
est rationnellement équivalent à zéro sur $X$. La somme de droite
a un sens car $\{W_j\}$ est localement finie dans $X$ d'après le
Lemme \ref{lemma-inverse-image-locally-finite}.

\medskip\noindent
Considérons les produits fibrés
$$
i'_j : W'_j = W_j \times_Y X \longrightarrow X.
$$
et notons $f_j : W'_j \to W_j$ la première projection.
D'après le Lemme \ref{lemma-flat-pullback-proper-pushforward},
nous pouvons écrire la somme ci-dessus sous la forme
$$
\sum i'_{j, *}(f_j^*\text{div}(g_j))
$$
D'après le Lemme \ref{lemma-prepare-flat-pullback-rational-equivalence},
nous voyons que chaque $f_j^*\text{div}(g_j)$ est rationnellement équivalent
à zéro sur $W'_j$. Ainsi, chaque $i'_{j, *}(f_j^*\text{div}(g_j))$
est rationnellement équivalent à zéro. Il en va alors de même pour
la somme affichée d'après la discussion de la
Remarque \ref{remark-infinite-sums-rational-equivalences}.
\end{proof}

\begin{lemma}
\label{lemma-proper-pushforward-rational-equivalence}
Soit $(S, \delta)$ comme dans la Situation \ref{situation-setup}.
Soient $X$, $Y$ des schémas localement de type fini sur $S$.
Soit $p : X \to Y$ un morphisme propre.
Supposons que $\alpha, \beta \in Z_k(X)$ soient rationnellement équivalents.
Alors $p_*\alpha$ est rationnellement équivalent à $p_*\beta$.
\end{lemma}

\begin{proof}
Que devons-nous montrer ? Supposons que l'on nous donne une famille
$$
i_j : W_j \longrightarrow X
$$
d'immersions fermées, avec chaque $W_j$ intègre de $\delta$-dimension $k + 1$,
et des fonctions rationnelles $f_j \in R(W_j)^*$.
Supposons de plus que
la famille $\{i_j(W_j)\}_{j \in J}$ soit localement finie sur $X$.
Nous devons alors montrer que
$$
p_*\left(\sum i_{j, *}\text{div}(f_j)\right)
$$
est rationnellement équivalent à zéro sur $X$.

\medskip\noindent
Notons que la somme est égale à
$$
\sum p_*i_{j, *}\text{div}(f_j).
$$
Soit $W'_j \subset Y$ le sous-schéma fermé intègre qui est l'
image de $p \circ i_j$. La famille $\{W'_j\}$ est localement finie
dans $Y$ d'après le Lemme \ref{lemma-quasi-compact-locally-finite}.
Il suffit donc de montrer, pour un $j$ donné, que soit
$p_*i_{j, *}\text{div}(f_j) = 0$, soit il
est égal à $i'_{j, *}\text{div}(g_j)$ pour un certain $g_j \in R(W'_j)^*$.

\medskip\noindent
Les arguments ci-dessus nous ramènent donc au cas d'un seul
sous-schéma fermé intègre $W \subset X$ de $\delta$-dimension $k + 1$.
Soit $f \in R(W)^*$. Soit $W' = p(W)$ comme ci-dessus.
Nous obtenons un diagramme commutatif de morphismes
$$
\xymatrix{
W \ar[r]_i \ar[d]_{p'} & X \ar[d]^p \\
W' \ar[r]^{i'} & Y
}
$$
Notons que $p_*i_*\text{div}(f) = i'_*(p')_*\text{div}(f)$ d'après le
Lemme \ref{lemma-compose-pushforward}. Comme expliqué ci-dessus,
nous devons montrer que $(p')_*\text{div}(f)$
est le diviseur d'une fonction rationnelle sur $W'$, ou zéro.
Il y a trois cas à distinguer.

\medskip\noindent
Le cas $\dim_\delta(W') < k$. Dans ce cas, automatiquement,
$(p')_*\text{div}(f) = 0$ et il n'y a rien à démontrer.

\medskip\noindent
Le cas $\dim_\delta(W') = k$. Montrons que $(p')_*\text{div}(f) = 0$
dans ce cas. Soit $\eta \in W'$ le point générique.
Notons que $c : W_\eta \to \Spec(K)$
est une courbe intègre propre sur $K = \kappa(\eta)$
dont le corps des fonctions $K(W_\eta)$ s'identifie à $R(W)$.
Voici un diagramme
$$
\xymatrix{
W_\eta \ar[r] \ar[d]_c & W \ar[d]^{p'} \\
\Spec(K) \ar[r] & W'
}
$$
Notons $f_\eta \in K(W_\eta)^*$ la fonction rationnelle
correspondant à $f \in R(W)^*$.
De plus, les points fermés $\xi$ de $W_\eta$ correspondent $1 - 1$ aux
sous-schémas intègres fermés $Z = Z_\xi \subset W$ de $\delta$-dimension $k$
avec $p'(Z) = W'$. Notons que la multiplicité
de $Z_\xi$ dans $\text{div}(f)$ est égale à
$\text{ord}_{\mathcal{O}_{W_\eta, \xi}}(f_\eta)$, simplement parce que les
anneaux locaux $\mathcal{O}_{W_\eta, \xi}$ et $\mathcal{O}_{W, \xi}$
sont identifiés (comme sous-anneaux de leurs corps des fractions).
Nous voyons donc que la multiplicité de $[W']$ dans
$(p')_*\text{div}(f)$ est égale à la multiplicité de
$[\Spec(K)]$ dans $c_*\text{div}(f_\eta)$.
D'après le Lemme \ref{lemma-curve-principal-divisor}, celle-ci est nulle.

\medskip\noindent
Le cas $\dim_\delta(W') = k + 1$. Dans ce cas, le
Lemme \ref{lemma-proper-pushforward-alteration} s'applique,
et nous voyons qu'en effet $p'_*\text{div}(f) = \text{div}(g)$
pour un certain $g \in R(W')^*$, comme souhaité.
\end{proof}















\section{Équivalence rationnelle et droite projective}
\label{section-different-rational-equivalence}

\noindent
Soit $(S, \delta)$ comme dans la Situation \ref{situation-setup}.
Soit $X$ un schéma localement de type fini sur $S$.
Étant donné un sous-schéma fermé quelconque
$Z \subset X \times_S \mathbf{P}^1_S = X \times \mathbf{P}^1$,
nous notons $Z_0$, resp.\ $Z_\infty$, le sous-schéma fermé schématique
$Z_0 = \text{pr}_2^{-1}(D_0)$,
resp.\ $Z_\infty = \text{pr}_2^{-1}(D_\infty)$.
Ici $D_0$, $D_\infty$ sont comme dans (\ref{equation-zero-infty}).

\begin{lemma}
\label{lemma-rational-equivalence-family}
Soit $(S, \delta)$ comme dans la Situation \ref{situation-setup}.
Soit $X$ un schéma localement de type fini sur $S$.
Soit $W \subset X \times_S \mathbf{P}^1_S$ un sous-schéma fermé
intègre de $\delta$-dimension $k + 1$.
Supposons $W \not = W_0$ et $W \not = W_\infty$. Alors
\begin{enumerate}
\item $W_0$, $W_\infty$ sont des diviseurs de Cartier effectifs de $W$,
\item $W_0$, $W_\infty$ peuvent être considérés comme des sous-schémas fermés
de $X$ et
$$
[W_0]_k \sim_{rat} [W_\infty]_k,
$$
\item pour toute famille localement finie de
sous-schémas fermés intègres
$W_i \subset X \times_S \mathbf{P}^1_S$
de $\delta$-dimension $k + 1$ avec $W_i \not = (W_i)_0$ et
$W_i \not = (W_i)_\infty$, nous avons
$\sum ([(W_i)_0]_k - [(W_i)_\infty]_k) \sim_{rat} 0$
sur $X$, et
\item pour tout $\alpha \in Z_k(X)$ avec $\alpha \sim_{rat} 0$,
il existe une famille localement finie de
sous-schémas fermés intègres $W_i \subset X \times_S \mathbf{P}^1_S$
comme ci-dessus telle que $\alpha = \sum ([(W_i)_0]_k - [(W_i)_\infty]_k)$.
\end{enumerate}
\end{lemma}

\begin{proof}
La partie (1) résulte du
Lemme \ref{divisors-lemma-pullback-effective-Cartier-defined} de Diviseurs,
puisque, par hypothèse, le point générique
de $W$ n'est envoyé ni dans $D_0$ ni dans $D_\infty$ par la projection
$X \times_S \mathbf{P}^1_S \to \mathbf{P}^1_S$.

\medskip\noindent
Comme $X \times_S D_0 \to X$ est une immersion fermée, nous voyons que $W_0$
est isomorphe à un sous-schéma fermé de $X$. Il en va de même pour $W_\infty$.
Le morphisme $p : W \to X$ est propre comme composée de
l'immersion fermée $W \to X \times_S \mathbf{P}^1_S$ et du
morphisme propre $X \times_S \mathbf{P}^1_S \to X$. D'après le
Lemme \ref{lemma-rational-function}, nous avons
$[W_0]_k \sim_{rat} [W_\infty]_k$ comme cycles sur $W$. Ainsi la partie (2) résulte du
Lemme \ref{lemma-proper-pushforward-rational-equivalence}, puisque manifestement
$p_*[W_0]_k = [W_0]_k$ et de même pour $W_\infty$.

\medskip\noindent
Le seul contenu de l'assertion (3), compte tenu des parties (1) et (2), est que
la famille $\{(W_i)_0, (W_i)_\infty\}$ est une famille localement finie
de sous-schémas fermés de $X$. C'est clair.

\medskip\noindent
Supposons que $\alpha \sim_{rat} 0$.
Par définition, cela signifie qu'il existe des sous-schémas fermés intègres
$V_i \subset X$ de $\delta$-dimension $k + 1$ et des fonctions rationnelles
$f_i \in R(V_i)^*$ tels que la famille
$\{V_i\}_{i \in I}$ soit localement finie dans $X$ et tels que
$\alpha = \sum (V_i \to X)_*\text{div}(f_i)$.
Soit
$$
W_i \subset V_i \times_S \mathbf{P}^1_S \subset X \times_S \mathbf{P}^1_S
$$
l'adhérence du graphe de l'application rationnelle $f_i$, comme dans le
Lemme \ref{lemma-rational-function}.
Alors $(V_i \to X)_*\text{div}(f_i)$
est égal à $[(W_i)_0]_k - [(W_i)_\infty]_k$ d'après ce même lemme.
Le résultat est donc clair.
\end{proof}

\begin{lemma}
\label{lemma-closed-subscheme-cross-p1}
Soit $(S, \delta)$ comme dans la Situation \ref{situation-setup}.
Soit $X$ un schéma localement de type fini sur $S$.
Soit $Z$ un sous-schéma fermé de $X \times \mathbf{P}^1$.
Supposons
\begin{enumerate}
\item $\dim_\delta(Z) \leq k + 1$,
\item $\dim_\delta(Z_0) \leq k$, $\dim_\delta(Z_\infty) \leq k$, et
\item pour tout point immergé $\xi$ (Diviseurs, Définition
\ref{divisors-definition-embedded}) de $Z$, soit
$\xi \not \in Z_0 \cup Z_\infty$, soit $\delta(\xi) < k$.
\end{enumerate}
Alors $[Z_0]_k \sim_{rat} [Z_\infty]_k$ comme $k$-cycles sur $X$.
\end{lemma}

\begin{proof}
Soit $\{W_i\}_{i \in I}$ la famille des composantes
irréductibles de $Z$ qui sont de $\delta$-dimension $k + 1$.
Écrivons
$$
[Z]_{k + 1} = \sum n_i[W_i]
$$
avec $n_i > 0$, conformément à la définition. Notons que $\{W_i\}$
est une famille localement finie de fermés de
$X \times_S \mathbf{P}^1_S$ d'après le
Lemme \ref{divisors-lemma-components-locally-finite} de Diviseurs.
Nous affirmons que
$$
[Z_0]_k = \sum n_i[(W_i)_0]_k
$$
et de même pour $[Z_\infty]_k$. Si nous le démontrons, alors le lemme
résulte du Lemme \ref{lemma-rational-equivalence-family}.

\medskip\noindent
Soit $Z' \subset X$ un sous-schéma fermé intègre de $\delta$-dimension $k$.
Pour démontrer l'égalité ci-dessus, il suffit de montrer que le coefficient $n$
de $[Z']$ dans $[Z_0]_k$ est égal au coefficient $m$ de
$[Z']$ dans $\sum n_i[(W_i)_0]_k$. Soit $\xi' \in Z'$ le point générique.
Posons $\xi = (\xi', 0) \in  X \times_S \mathbf{P}^1_S$.
Considérons l'anneau local $A = \mathcal{O}_{X \times_S \mathbf{P}^1_S, \xi}$.
Soit $I \subset A$ l'idéal définissant $Z$, autrement dit tel que
$A/I = \mathcal{O}_{Z, \xi}$. Soit $t \in A$ l'élément définissant
$X \times_S D_0$ (c'est-à-dire le relevé de la coordonnée de $\mathbf{P}^1$ en zéro).
Par notre choix de $\xi' \in Z'$, nous avons $\delta(\xi) = k$,
et donc $\dim(A/I) = 1$. Puisque $\xi$ n'est pas un point immergé d'après
l'hypothèse (3), nous voyons que $A/I$ est de Cohen-Macaulay. Puisque $\dim_\delta(Z_0)
= k$, nous voyons que $\dim(A/(t, I)) = 0$, ce qui implique que $t$
est un non-diviseur de zéro dans $A/I$. Enfin, les sous-schémas fermés irréductibles
$W_i$ passant par $\xi$ correspondent aux idéaux premiers minimaux
$I \subset \mathfrak q_i$ au-dessus de $I$. Les multiplicités $n_i$ correspondent
aux longueurs $\text{longueur}_{A_{\mathfrak q_i}}(A/I)_{\mathfrak q_i}$.
Nous voyons donc que
$$
n = \text{longueur}_A(A/(t, I))
$$
et
$$
m = \sum
\text{longueur}_A(A/(t, \mathfrak q_i))
\text{longueur}_{A_{\mathfrak q_i}}(A/I)_{\mathfrak q_i}
$$
Le résultat résulte donc du
Lemme \ref{lemma-additivity-divisors-restricted}.
\end{proof}

\begin{lemma}
\label{lemma-coherent-sheaf-cross-p1}
Soit $(S, \delta)$ comme dans la Situation \ref{situation-setup}.
Soit $X$ un schéma localement de type fini sur $S$.
Soit $\mathcal{F}$ un faisceau cohérent sur $X \times \mathbf{P}^1$.
Soient $i_0, i_\infty : X \to X \times \mathbf{P}^1$ les immersions fermées
telles que $i_t(x) = (x, t)$. Notons $\mathcal{F}_0 = i_0^*\mathcal{F}$ et
$\mathcal{F}_\infty = i_\infty^*\mathcal{F}$.
Supposons
\begin{enumerate}
\item $\dim_\delta(\text{Supp}(\mathcal{F})) \leq k + 1$,
\item $\dim_\delta(\text{Supp}(\mathcal{F}_0)) \leq k$,
$\dim_\delta(\text{Supp}(\mathcal{F}_\infty)) \leq k$, et
\item pour tout point associé immergé $\xi$ de $\mathcal{F}$, soit
$\xi \not \in (X \times \mathbf{P}^1)_0 \cup (X \times \mathbf{P}^1)_\infty$,
soit $\delta(\xi) < k$.
\end{enumerate}
Alors $[\mathcal{F}_0]_k \sim_{rat} [\mathcal{F}_\infty]_k$ comme $k$-cycles sur $X$.
\end{lemma}

\begin{proof}
Soit $\{W_i\}_{i \in I}$ la famille des composantes
irréductibles de $\text{Supp}(\mathcal{F})$
qui sont de $\delta$-dimension $k + 1$.
Écrivons
$$
[\mathcal{F}]_{k + 1} = \sum n_i[W_i]
$$
avec $n_i > 0$, conformément à la définition. Notons que $\{W_i\}$
est une famille localement finie de fermés de
$X \times_S \mathbf{P}^1_S$ d'après le Lemme \ref{lemma-length-finite}.
Nous affirmons que
$$
[\mathcal{F}_0]_k = \sum n_i[(W_i)_0]_k
$$
et de même pour $[\mathcal{F}_\infty]_k$. Si nous le démontrons, alors le lemme
résulte du Lemme \ref{lemma-rational-equivalence-family}.

\medskip\noindent
Soit $Z' \subset X$ un sous-schéma fermé intègre de $\delta$-dimension $k$.
Pour démontrer l'égalité ci-dessus, il suffit de montrer que le coefficient $n$
de $[Z']$ dans $[\mathcal{F}_0]_k$ est égal au coefficient $m$ de
$[Z']$ dans $\sum n_i[(W_i)_0]_k$. Soit $\xi' \in Z'$ le point générique.
Posons $\xi = (\xi', 0) \in  X \times_S \mathbf{P}^1_S$.
Considérons l'anneau local $A = \mathcal{O}_{X \times_S \mathbf{P}^1_S, \xi}$.
Considérons $M = \mathcal{F}_\xi$ comme un $A$-module.
Soit $t \in A$ l'élément définissant $X \times_S D_0$
(c'est-à-dire le relevé de la coordonnée de $\mathbf{P}^1$ en zéro).
Par notre choix de $\xi' \in Z'$, nous avons $\delta(\xi) = k$,
et donc $\dim(\text{Supp}(M)) = 1$. Puisque $\xi$ n'est pas un point associé
de $\mathcal{F}$ d'après l'hypothèse (3), nous voyons que $M$ est un module de Cohen-Macaulay.
Puisque $\dim_\delta(\text{Supp}(\mathcal{F}_0)) = k$,
nous voyons que $\dim(\text{Supp}(M/tM)) = 0$, ce qui implique que $t$
est un non-diviseur de zéro dans $M$. Enfin, les sous-schémas fermés irréductibles
$W_i$ passant par $\xi$ correspondent aux idéaux premiers minimaux
$\mathfrak q_i$ de $\text{Ass}(M)$. Les multiplicités $n_i$ correspondent
aux longueurs $\text{longueur}_{A_{\mathfrak q_i}}M_{\mathfrak q_i}$.
Nous voyons donc que
$$
n = \text{longueur}_A(M/tM)
$$
et
$$
m = \sum
\text{longueur}_A(A/(t, \mathfrak q_i)A)
\text{longueur}_{A_{\mathfrak q_i}}M_{\mathfrak q_i}
$$
Le résultat résulte donc du
Lemme \ref{lemma-additivity-divisors-restricted}.
\end{proof}







\section{Groupes de Chow et enveloppes}
\label{section-envelopes}

\noindent
Voici la définition.

\begin{definition}
\label{definition-envelope}
\begin{reference}
\cite[Definition 18.3]{F}
\end{reference}
Soit $X$ un schéma. Une {\it enveloppe} est un morphisme propre $f : Y \to X$
qui est complètement décomposé
(Plus sur les morphismes, Définition \ref{more-morphisms-definition-cd-morphism}).
\end{definition}

\noindent
La suite exacte du Lemme \ref{lemma-envelope}
est la motivation principale de la définition.

\begin{lemma}
\label{lemma-composition-envelope}
Soit $(S, \delta)$ comme dans la Situation \ref{situation-setup}.
Soit $X$ un schéma localement de type fini sur $S$.
Si $f : Y \to X$ et $g : Z \to Y$ sont des enveloppes, alors
$f \circ g$ est une enveloppe.
\end{lemma}

\begin{proof}
Ceci résulte du Lemme \ref{morphisms-lemma-composition-proper} de Morphismes
et du Lemme \ref{more-morphisms-lemma-composition-cd} de Plus sur les morphismes.
\end{proof}

\begin{lemma}
\label{lemma-base-change-envelope}
Soit $(S, \delta)$ comme dans la Situation \ref{situation-setup}.
Soit $X' \to X$ un morphisme de schémas localement de type fini sur $S$.
Si $f : Y \to X$ est une enveloppe, alors le changement de base $f' : Y' \to X'$
de $f$ est aussi une enveloppe.
\end{lemma}

\begin{proof}
Ceci résulte du Lemme \ref{morphisms-lemma-base-change-proper} de Morphismes
et du Lemme \ref{more-morphisms-lemma-base-change-cd} de Plus sur les morphismes.
\end{proof}

\begin{lemma}
\label{lemma-envelope}
Soit $(S, \delta)$ comme dans la Situation \ref{situation-setup}.
Soit $X$ un schéma localement de type fini sur $S$.
Soit $f : Y \to X$ une enveloppe. Alors
nous avons une suite exacte
$$
\CH_k(Y \times_X Y) \xrightarrow{p_* - q_*}
\CH_k(Y) \xrightarrow{f_*}
\CH_k(X) \to 0
$$
pour tout $k \in \mathbf{Z}$. Ici $p, q : Y \times_X Y \to Y$ sont
les projections.
\end{lemma}

\begin{proof}
Puisque $f$ est une enveloppe, $f$ est propre, et l'image directe des
cycles et des classes de cycles est donc définie ; voir les
sections \ref{section-proper-pushforward} et \ref{section-push-pull}.
De même, les morphismes $p$ et $q$ sont propres comme changements de base de $f$.
La composée des flèches est nulle, car
$f_* \circ p_* = (p \circ f)_* = (q \circ f)_* = f_* \circ q_*$ ; voir le
Lemme \ref{lemma-compose-pushforward}.

\medskip\noindent
Montrons que $f_* : Z_k(Y) \to Z_k(X)$ est surjectif.
En effet, supposons que nous ayons $\alpha = \sum n_i[Z_i] \in Z_k(X)$,
où les $Z_i \subset X$ forment une famille localement finie de sous-schémas
fermés intègres. Soit $x_i \in Z_i$ le point générique.
Puisque $f$ est une enveloppe et donc complètement décomposé,
il existe un point $y_i \in Y$ avec $f(y_i) = x_i$
et avec $\kappa(y_i)/\kappa(x_i)$ triviale. Soit $W_i \subset Y$
le sous-schéma fermé intègre de point générique $y_i$.
Comme $f$ est fermé, nous voyons que $f(W_i) = Z_i$.
Il en résulte que la famille de sous-schémas fermés $W_i$ est localement finie
sur $Y$. Comme $\kappa(y_i)/\kappa(x_i)$ est triviale, nous voyons
que $\dim_\delta(W_i) = \dim_\delta(Z_i) = k$. Ainsi
$\beta = \sum n_i[W_i]$ appartient à $Z_k(Y)$. Enfin, puisque
$\kappa(y_i)/\kappa(x_i)$ est triviale, le degré du morphisme dominant
$f|_{W_i} : W_i \to Z_i$ est $1$, et nous concluons
que $f_*\beta = \alpha$.

\medskip\noindent
Puisque $f_* : Z_k(Y) \to Z_k(X)$ est surjectif, l'application
$f_* : \CH_k(Y) \to \CH_k(X)$ est a fortiori surjective.

\medskip\noindent
Soit $\beta \in Z_k(Y)$ un élément tel que $f_*\beta$ soit nul dans
$\CH_k(X)$. Cela signifie que nous pouvons trouver une famille localement finie de
sous-schémas fermés intègres $Z_j \subset X$ avec $\dim_\delta(Z_j) = k + 1$
et des $f_j \in R(Z_j)^*$ tels que
$$
f_*\beta = \sum (Z_j \to X)_*\text{div}(f_j)
$$
comme cycles, où $i_j : Z_j \to X$ est l'immersion fermée donnée.
En raisonnant exactement comme ci-dessus, nous pouvons trouver une famille
localement finie de sous-schémas fermés intègres $W_j \subset Y$
avec $f(W_j) = Z_j$ et telle que $W_j \to Z_j$ soit birationnel, c'est-à-dire
induise un isomorphisme $R(Z_j) = R(W_j)$. Notons $g_j \in R(W_j)^*$
l'élément correspondant à $f_j$. Observons que $W_j \to Z_j$
est propre et que $(W_j \to Z_j)_*\text{div}(g_j) = \text{div}(f_j)$
comme cycles sur $Z_j$. Il s'ensuit que, si nous remplaçons
$\beta$ par le cycle rationnellement équivalent
$$
\beta' = \beta - \sum (W_j \to Y)_*\text{div}(g_j)
$$
nous obtenons $f_*\beta' = 0$.
(On utilise ici le Lemme \ref{lemma-compose-pushforward}.)
Ainsi, pour achever la démonstration
du lemme, il suffit de démontrer l'affirmation du paragraphe suivant.

\medskip\noindent
Affirmation : si $\beta \in Z_k(Y)$ et $f_*\beta = 0$, alors
$\beta = \delta + p_*\gamma - q_*\gamma$ dans $Z_k(Y)$ pour un certain
$\gamma \in Z_k(Y \times_X Y)$. En effet, écrivons $\beta = \sum_{j \in J} n_j[W_j]$,
où $\{W_j\}_{j \in J}$ est une famille localement finie de sous-schémas fermés
intègres de $Y$ avec $\dim_\delta(W_j) = k$.
Fixons un sous-schéma fermé intègre $Z \subset X$. Considérons le sous-ensemble
$J_Z = \{j \in J : f(W_j) = Z\}$. C'est un ensemble fini. Il y a trois
cas :
\begin{enumerate}
\item $J_Z = \emptyset$. Dans ce cas, nous posons $\gamma_Z = 0$.
\item $J_Z \not = \emptyset$ et $\dim_\delta(Z) = k$.
La condition $f_*\beta = 0$ implique, en considérant le
coefficient de $Z$, que $\sum_{j \in J_Z} n_j\deg(W_j/Z) = 0$.
Dans ce cas, nous choisissons un sous-schéma fermé intègre $W \subset Y$
qui s'envoie birationnellement sur $Z$ (voir ci-dessus). En considérant les points
génériques, nous voyons que $W_j \times_Z W$ possède une unique composante
irréductible $W'_j \subset W_j \times_Z W \subset Y \times_X Y$
qui s'envoie birationnellement sur $W_j$. Alors $W'_j \to W$ est dominant
et $\deg(W'_j/W) = \deg(W_j/W)$. Ainsi, si nous posons
$\gamma_Z = \sum_{j \in J_Z} n_j[W'_j]$,
alors nous voyons que
$p_*\gamma_Z = \sum_{j \in J_Z} n_j[W_j]$ et
$q_*\gamma_Z = \sum_{j \in J_Z} n_j\deg(W'_j/W)[W] = 0$.
\item $J_Z \not = \emptyset$ et $\dim_\delta(Z) < k$.
Dans ce cas, nous choisissons un sous-schéma fermé intègre $W \subset Y$
qui s'envoie birationnellement sur $Z$ (voir ci-dessus). En considérant les points
génériques, nous voyons que $W_j \times_Z W$ possède une unique composante
irréductible $W'_j \subset W_j \times_Z W \subset Y \times_X Y$
qui s'envoie birationnellement sur $W_j$. Alors $W'_j \to W$ est dominant
et $k = \dim_\delta(W'_j) > \dim_\delta(W) = \dim_\delta(Z)$.
Ainsi, si nous posons $\gamma_Z = \sum_{j \in J_Z} n_j[W'_j]$,
alors nous voyons que
$p_*\gamma_Z = \sum_{j \in J_Z} n_j[W_j]$ et
$q_*\gamma_Z = 0$.
\end{enumerate}
Puisque la famille de sous-schémas fermés intègres $\{f(W_j)\}$
est localement finie sur $X$
(Lemme \ref{lemma-quasi-compact-locally-finite}),
nous voyons que le $k$-cycle
$$
\gamma = \sum\nolimits_{Z \subset X\text{ fermé intègre}} \gamma_Z
$$
sur $Y \times_X Y$ est bien défini. Nos calculs ci-dessus montrent que
$p_*\gamma_Z = \beta$ et $q_*\gamma_Z = 0$, ce qui implique
ce que nous voulions démontrer.
\end{proof}













\section{Groupes de Chow et groupes K}
\label{section-chow-and-K}

\noindent
Dans cette section, nous allons comparer le $K_0$ de la
catégorie des faisceaux cohérents aux groupes de Chow.

\medskip\noindent
Soit $(S, \delta)$ comme dans la Situation \ref{situation-setup}.
Soit $X$ un schéma localement de type fini sur $S$.
Nous notons $\textit{Coh}(X) = \textit{Coh}(\mathcal{O}_X)$
la catégorie des faisceaux cohérents sur $X$.
C'est une catégorie abélienne ; voir le
Lemme \ref{coherent-lemma-coherent-abelian-Noetherian} de Cohomologie des schémas.
Pour tout $k \in \mathbf{Z}$, nous notons $\textit{Coh}_{\leq k}(X)$
la sous-catégorie pleine de $\textit{Coh}(X)$
formée des faisceaux cohérents $\mathcal{F}$
tels que $\dim_\delta(\text{Supp}(\mathcal{F})) \leq k$.

\begin{lemma}
\label{lemma-Serre-subcategories}
Soit $(S, \delta)$ comme dans la Situation \ref{situation-setup}.
Soit $X$ un schéma localement de type fini sur $S$.
Les catégories $\textit{Coh}_{\leq k}(X)$ sont des sous-catégories de Serre
de la catégorie abélienne $\textit{Coh}(X)$.
\end{lemma}

\begin{proof}
La définition d'une sous-catégorie de Serre est la
Définition \ref{homology-definition-serre-subcategory} d'Homologie.
La démonstration du lemme est immédiate et omise.
\end{proof}

\begin{lemma}
\label{lemma-cycles-k-group}
Soit $(S, \delta)$ comme dans la Situation \ref{situation-setup}.
Soit $X$ un schéma localement de type fini sur $S$.
Les applications
$$
Z_k(X)
\longrightarrow
K_0(\textit{Coh}_{\leq k}(X)/\textit{Coh}_{\leq k - 1}(X)),
\quad
\sum n_Z[Z] \mapsto
\left[\bigoplus\nolimits_{n_Z > 0} \mathcal{O}_Z^{\oplus n_Z}\right]
-
\left[\bigoplus\nolimits_{n_Z < 0} \mathcal{O}_Z^{\oplus -n_Z}\right]
$$
et
$$
K_0(\textit{Coh}_{\leq k}(X)/\textit{Coh}_{\leq k - 1}(X))
\longrightarrow
Z_k(X),\quad
\mathcal{F} \longmapsto [\mathcal{F}]_k
$$
sont des isomorphismes réciproques l'un de l'autre.
\end{lemma}

\begin{proof}
Notons que, si $\sum n_Z[Z]$ appartient à $Z_k(X)$, alors
les sommes directes
$\bigoplus\nolimits_{n_Z > 0} \mathcal{O}_Z^{\oplus n_Z}$ et
$\bigoplus\nolimits_{n_Z < 0} \mathcal{O}_Z^{\oplus -n_Z}$
sont des faisceaux cohérents sur $X$, puisque la famille $\{Z \mid n_Z > 0\}$
est localement finie sur $X$.
L'application $\mathcal{F} \to [\mathcal{F}]_k$ est additive
sur $\textit{Coh}_{\leq k}(X)$ ; voir le
Lemme \ref{lemma-additivity-sheaf-cycle}. Et $[\mathcal{F}]_k = 0$
si $\mathcal{F} \in \textit{Coh}_{\leq k - 1}(X)$. D'après la partie (1)
du Lemme \ref{homology-lemma-serre-subcategory-K-groups} d'Homologie,
cela implique que la deuxième application est elle aussi bien définie.
Il est clair que la composée de la première application avec la deuxième
est l'identité.

\medskip\noindent
Réciproquement, supposons que nous partions d'un faisceau cohérent $\mathcal{F}$
sur $X$. Écrivons $[\mathcal{F}]_k = \sum_{i \in I} n_i[Z_i]$
avec $n_i > 0$ et $Z_i \subset X$, $i \in I$,
des sous-schémas fermés intègres deux à deux distincts de $\delta$-dimension $k$.
Nous devons montrer que
$$
[\mathcal{F}] = [\bigoplus\nolimits_{i \in I} \mathcal{O}_{Z_i}^{\oplus n_i}]
$$
dans $K_0(\textit{Coh}_{\leq k}(X)/\textit{Coh}_{\leq k - 1}(X))$.
Notons $\xi_i \in Z_i$ le point générique.
Si nous posons
$$
\mathcal{F}' = \Ker(\mathcal{F} \to \bigoplus \xi_{i, *}\mathcal{F}_{\xi_i})
$$
alors $\mathcal{F}'$ est le sous-module cohérent maximal de $\mathcal{F}$
dont le support est de dimension $\leq k - 1$. En particulier, $\mathcal{F}$
et $\mathcal{F}/\mathcal{F}'$ ont la même classe dans
$K_0(\textit{Coh}_{\leq k}(X)/\textit{Coh}_{\leq k - 1}(X))$.
Ainsi, après avoir remplacé $\mathcal{F}$ par $\mathcal{F}/\mathcal{F}'$,
nous pouvons supposer, et supposons, que le noyau $\mathcal{F}'$ affiché
ci-dessus est nul.

\medskip\noindent
Pour tout $i \in I$, choisissons une filtration
$$
\mathcal{F}_{\xi_i} = \mathcal{F}_i^0 \supset \mathcal{F}_i^1 \supset
\ldots \supset \mathcal{F}_i^{n_i} = 0
$$
telle que les quotients successifs soient de dimension $1$ sur le corps
résiduel en $\xi_i$. C'est possible puisque la longueur de $\mathcal{F}_{\xi_i}$
sur $\mathcal{O}_{X, \xi_i}$ est $n_i$.
Pour $p > n_i$, posons $\mathcal{F}_i^p = 0$. Pour $p \geq 0$, nous notons
$$
\mathcal{F}^p =
\Ker\left(\mathcal{F} \longrightarrow \bigoplus
\xi_{i, *}(\mathcal{F}_{\xi_i}/\mathcal{F}_i^p)\right)
$$
Alors $\mathcal{F}^p$ est cohérent, $\mathcal{F}^0 = \mathcal{F}$, et
$\mathcal{F}^p/\mathcal{F}^{p + 1}$ est isomorphe à un
$\mathcal{O}_{Z_i}$-module libre de rang $1$ (si $n_i > p$), ou à $0$
(si $n_i \leq p$), dans un voisinage ouvert de $\xi_i$. De plus,
$\mathcal{F}' = \bigcap \mathcal{F}^p = 0$. Puisque tout ouvert quasi-compact
$U \subset X$ ne contient qu'un nombre fini de $\xi_i$,
nous concluons que $\mathcal{F}^p|_U$ est nul pour $p \gg 0$.
Ainsi $\bigoplus_{p \geq 0} \mathcal{F}^p$ est un
$\mathcal{O}_X$-module cohérent. Considérons les suites exactes courtes
$$
0 \to
\bigoplus\nolimits_{p > 0} \mathcal{F}^p \to
\bigoplus\nolimits_{p \geq 0} \mathcal{F}^p \to
\bigoplus\nolimits_{p > 0} \mathcal{F}^p/\mathcal{F}^{p + 1} \to 0
$$
et
$$
0 \to
\bigoplus\nolimits_{p > 0} \mathcal{F}^p \to
\bigoplus\nolimits_{p \geq 0} \mathcal{F}^p \to
\mathcal{F} \to 0
$$
de $\mathcal{O}_X$-modules cohérents. Cela montre déjà que
$$
[\mathcal{F}] = [\bigoplus \mathcal{F}^p/\mathcal{F}^{p + 1}]
$$
dans $K_0(\textit{Coh}_{\leq k}(X)/\textit{Coh}_{\leq k - 1}(X))$.
Ensuite, pour tout $p \geq 0$ et tout $i \in I$ tel que $n_i > p$,
nous choisissons un faisceau d'idéaux non nul $\mathcal{I}_{i, p} \subset \mathcal{O}_{Z_i}$
et une application $\mathcal{I}_{i, p} \to \mathcal{F}^p/\mathcal{F}^{p + 1}$ sur $X$
qui est un isomorphisme sur le voisinage ouvert de $\xi_i$
mentionné ci-dessus. C'est possible d'après le
Lemme \ref{coherent-lemma-extend-coherent} de Cohomologie des schémas.
Nous considérons alors la suite exacte courte
$$
0 \to
\bigoplus\nolimits_{p \geq 0, i \in I, n_i > p} \mathcal{I}_{i, p}
\to
\bigoplus \mathcal{F}^p/\mathcal{F}^{p + 1} \to
\mathcal{Q} \to 0
$$
et la suite exacte courte
$$
0 \to
\bigoplus\nolimits_{p \geq 0, i \in I, n_i > p} \mathcal{I}_{i, p}
\to
\bigoplus\nolimits_{p \geq 0, i \in I, n_i > p} \mathcal{O}_{Z_i}
\to
\mathcal{Q}' \to 0
$$
Observons que $\mathcal{Q}$ et $\mathcal{Q}'$ sont tous deux nuls dans un voisinage
des points $\xi_i$ et qu'ils sont supportés sur $\bigcup Z_i$.
Ainsi $\mathcal{Q}$ et $\mathcal{Q}'$ appartiennent à
$\textit{Coh}_{\leq k - 1}(X)$.
Puisque
$$
\bigoplus\nolimits_{i \in I} \mathcal{O}_{Z_i}^{\oplus n_i} \cong
\bigoplus\nolimits_{p \geq 0, i \in I, n_i > p} \mathcal{O}_{Z_i}
$$
ceci achève la démonstration.
\end{proof}

\begin{lemma}
\label{lemma-finite-cycles-k-group}
Soit $\pi : X \to Y$ un morphisme fini de schémas localement de type fini
sur $(S, \delta)$ comme dans la Situation \ref{situation-setup}. Alors
$\pi_* : \textit{Coh}(X) \to \textit{Coh}(Y)$ est un foncteur exact
qui envoie $\textit{Coh}_{\leq k}(X)$ dans $\textit{Coh}_{\leq k}(Y)$
et induit des homomorphismes sur les $K_0$ de ces catégories et de
leurs quotients. Les applications du Lemme \ref{lemma-cycles-k-group}
s'insèrent dans un diagramme commutatif
$$
\xymatrix{
Z_k(X) \ar[d]^{\pi_*} \ar[r] &
K_0(\textit{Coh}_{\leq k}(X)/\textit{Coh}_{\leq k - 1}(X))
\ar[d]^{\pi_*} \ar[r] &
Z_k(X) \ar[d]^{\pi_*} \\
Z_k(Y) \ar[r] &
K_0(\textit{Coh}_{\leq k}(Y)/\textit{Coh}_{\leq k - 1}(Y)) \ar[r] &
Z_k(Y)
}
$$
\end{lemma}

\begin{proof}
Un morphisme fini est affine ; l'image directe des modules quasi-cohérents
par $\pi$ est donc un foncteur exact d'après le
Lemme \ref{coherent-lemma-relative-affine-vanishing} de Cohomologie des schémas.
Un morphisme fini est propre ; $\pi_*$ envoie donc les faisceaux cohérents
sur des faisceaux cohérents ; voir la Proposition
\ref{coherent-proposition-proper-pushforward-coherent} de Cohomologie des schémas.
L'assertion sur les dimensions des supports est claire.
La commutativité à droite résulte immédiatement du
Lemme \ref{lemma-cycle-push-sheaf}.
Puisque les flèches horizontales sont des bijections, nous constatons que
nous avons aussi la commutativité à gauche.
\end{proof}

\begin{lemma}
\label{lemma-from-chow-to-K}
Soit $X$ un schéma localement de type fini sur $(S, \delta)$
comme dans la Situation \ref{situation-setup}. Il existe une application canonique
$$
\CH_k(X)
\longrightarrow
K_0(\textit{Coh}_{\leq k + 1}(X)/\textit{Coh}_{\leq k - 1}(X))
$$
induite par l'application
$Z_k(X) \to K_0(\textit{Coh}_{\leq k}(X)/\textit{Coh}_{\leq k - 1}(X))$
du Lemme \ref{lemma-cycles-k-group}.
\end{lemma}

\begin{proof}
Nous devons montrer qu'un élément $\alpha$ de $Z_k(X)$ qui est rationnellement
équivalent à zéro est envoyé sur zéro dans
$K_0(\textit{Coh}_{\leq k + 1}(X)/\textit{Coh}_{\leq k - 1}(X))$.
Écrivons $\alpha = \sum (i_j)_*\text{div}(f_j)$ comme dans la
Définition \ref{definition-rational-equivalence}.
Observons que
$$
\pi = \coprod i_j : W = \coprod W_j \longrightarrow X
$$
est un morphisme fini, car chaque $i_j : W_j \to X$ est une immersion fermée
et la famille des $W_j$ est localement finie dans $X$. Nous pouvons donc utiliser le
Lemme \ref{lemma-finite-cycles-k-group} pour nous ramener au cas de $W$.
Puisque $W$ est une union disjointe de schémas intègres, nous nous ramenons
au cas étudié dans le paragraphe suivant.

\medskip\noindent
Supposons $X$ intègre de $\delta$-dimension $k + 1$.
Soit $f$ une fonction rationnelle non nulle sur $X$.
Soit $\alpha = \text{div}(f)$. Nous devons montrer que
$\alpha$ est envoyé sur zéro dans
$K_0(\textit{Coh}_{\leq k + 1}(X)/\textit{Coh}_{\leq k - 1}(X))$.
Soit $\mathcal{I} \subset \mathcal{O}_X$ l'idéal des dénominateurs
de $f$ ; voir Diviseurs, Définition
\ref{divisors-definition-regular-meromorphic-ideal-denominators}.
Nous avons alors des suites exactes courtes
$$
0 \to \mathcal{I} \to \mathcal{O}_X \to \mathcal{O}_X/\mathcal{I} \to 0
$$
et
$$
0 \to \mathcal{I} \xrightarrow{f} \mathcal{O}_X \to
\mathcal{O}_X/f\mathcal{I} \to 0
$$
Voir Diviseurs, Lemme
\ref{divisors-lemma-regular-meromorphic-ideal-denominators}.
Nous affirmons que
$$
[\mathcal{O}_X/\mathcal{I}]_k - [\mathcal{O}_X/f\mathcal{I}]_k =
\text{div}(f)
$$
Cette affirmation implique que l'élément $\alpha = \text{div}(f)$ est représenté par
$[\mathcal{O}_X/\mathcal{I}] - [\mathcal{O}_X/f\mathcal{I}]$
dans $K_0(\textit{Coh}_{\leq k}(X)/\textit{Coh}_{\leq k - 1}(X))$.
Les suites exactes courtes montrent alors que cet élément s'envoie sur
zéro dans $K_0(\textit{Coh}_{\leq k + 1}(X)/\textit{Coh}_{\leq k - 1}(X))$.

\medskip\noindent
Pour démontrer l'affirmation, soit $Z \subset X$ un sous-schéma fermé intègre
de $\delta$-dimension $k$, et soit $\xi \in Z$ son point générique.
Alors $I = \mathcal{I}_\xi \subset A = \mathcal{O}_{X, \xi}$
est un idéal tel que $fI \subset A$. Or le coefficient de
$[Z]$ dans $\text{div}(f)$ est $\text{ord}_A(f)$. (Bien sûr, comme d'habitude,
nous identifions le corps des fonctions de $X$ au corps des fractions de $A$.)
D'autre part, le coefficient de $[Z]$ dans
$[\mathcal{O}_X/\mathcal{I}] - [\mathcal{O}_X/f\mathcal{I}]$
est
$$
\text{longueur}_A(A/I) - \text{longueur}_A(A/fI)
$$
À l'aide de la fonction distance de la
Définition \ref{algebra-definition-distance} d'Algèbre,
nous pouvons récrire ceci sous la forme
$$
d(A, I) - d(A, fI) = d(I, fI) = \text{ord}_A(f)
$$
Les égalités sont vraies d'après les Lemmes
\ref{algebra-lemma-properties-distance-function} et
\ref{algebra-lemma-order-vanishing-determinant} d'Algèbre.
(L'utilisation de ces lemmes n'est pas nécessaire, mais commode.)
\end{proof}

\begin{remark}
\label{remark-good-cases-K-A}
Soit $(S, \delta)$ comme dans la Situation \ref{situation-setup}.
Soit $X$ un schéma localement de type fini sur $S$.
Nous verrons plus loin (dans le Lemme \ref{lemma-cycles-rational-equivalence-K-group})
que l'application
$$
\CH_k(X)
\longrightarrow
K_0(\textit{Coh}_{k + 1}(X)/\textit{Coh}_{\leq k - 1}(X))
$$
du Lemme \ref{lemma-from-chow-to-K} est injective.
En la composant avec l'application canonique
$$
K_0(\textit{Coh}_{k + 1}(X)/\textit{Coh}_{\leq k - 1}(X))
\longrightarrow
K_0(\textit{Coh}(X)/\textit{Coh}_{\leq k - 1}(X))
$$
nous obtenons une application canonique
$$
\CH_k(X)
\longrightarrow
K_0(\textit{Coh}(X)/\textit{Coh}_{\leq k - 1}(X)).
$$
Nous n'avons trouvé dans la littérature aucun énoncé ni aucune conjecture
indiquant si cette application devrait être injective ou non.
Il semble raisonnable de s'attendre à ce que le noyau de cette application soit de torsion.
Nous reviendrons sur cette question (insérer une référence ultérieure).
\end{remark}

\begin{lemma}
\label{lemma-K-coherent-supported-on-closed}
Soit $X$ un schéma localement noethérien. Soit $Z \subset X$ un sous-schéma
fermé. Notons $\textit{Coh}_Z(X) \subset \textit{Coh}(X)$
la sous-catégorie de Serre des $\mathcal{O}_X$-modules cohérents dont le
support ensembliste est contenu dans $Z$. Alors le foncteur exact d'inclusion
$\textit{Coh}(Z) \to \textit{Coh}_Z(X)$ induit
un isomorphisme
$$
K'_0(Z) = K_0(\textit{Coh}(Z)) \longrightarrow K_0(\textit{Coh}_Z(X))
$$
\end{lemma}

\begin{proof}
Soit $\mathcal{F}$ un objet de $\textit{Coh}_Z(X)$.
Soit $\mathcal{I} \subset \mathcal{O}_X$ le faisceau quasi-cohérent
d'idéaux de $Z$. Considérons la filtration décroissante
$$
\ldots \subset
\mathcal{F}^p = \mathcal{I}^p \mathcal{F} \subset
\mathcal{F}^{p - 1} \subset \ldots \subset \mathcal{F}^0 = \mathcal{F}
$$
Exactement comme dans la démonstration du Lemme \ref{lemma-from-chow-to-K}, cette filtration
est localement finie, et donc
$\bigoplus_{p \geq 0} \mathcal{F}^p$,
$\bigoplus_{p \geq 1} \mathcal{F}^p$ et
$\bigoplus_{p \geq 0} \mathcal{F}^p/\mathcal{F}^{p + 1}$
sont des $\mathcal{O}_X$-modules cohérents supportés sur $Z$.
Nous obtenons ainsi
$$
[\mathcal{F}] =
[\bigoplus\nolimits_{p \geq 0} \mathcal{F}^p/\mathcal{F}^{p + 1}]
$$
dans $K_0(\textit{Coh}_Z(X))$, exactement comme dans la démonstration du
Lemme \ref{lemma-from-chow-to-K}. Puisque le module cohérent
$\bigoplus_{p \geq 0} \mathcal{F}^p/\mathcal{F}^{p + 1}$
est annulé par $\mathcal{I}$, nous concluons que
$[\mathcal{F}]$ appartient à l'image. En fait, nous affirmons que l'application
$$
\mathcal{F} \longmapsto 
c(\mathcal{F}) =
[\bigoplus\nolimits_{p \geq 0} \mathcal{F}^p/\mathcal{F}^{p + 1}]
$$
se factorise par $K_0(\textit{Coh}_Z(X))$ et est réciproque de
l'application de l'énoncé du lemme. Pour le voir, il nous
suffit de montrer que, si
$$
0 \to \mathcal{F} \to \mathcal{G} \to \mathcal{H} \to 0
$$
est une suite exacte courte dans $\textit{Coh}_Z(X)$, alors nous
obtenons $c(\mathcal{G}) = c(\mathcal{F}) + c(\mathcal{H})$.
Observons que, pour tout $q \geq 0$, nous avons une suite exacte courte
$$
0 \to
(\mathcal{F} \cap \mathcal{I}^q\mathcal{G})/
(\mathcal{F} \cap \mathcal{I}^{q + 1}\mathcal{G}) \to
\mathcal{G}^q/\mathcal{G}^{q + 1} \to
\mathcal{H}^q/\mathcal{H}^{q + 1} \to 0
$$
Pour $p, q \geq 0$, considérons le sous-module cohérent
$$
\mathcal{F}^{p, q} = \mathcal{I}^p\mathcal{F} \cap \mathcal{I}^q\mathcal{G}
$$
En raisonnant exactement comme ci-dessus et en utilisant que les filtrations
$\mathcal{F}^p = \mathcal{I}^p\mathcal{F}$ et
$\mathcal{F} \cap \mathcal{I}^q\mathcal{G}$ sont localement finies,
nous obtenons
$$
[\bigoplus\nolimits_{p \geq 0} \mathcal{F}^p/\mathcal{F}^{p + 1}] =
[\bigoplus\nolimits_{p, q \geq 0}
\mathcal{F}^{p, q}/(\mathcal{F}^{p + 1, q} + \mathcal{F}^{p, q + 1})] =
[\bigoplus\nolimits_{q \geq 0}
(\mathcal{F} \cap \mathcal{I}^q\mathcal{G})/
(\mathcal{F} \cap \mathcal{I}^{q + 1}\mathcal{G})]
$$
dans $K_0(\textit{Coh}(Z))$. Avec les suites exactes ci-dessus, nous obtenons
le résultat souhaité. Certains détails sont omis.
\end{proof}













\section{Le diviseur associé à un faisceau inversible}
\label{section-divisor-invertible-sheaf}

\noindent
La définition suivante est l'analogue de la
Définition \ref{divisors-definition-divisor-invertible-sheaf} de Diviseurs
dans notre cadre actuel.

\begin{definition}
\label{definition-divisor-invertible-sheaf}
Soit $(S, \delta)$ comme dans la Situation \ref{situation-setup}.
Soit $X$ localement de type fini sur $S$. Supposons $X$
intègre et $n = \dim_\delta(X)$.
Soit $\mathcal{L}$ un $\mathcal{O}_X$-module inversible.
\begin{enumerate}
\item Pour toute section méromorphe non nulle $s$ de $\mathcal{L}$,
nous définissons le {\it diviseur de Weil associé à $s$} comme le
$(n - 1)$-cycle
$$
\text{div}_\mathcal{L}(s) =
\sum \text{ord}_{Z, \mathcal{L}}(s) [Z]
$$
défini dans Diviseurs, Définition
\ref{divisors-definition-divisor-invertible-sheaf}.
Ceci a un sens parce que les diviseurs de Weil sont de $\delta$-dimension $n - 1$
d'après le Lemme \ref{lemma-divisor-delta-dimension}.
\item Nous définissons le {\it diviseur de Weil associé à $\mathcal{L}$} par
$$
c_1(\mathcal{L}) \cap [X] =
\text{classe de }\text{div}_\mathcal{L}(s) \in \CH_{n - 1}(X)
$$
où $s$ est une section méromorphe non nulle quelconque de $\mathcal{L}$ sur
$X$. Ceci est bien défini d'après le
Lemme \ref{divisors-lemma-divisor-meromorphic-well-defined} de Diviseurs.
\end{enumerate}
\end{definition}

\noindent
Soient $X$ et $S$ comme dans la Définition \ref{definition-divisor-invertible-sheaf}
ci-dessus. Posons $n = \dim_\delta(X)$. Il est clair d'après les définitions que
$Cl(X) = \CH_{n - 1}(X)$, où $Cl(X)$ est le groupe des classes de diviseurs de Weil de $X$
défini dans Diviseurs, Définition \ref{divisors-definition-class-group}.
L'application
$$
\Pic(X) \longrightarrow \CH_{n - 1}(X), \quad
\mathcal{L} \longmapsto c_1(\mathcal{L}) \cap [X]
$$
est la même que l'application $\Pic(X) \to Cl(X)$ construite dans
Diviseurs, Équation (\ref{divisors-equation-c1}), pour les schémas intègres
localement noethériens quelconques. En particulier, cette application
est un homomorphisme de groupes abéliens ; elle est injective si $X$ est
un schéma normal, et un isomorphisme si tous les anneaux locaux de $X$
sont factoriels. Voir Diviseurs, Lemmes \ref{divisors-lemma-normal-c1-injective} et
\ref{divisors-lemma-local-rings-UFD-c1-bijective}.
Dans certains cas, il est facile de calculer le
diviseur de Weil associé à un faisceau inversible.

\begin{lemma}
\label{lemma-compute-c1}
Soit $(S, \delta)$ comme dans la Situation \ref{situation-setup}.
Soit $X$ localement de type fini sur $S$. Supposons $X$
intègre et $n = \dim_\delta(X)$.
Soit $\mathcal{L}$ un $\mathcal{O}_X$-module inversible.
Soit $s \in \Gamma(X, \mathcal{L})$ une section globale non nulle.
Alors
$$
\text{div}_\mathcal{L}(s) = [Z(s)]_{n - 1}
$$
dans $Z_{n - 1}(X)$ et
$$
c_1(\mathcal{L}) \cap [X] = [Z(s)]_{n - 1}
$$
dans $\CH_{n - 1}(X)$.
\end{lemma}

\begin{proof}
Soit $Z \subset X$ un sous-schéma fermé intègre de
$\delta$-dimension $n - 1$. Soit $\xi \in Z$ son point
générique. Choisissons un générateur $s_\xi \in \mathcal{L}_\xi$.
Écrivons $s = fs_\xi$ pour un certain $f \in \mathcal{O}_{X, \xi}$.
Par définition de $Z(s)$, voir la
Définition \ref{divisors-definition-zero-scheme-s} de Diviseurs,
nous voyons que $Z(s)$ est défini par un faisceau quasi-cohérent
d'idéaux $\mathcal{I} \subset \mathcal{O}_X$ tel
que $\mathcal{I}_\xi = (f)$. Ainsi
$\text{longueur}_{\mathcal{O}_{X, \xi}}(\mathcal{O}_{Z(s), \xi})
=
\text{longueur}_{\mathcal{O}_{X, \xi}}(\mathcal{O}_{X, \xi}/(f))
=
\text{ord}_{\mathcal{O}_{X, \xi}}(f)$, comme souhaité.
\end{proof}

\noindent
Le lemme suivant sera supplanté par le Lemme plus général
\ref{lemma-flat-pullback-cap-c1}.

\begin{lemma}
\label{lemma-flat-pullback-divisor-invertible-sheaf}
Soit $(S, \delta)$ comme dans la Situation \ref{situation-setup}.
Soient $X$, $Y$ localement de type fini sur $S$. Supposons $X$, $Y$
intègres et $n = \dim_\delta(Y)$.
Soit $\mathcal{L}$ un $\mathcal{O}_Y$-module inversible.
Soit $f : X \to Y$ un morphisme plat de dimension relative $r$. Alors
$$
f^*(c_1(\mathcal{L}) \cap [Y]) = c_1(f^*\mathcal{L}) \cap [X]
$$
dans $\CH_{n + r - 1}(X)$.
\end{lemma}

\begin{proof}
Soit $s$ une section méromorphe non nulle de $\mathcal{L}$.
Nous allons montrer qu'en fait
$f^*\text{div}_\mathcal{L}(s) = \text{div}_{f^*\mathcal{L}}(f^*s)$,
et donc que le lemme est vrai.
Pour le voir, soit $\xi \in Y$ un point, et soit $s_\xi \in \mathcal{L}_\xi$
un générateur. Écrivons $s = gs_\xi$ avec $g \in R(Y)^*$.
Il existe alors un voisinage ouvert $V \subset Y$ de $\xi$
tel que $s_\xi \in \mathcal{L}(V)$ et que $s_\xi$ engendre
$\mathcal{L}|_V$. Nous voyons donc que
$$
\text{div}_\mathcal{L}(s)|_V = \text{div}_Y(g)|_V.
$$
De la même manière exactement, puisque $f^*s_\xi$ engendre $f^*\mathcal{L}$
sur $f^{-1}(V)$ et puisque $f^*s = g f^*s_\xi$, nous
avons aussi
$$
\text{div}_\mathcal{L}(f^*s)|_{f^{-1}(V)}
=
\text{div}_X(g)|_{f^{-1}(V)}.
$$
L'égalité souhaitée de cycles sur $f^{-1}(V)$ résulte donc du
résultat correspondant pour les images réciproques de diviseurs principaux ; voir le
Lemme \ref{lemma-flat-pullback-principal-divisor}.
\end{proof}



\section{Intersection par un faisceau inversible}
\label{section-intersecting-with-divisors}

\noindent
Dans cette section, nous étudions la construction suivante.

\begin{definition}
\label{definition-cap-c1}
Soit $(S, \delta)$ comme dans la Situation \ref{situation-setup}.
Soit $X$ localement de type fini sur $S$.
Soit $\mathcal{L}$ un $\mathcal{O}_X$-module inversible.
Nous définissons, pour tout entier $k$, une opération
$$
c_1(\mathcal{L}) \cap - :
Z_{k + 1}(X) \to \CH_k(X)
$$
appelée {\it intersection avec la première classe de Chern de $\mathcal{L}$}.
\begin{enumerate}
\item Étant donné un sous-schéma fermé intègre $i : W \to X$ avec
$\dim_\delta(W) = k + 1$, nous définissons
$$
c_1(\mathcal{L}) \cap [W] = i_*(c_1({i^*\mathcal{L}}) \cap [W])
$$
où le membre de droite est défini dans la
Définition \ref{definition-divisor-invertible-sheaf}.
\item Pour un $(k + 1)$-cycle général $\alpha = \sum n_i [W_i]$, nous posons
$$
c_1(\mathcal{L}) \cap \alpha = \sum n_i c_1(\mathcal{L}) \cap [W_i]
$$
\end{enumerate}
\end{definition}

\noindent
Écrivons chaque $c_1(\mathcal{L}) \cap W_i = \sum_j n_{i, j} [Z_{i, j}]$,
où $\{Z_{i, j}\}_j$ est une somme localement finie
de sous-schémas fermés intègres de $W_i$. Puisque $\{W_i\}$ est une famille
localement finie de sous-schémas fermés intègres de $X$, il en résulte
facilement que $\{Z_{i, j}\}_{i, j}$ est une famille localement finie
de sous-schémas fermés de $X$. Ainsi
$c_1(\mathcal{L}) \cap \alpha = \sum n_in_{i, j}[Z_{i, j}]$
est un cycle. Une autre manière, plus commode, de voir cela
est d'observer que le morphisme $\coprod W_i \to X$ est
propre. Ainsi $c_1(\mathcal{L}) \cap \alpha$ peut être considéré
comme l'image directe d'une classe de $\CH_k(\coprod W_i) = \prod \CH_k(W_i)$.
Cela explique aussi pourquoi le résultat est bien défini modulo équivalence
rationnelle sur $X$.

\medskip\noindent
L'objectif principal des quelques sections suivantes est de montrer que l'intersection avec
$c_1(\mathcal{L})$ se factorise par l'équivalence rationnelle.
Ce n'est pas une trivialité.

\begin{lemma}
\label{lemma-c1-cap-additive}
Soit $(S, \delta)$ comme dans la Situation \ref{situation-setup}.
Soit $X$ localement de type fini sur $S$.
Soient $\mathcal{L}$, $\mathcal{N}$ des faisceaux inversibles sur $X$.
Alors
$$
c_1(\mathcal{L}) \cap \alpha  + c_1(\mathcal{N}) \cap \alpha =
c_1(\mathcal{L} \otimes_{\mathcal{O}_X} \mathcal{N}) \cap \alpha
$$
dans $\CH_k(X)$ pour tout $\alpha \in Z_{k + 1}(X)$. De plus,
$c_1(\mathcal{O}_X) \cap \alpha = 0$ pour tout $\alpha$.
\end{lemma}

\begin{proof}
L'additivité résulte directement du
Lemme \ref{divisors-lemma-c1-additive} de Diviseurs
et des définitions. Pour voir que $c_1(\mathcal{O}_X) \cap \alpha = 0$,
considérons la section $1 \in \Gamma(X, \mathcal{O}_X)$. Sa restriction
est une section partout non nulle sur tout sous-schéma fermé intègre
$W \subset X$. Ainsi $c_1(\mathcal{O}_X) \cap [W] = 0$, comme souhaité.
\end{proof}

\noindent
Rappelons que $Z(s) \subset X$ désigne le schéma des zéros d'une section globale
$s$ d'un faisceau inversible sur un schéma $X$ ; voir la
Définition \ref{divisors-definition-zero-scheme-s} de Diviseurs.

\begin{lemma}
\label{lemma-prepare-geometric-cap}
Soit $(S, \delta)$ comme dans la Situation \ref{situation-setup}.
Soit $Y$ localement de type fini sur $S$.
Soit $\mathcal{L}$ un $\mathcal{O}_Y$-module inversible.
Soit $s \in \Gamma(Y, \mathcal{L})$.
Supposons
\begin{enumerate}
\item $\dim_\delta(Y) \leq k + 1$,
\item $\dim_\delta(Z(s)) \leq k$, et
\item pour tout point générique $\xi$ d'une composante irréductible de
$Z(s)$ de $\delta$-dimension $k$, la multiplication par $s$
induit une injection $\mathcal{O}_{Y, \xi} \to \mathcal{L}_\xi$.
\end{enumerate}
Écrivons $[Y]_{k + 1} = \sum n_i[Y_i]$, où les $Y_i \subset Y$ sont les
composantes irréductibles de $Y$ de $\delta$-dimension $k + 1$.
Posons $s_i = s|_{Y_i} \in \Gamma(Y_i, \mathcal{L}|_{Y_i})$. Alors
\begin{equation}
\label{equation-equal-as-cycles}
[Z(s)]_k =  \sum n_i[Z(s_i)]_k
\end{equation}
comme $k$-cycles sur $Y$.
\end{lemma}

\begin{proof}
Soit $Z \subset Y$ un sous-schéma fermé intègre de
$\delta$-dimension $k$. Soit $\xi \in Z$ son point générique.
Nous voulons comparer le coefficient $n$ de $[Z]$ dans l'expression
$\sum n_i[Z(s_i)]_k$ au coefficient $m$ de $[Z]$ dans
l'expression $[Z(s)]_k$. Choisissons un générateur $s_\xi \in \mathcal{L}_\xi$.
Écrivons $A = \mathcal{O}_{Y, \xi}$, $L = \mathcal{L}_\xi$.
Alors $L = As_\xi$. Écrivons $s = f s_\xi$ pour un certain (unique) $f \in A$.
L'hypothèse (3) signifie que $f : A \to A$ est injectif.
Puisque $\dim_\delta(Y) \leq k + 1$ et $\dim_\delta(Z) = k$,
nous avons $\dim(A) = 0$ ou $1$. Nous avons
$$
m = \text{longueur}_A(A/(f))
$$
qui est finie dans les deux cas.

\medskip\noindent
Si $\dim(A) = 0$, alors l'injectivité de $f : A \to A$
implique que $f \in A^*$. Ainsi, dans ce cas, $m$ est nul.
De plus, la condition $\dim(A) = 0$ signifie que $\xi$
n'appartient à aucune composante irréductible de $\delta$-dimension
$k + 1$, c'est-à-dire que $n = 0$ également.

\medskip\noindent
Supposons maintenant $\dim(A) = 1$.
Puisque $A$ est un anneau local noethérien, il possède un nombre fini
d'idéaux premiers minimaux $\mathfrak q_1, \ldots, \mathfrak q_t$.
Ceux-ci correspondent bijectivement aux $Y_i$ passant par $\xi'$.
De plus, $n_i = \text{longueur}_{A_{\mathfrak q_i}}(A_{\mathfrak q_i})$.
La multiplicité de $[Z]$ dans $[Z(s_i)]_k$ est aussi
$\text{longueur}_A(A/(f, \mathfrak q_i))$.
L'équation à démontrer dans ce cas est donc
$$
\text{longueur}_A(A/(f))
=
\sum \text{longueur}_{A_{\mathfrak q_i}}(A_{\mathfrak q_i})
\text{longueur}_A(A/(f, \mathfrak q_i))
$$
qui résulte du
Lemme \ref{lemma-additivity-divisors-restricted}.
\end{proof}

\noindent
Le lemme suivant est un résultat utile pour calculer le produit d'intersection
de $c_1$ d'un faisceau inversible avec le cycle associé
à un sous-schéma fermé.
Rappelons que $Z(s) \subset X$ désigne le schéma des zéros d'une section globale
$s$ d'un faisceau inversible sur un schéma $X$ ; voir la
Définition \ref{divisors-definition-zero-scheme-s} de Diviseurs.

\begin{lemma}
\label{lemma-geometric-cap}
Soit $(S, \delta)$ comme dans la Situation \ref{situation-setup}.
Soit $X$ localement de type fini sur $S$.
Soit $\mathcal{L}$ un $\mathcal{O}_X$-module inversible.
Soit $Y \subset X$ un sous-schéma fermé.
Soit $s \in \Gamma(Y, \mathcal{L}|_Y)$.
Supposons
\begin{enumerate}
\item $\dim_\delta(Y) \leq k + 1$,
\item $\dim_\delta(Z(s)) \leq k$, et
\item pour tout point générique $\xi$ d'une composante irréductible de
$Z(s)$ de $\delta$-dimension $k$, la multiplication par $s$
induit une injection
$\mathcal{O}_{Y, \xi} \to (\mathcal{L}|_Y)_\xi$\footnote{Par exemple,
ceci est vrai si $s$ est une section régulière de $\mathcal{L}|_Y$.}.
\end{enumerate}
Alors
$$
c_1(\mathcal{L}) \cap [Y]_{k + 1} = [Z(s)]_k
$$
dans $\CH_k(X)$.
\end{lemma}

\begin{proof}
Écrivons
$$
[Y]_{k + 1} = \sum n_i[Y_i]
$$
où les $Y_i \subset Y$ sont les composantes irréductibles de
$Y$ de $\delta$-dimension $k + 1$ et $n_i > 0$.
Par hypothèse, la restriction
$s_i = s|_{Y_i} \in \Gamma(Y_i, \mathcal{L}|_{Y_i})$ n'est pas
nulle, et est donc une section régulière. D'après le Lemme \ref{lemma-compute-c1},
nous voyons que $[Z(s_i)]_k$ représente $c_1(\mathcal{L}|_{Y_i})$.
Ainsi, par définition,
$$
c_1(\mathcal{L}) \cap [Y]_{k + 1} = \sum n_i[Z(s_i)]_k
$$
Le résultat résulte donc du Lemme \ref{lemma-prepare-geometric-cap}.
\end{proof}




\section{Intersection par un faisceau inversible et images directe et réciproque}
\label{section-intersecting-with-divisors-push-pull}

\noindent
Dans cette section, nous démontrons que l'opération $c_1(\mathcal{L}) \cap -$
commute avec l'image réciproque plate et l'image directe propre.

\begin{lemma}
\label{lemma-prepare-flat-pullback-cap-c1}
Soit $(S, \delta)$ comme dans la Situation \ref{situation-setup}.
Soient $X$, $Y$ localement de type fini sur $S$.
Soit $f : X \to Y$ un morphisme plat de dimension relative $r$.
Soit $\mathcal{L}$ un faisceau inversible sur $Y$.
Supposons $Y$ intègre et $n = \dim_\delta(Y)$.
Soit $s$ une section méromorphe non nulle de $\mathcal{L}$.
Alors nous avons
$$
f^*\text{div}_\mathcal{L}(s) = \sum n_i\text{div}_{f^*\mathcal{L}|_{X_i}}(s_i)
$$
dans $Z_{n + r - 1}(X)$. Ici, la somme porte sur les composantes
irréductibles $X_i \subset X$ de $\delta$-dimension $n + r$,
la section $s_i = f|_{X_i}^*(s)$ est l'image réciproque de $s$, et
$n_i = m_{X_i, X}$ est la multiplicité de $X_i$ dans $X$.
\end{lemma}

\begin{proof}
Pour démontrer cette égalité de cycles, nous pouvons travailler localement sur $Y$.
Nous pouvons donc supposer $Y$ affine et $s = p/q$ pour des sections non nulles
$p \in \Gamma(Y, \mathcal{L})$ et $q \in \Gamma(Y, \mathcal{O})$.
Si nous pouvons démontrer à la fois
$$
f^*\text{div}_\mathcal{L}(p) =
\sum n_i\text{div}_{f^*\mathcal{L}|_{X_i}}(p_i)
\quad\text{et}\quad
f^*\text{div}_\mathcal{O}(q) =
\sum n_i\text{div}_{\mathcal{O}_{X_i}}(q_i)
$$
(avec les notations évidentes), alors nous avons le résultat par
additivité ; voir le Lemme \ref{divisors-lemma-c1-additive} de Diviseurs.
Nous pouvons donc supposer que $s \in \Gamma(Y, \mathcal{L})$.
Dans ce cas, nous pouvons appliquer l'égalité
(\ref{equation-equal-as-cycles}) pour voir que
$$
[Z(f^*(s))]_{k + r - 1} =
\sum n_i\text{div}_{f^*\mathcal{L}|_{X_i}}(s_i)
$$
où $f^*(s) \in f^*\mathcal{L}$ désigne l'image réciproque de $s$ sur $X$.
D'autre part, nous avons
$$
f^*\text{div}_\mathcal{L}(s) = f^*[Z(s)]_{k - 1}
= [f^{-1}(Z(s))]_{k + r - 1},
$$
d'après les Lemmes \ref{lemma-compute-c1} et \ref{lemma-pullback-coherent}.
Puisque $Z(f^*(s)) = f^{-1}(Z(s))$, nous avons le résultat.
\end{proof}

\begin{lemma}
\label{lemma-flat-pullback-cap-c1}
Soit $(S, \delta)$ comme dans la Situation \ref{situation-setup}.
Soient $X$, $Y$ localement de type fini sur $S$.
Soit $f : X \to Y$ un morphisme plat de dimension relative $r$.
Soit $\mathcal{L}$ un faisceau inversible sur $Y$.
Soit $\alpha$ un $k$-cycle sur $Y$.
Alors
$$
f^*(c_1(\mathcal{L}) \cap \alpha) = c_1(f^*\mathcal{L}) \cap f^*\alpha
$$
dans $\CH_{k + r - 1}(X)$.
\end{lemma}

\begin{proof}
Écrivons $\alpha = \sum n_i[W_i]$. Nous allons montrer que
$$
f^*(c_1(\mathcal{L}) \cap [W_i]) = c_1(f^*\mathcal{L}) \cap f^*[W_i]
$$
dans $\CH_{k + r - 1}(X)$ en produisant une équivalence rationnelle
sur le sous-schéma fermé $f^{-1}(W_i)$ de $X$.
D'après la discussion de la
Remarque \ref{remark-infinite-sums-rational-equivalences},
ceci démontrera l'égalité du lemme.

\medskip\noindent
Soit $W \subset Y$ un sous-schéma fermé intègre de $\delta$-dimension $k$.
Considérons le sous-schéma fermé $W' = f^{-1}(W) = W \times_Y X$,
de sorte que nous avons le diagramme cartésien
$$
\xymatrix{
W' \ar[r] \ar[d]_h & X \ar[d]^f \\
W \ar[r] & Y
}
$$
Nous devons montrer que
$f^*(c_1(\mathcal{L}) \cap [W]) = c_1(f^*\mathcal{L}) \cap f^*[W]$.
Choisissons une section méromorphe non nulle $s$ de $\mathcal{L}|_W$.
Soient $W'_i \subset W'$ les composantes irréductibles de
$\delta$-dimension $k + r$. Écrivons $[W']_{k + r} = \sum n_i[W'_i]$,
où $n_i$ est la multiplicité de $W'_i$ dans $W'$, conformément à la définition.
Ainsi $f^*[W] = \sum n_i[W'_i]$ dans $Z_{k + r}(X)$.
Puisque chaque $W'_i \to W$ est dominant, nous
voyons que $s_i = s|_{W'_i}$ est une section méromorphe non nulle pour
tout $i$. D'après le Lemme \ref{lemma-prepare-flat-pullback-cap-c1},
nous avons l'égalité de cycles suivante
$$
h^*\text{div}_{\mathcal{L}|_W}(s) =
\sum n_i\text{div}_{f^*\mathcal{L}|_{W'_i}}(s_i)
$$
dans $Z_{k + r - 1}(W')$. Ceci achève la démonstration, puisque
le membre de gauche est un cycle sur $W'$ dont l'image directe est
$f^*(c_1(\mathcal{L}) \cap [W])$ dans $\CH_{k + r - 1}(X)$,
et le membre de droite est un cycle sur $W'$ dont l'image directe est
$c_1(f^*\mathcal{L}) \cap f^*[W]$ dans $\CH_{k + r - 1}(X)$.
\end{proof}

\begin{lemma}
\label{lemma-equal-c1-as-cycles}
Soit $(S, \delta)$ comme dans la Situation \ref{situation-setup}.
Soient $X$, $Y$ localement de type fini sur $S$.
Soit $f : X \to Y$ un morphisme propre.
Soit $\mathcal{L}$ un faisceau inversible sur $Y$.
Soit $s$ une section méromorphe non nulle, encore notée $s$, de $\mathcal{L}$ sur $Y$.
Supposons $X$, $Y$ intègres, $f$ dominant et $\dim_\delta(X) = \dim_\delta(Y)$.
Alors
$$
f_*\left(\text{div}_{f^*\mathcal{L}}(f^*s)\right) =
[R(X) : R(Y)]\text{div}_\mathcal{L}(s).
$$
comme cycles sur $Y$. En particulier,
$$
f_*(c_1(f^*\mathcal{L}) \cap [X]) =
[R(X) : R(Y)] c_1(\mathcal{L}) \cap [Y] =
c_1(\mathcal{L}) \cap f_*[X]
$$
\end{lemma}

\begin{proof}
La dernière égalité résulte de la première puisque $f_*[X] = [R(X) : R(Y)][Y]$
par définition. Il se trouve que nous pouvons réutiliser le
Lemme \ref{lemma-proper-pushforward-alteration}
pour le démontrer. En effet, puisque nous cherchons à démontrer une égalité
de cycles, nous pouvons travailler localement sur $Y$. Nous pouvons donc supposer
que $\mathcal{L} = \mathcal{O}_Y$. Dans ce cas, $s$
correspond à une fonction rationnelle $g \in R(Y)$, et
nous cherchons simplement à démontrer
$$
f_*\left(\text{div}_X(g)\right) =
[R(X) : R(Y)]\text{div}_Y(g).
$$
En comparant avec le résultat du
Lemme \ref{lemma-proper-pushforward-alteration} précité,
nous voyons que c'est vrai puisque
$\text{Nm}_{R(X)/R(Y)}(g) = g^{[R(X) : R(Y)]}$,
car $g \in R(Y)^*$.
\end{proof}

\begin{lemma}
\label{lemma-pushforward-cap-c1}
Soit $(S, \delta)$ comme dans la Situation \ref{situation-setup}.
Soient $X$, $Y$ localement de type fini sur $S$.
Soit $p : X \to Y$ un morphisme propre.
Soit $\alpha \in Z_{k + 1}(X)$.
Soit $\mathcal{L}$ un faisceau inversible sur $Y$.
Alors
$$
p_*(c_1(p^*\mathcal{L}) \cap \alpha) = c_1(\mathcal{L}) \cap p_*\alpha
$$
dans $\CH_k(Y)$.
\end{lemma}

\begin{proof}
Supposons que $p$ ait la propriété que, pour tout sous-schéma fermé
intègre $W \subset X$, l'application $p|_W : W \to Y$
soit une immersion fermée. Alors, par définition du produit
avec $c_1(\mathcal{L})$, le lemme est vrai.

\medskip\noindent
Nous allons utiliser cette remarque pour nous ramener à un cas particulier. En effet,
écrivons $\alpha = \sum n_i[W_i]$ avec $n_i \not = 0$ et les $W_i$ deux à deux
distincts. Soit $W'_i \subset Y$ l'image de $W_i$ (comme sous-schéma
fermé intègre). Considérons le diagramme
$$
\xymatrix{
X' = \coprod W_i \ar[r]_-q \ar[d]_{p'} & X \ar[d]^p \\
Y' = \coprod W'_i \ar[r]^-{q'} & Y.
}
$$
Puisque $\{W_i\}$ est localement finie sur $X$ et que $p$ est propre,
nous voyons que $\{W'_i\}$ est localement finie sur $Y$ et que
$q, q', p'$ sont aussi des morphismes propres.
Nous pouvons aussi considérer $\sum n_i[W_i]$ comme un $k$-cycle
$\alpha' \in Z_k(X')$. Il est clair que $q_*\alpha' = \alpha$.
Nous avons
$q_*(c_1(q^*p^*\mathcal{L}) \cap \alpha')
= c_1(p^*\mathcal{L}) \cap q_*\alpha'$
et
$(q')_*(c_1((q')^*\mathcal{L}) \cap p'_*\alpha') =
c_1(\mathcal{L}) \cap q'_*p'_*\alpha'$ d'après la remarque initiale
de la démonstration. Il suffit donc de démontrer le lemme
pour le morphisme $p'$ et le cycle $\sum n_i[W_i]$.
Il est clair que cela signifie que nous pouvons supposer $X$, $Y$ intègres,
$f : X \to Y$ dominant et $\alpha = [X]$.
Dans ce cas, le résultat résulte du
Lemme \ref{lemma-equal-c1-as-cycles}.
\end{proof}






\section{La formule clé}
\label{section-key}

\noindent
Soit $(S, \delta)$ comme dans la Situation \ref{situation-setup}.
Soit $X$ localement de type fini sur $S$. Supposons
$X$ intègre et $\dim_\delta(X) = n$.
Soient $\mathcal{L}$ et $\mathcal{N}$ des faisceaux inversibles sur $X$.
Soit $s$ une section méromorphe non nulle de $\mathcal{L}$ et
soit $t$ une section méromorphe non nulle de $\mathcal{N}$.
Soit $Z_i \subset X$, $i \in I$, une famille localement finie de fermés
irréductibles de codimension $1$ ayant la propriété suivante :
si $Z \not \in \{Z_i\}$ a pour point générique $\xi$, alors $s$ est un générateur
de $\mathcal{L}_\xi$ et $t$ est un générateur de $\mathcal{N}_\xi$.
Une telle famille existe d'après le
Lemme \ref{divisors-lemma-divisor-meromorphic-locally-finite} de Diviseurs.
Alors
$$
\text{div}_\mathcal{L}(s) = \sum \text{ord}_{Z_i, \mathcal{L}}(s) [Z_i]
$$
et, de même,
$$
\text{div}_\mathcal{N}(t) = \sum \text{ord}_{Z_i, \mathcal{N}}(t) [Z_i]
$$
En développant davantage les définitions, choisissons, pour chaque $i$, des générateurs
$s_i \in \mathcal{L}_{\xi_i}$ et $t_i \in \mathcal{N}_{\xi_i}$,
où $\xi_i$ est le point générique de $Z_i$. Nous pouvons alors écrire
$$
s = f_i s_i
\quad\text{et}\quad
t = g_i t_i
$$
Posons $B_i = \mathcal{O}_{X, \xi_i}$. Alors, par définition,
$$
\text{ord}_{Z_i, \mathcal{L}}(s) = \text{ord}_{B_i}(f_i)
\quad\text{et}\quad
\text{ord}_{Z_i, \mathcal{N}}(t) = \text{ord}_{B_i}(g_i)
$$
Puisque $t_i$ est un générateur de $\mathcal{N}_{\xi_i}$, nous voyons que
son image dans la fibre $\mathcal{N}_{\xi_i} \otimes \kappa(\xi_i)$
est une section méromorphe non nulle de $\mathcal{N}|_{Z_i}$. Nous noterons
cette image $t_i|_{Z_i}$. Il résulte de nos définitions que
$$
c_1(\mathcal{N}) \cap \text{div}_\mathcal{L}(s) =
\sum \text{ord}_{B_i}(f_i)
(Z_i \to X)_*\text{div}_{\mathcal{N}|_{Z_i}}(t_i|_{Z_i})
$$
et, de même,
$$
c_1(\mathcal{L}) \cap \text{div}_\mathcal{N}(t) =
\sum \text{ord}_{B_i}(g_i)
(Z_i \to X)_*\text{div}_{\mathcal{L}|_{Z_i}}(s_i|_{Z_i})
$$
dans $\CH_{n - 2}(X)$. Nous allons trouver une équivalence rationnelle entre
ces deux cycles. Pour cela, nous considérons le symbole modéré
$$
\partial_{B_i}(f_i, g_i) \in \kappa(\xi_i)^*
$$
voir la section \ref{section-tame-symbol}.

\begin{lemma}[Formule fondamentale]
\label{lemma-key-formula}
Dans la situation ci-dessus, le cycle
$$
\sum
(Z_i \to X)_*\left(
\text{ord}_{B_i}(f_i) \text{div}_{\mathcal{N}|_{Z_i}}(t_i|_{Z_i}) -
\text{ord}_{B_i}(g_i) \text{div}_{\mathcal{L}|_{Z_i}}(s_i|_{Z_i}) \right)
$$
est égal au cycle
$$
\sum (Z_i \to X)_*\text{div}(\partial_{B_i}(f_i, g_i))
$$
\end{lemma}

\begin{proof}
Examinons d'abord ce qui se passe si l'on remplace $s_i$ par $us_i$,
où $u$ est une unité de $B_i$. Alors $f_i$ est remplacé par $u^{-1} f_i$.
La première partie de la première expression du lemme reste donc inchangée,
et l'on ajoute à la seconde partie
$$
-\text{ord}_{B_i}(g_i)\text{div}(u|_{Z_i})
$$
(où $u|_{Z_i}$ est l'image de $u$ dans le corps résiduel), d'après
Diviseurs, lemme \ref{divisors-lemma-divisor-meromorphic-well-defined},
tandis que l'on ajoute à la seconde expression
$$
\text{div}(\partial_{B_i}(u^{-1}, g_i))
$$
par bilinéarité du symbole modéré. Ces termes coïncident d'après la propriété
(\ref{item-normalization}) du symbole modéré.

\medskip\noindent
Soit $Z \subset X$ une partie fermée irréductible telle que $\dim_\delta(Z) = n - 2$.
Pour montrer que les coefficients de $Z$ dans les deux cycles du lemme
coïncident, on peut effectuer un remplacement $s_i \mapsto us_i$ comme au
paragraphe précédent. De la même manière, on montre que l'on peut effectuer un remplacement
$t_i \mapsto vt_i$, où $v$ est une unité de $B_i$.

\medskip\noindent
Puisqu'il s'agit de démontrer une égalité de cycles, on peut raisonner coefficient
par coefficient. Choisissons donc une partie fermée irréductible $Z \subset X$
telle que $\dim_\delta(Z) = n - 2$, et comparons les coefficients. Soit $\xi \in Z$
son point générique, et posons $A = \mathcal{O}_{X, \xi}$. C'est un anneau local
noethérien intègre de dimension $2$. Choisissons des générateurs $\sigma$ et $\tau$
de $\mathcal{L}_\xi$ et $\mathcal{N}_\xi$. Après avoir restreint $X$, on peut
supposer, et l'on suppose, que $\sigma$ et $\tau$ définissent des trivialisations
des faisceaux inversibles $\mathcal{L}$ et $\mathcal{N}$ sur tout $X$.
La famille $Z_i$ étant localement
finie, on peut, après avoir restreint $X$, supposer $Z \subset Z_i$ pour tout $i \in I$
et $I$ fini. Alors $\xi_i$ correspond à un idéal premier
$\mathfrak q_i \subset A$ de hauteur $1$.
On peut écrire $s_i = a_i \sigma$ et $t_i = b_i \tau$,
où $a_i$ et $b_i$ sont des unités de $A_{\mathfrak q_i}$.
D'après les remarques précédentes, il suffit de démontrer le lemme lorsque
$a_i = b_i = 1$ pour tout $i$.

\medskip\noindent
Supposons $a_i = b_i = 1$ pour tout $i$. Alors la première expression du
lemme est nulle, car nous avons choisi $\sigma$ et $\tau$ comme sections
trivialisantes. Écrivons $s = f\sigma$ et $t = g \tau$, avec $f$ et $g$ dans le
corps des fractions de $A$. Le paragraphe précédent nous a ramenés au cas où
$f_i = f$ et $g_i = g$ pour tout $i$. De plus, pour tout idéal premier de hauteur $1$
$\mathfrak q$ de $A$ qui n'appartient pas à $\{\mathfrak q_i\}$,
$f$ et $g$ sont des unités de $A_\mathfrak q$ (d'après le choix de
la famille $\{Z_i\}$ dans la discussion précédant le lemme). Ainsi,
le coefficient de $Z$ dans la seconde expression du lemme est
$$
\sum\nolimits_i \text{ord}_{A/\mathfrak q_i}(\partial_{B_i}(f, g))
$$
qui est nul d'après le lemme fondamental \ref{lemma-milnor-gersten-low-degree}.
\end{proof}

\begin{remark}
\label{remark-higher-chow-pointwise}
Soit $(S, \delta)$ comme dans la situation \ref{situation-setup}.
Soit $X$ localement de type fini sur $S$. Soit $k \in \mathbf{Z}$.
Nous affirmons qu'il existe un complexe
$$
\bigoplus\nolimits_{\delta(x) = k + 2}' K_2^M(\kappa(x))
\xrightarrow{\partial}
\bigoplus\nolimits_{\delta(x) = k + 1}' K_1^M(\kappa(x))
\xrightarrow{\partial}
\bigoplus\nolimits_{\delta(x) = k}' K_0^M(\kappa(x))
$$
Nous utilisons ici les notations et conventions introduites dans la
remarque \ref{remark-chow-group-pointwise}, avec en outre les précisions suivantes :
\begin{enumerate}
\item $K_2^M(\kappa(x))$ est la composante de degré $2$ de
la K-théorie de Milnor du corps résiduel $\kappa(x)$ du point
$x \in X$ (voir la remarque \ref{remark-gersten-complex-milnor}) ; c'est
le quotient de $\kappa(x)^* \otimes_\mathbf{Z} \kappa(x)^*$
par le sous-groupe engendré par les éléments de la forme
$\lambda \otimes (1 - \lambda)$, pour
$\lambda \in \kappa(x) \setminus \{0, 1\}$, et
\item la première différentielle $\partial$ est définie comme suit :
étant donné un élément $\xi = \sum_x \alpha_x$ du premier terme,
on pose
$$
\partial(\xi) = \sum\nolimits_{x \leadsto x',\ \delta(x') = k + 1}
\partial_{\mathcal{O}_{W_x, x'}}(\alpha_x)
$$
où
$\partial_{\mathcal{O}_{W_x, x'}} : K_2^M(\kappa(x)) \to K_1^M(\kappa(x))$
est le symbole modéré construit dans la section \ref{section-tame-symbol}.
\end{enumerate}
Nous affirmons que l'on obtient un complexe, c'est-à-dire que $\partial \circ \partial = 0$.
Pour le voir, il suffit de prendre un élément $\xi$ comme ci-dessus et un point
$x'' \in X$ tel que $\delta(x'') = k$, et de vérifier que le coefficient de
$x''$ dans l'élément $\partial(\partial(\xi))$ est nul.
Comme $\xi = \sum \alpha_x$ est une somme localement finie, on
peut en fait supposer, par additivité, que $\xi = \alpha_x$ pour
un certain $x \in X$ tel que $\delta(x) = k + 2$ et $\alpha_x \in K_2^M(\kappa(x))$.
Par linéarité, on peut encore supposer que $\alpha_x = f \otimes g$ pour
certains $f, g \in \kappa(x)^*$. Notons $W \subset X$ le sous-schéma fermé
intègre de point générique $x$. Si $x'' \not \in W$, il est
immédiatement clair que le coefficient de $x$ dans $\partial(\partial(\xi))$
est nul. Si $x'' \in W$, on voit que le coefficient de $x''$
dans $\partial(\partial(x))$ est égal à
$$
\sum\nolimits_{x \leadsto x' \leadsto x'',\ \delta(x') = k + 1}
\text{ord}_{\mathcal{O}_{\overline{\{x'\}}, x''}}(
\partial_{\mathcal{O}_{W, x'}}(f, g))
$$
Le lemme algébrique fondamental \ref{lemma-milnor-gersten-low-degree}
affirme précisément que cette somme est nulle.
\end{remark}

\begin{remark}
\label{remark-higher-chow}
Soit $(S, \delta)$ comme dans la situation \ref{situation-setup}.
Soit $X$ localement de type fini sur $S$. Soit $k \in \mathbf{Z}$.
Le complexe de la remarque \ref{remark-higher-chow-pointwise} et la
présentation de $\CH_k(X)$ dans la remarque \ref{remark-chow-group-pointwise}
suggèrent de définir un premier groupe de Chow supérieur par
$$
\CH^M_k(X, 1) =
H_1(\text{le complexe de la remarque \ref{remark-higher-chow-pointwise}})
$$
Nous utilisons l'exposant ${}^M$ pour distinguer notre notation de celle des
groupes de Chow supérieurs définis dans la littérature, par exemple dans les articles
de Spencer Bloch (\cite{Bloch} et \cite{Bloch-moving}).
Soit $U \subset X$ un ouvert dont le complémentaire est $Y \subset X$ (muni de sa structure de
sous-schéma fermé réduit). On obtient alors une suite exacte courte scindée
$$
0 \to
\bigoplus\nolimits_{y \in Y, \delta(y) = k + i}' K_i^M(\kappa(y)) \to
\bigoplus\nolimits_{x \in X, \delta(x) = k + i}' K_i^M(\kappa(x)) \to
\bigoplus\nolimits_{u \in U, \delta(u) = k + i}' K_i^M(\kappa(u)) \to 0
$$
pour $i = 2, 1, 0$, compatible avec les opérateurs de bord des complexes
de la remarque \ref{remark-higher-chow-pointwise}. En appliquant le lemme du serpent
(voir Homologie, lemme \ref{homology-lemma-long-exact-sequence-chain}),
on obtient une suite exacte à six termes
$$
\CH^M_k(Y, 1) \to \CH^M_k(X, 1) \to \CH^M_k(U, 1) \to
\CH_k(Y) \to \CH_k(X) \to \CH_k(U) \to 0
$$
prolongeant la suite exacte canonique du lemme \ref{lemma-restrict-to-open}.
Au prix de quelques efforts, on peut aussi définir l'image réciproque plate et l'image directe propre
pour le premier groupe de Chow supérieur $\CH^M_k(X, 1)$. Nous y reviendrons
plus loin (insérer ici un renvoi ultérieur).
\end{remark}






\section{Intersection avec un faisceau inversible et équivalence rationnelle}
\label{section-commutativity}

\noindent
L'application du lemme fondamental donne les propriétés essentielles de l'intersection
avec les faisceaux inversibles. Nous verrons en particulier que
$c_1(\mathcal{L}) \cap -$ passe au quotient par l'équivalence rationnelle et
que ces opérations commutent pour différents faisceaux inversibles.

\begin{lemma}
\label{lemma-commutativity-on-integral}
Soit $(S, \delta)$ comme dans la situation \ref{situation-setup}.
Soit $X$ localement de type fini sur $S$.
Supposons $X$ intègre et $\dim_\delta(X) = n$.
Soient $\mathcal{L}$, $\mathcal{N}$ des faisceaux inversibles sur $X$.
Choisissons une section méromorphe non nulle $s$ de $\mathcal{L}$
et une section méromorphe non nulle $t$ de $\mathcal{N}$.
Posons $\alpha = \text{div}_\mathcal{L}(s)$ et
$\beta = \text{div}_\mathcal{N}(t)$.
Alors
$$
c_1(\mathcal{N}) \cap \alpha
=
c_1(\mathcal{L}) \cap \beta
$$
dans $\CH_{n - 2}(X)$.
\end{lemma}

\begin{proof}
Cela découle immédiatement du lemme fondamental \ref{lemma-key-formula}
et de la discussion qui le précède.
\end{proof}

\begin{lemma}
\label{lemma-factors}
Soit $(S, \delta)$ comme dans la situation \ref{situation-setup}.
Soit $X$ localement de type fini sur $S$.
Soit $\mathcal{L}$ un faisceau inversible sur $X$.
L'opération $\alpha \mapsto c_1(\mathcal{L}) \cap \alpha$
passe au quotient par l'équivalence rationnelle et donne une opération
$$
c_1(\mathcal{L}) \cap - : \CH_{k + 1}(X) \to \CH_k(X)
$$
\end{lemma}

\begin{proof}
Soit $\alpha \in Z_{k + 1}(X)$, avec $\alpha \sim_{rat} 0$.
Il faut montrer que $c_1(\mathcal{L}) \cap \alpha$,
au sens de la définition \ref{definition-cap-c1}, est nul.
D'après la définition \ref{definition-rational-equivalence}, il
existe une famille localement finie $\{W_j\}$ de sous-schémas fermés
intègres tels que $\dim_\delta(W_j) = k + 2$, et des fonctions rationnelles
$f_j \in R(W_j)^*$ telles que
$$
\alpha = \sum (i_j)_*\text{div}_{W_j}(f_j)
$$
Remarquons que $p : \coprod W_j \to X$ est un morphisme propre ;
on a donc $\alpha = p_*\alpha'$, où $\alpha' \in Z_{k + 1}(\coprod W_j)$
est la somme des diviseurs principaux $\text{div}_{W_j}(f_j)$.
D'après le lemme \ref{lemma-pushforward-cap-c1}, on a
$c_1(\mathcal{L}) \cap \alpha = p_*(c_1(p^*\mathcal{L}) \cap \alpha')$.
Il suffit donc de montrer que chacun des
$c_1(\mathcal{L}|_{W_j}) \cap \text{div}_{W_j}(f_j)$ est nul.
Autrement dit, on peut supposer $X$ intègre et
$\alpha = \text{div}_X(f)$ pour un certain $f \in R(X)^*$.

\medskip\noindent
Supposons $X$ intègre et $\alpha = \text{div}_X(f)$ pour un certain $f \in R(X)^*$.
On peut considérer $f$ comme une section méromorphe régulière du faisceau
inversible $\mathcal{N} = \mathcal{O}_X$. Choisissons une section méromorphe
$s$ de $\mathcal{L}$ et notons $\beta = \text{div}_\mathcal{L}(s)$.
D'après le lemme \ref{lemma-commutativity-on-integral},
on en conclut que
$$
c_1(\mathcal{L}) \cap \alpha = c_1(\mathcal{O}_X) \cap \beta.
$$
Or le lemme \ref{lemma-c1-cap-additive} montre que le membre de droite
est nul dans $\CH_k(X)$, comme voulu.
\end{proof}

\noindent
Soit $(S, \delta)$ comme dans la situation \ref{situation-setup}.
Soit $X$ localement de type fini sur $S$.
Soit $\mathcal{L}$ un faisceau inversible sur $X$.
Nous noterons
$$
c_1(\mathcal{L}) \cap - : \CH_{k + 1}(X) \to \CH_k(X)
$$
l'opération $c_1(\mathcal{L}) \cap - $. Cette notation a un sens d'après le
lemme \ref{lemma-factors}. Nous noterons $c_1(\mathcal{L})^s \cap -$
la $s$-ième itérée de cette opération pour tout $s \geq 0$.

\begin{lemma}
\label{lemma-cap-commutative}
Soit $(S, \delta)$ comme dans la situation \ref{situation-setup}.
Soit $X$ localement de type fini sur $S$.
Soient $\mathcal{L}$, $\mathcal{N}$ des faisceaux inversibles sur $X$.
Pour tout $\alpha \in \CH_{k + 2}(X)$, on a
$$
c_1(\mathcal{L}) \cap c_1(\mathcal{N}) \cap \alpha
=
c_1(\mathcal{N}) \cap c_1(\mathcal{L}) \cap \alpha
$$
comme éléments de $\CH_k(X)$.
\end{lemma}

\begin{proof}
Écrivons $\alpha = \sum m_j[Z_j]$ pour une famille localement finie
de sous-schémas fermés intègres $Z_j \subset X$
tels que $\dim_\delta(Z_j) = k + 2$.
Considérons le morphisme propre $p : \coprod Z_j \to X$.
Posons $\alpha' = \sum m_j[Z_j]$, considéré comme un $(k + 2)$-cycle sur
$\coprod Z_j$. Plusieurs applications du
lemme \ref{lemma-pushforward-cap-c1} donnent
$c_1(\mathcal{L}) \cap c_1(\mathcal{N}) \cap \alpha
= p_*(c_1(p^*\mathcal{L}) \cap c_1(p^*\mathcal{N}) \cap \alpha')$
et
$c_1(\mathcal{N}) \cap c_1(\mathcal{L}) \cap \alpha
= p_*(c_1(p^*\mathcal{N}) \cap c_1(p^*\mathcal{L}) \cap \alpha')$.
Il suffit donc de démontrer la formule lorsque $X$ est intègre
et $\alpha = [X]$. Dans ce cas, le résultat découle
du lemme \ref{lemma-commutativity-on-integral} et des définitions.
\end{proof}






\section{Homomorphismes de Gysin}
\label{section-intersecting-effective-Cartier}

\noindent
Dans cette section, nous définissons l'application de Gysin pour le lieu des zéros $D$ d'une
section d'un faisceau inversible. Le cas où $D$
est un diviseur de Cartier effectif est intéressant, mais la généralisation à un $D$ quelconque
donne la souplesse nécessaire pour formuler diverses compatibilités ; voir la
remarque \ref{remark-pullback-pairs} et les
lemmes \ref{lemma-closed-in-X-gysin}, \ref{lemma-gysin-flat-pullback} et
\ref{lemma-gysin-commutes-gysin}.
Ces résultats se généralisent aux sous-schémas fermés localement principaux
munis d'un fibré normal virtuel
(remarque \ref{remark-generalize-to-virtual}) ou aux
pseudo-diviseurs (remarque \ref{remark-generalize-to-pseudo-divisor}).

\medskip\noindent
Rappelons que les diviseurs de Cartier effectifs correspondent $1$ à $1$ aux
classes d'isomorphisme de couples $(\mathcal{L}, s)$, où $\mathcal{L}$
est un faisceau inversible et $s$ une section globale régulière ; voir
Diviseurs, lemme \ref{divisors-lemma-characterize-OD}.
Si $D$ correspond à $(\mathcal{L}, s)$, alors
$\mathcal{L} = \mathcal{O}_X(D)$. On gardera ce fait à l'esprit
dans la lecture de cette section.

\begin{definition}
\label{definition-gysin-homomorphism}
Soit $(S, \delta)$ comme dans la situation \ref{situation-setup}.
Soit $X$ localement de type fini sur $S$.
Soit $(\mathcal{L}, s)$ un couple formé d'un faisceau
inversible et d'une section globale $s \in \Gamma(X, \mathcal{L})$.
Soit $D = Z(s)$ le schéma des zéros de $s$, et
notons $i : D \to X$ l'immersion fermée.
Nous définissons, pour tout entier $k$, un {\it homomorphisme de Gysin}
$$
i^* : Z_{k + 1}(X) \to \CH_k(D).
$$
par les règles suivantes :
\begin{enumerate}
\item Étant donné un sous-schéma fermé intègre $W \subset X$ tel que
$\dim_\delta(W) = k + 1$, on définit
\begin{enumerate}
\item si $W \not \subset D$, alors $i^*[W] = [D \cap W]_k$, considéré comme
$k$-cycle sur $D$, et
\item si $W \subset D$, alors
$i^*[W] = i'_*(c_1(\mathcal{L}|_W) \cap [W])$,
où $i' : W \to D$ est l'immersion fermée induite.
\end{enumerate}
\item Pour un $(k + 1)$-cycle quelconque $\alpha = \sum n_j[W_j]$,
on pose
$$
i^*\alpha = \sum n_j i^*[W_j]
$$
\item Si $D$ est un diviseur de Cartier effectif, on note
$D \cdot \alpha = i_*i^*\alpha$ l'image directe de la classe $i^*\alpha$,
considérée comme classe sur $X$.
\end{enumerate}
\end{definition}

\noindent
En fait, comme nous le verrons plus loin, cet homomorphisme de Gysin $i^*$ peut être considéré
comme un exemple d'image réciproque non plate. Nous appellerons donc parfois, de manière informelle,
la classe $i^*\alpha$ l'{\it image réciproque} de la classe $\alpha$.

\begin{remark}
\label{remark-generalize-to-virtual}
Soit $X$ un schéma localement de type fini sur $S$ comme dans la
situation \ref{situation-setup}. Soit $(D, \mathcal{N}, \sigma)$
un triplet formé d'un sous-schéma fermé localement principal (Diviseurs, définition
\ref{divisors-definition-effective-Cartier-divisor})
$i : D \to X$, d'un $\mathcal{O}_D$-module inversible $\mathcal{N}$ et
d'une surjection $\sigma : \mathcal{N}^{\otimes -1} \to i^*\mathcal{I}_D$
de $\mathcal{O}_D$-modules\footnote{Cette condition assure que, si
$D$ est un diviseur de Cartier effectif, alors $\mathcal{N} = \mathcal{O}_X(D)|_D$.}.
Il faut ici considérer $\mathcal{N}$ comme
un {\it fibré normal virtuel de $D$ dans $X$}. La construction de
$i^* : Z_{k + 1}(X) \to \CH_k(D)$ dans la
définition \ref{definition-gysin-homomorphism}
se généralise à de tels triplets ; voir la
section \ref{section-gysin-higher-codimension}.
\end{remark}

\begin{remark}
\label{remark-generalize-to-pseudo-divisor}
Soit $X$ un schéma localement de type fini sur $S$ comme dans la
situation \ref{situation-setup}. Dans \cite{F}, un {\it pseudo-diviseur} sur $X$
est défini comme un triplet $D = (\mathcal{L}, Z, s)$, où $\mathcal{L}$
est un $\mathcal{O}_X$-module inversible, $Z \subset X$ une partie fermée,
et $s \in \Gamma(X \setminus Z, \mathcal{L})$ une section qui ne s'annule
nulle part. Comme ci-dessus, on peut définir pour tout $\alpha$
dans $\CH_{k + 1}(X)$ un produit $D \cdot \alpha$ dans $\CH_k(Z \cap |\alpha|)$,
où $|\alpha|$ est le support de $\alpha$.
\end{remark}

\begin{lemma}
\label{lemma-support-cap-effective-Cartier}
Soit $(S, \delta)$ comme dans la situation \ref{situation-setup}. Soit $X$ localement
de type fini sur $S$. Soit $(\mathcal{L}, s, i : D \to X)$ comme dans la
définition \ref{definition-gysin-homomorphism}. Soit $\alpha$ un
$(k + 1)$-cycle sur $X$. Alors $i_*i^*\alpha = c_1(\mathcal{L}) \cap \alpha$
dans $\CH_k(X)$. En particulier, si $D$ est un diviseur de Cartier effectif, alors
$D \cdot \alpha = c_1(\mathcal{O}_X(D)) \cap \alpha$.
\end{lemma}

\begin{proof}
Écrivons $\alpha = \sum n_j[W_j]$, où $i_j : W_j \to X$ sont des sous-schémas fermés
intègres tels que $\dim_\delta(W_j) = k$.
Puisque $D$ est le schéma des zéros de $s$, on voit que $D \cap W_j$ est le schéma des zéros
de la restriction $s|_{W_j}$. Ainsi, pour tout $j$ tel que
$W_j \not \subset D$, on a
$c_1(\mathcal{L}) \cap [W_j] = [D \cap W_j]_k$
d'après le lemme \ref{lemma-geometric-cap}. On a donc
$$
c_1(\mathcal{L}) \cap \alpha
=
\sum\nolimits_{W_j \not \subset D} n_j[D \cap W_j]_k
+
\sum\nolimits_{W_j \subset D}
n_j i_{j, *}(c_1(\mathcal{L})|_{W_j}) \cap [W_j])
$$
dans $\CH_k(X)$, d'après la définition \ref{definition-cap-c1}.
Le membre de droite coïncide terme à terme avec l'image directe de la classe
$i^*\alpha$ sur $D$ de la définition \ref{definition-gysin-homomorphism}.
On obtient donc le résultat.
\end{proof}

\begin{lemma}
\label{lemma-easy-gysin}
Soit $(S, \delta)$ comme dans la situation \ref{situation-setup}.
Soit $X$ localement de type fini sur $S$.
Soit $(\mathcal{L}, s, i : D \to X)$ comme dans la
définition \ref{definition-gysin-homomorphism}.
\begin{enumerate}
\item Soit $Z \subset X$ un sous-schéma fermé tel
que $\dim_\delta(Z) \leq k + 1$ et tel que
$D \cap Z$ soit un diviseur de Cartier effectif sur $Z$. Alors
$i^*[Z]_{k + 1} = [D \cap Z]_k$.
\item Soit $\mathcal{F}$ un faisceau cohérent sur $X$
tel que $\dim_\delta(\text{Supp}(\mathcal{F})) \leq k + 1$ et que
$s : \mathcal{F} \to \mathcal{F} \otimes_{\mathcal{O}_X} \mathcal{L}$
soit injectif. Alors
$$
i^*[\mathcal{F}]_{k + 1} = [i^*\mathcal{F}]_k
$$
dans $\CH_k(D)$.
\end{enumerate}
\end{lemma}

\begin{proof}
Soit $Z \subset X$ comme dans (1). Posons $\mathcal{F} = \mathcal{O}_Z$.
L'hypothèse selon laquelle $D \cap Z$ est un diviseur de Cartier effectif
équivaut à l'injectivité de
$s : \mathcal{F} \to \mathcal{F} \otimes_{\mathcal{O}_X} \mathcal{L}$.
De plus, $[Z]_{k + 1} = [\mathcal{F}]_{k + 1}]$
et $[D \cap Z]_k = [\mathcal{O}_{D \cap Z}]_k = [i^*\mathcal{F}]_k$.
Voir le lemme \ref{lemma-cycle-closed-coherent}.
L'assertion (1) découle donc de (2).

\medskip\noindent
Écrivons $[\mathcal{F}]_{k + 1} = \sum m_j[W_j]$, avec $m_j > 0$
et des sous-schémas fermés intègres deux à deux distincts $W_j \subset X$
de $\delta$-dimension $k + 1$. L'hypothèse d'injectivité de
$s : \mathcal{F} \to \mathcal{F} \otimes_{\mathcal{O}_X} \mathcal{L}$
entraîne que $W_j \not \subset D$ pour tout $j$.
Par définition, on a
$$
i^*[\mathcal{F}]_{k + 1} = \sum m_j [D \cap W_j]_k.
$$
Nous affirmons que
$$
\sum [D \cap W_j]_k = [i^*\mathcal{F}]_k
$$
comme cycles.
Soit $Z \subset D$ un sous-schéma fermé intègre de $\delta$-dimension
$k$. Soit $\xi \in Z$ son point générique. Soit $A = \mathcal{O}_{X, \xi}$.
Soit $M = \mathcal{F}_\xi$. Soit $f \in A$ un élément engendrant
l'idéal de $D$, c'est-à-dire tel que $\mathcal{O}_{D, \xi} = A/fA$.
Par hypothèse, $\dim(\text{Supp}(M)) = 1$,
l'application $f : M \to M$ est injective, et
$\text{longueur}_A(M/fM) < \infty$. De plus, $\text{longueur}_A(M/fM)$
est le coefficient de $[Z]$ dans $[i^*\mathcal{F}]_k$. D'autre
part, soient $\mathfrak q_1, \ldots, \mathfrak q_t$ les idéaux premiers minimaux
du support de $M$. Alors
$$
\sum
\text{longueur}_{A_{\mathfrak q_i}}(M_{\mathfrak q_i})
\text{ord}_{A/\mathfrak q_i}(f)
$$
est le coefficient de $[Z]$ dans $\sum [D \cap W_j]_k$.
L'égalité résulte donc du
lemme \ref{lemma-additivity-divisors-restricted}.
\end{proof}

\begin{remark}
\label{remark-gysin-on-cycles}
Soient $X \to S$, $\mathcal{L}$, $s$, $i : D \to X$ comme dans la
définition \ref{definition-gysin-homomorphism}, et supposons
que $\mathcal{L}|_D \cong \mathcal{O}_D$. Dans ce cas, on
peut définir une application canonique $i^* : Z_{k + 1}(X) \to Z_k(D)$
sur les cycles, en imposant $i^*[W] = 0$ lorsque $W \subset D$
est un sous-schéma fermé intègre.
Cette possibilité nous sera utile par la suite.
\end{remark}

\begin{remark}
\label{remark-pullback-pairs}
Soit $f : X' \to X$ un morphisme de schémas localement de type fini sur $S$
comme dans la situation \ref{situation-setup}. Soit $(\mathcal{L}, s, i : D \to X)$
un triplet comme dans la définition \ref{definition-gysin-homomorphism}.
On peut alors poser $\mathcal{L}' = f^*\mathcal{L}$, $s' = f^*s$ et
$D' = X' \times_X D = Z(s')$. Cela donne un diagramme commutatif
$$
\xymatrix{
D' \ar[d]_g \ar[r]_{i'} & X' \ar[d]^f \\
D \ar[r]^i & X
}
$$
et l'on peut étudier diverses compatibilités entre $i^*$ et $(i')^*$.
\end{remark}

\begin{lemma}
\label{lemma-closed-in-X-gysin}
Soit $(S, \delta)$ comme dans la situation \ref{situation-setup}.
Soit $f : X' \to X$ un morphisme propre de schémas
localement de type fini sur $S$.
Soit $(\mathcal{L}, s, i : D \to X)$ comme dans la
définition \ref{definition-gysin-homomorphism}.
Formons le diagramme
$$
\xymatrix{
D' \ar[d]_g \ar[r]_{i'} & X' \ar[d]^f \\
D \ar[r]^i & X
}
$$
comme dans la remarque \ref{remark-pullback-pairs}.
Pour tout $(k + 1)$-cycle $\alpha'$ sur $X'$, on a
$i^*f_*\alpha' = g_*(i')^*\alpha'$ dans $\CH_k(D)$
(ce qui a un sens, puisque $f_*$ est défini au niveau des cycles).
\end{lemma}

\begin{proof}
Supposons $\alpha = [W']$ pour un sous-schéma fermé intègre
$W' \subset X'$. Soit $W = f(W') \subset X$. Si $W' \not \subset D'$,
alors $W \not \subset D$, et l'on voit que
$$
[W' \cap D']_k = \text{div}_{\mathcal{L}'|_{W'}}({s'|_{W'}})
\quad\text{et}\quad
[W \cap D]_k = \text{div}_{\mathcal{L}|_W}(s|_W)
$$
et donc que l'image du premier cycle par $f_*$ est égale au second cycle, d'après le
lemme \ref{lemma-equal-c1-as-cycles}. L'égalité
est donc vraie comme égalité de cycles. Si $W' \subset D'$, alors
$W \subset D$ et $f_*(c_1(\mathcal{L}|_{W'}) \cap [W'])$
est égal à $c_1(\mathcal{L}|_W) \cap [W]$ dans $\CH_k(W)$, d'après la seconde
assertion du lemme \ref{lemma-equal-c1-as-cycles}.
D'après la remarque \ref{remark-infinite-sums-rational-equivalences},
le résultat s'ensuit pour un $\alpha'$ quelconque.
\end{proof}

\begin{lemma}
\label{lemma-gysin-flat-pullback}
Soit $(S, \delta)$ comme dans la situation \ref{situation-setup}. Soit $f : X' \to X$
un morphisme plat de dimension relative $r$ entre schémas localement de type fini
sur $S$. Soit $(\mathcal{L}, s, i : D \to X)$ comme dans la
définition \ref{definition-gysin-homomorphism}. Formons le diagramme
$$
\xymatrix{
D' \ar[d]_g \ar[r]_{i'} & X' \ar[d]^f \\
D \ar[r]^i & X
}
$$
comme dans la remarque \ref{remark-pullback-pairs}.
Pour tout $(k + 1)$-cycle $\alpha$ sur $X$, on a
$(i')^*f^*\alpha = g^*i^*\alpha$ dans $\CH_{k + r}(D')$
(ce qui a un sens, puisque $f^*$ est défini au niveau des cycles).
\end{lemma}

\begin{proof}
Supposons $\alpha = [W]$ pour un sous-schéma fermé intègre
$W \subset X$. Soit $W' = f^{-1}(W) \subset X'$. Si $W \not \subset D$,
alors $W' \not \subset D'$, et l'on voit que
$$
W' \cap D' = g^{-1}(W \cap D)
$$
comme sous-schémas fermés de $D'$. L'égalité
est donc vraie comme égalité de cycles ; voir le lemme \ref{lemma-pullback-coherent}.
Si $W \subset D$, alors $W' \subset D'$ et $W' = g^{-1}(W)$,
avec $[W']_{k + 1 + r} = g^*[W]$, et l'égalité est vraie dans
$\CH_{k + r}(D')$ d'après le lemme \ref{lemma-flat-pullback-cap-c1}.
D'après la remarque \ref{remark-infinite-sums-rational-equivalences},
le résultat s'ensuit pour un $\alpha'$ quelconque.
\end{proof}








\section{Homomorphismes de Gysin et équivalence rationnelle}
\label{section-gysin}

\noindent
Dans cette section, nous utilisons la formule fondamentale pour montrer que l'homomorphisme
de Gysin passe au quotient par l'équivalence rationnelle. Nous démontrons aussi une importante
propriété de commutativité.

\begin{lemma}
\label{lemma-gysin-factors-general}
Soit $(S, \delta)$ comme dans la situation \ref{situation-setup}.
Soit $X$ localement de type fini sur $S$.
Supposons $X$ intègre et $n = \dim_\delta(X)$.
Soit $i : D \to X$ un diviseur de Cartier effectif.
Soit $\mathcal{N}$ un $\mathcal{O}_X$-module inversible,
et soit $t$ une section méromorphe non nulle de $\mathcal{N}$.
Alors $i^*\text{div}_\mathcal{N}(t) = c_1(\mathcal{N}|_D) \cap [D]_{n - 1}$
dans $\CH_{n - 2}(D)$.
\end{lemma}

\begin{proof}
Écrivons $\text{div}_\mathcal{N}(t) = \sum \text{ord}_{Z_i, \mathcal{N}}(t)[Z_i]$
pour des sous-schémas fermés intègres $Z_i \subset X$ de $\delta$-dimension
$n - 1$. On peut supposer que la famille $\{Z_i\}$ est localement
finie, que $t \in \Gamma(U, \mathcal{N}|_U)$ est un générateur,
où $U = X \setminus \bigcup Z_i$, et que toute composante irréductible
de $D$ est l'un des $Z_i$ ; voir
Diviseurs, lemmes \ref{divisors-lemma-components-locally-finite},
\ref{divisors-lemma-divisor-locally-finite} et
\ref{divisors-lemma-divisor-meromorphic-locally-finite}.

\medskip\noindent
Posons $\mathcal{L} = \mathcal{O}_X(D)$. Notons
$s \in \Gamma(X, \mathcal{O}_X(D)) = \Gamma(X, \mathcal{L})$
la section canonique. Nous allons appliquer la discussion de la
section \ref{section-key} à la situation présente.
Pour chaque $i$, soit $\xi_i \in Z_i$ son point générique. Soit
$B_i = \mathcal{O}_{X, \xi_i}$. Pour chaque $i$, choisissons des générateurs
$s_i \in \mathcal{L}_{\xi_i}$ et $t_i \in \mathcal{N}_{\xi_i}$
sur $B_i$, en imposant toutefois $s_i = s$ si $Z_i \not \subset D$.
Écrivons $s = f_i s_i$ et $t = g_i t_i$, avec $f_i, g_i \in B_i$.
Alors $\text{ord}_{Z_i, \mathcal{N}}(t) = \text{ord}_{B_i}(g_i)$.
D'autre part, on a $f_i \in B_i$ et
$$
[D]_{n - 1} = \sum \text{ord}_{B_i}(f_i)[Z_i]
$$
d'après le choix des $s_i$. Nous affirmons que
$$
i^*\text{div}_\mathcal{N}(t) =
\sum \text{ord}_{B_i}(g_i) \text{div}_{\mathcal{L}|_{Z_i}}(s_i|_{Z_i})
$$
comme cycles. Plus précisément, le membre de droite est un cycle
représentant le membre de gauche. Cela résulte en effet de notre
formule pour $\text{div}_\mathcal{N}(t)$ et du fait que
$\text{div}_{\mathcal{L}|_{Z_i}}(s_i|_{Z_i}) = [Z(s_i|_{Z_i})]_{n - 2} =
[Z_i \cap D]_{n - 2}$ lorsque $Z_i \not \subset D$, car dans
ce cas $s_i|_{Z_i} = s|_{Z_i}$ est une section régulière ; voir le
lemme \ref{lemma-compute-c1}. De même,
$$
c_1(\mathcal{N}) \cap [D]_{n - 1} =
\sum \text{ord}_{B_i}(f_i) \text{div}_{\mathcal{N}|_{Z_i}}(t_i|_{Z_i})
$$
La formule fondamentale (lemme \ref{lemma-key-formula}) donne l'égalité
$$
\sum \left(
\text{ord}_{B_i}(f_i) \text{div}_{\mathcal{N}|_{Z_i}}(t_i|_{Z_i}) -
\text{ord}_{B_i}(g_i) \text{div}_{\mathcal{L}|_{Z_i}}(s_i|_{Z_i}) \right) =
\sum \text{div}_{Z_i}(\partial_{B_i}(f_i, g_i))
$$
de cycles. Si $Z_i \not \subset D$, alors $f_i = 1$, et donc
$\text{div}_{Z_i}(\partial_{B_i}(f_i, g_i)) = 0$. On obtient ainsi une équivalence
rationnelle entre les cycles particuliers représentant
$i^*\text{div}_\mathcal{N}(t)$ et $c_1(\mathcal{N}) \cap [D]_{n - 1}$
sur $D$. Cela achève la démonstration.
\end{proof}

\begin{lemma}
\label{lemma-gysin-factors}
Soit $(S, \delta)$ comme dans la situation \ref{situation-setup}.
Soit $X$ localement de type fini sur $S$.
Soit $(\mathcal{L}, s, i : D \to X)$ comme dans la
définition \ref{definition-gysin-homomorphism}.
L'homomorphisme de Gysin passe au quotient par l'équivalence rationnelle et
donne une application $i^* : \CH_{k + 1}(X) \to \CH_k(D)$.
\end{lemma}

\begin{proof}
Soit $\alpha \in Z_{k + 1}(X)$, et supposons $\alpha \sim_{rat} 0$.
Cela signifie qu'il existe une famille localement finie de sous-schémas
fermés intègres $W_j \subset X$ de $\delta$-dimension $k + 2$
et des $f_j \in R(W_j)^*$ tels que
$\alpha = \sum i_{j, *}\text{div}_{W_j}(f_j)$.
Posons $X' = \coprod W_i$ et considérons le diagramme
$$
\xymatrix{
D' \ar[d]_q \ar[r]_{i'} & X' \ar[d]^p \\
D \ar[r]^i & X
}
$$
de la remarque \ref{remark-pullback-pairs}. Puisque $X' \to X$ est propre,
on a $i^*p_* = q_*(i')^*$ d'après le lemme \ref{lemma-closed-in-X-gysin}.
Comme on sait que $q_*$ passe au quotient par l'équivalence rationnelle
(lemme \ref{lemma-proper-pushforward-rational-equivalence}), il suffit
de démontrer le résultat pour $\alpha' = \sum \text{div}_{W_j}(f_j)$
sur $X'$. Cela nous ramène clairement au cas où $X$ est intègre
et $\alpha = \text{div}(f)$ pour un certain $f \in R(X)^*$.

\medskip\noindent
Supposons $X$ intègre et $\alpha = \text{div}(f)$ pour un certain $f \in R(X)^*$.
Si $X = D$, on voit que $i^*\alpha$ est égal
à $c_1(\mathcal{L}) \cap \alpha$.
Cette classe est rationnellement équivalente à zéro d'après le lemme \ref{lemma-factors}.
Si $D \not = X$, on voit que $i^*\text{div}_X(f)$ est égal à
$c_1(\mathcal{O}_D) \cap [D]_{n - 1}$ dans $\CH_{n - 2}(D)$ d'après le
lemme \ref{lemma-gysin-factors-general}. Bien entendu,
le produit cap avec $c_1(\mathcal{O}_D)$ est l'application nulle
(lemme \ref{lemma-c1-cap-additive}).
\end{proof}

\begin{lemma}
\label{lemma-gysin-back}
Soit $(S, \delta)$ comme dans la situation \ref{situation-setup}. Soit $X$ localement
de type fini sur $S$. Soit $(\mathcal{L}, s, i : D \to X)$ comme dans la
définition \ref{definition-gysin-homomorphism}. Alors
$i^*i_* : \CH_k(D) \to \CH_{k - 1}(D)$ envoie $\alpha$ sur
$c_1(\mathcal{L}|_D) \cap \alpha$.
\end{lemma}

\begin{proof}
Cela découle immédiatement de la définition de $i_*$ sur les cycles
et de celle de $i^*$ donnée dans la
définition \ref{definition-gysin-homomorphism}.
\end{proof}

\begin{lemma}
\label{lemma-gysin-commutes-cap-c1}
Soit $(S, \delta)$ comme dans la situation \ref{situation-setup}. Soit $X$
localement de type fini sur $S$. Soit $(\mathcal{L}, s, i : D \to X)$
un triplet comme dans la définition \ref{definition-gysin-homomorphism}.
Soit $\mathcal{N}$ un $\mathcal{O}_X$-module inversible.
Alors $i^*(c_1(\mathcal{N}) \cap \alpha) = c_1(i^*\mathcal{N}) \cap i^*\alpha$
dans $\CH_{k - 2}(D)$ pour tout $\alpha \in \CH_k(X)$.
\end{lemma}

\begin{proof}
Par une démonstration identique à celle du lemme \ref{lemma-gysin-factors},
cela résulte des lemmes
\ref{lemma-pushforward-cap-c1},
\ref{lemma-cap-commutative} et
\ref{lemma-gysin-factors-general}.
\end{proof}

\begin{lemma}
\label{lemma-gysin-commutes-gysin}
Soit $(S, \delta)$ comme dans la situation \ref{situation-setup}. Soit $X$ localement
de type fini sur $S$. Soient $(\mathcal{L}, s, i : D \to X)$ et
$(\mathcal{L}', s', i' : D' \to X)$ deux triplets comme dans la
définition \ref{definition-gysin-homomorphism}. Alors le diagramme
$$
\xymatrix{
\CH_k(X) \ar[r]_{i^*} \ar[d]_{(i')^*} & \CH_{k - 1}(D) \ar[d]^{j^*} \\
\CH_{k - 1}(D') \ar[r]^{(j')^*} & \CH_{k - 2}(D \cap D')
}
$$
est commutatif, chacune des applications étant une application de Gysin.
\end{lemma}

\begin{proof}
Notons $j : D \cap D' \to D$ et $j' : D \cap D' \to D'$ les immersions
fermées correspondant à $(\mathcal{L}|_{D'}, s|_{D'}$ et
$(\mathcal{L}'_D, s|_D)$. Il faut montrer que
$(j')^*i^*\alpha = j^* (i')^*\alpha$ pour tout $\alpha \in \CH_k(X)$.
Soit $W \subset X$ un sous-schéma fermé intègre de dimension $k$.
Démontrons l'égalité lorsque $\alpha = [W]$. Nous la déduirons
de la formule fondamentale.

\medskip\noindent
Soit $\sigma$ une section méromorphe non nulle de $\mathcal{L}|_W$,
que l'on impose égale à $s|_W$ si $W \not \subset D$.
Soit $\sigma'$ une section méromorphe non nulle de $\mathcal{L}'|_W$,
que l'on impose égale à $s'|_W$ si $W \not \subset D'$.
Écrivons
$$
\text{div}_{\mathcal{L}|_W}(\sigma) =
\sum \text{ord}_{Z_i, \mathcal{L}|_W}(\sigma)[Z_i] = \sum n_i[Z_i]
$$
et, de même,
$$
\text{div}_{\mathcal{L}'|_W}(\sigma') =
\sum \text{ord}_{Z_i, \mathcal{L}'|_W}(\sigma')[Z_i] = \sum n'_i[Z_i]
$$
comme dans la discussion de la section \ref{section-key}.
On voit alors que $Z_i \subset D$ si $n_i \not = 0$, et
$Z'_i \subset D'$ si $n'_i \not = 0$. Pour chaque $i$, soit $\xi_i \in Z_i$
le point générique. Comme dans la section \ref{section-key}, choisissons
pour chaque $i$ un élément
$\sigma_i \in \mathcal{L}_{\xi_i}$, resp.\ $\sigma'_i \in \mathcal{L}'_{\xi_i}$,
qui soit un générateur sur $B_i = \mathcal{O}_{W, \xi_i}$
et qui soit égal à l'image de
$s$, resp.\ $s'$ si $Z_i \not \subset D$, resp.\ $Z_i \not \subset D'$.
Écrivons $\sigma = f_i \sigma_i$ et $\sigma' = f'_i\sigma'_i$, de sorte que
$n_i = \text{ord}_{B_i}(f_i)$ et
$n'_i = \text{ord}_{B_i}(f'_i)$.
Il résulte de nos définitions que
$$
(j')^*i^*[W] =
\sum \text{ord}_{B_i}(f_i) \text{div}_{\mathcal{L}'|_{Z_i}}(\sigma'_i|_{Z_i})
$$
comme cycles, et que
$$
j^*(i')^*[W] =
\sum \text{ord}_{B_i}(f'_i) \text{div}_{\mathcal{L}|_{Z_i}}(\sigma_i|_{Z_i})
$$
La formule fondamentale (lemme \ref{lemma-key-formula}) donne maintenant l'égalité
$$
\sum \left(
\text{ord}_{B_i}(f_i) \text{div}_{\mathcal{L}'|_{Z_i}}(\sigma'_i|_{Z_i}) -
\text{ord}_{B_i}(f'_i) \text{div}_{\mathcal{L}|_{Z_i}}(\sigma_i|_{Z_i})
\right) =
\sum \text{div}_{Z_i}(\partial_{B_i}(f_i, f'_i))
$$
de cycles. Remarquons que $\text{div}_{Z_i}(\partial_{B_i}(f_i, f'_i)) = 0$ si
$Z_i \not \subset D \cap D'$, car, dans ce cas, soit $f_i = 1$,
soit $f'_i = 1$. On obtient ainsi une équivalence rationnelle entre nos cycles
particuliers représentant $(j')^*i^*[W]$ et $j^*(i')^*[W]$ sur $D \cap D' \cap W$.
D'après la remarque \ref{remark-infinite-sums-rational-equivalences},
le résultat s'ensuit pour un $\alpha$ quelconque.
\end{proof}
















\section{Diviseurs de Cartier effectifs relatifs}
\label{section-relative-effective-cartier}

\noindent
Les diviseurs de Cartier effectifs relatifs sont définis et étudiés
dans Diviseurs, section \ref{divisors-section-effective-Cartier-morphisms}.
Pour établir les résultats fondamentaux sur les classes de Chern des fibrés vectoriels,
nous n'avons besoin que du cas où le schéma ambiant et le diviseur de Cartier
effectif sont tous deux plats sur la base.

\begin{lemma}
\label{lemma-relative-effective-cartier}
Soit $(S, \delta)$ comme dans la situation \ref{situation-setup}.
Soient $X$, $Y$ localement de type fini sur $S$.
Soit $p : X \to Y$ un morphisme plat de dimension relative $r$.
Soit $i : D \to X$ un diviseur de Cartier effectif relatif
(Diviseurs, définition
\ref{divisors-definition-relative-effective-Cartier-divisor}).
Soit $\mathcal{L} = \mathcal{O}_X(D)$.
Pour tout $\alpha \in \CH_{k + 1}(Y)$, on a
$$
i^*p^*\alpha = (p|_D)^*\alpha
$$
dans $\CH_{k + r}(D)$ et
$$
c_1(\mathcal{L}) \cap p^*\alpha = i_* ((p|_D)^*\alpha)
$$
dans $\CH_{k + r}(X)$.
\end{lemma}

\begin{proof}
Soit $W \subset Y$ un sous-schéma fermé intègre de $\delta$-dimension
$k + 1$. D'après Diviseurs, lemme \ref{divisors-lemma-relative-Cartier},
on voit que $D \cap p^{-1}W$ est un diviseur de Cartier
effectif sur $p^{-1}W$. Le lemme \ref{lemma-easy-gysin}
donne la première égalité dans
$$
i^*[p^{-1}W]_{k + r + 1} =
[D \cap p^{-1}W]_{k + r} =
[(p|_D)^{-1}(W)]_{k + r}.
$$
et la seconde découle de l'égalité $D \cap p^{-1}(W) = (p|_D)^{-1}(W)$ de schémas.
Puisque, par définition, $p^*[W] = [p^{-1}W]_{k + r + 1}$, on voit que
$i^*p^*[W] = (p|_D)^*[W]$ comme cycles. Si $\alpha = \sum m_j[W_j]$ est un
$k + 1$-cycle quelconque, on obtient
$i^*\alpha = \sum m_j i^*p^*[W_j] = \sum m_j(p|_D)^*[W_j]$ comme cycles.
Cela démontre la première égalité. Pour en déduire la seconde,
on applique le lemme \ref{lemma-support-cap-effective-Cartier}.
\end{proof}








\section{Fibrés affines}
\label{section-affine-vector}

\noindent
Pour un fibré affine, l'application d'image réciproque est surjective sur les groupes de Chow.

\begin{lemma}
\label{lemma-pullback-affine-fibres-surjective}
Soit $(S, \delta)$ comme dans la situation \ref{situation-setup}.
Soient $X$, $Y$ localement de type fini sur $S$.
Soit $f : X \to Y$ un morphisme plat de dimension relative $r$.
Supposons que, pour tout $y \in Y$, il existe un voisinage ouvert
$U \subset Y$ tel que $f|_{f^{-1}(U)} : f^{-1}(U) \to U$
s'identifie au morphisme $U \times \mathbf{A}^r \to U$.
Alors $f^* : \CH_k(Y) \to \CH_{k + r}(X)$ est surjectif pour tout
$k \in \mathbf{Z}$.
\end{lemma}

\begin{proof}
Soit $\alpha \in \CH_{k + r}(X)$.
Écrivons $\alpha = \sum m_j[W_j]$, avec $m_j \not = 0$ et
des sous-schémas fermés intègres $W_j$ deux à deux distincts de
$\delta$-dimension $k + r$. Alors la famille $\{W_j\}$
est localement finie dans $X$. Pour tout ouvert quasi-compact
$V \subset Y$, on voit que $f^{-1}(V) \cap W_j$
n'est non vide que pour un nombre fini de $j$. La
famille $Z_j = \overline{f(W_j)}$ des adhérences
des images est donc une famille localement finie de sous-schémas
fermés intègres de $Y$.

\medskip\noindent
Considérons les diagrammes de produit fibré
$$
\xymatrix{
f^{-1}(Z_j) \ar[r] \ar[d]_{f_j} & X \ar[d]^f \\
Z_j \ar[r] & Y
}
$$
Supposons que $[W_j] \in Z_{k + r}(f^{-1}(Z_j))$
soit rationnellement équivalent à $f_j^*\beta_j$ pour un
$k$-cycle $\beta_j \in \CH_k(Z_j)$. Alors
$\beta = \sum m_j \beta_j$ sera un $k$-cycle sur $Y$,
et $f^*\beta = \sum m_j f_j^*\beta_j$ sera rationnellement
équivalent à $\alpha$ (voir la
remarque \ref{remark-infinite-sums-rational-equivalences}).
Cela nous ramène au cas où $Y$ est intègre et
$\alpha = [W]$ pour un sous-schéma fermé intègre
de $X$ dominant $Y$. En particulier, on peut
supposer que $d = \dim_\delta(Y) < \infty$.

\medskip\noindent
On peut donc raisonner par récurrence sur $d = \dim_\delta(Y)$.
Si $d < k$, alors $\CH_{k + r}(X) = 0$ et le lemme est vrai.
Par hypothèse, il existe un ouvert dense $V \subset Y$ tel
que $f^{-1}(V) \cong V \times \mathbf{A}^r$ comme schémas sur $V$.
Supposons que l'on puisse montrer que $\alpha|_{f^{-1}(V)} = f^*\beta$
pour un certain $\beta \in Z_k(V)$. D'après le lemme \ref{lemma-exact-sequence-open},
on voit que
$\beta = \beta'|_V$ pour un certain $\beta' \in Z_k(Y)$.
D'après la suite exacte
$\CH_k(f^{-1}(Y \setminus V)) \to \CH_k(X) \to \CH_k(f^{-1}(V))$
du lemme \ref{lemma-restrict-to-open},
on voit que $\alpha - f^*\beta'$ provient
d'un cycle $\alpha' \in \CH_{k + r}(f^{-1}(Y \setminus V))$.
Comme $\dim_\delta(Y \setminus V) < d$, la récurrence sur
$d$ donne le résultat.

\medskip\noindent
On peut donc supposer que $X = Y \times \mathbf{A}^r$.
Dans ce cas, on peut factoriser $f$ sous la forme
$$
X = Y \times \mathbf{A}^r \to
Y \times \mathbf{A}^{r - 1} \to \ldots \to
Y \times \mathbf{A}^1 \to Y.
$$
Il suffit donc de traiter le cas $r = 1$. L'argument du
deuxième paragraphe de la démonstration nous ramène au cas où
$\alpha = [W]$, $Y$ est intègre et $W \to Y$ dominant.
On peut à nouveau raisonner par récurrence sur $d = \dim_\delta(Y)$.
Si $W = Y \times \mathbf{A}^1$, alors $[W] = f^*[Y]$.
Enfin, si $W \subset Y \times \mathbf{A}^1$ est une inclusion stricte,
alors $W \to Y$ induit une extension finie de corps $R(W)/R(Y)$.
Soit $P(T) \in R(Y)[T]$ le polynôme unitaire irréductible tel
que la fibre générique de $W \to Y$ soit définie par $P$ dans
$\mathbf{A}^1_{R(Y)}$. Soit $V \subset Y$ un ouvert non vide tel
que $P \in \Gamma(V, \mathcal{O}_Y)[T]$ et tel que
$W \cap f^{-1}(V)$ soit encore défini par $P$. On voit alors que
$\alpha|_{f^{-1}(V)} \sim_{rat} 0$, et donc que $\alpha \sim_{rat} \alpha'$
pour un cycle $\alpha'$ sur $(Y \setminus V) \times \mathbf{A}^1$.
La récurrence sur la dimension donne le résultat.
\end{proof}

\begin{lemma}
\label{lemma-linebundle}
Soit $(S, \delta)$ comme dans la situation \ref{situation-setup}.
Soit $X$ localement de type fini sur $S$.
Soit $\mathcal{L}$ un $\mathcal{O}_X$-module inversible.
Soit
$$
p :
L = \underline{\Spec}(\text{Sym}^*(\mathcal{L}))
\longrightarrow
X
$$
le fibré vectoriel associé sur $X$.
Alors $p^* : \CH_k(X) \to \CH_{k + 1}(L)$ est un isomorphisme pour tout $k$.
\end{lemma}

\begin{proof}
Pour la surjectivité, voir le lemme \ref{lemma-pullback-affine-fibres-surjective}.
Soit $o : X \to L$ la section nulle de $L \to X$, c'est-à-dire le morphisme
correspondant à la surjection $\text{Sym}^*(\mathcal{L}) \to \mathcal{O}_X$
qui envoie $\mathcal{L}^{\otimes n}$ sur zéro pour tout $n > 0$.
Alors $p \circ o = \text{id}_X$, et $o(X)$ est un diviseur de Cartier
effectif sur $L$. Le lemme \ref{lemma-relative-effective-cartier}
donne donc $o^* \circ p^* = \text{id}$, et l'on en conclut que $p^*$ est
également injectif.
\end{proof}

\begin{remark}
\label{remark-when-isomorphism}
Nous verrons plus loin (lemme \ref{lemma-vectorbundle}) que, si $X$ est un
fibré vectoriel de rang $r$ sur $Y$, alors l'application d'image réciproque
$\CH_k(Y) \to \CH_{k + r}(X)$
est un isomorphisme. Cela est vrai dès que $X \to Y$ satisfait
aux hypothèses du lemme \ref{lemma-pullback-affine-fibres-surjective} ; voir
\cite[Lemme 2.2]{Totaro-group}. Nous esquisserons une démonstration dans la
remarque \ref{remark-higher-chow-isomorphism} à l'aide des groupes de Chow supérieurs.
\end{remark}

\begin{lemma}
\label{lemma-linebundle-formulae}
Dans la situation du lemme \ref{lemma-linebundle}, notons $o : X \to L$
la section nulle (voir la démonstration du lemme). On a alors
\begin{enumerate}
\item $o(X)$ est le schéma des zéros d'une section globale régulière de
$p^*\mathcal{L}^{\otimes -1}$,
\item $o_* : \CH_k(X) \to \CH_k(L)$, puisque $o$ est une immersion fermée,
\item $o^* : \CH_{k + 1}(L) \to \CH_k(X)$, puisque $o(X)$
est un diviseur de Cartier effectif,
\item $o^* p^* : \CH_k(X) \to \CH_k(X)$ est l'application identique,
\item $o_*\alpha = - p^*(c_1(\mathcal{L}) \cap \alpha)$ pour tout
$\alpha \in \CH_k(X)$, et
\item $o^* o_* : \CH_k(X) \to \CH_{k - 1}(X)$ est égal à l'application
$\alpha \mapsto - c_1(\mathcal{L}) \cap \alpha$.
\end{enumerate}
\end{lemma}

\begin{proof}
Puisque $p_*\mathcal{O}_L = \text{Sym}^*(\mathcal{L})$, on a
$p_*(p^*\mathcal{L}^{\otimes -1}) =
\text{Sym}^*(\mathcal{L}) \otimes_{\mathcal{O}_X} \mathcal{L}^{\otimes -1}$
d'après la formule de projection
(Cohomologie, lemme \ref{cohomology-lemma-projection-formula}),
et la section mentionnée dans (1) est
la trivialisation canonique $\mathcal{O}_X \to
\mathcal{L} \otimes_{\mathcal{O}_X} \mathcal{L}^{\otimes -1}$.
Nous omettons la démonstration du fait que le lieu des zéros de
cette section est précisément $o(X)$. Cela démontre (1).

\medskip\noindent
Nous avons établi les assertions (2), (3) et (4) au cours de la démonstration du
lemme \ref{lemma-linebundle}. Bien entendu, (4) est la première
formule du lemme \ref{lemma-relative-effective-cartier}.

\medskip\noindent
L'assertion (5) découle de la seconde formule du
lemme \ref{lemma-relative-effective-cartier},
de l'additivité du produit cap avec $c_1$ (lemme \ref{lemma-c1-cap-additive})
et du fait que ce produit cap avec $c_1$ commute à l'image réciproque plate
(lemme \ref{lemma-flat-pullback-cap-c1}).

\medskip\noindent
L'assertion (6) découle du lemme \ref{lemma-gysin-back}
et du fait que $o^*p^*\mathcal{L} = \mathcal{L}$.
\end{proof}

\begin{lemma}
\label{lemma-decompose-section}
Soit $Y$ un schéma. Soient $\mathcal{L}_i$, $i = 1, 2$ des
$\mathcal{O}_Y$-modules inversibles. Soit $s$ une section globale de
$\mathcal{L}_1 \otimes_{\mathcal{O}_Y} \mathcal{L}_2$.
Notons $i : D \to Y$ le schéma des zéros de $s$.
Il existe alors un diagramme commutatif
$$
\xymatrix{
D_1 \ar[r]_{i_1} \ar[d]_{p_1} &
L \ar[d]^p &
D_2 \ar[l]^{i_2} \ar[d]^{p_2} \\
D \ar[r]^i &
Y &
D \ar[l]_i
}
$$
et des sections $s_i$ de $p^*\mathcal{L}_i$ telles que
les assertions suivantes soient vérifiées :
\begin{enumerate}
\item $p^*s = s_1 \otimes s_2$,
\item $p$ est de type fini et plat de dimension relative $1$,
\item $D_i$ est le schéma des zéros de $s_i$,
\item $D_i \cong
\underline{\Spec}(\text{Sym}^*(\mathcal{L}_{3 - i}^{\otimes -1})|_D))$
sur $D$ pour $i = 1, 2$,
\item $p^{-1}D = D_1 \cup D_2$ (réunion schématique),
\item $D_1 \cap D_2$ (intersection schématique) s'envoie
isomorphiquement sur $D$, et
\item $D_1 \cap D_2 \to D_i$
est la section nulle du fibré en droites $D_i \to D$ pour $i = 1, 2$.
\end{enumerate}
De plus, la formation de ce diagramme et des sections $s_i$
commute à tout changement de base.
\end{lemma}

\begin{proof}
Soit $p : X \to Y$ le spectre relatif du faisceau quasi-cohérent
de $\mathcal{O}_Y$-algèbres
$$
\mathcal{A} =
\left(\bigoplus\nolimits_{a_1, a_2 \geq 0}
\mathcal{L}_1^{\otimes -a_1} \otimes_{\mathcal{O}_Y}
\mathcal{L}_2^{\otimes -a_2}\right)/\mathcal{J}
$$
où $\mathcal{J}$ est l'idéal engendré par les sections locales de
la forme $st - t$, pour $t$ section locale de l'un des facteurs directs
$\mathcal{L}_1^{\otimes -a_1} \otimes \mathcal{L}_2^{\otimes -a_2}$
avec $a_1, a_2 > 0$. Les sections $s_i$, considérées comme applications
$p^*\mathcal{L}_i^{\otimes -1} \to \mathcal{O}_X$, sont définies comme les adjointes
des applications $\mathcal{L}_i^{\otimes -1} \to \mathcal{A} = p_*\mathcal{O}_X$.
Pour tout $y \in Y$, on peut choisir un
ouvert affine $V \subset Y$, disons $V = \Spec(A)$, contenant $y$, et des
trivialisations $z_i : \mathcal{O}_V \to \mathcal{L}_i^{\otimes -1}|_V$.
Remarquons que $f = s(z_1z_2) \in A$ définit le sous-schéma fermé $D$.
Alors, clairement,
$$
p^{-1}(V) = \Spec(B[z_1, z_2]/(z_1 z_2 - f))
$$
Puisque $D_i$ est défini par $z_i$, tout est clair.
\end{proof}

\begin{lemma}
\label{lemma-decompose-section-formulae}
Dans la situation du lemme \ref{lemma-decompose-section},
supposons $Y$ localement de type fini sur $(S, \delta)$ comme dans la
situation \ref{situation-setup}. Alors on a
$i_1^*p^*\alpha = p_1^*i^*\alpha$
dans $\CH_k(D_1)$ pour tout $\alpha \in \CH_k(Y)$.
\end{lemma}

\begin{proof}
Soit $W \subset Y$ un sous-schéma fermé intègre de $\delta$-dimension $k$.
Distinguons deux cas.

\medskip\noindent
Supposons $W \subset D$. Alors
$i^*[W] = c_1(\mathcal{L}_1) \cap [W] + c_1(\mathcal{L}_2) \cap [W]$
dans $\CH_{k - 1}(D)$, d'après notre définition des homomorphismes de Gysin et
l'additivité du lemme \ref{lemma-c1-cap-additive}.
Ainsi $p_1^*i^*[W] =
p_1^*(c_1(\mathcal{L}_1) \cap [W]) + p_1^*(c_1(\mathcal{L}_2) \cap [W])$.
D'autre part, on a
$p^*[W] = [p^{-1}(W)]_{k + 1}$ par construction de l'image réciproque plate.
De plus, $p^{-1}(W) = W_1 \cup W_2$ (au sens schématique),
où $W_i = p_i^{-1}(W)$ est un fibré en droites sur $W$
d'après le lemme (puisque la formation du diagramme commute au changement de base).
Alors $[p^{-1}(W)]_{k + 1} = [W_1] + [W_2]$, car les $W_i$ sont des sous-schémas fermés
intègres de $L$ de $\delta$-dimension $k + 1$. Ainsi
\begin{align*}
i_1^*p^*[W]
& =
i_1^*[p^{-1}(W)]_{k + 1} \\
& =
i_1^*([W_1] + [W_2]) \\
& =
c_1(p_1^*\mathcal{L}_1) \cap [W_1] + [W_1 \cap W_2]_k \\
& =
c_1(p_1^*\mathcal{L}_1) \cap p_1^*[W] + [W_1 \cap W_2]_k \\
& =
p_1^*(c_1(\mathcal{L}_1) \cap [W]) + [W_1 \cap W_2]_k
\end{align*}
par construction des homomorphismes de Gysin, par définition de l'image réciproque plate
(pour la deuxième égalité) et par compatibilité de $c_1 \cap -$
avec l'image réciproque plate (lemme \ref{lemma-flat-pullback-cap-c1}).
Puisque $W_1 \cap W_2$ est la section nulle du fibré en droites
$W_1 \to W$, le lemme \ref{lemma-linebundle-formulae} donne
$[W_1 \cap W_2]_k = p_1^*(c_1(\mathcal{L}_2) \cap [W])$.
On utilise ici le fait que $D_1$ est le fibré en droites
qui est le spectre relatif de l'inverse de $\mathcal{L}_2$.
On obtient donc la même expression que précédemment.

\medskip\noindent
Supposons $W \not \subset D$. Dans ce cas, $i_1^*p^*[W]$
et $p_1^*i^*[W]$ sont tous deux représentés par le $k - 1$-cycle associé
à l'image réciproque schématique de $W$ dans $D_1$.
\end{proof}

\begin{lemma}
\label{lemma-normal-cone-effective-Cartier}
Dans la situation \ref{situation-setup}, soit $X$ un schéma localement
de type fini sur $S$. Soit $(\mathcal{L}, s, i : D \to X)$
un triplet comme dans la définition \ref{definition-gysin-homomorphism}.
Il existe un diagramme commutatif
$$
\xymatrix{
D' \ar[r]_{i'} \ar[d]_p & X' \ar[d]^g \\
D \ar[r]^i & X
}
$$
tel que
\begin{enumerate}
\item $p$ et $g$ soient de type fini et plats de dimension relative $1$,
\item $p^* : \CH_k(D) \to \CH_{k + 1}(D')$ soit injectif pour tout $k$,
\item $D' \subset X'$ soit le schéma des zéros d'une section globale
$s' \in \Gamma(X', \mathcal{O}_{X'})$,
\item $p^*i^* = (i')^*g^*$ comme applications $\CH_k(X) \to \CH_k(D')$.
\end{enumerate}
De plus, ces propriétés restent vraies après tout changement de base
par un morphisme $Y \to X$ localement de type fini.
\end{lemma}

\begin{proof}
Remarquons que $(i')^*$ est défini, car on dispose du triplet
$(\mathcal{O}_{X'}, s', i' : D' \to X')$ comme dans la
définition \ref{definition-gysin-homomorphism}. L'énoncé a donc un sens.

\medskip\noindent
Posons $\mathcal{L}_1 = \mathcal{O}_X$, $\mathcal{L}_2 = \mathcal{L}$,
et appliquons le lemme \ref{lemma-decompose-section} à la section $s$ de
$\mathcal{L} = \mathcal{L}_1 \otimes_{\mathcal{O}_X} \mathcal{L}_2$.
Prenons $D' = D_1$. Les résultats découlent alors de ce lemme, du
lemme \ref{lemma-decompose-section-formulae},
et l'injectivité du
lemme \ref{lemma-linebundle}.
\end{proof}

\begin{remark}
\label{remark-higher-chow-isomorphism}
Soit $(S, \delta)$ comme dans la situation \ref{situation-setup}.
Soit $Y$ localement de type fini sur $S$. Soit $r \geq 0$. Soit
$f : X \to Y$ un morphisme de schémas. Supposons que tout $y \in Y$
soit contenu dans un ouvert $V \subset Y$ tel que
$f^{-1}(V) \cong V \times \mathbf{A}^r$ comme schémas sur $V$.
Dans cette remarque, nous esquissons une démonstration du fait que
$f^* : \CH_k(Y) \to \CH_{k + r}(X)$ est un isomorphisme.
D'abord, d'après le lemme \ref{lemma-pullback-affine-fibres-surjective},
l'application est surjective. Soit $\alpha \in \CH_k(Y)$ tel que $f^*\alpha = 0$.
Nous allons démontrer que $\alpha = 0$.

\medskip\noindent
Étape 1. On peut supposer que $\dim_\delta(Y) < \infty$. (En pratique,
cela est immédiat dans tous les cas, et nous suggérons donc de sauter cette étape.)
En effet, toute équivalence rationnelle témoignant de l'égalité $f^*\alpha = 0$ sur $X$
fait intervenir une famille localement finie de sous-schémas fermés intègres
de dimension $k + r + 1$. En prenant la réunion des adhérences de leurs images
dans $Y$, on obtient un sous-schéma fermé $Y' \subset Y$
tel que $\dim_\delta(Y') \leq k + r + 1$, que
$\alpha$ soit l'image d'un certain $\alpha' \in \CH_k(Y')$,
et que $(f')^*\alpha = 0$, où $f'$ est le changement de base de $f$ à $Y'$.

\medskip\noindent
Étape 2. Supposons $d = \dim_\delta(Y) < \infty$. On peut alors raisonner par récurrence sur
$d$. Si $d < k$, alors $\alpha = 0$ et le résultat est acquis ; c'est le cas initial
de la récurrence.
En général, notre hypothèse sur $f$ montre que l'on peut choisir un ouvert dense
$V \subset Y$ tel que $U = f^{-1}(V) = \mathbf{A}^r_V$.
Notons $Y' \subset Y$ le complémentaire de $V$, muni de sa structure de sous-schéma fermé
réduit, et posons $X' = f^{-1}(Y')$. Considérons
$$
\xymatrix{
\CH^M_{k + r}(U, 1) \ar[r] &
\CH_{k + r}(X') \ar[r] &
\CH_{k + r}(X) \ar[r] &
\CH_{k + r}(U) \ar[r] &
0 \\
\CH^M_k(V, 1) \ar[r] \ar[u] &
\CH_k(Y') \ar[r] \ar[u] &
\CH_k(Y) \ar[r] \ar[u] &
\CH_k(V) \ar[r] \ar[u] &
0
}
$$
Nous utilisons ici les premiers groupes de Chow supérieurs de $V$ et de $U$ et les
suites exactes à six termes construites dans la remarque \ref{remark-higher-chow},
ainsi que l'image réciproque plate pour ces groupes de Chow supérieurs et
sa compatibilité avec ces suites exactes à six termes.
Puisque $U = \mathbf{A}^r_V$, l'application verticale de droite est un isomorphisme.
L'application $\CH_k(Y') \to \CH_{k + r}(X')$ est bijective par récurrence sur $d$.
Pour achever le raisonnement, il suffit donc de montrer que
$$
\CH^M_k(V, 1) \longrightarrow \CH^M_{k + r}(U, 1)
$$
est surjective. En raisonnant comme dans la démonstration du
lemme \ref{lemma-pullback-affine-fibres-surjective},
on se ramène à l'étape 3 ci-dessous.

\medskip\noindent
Étape 3. Soit $F$ un corps. Alors $\CH^M_0(\mathbf{A}^1_F, 1) = 0$.
(Dans la démonstration du lemme cité ci-dessus, nous avons montré de manière analogue que
$\CH_0(\mathbf{A}^1_F) = 0$.) On a
$$
\CH^M_0(\mathbf{A}^1_F, 1) = \Coker\left(
\partial : K^M_2(F(T)) \longrightarrow
\bigoplus\nolimits_{\mathfrak p \subset F[T]\text{ maximal}}
\kappa(\mathfrak p)^*\right)
$$
L'argument classique pour établir la nullité du conoyau consiste à
montrer, par récurrence sur le degré de $\kappa(\mathfrak p)/F$,
que le facteur direct correspondant à $\mathfrak p$ est contenu dans l'image.
Si $\mathfrak p$ est engendré par le
polynôme unitaire irréductible $P(T) \in F[T]$ et si
$u \in \kappa(x)^*$ est la classe résiduelle d'un
$Q(T) \in F[T]$ tel que $\deg(Q) < \deg(P)$, on montre que
$\partial(Q, P)$ donne l'élément $u$ en $\mathfrak p$,
et éventuellement d'autres unités en des idéaux premiers divisant $Q$, qui
ont un degré inférieur. Cela achève l'esquisse de démonstration.
\end{remark}











\section{Théorie bivariante de l'intersection}
\label{section-bivariant}

\noindent
Afin de traiter convenablement les classes de Chern supérieures des fibrés
vectoriels, nous introduisons les classes de Chow bivariantes comme dans \cite{F}.
Notre définition diffère de celle de \cite{F} sur deux points :
(1) nous travaillons dans un cadre différent, et (2) nous exigeons seulement
que nos classes bivariantes commutent aux homomorphismes de Gysin pour
les schémas des zéros de sections de modules inversibles
(section \ref{section-intersecting-effective-Cartier}).
Nous verrons plus loin, dans le lemme \ref{lemma-gysin-commutes}, que nos
classes bivariantes commutent à tous les homomorphismes de Gysin en codimension supérieure,
et satisfont donc à toutes les propriétés exigées dans \cite{F} ; voir
aussi \cite[Théorème 17.1]{F}.

\begin{definition}
\label{definition-bivariant-class}
\begin{reference}
Analogue à \cite[Définition 17.1]{F}
\end{reference}
Soit $(S, \delta)$ comme dans la situation \ref{situation-setup}.
Soit $f : X \to Y$ un morphisme de schémas localement de type fini sur $S$.
Soit $p \in \mathbf{Z}$.
Une {\it classe bivariante $c$ de degré $p$ pour $f$} est donnée par une règle
qui associe à tout morphisme localement de type fini $Y' \to Y$
et à tout $k$ une application
$$
c \cap - : \CH_k(Y') \longrightarrow \CH_{k - p}(X')
$$
où $X' = Y' \times_Y X$, satisfaisant aux conditions suivantes :
\begin{enumerate}
\item si $Y'' \to Y'$ est propre, alors
$c \cap (Y'' \to Y')_*\alpha'' = (X'' \to X')_*(c \cap \alpha'')$
pour tout $\alpha''$ sur $Y''$, où $X'' = Y'' \times_Y X$,
\item si $Y'' \to Y'$ est plat, localement de type fini et de
dimension relative fixée, alors
$c \cap (Y'' \to Y')^*\alpha' = (X'' \to X')^*(c \cap \alpha')$
pour tout $\alpha'$ sur $Y'$, et
\item si $(\mathcal{L}', s', i' : D' \to Y')$ est comme dans la
définition \ref{definition-gysin-homomorphism},
avec pour image réciproque $(\mathcal{N}', t', j' : E' \to X')$ sur $X'$,
alors on a $c \cap (i')^*\alpha' = (j')^*(c \cap \alpha')$
pour tout $\alpha'$ sur $Y'$.
\end{enumerate}
L'ensemble des classes bivariantes de degré $p$ pour $f$ est
noté $A^p(X \to Y)$.
\end{definition}

\noindent
Soit $(S, \delta)$ comme dans la situation \ref{situation-setup}. Soient $X \to Y$
et $Y \to Z$
des morphismes de schémas localement de type fini sur $S$. Soit
$p \in \mathbf{Z}$. Il est clair que $A^p(X \to Y)$ est un groupe abélien.
De plus, on dispose clairement d'une composition bilinéaire
$$
A^p(X \to Y) \times A^q(Y \to Z) \to A^{p + q}(X \to Z)
$$
qui est associative.

\begin{lemma}
\label{lemma-flat-pullback-bivariant}
Soit $(S, \delta)$ comme dans la situation \ref{situation-setup}.
Soit $f : X \to Y$ un morphisme plat de dimension relative $r$
entre schémas localement de type fini sur $S$.
Alors la règle qui associe à $Y' \to Y$
$(f')^* : \CH_k(Y') \to \CH_{k + r}(X')$, où $X' = X \times_Y Y'$,
est une classe bivariante de degré $-r$.
\end{lemma}

\begin{proof}
Cela découle des
lemmes \ref{lemma-flat-pullback-rational-equivalence},
\ref{lemma-compose-flat-pullback},
\ref{lemma-flat-pullback-proper-pushforward} et
\ref{lemma-gysin-flat-pullback}.
\end{proof}

\begin{lemma}
\label{lemma-gysin-bivariant}
Soit $(S, \delta)$ comme dans la situation \ref{situation-setup}.
Soit $X$ localement de type fini sur $S$.
Soit $(\mathcal{L}, s, i : D \to X)$ un triplet comme dans la
définition \ref{definition-gysin-homomorphism}.
Alors la règle qui associe à $f : X' \to X$
$(i')^* : \CH_k(X') \to \CH_{k - 1}(D')$, où $D' = D \times_X X'$,
est une classe bivariante de degré $1$.
\end{lemma}

\begin{proof}
Cela découle des lemmes \ref{lemma-gysin-factors},
\ref{lemma-closed-in-X-gysin},
\ref{lemma-gysin-flat-pullback} et
\ref{lemma-gysin-commutes-gysin}.
\end{proof}

\begin{lemma}
\label{lemma-push-proper-bivariant}
Soit $(S, \delta)$ comme dans la situation \ref{situation-setup}.
Soient $f : X \to Y$ et $g : Y \to Z$ des morphismes de
schémas localement de type fini sur $S$.
Soit $c \in A^p(X \to Z)$, et supposons $f$ propre.
Alors la règle qui associe à $Z' \to Z$
$\alpha \longmapsto f'_*(c \cap \alpha)$
est une classe bivariante notée $f_* \circ c \in A^p(Y \to Z)$.
\end{lemma}

\begin{proof}
Cela découle des lemmes \ref{lemma-compose-pushforward},
\ref{lemma-flat-pullback-proper-pushforward} et
\ref{lemma-closed-in-X-gysin}.
\end{proof}

\begin{remark}
\label{remark-restriction-bivariant}
Soit $(S, \delta)$ comme dans la situation \ref{situation-setup}. Soient $X \to Y$
et $Y' \to Y$ des morphismes de schémas localement de type fini sur $S$.
Soit $X' = Y' \times_Y X$. On dispose alors d'une application de restriction évidente
$$
A^p(X \to Y) \longrightarrow A^p(X' \to Y'),\quad
c \longmapsto res(c)
$$
obtenue en considérant un schéma $Y''$ localement de type fini sur $Y'$
comme un schéma localement de type fini sur $Y$ et en posant
$res(c) \cap \alpha'' = c \cap \alpha''$ pour tout $\alpha'' \in \CH_k(Y'')$.
Cette opération de restriction est compatible avec les compositions de manière
évidente.
\end{remark}

\begin{remark}
\label{remark-bivariant-commute}
Soit $(S, \delta)$ comme dans la situation \ref{situation-setup}. Soit $X$
localement de type fini sur $S$. Pour $i = 1, 2$, soit $Z_i \to X$
un morphisme de schémas localement de type fini. Soient
$c_i \in A^{p_i}(Z_i \to X)$, $i = 1, 2$ des classes bivariantes.
Pour tout $\alpha \in \CH_k(X)$, on peut se demander si
$$
c_1 \cap c_2 \cap \alpha = c_2 \cap c_1 \cap \alpha
$$
dans $\CH_{k - p_1 - p_2}(Z_1 \times_X Z_2)$. Si cette égalité est vraie et le demeure
après tout changement de base par un morphisme $X' \to X$ localement de type fini, on dit
que $c_1$ et $c_2$ {\it commutent}. Bien entendu, cela revient à dire que
$$
res(c_1) \circ c_2 = res(c_2) \circ c_1
$$
dans $A^{p_1 + p_2}(Z_1 \times_X Z_2 \to X)$. Ici,
$res(c_1) \in A^{p_1}(Z_1 \times_X Z_2 \to Z_2)$ est la restriction de $c_1$
au sens de la remarque \ref{remark-restriction-bivariant} ; de même pour $res(c_2)$.
\end{remark}

\begin{example}
\label{example-gysin-commute}
Soit $(S, \delta)$ comme dans la situation \ref{situation-setup}. Soit $X$ localement
de type fini sur $S$. Soit $(\mathcal{L}, s, i : D \to X)$ un triplet comme dans la
définition \ref{definition-gysin-homomorphism}. Soit $Z \to X$ un morphisme
de schémas localement de type fini, et soit $c \in A^p(Z \to X)$ une
classe bivariante. Alors la classe de Gysin bivariante $c' \in A^1(D \to X)$ du
lemme \ref{lemma-gysin-bivariant} commute à $c$ au sens de la
remarque \ref{remark-bivariant-commute}. Il s'agit en effet d'une reformulation de la
condition (3) de la définition \ref{definition-bivariant-class}.
\end{example}

\begin{remark}
\label{remark-more-general-bivariant}
Il existe un type plus général de classe bivariante qui ne semble pas être
considéré dans la littérature. Supposons en effet donné un diagramme
$$
X \longrightarrow Z \longleftarrow Y
$$
de schémas localement de type fini sur $(S, \delta)$ comme dans la
situation \ref{situation-setup}. Soit $p \in \mathbf{Z}$.
On peut alors considérer une règle $c$ qui associe à tout $Z' \to Z$
localement de type fini des applications
$$
c \cap - : \CH_k(Y') \longrightarrow \CH_{k - p}(X')
$$
pour tout $k \in \mathbf{Z}$,
où $X' = X \times_Z Z'$ et $Y' = Z' \times_Z Y$, compatibles aux
\begin{enumerate}
\item images directes propres pour $Z'' \to Z'$ propre,
\item images réciproques plates pour $Z'' \to Z'$ plat
de dimension relative fixée, et
\item applications de Gysin pour $D' \subset Z'$ comme dans la
définition \ref{definition-gysin-homomorphism}.
\end{enumerate}
Nous omettons les formulations détaillées. Notons l'ensemble
de toutes ces opérations $A^p(X \to Z \leftarrow Y)$. Un exemple simple
de l'utilité de cette notion est fourni par un morphisme propre
$f : X_2 \to X_1$. Alors $f_*$ n'est pas une opération bivariante au sens de la
définition \ref{definition-bivariant-class}, mais l'est au sens
généralisé ci-dessus ; on a $f_* \in A^0(X_1 \to X_1 \leftarrow X_2)$.
\end{remark}





\section{Cohomologie de Chow et première classe de Chern}
\label{section-chow-cohomology}

\noindent
Nous nous intéresserons surtout à $A^p(X) = A^p(X \to X)$, qui désignera toujours
les classes de cohomologie bivariante pour $\text{id}_X$. C'est en effet là que
se trouveront les classes de Chern.

\begin{definition}
\label{definition-chow-cohomology}
Soit $(S, \delta)$ comme dans la situation \ref{situation-setup}.
Soit $X$ localement de type fini sur $S$. La {\it cohomologie de Chow}
de $X$ est la $\mathbf{Z}$-algèbre graduée $A^*(X)$ dont la composante de degré
$p$ est $A^p(X \to X)$.
\end{definition}

\noindent
Attention : il n'est pas clair que la structure de $\mathbf{Z}$-algèbre
sur $A^*(X)$ soit commutative, mais nous verrons que les classes de Chern appartiennent
à son centre.

\begin{remark}
\label{remark-pullback-cohomology}
Soit $(S, \delta)$ comme dans la situation \ref{situation-setup}.
Soit $f : Y' \to Y$ un morphisme de schémas localement de type fini sur $S$.
Comme cas particulier de la remarque \ref{remark-restriction-bivariant},
on dispose d'un homomorphisme canonique de $\mathbf{Z}$-algèbres $res : A^*(Y) \to A^*(Y')$.
Cet homomorphisme est souvent noté $f^*$ dans la littérature.
\end{remark}

\begin{lemma}
\label{lemma-cap-c1-bivariant}
Soit $(S, \delta)$ comme dans la situation \ref{situation-setup}.
Soit $X$ localement de type fini sur $S$.
Soit $\mathcal{L}$ un $\mathcal{O}_X$-module inversible.
Alors la règle qui associe à $f : X' \to X$
$c_1(f^*\mathcal{L}) \cap - : \CH_k(X') \to \CH_{k - 1}(X')$
est une classe bivariante de degré $1$.
\end{lemma}

\begin{proof}
Cela découle des lemmes \ref{lemma-factors},
\ref{lemma-pushforward-cap-c1},
\ref{lemma-flat-pullback-cap-c1} et
\ref{lemma-gysin-commutes-cap-c1}.
\end{proof}

\noindent
Le lemme précédent nous permet enfin de donner la définition suivante.

\begin{definition}
\label{definition-first-chern-class}
Soit $(S, \delta)$ comme dans la situation \ref{situation-setup}. Soit $X$
localement de type fini sur $S$. Soit $\mathcal{L}$ un
$\mathcal{O}_X$-module inversible. La {\it première classe de Chern}
$c_1(\mathcal{L}) \in A^1(X)$ de $\mathcal{L}$
est la classe bivariante du lemme \ref{lemma-cap-c1-bivariant}.
\end{definition}

\noindent
Pour les modules localement libres de rang fini, nous construisons les classes de Chern dans la
section \ref{section-intersecting-chern-classes}.
Montrons que $c_1(\mathcal{L})$ appartient au centre de $A^*(X)$.

\begin{lemma}
\label{lemma-c1-center}
Soit $(S, \delta)$ comme dans la situation \ref{situation-setup}.
Soit $X$ localement de type fini sur $S$.
Soit $\mathcal{L}$ un $\mathcal{O}_X$-module inversible.
Alors
\begin{enumerate}
\item $c_1(\mathcal{L}) \in A^1(X)$ appartient au centre de $A^*(X)$, et
\item si $f : X' \to X$ est localement de type fini et $c \in A^*(X' \to X)$,
alors $c \circ c_1(\mathcal{L}) = c_1(f^*\mathcal{L}) \circ c$.
\end{enumerate}
\end{lemma}

\begin{proof}
Bien entendu, (2) entraîne (1).
Soit $p : L \to X$ comme dans le lemme \ref{lemma-linebundle}, et soit $o : X \to L$
la section nulle. Notons $p' : L' \to X'$ et $o' : X' \to L'$
leurs changements de base. D'après le lemme \ref{lemma-linebundle-formulae}, on a
$$
p^*(c_1(\mathcal{L}) \cap \alpha) = - o_* \alpha
\quad\text{et}\quad
(p')^*(c_1(f^*\mathcal{L}) \cap \alpha') = - o'_* \alpha'
$$
Puisque $c$ est une classe bivariante, on a
\begin{align*}
(p')^*(c \cap c_1(\mathcal{L}) \cap \alpha)
& =
c \cap p^*(c_1(\mathcal{L}) \cap \alpha) \\
& =
- c \cap o_* \alpha \\
& =
- o'_*(c \cap \alpha) \\
& =
(p')^*(c_1(f^*\mathcal{L}) \cap c \cap \alpha)
\end{align*}
Puisque $(p')^*$ est injectif d'après l'un des lemmes cités ci-dessus, on obtient
$c \cap c_1(\mathcal{L}) \cap \alpha =
c_1(f^*\mathcal{L}) \cap c \cap \alpha$.
Il en est de même après tout changement de base par un morphisme $Y \to X$ localement de type fini,
d'où l'égalité de classes bivariantes énoncée dans (2).
\end{proof}

\begin{lemma}
\label{lemma-vanish-above-dimension}
Soit $(S, \delta)$ comme dans la situation \ref{situation-setup}. Soit $X$ un
schéma de type fini sur $S$ possédant un faisceau inversible ample.
Supposons $d = \dim(X) < \infty$ (il s'agit bien ici de la dimension et
non de la $\delta$-dimension).
Alors, pour tous faisceaux inversibles $\mathcal{L}_1, \ldots, \mathcal{L}_{d + 1}$
sur $X$, on a
$c_1(\mathcal{L}_1) \circ \ldots \circ c_1(\mathcal{L}_{d + 1}) = 0$
dans $A^{d + 1}(X)$.
\end{lemma}

\begin{proof}
Raisonnons par récurrence sur $d$. Le cas initial $d = 0$ est vrai, car
dans ce cas $X$ est un ensemble fini de points fermés, et tout module
inversible est donc trivial. Supposons $d > 0$. D'après Diviseurs, lemme
\ref{divisors-lemma-quasi-projective-Noetherian-pic-effective-Cartier},
on peut écrire $\mathcal{L}_{d + 1} \cong \mathcal{O}_X(D) \otimes
\mathcal{O}_X(D')^{\otimes -1}$ pour des diviseurs de Cartier effectifs
$D, D' \subset X$. Alors $c_1(\mathcal{L}_{d + 1})$ est la différence
de $c_1(\mathcal{O}_X(D))$ et de $c_1(\mathcal{O}_X(D'))$ ; on peut donc
supposer $\mathcal{L}_{d + 1} = \mathcal{O}_X(D)$ pour un
diviseur de Cartier effectif.

\medskip\noindent
Notons $i : D \to X$ le morphisme d'inclusion, et notons
$i^* \in A^1(D \to X)$ la classe bivariante donnée par
l'homomorphisme de Gysin comme dans le lemme \ref{lemma-gysin-bivariant}.
On a $i_* \circ i^* = c_1(\mathcal{L}_{d + 1})$
dans $A^1(X)$ d'après le lemme \ref{lemma-support-cap-effective-Cartier}
(et le lemme \ref{lemma-push-proper-bivariant}
pour donner un sens au membre de gauche).
Puisque $c_1(\mathcal{L}_i)$ commute
à $i_*$ et à $i^*$ (par définition des classes bivariantes),
on en conclut que
$$
c_1(\mathcal{L}_1) \circ \ldots \circ c_1(\mathcal{L}_{d + 1}) =
i_* \circ c_1(\mathcal{L}_1) \circ \ldots \circ c_1(\mathcal{L}_d) \circ i^* =
i_* \circ c_1(\mathcal{L}_1|_D) \circ \ldots \circ c_1(\mathcal{L}_d|_D)
\circ i^*
$$
On conclut donc par récurrence sur $d$. En effet, on a $\dim(D) < d$,
car aucun des points génériques de $X$ n'appartient à $D$.
\end{proof}

\begin{remark}
\label{remark-ring-loc-classes}
Soit $(S, \delta)$ comme dans la situation \ref{situation-setup}.
Soit $Z \to X$ une immersion fermée de schémas localement de
type fini sur $S$, et soit $p \geq 0$. Dans ce cadre, on définit
$$
A^{(p)}(Z \to X) =
\prod\nolimits_{i \leq p - 1} A^i(X) \times
\prod\nolimits_{i \geq p} A^i(Z \to X).
$$
Alors $A^{(p)}(Z \to X)$ est canoniquement muni d'une structure
d'algèbre graduée. En fait, plus généralement, on dispose d'une multiplication
$$
A^{(p)}(Z \to X) \times A^{(q)}(Z \to X)
\longrightarrow A^{(\max(p, q))}(Z \to X)
$$
Pour les définir, on définit des applications
\begin{align*}
A^i(Z \to X) \times A^j(X) & \to A^{i + j}(Z \to X) \\
A^i(X) \times A^j(Z \to X) & \to A^{i + j}(Z \to X) \\
A^i(Z \to X) \times A^j(Z \to X) & \to A^{i + j}(Z \to X)
\end{align*}
Pour la première, on utilise la composition des classes bivariantes.
Pour la deuxième, on utilise la restriction
$A^i(X) \to A^i(Z)$ (remarque \ref{remark-restriction-bivariant}) et la
composition $A^i(Z) \times A^j(Z \to X) \to A^{i + j}(Z \to X)$.
Pour la troisième, on envoie $(c, c')$ sur $res(c) \circ c'$,
où $res : A^i(Z \to X) \to A^i(Z)$ est l'application de restriction (voir la
remarque \ref{remark-restriction-bivariant}). Nous omettons la
vérification de l'associativité de ces multiplications en un sens approprié.
\end{remark}

\begin{remark}
\label{remark-res-push}
Soit $(S, \delta)$ comme dans la situation \ref{situation-setup}.
Soit $Z \to X$ une immersion fermée de schémas localement de
type fini sur $S$. Notons $res : A^p(Z \to X) \to A^p(Z)$
l'application de restriction de la remarque \ref{remark-restriction-bivariant}.
Pour $c \in A^p(Z \to X)$, on a
$res(c) \cap \alpha = c \cap i_*\alpha$ pour $\alpha \in \CH_*(Z)$.
En effet, $res(c) \cap \alpha = c \cap \alpha$,
et la compatibilité de $c$ avec l'image directe propre
donne $(Z \to Z)_*(c \cap \alpha) = c \cap (Z \to X)_*\alpha$.
\end{remark}






\section{Lemmes sur les classes bivariantes}
\label{section-bivariant-II}

\noindent
Dans cette section, nous démontrons quelques résultats élémentaires sur les
classes bivariantes. Voici un critère permettant de voir qu'une opération
passe au quotient par équivalence rationnelle.

\begin{lemma}
\label{lemma-factors-through-rational-equivalence}
\begin{reference}
Forme très faible de \cite[Théorème 17.1]{F}
\end{reference}
Soit $(S, \delta)$ comme dans la situation \ref{situation-setup}.
Soit $f : X \to Y$ un morphisme de schémas localement de type fini sur $S$.
Soit $p \in \mathbf{Z}$. Supposons donnée une règle
qui associe à tout morphisme localement de type fini $Y' \to Y$
et à tout $k$ une application
$$
c \cap - : Z_k(Y') \longrightarrow \CH_{k - p}(X')
$$
où $Y' = X' \times_X Y$, satisfaisant à la condition (3) de la
définition \ref{definition-bivariant-class}
dès que $\mathcal{L}'|_{D'} \cong \mathcal{O}_{D'}$. Alors
$c \cap -$ se factorise par l'équivalence rationnelle.
\end{lemma}

\begin{proof}
L'énoncé a un sens car, étant donné un triplet
$(\mathcal{L}, s, i : D \to X)$ comme dans la
définition \ref{definition-gysin-homomorphism}
tel que $\mathcal{L}|_D \cong \mathcal{O}_D$, l'opération
$i^*$ est définie au niveau des cycles, voir la
remarque \ref{remark-gysin-on-cycles}.
Soit $\alpha \in Z_k(X')$ un cycle rationnellement équivalent à zéro.
Nous devons montrer que $c \cap \alpha = 0$. D'après le
lemme \ref{lemma-rational-equivalence-family},
il existe un cycle $\beta \in Z_{k + 1}(X' \times \mathbf{P}^1)$
tel que $\alpha = i_0^*\beta - i_\infty^*\beta$,
où $i_0, i_\infty : X' \to X' \times \mathbf{P}^1$ sont les
immersions fermées de $X'$ au-dessus de $0, \infty$. Puisque ce sont
des exemples de diviseurs de Cartier effectifs à fibrés normaux
triviaux, on voit que $c \cap i_0^*\beta = j_0^*(c \cap \beta)$
et $c \cap i_\infty^*\beta = j_\infty^*(c \cap \beta)$,
où $j_0, j_\infty : Y' \to Y' \times \mathbf{P}^1$ sont les
immersions fermées comme ci-dessus. Comme
$j_0^*(c \cap \beta) \sim_{rat} j_\infty^*(c \cap \beta)$
(cela résulte du lemme \ref{lemma-rational-equivalence-family}), on conclut.
\end{proof}

\begin{lemma}
\label{lemma-bivariant-weaker}
\begin{reference}
Forme faible de \cite[Théorème 17.1]{F}
\end{reference}
Soit $(S, \delta)$ comme dans la situation \ref{situation-setup}.
Soit $f : X \to Y$ un morphisme de schémas localement de type fini sur $S$.
Soit $p \in \mathbf{Z}$. Supposons donnée une règle
qui associe à tout morphisme localement de type fini $Y' \to Y$
et à tout $k$ une application
$$
c \cap - : \CH_k(Y') \longrightarrow \CH_{k - p}(X')
$$
où $Y' = X' \times_X Y$, satisfaisant aux conditions (1), (2) de la
définition \ref{definition-bivariant-class}
et à la condition (3) dès que $\mathcal{L}'|_{D'} \cong \mathcal{O}_{D'}$. Alors
$c \cap -$ est une classe bivariante.
\end{lemma}

\begin{proof}
Soit $Y' \to Y$ un morphisme de schémas localement de type fini.
Soit $(\mathcal{L}', s', i' : D' \to Y')$ comme dans la
définition \ref{definition-gysin-homomorphism},
et soit $(\mathcal{N}', t', j' : E' \to X')$ son image réciproque sur $X'$.
Nous devons montrer que $c \cap (i')^*\alpha' = (j')^*(c \cap \alpha')$
pour tout $\alpha' \in \CH_k(Y')$.

\medskip\noindent
Notons $g : Y'' \to Y'$ le morphisme lisse de dimension relative
$1$, avec $i'' : D'' \to Y''$ et $p : D'' \to D'$,
construit dans le lemme \ref{lemma-normal-cone-effective-Cartier}.
(Attention : $D''$ n'est pas l'image réciproque entière de $D'$.)
Notons $f : X'' \to X'$ et $E'' \subset X''$
leurs changements de base par $X' \to Y'$. On a le diagramme
$$
\xymatrix{
& X'' \ar[rr] \ar'[d][dd]_h & & Y'' \ar[dd]^g \\
E'' \ar[rr] \ar[dd]_q \ar[ru]^{j''} & & D'' \ar[dd]^p \ar[ru]^{i''} & \\
& X' \ar'[r][rr] & & Y' \\
E' \ar[rr] \ar[ru]^{j'} & & D' \ar[ru]^{i'}
}
$$
D'après les propriétés données dans le lemme, nous savons que $\beta' = (i')^*\alpha'$
est l'unique élément de $\CH_{k - 1}(D')$ tel que
$p^*\beta' = (i'')^*g^*\alpha'$. De même, nous savons que
$\gamma' = (j')^*(c \cap \alpha')$ est l'unique élément de
$\CH_{k - 1 - p}(E')$ tel que $q^*\gamma' = (j'')^*h^*(c \cap \alpha')$.
Or nous savons que
$$
(j'')^*h^*(c \cap \alpha') =
(j'')^*(c \cap g^*\alpha') =
c \cap (i'')^*g^*\alpha'
$$
par nos hypothèses sur $c$ ; remarquons que la version modifiée de (3)
supposée dans l'énoncé du lemme s'applique à $i''$
et à son changement de base $j''$. Nous savons de même que
$$
q^*(c \cap \beta') = c \cap p^*\beta'
$$
On conclut que $\gamma' = c \cap \beta'$ par l'unicité signalée
ci-dessus.
\end{proof}

\noindent
Voici un critère pour qu'une classe bivariante soit nulle.

\begin{lemma}
\label{lemma-bivariant-zero}
Soit $(S, \delta)$ comme dans la situation \ref{situation-setup}.
Soit $f : X \to Y$ un morphisme de schémas localement de type fini sur $S$.
Soit $c \in A^p(X \to Y)$. Pour $Y'' \to Y' \to Y$, posons
$X'' = Y'' \times_Y X$ et $X' = Y' \times_Y X$.
Les assertions suivantes sont équivalentes :
\begin{enumerate}
\item $c$ est nulle,
\item $c \cap [Y'] = 0$ dans $\CH_*(X')$ pour tout schéma intègre $Y'$
localement de type fini sur $Y$, et
\item pour tout schéma intègre $Y'$ localement de type fini sur $Y$,
il existe un morphisme propre birationnel $Y'' \to Y'$ tel que
$c \cap [Y''] = 0$ dans $\CH_*(X'')$.
\end{enumerate}
\end{lemma}

\begin{proof}
Les implications (1) $\Rightarrow$ (2) $\Rightarrow$ (3) sont claires.
L'hypothèse (3) implique (2) car $(Y'' \to Y')_*[Y''] = [Y']$
et donc $c \cap [Y'] = (X'' \to X')_*(c \cap [Y''])$, puisque $c$
est une classe bivariante. Supposons (2).
Soit $Y' \to Y$ localement de type fini. Soit $\alpha \in \CH_k(Y')$.
Écrivons $\alpha = \sum n_i [Y'_i]$, où les $Y'_i \subset Y'$ forment une
famille localement finie de sous-schémas fermés intègres de $\delta$-dimension $k$.
On voit alors que $\alpha$ est l'image directe du cycle
$\alpha' = \sum n_i[Y'_i]$ sur $Y'' = \coprod Y'_i$ par le
morphisme propre $Y'' \to Y'$. Par les propriétés des classes bivariantes,
il suffit de prouver que $c \cap \alpha' = 0$ dans $\CH_{k - p}(X'')$.
On a $\CH_{k - p}(X'') = \prod \CH_{k - p}(X'_i)$, où
$X'_i = Y'_i \times_Y X$. Cela résulte immédiatement
des définitions. Les applications de projection
$\CH_{k - p}(X'') \to \CH_{k - p}(X'_i)$ sont données par l'image réciproque
plate. Puisque le produit cap avec $c$ commute à l'image réciproque
plate, on voit qu'il suffit de montrer que $c \cap [Y'_i]$
est nul dans $\CH_{k - p}(X'_i)$, ce qui est vrai par hypothèse.
\end{proof}

\begin{lemma}
\label{lemma-disjoint-decomposition-bivariant}
Soit $(S, \delta)$ comme dans la situation \ref{situation-setup}.
Soit $f : X \to Y$ un morphisme de schémas localement de type fini sur $S$.
Supposons données des décompositions en réunions disjointes
$X = \coprod_{i \in I} X_i$ et $Y = \coprod_{j \in J} Y_j$
par des sous-schémas ouverts et fermés,
ainsi qu'une application d'ensembles $a : I \to J$ telle que $f(X_i) \subset Y_{a(i)}$.
Alors
$$
A^p(X \to Y) = \prod\nolimits_{i \in I} A^p(X_i \to Y_{a(i)})
$$
\end{lemma}

\begin{proof}
Supposons donné un élément $(c_i) \in \prod_i A^p(X_i \to Y_{a(i)})$.
Étant donné $\beta \in \CH_k(Y)$, on peut alors l'envoyer sur l'élément de
$\CH_{k - p}(X)$ dont la restriction à $X_i$ est $c_i \cap \beta|_{Y_{a(i)}}$.
Cela fonctionne car $\CH_{k - p}(X) = \prod_i \CH_{k - p}(X_i)$.
La même construction fonctionne après tout changement de base $Y' \to Y$
localement de type fini, et l'on obtient $c \in A^p(X \to Y)$.
On obtient ainsi une application $\Psi$ du membre droit de la formule
vers son membre gauche.
Réciproquement, étant donnés $c \in A^p(X \to Y)$ et un élément
$\beta_i \in \CH_k(Y_{a(i)})$, on peut considérer l'élément
$(c \cap (Y_{a(i)} \to Y)_*\beta_i)|_{X_i}$ de $\CH_{k - p}(X_i)$.
La même chose fonctionne après tout changement de base $Y' \to Y$
localement de type fini, et l'on obtient $c_i \in A^p(X_i \to Y_{a(i)})$.
On obtient ainsi une application $\Phi$ du membre gauche
de la formule vers son membre droit.
Il est immédiat que $\Phi \circ \Psi = \text{id}$.
Réciproquement, supposons que $c \in A^p(X \to Y)$ et
$\beta \in \CH_k(Y)$. Écrivons $\Phi(c) = (c_i)$. Soit $j \in J$.
Comme $c$ commute à l'image réciproque plate, on a
$$
(c \cap \beta)|_{\coprod_{a(i) = j} X_i} =
c \cap \beta|_{Y_j}
$$
Comme $c$ commute à l'image directe propre, on a
$$
(\coprod\nolimits_{a(i) = j} X_i \to X)_*
((c \cap \beta)|_{\coprod_{a(i) = j} X_i})
=
c \cap (Y_j \to Y)_*\beta|_{Y_j}
$$
Le membre gauche est le cycle sur $X$ qui se restreint en $(c \cap \beta)|_{X_i}$
sur $X_i$ pour $i \in I$ tel que $a(i) = j$, et qui vaut $0$ sinon.
Le membre droit est un cycle sur $X$ dont la restriction à $X_i$
est $c_i \cap \beta|_{Y_j}$ pour $i \in I$ tel que $a(i) = j$.
Ainsi, $c \cap \beta = \Psi((c_i))$, comme voulu.
\end{proof}

\begin{remark}
\label{remark-completion-bivariant}
Soit $(S, \delta)$ comme dans la situation \ref{situation-setup}.
Soit $f : X \to Y$ un morphisme de schémas localement de type fini sur $S$.
Soient $X = \coprod_{i \in I} X_i$ et $Y = \coprod_{j \in J} Y_j$
les décompositions de $X$ et de $Y$ en leurs composantes connexes
(les composantes connexes sont ouvertes puisque $X$ et $Y$ sont localement noethériens ; voir
Topologie, Lemme \ref{topology-lemma-locally-Noetherian-locally-connected}, et
Propriétés, Lemme \ref{properties-lemma-Noetherian-topology}).
Soit $a(i) \in J$ l'indice tel que $f(X_i) \subset Y_{a(i)}$.
Alors $A^p(X \to Y) = \prod A^p(X_i \to Y_{a(i)})$ d'après le
Lemme \ref{lemma-disjoint-decomposition-bivariant}.
Dans ce cadre, il est commode de poser
$$
A^*(X \to Y)^\wedge = \prod\nolimits_i A^*(X_i \to Y_{a(i)})
$$
Ce groupe bivariant « complété » est le sous-ensemble
$$
A^*(X \to Y)^\wedge \quad\subset\quad \prod\nolimits_{p \geq 0} A^p(X \to Y)
$$
formé des éléments $c = (c_0, c_1, c_2, \ldots)$ tels que,
pour toute composante connexe $X_i$, l'image de $c_p$ dans
$A^p(X_i \to Y_{a(i)})$ soit nulle pour presque tout $p$.
Si $Y \to Z$ est un second morphisme, alors la
composition $A^*(X \to Y) \times A^*(Y \to Z) \to A^*(X \to Z)$
s'étend en une composition
$A^*(X \to Y)^\wedge \times A^*(Y \to Z)^\wedge \to A^*(X \to Z)^\wedge$
des complétés.
On appelle parfois $A^*(X)^\wedge = A^*(X \to X)^\wedge$
l'{\it anneau de cohomologie bivariante complété} de $X$.
\end{remark}

\begin{lemma}
\label{lemma-envelope-bivariant}
Soit $(S, \delta)$ comme dans la situation \ref{situation-setup}.
Soit $f : X \to Y$ un morphisme de schémas localement de type fini sur $S$.
Soit $g : Y' \to Y$ une enveloppe (Définition \ref{definition-envelope})
et notons $X' = Y' \times_Y X$. Soit $p \in \mathbf{Z}$ et soit
$c' \in A^p(X' \to Y')$. Si les deux restrictions
$$
res_1(c') = res_2(c') \in A^p(X' \times_X X' \to Y' \times_Y Y')
$$
sont égales (voir la démonstration), alors il existe un unique $c \in A^p(X \to Y)$
dont la restriction $res(c) = c'$ dans $A^p(X' \to Y')$.
\end{lemma}

\begin{proof}
On a un diagramme commutatif
$$
\xymatrix{
X' \times_X X' \ar[d]^{f''} \ar@<1ex>[r]^-a \ar@<-1ex>[r]_-b &
X' \ar[d]^{f'} \ar[r]_h &
X \ar[d]^f \\
Y' \times_Y Y' \ar@<1ex>[r]^-p \ar@<-1ex>[r]_-q &
Y' \ar[r]^g &
Y
}
$$
L'élément $res_1(c')$ est la restriction (voir
Remarque \ref{remark-restriction-bivariant}) de $c'$ pour le carré cartésien
de morphismes $a, f', p, f''$, et l'élément $res_2(c')$ est la restriction
de $c'$ pour le carré cartésien de morphismes $b, f', q, f''$.
Supposons $res_1(c') = res_2(c')$ et soit $\beta \in \CH_k(Y)$.
D'après le Lemme \ref{lemma-envelope}, on peut trouver un $\beta' \in \CH_k(Y')$
tel que $g_*\beta' = \beta$. On pose alors
$$
c \cap \beta = h_*(c' \cap \beta')
$$
Pour voir que ceci ne dépend pas du choix de
$\beta'$, il suffit de montrer que
$h_*(c' \cap (p_*\gamma - q_*\gamma))$ est nul
pour $\gamma \in \CH_k(Y' \times_Y Y')$.
Puisque $c'$ est une classe bivariante, on a
$$
h_*(c' \cap (p_*\gamma - q_*\gamma)) =
h_*(a_*(c' \cap \gamma) - b_*(c' \cap \gamma)) = 0
$$
la dernière égalité résultant de $h_* \circ a_* = h_* \circ b_*$,
car $h \circ a = h \circ b$.

\medskip\noindent
Remarquons que notre choix de $c \cap \beta$ est imposé
par la condition $res(c) = c'$ et par la compatibilité
des classes bivariantes avec l'image directe propre.

\medskip\noindent
Bien entendu, pour définir la classe bivariante $c$, il nous faut
construire des applications $c \cap -: \CH_k(Y_1) \to \CH_{k + p}(Y_1 \times_Y X)$
pour tout morphisme $Y_1 \to Y$ localement de type fini satisfaisant aux
conditions énumérées dans la Définition \ref{definition-bivariant-class}.
Notons $Y'_1 = Y' \times_Y Y_1$, $X_1 = X \times_Y Y_1$.
Le morphisme $Y'_1 \to Y_1$ est une enveloppe d'après le
Lemme \ref{lemma-base-change-envelope}. Nous pouvons donc utiliser le
diagramme obtenu par changement de base
$$
\xymatrix{
X'_1 \times_{X_1} X'_1 \ar[d]^{f''_1} \ar@<1ex>[r]^-{a_1} \ar@<-1ex>[r]_-{b_1} &
X'_1 \ar[d]^{f'_1} \ar[r]_{h_1} &
X_1 \ar[d]^{f_1} \\
Y'_1 \times_{Y_1} Y'_1 \ar@<1ex>[r]^-{p_1} \ar@<-1ex>[r]_-{q_1} &
Y'_1 \ar[r]^{g_1} &
Y_1
}
$$
et les mêmes arguments pour obtenir une application bien définie
$c \cap - : \CH_k(Y_1) \to \CH_{k + p}(X_1)$ comme ci-dessus.

\medskip\noindent
Il nous faut ensuite vérifier les conditions (1), (2) et (3) de la
Définition \ref{definition-bivariant-class} pour $c$.
Par exemple, supposons que $t : Y_2 \to Y_1$ soit un morphisme propre
de schémas localement de type fini sur $Y$. Notons comme ci-dessus
les changements de base du premier diagramme à $Y_1$, resp.\ à $Y_2$,
par les indices ${}_1$, resp.\ ${}_2$. Notons $t' : Y'_2 \to Y'_1$,
$s : X_2 \to X_1$ et $s' : X'_2 \to X'_1$ les changements de base de $t$
à $Y'$, $X$ et $X'$. Nous devons montrer que
$$
s_*(c \cap \beta_2) = c \cap t_*\beta_2
$$
pour $\beta_2 \in \CH_k(Y_2)$. Choisissons $\beta'_2 \in \CH_k(Y'_2)$
tel que $g_{2, *}\beta'_2 = \beta_2$. Puisque $c'$ est une classe bivariante
et que les diagrammes
$$
\vcenter{
\xymatrix{
X'_2 \ar[d]_{s'} \ar[r]_{h_2} & X_2 \ar[d]^s \\
X'_1 \ar[r]^{h_1} & X_1
}
}
\quad\text{et}\quad
\vcenter{
\xymatrix{
X'_2 \ar[d]_{s'} \ar[r]_{f'_2} & Y'_2 \ar[d]^{t'} \\
X'_2 \ar[r]^{f'_1} & Y'_1
}
}
$$
sont cartésiens, on a
$$
s_*(c \cap \beta_2) =
s_*(h_{2, *}(c' \cap \beta'_2)) =
h_{1, *}s'_*(c' \cap \beta'_2) =
h_{1, *}(c' \cap (t'_*\beta'_2))
$$
et l'expression finale calcule $c \cap t_*\beta_2$ par construction :
$t'_*\beta'_2 \in \CH_k(Y'_1)$ est une classe dont l'image par $g_{1, *}$ est
$t_*\beta_2$. Ceci démontre la condition (1).
Les autres conditions se démontrent de la même manière,
et nous omettons les arguments détaillés.
\end{proof}















\section{Formule du fibré projectif}
\label{section-projective-space-bundle-formula}

\noindent
Soit $(S, \delta)$ comme dans la situation \ref{situation-setup}.
Soit $X$ localement de type fini sur $S$.
Considérons un $\mathcal{O}_X$-module localement libre de type fini
$\mathcal{E}$ de rang $r$.
Notre convention est que le {\it fibré projectif associé à
$\mathcal{E}$} est le morphisme
$$
\xymatrix{
\mathbf{P}(\mathcal{E}) =
\underline{\text{Proj}}_X(\text{Sym}^*(\mathcal{E}))
\ar[r]^-\pi
& X
}
$$
au-dessus de $X$, avec
$\mathcal{O}_{\mathbf{P}(\mathcal{E})}(1)$ normalisé de sorte que
$\pi_*(\mathcal{O}_{\mathbf{P}(\mathcal{E})}(1)) = \mathcal{E}$.
En particulier, il existe une surjection
$\pi^*\mathcal{E} \to \mathcal{O}_{\mathbf{P}(\mathcal{E})}(1)$.
Nous dirons informellement « soit $(\pi : P \to X, \mathcal{O}_P(1))$
le fibré projectif associé à $\mathcal{E}$ » pour désigner
la situation où $P = \mathbf{P}(\mathcal{E})$ et
$\mathcal{O}_P(1) = \mathcal{O}_{\mathbf{P}(\mathcal{E})}(1)$.

\begin{lemma}
\label{lemma-cap-projective-bundle}
Soit $(S, \delta)$ comme dans la situation \ref{situation-setup}.
Soit $X$ localement de type fini sur $S$.
Soit $\mathcal{E}$ un $\mathcal{O}_X$-module localement libre de type fini
$\mathcal{E}$ de rang $r$. Soit $(\pi : P \to X, \mathcal{O}_P(1))$
le fibré projectif associé à $\mathcal{E}$.
Pour tout $\alpha \in \CH_k(X)$, l'élément
$$
\pi_*\left(
c_1(\mathcal{O}_P(1))^s \cap \pi^*\alpha
\right)
\in
\CH_{k + r - 1 - s}(X)
$$
est $0$ si $s < r - 1$ et est égal à $\alpha$ lorsque $s = r - 1$.
\end{lemma}

\begin{proof}
Soit $Z \subset X$ un sous-schéma fermé intègre de $\delta$-dimension $k$.
Remarquons que $\pi^*[Z] = [\pi^{-1}(Z)]$, car $\pi^{-1}(Z)$ est intègre de
$\delta$-dimension $r - 1$.
Si $s < r - 1$, alors, par construction,
$c_1(\mathcal{O}_P(1))^s \cap \pi^*[Z]$
est représenté par un $(k + r - 1 - s)$-cycle supporté sur
$\pi^{-1}(Z)$. L'image directe de ce cycle
est donc nulle pour des raisons de dimension.

\medskip\noindent
Soit $s = r - 1$. L'argument donné ci-dessus montre que
$\pi_*(c_1(\mathcal{O}_P(1))^s \cap \pi^*\alpha) = n [Z]$
pour un certain $n \in \mathbf{Z}$. Nous voulons montrer que $n = 1$.
Pour les mêmes raisons de dimension
que ci-dessus, il suffit de démontrer ce résultat après avoir remplacé $X$ par
$X \setminus T$, où $T \subset Z$ est une partie fermée propre.
Soit $\xi$ le point générique de $Z$.
On peut choisir des éléments $e_1, \ldots, e_{r - 1} \in \mathcal{E}_\xi$
qui font partie d'une base de $\mathcal{E}_\xi$.
Ils donnent des sections rationnelles $s_1, \ldots, s_{r - 1}$
de $\mathcal{O}_P(1)|_{\pi^{-1}(Z)}$ dont le lieu commun des zéros
est l'adhérence de l'image d'une section rationnelle de
$\mathbf{P}(\mathcal{E}|_Z) \to Z$, réunie à une partie fermée dont le
support s'envoie dans une partie fermée propre $T$ de $Z$.
Après avoir retiré $T$ de $X$ (et, de façon correspondante, $\pi^{-1}(T)$
de $P$), on voit que $s_1, \ldots, s_n$ forment une suite
de sections globales
$s_i \in \Gamma(\pi^{-1}(Z), \mathcal{O}_{\pi^{-1}(Z)}(1))$
dont le lieu commun des zéros est l'image d'une section $Z \to \pi^{-1}(Z)$.
On voit donc successivement que
\begin{eqnarray*}
\pi^*[Z] & = & [\pi^{-1}(Z)] \\
c_1(\mathcal{O}_P(1)) \cap \pi^*[Z] & = & [Z(s_1)] \\
c_1(\mathcal{O}_P(1))^2 \cap \pi^*[Z] & = & [Z(s_1) \cap Z(s_2)] \\
\ldots & = & \ldots \\
c_1(\mathcal{O}_P(1))^{r - 1} \cap \pi^*[Z] & = &
[Z(s_1) \cap \ldots \cap Z(s_{r - 1})]
\end{eqnarray*}
par applications répétées du Lemme \ref{lemma-geometric-cap}.
Puisque l'image directe par $\pi$ de l'image d'une
section de $\pi$ au-dessus de $Z$ est clairement $[Z]$, on obtient le résultat
lorsque $\alpha = [Z]$. Nous omettons la vérification que ces
arguments donnent le résultat pour un cycle général $\alpha = \sum n_j [Z_j]$.
\end{proof}

\begin{lemma}[Formule du fibré projectif]
\label{lemma-chow-ring-projective-bundle}
Soit $(S, \delta)$ comme dans la situation \ref{situation-setup}.
Soit $X$ localement de type fini sur $S$.
Soit $\mathcal{E}$ un $\mathcal{O}_X$-module localement libre de type fini
$\mathcal{E}$ de rang $r$. Soit $(\pi : P \to X, \mathcal{O}_P(1))$
le fibré projectif associé à $\mathcal{E}$.
L'application
$$
\bigoplus\nolimits_{i = 0}^{r - 1}
\CH_{k + i}(X)
\longrightarrow
\CH_{k + r - 1}(P),
$$
$$
(\alpha_0, \ldots, \alpha_{r-1})
\longmapsto
\pi^*\alpha_0 +
c_1(\mathcal{O}_P(1)) \cap \pi^*\alpha_1
+ \ldots +
c_1(\mathcal{O}_P(1))^{r - 1} \cap \pi^*\alpha_{r-1}
$$
est un isomorphisme.
\end{lemma}

\begin{proof}
Fixons $k \in \mathbf{Z}$. Montrons d'abord que l'application est injective.
Supposons que $(\alpha_0, \ldots, \alpha_{r - 1})$ soit un élément
du membre gauche qui s'envoie sur zéro.
D'après le Lemme \ref{lemma-cap-projective-bundle}, on voit que
$$
0 = \pi_*(\pi^*\alpha_0 +
c_1(\mathcal{O}_P(1)) \cap \pi^*\alpha_1
+ \ldots +
c_1(\mathcal{O}_P(1))^{r - 1} \cap \pi^*\alpha_{r-1})
= \alpha_{r - 1}
$$
Ensuite, on voit que
$$
0 = \pi_*(c_1(\mathcal{O}_P(1)) \cap (\pi^*\alpha_0 +
c_1(\mathcal{O}_P(1)) \cap \pi^*\alpha_1
+ \ldots +
c_1(\mathcal{O}_P(1))^{r - 2} \cap \pi^*\alpha_{r - 2}))
= \alpha_{r - 2}
$$
et ainsi de suite. L'application est donc injective.

\medskip\noindent
Il reste à montrer que l'application est surjective.
Soient $X_i$, $i \in I$, les composantes irréductibles de $X$.
Alors $P_i = \mathbf{P}(\mathcal{E}|_{X_i})$, $i \in I$,
sont les composantes irréductibles de $P$. Considérons le
diagramme commutatif
$$
\xymatrix{
\coprod P_i \ar[d]_{\coprod \pi_i} \ar[r]_p & P \ar[d]^\pi \\
\coprod X_i \ar[r]^q & X
}
$$
Remarquons que $p_*$ est surjective. Si $\beta \in \CH_k(\coprod X_i)$,
alors $\pi^* q_* \beta = p_*(\coprod \pi_i)^* \beta$ ; voir le
Lemme \ref{lemma-flat-pullback-proper-pushforward}. Il en va de même pour
le produit cap avec $c_1(\mathcal{O}(1))$, d'après le
Lemme \ref{lemma-pushforward-cap-c1}.
Par conséquent, si l'application du lemme est surjective pour chacun
des morphismes $\pi_i : P_i \to X_i$, alors elle est
surjective pour $\pi : P \to X$. On peut donc supposer $X$ irréductible.
Ainsi, $\dim_\delta(X) < \infty$, et l'on peut en particulier raisonner
par récurrence sur $\dim_\delta(X)$.

\medskip\noindent
Le résultat est clair si $\dim_\delta(X) < k$.
Soit $\alpha \in \CH_{k + r - 1}(P)$.
Pour tout sous-schéma localement fermé $T \subset X$, notons
$\gamma_T : \bigoplus \CH_{k + i}(T) \to \CH_{k + r - 1}(\pi^{-1}(T))$
l'application
$$
\gamma_T(\alpha_0, \ldots, \alpha_{r - 1})
= \pi^*\alpha_0 + \ldots +
c_1(\mathcal{O}_{\pi^{-1}(T)}(1))^{r - 1} \cap \pi^*\alpha_{r - 1}.
$$
Supposons que, pour un certain ouvert non vide $U \subset X$, on ait
$\alpha|_{\pi^{-1}(U)} = \gamma_U(\alpha_0, \ldots, \alpha_{r - 1})$.
On peut alors choisir des relèvements $\alpha'_i \in \CH_{k + i}(X)$, et l'on
voit que $\alpha - \gamma_X(\alpha'_0, \ldots, \alpha'_{r - 1})$
est, d'après le Lemme \ref{lemma-restrict-to-open},
rationnellement équivalent à un $k$-cycle sur $P_Y = \mathbf{P}(\mathcal{E}|_Y)$,
où $Y = X \setminus U$ est muni de sa structure de sous-schéma fermé réduit.
Remarquons que $\dim_\delta(Y) < \dim_\delta(X)$.
Par récurrence, le résultat vaut
pour $P_Y \to Y$, et il vaut donc pour $\alpha$.
On peut par conséquent remplacer $X$ par n'importe quel ouvert non vide de $X$.

\medskip\noindent
En particulier, on peut supposer que $\mathcal{E} \cong \mathcal{O}_X^{\oplus r}$.
Dans ce cas, $\mathbf{P}(\mathcal{E}) = X \times \mathbf{P}^{r - 1}$.
Utilisons la stratification
$$
\mathbf{P}^{r - 1} = \mathbf{A}^{r - 1}
\amalg \mathbf{A}^{r - 2}
\amalg \ldots
\amalg \mathbf{A}^0
$$
L'adhérence de chaque strate est un $\mathbf{P}^{r - 1 - i}$ qui représente
$c_1(\mathcal{O}(1))^i \cap [\mathbf{P}^{r - 1}]$.
Ainsi, $P$ possède une stratification analogue
$$
P = U^{r - 1} \amalg U^{r - 2} \amalg \ldots \amalg U^0
$$
Soit $P^i$ l'adhérence de $U^i$. Soit $\pi^i : P^i \to X$
la restriction de $\pi$ à $P^i$.
Soit $\alpha \in \CH_{k + r - 1}(P)$. D'après le
Lemme \ref{lemma-pullback-affine-fibres-surjective},
on peut écrire $\alpha|_{U^{r - 1}} = \pi^*\alpha_0|_{U^{r - 1}}$
pour un certain $\alpha_0 \in \CH_k(X)$. La différence
$\alpha - \pi^*\alpha_0$ est donc l'image d'un certain
$\alpha' \in \CH_{k + r - 1}(P^{r - 2})$.
En appliquant de nouveau le Lemme \ref{lemma-pullback-affine-fibres-surjective},
on peut écrire
$\alpha'|_{U^{r - 2}} = (\pi^{r - 2})^*\alpha_1|_{U^{r - 2}}$
pour un certain $\alpha_1 \in \CH_{k + 1}(X)$.
D'après le Lemme \ref{lemma-relative-effective-cartier},
on voit que l'image de $(\pi^{r - 2})^*\alpha_1$
représente $c_1(\mathcal{O}_P(1)) \cap \pi^*\alpha_1$.
On voit également que
$\alpha - \pi^*\alpha_0 - c_1(\mathcal{O}_P(1)) \cap \pi^*\alpha_1$
est l'image d'un certain $\alpha'' \in \CH_{k + r - 1}(P^{r - 3})$.
Et ainsi de suite.
\end{proof}

\begin{lemma}
\label{lemma-vectorbundle}
Soit $(S, \delta)$ comme dans la situation \ref{situation-setup}.
Soit $X$ localement de type fini sur $S$.
Soit $\mathcal{E}$ un faisceau localement libre de type fini de rang $r$ sur $X$.
Soit
$$
p :
E = \underline{\Spec}(\text{Sym}^*(\mathcal{E}))
\longrightarrow
X
$$
le fibré vectoriel associé au-dessus de $X$.
Alors $p^* : \CH_k(X) \to \CH_{k + r}(E)$ est un isomorphisme pour tout $k$.
\end{lemma}

\begin{proof}
(Pour le cas des fibrés en droites, voir le Lemme \ref{lemma-linebundle}.)
Pour la surjectivité, voir le Lemme \ref{lemma-pullback-affine-fibres-surjective}.
Soit $(\pi  : P \to X, \mathcal{O}_P(1))$
le fibré projectif associé
au faisceau localement libre de type fini $\mathcal{E} \oplus \mathcal{O}_X$.
Soit $s \in \Gamma(P, \mathcal{O}_P(1))$ correspondant à la section globale
$(0, 1) \in \Gamma(X, \mathcal{E} \oplus \mathcal{O}_X)$.
Soit $D = Z(s) \subset P$. Remarquons que
$(\pi|_D : D \to X , \mathcal{O}_P(1)|_D)$
est le fibré projectif associé
à $\mathcal{E}$. On note $\pi_D = \pi|_D$ et
$\mathcal{O}_D(1) = \mathcal{O}_P(1)|_D$.
De plus, $D$ est un diviseur de Cartier effectif
sur $P$. Ainsi, $\mathcal{O}_P(D) = \mathcal{O}_P(1)$
(voir Diviseurs, Lemme \ref{divisors-lemma-characterize-OD}).
Il existe aussi un isomorphisme
$E \cong P \setminus D$. Notons $j : E \to P$
l'immersion ouverte correspondante.
Pour l'injectivité, on utilise que le noyau de
$$
j^* :
\CH_{k + r}(P)
\longrightarrow
\CH_{k + r}(E)
$$
est formé des cycles supportés dans le diviseur de Cartier effectif $D$ ;
voir le Lemme \ref{lemma-restrict-to-open}. Ainsi, si $p^*\alpha = 0$, alors
$\pi^*\alpha = i_*\beta$ pour un certain $\beta \in \CH_{k + r}(D)$.
D'après le Lemme \ref{lemma-chow-ring-projective-bundle}, on peut écrire
$$
\beta = \pi_D^*\beta_0 +
\ldots + c_1(\mathcal{O}_D(1))^{r - 1} \cap \pi_D^* \beta_{r - 1}.
$$
pour certains $\beta_i \in \CH_{k + i}(X)$.
D'après les Lemmes \ref{lemma-relative-effective-cartier}
et \ref{lemma-pushforward-cap-c1},
cela implique
$$
\pi^*\alpha = i_*\beta =
c_1(\mathcal{O}_P(1)) \cap \pi^*\beta_0 +
\ldots +
c_1(\mathcal{O}_D(1))^r \cap \pi^*\beta_{r - 1}.
$$
Puisque le rang de $\mathcal{E} \oplus \mathcal{O}_X$ est $r + 1$,
cela contredit le Lemme \ref{lemma-pushforward-cap-c1}, à moins que tous les
$\alpha$ et tous les $\beta_i$ ne soient nuls.
\end{proof}








\section{Les classes de Chern d'un fibré vectoriel}
\label{section-chern-classes-vector-bundles}

\noindent
Nous pouvons utiliser la formule du fibré en espaces projectifs pour définir les
classes de Chern d'un fibré vectoriel de rang $r$ au moyen du développement
de $c_1(\mathcal{O}(1))^r$ suivant les puissances inférieures, voir
la formule (\ref{equation-chern-classes}).
La raison des signes sera expliquée plus loin.

\begin{definition}
\label{definition-chern-classes}
Soit $(S, \delta)$ comme dans la Situation \ref{situation-setup}.
Soit $X$ localement de type fini sur $S$.
Supposons $X$ intègre et $n = \dim_\delta(X)$.
Soit $\mathcal{E}$ un faisceau localement libre de type fini de rang $r$
sur $X$. Soit $(\pi : P \to X, \mathcal{O}_P(1))$ le fibré en espaces
projectifs associé à $\mathcal{E}$.
\begin{enumerate}
\item D'après le Lemme \ref{lemma-chow-ring-projective-bundle}, il existe des
éléments $c_i \in \CH_{n - i}(X)$, $i = 0, \ldots, r$
tels que $c_0 = [X]$, et
\begin{equation}
\label{equation-chern-classes}
\sum\nolimits_{i = 0}^r
(-1)^i c_1(\mathcal{O}_P(1))^i \cap \pi^*c_{r - i}
= 0.
\end{equation}
\item Avec les notations ci-dessus, nous posons
$c_i(\mathcal{E}) \cap [X] = c_i$
comme élément de $\CH_{n - i}(X)$.
Nous les appelons les {\it classes de Chern de $\mathcal{E}$ sur $X$}.
\item La {\it classe de Chern totale de $\mathcal{E}$ sur $X$}
est la somme
$$
c({\mathcal E}) \cap [X] =
c_0({\mathcal E}) \cap [X]
+ c_1({\mathcal E}) \cap [X] + \ldots
+ c_r({\mathcal E}) \cap [X]
$$
qui est un élément de
$\CH_*(X) = \bigoplus_{k \in \mathbf{Z}} \CH_k(X)$.
\end{enumerate}
\end{definition}

\noindent
Vérifions que cela ne donne pas une notion nouvelle lorsque le
fibré vectoriel est de rang $1$.

\begin{lemma}
\label{lemma-first-chern-class}
Soit $(S, \delta)$ comme dans la Situation \ref{situation-setup}.
Soit $X$ localement de type fini sur $S$.
Supposons $X$ intègre et $n = \dim_\delta(X)$.
Soit $\mathcal{L}$ un $\mathcal{O}_X$-module inversible.
La première classe de Chern de $\mathcal{L}$ sur $X$ de la
Définition \ref{definition-chern-classes}
est égale au diviseur de Weil associé à $\mathcal{L}$
par la Définition \ref{definition-divisor-invertible-sheaf}.
\end{lemma}

\begin{proof}
Dans cette démonstration, nous employons $c_1(\mathcal{L}) \cap [X]$ pour désigner la
construction de la Définition \ref{definition-divisor-invertible-sheaf}.
Comme $\mathcal{L}$ est de rang $1$, nous avons
$\mathbf{P}(\mathcal{L}) = X$ et
$\mathcal{O}_{\mathbf{P}(\mathcal{L})}(1) = \mathcal{L}$
d'après nos normalisations. Ainsi, (\ref{equation-chern-classes})
s'écrit
$$
(-1)^1 c_1(\mathcal{L}) \cap c_0 + (-1)^0 c_1 = 0
$$
Puisque $c_0 = [X]$, nous concluons que $c_1 = c_1(\mathcal{L}) \cap [X]$,
comme voulu.
\end{proof}

\begin{remark}
\label{remark-equation-signs}
Nous pourrions aussi récrire l'équation \ref{equation-chern-classes} sous la forme
\begin{equation}
\label{equation-signs}
\sum\nolimits_{i = 0}^r
c_1(\mathcal{O}_P(-1))^i \cap \pi^*c_{r - i}
= 0.
\end{equation}
mais il nous paraît plus facile de travailler avec le faisceau quotient
tautologique $\mathcal{O}_P(1)$ plutôt qu'avec
son dual.
\end{remark}




\section{Intersection avec les classes de Chern}
\label{section-intersecting-chern-classes}

\noindent
Dans cette section, nous définissons les classes de Chern des fibrés vectoriels sur $X$ comme
classes bivariantes sur $X$, voir le Lemme \ref{lemma-cap-cp-bivariant}
et la discussion qui le suit. Notre construction suit le procédé familier
consistant à définir d'abord l'opération sur les cycles premiers, puis à
sommer. Dans le Lemme \ref{lemma-determine-intersections}, nous montrons
que le résultat est déterminé par la formule usuelle sur le fibré
projectif associé. Ensuite, nous montrons que l'intersection avec les classes de Chern
passe à l'équivalence rationnelle, commute à l'image directe propre,
commute à l'image inverse par un morphisme plat et commute aux homomorphismes de Gysin pour les
inclusions de diviseurs de Cartier effectifs. On aurait pu éviter ces lemmes
en utilisant directement la caractérisation du
Lemme \ref{lemma-determine-intersections} et le
Lemme \ref{lemma-push-proper-bivariant} ; le lecteur qui souhaite
en voir les détails pourra consulter
Groupes de Chow des espaces, Lemme \ref{spaces-chow-lemma-segre-classes}.

\begin{definition}
\label{definition-cap-chern-classes}
Soit $(S, \delta)$ comme dans la Situation \ref{situation-setup}.
Soit $X$ localement de type fini sur $S$.
Soit $\mathcal{E}$ un faisceau localement libre de type fini de rang $r$ sur $X$.
Nous définissons, pour tout entier $k$ et tout $0 \leq j \leq r$,
une opération
$$
c_j(\mathcal{E}) \cap - : Z_k(X) \to \CH_{k - j}(X)
$$
appelée {\it intersection avec la $j$-ième classe de Chern de $\mathcal{E}$}.
\begin{enumerate}
\item Étant donné un sous-schéma fermé intègre $i : W \to X$ de $\delta$-dimension
$k$, nous définissons
$$
c_j(\mathcal{E}) \cap [W] = i_*(c_j({i^*\mathcal{E}}) \cap [W])
\in
\CH_{k - j}(X)
$$
où $c_j({i^*\mathcal{E}}) \cap [W]$ est défini comme dans la
Définition \ref{definition-chern-classes}.
\item Pour un $k$-cycle quelconque $\alpha = \sum n_i [W_i]$, nous posons
$$
c_j(\mathcal{E}) \cap \alpha = \sum n_i c_j(\mathcal{E}) \cap [W_i]
$$
\end{enumerate}
\end{definition}

\noindent
Si $\mathcal{E}$ est de rang $1$, cela coïncide avec notre
définition précédente (Définition \ref{definition-cap-c1})
d'après le Lemme \ref{lemma-first-chern-class}.

\begin{lemma}
\label{lemma-determine-intersections}
Soit $(S, \delta)$ comme dans la Situation \ref{situation-setup}.
Soit $X$ localement de type fini sur $S$.
Soit $\mathcal{E}$ un faisceau localement libre de type fini de rang $r$ sur $X$.
Soit $(\pi : P \to X, \mathcal{O}_P(1))$ le fibré projectif
associé à $\mathcal{E}$.
Pour $\alpha \in Z_k(X)$, les éléments
$c_j(\mathcal{E}) \cap \alpha$ sont les uniques éléments
$\alpha_j$ de $\CH_{k - j}(X)$
tels que $\alpha_0 = \alpha$ et que
$$
\sum\nolimits_{i = 0}^r
(-1)^i c_1(\mathcal{O}_P(1))^i \cap
\pi^*(\alpha_{r - i}) = 0
$$
dans le groupe de Chow de $P$.
\end{lemma}

\begin{proof}
L'unicité de $\alpha_0, \ldots, \alpha_r$ tels que
$\alpha_0 = \alpha$ et que
l'équation affichée soit satisfaite résulte de
la formule du fibré en espaces projectifs,
Lemme \ref{lemma-chow-ring-projective-bundle}.
Par définition, l'identité est satisfaite pour $\alpha = [W]$, où $W$
est un sous-schéma fermé intègre de $X$.
Pour un $k$-cycle quelconque $\alpha$ sur $X$, écrivons
$\alpha = \sum n_a[W_a]$ avec $n_a \not = 0$, et
$i_a : W_a \to X$ des sous-schémas fermés intègres deux à deux distincts.
Alors la famille $\{W_a\}$ est localement finie sur $X$.
Posons $P_a = \pi^{-1}(W_a) = \mathbf{P}(\mathcal{E}|_{W_a})$.
Notons $i'_a : P_a \to P$ les immersions fermées correspondantes.
Considérons le diagramme de produit fibré
$$
\xymatrix{
P' \ar@{=}[r] \ar[d]_{\pi'} &
\coprod P_a \ar[d]_{\coprod \pi_a} \ar[r]_{\coprod i'_a} &
P \ar[d]^\pi \\
X' \ar@{=}[r] &
\coprod W_a \ar[r]^{\coprod i_a} &
X
}
$$
Le morphisme $p : X' \to X$ est propre. De plus,
$\pi' : P' \to X'$, muni du faisceau inversible
$\mathcal{O}_{P'}(1) = \coprod \mathcal{O}_{P_a}(1)$,
qui est également l'image inverse de $\mathcal{O}_P(1)$,
est le fibré projectif associé à
$\mathcal{E}' = p^*\mathcal{E}$. Par définition,
$$
c_j(\mathcal{E}) \cap [\alpha]
=
\sum i_{a, *}(c_j(\mathcal{E}|_{W_a}) \cap [W_a]).
$$
Écrivons $\beta_{a, j} = c_j(\mathcal{E}|_{W_a}) \cap [W_a]$,
qui est un élément de $\CH_{k - j}(W_a)$. Nous avons
$$
\sum\nolimits_{i = 0}^r
(-1)^i c_1(\mathcal{O}_{P_a}(1))^i \cap \pi_a^*(\beta_{a, r - i}) = 0
$$
pour tout $a$, par définition. Il est donc clair que nous avons
$$
\sum\nolimits_{i = 0}^r
(-1)^i c_1(\mathcal{O}_{P'}(1))^i \cap (\pi')^*(\beta_{r - i}) = 0
$$
avec $\beta_j = \sum n_a\beta_{a, j} \in \CH_{k - j}(X')$. Notons
$p' : P' \to P$ le morphisme $\coprod i'_a$.
Nous avons $\pi^*p_*\beta_j = p'_*(\pi')^*\beta_j$
d'après le Lemme \ref{lemma-flat-pullback-proper-pushforward}.
Par la formule de projection du Lemme \ref{lemma-pushforward-cap-c1},
nous concluons que
$$
\sum\nolimits_{i = 0}^r
(-1)^i c_1(\mathcal{O}_P(1))^i \cap \pi^*(p_*\beta_j) = 0
$$
Puisque $p_*\beta_j$ est un représentant de $c_j(\mathcal{E}) \cap \alpha$,
le résultat s'ensuit.
\end{proof}

\noindent
Nous utiliserons systématiquement cette caractérisation des classes de Chern
pour en démontrer bien d'autres propriétés.

\begin{lemma}
\label{lemma-cap-chern-class-factors-rational-equivalence}
Soit $(S, \delta)$ comme dans la Situation \ref{situation-setup}.
Soit $X$ localement de type fini sur $S$.
Soit $\mathcal{E}$ un faisceau localement libre de type fini de rang $r$ sur $X$.
Si $\alpha \sim_{rat} \beta$ sont des $k$-cycles rationnellement équivalents
sur $X$, alors $c_j(\mathcal{E}) \cap \alpha = c_j(\mathcal{E}) \cap \beta$
dans $\CH_{k - j}(X)$.
\end{lemma}

\begin{proof}
D'après le Lemme \ref{lemma-determine-intersections}, les éléments
$\alpha_j = c_j(\mathcal{E}) \cap \alpha$, $j \geq 1$, et
$\beta_j = c_j(\mathcal{E}) \cap \beta$, $j \geq 1$, sont uniquement déterminés
par la {\it même} équation dans le groupe de Chow du fibré
projectif associé à $\mathcal{E}$. (Cela repose bien entendu sur le fait que
l'image inverse par un morphisme plat est compatible avec l'équivalence rationnelle, voir le
Lemme \ref{lemma-flat-pullback-rational-equivalence}.) Ils sont donc égaux.
\end{proof}

\noindent
Autrement dit, l'intersection avec les classes de Chern des
faisceaux localement libres de type fini passe à l'équivalence rationnelle
et donne des applications
$$
c_j(\mathcal{E}) \cap - : \CH_k(X) \to \CH_{k - j}(X).
$$
Notre tâche suivante consiste à montrer que les classes de Chern sont des classes bivariantes, voir
la Définition \ref{definition-bivariant-class}.

\begin{lemma}
\label{lemma-pushforward-cap-cj}
Soit $(S, \delta)$ comme dans la Situation \ref{situation-setup}.
Soient $X$, $Y$ localement de type fini sur $S$.
Soit $\mathcal{E}$ un faisceau localement libre de type fini de rang $r$ sur $X$.
Soit $p : X \to Y$ un morphisme propre.
Soit $\alpha$ un $k$-cycle sur $X$.
Soit $\mathcal{E}$ un faisceau localement libre de type fini sur $Y$.
Alors
$$
p_*(c_j(p^*\mathcal{E}) \cap \alpha) = c_j(\mathcal{E}) \cap p_*\alpha
$$
\end{lemma}

\begin{proof}
Soit $(\pi : P \to Y, \mathcal{O}_P(1))$ le fibré projectif associé
à $\mathcal{E}$. Alors $P_X = X \times_Y P$ est le fibré projectif associé
à $p^*\mathcal{E}$ et $\mathcal{O}_{P_X}(1)$ est l'image inverse de
$\mathcal{O}_P(1)$. Écrivons $\alpha_j = c_j(p^*\mathcal{E}) \cap \alpha$, de sorte que
$\alpha_0 = \alpha$. D'après le Lemme \ref{lemma-determine-intersections}, nous avons
$$
\sum\nolimits_{i = 0}^r
(-1)^i c_1(\mathcal{O}_P(1))^i \cap
\pi_X^*(\alpha_{r - i}) = 0
$$
dans le groupe de Chow de $P_X$. Considérons le diagramme de produit fibré
$$
\xymatrix{
P_X \ar[r]_-{p'} \ar[d]_{\pi_X} & P \ar[d]^\pi \\
X \ar[r]^p & Y
}
$$
Appliquons l'image directe propre $p'_*$,
(Lemme \ref{lemma-proper-pushforward-rational-equivalence})
à l'égalité affichée ci-dessus. En utilisant les
Lemmes \ref{lemma-pushforward-cap-c1} et
\ref{lemma-flat-pullback-proper-pushforward}, nous obtenons
$$
\sum\nolimits_{i = 0}^r
(-1)^i c_1(\mathcal{O}_P(1))^i \cap
\pi^*(p_*\alpha_{r - i}) = 0
$$
dans le groupe de Chow de $P$. Par la caractérisation du
Lemme \ref{lemma-determine-intersections}, nous concluons.
\end{proof}

\begin{lemma}
\label{lemma-flat-pullback-cap-cj}
Soit $(S, \delta)$ comme dans la Situation \ref{situation-setup}.
Soient $X$, $Y$ localement de type fini sur $S$.
Soit $\mathcal{E}$ un faisceau localement libre de type fini de rang $r$ sur $Y$.
Soit $f : X \to Y$ un morphisme plat de dimension relative $r$.
Soit $\alpha$ un $k$-cycle sur $Y$.
Alors
$$
f^*(c_j(\mathcal{E}) \cap \alpha) = c_j(f^*\mathcal{E}) \cap f^*\alpha
$$
\end{lemma}

\begin{proof}
Écrivons $\alpha_j = c_j(\mathcal{E}) \cap \alpha$, de sorte que $\alpha_0 = \alpha$.
D'après le Lemme \ref{lemma-determine-intersections}, nous avons
$$
\sum\nolimits_{i = 0}^r
(-1)^i c_1(\mathcal{O}_P(1))^i \cap
\pi^*(\alpha_{r - i}) = 0
$$
dans le groupe de Chow du fibré projectif
$(\pi : P \to Y, \mathcal{O}_P(1))$
associé à $\mathcal{E}$. Considérons le diagramme de produit fibré
$$
\xymatrix{
P_X = \mathbf{P}(f^*\mathcal{E}) \ar[r]_-{f'} \ar[d]_{\pi_X} &
P \ar[d]^\pi \\
X \ar[r]^f & Y
}
$$
Notons que $\mathcal{O}_{P_X}(1)$ est l'image inverse de $\mathcal{O}_P(1)$.
Appliquons l'image inverse plate $(f')^*$,
(Lemme \ref{lemma-flat-pullback-rational-equivalence}) à l'équation
affichée ci-dessus. D'après les Lemmes \ref{lemma-flat-pullback-cap-c1} et
\ref{lemma-compose-flat-pullback}, nous voyons que
$$
\sum\nolimits_{i = 0}^r
(-1)^i c_1(\mathcal{O}_{P_X}(1))^i \cap
\pi_X^*(f^*\alpha_{r - i}) = 0
$$
dans le groupe de Chow de $P_X$. Par la caractérisation du
Lemme \ref{lemma-determine-intersections}, nous concluons.
\end{proof}

\begin{lemma}
\label{lemma-cap-chern-class-commutes-with-gysin}
Soit $(S, \delta)$ comme dans la Situation \ref{situation-setup}.
Soit $X$ localement de type fini sur $S$.
Soit $\mathcal{E}$ un faisceau localement libre de type fini de rang $r$ sur $X$.
Soit $(\mathcal{L}, s, i : D \to X)$ comme dans la
Définition \ref{definition-gysin-homomorphism}.
Alors $c_j(\mathcal{E}|_D) \cap i^*\alpha = i^*(c_j(\mathcal{E}) \cap \alpha)$
pour tout $\alpha \in \CH_k(X)$.
\end{lemma}

\begin{proof}
Écrivons $\alpha_j = c_j(\mathcal{E}) \cap \alpha$, de sorte que $\alpha_0 = \alpha$.
D'après le Lemme \ref{lemma-determine-intersections}, nous avons
$$
\sum\nolimits_{i = 0}^r
(-1)^i c_1(\mathcal{O}_P(1))^i \cap
\pi^*(\alpha_{r - i}) = 0
$$
dans le groupe de Chow du fibré projectif
$(\pi : P \to X, \mathcal{O}_P(1))$
associé à $\mathcal{E}$. Considérons le diagramme de produit fibré
$$
\xymatrix{
P_D = \mathbf{P}(\mathcal{E}|_D) \ar[r]_-{i'} \ar[d]_{\pi_D} &
P \ar[d]^\pi \\
D \ar[r]^i & X
}
$$
Notons que $\mathcal{O}_{P_D}(1)$ est l'image inverse de $\mathcal{O}_P(1)$.
Appliquons l'homomorphisme de Gysin $(i')^*$ (Lemme \ref{lemma-gysin-factors}) à
l'équation affichée ci-dessus.
En appliquant les Lemmes \ref{lemma-gysin-commutes-cap-c1} et
\ref{lemma-gysin-flat-pullback}, nous obtenons
$$
\sum\nolimits_{i = 0}^r
(-1)^i c_1(\mathcal{O}_{P_D}(1))^i \cap
\pi_D^*(i^*\alpha_{r - i}) = 0
$$
dans le groupe de Chow de $P_D$.
Par la caractérisation du Lemme \ref{lemma-determine-intersections},
nous concluons.
\end{proof}

\noindent
Nous disposons maintenant de suffisamment d'éléments pour démontrer que
l'intersection avec les classes de Chern définit une classe bivariante.

\begin{lemma}
\label{lemma-cap-cp-bivariant}
Soit $(S, \delta)$ comme dans la Situation \ref{situation-setup}.
Soit $X$ localement de type fini sur $S$.
Soit $\mathcal{E}$ un $\mathcal{O}_X$-module localement libre
de rang $r$. Soit $0 \leq p \leq r$.
Alors la règle qui, à $f : X' \to X$, associe
$c_p(f^*\mathcal{E}) \cap - : \CH_k(X') \to \CH_{k - p}(X')$
est une classe bivariante de degré $p$.
\end{lemma}

\begin{proof}
Cela résulte immédiatement des Lemmes
\ref{lemma-cap-chern-class-factors-rational-equivalence},
\ref{lemma-pushforward-cap-cj},
\ref{lemma-flat-pullback-cap-cj}, et
\ref{lemma-cap-chern-class-commutes-with-gysin},
ainsi que de la Définition \ref{definition-bivariant-class}.
\end{proof}

\noindent
Ce lemme nous permet de définir comme suit les classes de Chern d'un module
localement libre de type fini.

\begin{definition}
\label{definition-chern-classes-final}
Soit $(S, \delta)$ comme dans la Situation \ref{situation-setup}.
Soit $X$ localement de type fini sur $S$.
Soit $\mathcal{E}$ un $\mathcal{O}_X$-module localement libre
de rang $r$. Pour $i = 0, \ldots, r$, la {\it $i$-ième classe de Chern}
de $\mathcal{E}$ est la classe bivariante
$c_i(\mathcal{E}) \in A^i(X)$ de degré $i$
construite dans le Lemme \ref{lemma-cap-cp-bivariant}. La
{\it classe de Chern totale} de $\mathcal{E}$ est la somme formelle
$$
c(\mathcal{E}) = 
c_0(\mathcal{E}) + c_1(\mathcal{E}) + \ldots + c_r(\mathcal{E})
$$
considérée comme une classe bivariante non homogène sur $X$.
\end{definition}

\noindent
D'après la remarque qui suit la Définition \ref{definition-cap-chern-classes},
si $\mathcal{E}$ est inversible, cette définition coïncide avec la
Définition \ref{definition-first-chern-class}.
Nous voyons ensuite que les classes de Chern appartiennent au centre de l'anneau de
cohomologie bivariante de Chow $A^*(X)$.

\begin{lemma}
\label{lemma-cap-commutative-chern}
Soit $(S, \delta)$ comme dans la Situation \ref{situation-setup}.
Soit $X$ localement de type fini sur $S$.
Soit $\mathcal{E}$ un $\mathcal{O}_X$-module localement libre de rang $r$.
Alors
\begin{enumerate}
\item $c_j(\mathcal{E}) \in A^j(X)$ appartient au centre de $A^*(X)$, et
\item si $f : X' \to X$ est localement de type fini et $c \in A^*(X' \to X)$,
alors $c \circ c_j(\mathcal{E}) = c_j(f^*\mathcal{E}) \circ c$.
\end{enumerate}
En particulier, si $\mathcal{F}$ est un second
$\mathcal{O}_X$-module localement libre sur $X$ de rang $s$, alors
$$
c_i(\mathcal{E}) \cap c_j(\mathcal{F}) \cap \alpha
=
c_j(\mathcal{F}) \cap c_i(\mathcal{E}) \cap \alpha
$$
comme éléments de $\CH_{k - i - j}(X)$ pour tout $\alpha \in \CH_k(X)$.
\end{lemma}

\begin{proof}
Il est immédiat que (2) implique (1).
Soit $\alpha \in \CH_k(X)$. Écrivons $\alpha_j = c_j(\mathcal{E}) \cap \alpha$, de sorte que
$\alpha_0 = \alpha$. D'après le Lemme \ref{lemma-determine-intersections}, nous avons
$$
\sum\nolimits_{i = 0}^r
(-1)^i c_1(\mathcal{O}_P(1))^i \cap
\pi^*(\alpha_{r - i}) = 0
$$
dans le groupe de Chow du fibré projectif
$(\pi : P \to Y, \mathcal{O}_P(1))$
associé à $\mathcal{E}$. Notons $\pi' : P' \to X'$ le changement de base
de $\pi$ par $f$. En utilisant le Lemme \ref{lemma-c1-center} et
les propriétés des classes bivariantes, nous obtenons
\begin{align*}
0 & = c \cap \left(\sum\nolimits_{i = 0}^r
(-1)^i c_1(\mathcal{O}_P(1))^i \cap
\pi^*(\alpha_{r - i})\right) \\
& =
\sum\nolimits_{i = 0}^r
(-1)^i c_1(\mathcal{O}_{P'}(1))^i \cap
(\pi')^*(c \cap \alpha_{r - i})
\end{align*}
dans le groupe de Chow de $P'$ (calcul omis).
Nous voyons donc que $c \cap \alpha_j$ est
égal à $c_j(f^*\mathcal{E}) \cap (c \cap \alpha)$ par la caractérisation
du Lemme \ref{lemma-determine-intersections}.
Cela démontre le lemme.
\end{proof}

\begin{remark}
\label{remark-extend-to-finite-locally-free}
Soit $(S, \delta)$ comme dans la Situation \ref{situation-setup}.
Soit $X$ localement de type fini sur $S$.
Soit $\mathcal{E}$ un $\mathcal{O}_X$-module localement libre de type fini.
Si le rang de $\mathcal{E}$ n'est pas constant, nous pouvons
encore définir les classes de Chern de $\mathcal{E}$. En effet, dans ce
cas, nous pouvons écrire
$$
X = X_0 \amalg X_1 \amalg X_2 \amalg \ldots
$$
où $X_r \subset X$ est le sous-espace ouvert et fermé sur lequel
le rang de $\mathcal{E}$ est $r$. D'après le
Lemme \ref{lemma-disjoint-decomposition-bivariant},
nous avons $A^p(X) = \prod A^p(X_r)$.
Nous pouvons donc définir $c_p(\mathcal{E})$ comme le
produit des classes $c_p(\mathcal{E}|_{X_r})$ dans $A^p(X_r)$.
Explicitement, si $X' \to X$ est un morphisme localement de type fini,
nous obtenons par image inverse une décomposition correspondante de $X'$,
et nous trouvons que
$$
\CH_*(X') = \prod\nolimits_{r \geq 0} \CH_*(X'_r)
$$
d'après nos définitions. Alors $c_p(\mathcal{E}) \in A^p(X)$
est la classe bivariante qui préserve ces décompositions
en produits directs et agit par les opérations déjà définies
$c_i(\mathcal{E}|_{X_r}) \cap -$
sur les facteurs. Observons que, dans ce cadre, il peut arriver
que $c_p(\mathcal{E})$ soit non nul pour une infinité de $p$.
Il s'ensuit que la classe de Chern totale est un élément
$$
c(\mathcal{E}) =
c_0(\mathcal{E}) + c_1(\mathcal{E}) + c_2(\mathcal{E}) + \ldots
\in A^*(X)^\wedge
$$
de l'anneau de cohomologie bivariante complété, voir la
Remarque \ref{remark-completion-bivariant}.
Dans ce cadre, nous définissons le « rang » de $\mathcal{E}$
comme l'élément $r(\mathcal{E}) \in A^0(X)$,
à savoir l'opération bivariante qui envoie $(\alpha_r) \in \prod \CH_*(X'_r)$
sur $(r\alpha_r) \in \prod \CH_*(X'_r)$.
Notons qu'il reste vrai que $c_p(\mathcal{E})$ et $r(\mathcal{E})$
appartiennent au centre de $A^*(X)$.
\end{remark}

\begin{remark}
\label{remark-top-chern-class}
Soit $(S, \delta)$ comme dans la Situation \ref{situation-setup}. Soit $X$
localement de type fini sur $S$. Soit $\mathcal{E}$ un
$\mathcal{O}_X$-module localement libre de type fini. En général,
nous écrivons $X = \coprod X_r$ comme dans la
Remarque \ref{remark-extend-to-finite-locally-free}.
Si seul un nombre fini des $X_r$ sont non vides, alors
nous pouvons poser
$$
c_{top}(\mathcal{E}) = \sum\nolimits_r c_r(\mathcal{E}|_{X_r})
\in A^*(X) = \bigoplus A^*(X_r)
$$
où l'égalité est le Lemme \ref{lemma-disjoint-decomposition-bivariant}.
Si une infinité des $X_r$ sont non vides, nous emploierons la même
notation pour désigner
$$
c_{top}(\mathcal{E}) = \prod c_r(\mathcal{E}|_{X_r})
\in \prod A^r(X_r) \subset A^*(X)^\wedge
$$
voir la Remarque \ref{remark-completion-bivariant} pour les notations.
\end{remark}











\section{Relations polynomiales entre classes de Chern}
\label{section-relations-chern-classes}

\noindent
Soit $(S, \delta)$ comme dans la Situation \ref{situation-setup}. Soit $X$ localement
de type fini sur $S$. Les $\mathcal{E}_i$ forment une famille finie de faisceaux
localement libres de type fini sur $X$. D'après le Lemme \ref{lemma-cap-commutative-chern},
nous voyons que les classes de Chern
$$
c_j(\mathcal{E}_i) \in A^*(X)
$$
engendrent une $\mathbf{Z}$-sous-algèbre commutative (et même centrale) de
l'algèbre de cohomologie de Chow $A^*(X)$.
Nous pouvons donc préciser ce que signifie qu'un polynôme en ces classes de Chern
est nul, ou que deux polynômes sont égaux. Par exemple, dire que
$c_1(\mathcal{E}_1)^5 + c_2(\mathcal{E}_2)c_3(\mathcal{E}_3) = 0$
signifie que les opérations
$$
\CH_k(Y) \longrightarrow \CH_{k - 5}(Y), \quad
\alpha \longmapsto
c_1(\mathcal{E}_1)^5 \cap \alpha +
c_2(\mathcal{E}_2) \cap c_3(\mathcal{E}_3) \cap \alpha
$$
sont nulles pour tout morphisme $f : Y \to X$ localement de type fini.
D'après le Lemme \ref{lemma-bivariant-zero},
cela équivaut à exiger que, pour tout morphisme
$f : Y \to X$  où $Y$ est un schéma intègre
localement de type fini sur $S$, le cycle
$$
c_1(\mathcal{E}_1)^5 \cap [Y] +
c_2(\mathcal{E}_2) \cap c_3(\mathcal{E}_3) \cap [Y]
$$
soit nul dans $\CH_{\dim(Y) - 5}(Y)$.

\medskip\noindent
Un exemple particulier est la relation
$$
c_1(\mathcal{L} \otimes_{\mathcal{O}_X} \mathcal{N})
=
c_1(\mathcal{L}) + c_1(\mathcal{N})
$$
démontrée dans le Lemme \ref{lemma-c1-cap-additive}.
Plus généralement, voici ce qui se passe lorsque l'on tensorise un
faisceau localement libre quelconque par un faisceau inversible.

\begin{lemma}
\label{lemma-chern-classes-E-tensor-L}
Soit $(S, \delta)$ comme dans la Situation \ref{situation-setup}.
Soit $X$ localement de type fini sur $S$.
Soit $\mathcal{E}$ un faisceau localement libre de type fini de
rang $r$ sur $X$. Soit $\mathcal{L}$ un faisceau inversible
sur $X$. Alors nous avons
\begin{equation}
\label{equation-twist}
c_i({\mathcal E} \otimes {\mathcal L})
=
\sum\nolimits_{j = 0}^i
\binom{r - i + j}{j} c_{i - j}({\mathcal E}) c_1({\mathcal L})^j
\end{equation}
dans $A^*(X)$.
\end{lemma}

\begin{proof}
Cela doit être vrai pour tout triplet $(X, \mathcal{E}, \mathcal{L})$.
En particulier, cela doit être vrai lorsque $X$ est intègre et, d'après le
Lemme \ref{lemma-bivariant-zero},
il suffit de démontrer
que cela vaut après intersection avec $[X]$ pour un tel $X$. Supposons donc
que $X$ est intègre. Soit $(\pi : P \to X, \mathcal{O}_P(1))$,
resp.\ $(\pi' : P' \to X, \mathcal{O}_{P'}(1))$, le
fibré en espaces projectifs associé à $\mathcal{E}$,
resp.\ à $\mathcal{E} \otimes \mathcal{L}$. Considérons le morphisme canonique
$$
\xymatrix{
P \ar[rd]_\pi \ar[rr]_g & & P' \ar[ld]^{\pi'} \\
& X &
}
$$
voir Constructions, Lemme \ref{constructions-lemma-twisting-and-proj}.
Il possède la propriété
$g^*\mathcal{O}_{P'}(1)
= \mathcal{O}_P(1) \otimes \pi^* {\mathcal L}$.
Cela signifie que nous avons
$$
\sum\nolimits_{i = 0}^r
(-1)^i
(\xi + x)^i \cap \pi^*(c_{r - i}(\mathcal{E} \otimes \mathcal{L}) \cap [X])
=
0
$$
dans $\CH_*(P)$, où $\xi$ représente
$c_1(\mathcal{O}_P(1))$ et $x$
représente $c_1(\pi^*\mathcal{L})$. Un calcul algébrique élémentaire montre que cela
équivaut à
$$
\sum\nolimits_{i = 0}^r
(-1)^i \xi^i \left(
\sum\nolimits_{j = i}^r
(-1)^{j - i}
\binom{j}{i}
x^{j - i} \cap
\pi^*(c_{r - j}(\mathcal{E} \otimes \mathcal{L}) \cap [X])
\right)
=
0
$$
En comparant avec
l'Équation (\ref{equation-chern-classes}), il en résulte que
$$
c_{r - i}(\mathcal{E}) \cap [X] =
\sum\nolimits_{j = i}^r
\binom{j}{i}
(-c_1(\mathcal{L}))^{j - i} \cap
c_{r - j}(\mathcal{E} \otimes \mathcal{L}) \cap [X]
$$
En remaniant ceci (en éliminant les signes moins et en renumérotant), nous obtenons
la relation voulue.
\end{proof}

\noindent
Quelques cas particuliers de (\ref{equation-twist}) sont
\begin{align*}
c_1(\mathcal{E} \otimes \mathcal{L})
& =
c_1(\mathcal{E}) +
r c_1(\mathcal{L}) \\
c_2(\mathcal{E} \otimes \mathcal{L})
& =
c_2(\mathcal{E}) +
(r - 1) c_1(\mathcal{E}) c_1(\mathcal{L}) +
\binom{r}{2} c_1(\mathcal{L})^2 \\
c_3(\mathcal{E} \otimes \mathcal{L})
& =
c_3(\mathcal{E}) +
(r - 2) c_2(\mathcal{E})c_1(\mathcal{L}) +
\binom{r - 1}{2} c_1(\mathcal{E})c_1(\mathcal{L})^2 +
\binom{r}{3} c_1(\mathcal{L})^3
\end{align*}








\section{Additivité des classes de Chern}
\label{section-additivity-chern-classes}

\noindent
Tous les lemmes préliminaires découlent trivialement du
résultat final.

\begin{lemma}
\label{lemma-get-rid-of-trivial-subbundle}
Soit $(S, \delta)$ comme dans la Situation \ref{situation-setup}.
Soit $X$ localement de type fini sur $S$.
Soient $\mathcal{E}$, $\mathcal{F}$ des faisceaux localement libres de type fini
sur $X$, de rangs $r$, $r - 1$, qui s'insèrent dans une suite
exacte courte
$$
0 \to \mathcal{O}_X \to \mathcal{E} \to \mathcal{F} \to 0
$$
Alors nous avons
$$
c_r(\mathcal{E}) = 0, \quad
c_j(\mathcal{E}) = c_j(\mathcal{F}), \quad j = 0, \ldots, r - 1
$$
dans $A^*(X)$.
\end{lemma}

\begin{proof}
D'après le Lemme \ref{lemma-bivariant-zero},
il suffit de montrer que, si $X$ est intègre,
alors $c_j(\mathcal{E}) \cap [X] = c_j(\mathcal{F}) \cap [X]$.
Soit $(\pi : P \to X, \mathcal{O}_P(1))$,
resp.\ $(\pi' : P' \to X, \mathcal{O}_{P'}(1))$, le
fibré en espaces projectifs associé à $\mathcal{E}$, resp.\ à $\mathcal{F}$.
La surjection $\mathcal{E} \to \mathcal{F}$ donne lieu
à une immersion fermée
$$
i : P' \longrightarrow P
$$
sur $X$. De plus, l'élément
$1 \in \Gamma(X, \mathcal{O}_X) \subset \Gamma(X, \mathcal{E})$
donne lieu à une section globale $s \in \Gamma(P, \mathcal{O}_P(1))$
dont le lieu des zéros est exactement $P'$. Ainsi, $P'$ est un diviseur de Cartier effectif
sur $P$ tel que $\mathcal{O}_P(P') \cong \mathcal{O}_P(1)$.
Nous voyons donc que
$$
c_1(\mathcal{O}_P(1)) \cap \pi^*\alpha = i_*((\pi')^*\alpha)
$$
pour toute classe de cycle $\alpha$ sur $X$, d'après le
Lemme \ref{lemma-relative-effective-cartier}.
D'après le Lemme \ref{lemma-determine-intersections}, nous voyons que
$\alpha_j = c_j(\mathcal{F}) \cap [X]$, $j = 0, \ldots, r - 1$,
satisfont
$$
\sum\nolimits_{j = 0}^{r - 1} (-1)^jc_1(\mathcal{O}_{P'}(1))^j
\cap (\pi')^*\alpha_j = 0
$$
En prenant l'image directe dans $P$ et en utilisant la remarque ci-dessus ainsi que le
Lemme \ref{lemma-pushforward-cap-c1}, nous obtenons
$$
\sum\nolimits_{j = 0}^{r - 1}
(-1)^j c_1(\mathcal{O}_P(1))^{j + 1}
\cap \pi^*\alpha_j = 0
$$
Par l'unicité du Lemme \ref{lemma-determine-intersections},
nous concluons que
$c_r(\mathcal{E}) \cap [X] = 0$ et
$c_j(\mathcal{E}) \cap [X] = \alpha_j = c_j(\mathcal{F}) \cap [X]$
pour $j = 0, \ldots, r - 1$. Le lemme est donc vrai.
\end{proof}

\begin{lemma}
\label{lemma-additivity-invertible-subsheaf}
Soit $(S, \delta)$ comme dans la Situation \ref{situation-setup}.
Soit $X$ localement de type fini sur $S$.
Soient $\mathcal{E}$, $\mathcal{F}$ des faisceaux localement libres de type fini
sur $X$, de rangs $r$, $r - 1$, qui s'insèrent dans une suite
exacte courte
$$
0 \to \mathcal{L} \to \mathcal{E} \to \mathcal{F} \to 0
$$
où $\mathcal{L}$ est un faisceau inversible.
Alors
$$
c(\mathcal{E}) = c(\mathcal{L}) c(\mathcal{F})
$$
dans $A^*(X)$.
\end{lemma}

\begin{proof}
Cette relation dit simplement que
$c_i(\mathcal{E}) = c_i(\mathcal{F}) + c_1(\mathcal{L})c_{i - 1}(\mathcal{F})$.
D'après le Lemme \ref{lemma-get-rid-of-trivial-subbundle},
nous avons $c_j(\mathcal{E} \otimes \mathcal{L}^{\otimes -1})
= c_j(\mathcal{F} \otimes \mathcal{L}^{\otimes -1})$ pour
$j = 0, \ldots, r$ étaient nous posons
$c_r(\mathcal{F} \otimes \mathcal{L}^{-1}) = 0$ par convention.
En appliquant le Lemme \ref{lemma-chern-classes-E-tensor-L}, nous déduisons
$$
\sum_{j = 0}^i
\binom{r - i + j}{j} (-1)^j c_{i - j}({\mathcal E}) c_1({\mathcal L})^j
=
\sum_{j = 0}^i
\binom{r - 1 - i + j}{j} (-1)^j c_{i - j}({\mathcal F}) c_1({\mathcal L})^j
$$
En posant
$c_i(\mathcal{E}) = c_i(\mathcal{F}) + c_1(\mathcal{L})c_{i - 1}(\mathcal{F})$,
on obtient une « solution » de cette équation. Le lemme s'ensuit si nous montrons
que c'est la seule solution possible. Nous omettons la vérification.
\end{proof}

\begin{lemma}
\label{lemma-additivity-chern-classes}
Soit $(S, \delta)$ comme dans la Situation \ref{situation-setup}.
Soit $X$ un schéma localement de type fini sur $S$.
Supposons que ${\mathcal E}$ s'insère dans une
suite exacte
$$
0
\to
{\mathcal E}_1
\to
{\mathcal E}
\to
{\mathcal E}_2
\to
0
$$
de faisceaux localement libres de type fini $\mathcal{E}_i$ de rang $r_i$.
Les classes de Chern totales satisfont
$$
c({\mathcal E}) = c({\mathcal E}_1) c({\mathcal E}_2)
$$
dans $A^*(X)$.
\end{lemma}

\begin{proof}
D'après le Lemme \ref{lemma-bivariant-zero}, nous pouvons supposer $X$ intègre
et nous devons montrer l'identité après intersection avec $[X]$.
Raisonnons par récurrence sur $r_1$. Le cas $r_1 = 1$ est le
Lemme \ref{lemma-additivity-invertible-subsheaf}.
Supposons $r_1 > 1$. Soit $(\pi : P \to X, \mathcal{O}_P(1))$
le fibré en espaces projectifs associé à $\mathcal{E}_1$. Notons que
\begin{enumerate}
\item $\pi^* : \CH_*(X) \to \CH_*(P)$ est injective, et
\item $\pi^*\mathcal{E}_1$ s'insère dans une suite exacte courte
$0 \to \mathcal{F} \to \pi^*\mathcal{E}_1 \to \mathcal{L} \to 0$,
où $\mathcal{L}$ est inversible.
\end{enumerate}
La première assertion résulte de la formule du fibré en espaces projectifs
et la seconde de la définition d'un fibré en espaces projectifs.
(En fait, $\mathcal{L} = \mathcal{O}_P(1)$.)
Soit $Q = \pi^*\mathcal{E}/\mathcal{F}$, qui s'insère dans une
suite exacte $0 \to \mathcal{L} \to Q \to \pi^*\mathcal{E}_2 \to 0$.
Par récurrence, nous avons
\begin{eqnarray*}
c(\pi^*\mathcal{E}) \cap [P]
& = &
c(\mathcal{F}) \cap c(\pi^*\mathcal{E}/\mathcal{F}) \cap [P] \\
& = &
c(\mathcal{F}) \cap c(\mathcal{L}) \cap c(\pi^*\mathcal{E}_2) \cap [P] \\
& = &
c(\pi^*\mathcal{E}_1) \cap c(\pi^*\mathcal{E}_2) \cap [P]
\end{eqnarray*}
Puisque $[P] = \pi^*[X]$, nous
concluons par le Lemme \ref{lemma-flat-pullback-cap-cj}.
\end{proof}

\begin{lemma}
\label{lemma-chern-filter-by-linebundles}
Soit $(S, \delta)$ comme dans la Situation \ref{situation-setup}.
Soit $X$ localement de type fini sur $S$.
Soient ${\mathcal L}_i$, $i = 1, \ldots, r$, des
$\mathcal{O}_X$-modules inversibles sur $X$.
Soit $\mathcal{E}$ un module localement libre de rang
$\mathcal{O}_X$-module muni d'une filtration
$$
0 = \mathcal{E}_0 \subset \mathcal{E}_1 \subset \mathcal{E}_2
\subset \ldots \subset \mathcal{E}_r = \mathcal{E}
$$
telle que $\mathcal{E}_i/\mathcal{E}_{i - 1} \cong \mathcal{L}_i$.
Posons $c_1({\mathcal L}_i) = x_i$. Alors
$$
c(\mathcal{E})
=
\prod\nolimits_{i = 1}^r (1 + x_i)
$$
dans $A^*(X)$.
\end{lemma}

\begin{proof}
Appliquer le Lemme \ref{lemma-additivity-invertible-subsheaf} et procéder par récurrence.
\end{proof}




\section{Degrés des zéro-cycles}
\label{section-degree-zero-cycles}

\noindent
Nous commençons par définir le degré d'un zéro-cycle sur un schéma propre sur un corps.
Une méthode consiste à le définir directement comme dans le
Lemme \ref{lemma-spell-out-degree-zero-cycle}, puis à montrer
qu'il est bien défini à l'aide du
Lemme \ref{lemma-curve-principal-divisor}.
Nous le définissons plutôt comme suit.

\begin{definition}
\label{definition-degree-zero-cycle}
Soit $k$ un corps (Exemple \ref{example-field}). Soit $p : X \to \Spec(k)$
propre. Le {\it degré d'un zéro-cycle} sur $X$ est donné par l'image
directe propre
$$
p_* : \CH_0(X) \to \CH_0(\Spec(k))
$$
(Lemme \ref{lemma-proper-pushforward-rational-equivalence})
composée avec l'isomorphisme naturel $\CH_0(\Spec(k)) = \mathbf{Z}$
qui envoie $[\Spec(k)]$ sur $1$. Notation : $\deg(\alpha)$.
\end{definition}

\noindent
Explicitons davantage cette définition.

\begin{lemma}
\label{lemma-spell-out-degree-zero-cycle}
Soit $k$ un corps. Soit $X$ propre sur $k$. Soit $\alpha = \sum n_i[Z_i]$
un élément de $Z_0(X)$. Alors
$$
\deg(\alpha) = \sum n_i\deg(Z_i)
$$
où $\deg(Z_i)$ est le degré de $Z_i \to \Spec(k)$, c'est-à-dire
$\deg(Z_i) = \dim_k \Gamma(Z_i, \mathcal{O}_{Z_i})$.
\end{lemma}

\begin{proof}
C'est la définition de l'image directe propre
(Définition \ref{definition-proper-pushforward}).
\end{proof}

\noindent
Nous établissons ensuite le lien avec les degrés des fibrés vectoriels
sur les schémas propres de dimension $1$ sur des corps, tels qu'ils sont définis dans
Variétés, section \ref{varieties-section-divisors-curves}.

\begin{lemma}
\label{lemma-degree-vector-bundle}
Soit $k$ un corps. Soit $X$ un schéma propre sur $k$ de dimension $\leq 1$.
Soit $\mathcal{E}$ un $\mathcal{O}_X$-module localement libre de type fini de
rang constant. Alors
$$
\deg(\mathcal{E}) = \deg(c_1(\mathcal{E}) \cap [X]_1)
$$
où le membre de gauche est défini dans
Variétés, Définition \ref{varieties-definition-degree-invertible-sheaf}.
\end{lemma}

\begin{proof}
Soient $C_i \subset X$, $i = 1, \ldots, t$, les composantes irréductibles
de dimension $1$, munies de leur structure de schéma induite réduite, et soit $m_i$ la
multiplicité de $C_i$ dans $X$. Alors $[X]_1 = \sum m_i[C_i]$ et
$c_1(\mathcal{E}) \cap [X]_1$ est la somme des images directes des cycles
$m_i c_1(\mathcal{E}|_{C_i}) \cap [C_i]$. Puisque nous avons une décomposition analogue
du degré de $\mathcal{E}$ par
Variétés, Lemme \ref{varieties-lemma-degree-in-terms-of-components},
il suffit de démontrer le lemme lorsque $X$ est une courbe propre sur $k$.

\medskip\noindent
Supposons que $X$ soit une courbe propre sur $k$.
Par Diviseurs, Lemme \ref{divisors-lemma-filter-after-modification},
il existe une modification $f : X' \to X$ telle que $f^*\mathcal{E}$
possède une filtration dont les quotients successifs sont des
$\mathcal{O}_{X'}$-modules inversibles. Puisque $f_*[X']_1 = [X]_1$, nous déduisons
du Lemme \ref{lemma-pushforward-cap-cj} que
$$
\deg(c_1(\mathcal{E}) \cap [X]_1) = \deg(c_1(f^*\mathcal{E}) \cap [X']_1)
$$
Puisque nous avons une relation analogue pour le degré par
Variétés, Lemme \ref{varieties-lemma-degree-birational-pullback},
nous nous ramenons au cas où $\mathcal{E}$ possède une filtration dont les
quotients successifs sont des $\mathcal{O}_X$-modules inversibles.
Dans ce cas, nous pouvons utiliser l'additivité du degré
(Variétés, Lemme \ref{varieties-lemma-degree-additive})
et celle des premières classes de Chern (Lemme \ref{lemma-additivity-chern-classes})
pour nous ramener au cas traité dans le paragraphe suivant.

\medskip\noindent
Supposons que $X$ soit une courbe propre sur $k$ et que $\mathcal{E}$ soit un
$\mathcal{O}_X$-module inversible. Par
Diviseurs, Lemme
\ref{divisors-lemma-quasi-projective-Noetherian-pic-effective-Cartier},
nous voyons que $\mathcal{E}$ est isomorphe à
$\mathcal{O}_X(D) \otimes \mathcal{O}_X(D')^{\otimes -1}$
pour certains diviseurs de Cartier effectifs $D, D'$ sur $X$ (cela utilise aussi
que $X$ est projectif, voir par exemple
Variétés, Lemme \ref{varieties-lemma-dim-1-proper-projective}).
Par additivité du degré relativement au produit tensoriel de faisceaux inversibles
(Variétés, Lemme \ref{varieties-lemma-degree-tensor-product})
et additivité de $c_1$ relativement au produit tensoriel de faisceaux inversibles
(Lemme \ref{lemma-c1-cap-additive} ou \ref{lemma-chern-classes-E-tensor-L}),
nous nous ramenons au cas $\mathcal{E} = \mathcal{O}_X(D)$.
Dans ce cas, le membre de gauche donne $\deg(D)$
(Variétés, Lemme \ref{varieties-lemma-degree-effective-Cartier-divisor})
et le membre de droite donne $\deg([D]_0)$ par le
Lemme \ref{lemma-geometric-cap}.
Puisque
$$
[D]_0 = \sum\nolimits_{x \in D}
\text{longueur}_{\mathcal{O}_{X, x}}(\mathcal{O}_{D, x}) [x] =
\sum\nolimits_{x \in D}
\text{longueur}_{\mathcal{O}_{D, x}}(\mathcal{O}_{D, x}) [x]
$$
par définition, nous voyons que
$$
\deg([D]_0) = \sum\nolimits_{x \in D}
\text{longueur}_{\mathcal{O}_{D, x}}(\mathcal{O}_{D, x}) [\kappa(x) : k] =
\dim_k \Gamma(D, \mathcal{O}_D) = \deg(D)
$$
L'avant-dernière égalité résulte de
Algèbre, Lemme \ref{algebra-lemma-pushdown-module},
en utilisant que $D$ est affine.
\end{proof}

\noindent
Enfin, nous pouvons tout relier aux intersections numériques
définies dans Variétés, section \ref{varieties-section-num}.

\begin{lemma}
\label{lemma-degrees-and-numerical-intersections}
Soit $k$ un corps. Soit $X$ un schéma propre sur $k$.
Soit $Z \subset X$ un sous-schéma fermé de dimension $d$.
Soient $\mathcal{L}_1, \ldots, \mathcal{L}_d$ des
$\mathcal{O}_X$-modules inversibles. Alors
$$
(\mathcal{L}_1 \cdots \mathcal{L}_d \cdot Z) =
\deg(
c_1(\mathcal{L}_1) \cap \ldots \cap c_1(\mathcal{L}_d) \cap [Z]_d)
$$
où le membre de gauche est défini dans
Variétés, Définition \ref{varieties-definition-intersection-number}.
En particulier,
$$
\deg_\mathcal{L}(Z) = \deg(c_1(\mathcal{L})^d \cap [Z]_d)
$$
si $\mathcal{L}$ est un $\mathcal{O}_X$-module inversible ample.
\end{lemma}

\begin{proof}
Nous allons le démontrer par récurrence sur $d$. Si $d = 0$, le résultat est
vrai par Variétés, Lemme \ref{varieties-lemma-chi-tensor-finite}.
Supposons $d > 0$.

\medskip\noindent
Soient $Z_i \subset Z$, $i = 1, \ldots, t$, les composantes irréductibles
de dimension $d$, munies de leur structure de schéma induite réduite, et soit $m_i$ la
multiplicité de $Z_i$ dans $Z$. Alors $[Z]_d = \sum m_i[Z_i]$ et
$c_1(\mathcal{L}_1) \cap \ldots \cap c_1(\mathcal{L}_d) \cap [Z]_d$
est la somme des cycles
$m_i c_1(\mathcal{L}_1) \cap \ldots \cap c_1(\mathcal{L}_d) \cap [Z_i]$.
Puisque nous avons une décomposition analogue pour
$(\mathcal{L}_1 \cdots \mathcal{L}_d \cdot Z)$ par
Variétés, Lemme \ref{varieties-lemma-numerical-polynomial-leading-term},
il suffit de démontrer le lemme dans le cas où $Z = X$
est une variété propre de dimension $d$ sur $k$.

\medskip\noindent
Par le lemme de Chow, il existe un morphisme propre birationnel $f : Y \to X$
tel que $Y$ soit H-projectif sur $k$. Voir Cohomologie des schémas, Lemme
\ref{coherent-lemma-chow-Noetherian} et Remarque
\ref{coherent-remark-chow-Noetherian}. Alors
$$
(f^*\mathcal{L}_1 \cdots f^*\mathcal{L}_d \cdot Y) =
(\mathcal{L}_1 \cdots \mathcal{L}_d \cdot X)
$$
par Variétés, Lemme \ref{varieties-lemma-intersection-number-and-pullback},
et nous avons
$$
f_*(c_1(f^*\mathcal{L}_1) \cap \ldots \cap c_1(f^*\mathcal{L}_d) \cap [Y]) =
c_1(\mathcal{L}_1) \cap \ldots \cap c_1(\mathcal{L}_d) \cap [X]
$$
par le Lemme \ref{lemma-pushforward-cap-c1}. Ainsi, nous pouvons remplacer $X$ par $Y$
et supposer que $X$ est projectif sur $k$.

\medskip\noindent
Si $X$ est une variété projective propre de dimension $d$, nous pouvons
écrire $\mathcal{L}_1 = \mathcal{O}_X(D) \otimes \mathcal{O}_X(D')^{\otimes -1}$
pour certains diviseurs de Cartier effectifs $D, D' \subset X$
par Diviseurs, Lemme
\ref{divisors-lemma-quasi-projective-Noetherian-pic-effective-Cartier}.
Par additivité des deux membres de l'équation
(Variétés, Lemme \ref{varieties-lemma-intersection-number-additive} et
Lemme \ref{lemma-c1-cap-additive}),
nous nous ramenons au cas $\mathcal{L}_1 = \mathcal{O}_X(D)$ pour un
diviseur de Cartier effectif $D$.
Par Variétés, Lemme
\ref{varieties-lemma-numerical-intersection-effective-Cartier-divisor},
nous avons
$$
(\mathcal{L}_1 \cdots \mathcal{L}_d \cdot X) =
(\mathcal{L}_2 \cdots \mathcal{L}_d \cdot D)
$$
et, par le Lemme \ref{lemma-geometric-cap}, nous avons
$$
c_1(\mathcal{L}_1) \cap \ldots \cap c_1(\mathcal{L}_d) \cap [X] =
c_1(\mathcal{L}_2) \cap \ldots \cap c_1(\mathcal{L}_d) \cap [D]_{d - 1}
$$
Nous obtenons ainsi le résultat par notre hypothèse de récurrence.
\end{proof}








\section{Cycles de codimension donnée}
\label{section-cycles-codimension}

\noindent
Dans certains cas, le groupe abélien de tous les cycles possède une seconde
graduation donnée par la codimension.

\begin{lemma}
\label{lemma-locally-equidimensional}
Soit $(S, \delta)$ comme dans la Situation \ref{situation-setup}.
Soit $X$ localement de type fini sur $S$. Écrivons
$\delta = \delta_{X/S}$ comme dans la section \ref{section-setup}.
Les assertions suivantes sont équivalentes :
\begin{enumerate}
\item Il existe une décomposition $X = \coprod_{n \in \mathbf{Z}} X_n$
en sous-schémas ouverts et fermés telle que $\delta(\xi) = n$ dès que
$\xi \in X_n$ est un point générique d'une composante irréductible de $X_n$.
\item Pour tout $x \in X$, il existe un voisinage ouvert $U \subset X$
de $x$ et un entier $n$ tels que $\delta(\xi) = n$ dès que
$\xi \in U$ est un point générique d'une composante irréductible de $U$.
\item Pour tout $x \in X$, il existe un entier $n_x$ tel que
$\delta(\xi) = n_x$ pour tout point générique $\xi$ d'une composante
irréductible de $X$ contenant $x$.
\end{enumerate}
Les conditions sont satisfaites si $X$ est soit
normal, soit de Cohen--Macaulay\footnote{En fait, il suffit que
$X$ soit $(S_2)$. Comparer avec Cohomologie locale, Lemme
\ref{local-cohomology-lemma-catenary-S2-equidimensional}.}.
\end{lemma}

\begin{proof}
Il est clair que (1) $\Rightarrow$ (2) $\Rightarrow$ (3).
Réciproquement, si (3) est vérifié, alors nous posons $X_n = \{x \in X \mid n_x = n\}$
et nous obtenons une décomposition comme dans (1). En effet, $X_n$ est ouvert car,
pour $x$ donné, l'union des composantes irréductibles de $X$ passant par $x$
privée de l'union des composantes irréductibles de $X$ ne passant pas par $x$
est un voisinage ouvert de $x$. Si $X$ est normal, alors $X$ est une
réunion disjointe de schémas intègres
(Propriétés, Lemme \ref{properties-lemma-normal-locally-Noetherian})
et les propriétés sont donc vérifiées.
Si $X$ est de Cohen--Macaulay, alors
$\delta' : X \to \mathbf{Z}$, $x \mapsto -\dim(\mathcal{O}_{X, x})$
est une fonction de dimension sur $X$ (voir Exemple \ref{example-CM-irreducible}).
Puisque $\delta - \delta'$ est localement constante
(Topologie, Lemme \ref{topology-lemma-dimension-function-unique})
et puisque $\delta'(\xi) = 0$ pour tout point générique $\xi$ de $X$,
nous voyons que (2) est vérifié.
\end{proof}

\noindent
Soit $(S, \delta)$ comme dans la Situation \ref{situation-setup}. Soit $X$
localement de type fini sur $S$ et satisfaisant aux conditions équivalentes
du Lemme \ref{lemma-locally-equidimensional}. Pour un sous-schéma fermé
intègre $Z \subset X$, nous avons la codimension $\text{codim}(Z, X)$
de $Z$ dans $X$, voir Topologie, Définition \ref{topology-definition-codimension}.
Nous définissons un {\it cycle de codimension $p$} comme un cycle $\alpha = \sum n_Z[Z]$
sur $X$ tel que $n_Z \not = 0 \Rightarrow \text{codim}(Z, X) = p$.
Le groupe abélien de tous les cycles de codimension $p$ est noté $Z^p(X)$.
Soit $X = \coprod X_n$ la décomposition donnée dans la
partie (1) du Lemme \ref{lemma-locally-equidimensional}.
Rappelant que nos cycles sont définis comme des sommes localement finies, il est clair que
$$
Z^p(X) = \prod\nolimits_n Z_{n - p}(X_n)
$$
De plus, nous voyons que $\prod_p Z^p(X) = \prod_k Z_k(X)$. Nous pourrions maintenant définir
l'équivalence rationnelle des cycles de codimension $p$ sur $X$ exactement de la même
manière que précédemment et, en fait, redévelopper toute la théorie à partir de zéro
pour les cycles de codimension donnée lorsque $X$ est comme dans le
Lemme \ref{lemma-locally-equidimensional}. Nous nous contentons plutôt de
définir le {\it groupe de Chow des cycles de codimension $p$} par
$$
\CH^p(X) = \prod\nolimits_n \CH_{n - p}(X_n)
$$
Comme précédemment, nous avons $\prod_p \CH^p(X) = \prod_k \CH_k(X)$.
Si $X$ est quasi-compact, alors le produit dans la formule est fini
(et c'est donc une somme directe) et nous avons
$\bigoplus_p \CH^p(X) = \bigoplus_k \CH_k(X)$. Si $X$ est quasi-compact
et de dimension finie, alors seul un nombre fini de ces groupes
est non nul.

\medskip\noindent
De nombreuses constructions et de nombreux résultats pour les groupes de Chow démontrés ci-dessus ont
des analogues naturels pour les groupes de Chow $\CH^*(X)$. Chacun s'obtient
en décomposant les schémas pertinents en morceaux « équidimensionnels »
comme dans le Lemme \ref{lemma-locally-equidimensional},
puis en appliquant les résultats déjà démontrés aux
facteurs de la décomposition en produit donnée ci-dessus.
Énumérons-en quelques-uns.
\begin{enumerate}
\item Si $f : X \to Y$ est un morphisme plat de schémas localement de type fini
sur $S$ et si $X$ et $Y$ satisfont aux conditions équivalentes du
Lemme \ref{lemma-locally-equidimensional}, alors l'image inverse plate détermine une application
$$
f^* : \CH^p(Y) \to \CH^p(X)
$$
\item Si $f : X \to Y$ est un morphisme de schémas localement de type fini
sur $S$ et si $X$ et $Y$ satisfont aux conditions équivalentes du
Lemme \ref{lemma-locally-equidimensional}, disons que $f$ est de
{\it codimension} $r \in \mathbf{Z}$ si, pour toute paire de composantes irréductibles
$Z \subset X$, $W \subset Y$ telles que $f(Z) \subset W$, nous avons
$\dim_\delta(W) - \dim_\delta(Z) = r$.
\item Si $f : X \to Y$ est un morphisme propre de schémas localement de type fini
sur $S$, si $X$ et $Y$ satisfont aux conditions équivalentes du
Lemme \ref{lemma-locally-equidimensional} et si $f$ est de codimension $r$,
alors l'image directe propre est une application
$$
f_* : \CH^p(X) \to \CH^{p + r}(Y)
$$
\item Si $f : X \to Y$ est un morphisme de schémas localement de type fini sur $S$,
si $X$ et $Y$ satisfont aux conditions équivalentes du
Lemme \ref{lemma-locally-equidimensional}, si $f$ est de codimension $r$
et si $c \in A^q(X \to Y)$, alors $c$ induit des applications
$$
c \cap -  : \CH^p(Y) \to \CH^{p + q - r}(X)
$$
\item Si $X$ est un schéma localement de type fini sur $S$
satisfaisant aux conditions équivalentes du
Lemme \ref{lemma-locally-equidimensional} et si
$\mathcal{L}$ est un $\mathcal{O}_X$-module inversible,
alors
$$
c_1(\mathcal{L}) \cap - : \CH^p(X) \to \CH^{p + 1}(X)
$$
\item Si $X$ est un schéma localement de type fini sur $S$
satisfaisant aux conditions équivalentes du
Lemme \ref{lemma-locally-equidimensional} et si
$\mathcal{E}$ est un $\mathcal{O}_X$-module localement libre de type fini,
alors
$$
c_i(\mathcal{E}) \cap - : \CH^p(X) \to \CH^{p + i}(X)
$$
\end{enumerate}
Attention : la propriété, pour un morphisme, d'être de codimension $r$
n'est pas préservée par changement de base.

\begin{remark}
\label{remark-fundamental-class}
Soit $(S, \delta)$ comme dans la Situation \ref{situation-setup}.
Soit $X$ localement de type fini sur $S$ et satisfaisant aux
conditions équivalentes du Lemme \ref{lemma-locally-equidimensional}.
Soit $X = \coprod X_n$ la décomposition en sous-schémas ouverts et fermés
telle que toute composante irréductible de $X_n$ soit de
dimension $\delta$ égale à $n$. Dans cette situation, nous posons parfois
$$
[X] = \sum\nolimits_n [X_n]_n \in \CH^0(X)
$$
Cette classe est une sorte de « classe fondamentale » de $X$ dans la théorie de Chow.
\end{remark}




\section{Le principe de scindage}
\label{section-splitting-principle}

\noindent
Dans notre cadre, il n'est pas si facile de dire ce que le principe de scindage
dit/est exactement. En voici une formulation possible.

\begin{lemma}
\label{lemma-splitting-principle}
Soit $(S, \delta)$ comme dans la Situation \ref{situation-setup}. Soit $X$ localement
de type fini sur $S$. Les $\mathcal{E}_i$ forment une collection finie de
$\mathcal{O}_X$-modules localement libres de rang $r_i$. Il existe un morphisme
projectif plat $\pi : P \to X$ de dimension relative $d$ tel que
\begin{enumerate}
\item pour tout morphisme $f : Y \to X$, l'application
$\pi_Y^* : \CH_*(Y) \to \CH_{* + d}(Y \times_X P)$ est injective, et
\item chaque $\pi^*\mathcal{E}_i$ possède une filtration
dont les quotients successifs $\mathcal{L}_{i, 1}, \ldots, \mathcal{L}_{i, r_i}$
sont des ${\mathcal O}_P$-modules inversibles.
\end{enumerate}
De plus, lorsque (1) est vérifié, l'application de restriction $A^*(X) \to A^*(P)$
(Remarque \ref{remark-pullback-cohomology}) est injective.
\end{lemma}

\begin{proof}
Nous pouvons supposer $r_i \geq 1$ pour tout $i$. Nous démontrerons le lemme par récurrence
sur $\sum (r_i - 1)$. Si cet entier vaut $0$, alors $\mathcal{E}_i$
est inversible pour tout $i$ et nous concluons en prenant $\pi = \text{id}_X$.
Sinon, nous pouvons choisir un $i$ tel que $r_i > 1$ et considérer le
morphisme $\pi_i : P_i = \mathbf{P}(\mathcal{E}_i) \to X$.
Nous avons une suite exacte courte
$$
0 \to \mathcal{F} \to \pi_i^*\mathcal{E}_i \to \mathcal{O}_{P_i}(1) \to 0
$$
de $\mathcal{O}_{P_i}$-modules localement libres de type fini de rangs $r_i - 1$,
$r_i$ et $1$. Observons que $\pi_i^*$ est injective sur les groupes de Chow
après tout changement de base, par la formule du fibré projectif
(Lemme \ref{lemma-chow-ring-projective-bundle}).
Par l'hypothèse de récurrence appliquée aux
$\mathcal{O}_{P_i}$-modules localement libres de type fini $\mathcal{F}$ et $\pi_{i'}^*\mathcal{E}_{i'}$
pour $i' \not = i$, nous trouvons un morphisme $\pi : P \to P_i$ ayant
les propriétés énoncées dans le lemme. Alors la composée
$\pi_i \circ \pi : P \to X$ convient. Certains détails sont omis.
\end{proof}

\begin{remark}
\label{remark-the-proof-shows-more}
La démonstration du Lemme \ref{lemma-splitting-principle}
montre que le morphisme $\pi : P \to X$ possède les propriétés supplémentaires
suivantes :
\begin{enumerate}
\item $\pi$ est une composée finie de fibrés en espaces projectifs
associés à des modules localement libres de rang fini constant, et
\item pour tout $\alpha \in \CH_k(X)$, nous avons
$\alpha = \pi_*(\xi_1 \cap \ldots \cap \xi_d \cap \pi^*\alpha)$,
où $\xi_i$ est la première classe de Chern d'un
$\mathcal{O}_P$-module inversible.
\end{enumerate}
La seconde observation résulte de la première et du
Lemme \ref{lemma-cap-projective-bundle}.
Nous ajouterons ici d'autres observations selon les besoins.
\end{remark}

\noindent
Soient $(S, \delta)$, $X$ et $\mathcal{E}_i$ comme dans le
Lemme \ref{lemma-splitting-principle}.
Le {\it principe de scindage} désigne l'usage consistant à écrire symboliquement
$$
c(\mathcal{E}_i) = \prod (1 + x_{i, j})
$$
Les symboles $x_{i, 1}, \ldots, x_{i, r_i}$ sont appelés les {\it racines de Chern}
de $\mathcal{E}_i$. Autrement dit, la $p$-ième classe de Chern de $\mathcal{E}_i$
est la $p$-ième fonction symétrique élémentaire des racines de Chern.
L'utilité du principe de scindage vient de l'assertion selon laquelle,
pour démontrer une relation polynomiale entre les classes de Chern des
$\mathcal{E}_i$, il suffit de démontrer la relation correspondante entre les
racines de Chern.

\medskip\noindent
En effet, soit $\pi : P \to X$ comme dans le Lemme \ref{lemma-splitting-principle}.
Rappelons qu'il existe un morphisme canonique de $\mathbf{Z}$-algèbres
$\pi^* : A^*(X) \to A^*(P)$, voir
Remarque \ref{remark-pullback-cohomology}. L'injectivité de $\pi_Y^*$
sur les groupes de Chow pour tout $Y$ sur $X$ implique que l'application
$\pi^* : A^*(X) \to A^*(P)$ est injective (détails omis).
Nous avons
$$
\pi^*c(\mathcal{E}_i) = \prod (1 + c_1(\mathcal{L}_{i, j}))
$$
par le Lemme \ref{lemma-chern-filter-by-linebundles}. Nous pouvons donc considérer les
racines de Chern $x_{i, j}$ comme les éléments $c_1(\mathcal{L}_{i, j}) \in A^*(P)$
et l'équation affichée comme ayant lieu dans $A^*(P)$ après
application du morphisme injectif $\pi^* : A^*(X) \to A^*(P)$ au
membre de gauche de l'équation.

\medskip\noindent
Pour voir comment cela fonctionne, il est préférable de donner quelques exemples.

\begin{lemma}
\label{lemma-chern-classes-dual}
Dans la Situation \ref{situation-setup}, soit $X$ localement de type fini sur $S$.
Soit $\mathcal{E}$ un $\mathcal{O}_X$-module localement libre de type fini
de dual $\mathcal{E}^\vee$. Alors
$$
c_i(\mathcal{E}^\vee) = (-1)^i c_i(\mathcal{E})
$$
dans $A^i(X)$.
\end{lemma}

\begin{proof}
Choisissons un morphisme $\pi : P \to X$ comme dans le
Lemme \ref{lemma-splitting-principle}.
Par l'injectivité de $\pi^*$ (après tout changement de base),
il suffit de démontrer la relation entre
les classes de Chern de $\mathcal{E}$ et de $\mathcal{E}^\vee$
après image inverse sur $P$. Nous pouvons donc supposer qu'il
existe des $\mathcal{O}_X$-modules inversibles
${\mathcal L}_i$, $i = 1, \ldots, r$,
et une filtration
$$
0 = \mathcal{E}_0 \subset \mathcal{E}_1 \subset \mathcal{E}_2
\subset \ldots \subset \mathcal{E}_r = \mathcal{E}
$$
telle que $\mathcal{E}_i/\mathcal{E}_{i - 1} \cong \mathcal{L}_i$.
Nous obtenons alors la filtration duale
$$
0 = \mathcal{E}_r^\perp \subset \mathcal{E}_1^\perp \subset \mathcal{E}_2^\perp
\subset \ldots \subset \mathcal{E}_0^\perp = \mathcal{E}^\vee
$$
telle que $\mathcal{E}_{i - 1}^\perp/\mathcal{E}_i^\perp \cong
\mathcal{L}_i^{\otimes -1}$.
Posons $x_i = c_1(\mathcal{L}_i)$.
Alors $c_1(\mathcal{L}_i^{\otimes -1}) = - x_i$
par le Lemme \ref{lemma-c1-cap-additive}.
Par le Lemme \ref{lemma-chern-filter-by-linebundles},
nous avons
$$
c(\mathcal{E}) = \prod\nolimits_{i = 1}^r (1 + x_i)
\quad\text{et}\quad
c(\mathcal{E}^\vee) = \prod\nolimits_{i = 1}^r (1 - x_i)
$$
dans $A^*(X)$. Le résultat résulte d'un calcul formel
que nous omettons.
\end{proof}

\begin{lemma}
\label{lemma-chern-classes-tensor-product}
Dans la Situation \ref{situation-setup}, soit $X$ localement de type fini sur $S$.
Soient $\mathcal{E}$ et $\mathcal{F}$ deux
$\mathcal{O}_X$-modules localement libres de type fini, de rangs $r$ et $s$. Alors nous avons
$$
c_1(\mathcal{E} \otimes \mathcal{F})
=
r c_1(\mathcal{F}) + s c_1(\mathcal{E})
$$
$$
c_2(\mathcal{E} \otimes \mathcal{F})
=
r c_2(\mathcal{F}) + s c_2(\mathcal{E}) +
{r \choose 2} c_1(\mathcal{F})^2 +
(rs - 1) c_1(\mathcal{F})c_1(\mathcal{E}) +
{s \choose 2} c_1(\mathcal{E})^2
$$
et ainsi de suite dans $A^*(X)$.
\end{lemma}

\begin{proof}
En raisonnant exactement comme dans la démonstration du Lemme \ref{lemma-chern-classes-dual},
nous pouvons supposer que nous avons
des $\mathcal{O}_X$-modules inversibles
${\mathcal L}_i$, $i = 1, \ldots, r$,
${\mathcal N}_i$, $i = 1, \ldots, s$,
des filtrations
$$
0 = \mathcal{E}_0 \subset \mathcal{E}_1 \subset \mathcal{E}_2
\subset \ldots \subset \mathcal{E}_r = \mathcal{E}
\quad\text{et}\quad
0 = \mathcal{F}_0 \subset \mathcal{F}_1 \subset \mathcal{F}_2
\subset \ldots \subset \mathcal{F}_s = \mathcal{F}
$$
telles que $\mathcal{E}_i/\mathcal{E}_{i - 1} \cong \mathcal{L}_i$
et telles que $\mathcal{F}_j/\mathcal{F}_{j - 1} \cong \mathcal{N}_j$.
En ordonnant les paires $(i, j)$ lexicographiquement,
nous obtenons une filtration
$$
0 \subset \ldots \subset
\mathcal{E}_i \otimes \mathcal{F}_j
+
\mathcal{E}_{i - 1} \otimes \mathcal{F}
\subset \ldots \subset \mathcal{E} \otimes \mathcal{F}
$$
de quotients successifs
$$
\mathcal{L}_1 \otimes \mathcal{N}_1,
\mathcal{L}_1 \otimes \mathcal{N}_2,
\ldots,
\mathcal{L}_1 \otimes \mathcal{N}_s,
\mathcal{L}_2 \otimes \mathcal{N}_1,
\ldots,
\mathcal{L}_r \otimes \mathcal{N}_s
$$
Par le Lemme \ref{lemma-chern-filter-by-linebundles},
nous avons
$$
c(\mathcal{E}) = \prod (1 + x_i),
\quad
c(\mathcal{F}) = \prod (1 + y_j),
\quad\text{et}\quad
c(\mathcal{E} \otimes \mathcal{F}) = \prod (1 + x_i + y_j),
$$
dans $A^*(X)$. Le résultat résulte d'un calcul formel
que nous omettons.
\end{proof}

\begin{remark}
\label{remark-equalities-nonconstant-rank}
Les égalités démontrées ci-dessus restent vraies même lorsque nous travaillons avec
des modules localement libres de type fini
sur $\mathcal{O}_X$ dont le rang peut ne pas être constant.
En fait, nous pouvons travailler avec des polynômes en le rang et les
classes de Chern comme suit. Considérons l'anneau gradué de polynômes
$\mathbf{Z}[r, c_1, c_2, c_3, \ldots]$,
où $r$ est de degré $0$ et $c_i$ de degré $i$. Soit
$$
P \in \mathbf{Z}[r, c_1, c_2, c_3, \ldots]
$$
un polynôme homogène de degré $p$. Alors, pour tout
$\mathcal{O}_X$-module localement libre de type fini $\mathcal{E}$ sur $X$, nous pouvons considérer
$$
P(\mathcal{E}) =
P(r(\mathcal{E}), c_1(\mathcal{E}), c_2(\mathcal{E}), c_3(\mathcal{E}), \ldots)
\in A^p(X)
$$
voir la Remarque \ref{remark-extend-to-finite-locally-free} pour les notations et
conventions. Pour démontrer des relations entre ces polynômes (pour plusieurs
modules localement libres de type fini), nous pouvons travailler localement sur $X$ et utiliser le principe de
scindage comme ci-dessus. Par exemple, nous affirmons que
$$
c_2(\SheafHom_{\mathcal{O}_X}(\mathcal{E}, \mathcal{E})) =
P(\mathcal{E})
$$
où $P = 2rc_2 - (r - 1)c_1^2$.
En effet, puisque $\SheafHom_{\mathcal{O}_X}(\mathcal{E}, \mathcal{E}) =
\mathcal{E} \otimes \mathcal{E}^\vee$, cela résulte facilement des
Lemmes \ref{lemma-chern-classes-dual} et
\ref{lemma-chern-classes-tensor-product}
ci-dessus, en décomposant $X$ en morceaux où le rang
de $\mathcal{E}$ est constant comme dans la
Remarque \ref{remark-extend-to-finite-locally-free}.
\end{remark}

\begin{example}
\label{example-power-sum}
Pour tout $p \geq 1$, il existe un unique polynôme homogène
$P_p \in \mathbf{Z}[c_1, c_2, c_3, \ldots]$ de degré $p$
tel que, pour tout $n \geq p$, nous ayons
$$
P_p(s_1, s_2, \ldots, s_p) = \sum x_i^p
$$
dans $\mathbf{Z}[x_1, \ldots, x_n]$, où $s_1, \ldots, s_p$ sont les
polynômes symétriques élémentaires en $x_1, \ldots, x_n$, donc
$$
s_i = \sum\nolimits_{1 \leq j_1 < \ldots < j_i \leq n}
x_{j_1}x_{j_2} \ldots x_{j_i}
$$
L'existence de $P_p$ résulte du fait bien connu que
les fonctions symétriques élémentaires engendrent l'anneau de
toutes les fonctions symétriques sur les entiers. Une autre manière de
caractériser $P_p \in \mathbf{Z}[c_1, c_2, c_3, \ldots]$ est que nous avons
$$
\log(1 + c_1 + c_2 + c_3 + \ldots) =
\sum\nolimits_{p \geq 1} (-1)^{p - 1}\frac{P_p}{p}
$$
comme série formelle. Cela est clair en écrivant
$1 + c_1 + c_2 + \ldots = \prod (1 + x_i)$ et en appliquant
la série de la fonction logarithme. En développant le membre de
gauche, nous obtenons
\begin{align*}
& (c_1 + c_2 + \ldots) - (1/2)(c_1 + c_2 + \ldots)^2 +
(1/3)(c_1 + c_2 + \ldots)^3 - \ldots \\
& =
c_1 + (c_2 - (1/2)c_1^2) + (c_3 - c_1c_2 + (1/3)c_1^3) + \ldots
\end{align*}
De cette manière, nous trouvons
\begin{align*}
P_1 & = c_1, \\
P_2 & = c_1^2 - 2c_2, \\
P_3 & = c_1^3 - 3c_1c_2 + 3c_3, \\
P_4 & = c_1^4 - 4c_1^2c_2 + 4c_1c_3 + 2c_2^2 - 4c_4,
\end{align*}
et ainsi de suite. Puisque les classes de Chern d'un
$\mathcal{O}_X$-module localement libre de type fini $\mathcal{E}$ sont les polynômes
symétriques élémentaires en les racines de Chern $x_i$, nous voyons que
$$
P_p(\mathcal{E}) = \sum x_i^p
$$
Pour plus de commodité, nous posons $P_0 = r$ dans $\mathbf{Z}[r, c_1, c_2, c_3, \ldots]$
de sorte que $P_0(\mathcal{E}) = r(\mathcal{E})$ comme classe bivariante
(comme dans les Remarques \ref{remark-extend-to-finite-locally-free} et
\ref{remark-equalities-nonconstant-rank}).
\end{example}






\section{Classes de Chern et sections}
\label{section-top-chern-class}

\noindent
Une brève section dont le résultat principal est que nous pouvons calculer la
classe de Chern de degré maximal d'un module localement libre de type fini au moyen du
lieu des zéros d'une « section régulière ».

\medskip\noindent
Soit $(S, \delta)$ comme dans la Situation \ref{situation-setup}. Soit $X$ un schéma
localement de type fini sur $S$. Soit $\mathcal{E}$ un
$\mathcal{O}_X$-module localement libre de type fini. Soit $f : X' \to X$ localement de type fini. Soit
$$
s \in \Gamma(X', f^*\mathcal{E})
$$
une section globale de l'image inverse de $\mathcal{E}$ sur $X'$. Soit
$Z(s) \subset X'$ le schéma des zéros de $s$. Plus précisément, nous définissons
$Z(s)$ comme le sous-schéma fermé dont le faisceau quasi-cohérent
d'idéaux est l'image de l'application
$s : f^*\mathcal{E}^\vee \to \mathcal{O}_{X'}$.

\begin{lemma}
\label{lemma-top-chern-class}
Dans la situation décrite juste ci-dessus, supposons $\dim_\delta(X') = n$,
que $f^*\mathcal{E}$ soit de rang constant $r$, que
$\dim_\delta(Z(s)) \leq n - r$, et que, pour tout point générique
$\xi \in Z(s)$ avec $\delta(\xi) = n - r$, l'idéal de $Z(s)$
dans $\mathcal{O}_{X', \xi}$ soit engendré par une suite régulière
de longueur $r$. Alors
$$
c_r(\mathcal{E}) \cap [X']_n = [Z(s)]_{n - r}
$$
dans $\CH_*(X')$.
\end{lemma}

\begin{proof}
Puisque $c_r(\mathcal{E})$ est une classe bivariante
(Lemme \ref{lemma-cap-cp-bivariant}),
nous pouvons supposer $X = X'$ et nous devons montrer que
$c_r(\mathcal{E}) \cap [X]_n = [Z(s)]_{n - r}$ dans $\CH_{n - r}(X)$.
Nous démontrerons le lemme par récurrence sur $r \geq 0$. (Le cas
$r = 0$ est trivial.) Le cas $r = 1$
est traité par le Lemme \ref{lemma-geometric-cap}. Supposons $r > 1$.

\medskip\noindent
Soit $\pi : P \to X$ le fibré en espaces projectifs associé à
$\mathcal{E}$ et considérons la suite exacte courte
$$
0 \to \mathcal{E}' \to \pi^*\mathcal{E} \to \mathcal{O}_P(1) \to 0
$$
Par la formule du fibré en espaces projectifs
(Lemme \ref{lemma-chow-ring-projective-bundle}),
il suffit de démontrer l'égalité après image inverse par $\pi$.
Observons que $\pi^{-1}Z(s) = Z(\pi^*s)$ est de dimension $\delta$
$\leq n - 1$ et que l'hypothèse sur les suites régulières aux
points génériques de dimension $\delta$ égale à $n - 1$ est vérifiée par
image inverse plate, voir
Algèbre, Lemme \ref{algebra-lemma-flat-increases-depth}.
Soit $t \in \Gamma(P, \mathcal{O}_P(1))$ l'image de $\pi^*s$.
Nous affirmons que
$$
[Z(t)]_{n + r - 2} = c_1(\mathcal{O}_P(1)) \cap [P]_{n + r - 1}
$$
En admettant cette assertion, nous terminons la démonstration comme suit.
La restriction $\pi^*s|_{Z(t)}$ s'envoie sur zéro dans
$\mathcal{O}_P(1)|_{Z(t)}$ et provient donc d'un unique
élément $s' \in \Gamma(Z(t), \mathcal{E}'|_{Z(t)})$.
Notons que $Z(s') = Z(\pi^*s)$ comme sous-schémas fermés de $P$.
Si $\xi \in Z(s')$ est un point générique avec $\delta(\xi) = n - 1$,
alors l'idéal de $Z(s')$ dans $\mathcal{O}_{Z(t), \xi}$
peut être engendré par une suite régulière de longueur $r - 1$ : il est engendré par
$r - 1$ éléments qui sont les images de $r - 1$ éléments de
$\mathcal{O}_{P, \xi}$ qui, avec un générateur de
l'idéal de $Z(t)$ dans $\mathcal{O}_{P, \xi}$, forment une suite régulière
de longueur $r$ dans $\mathcal{O}_{P, \xi}$. Nous pouvons donc appliquer
l'hypothèse de récurrence à $s'$ sur $Z(t)$ pour obtenir
$c_{r - 1}(\mathcal{E}') \cap [Z(t)]_{n + r - 2} = [Z(s')]_{n - 1}$.
En combinant tout ce qui précède, nous obtenons
\begin{align*}
c_r(\pi^*\mathcal{E}) \cap [P]_{n + r - 1}
& =
c_{r - 1}(\mathcal{E}') \cap c_1(\mathcal{O}_P(1)) \cap [P]_{n + r - 1} \\
& =
c_{r - 1}(\mathcal{E}') \cap [Z(t)]_{n + r - 2} \\
& =
[Z(s')]_{n - 1} \\
& = [Z(\pi^*s)]_{n - 1}
\end{align*}
ce qui est précisément ce qu'il fallait montrer.

\medskip\noindent
Démonstration de l'assertion. Elle résultera d'une application du
Lemme \ref{lemma-geometric-cap}, déjà utilisé.
Nous avons $\pi^{-1}(Z(s)) = Z(\pi^*s) \subset Z(t)$.
D'autre part, pour $x \in X$, si $P_x \subset Z(t)$, alors
$t|_{P_x} = 0$, ce qui implique que $s$ est nul dans la fibre
$\mathcal{E} \otimes \kappa(x)$, ce qui implique $x \in Z(s)$.
Il s'ensuit que $\dim_\delta(Z(t)) \leq n + (r - 1) - 1$.
Enfin, soit $\xi \in Z(t)$ un point générique tel que
$\delta(\xi) = n + r - 2$. Si $\xi$ n'est pas le point générique
de la fibre de $P \to X$, il est immédiat que
toute équation locale de $Z(t)$ est un non-diviseur de zéro dans $\mathcal{O}_{P, \xi}$
(car nous pouvons le vérifier sur la fibre par
Algèbre, Lemme \ref{algebra-lemma-grothendieck}).
Si $\xi$ est le point générique d'une fibre, alors $x = \pi(\xi) \in Z(s)$
et $\delta(x) = n + r - 2 - (r - 1) = n - 1$. Cela contredit
$\dim_\delta(Z(s)) \leq n - r$ puisque $r > 1$,
de sorte que ce cas ne se produit pas.
\end{proof}

\begin{lemma}
\label{lemma-easy-virtual-class}
Soit $(S, \delta)$ comme dans la Situation \ref{situation-setup}. Soit $X$
un schéma localement de type fini sur $S$. Soit
$$
0 \to \mathcal{N}' \to \mathcal{N} \to \mathcal{E} \to 0
$$
une suite exacte courte de $\mathcal{O}_X$-modules localement libres de type fini.
Considérons l'immersion fermée
$$
i :
N' = \underline{\Spec}_X(\text{Sym}((\mathcal{N}')^\vee))
\longrightarrow
N = \underline{\Spec}_X(\text{Sym}(\mathcal{N}^\vee))
$$
Pour $\alpha \in \CH_k(X)$, nous avons
$$
i_*(p')^*\alpha = p^*(c_{top}(\mathcal{E}) \cap \alpha)
$$
où $p' : N' \to X$ et $p : N \to X$ sont les morphismes structuraux.
\end{lemma}

\begin{proof}
Ici, $c_{top}(\mathcal{E})$ est la classe bivariante définie dans la
Remarque \ref{remark-top-chern-class}. Par sa définition même, pour
vérifier la formule, nous pouvons supposer que $\mathcal{E}$
est de rang constant. Nous pouvons de même supposer que $\mathcal{N}'$ et
$\mathcal{N}$ sont de rangs constants, disons $r'$ et $r$, de sorte que
$\mathcal{E}$ est de rang $r - r'$ et
$c_{top}(\mathcal{E}) = c_{r - r'}(\mathcal{E})$.
Observons que $p^*\mathcal{E}$ possède une section canonique
$$
s \in \Gamma(N, p^*\mathcal{E}) = \Gamma(X, p_*p^*\mathcal{E}) =
\Gamma(X, \mathcal{E} \otimes_{\mathcal{O}_X} \text{Sym}(\mathcal{N}^\vee)
\supset \Gamma(X, \SheafHom(\mathcal{N}, \mathcal{E}))
$$
correspondant à la surjection $\mathcal{N} \to \mathcal{E}$ donnée
dans l'énoncé du lemme. Le schéma d'annulation de cette section
est exactement $N' \subset N$. Soit $Y \subset X$ un sous-schéma fermé
intègre de dimension $\delta$ égale à $n$. Alors nous avons
\begin{enumerate}
\item $p^*[Y] = [p^{-1}(Y)]$ puisque $p^{-1}(Y)$ est intègre de
dimension $\delta$ égale à $n + r$,
\item $(p')^*[Y] = [(p')^{-1}(Y)]$ puisque $(p')^{-1}(Y)$ est intègre de
dimension $\delta$ égale à $n + r'$,
\item la restriction de $s$ à $p^{-1}Y$ a pour schéma d'annulation
$(p')^{-1}Y$ et l'immersion fermée $(p')^{-1}Y \to p^{-1}Y$
est une immersion régulière (localement découpée par une suite régulière).
\end{enumerate}
Nous concluons que
$$
(p')^*[Y] = c_{r - r'}(p^*\mathcal{E}) \cap p^*[Y]
\quad\text{dans}\quad \CH_*(N)
$$
par le Lemme \ref{lemma-top-chern-class}. Cela démontre le lemme.
\end{proof}










\section{Le caractère de Chern et les produits tensoriels}
\label{section-chern-classes-tensor}

\noindent
Soit $(S, \delta)$ comme dans la Situation \ref{situation-setup}.
Soit $X$ localement de type fini sur $S$.
Nous définissons le {\it caractère de Chern} d'un
$\mathcal{O}_X$-module localement libre de type fini comme l'expression formelle
$$
ch({\mathcal E}) = \sum\nolimits_{i=1}^r e^{x_i}
$$
si les $x_i$ sont les racines de Chern de ${\mathcal E}$. En écrivant cette
expression comme un polynôme en les classes de Chern, nous obtenons
\begin{align*}
ch(\mathcal{E})
& =
r(\mathcal{E})
+
c_1(\mathcal{E}) +
\frac{1}{2}(c_1(\mathcal{E})^2 - 2c_2(\mathcal{E}))
+
\frac{1}{6}(c_1(\mathcal{E})^3 - 3c_1(\mathcal{E})c_2(\mathcal{E}) + 3c_3(\mathcal{E})) \\
& \quad\quad +
\frac{1}{24}(c_1(\mathcal{E})^4 - 4c_1(\mathcal{E})^2c_2(\mathcal{E}) + 4c_1(\mathcal{E})c_3(\mathcal{E}) + 2c_2(\mathcal{E})^2 - 4c_4(\mathcal{E}))
+
\ldots \\
& =
\sum\nolimits_{p = 0, 1, 2, \ldots} \frac{P_p(\mathcal{E})}{p!}
\end{align*}
où $P_p$ sont les polynômes en les classes de Chern de
l'Exemple \ref{example-power-sum}. La composante de degré $p$ de
ce qui précède est
$$
ch_p(\mathcal{E}) = \frac{P_p(\mathcal{E})}{p!} \in A^p(X) \otimes \mathbf{Q}
$$
Que signifie le fait que les coefficients soient des nombres rationnels ?
Eh bien, cela signifie simplement que nous considérons
$ch_p(\mathcal{E})$ comme un élément de $A^p(X) \otimes \mathbf{Q}$.

\begin{remark}
\label{remark-extend-chern-character-to-finite-locally-free}
Dans la discussion ci-dessus, nous avons défini les composantes du caractère de Chern
$ch_p(\mathcal{E}) \in A^p(X) \otimes \mathbf{Q}$
de $\mathcal{E}$ même si le rang de $\mathcal{E}$
n'est pas constant. Voir les Remarques \ref{remark-extend-to-finite-locally-free} et
\ref{remark-equalities-nonconstant-rank}. Ainsi, le caractère de Chern complet
de $\mathcal{E}$ est
un élément de $\prod_{p \geq 0} (A^p(X) \otimes \mathbf{Q})$. Si $X$
est quasi-compact et $\dim(X) < \infty$ (dimension ordinaire), alors on peut montrer,
en utilisant le Lemme \ref{lemma-vanish-above-dimension} et le principe de scindage,
que $ch(\mathcal{E}) \in A^*(X) \otimes \mathbf{Q}$.
\end{remark}

\begin{lemma}
\label{lemma-chern-character-additive}
Soit $(S, \delta)$ comme dans la Situation \ref{situation-setup}. Soit $X$ localement
de type fini sur $S$. Soit
$
0 \to \mathcal{E}_1 \to \mathcal{E} \to \mathcal{E}_2 \to 0
$
une suite exacte courte de $\mathcal{O}_X$-modules localement libres de type fini.
Alors nous avons l'égalité
$$
ch(\mathcal{E}) = ch(\mathcal{E}_1) + ch(\mathcal{E}_2)
$$
Plus précisément, nous avons
$P_p(\mathcal{E}) = P_p(\mathcal{E}_1) + P_p(\mathcal{E}_2)$
dans $A^p(X)$, où $P_p$ est comme dans l'Exemple \ref{example-power-sum}.
\end{lemma}

\begin{proof}
Il suffit de démontrer l'énoncé plus précis. Par la
section \ref{section-splitting-principle},
cela résulte du fait que si $x_{1, i}$, $i = 1, \ldots, r_1$,
et $x_{2, i}$, $i = 1, \ldots, r_2$, sont les
racines de Chern de $\mathcal{E}_1$ et $\mathcal{E}_2$, alors
$x_{1, 1}, \ldots, x_{1, r_1}, x_{2, 1}, \ldots, x_{2, r_2}$
sont les racines de Chern de $\mathcal{E}$. Le résultat découle donc
de notre choix de $P_p$ dans l'Exemple \ref{example-power-sum}.
\end{proof}

\begin{lemma}
\label{lemma-chern-character-multiplicative}
Soit $(S, \delta)$ comme dans la Situation \ref{situation-setup}. Soit $X$ localement
de type fini sur $S$. Soient $\mathcal{E}_1$ et $\mathcal{E}_2$
des $\mathcal{O}_X$-modules localement libres de type fini.
Alors nous avons l'égalité
$$
ch(\mathcal{E}_1 \otimes_{\mathcal{O}_X} \mathcal{E}_2) =
ch(\mathcal{E}_1) ch(\mathcal{E}_2)
$$
Plus précisément, nous avons
$$
P_p(\mathcal{E}_1 \otimes_{\mathcal{O}_X} \mathcal{E}_2) =
\sum\nolimits_{p_1 + p_2 = p}
{p \choose p_1} P_{p_1}(\mathcal{E}_1) P_{p_2}(\mathcal{E}_2)
$$
dans $A^p(X)$, où $P_p$ est comme dans l'Exemple \ref{example-power-sum}.
\end{lemma}

\begin{proof}
Il suffit de démontrer l'énoncé plus précis. Par la
section \ref{section-splitting-principle},
cela résulte du fait que si $x_{1, i}$, $i = 1, \ldots, r_1$,
et $x_{2, i}$, $i = 1, \ldots, r_2$, sont les
racines de Chern de $\mathcal{E}_1$ et $\mathcal{E}_2$, alors
$x_{1, i} + x_{2, j}$, $1 \leq i \leq r_1$, $1 \leq j \leq r_2$,
sont les racines de Chern de $\mathcal{E}_1 \otimes \mathcal{E}_2$.
Le résultat découle donc de la formule du binôme pour
$(x_{1, i} + x_{2, j})^p$ et de la
forme de nos polynômes $P_p$ dans l'Exemple \ref{example-power-sum}.
\end{proof}

\begin{lemma}
\label{lemma-chern-character-dual}
Dans la Situation \ref{situation-setup}, soit $X$ localement de type fini sur $S$.
Soit $\mathcal{E}$ un $\mathcal{O}_X$-module localement libre de type fini
de dual $\mathcal{E}^\vee$. Alors
$ch_i(\mathcal{E}^\vee) = (-1)^i ch_i(\mathcal{E})$ dans
$A^i(X) \otimes \mathbf{Q}$.
\end{lemma}

\begin{proof}
Cela résulte du résultat correspondant pour les classes de Chern
(Lemme \ref{lemma-chern-classes-dual}).
\end{proof}











\section{Classes de Chern et catégorie dérivée}
\label{section-pre-derived}

\noindent
Dans cette section, nous définissons la classe de Chern totale d'un objet parfait $E$
de la catégorie dérivée d'un schéma $X$, sous l'hypothèse que $E$
peut être représenté par un complexe fini de modules localement libres de type fini
sur une enveloppe de $X$.

\medskip\noindent
Soit $(S, \delta)$ comme dans la Situation \ref{situation-setup}.
Soit $X$ localement de type fini sur $S$. Soit
$$
\mathcal{E}^a \to \mathcal{E}^{a + 1} \to \ldots \to \mathcal{E}^b
$$
un complexe borné de $\mathcal{O}_X$-modules localement libres de type fini
de rang constant.
Nous définissons alors la {\it classe de Chern totale du complexe} par la formule
$$
c(\mathcal{E}^\bullet) = \prod\nolimits_{n = a, \ldots, b}
c(\mathcal{E}^n)^{(-1)^n} \in \prod\nolimits_{p \geq 0} A^p(X)
$$
Ici, l'inverse est l'inverse formel, donc
$$
(1 + c_1 + c_2 + c_3 + \ldots)^{-1} =
1 - c_1 + c_1^2 - c_2 - c_1^3 + 2c_1 c_2 - c_3 + \ldots
$$
Nous noterons $c_p(\mathcal{E}^\bullet) \in A^p(X)$
la partie de degré $p$ de $c(\mathcal{E}^\bullet)$.
Nous définissons de même le {\it caractère de Chern du complexe} par
la formule
$$
ch(\mathcal{E}^\bullet) = \sum\nolimits_{n = a, \ldots, b}
(-1)^n ch(\mathcal{E}^n) \in
\prod\nolimits_{p \geq 0} (A^p(X) \otimes \mathbf{Q})
$$
Nous noterons $ch_p(\mathcal{E}^\bullet) \in A^p(X) \otimes \mathbf{Q}$
la partie de degré $p$ de $ch(\mathcal{E}^\bullet)$.
Enfin, pour $P_p \in \mathbf{Z}[r, c_1, c_2, c_3, \ldots]$
comme dans l'Exemple \ref{example-power-sum}, nous définissons
$$
P_p(\mathcal{E}^\bullet) = \sum\nolimits_{n = a, \ldots, b}
(-1)^n P_p(\mathcal{E}^n)
$$
dans $A^p(X)$. Alors nous avons
$ch_p(\mathcal{E}^\bullet) = (1/p!)P_p(\mathcal{E}^\bullet)$
comme d'habitude. Le lemme suivant montre que ces constructions ne dépend
que de l'image du complexe dans la catégorie dérivée.

\begin{lemma}
\label{lemma-pre-derived-chern-class}
Soit $(S, \delta)$ comme dans la Situation \ref{situation-setup}.
Soit $X$ localement de type fini sur $S$. Soit $E \in D(\mathcal{O}_X)$
un objet tel qu'il existe un complexe localement borné
$\mathcal{E}^\bullet$ de $\mathcal{O}_X$-modules localement libres de type fini
représentant $E$. Alors une légère généralisation des constructions ci-dessus
$$
c(\mathcal{E}^\bullet) \in \prod\nolimits_{p \geq 0} A^p(X),\quad
ch(\mathcal{E}^\bullet) \in
\prod\nolimits_{p \geq 0} A^p(X) \otimes \mathbf{Q},\quad
P_p(\mathcal{E}^\bullet) \in A^p(X)
$$
sont indépendantes du choix du complexe $\mathcal{E}^\bullet$.
\end{lemma}

\begin{proof}
Nous le démontrons pour la classe de Chern totale ; les deux autres cas résultent
des mêmes arguments en utilisant le
Lemme \ref{lemma-chern-character-additive}
à la place du
Lemme \ref{lemma-additivity-chern-classes}.

\medskip\noindent
Comme dans la Remarque \ref{remark-extend-to-finite-locally-free}, afin
de définir la classe de Chern totale $c(\mathcal{E}^\bullet)$,
nous décomposons $X$ en sous-schémas ouverts et fermés
$$
X = \coprod\nolimits_{i \in I} X_i
$$
tels que le rang de $\mathcal{E}^n$ soit constant sur $X_i$ pour
tous $n$ et $i$. (Comme ces rangs sont des fonctions localement constantes
sur $X$, nous pouvons le faire.) Puisque $\mathcal{E}^\bullet$ est localement
borné, nous voyons que seul un nombre fini des faisceaux
$\mathcal{E}^n|_{X_i}$ sont non nuls pour $i$ fixé. Nous pouvons donc
définir
$$
c(\mathcal{E}^\bullet|_{X_i}) =
\prod\nolimits_n c(\mathcal{E}^n|_{X_i})^{(-1)^n}
\in \prod\nolimits_{p \geq 0} A^p(X_i)
$$
comme ci-dessus. Par le Lemme \ref{lemma-disjoint-decomposition-bivariant},
nous avons $A^p(X) = \prod_i A^p(X_i)$. Ainsi, pour tout $p \in \mathbf{Z}$,
nous avons un unique élément $c_p(\mathcal{E}^\bullet) \in A^p(X)$ dont la restriction
est $c_p(\mathcal{E}^\bullet|_{X_i})$ sur $X_i$ pour tout $i$.

\medskip\noindent
Supposons que nous ayons un second complexe localement borné
$\mathcal{F}^\bullet$ de $\mathcal{O}_X$-modules localement libres de type fini
représentant $E$.
Soit $g : Y \to X$ un morphisme localement de type fini avec $Y$ intègre.
Par le Lemme \ref{lemma-bivariant-zero}, il suffit de montrer que
avec $c(g^*\mathcal{E}^\bullet) \cap [Y]$ est le même que
$c(g^*\mathcal{F}^\bullet) \cap [Y]$, et il suffit même de le démontrer
après avoir remplacé $Y$ par un schéma intègre propre et birationnel
sur $Y$. Nous concluons alors d'abord que $g^*\mathcal{E}^\bullet$
et $g^*\mathcal{F}^\bullet$ sont des complexes bornés de
$\mathcal{O}_Y$-modules localement libres de type fini de rang constant. Ensuite, par
Plus sur la platitude, Lemme \ref{flat-lemma-blowup-complex-integral},
nous pouvons supposer que $H^i(Lg^*E)$ est parfait de dimension de Tor $\leq 1$
pour tout $i \in \mathbf{Z}$.
Cela nous ramène au cas traité dans le paragraphe suivant.

\medskip\noindent
Supposons que $X$ soit intègre, que $\mathcal{E}^\bullet$ et $\mathcal{F}^\bullet$
soient des complexes bornés de modules localement libres de type fini de rang constant, et que
$H^i(E)$ soit un $\mathcal{O}_X$-module parfait de dimension de Tor $\leq 1$
pour tout $i \in \mathbf{Z}$. Nous devons
montrer que $c(\mathcal{E}^\bullet) \cap [X]$ est le même que
$c(\mathcal{F}^\bullet) \cap [X]$. Notons
$d_\mathcal{E}^i : \mathcal{E}^i \to \mathcal{E}^{i + 1}$ et
$d_\mathcal{F}^i : \mathcal{F}^i \to \mathcal{F}^{i + 1}$
les différentielles de nos complexes. Par
Plus sur la platitude, Remarque \ref{flat-remark-when-you-have-a-complex},
nous savons que $\Im(d_\mathcal{E}^i)$, $\Ker(d_\mathcal{E}^i)$,
$\Im(d_\mathcal{F}^i)$ et $\Ker(d_\mathcal{F}^i)$
sont des $\mathcal{O}_X$-modules localement libres de type fini pour tout $i$.
Par additivité (Lemme \ref{lemma-additivity-chern-classes}), nous voyons que
$$
c(\mathcal{E}^\bullet) = \prod\nolimits_i
c(\Ker(d_\mathcal{E}^i))^{(-1)^i} c(\Im(d_\mathcal{E}^i))^{(-1)^i}
$$
et de même pour $\mathcal{F}^\bullet$. Puisque nous avons les
suites exactes courtes
$$
0 \to \Im(d_\mathcal{E}^i) \to \Ker(d_\mathcal{E}^i) \to H^i(E) \to 0
\quad\text{et}\quad
0 \to \Im(d_\mathcal{F}^i) \to \Ker(d_\mathcal{F}^i) \to H^i(E) \to 0
$$
nous nous ramenons au problème énoncé et résolu dans le paragraphe suivant.

\medskip\noindent
Supposons que $X$ soit intègre et que nous ayons deux suites exactes courtes
$$
0 \to \mathcal{E}' \to \mathcal{E} \to \mathcal{Q} \to 0
\quad\text{et}\quad
0 \to \mathcal{F}' \to \mathcal{F} \to \mathcal{Q} \to 0
$$
où $\mathcal{E}$, $\mathcal{E}'$, $\mathcal{F}$, $\mathcal{F}'$ sont
localement libres de type fini. Problème : montrer que
$c(\mathcal{E})c(\mathcal{E}')^{-1} \cap [X] =
c(\mathcal{F})c(\mathcal{F}')^{-1} \cap [X]$.
Pour ce faire, considérons la suite exacte courte
$$
0 \to \mathcal{G} \to \mathcal{E} \oplus \mathcal{F} \to \mathcal{Q} \to 0
$$
qui définit $\mathcal{G}$. Puisque $\mathcal{Q}$ est de dimension de Tor $\leq 1$,
nous voyons que $\mathcal{G}$ est localement libre de type fini. Une chasse au diagramme
montre que le noyau de la surjection $\mathcal{G} \to \mathcal{F}$
s'envoie isomorphiquement sur $\mathcal{E}'$ dans $\mathcal{E}$ et
que le noyau de la surjection $\mathcal{G} \to \mathcal{E}$ s'envoie
isomorphiquement sur $\mathcal{F}'$ dans $\mathcal{F}$. (En travaillant localement
sur des ouverts affines, cela résulte du lemme de Schanuel, ou lui est équivalent, voir
Algèbre, Lemme \ref{algebra-lemma-Schanuel}.)
Nous concluons que
$$
c(\mathcal{E})c(\mathcal{F}') = c(\mathcal{G}) =
c(\mathcal{F})c(\mathcal{E}')
$$
comme souhaité.
\end{proof}

\begin{lemma}
\label{lemma-defined-by-envelope}
Soit $(S, \delta)$ comme dans la Situation \ref{situation-setup}.
Soit $X$ localement de type fini sur $S$. Soit $E \in D(\mathcal{O}_X)$
un objet parfait. Supposons qu'il existe une enveloppe
$f : Y \to X$ (Définition \ref{definition-envelope})
telle que $Lf^*E$ soit isomorphe dans $D(\mathcal{O}_Y)$
à un complexe localement borné $\mathcal{E}^\bullet$ de
$\mathcal{O}_Y$-modules localement libres de type fini. Alors il existe uniques des classes bivariantes
$c(E) \in \prod_{p \geq 0} A^p(X)$,
$ch(E) \in \prod_{p \geq 0} A^p(X) \otimes \mathbf{Q}$ et
$P_p(E) \in A^p(X)$, indépendantes du choix de $f : Y \to X$
et de $\mathcal{E}^\bullet$, telles que les restrictions de ces classes
à $Y$ soient égales à $c(\mathcal{E}^\bullet)$,
$ch(\mathcal{E}^\bullet)$ et $P_p(\mathcal{E}^\bullet)$.
\end{lemma}

\begin{proof}
Fixons $p \in \mathbf{Z}$. Nous démontrerons le lemme pour la classe
de Chern $c_p(E) \in A^p(X)$ et omettrons les arguments pour les autres cas.

\medskip\noindent
Soit $g : T \to X$ un morphisme localement de type fini pour lequel
il existe un complexe localement borné $\mathcal{E}^\bullet$ de
$\mathcal{O}_T$-modules localement libres de type fini représentant $Lg^*E$ dans $D(\mathcal{O}_T)$.
La classe bivariante $c_p(\mathcal{E}^\bullet) \in A^p(T)$
est indépendante du choix de $\mathcal{E}^\bullet$ par le
Lemme \ref{lemma-pre-derived-chern-class}.
Notons $c_p(Lg^*E) \in A^p(T)$ cette classe.
Pour tout morphisme ultérieur $h : T' \to T$ qui est localement
de type fini, en posant $g' = g \circ h$, nous voyons que
$L(g')^*E = L(g \circ h)^*E = Lh^*Lg^*E$ est représenté par
$h^*\mathcal{E}^\bullet$ dans $D(\mathcal{O}_{T'})$.
Nous concluons que $c_p(L(g')^*E)$ a un sens et est égal à la
restriction (Remarque \ref{remark-restriction-bivariant})
de $c_p(Lg^*E)$ à $T'$ (strictement parlant, cela requiert une application du
Lemme \ref{lemma-cap-cp-bivariant}).

\medskip\noindent
Soient $f : Y \to X$ et $\mathcal{E}^\bullet$ comme dans l'énoncé
du lemme. Nous obtenons une classe bivariante $c_p(E) \in A^p(X)$ par
une application du Lemme \ref{lemma-envelope-bivariant}
à $f : Y \to X$ et à la classe $c' = c_p(Lf^*E)$ que nous avons construite
dans le paragraphe précédent. L'hypothèse du lemme est
satisfaite car, d'après la discussion du paragraphe précédent, nous avons
$res_1(c') = c_p(Lg^*E) = res_2(c')$, où
$g = f \circ p = f \circ q : Y \times_X Y \to X$.

\medskip\noindent
Enfin, supposons que $f' : Y' \to X$ soit une seconde enveloppe
telle que $L(f')^*E$ soit représenté par un complexe borné de
$\mathcal{O}_{Y'}$-modules localement libres de type fini. Il s'ensuit
que les restrictions de $c_p(Lf^*E)$ et $c_p(L(f')^*E)$
à $Y \times_X Y'$ sont égales. Puisque $Y \times_X Y' \to X$
est une enveloppe (Lemmes \ref{lemma-base-change-envelope} et
\ref{lemma-composition-envelope}), nous voyons que nos deux candidats pour $c_p(E)$
coïncident par l'unicité dans le Lemme \ref{lemma-envelope-bivariant}.
\end{proof}

\begin{definition}
\label{definition-defined-on-perfect}
Soit $(S, \delta)$ comme dans la Situation \ref{situation-setup}.
Soit $X$ localement de type fini sur $S$. Soit $E \in D(\mathcal{O}_X)$
un objet parfait.
\begin{enumerate}
\item Nous disons que les {\it classes de Chern de $E$ sont définies}\footnote{Voir le
Lemme \ref{lemma-chern-classes-defined} pour quelques critères.} s'il existe
une enveloppe $f : Y \to X$ telle que $Lf^*E$ soit isomorphe dans
$D(\mathcal{O}_Y)$ à un complexe localement borné de
$\mathcal{O}_Y$-modules localement libres de type fini.
\item Si les classes de Chern de $E$ sont définies, alors nous définissons
$$
c(E) \in \prod\nolimits_{p \geq 0} A^p(X),\quad
ch(E) \in
\prod\nolimits_{p \geq 0} A^p(X) \otimes \mathbf{Q},\quad
P_p(E) \in A^p(X)
$$
par une application du Lemme \ref{lemma-defined-by-envelope}.
\end{enumerate}
\end{definition}

\noindent
Cette définition s'applique dans beaucoup, mais non dans toutes, les situations envisagées
dans ce chapitre, voir le Lemme \ref{lemma-chern-classes-defined}.
Peut-être existe-t-il une construction élémentaire de ces classes bivariantes pour
$E/X/(S,\delta)$ général comme dans la définition ; nous l'ignorons.

\begin{lemma}
\label{lemma-chern-classes-defined}
Soit $(S, \delta)$ comme dans la Situation \ref{situation-setup}.
Soit $X$ localement de type fini sur $S$. Soit $E \in D(\mathcal{O}_X)$
un objet parfait. Si l'une des conditions suivantes est satisfaite, alors
les classes de Chern de $E$ sont définies :
\begin{enumerate}
\item il existe une enveloppe $f : Y \to X$ telle que $Lf^*E$
soit isomorphe dans $D(\mathcal{O}_Y)$ à un complexe localement borné de
$\mathcal{O}_Y$-modules localement libres de type fini,
\item $E$ peut être représenté par un complexe borné de
$\mathcal{O}_X$-modules localement libres de type fini,
\item les composantes irréductibles de $X$ sont quasi-compactes,
\item $X$ est quasi-compact,
\item il existe un morphisme $X \to X'$ de schémas localement de type fini
sur $S$ tel que $E$ soit l'image inverse d'un objet parfait $E'$ sur $X'$
dont les classes de Chern sont définies, ou
\item ajouter davantage ici.
\end{enumerate}
\end{lemma}

\begin{proof}
La condition (1) est exactement la partie (1) de la Définition \ref{definition-defined-on-perfect}.
La condition (2) implique (1).

\medskip\noindent
Comme dans (3), supposons que les composantes irréductibles $X_i$ de $X$ soient quasi-compactes.
Nous considérons $X_i$ comme un sous-schéma fermé intègre réduit de $X$.
Le morphisme $\coprod X_i \to X$ est une enveloppe. Pour tout $i$, il
existe une enveloppe $X'_i \to X_i$ telle que $X'_i$ possède une famille ample
de modules inversibles, voir Plus sur les morphismes, Proposition
\ref{more-morphisms-proposition-envelope-with-resolution-property}.
Observons que $f : Y = \coprod X'_i \to X$ est une enveloppe ; un petit détail est omis.
Par Catégories dérivées de schémas, Lemme
\ref{perfect-lemma-resolution-property-ample-family},
chaque $X'_i$ possède la propriété de résolution.
Ainsi, l'objet parfait $L(f|_{X'_i})^*E$ de $D(\mathcal{O}_{X'_i})$
peut être représenté par un complexe borné
de $\mathcal{O}_{X'_i}$-modules localement libres de type fini, voir
Catégories dérivées de schémas, Lemme
\ref{perfect-lemma-resolution-property-perfect-complex}.
Cela démontre que (3) implique (1).

\medskip\noindent
La partie (4) implique (3).

\medskip\noindent
Soient $g : X \to X'$ et $E'$ comme dans la partie (5). Alors il existe une
enveloppe $f' : Y' \to X'$ telle que $L(f')^*E'$ soit représenté par un
complexe localement borné $(\mathcal{E}')^\bullet$ de $\mathcal{O}_{Y'}$-modules.
Alors le changement de base $f : Y \to X$ est une enveloppe par le
Lemme \ref{lemma-base-change-envelope}. De plus, la tiragee en arrière
$\mathcal{E}^\bullet = g^*(\mathcal{E}')^\bullet$ représente $Lf^*E$
et nous voyons que les classes de Chern de $E$ sont définies.
\end{proof}

\begin{lemma}
\label{lemma-chern-classes-computed}
Soit $(S, \delta)$ comme dans la Situation \ref{situation-setup}.
Soit $X$ localement de type fini sur $S$. Soit $E \in D(\mathcal{O}_X)$
un objet parfait. Supposons que les classes de Chern de $E$ soient définies.
Pour $g : W \to X$ localement de type fini avec $W$ intègre, il existe
un diagramme commutatif
$$
\xymatrix{
W' \ar[rd]_{g'} \ar[rr]_b & & W \ar[ld]^g \\
& X
}
$$
avec $W'$ intègre et $b : W' \to W$ propre birationnel tel que $L(g')^*E$
soit représenté par un complexe borné $\mathcal{E}^\bullet$ de
$\mathcal{O}_{W'}$-modules localement libres de rang constant et que nous ayons
$res(c_p(E)) = c_p(\mathcal{E}^\bullet)$ dans $A^p(W')$.
\end{lemma}

\begin{proof}
Choisissons une enveloppe $f : Y \to X$ telle que $Lf^*E$ soit isomorphe dans
$D(\mathcal{O}_Y)$ à un complexe localement borné $\mathcal{E}^\bullet$
de $\mathcal{O}_Y$-modules localement libres de type fini. Le changement de base
$Y \times_X W \to W$ de $f$ est une enveloppe par le
Lemme \ref{lemma-base-change-envelope}. Choisissons un point
$\xi \in Y \times_X W$ s'envoyant sur le point générique de $W$
et ayant le même corps résiduel. Considérons le sous-schéma fermé intègre
$W' \subset Y \times_X W$ de point générique $\xi$. La restriction
de la projection $Y \times_X W \to W$ à $W'$ est un morphisme propre
birationnel $b : W' \to W$. Posons $g' = g \circ b$. Enfin, considérons le complexe
image inverse $(W' \to Y)^*\mathcal{E}^\bullet$. C'est un complexe localement borné
de modules localement libres de type fini sur $W'$. Puisque $W'$ est intègre,
il s'ensuit qu'il est borné et que ses termes sont de rang constant.
Enfin, par construction, $(W' \to Y)^*\mathcal{E}^\bullet$
représente $L(g')^*E$ et, par construction, sa $p$-ième classe de Chern
donne la restriction de $c_p(E)$ par $W' \to X$. Cela achève la démonstration.
\end{proof}

\begin{lemma}
\label{lemma-commutative-chern-perfect}
Soit $(S, \delta)$ comme dans la Situation \ref{situation-setup}.
Soit $X$ localement de type fini sur $S$.
Soit $E \in D(\mathcal{O}_X)$ parfait. Si les classes de Chern
de $E$ sont définies, alors
\begin{enumerate}
\item $c_p(E)$ appartient au centre de l'algèbre $A^*(X)$, et
\item si $g : X' \to X$ est localement de type fini et $c \in A^*(X' \to X)$,
alors $c \circ c_p(E) = c_p(Lg^*E) \circ c$.
\end{enumerate}
\end{lemma}

\begin{proof}
La partie (1) résulte immédiatement de la partie (2). Soient $g : X' \to X$ et
$c \in A^*(X' \to X)$ comme dans (2). Pour montrer que
$c \circ c_p(E) - c_p(Lg^*E) \circ c = 0$, nous utilisons le critère du
Lemme \ref{lemma-bivariant-zero}. Nous pouvons donc supposer que $X$
est intègre et, par le Lemme \ref{lemma-chern-classes-computed},
nous pouvons même supposer que $E$ est représenté
par un complexe borné $\mathcal{E}^\bullet$
de $\mathcal{O}_X$-modules localement libres de type fini de rang constant.
Nous devons alors montrer que
$$
c \cap c_p(\mathcal{E}^\bullet) \cap [X] =
c_p(\mathcal{E}^\bullet) \cap c \cap [X] 
$$
dans $\CH_*(X')$. Cela résulte immédiatement du
Lemme \ref{lemma-cap-commutative-chern} et de la construction
de $c_p(\mathcal{E}^\bullet)$ comme polynôme en les
classes de Chern des modules localement libres $\mathcal{E}^n$.
\end{proof}

\begin{lemma}
\label{lemma-additivity-on-perfect}
Soit $(S, \delta)$ comme dans la Situation \ref{situation-setup}.
Soit $X$ localement de type fini sur $S$. Soit
$$
E_1 \to E_2 \to E_3 \to E_1[1]
$$
un triangle distingué d'objets parfaits dans $D(\mathcal{O}_X)$.
Si l'une des conditions suivantes est satisfaite :
\begin{enumerate}
\item il existe une enveloppe $f : Y \to X$ telle que
$Lf^*E_1 \to Lf^*E_2$ puisse être représenté par une application de complexes
localement bornés de $\mathcal{O}_Y$-modules localement libres de type fini,
\item $E_1 \to E_2$ peut être représenté par une application de complexes localement bornés
de $\mathcal{O}_X$-modules localement libres de type fini,
\item les composantes irréductibles de $X$ sont quasi-compactes,
\item $X$ est quasi-compact, ou
\item ajouter davantage ici,
\end{enumerate}
alors les classes de Chern de $E_1$, $E_2$, $E_3$ sont définies et nous avons
$c(E_2) = c(E_1) c(E_3)$, $ch(E_2) = ch(E_1) + ch(E_3)$ et
$P_p(E_2) = P_p(E_1) + P_p(E_3)$.
\end{lemma}

\begin{proof}
Soit $f : Y \to X$ une enveloppe et soit
$\alpha^\bullet : \mathcal{E}_1^\bullet \to \mathcal{E}_2^\bullet$
une application de complexes localement bornés de
$\mathcal{O}_Y$-modules localement libres de type fini représentant $Lf^*E_1 \to Lf^*E_2$.
Alors le cône $C(\alpha)^\bullet$ représente $Lf^*E_3$. Puisque
$C(\alpha)^n = \mathcal{E}_2^n \oplus \mathcal{E}_1^{n + 1}$,
nous voyons que $C(\alpha)^\bullet$ est un complexe localement borné
de $\mathcal{O}_Y$-modules localement libres de type fini.
Nous concluons que les classes de Chern de $E_1$, $E_2$, $E_3$
sont définies. De plus, rappelons que $c_p(E_1)$ est défini
comme l'unique élément de $A^p(X)$ dont la restriction est
$c_p(\mathcal{E}_1^\bullet)$ dans $A^p(Y)$. De même pour
$E_2$ et $E_3$. Il suffit donc
de démontrer $c(\mathcal{E}_2^\bullet) =
c(\mathcal{E}_1^\bullet) c(C(\alpha)^\bullet)$
dans $\prod_{p \geq 0} A^p(Y)$.
À son tour, il suffit de le démontrer après restriction
à une composante connexe de $Y$. Nous pouvons donc supposer que
les complexes $\mathcal{E}_1^\bullet$ et $\mathcal{E}_2^\bullet$
sont des complexes bornés de $\mathcal{O}_Y$-modules localement libres de type fini
de rang fixé. Dans ce cas, l'égalité souhaitée résulte de la
multiplicativité du Lemme \ref{lemma-additivity-chern-classes}.
Dans le cas de $ch$ ou $P_p$, nous utilisons les
Lemmes \ref{lemma-chern-character-additive}.

\medskip\noindent
Dans le paragraphe précédent, nous avons vu que le lemme est vrai si
la condition (1) est satisfaite. Puisque (2) implique (1), cela traite
le deuxième cas. Supposons (3). En raisonnant exactement comme dans la démonstration du
Lemme \ref{lemma-chern-classes-defined}, nous trouvons une enveloppe
$f : Y \to X$ tel que $Y$ soit une réunion disjointe $Y = \coprod Y_i$
de schémas quasi-compacts (et quasi-séparés), chacun possédant la
propriété de résolution. Nous pouvons alors représenter la restriction de
$Lf^*E_1 \to Lf^*E_2$ à $Y_i$ par un morphisme de complexes bornés
de modules localement libres de rang fini ; voir
Catégories dérivées de schémas, proposition
\ref{perfect-proposition-perfect-resolution-property}.
Nous voyons ainsi que la condition (3) implique la condition (1).
Bien entendu, la condition (4) implique la condition (3), ce qui achève
la démonstration.
\end{proof}

\begin{remark}
\label{remark-splitting-principle-perfect}
Les classes de Chern d'un complexe parfait, lorsqu'elles sont définies, satisfont une sorte de
principe de scindage. Plus précisément, supposons que $(S, \delta), X, E$ soient comme dans la
définition \ref{definition-defined-on-perfect}
et que les classes de Chern de $E$ soient définies.
Supposons que nous voulions démontrer une relation entre les classes bivariantes
$c_p(E)$, $P_p(E)$ et $ch_p(E)$. Pour cela, nous pouvons choisir une
enveloppe $f : Y \to X$ et un complexe localement borné
$\mathcal{E}^\bullet$ de $\mathcal{O}_X$-modules localement libres de rang fini
représentant $E$. Par l'unicité dans le lemme \ref{lemma-defined-by-envelope},
il suffit de démontrer la relation voulue entre les classes bivariantes
$c_p(\mathcal{E}^\bullet)$, $P_p(\mathcal{E}^\bullet)$ et
$ch_p(\mathcal{E}^\bullet)$. Nous pouvons donc remplacer $X$ par une composante
connexe de $Y$ et supposer que $E$ est représenté par un complexe borné
$\mathcal{E}^\bullet$ de modules localement libres de rang constant et fini.
En utilisant le principe de scindage
(lemme \ref{lemma-splitting-principle}), nous pouvons supposer que chaque
$\mathcal{E}^i$ possède une filtration dont les quotients
successifs $\mathcal{L}_{i, j}$ sont des modules inversibles.
En posant $x_{i, j} = c_1(\mathcal{L}_{i, j})$, nous voyons que
$$
c(E) =
\prod\nolimits_{i\text{ pair}} (1 + x_{i, j})
\prod\nolimits_{i\text{ impair}} (1 + x_{i, j})^{-1}
$$
et
$$
P_p(E) =  \sum\nolimits_{i\text{ pair}} (x_{i, j})^p -
\sum\nolimits_{i\text{ impair}} (x_{i, j})^p
$$
En prenant formellement le logarithme de l'expression de $c(E)$ ci-dessus,
nous obtenons
$$
\log(c(E)) = \sum (-1)^{p - 1}\frac{P_p(E)}{p}
$$
La construction des polynômes $P_p$ dans
l'exemple \ref{example-power-sum} montre que $P_p(E)$
est exactement la même expression dans les classes de Chern de $E$
que dans le cas des fibrés vectoriels ; autrement dit, nous avons
\begin{align*}
P_1(E) & = c_1(E), \\
P_2(E) & = c_1(E)^2 - 2c_2(E), \\
P_3(E) & = c_1(E)^3 - 3c_1(E)c_2(E) + 3c_3(E), \\
P_4(E) & = c_1(E)^4 - 4c_1(E)^2c_2(E) + 4c_1(E)c_3(E) + 2c_2(E)^2 - 4c_4(E),
\end{align*}
et ainsi de suite. D'autre part, la classe bivariante $P_0(E) = r(E) = ch_0(E)$
ne peut être retrouvée à partir de la classe de Chern $c(E)$ de $E$ ; la classe de Chern
ne contient pas l'information sur le rang du complexe.
\end{remark}

\begin{lemma}
\label{lemma-chern-classes-perfect-dual}
Dans la situation \ref{situation-setup}, soit $X$ localement de type fini sur $S$.
Soit $E \in D(\mathcal{O}_X)$ un objet parfait dont les classes de Chern sont
définies. Alors $c_i(E^\vee) = (-1)^i c_i(E)$, $P_i(E^\vee) = (-1)^iP_i(E)$
et $ch_i(E^\vee) = (-1)^ich_i(E)$ dans $A^i(X)$.
\end{lemma}

\begin{proof}
Première démonstration : raisonnons comme dans la démonstration du
lemme \ref{lemma-commutative-chern-perfect}
pour nous ramener au cas où $E$ est représenté
par un complexe borné de modules localement libres de rang
constant et fini, puis appliquons le lemme \ref{lemma-chern-classes-dual}.
Seconde démonstration : utilisons le principe de scindage exposé
dans la remarque \ref{remark-splitting-principle-perfect}
et le fait que les racines de Chern de $E^\vee$ sont les opposées
des racines de Chern de $E$.
\end{proof}

\begin{lemma}
\label{lemma-chern-class-perfect-tensor-invertible}
Dans la situation \ref{situation-setup}, soit $X$ localement de type fini sur $S$.
Soit $E$ un objet parfait de $D(\mathcal{O}_X)$ dont les classes de Chern
sont définies.
Soit $\mathcal{L}$ un $\mathcal{O}_X$-module inversible. Alors
$$
c_i(E \otimes \mathcal{L}) =
\sum\nolimits_{j = 0}^i
\binom{r - i + j}{j} c_{i - j}(E) c_1(\mathcal{L})^j
$$
pourvu que $E$ soit de rang constant $r \in \mathbf{Z}$.
\end{lemma}

\begin{proof}
Lorsque $E$ est localement libre de rang $r$, c'est le
lemme \ref{lemma-chern-classes-E-tensor-L}. Le lecteur peut déduire
le lemme de ce cas particulier par un calcul formel.
On peut aussi utiliser le principe de scindage de la
remarque \ref{remark-splitting-principle-perfect}.
Dans ce cas, on est ramené à démontrer le fait
d'algèbre suivant : si nous écrivons formellement
$$
\frac{\prod_{a = 1, \ldots, n} (1 + x_a)}{\prod_{n = 1, \ldots, m} (1 + y_b)}
= 1 + c_1 + c_2 + c_3 + \ldots
$$
où $c_i$ est homogène de degré $i$
dans $\mathbf{Z}[x_i, y_j]$, alors nous avons
$$
\frac{\prod_{a = 1, \ldots, n} (1 + x_a + t)}{\prod_{b = 1, \ldots, m} (1 + y_b + t)}
= \sum\nolimits_{i \geq 0} \sum\nolimits_{j = 0}^i
\binom{r - i + j}{j} c_{i - j} t^j
$$
où $r = n - m$. Nous omettons les détails.
\end{proof}

\begin{lemma}
\label{lemma-chern-classes-perfect-tensor-product}
Dans la situation \ref{situation-setup}, soit $X$ localement de type fini sur $S$.
Soient $E$ et $F$ des objets parfaits de $D(\mathcal{O}_X)$ dont les classes de Chern
sont définies. Alors nous avons
$$
c_1(E \otimes_{\mathcal{O}_X}^\mathbf{L} F) =
r(E) c_1(\mathcal{F}) + r(F) c_1(\mathcal{E})
$$
et, pour $c_2(E \otimes_{\mathcal{O}_X}^\mathbf{L} F)$, l'expression est
$$
r(E) c_2(F) + r(F) c_2(E) + {r(E) \choose 2} c_1(F)^2 +
(r(E)r(F) - 1) c_1(F)c_1(E) + {r(F) \choose 2} c_1(E)^2
$$
et ainsi de suite pour les classes de Chern supérieures dans $A^*(X)$. De même, nous avons
$ch(E \otimes_{\mathcal{O}_X}^\mathbf{L} F) = ch(E) ch(F)$
dans $A^*(X) \otimes \mathbf{Q}$. Plus précisément, nous avons
$$
P_p(E \otimes_{\mathcal{O}_X}^\mathbf{L} F) = \sum\nolimits_{p_1 + p_2 = p}
{p \choose p_1} P_{p_1}(E) P_{p_2}(F)
$$
dans $A^p(X)$.
\end{lemma}

\begin{proof}
Après avoir choisi une enveloppe $f : Y \to X$ telle que $Lf^*E$ et $Lf^*F$
puissent être représentés par des complexes localement bornés de
$\mathcal{O}_X$-modules localement libres de rang fini, ceci résulte, par calcul, du
résultat correspondant pour les fibrés vectoriels dans les
lemmes \ref{lemma-chern-classes-tensor-product} et
\ref{lemma-chern-character-multiplicative}.
Une meilleure démonstration consiste probablement à utiliser le principe de scindage comme dans la
remarque \ref{remark-splitting-principle-perfect}
et à réduire le lemme à des calculs dans des anneaux de polynômes 
que nous décrivons au paragraphe suivant.

\medskip\noindent
Soit $A$ un anneau commutatif (pour nous, ce sera le sous-anneau de l'anneau de Chow
bivariant de $X$ engendré par les classes de Chern).
Soit $S$ un ensemble fini muni d'applications $\epsilon : S \to \{\pm 1\}$
et $f : S \to A$. Définissons
$$
P_p(S, f , \epsilon) = \sum\nolimits_{s \in S} \epsilon(s) f(s)^p
$$
dans $A$. Étant donné un second triplet $(S', \epsilon', f')$,
l'égalité qu'il faut démontrer pour $P_p$ est
$$
P_p(S \times S', f + f' , \epsilon \epsilon') = 
\sum\nolimits_{p_1 + p_2 = p}
{p \choose p_1} P_{p_1}(S, f, \epsilon) P_{p_2}(S', f', \epsilon')
$$
Pour le voir, on se ramène à l'anneau de polynômes en les variables
$S \amalg S'$ et l'on montre que chaque terme $f(s)^if'(s')^j$ apparaît
dans les deux membres avec le même coefficient.
Pour vérifier les formules de $c_1(E \otimes_{\mathcal{O}_X}^\mathbf{L} F)$
et $c_2(E \otimes_{\mathcal{O}_X}^\mathbf{L} F)$, nous utilisons le principe de scindage
pour nous ramener à vérifier ces formules dans un anneau sans torsion.
Nous utilisons ensuite la relation entre $P_j(E)$ et $c_i(E)$ démontrée
dans la remarque \ref{remark-splitting-principle-perfect}. Par exemple,
$$
c_1(E \otimes F) = P_1(E \otimes F) = r(F)P_1(E) + r(E)P_1(F) =
r(F)c_1(E) + r(E)c_1(F)
$$
l'égalité du milieu résultant de $r(E) = P_0(E)$ par définition. De même, nous avons
\begin{align*}
& 2c_2(E \otimes F) \\
& = c_1(E \otimes F)^2 - P_2(E \otimes F) \\
& =
(r(F)c_1(E) + r(E)c_1(F))^2 -
r(F)P_2(E) - P_1(E)P_1(F) - r(E)P_2(F) \\
& =
(r(F)c_1(E) + r(E)c_1(F))^2 -
r(F)(c_1(E)^2 - 2c_2(E)) - c_1(E)c_1(F) - \\
& \quad r(E)(c_1(F)^2 - 2c_2(F))
\end{align*}
ce que le lecteur peut vérifier être égal à la formule de l'énoncé
du lemme à un facteur $2$ près.
\end{proof}







\section{Un cas élémentaire de classes de Chern localisées}
\label{section-preparation-localized-chern}

\noindent
Dans cette section, nous étudions quelques propriétés des classes bivariantes
construites dans le lemme suivant ; la plupart de ces propriétés résultent
immédiatement de la caractérisation donnée dans le lemme. Nous recommandons vivement au
lecteur de sauter le reste de la section.

\begin{lemma}
\label{lemma-silly}
Soit $(S, \delta)$ comme dans la situation \ref{situation-setup}. Soit $X$
localement de type fini sur $S$. Soient $i_j : X_j \to X$, $j = 1, 2$,
des immersions fermées telles que $X = X_1 \cup X_2$ ensemblistement. Soit
$E_2 \in D(\mathcal{O}_{X_2})$ un objet parfait. Supposons que
\begin{enumerate}
\item les classes de Chern de $E_2$ soient définies ;
\item la restriction $E_2|_{X_1 \cap X_2}$ soit nulle,
resp.\ isomorphe à un $\mathcal{O}_{X_1 \cap X_2}$-module localement libre de rang fini
de rang $< p$ placé en degré cohomologique $0$.
\end{enumerate}
Il existe alors une classe bivariante canonique
$$
P'_p(E_2),\text{ resp. }c'_p(E_2) \in A^p(X_2 \to X)
$$
caractérisée par la propriété
$$
P'_p(E_2) \cap i_{2, *} \alpha_2 = P_p(E_2) \cap \alpha_2
\quad\text{et}\quad
P'_p(E_2) \cap i_{1, *} \alpha_1 = 0,
$$
respectivement
$$
c'_p(E_2) \cap i_{2, *} \alpha_2 = c_p(E_2) \cap \alpha_2
\quad\text{et}\quad
c'_p(E_2) \cap i_{1, *} \alpha_1 = 0
$$
pour $\alpha_i \in \CH_k(X_i)$, et de même après tout changement de base
$X' \to X$ localement de type fini.
\end{lemma}

\begin{proof}
Nous allons utiliser les résultats de la section \ref{section-pre-derived}
sans autre mention.

\medskip\noindent
Supposons que $E_2|_{X_1 \cap X_2}$ soit nul.
Considérons un morphisme de schémas $X' \to X$
localement de type fini et notons $i'_j : X'_j \to X'$ le
changement de base de $i_j$. D'après le lemme \ref{lemma-exact-sequence-closed-chow},
nous pouvons écrire tout élément $\alpha' \in \CH_k(X')$ sous la forme
$i'_{1, *}\alpha'_1 + i'_{2, *}\alpha'_2$, où
$\alpha'_2 \in \CH_k(X'_2)$
est bien défini à un élément de l'image de l'image directe
par $X'_1 \cap X'_2 \to X'_2$ près. Nous pouvons alors poser
$P'_p(E_2) \cap \alpha' = P_p(E_2) \cap \alpha'_2 \in \CH_{k - p}(X'_2)$. Ceci
est bien défini puisque, par hypothèse, la restriction de $E_2$
à $X_1 \cap X_2$ est nulle.

\medskip\noindent
Si $E_2|_{X_1 \cap X_2}$ est isomorphe à un module localement libre de rang fini
sur $\mathcal{O}_{X_1 \cap X_2}$, de rang $< p$ et placé en
degré cohomologique $0$, alors $c_p(E_2|_{X_1 \cap X_2}) = 0$
pour des raisons de rang et nous pouvons raisonner exactement de la même manière.
\end{proof}

\begin{lemma}
\label{lemma-silly-independent}
Dans le lemme \ref{lemma-silly}, la classe bivariante
$P'_p(E_2)$, resp.\ $c'_p(E_2)$ dans $A^p(X_2 \to X)$
ne dépend pas du choix de $X_1$.
\end{lemma}

\begin{proof}
Supposons que $X_1' \subset X$ soit un autre sous-schéma fermé tel que
$X = X'_1 \cup X_2$ ensemblistement et que la restriction
$E_2|_{X'_1 \cap X_2}$ soit nulle, resp.\ isomorphe à un
$\mathcal{O}_{X'_1 \cap X_2}$-module localement libre de rang fini
de rang $< p$ placé en degré cohomologique $0$.
Alors $X = (X_1 \cap X'_1) \cup X_2$. Nous pouvons donc écrire
tout élément $\alpha \in \CH_k(X)$ sous la forme $i_*\beta + i_{2, *}\alpha_2$, avec
$\alpha_2 \in \CH_k(X'_2)$ et $\beta \in \CH_k(X_1 \cap X'_1)$.
Il est donc clair que
$P'_p(E_2) \cap \alpha = P_p(E_2) \cap \alpha_2 \in \CH_{k - p}(X_2)$,
resp.\  $c'_p(E_2) \cap \alpha = c_p(E_2) \cap \alpha_2 \in \CH_{k - p}(X_2)$,
ne dépend pas de l'emploi de $X_1$ ou de $X'_1$. Il en va de même
après tout changement de base.
\end{proof}

\begin{lemma}
\label{lemma-base-change-silly}
Dans le lemme \ref{lemma-silly}, soit $X' \to X$ un morphisme
localement de type fini. Notons $X' = X'_1 \cup X'_2$
et $E'_2 \in D(\mathcal{O}_{X'_2})$ les images inverses sur $X'$.
Alors la classe $P'_p(E_2')$, resp.\ $c'_p(E_2')$ dans
$A^p(X_2' \to X')$, construite dans le lemme \ref{lemma-silly} à partir de
$X' = X'_1 \cup X'_2$ et de $E_2'$, est la restriction
(remarque \ref{remark-restriction-bivariant})
de la classe $P'_p(E_2)$, resp.\ $c'_p(E_2)$ dans $A^p(X_2 \to X)$.
\end{lemma}

\begin{proof}
Cela résulte immédiatement de la caractérisation de ces classes dans le
lemme \ref{lemma-silly}.
\end{proof}

\begin{lemma}
\label{lemma-silly-silly}
Dans le lemme \ref{lemma-silly}, supposons que $E_2$ soit la restriction d'un objet
parfait $E \in D(\mathcal{O}_X)$ tel que $E|_{X_1}$ soit nul,
resp.\ isomorphe à un $\mathcal{O}_{X_1}$-module localement libre de rang fini
de rang $< p$ placé en degré cohomologique $0$.
Si les classes de Chern de $E$ sont définies, alors
$i_{2, *} \circ P'_p(E_2) = P_p(E)$,
resp.\ $i_{2, *} \circ c'_p(E_2) = c_p(E)$
(où $\circ$ est comme dans le lemme \ref{lemma-push-proper-bivariant}).
\end{lemma}

\begin{proof}
Supposons d'abord que $E|_{X_1}$ soit nul.
Avec les notations de la démonstration du lemme \ref{lemma-silly},
le lemme résulte dans ce cas de
\begin{align*}
P_p(E) \cap \alpha'
& =
i'_{1, *}(P_p(E) \cap \alpha'_1) +
i'_{2, *}(P_p(E) \cap \alpha'_2) \\
& =
i'_{1, *}(P_p(E|_{X_1}) \cap \alpha'_1) +
i'_{2, *}(P'_p(E_2) \cap \alpha') \\
& =
i'_{2, *}(P'_p(E_2) \cap \alpha')
\end{align*}
Le cas où $E|_{X_1}$ est isomorphe à un
$\mathcal{O}_{X_1}$-module localement libre de rang fini, de rang $< p$ et placé en degré cohomologique $0$,
est analogue.
\end{proof}

\begin{lemma}
\label{lemma-silly-shrink}
Dans le lemme \ref{lemma-silly}, supposons donnés des sous-schémas fermés
$X'_2 \subset X_2$ et $X_1 \subset X'_1 \subset X$ tels que
$X = X'_1 \cup X'_2$ ensemblistement. Supposons que $E_2|_{X'_1 \cap X_2}$
soit nul, resp.\ isomorphe à un module localement libre de rang fini
de rang $< p$ placé en degré $0$. Alors nous avons
$(X'_2 \to X_2)_* \circ P'_p(E_2|_{X'_2}) = P'_p(E_2)$,
resp.\  $(X'_2 \to X_2)_* \circ c'_p(E_2|_{X'_2}) = c_p(E_2)$
(où $\circ$ est comme dans le lemme \ref{lemma-push-proper-bivariant}).
\end{lemma}

\begin{proof}
Cela résulte immédiatement de la caractérisation de ces classes
dans le lemme \ref{lemma-silly}.
\end{proof}

\begin{lemma}
\label{lemma-silly-commutes}
Dans le lemme \ref{lemma-silly}, soit $f : Y \to X$ localement de type fini
et soit $c \in A^*(Y \to X)$. Alors
$$
c \circ P'_p(E_2) = P'_p(Lf_2^*E_2) \circ c
\quad\text{resp.}\quad
c \circ c'_p(E_2) = c'_p(Lf_2^*E_2) \circ c
$$
dans $A^*(Y_2 \to Y)$, où $f_2 : Y_2 \to X_2$ est le changement de base de $f$.
\end{lemma}

\begin{proof}
Soit $\alpha \in \CH_k(X)$. Nous pouvons écrire
$$
\alpha = \alpha_1 + \alpha_2
$$
avec $\alpha_i \in \CH_k(X_i)$ ; nous omettons les images directes
par les immersions fermées $X_i \to X$. Le lecteur vérifie alors que
$c'_p(E_2) \cap \alpha = c_p(E_2) \cap \alpha_2$,
$c \cap c'_p(E_2) \cap \alpha = c \cap c_p(E_2) \cap \alpha_2$,
$c \cap \alpha = c \cap \alpha_1 + c \cap \alpha_2$, et
$c'_p(Lf_2^*E_2) \cap c \cap \alpha = c_p(Lf_2^*E_2) \cap c \cap \alpha_2$.
Nous concluons par le lemme \ref{lemma-commutative-chern-perfect}.
\end{proof}

\begin{lemma}
\label{lemma-silly-compose}
Dans le lemme \ref{lemma-silly}, supposons que $E_2|_{X_1 \cap X_2}$ soit nul. Alors
\begin{align*}
P'_1(E_2) & = c'_1(E_2), \\
P'_2(E_2) & = c'_1(E_2)^2 - 2c'_2(E_2), \\
P'_3(E_2) & = c'_1(E_2)^3 - 3c'_1(E_2)c'_2(E_2) + 3c'_3(E_2), \\
P'_4(E_2) & = c'_1(E_2)^4 - 4c'_1(E_2)^2c'_2(E_2) +
4c'_1(E_2)c'_3(E_2) + 2c'_2(E_2)^2 - 4c'_4(E_2),
\end{align*}
et ainsi de suite, la multiplication étant celle de la remarque \ref{remark-ring-loc-classes}.
\end{lemma}

\begin{proof}
L'énoncé a un sens car le faisceau nul est de rang $< 1$ et
les classes $c'_p(E_2)$ sont donc définies pour tout $p \geq 1$. Les égalités
résultent immédiatement de la caractérisation des classes construites
par le lemme \ref{lemma-silly} et du résultat correspondant pour
l'intersection avec les classes de Chern de $E_2$ donné dans la
remarque \ref{remark-splitting-principle-perfect}.
\end{proof}

\begin{lemma}
\label{lemma-silly-sum-c}
Soit $(S, \delta)$ comme dans la situation \ref{situation-setup}. Soit $X$
localement de type fini sur $S$. Soient $i_j : X_j \to X$, $j = 1, 2$,
des immersions fermées telles que $X = X_1 \cup X_2$ ensemblistement. Soient
$E, F \in D(\mathcal{O}_X)$ des objets parfaits. Supposons que
\begin{enumerate}
\item les classes de Chern de $E$ et de $F$ soient définies ;
\item les restrictions $E|_{X_1 \cap X_2}$ et $F|_{X_1 \cap X_2}$
soient isomorphes à deux $\mathcal{O}_{X_1}$-modules localement libres,
de rangs respectivement $< p$ et $< q$, placés en degré cohomologique $0$.
\end{enumerate}
Avec la notation de la remarque \ref{remark-ring-loc-classes}, posons
$$
c^{(p)}(E) = 1 + c_1(E) + \ldots + c_{p - 1}(E) +
c'_p(E|_{X_2}) + c'_{p + 1}(E|_{X_2}) + \ldots \in A^{(p)}(X_2 \to X)
$$
où $c'_p(E|_{X_2})$ est comme dans le lemme \ref{lemma-silly}. De même
pour $c^{(q)}(F)$ et $c^{(p + q)}(E \oplus F)$.
Alors $c^{(p + q)}(E \oplus F) = c^{(p)}(E)c^{(q)}(F)$
dans $A^{(p + q)}(X_2 \to X)$.
\end{lemma}

\begin{proof}
Cela résulte immédiatement de la caractérisation des classes dans le
lemme \ref{lemma-silly} et de l'additivité dans le
lemme \ref{lemma-additivity-on-perfect}.
\end{proof}

\begin{lemma}
\label{lemma-silly-sum-P}
Soit $(S, \delta)$ comme dans la situation \ref{situation-setup}. Soit $X$
localement de type fini sur $S$. Soient $i_j : X_j \to X$, $j = 1, 2$,
des immersions fermées telles que $X = X_1 \cup X_2$ ensemblistement. Soient
$E, F \in D(\mathcal{O}_{X_2})$ des objets parfaits. Supposons que
\begin{enumerate}
\item les classes de Chern de $E$ et de $F$ soient définies ;
\item les restrictions $E|_{X_1 \cap X_2}$ et $F|_{X_1 \cap X_2}$ soient nulles,
\end{enumerate}
Notons $P'_p(E), P'_p(F), P'_p(E \oplus F) \in A^p(X_2 \to X)$, pour $p \geq 0$,
les classes construites dans le lemme \ref{lemma-silly}. Alors
$P'_p(E \oplus F) = P'_p(E) + P'_p(F)$.
\end{lemma}

\begin{proof}
Cela résulte immédiatement de la caractérisation des classes dans le
lemme \ref{lemma-silly} et de l'additivité dans le
lemme \ref{lemma-additivity-on-perfect}.
\end{proof}

\begin{lemma}
\label{lemma-silly-tensor-invertible}
Dans le lemme \ref{lemma-silly}, supposons que $E_2$ soit de rang constant $0$.
Soit $\mathcal{L}$ un $\mathcal{O}_X$-module inversible. Alors
$$
c'_i(E_2 \otimes \mathcal{L}) =
\sum\nolimits_{j = 0}^i
\binom{- i + j}{j} c'_{i - j}(E_2) c_1(\mathcal{L})^j
$$
\end{lemma}

\begin{proof}
L'hypothèse sur le rang implique que $E_2|_{X_1 \cap X_2}$ est nul.
Ainsi $c'_i(E_2)$ est défini pour tout $i \geq 1$ et l'énoncé
a un sens. L'égalité elle-même résulte
immédiatement du lemme \ref{lemma-chern-class-perfect-tensor-invertible}
et de la caractérisation de $c'_i$ dans le lemme \ref{lemma-silly}.
\end{proof}

\begin{lemma}
\label{lemma-silly-tensor-product}
Dans la situation \ref{situation-setup}, soit $X$ localement de type fini sur $S$.
Soit
$$
X = X_1 \cup X_2 = X'_1 \cup X'_2
$$
deux écritures de $X$ comme réunion ensembliste de sous-schémas fermés.
Soient $E$, $E'$ des objets parfaits de $D(\mathcal{O}_X)$
dont les classes de Chern sont définies.
Supposons que $E|_{X_1}$ et $E'|_{X'_1}$ soient nuls\footnote{Il existe vraisemblablement
une variante de ce lemme où l'on suppose seulement que ces restrictions sont
isomorphes à deux modules localement libres,
de rangs respectivement $< p$ et $< p'$.} pour $i = 1, 2$. Notons
\begin{enumerate}
\item $r = P'_0(E) \in A^0(X_2 \to X)$ et
$r' = P'_0(E') \in A^0(X'_2 \to X)$,
\item $\gamma_p = c'_p(E|_{X_2}) \in A^p(X_2 \to X)$ et
$\gamma'_p = c'_p(E'|_{X'_2}) \in A^p(X'_2 \to X)$,
\item $\chi_p = P'_p(E|_{X_2}) \in A^p(X_2 \to X)$ et
$\chi'_p = P'_p(E'|_{X'_2}) \in A^p(X'_2 \to X)$
\end{enumerate}
les classes construites dans le lemme \ref{lemma-silly}. Alors nous avons
$$
c'_1((E \otimes_{\mathcal{O}_X}^\mathbf{L} E')|_{X_2 \cap X'_2}) =
r \gamma'_1 + r' \gamma_1
$$
dans $A^1(X_2 \cap X'_2 \to X)$ et
$$
c'_2((E \otimes_{\mathcal{O}_X}^\mathbf{L} E')|_{X_2 \cap X'_2}) =
r \gamma'_2 + r' \gamma_2 + {r \choose 2} (\gamma'_1)^2 +
(rr' - 1) \gamma'_1\gamma_1 + {r' \choose 2} \gamma_1^2
$$
dans $A^2(X_2 \cap X'_2 \to X)$, et ainsi de suite pour les classes de Chern supérieures.
De même, nous avons
$$
P'_p((E \otimes_{\mathcal{O}_X}^\mathbf{L} E')|_{X_2 \cap X'_2}) =
\sum\nolimits_{p_1 + p_2 = p}
{p \choose p_1} \chi_{p_1} \chi'_{p_2}
$$
dans $A^p(X_2 \cap X'_2 \to X)$.
\end{lemma}

\begin{proof}
Observons d'abord que l'énoncé a un sens. En effet, nous avons
$X = (X_2 \cap X'_2) \cup Y$, où
$Y = (X_1 \cap X'_1) \cup (X_1 \cap X'_2) \cup (X_2 \cap X'_1)$,
et la restriction de l'objet $E \otimes_{\mathcal{O}_X}^\mathbf{L} E'$
à $Y$ est nulle.
Les égalités elles-mêmes résultent de la caractérisation
de nos classes dans le lemme \ref{lemma-silly}
et des égalités du lemme \ref{lemma-chern-classes-perfect-tensor-product}.
Nous omettons les détails.
\end{proof}







\section{Gysin à l'infini}
\label{section-gysin-at-infty}

\noindent
Cette section porte sur la classe bivariante construite dans le prochain
lemme. Nous recommandons vivement au lecteur de sauter le reste de la section.

\begin{lemma}
\label{lemma-gysin-at-infty}
Soit $(S, \delta)$ comme dans la situation \ref{situation-setup}.
Soit $X$ localement de type fini sur $S$. Soit
$b : W \to \mathbf{P}^1_X$ un morphisme propre de schémas
qui est un isomorphisme au-dessus de $\mathbf{A}^1_X$.
Notons $i_\infty : W_\infty \to W$ l'image inverse du diviseur
$D_\infty \subset \mathbf{P}^1_X$ dont le complémentaire est $\mathbf{A}^1_X$.
Il existe alors une classe bivariante canonique
$$
C \in A^0(W_\infty \to X)
$$
ayant la propriété suivante :
$i_{\infty, *}(C \cap \alpha) = i_{0, *}\alpha$
pour $\alpha \in \CH_k(X)$, et de même après tout changement de base par
$X' \to X$ localement de type fini.
\end{lemma}

\begin{proof}
Étant donné $\alpha \in \CH_k(X)$, il existe $\beta \in \CH_{k + 1}(W)$
dont la restriction est l'image inverse plate de $\alpha$ sur $b^{-1}(\mathbf{A}^1_X)$ ; voir
le lemme \ref{lemma-exact-sequence-open}.
Un second choix, encore noté $\beta$, diffère du choix initial $\beta$ par un cycle
à support dans $W_\infty$ ; voir
le lemme \ref{lemma-restrict-to-open}. Puisque le fibré normal du diviseur de Cartier
effectif $D_\infty \subset \mathbf{P}^1_X$ de
(\ref{equation-zero-infty}) est trivial,
l'homomorphisme de Gysin $i_\infty^*$ annule les classes de cycles
à support dans $W_\infty$ ; voir la remarque \ref{remark-gysin-on-cycles}.
Ainsi, poser $C \cap \alpha = i_\infty^*\beta$ donne une définition indépendante des choix.

\medskip\noindent
Puisque $W_\infty$ et $W_0 = X \times \{0\}$
sont les images inverses des diviseurs de Cartier effectifs rationnellement équivalents
$D_0, D_\infty$ dans $\mathbf{P}^1_X$, nous voyons que $i_\infty^*\beta$ et
$i_0^*\beta$ ont la même image comme classe de cycle sur $W$ ; en effet, tous deux
représentent la classe $c_1(\mathcal{O}_{\mathbf{P}^1_X}(1)) \cap \beta$ d'après le
lemme \ref{lemma-support-cap-effective-Cartier}. Par notre choix de
$\beta$, nous avons $i_0^*\beta = \alpha$ comme cycles sur
$W_0 = X \times \{0\}$ ; voir par exemple le
lemme \ref{lemma-relative-effective-cartier}.
Nous voyons donc que $i_{\infty, *}(C \cap \alpha) = i_{0, *}\alpha$,
comme l'affirme le lemme.

\medskip\noindent
Remarquons que les hypothèses sur $b$ sont préservées par tout changement de base
par $X' \to X$ localement de type fini. Nous obtenons donc une opération
$C \cap - : \CH_k(X') \to \CH_k(W'_\infty)$ par la même construction que ci-dessus.
Pour voir que cette famille d'opérations définit une classe bivariante,
considérons le diagramme
$$
\xymatrix{
& & & \CH_*(X) \ar[d]^{\text{image inverse plate}} \\
\CH_{* + 1}(W_\infty) \ar[r] \ar[rd]^0 &
\CH_{* + 1}(W) \ar[d]^{i_\infty^*} \ar[rr]^{\text{image inverse plate}} & &
\CH_{* + 1}(\mathbf{A}^1_X) \ar[r] \ar@{..>}[lld]^{C \cap -} &
0 \\
& \CH_*(W_\infty)
}
$$
pour $X$ comme indiqué, ainsi que le changement de base de ce diagramme pour tout $X' \to X$.
Nous savons que l'image inverse plate et $i_\infty^*$ sont des opérations bivariantes ; voir les
lemmes \ref{lemma-flat-pullback-bivariant} et \ref{lemma-gysin-bivariant}.
Un argument formel (faisant intervenir d'immenses diagrammes de schémas et de leurs
groupes de Chow) montre alors que la flèche pointillée est une opération bivariante.
\end{proof}

\begin{lemma}
\label{lemma-base-change-gysin-at-infty}
Dans le lemme \ref{lemma-gysin-at-infty}, soit $X' \to X$ un morphisme
localement de type fini. Notons $b' : W' \to \mathbf{P}^1_{X'}$
et $i'_\infty : W'_\infty \to W'$ les changements de base de $b$ et de $i_\infty$.
Alors la classe $C' \in A^0(W'_\infty \to X')$ construite comme dans le
lemme \ref{lemma-gysin-at-infty} à l'aide de $b'$ est la restriction
(remarque \ref{remark-restriction-bivariant}) de $C$.
\end{lemma}

\begin{proof}
Cela résulte immédiatement de la construction et du fait qu'un énoncé analogue
vaut pour l'image inverse plate et $i_\infty^*$.
\end{proof}

\begin{lemma}
\label{lemma-gysin-at-infty-independent}
Dans le lemme \ref{lemma-gysin-at-infty}, soit $g : W' \to W$ un morphisme propre
qui est un isomorphisme au-dessus de $\mathbf{A}^1_X$. Soient
$C' \in A^0(W'_\infty \to X)$ et $C \in A^0(W_\infty \to X)$
les classes construites dans le lemme \ref{lemma-gysin-at-infty}.
Alors $g_{\infty, *} \circ C' = C$ dans $A^0(W_\infty \to X)$.
\end{lemma}

\begin{proof}
Posons $b' = b \circ g : W' \to \mathbf{P}^1_X$. Notons
$i'_\infty : W'_\infty \to W'$ le morphisme d'inclusion.
Notons $g_\infty : W'_\infty \to W_\infty$ la restriction de $g$.
Étant donné $\alpha \in \CH_k(X)$, choisissons $\beta' \in \CH_{k + 1}(W')$
dont la restriction est l'image inverse plate de $\alpha$ sur $(b')^{-1}\mathbf{A}^1_X$.
Alors $\beta = g_*\beta' \in \CH_{k + 1}(W)$ se restreint en
l'image inverse plate de $\alpha$ sur $b^{-1}\mathbf{A}^1_X$.
Nous avons alors $i_\infty^*\beta = g_{\infty, *}(i'_\infty)^*\beta'$
d'après le lemme \ref{lemma-closed-in-X-gysin}.
Ceci, ainsi que le fait correspondant après changement de base par des
morphismes $X' \to X$ localement de type fini, équivaut
à l'assertion du lemme.
\end{proof}

\begin{lemma}
\label{lemma-homomorphism-pre}
Dans le lemme \ref{lemma-gysin-at-infty}, nous avons
$C \circ (W_\infty \to X)_* \circ i_\infty^* = i_\infty^*$.
\end{lemma}

\begin{proof}
Soit $\beta \in \CH_{k + 1}(W)$. Notons $i_0 : X = X \times \{0\} \to W$
l'immersion fermée de la fibre au-dessus de $0$ dans $\mathbf{P}^1$. Alors
$(W_\infty \to X)_* i_\infty^* \beta = i_0^*\beta$ dans $\CH_k(X)$, car
$i_{\infty, *}i_\infty^*\beta$ et $i_{0, *}i_0^*\beta$
représentent la même classe sur $W$ (par exemple d'après le
lemme \ref{lemma-support-cap-effective-Cartier})
et ont donc la même image directe sur $X$.
La restriction de $\beta$ à $b^{-1}(\mathbf{A}^1_X)$
se restreint en l'image inverse plate de
$i_0^*\beta = (W_\infty \to X)_* i_\infty^* \beta$, car nous pouvons le vérifier
après image inverse par $i_0$ ; voir les
lemmes \ref{lemma-linebundle} et \ref{lemma-linebundle-formulae}.
Nous pouvons donc utiliser $\beta$ pour calculer l'image de
$(W_\infty \to X)_*i_\infty^*\beta$ par $C$,
et nous obtenons le résultat voulu.
\end{proof}

\begin{lemma}
\label{lemma-gysin-at-infty-commutes}
Dans le lemme \ref{lemma-gysin-at-infty}, soit $f : Y \to X$ un morphisme
localement de type fini et soit $c \in A^*(Y \to X)$. Alors $C \circ c = c \circ C$
dans $A^*(W_\infty \times_X Y \to X)$.
\end{lemma}

\begin{proof}
Considérons le diagramme commutatif
$$
\xymatrix{
W_\infty \times_X Y \ar@{=}[r] &
W_{Y, \infty} \ar[r]_{i_{Y, \infty}} \ar[d] &
W_Y \ar[r]_{b_Y} \ar[d] &
\mathbf{P}^1_Y \ar[r]_{p_Y} \ar[d] &
Y \ar[d]^f \\
& W_\infty \ar[r]^{i_\infty} &
W \ar[r]^b &
\mathbf{P}^1_X \ar[r]^p &
X
}
$$
dont les carrés sont cartésiens. Pour un élément $\alpha \in \CH_k(X)$,
choisissons $\beta \in \CH_{k + 1}(W)$ dont la restriction à $b^{-1}(\mathbf{A}^1_X)$
est l'image inverse plate de $\alpha$. Alors $c \cap \beta$ est une classe
dans $\CH_*(W_Y)$ dont la restriction à $b_Y^{-1}(\mathbf{A}^1_Y)$
est l'image inverse plate de $c \cap \alpha$. Ensuite, nous avons
$$
i_{Y, \infty}^*(c \cap \beta) = c \cap i_\infty^*\beta
$$
car $c$ est une classe bivariante. Cela signifie exactement que
$C \cap c \cap \alpha = c \cap C \cap \alpha$. Le même argument
vaut après tout changement de base par $X' \to X$ localement de type fini.
Cela démontre le lemme.
\end{proof}





\section{Préparation aux classes de Chern localisées}
\label{section-preparation-localized-chern-II}

\noindent
Dans cette section, nous étudions quelques propriétés des classes bivariantes
construites dans le lemme suivant. Nous recommandons vivement au
lecteur de sauter le reste de la section.

\begin{lemma}
\label{lemma-localized-chern-pre}
Soit $(S, \delta)$ comme dans la situation \ref{situation-setup}. Soit $X$
localement de type fini sur $S$. Soit $Z \subset X$ un sous-schéma fermé.
Soit
$$
b : W \longrightarrow \mathbf{P}^1_X
$$
un morphisme propre de schémas. Soit $Q \in D(\mathcal{O}_W)$ un
objet parfait. Notons $W_\infty \subset W$ l'image inverse du diviseur
$D_\infty \subset \mathbf{P}^1_X$ dont le complémentaire est $\mathbf{A}^1_X$.
Nous supposons que
\begin{enumerate}
\item[(A0)] les classes de Chern de $Q$ sont définies
(section \ref{section-pre-derived}) ;
\item[(A1)] $b$ est un isomorphisme au-dessus de $\mathbf{A}^1_X$ ;
\item[(A2)] il existe un sous-schéma fermé $T \subset W_\infty$
contenant tous les points de $W_\infty$ situés au-dessus de $X \setminus Z$ tel que
$Q|_T$ soit nul, resp.\ isomorphe à un
$\mathcal{O}_T$-module localement libre de rang fini, de rang $< p$ et placé en degré cohomologique $0$.
\end{enumerate}
Il existe alors une classe bivariante canonique
$$
P'_p(Q),\text{ resp. }c'_p(Q) \in A^p(Z \to X)
$$
telle que
$(Z \to X)_* \circ P'_p(Q) = P_p(Q|_{X \times \{0\}})$,
resp.\ $(Z \to X)_* \circ c'_p(Q) = c_p(Q|_{X \times \{0\}})$.
\end{lemma}

\begin{proof}
Notons $E \subset W_\infty$ l'image inverse de $Z$. Alors
$W_\infty = T \cup E$ et $b$ induit un morphisme propre $E \to Z$.
Notons $C \in A^0(W_\infty \to X)$ la classe bivariante construite
dans le lemme \ref{lemma-gysin-at-infty}. Notons $P'_p(Q|_E)$, resp.\ $c'_p(Q|_E)$
dans $A^p(E \to W_\infty)$ la classe bivariante construite
dans le lemme \ref{lemma-silly}. Cela a un sens car
$(Q|_E)|_{T \cap E}$ est nul, resp.\ isomorphe à un
$\mathcal{O}_{E \cap T}$-module localement libre de rang fini, de rang $< p$ et placé en
degré cohomologique $0$ par l'hypothèse (A2). Nous définissons alors
$$
P'_p(Q) = (E \to Z)_* \circ P'_p(Q|_E) \circ C,\text{ resp. }
c'_p(Q) = (E \to Z)_* \circ c'_p(Q|_E) \circ C
$$
C'est une classe bivariante ; voir le lemme \ref{lemma-push-proper-bivariant}.
Puisque $E \to Z \to X$ est égal à $E \to W_\infty \to W \to X$, nous voyons que
\begin{align*}
(Z \to X)_* \circ c'_p(Q)
& =
(W \to X)_* \circ i_{\infty, *} \circ (E \to W_\infty)_*
\circ c'_p(Q|_E) \circ C \\
& =
(W \to X)_* \circ i_{\infty, *} \circ c_p(Q|_{W_\infty}) \circ C \\
& =
(W \to X)_* \circ c_p(Q) \circ i_{\infty, *} \circ C \\
& =
(W \to X)_*\circ c_p(Q) \circ i_{0, *} \\
& =
(W \to X)_* \circ i_{0, *} \circ c_p(Q|_{X \times \{0\}}) \\
& =
c_p(Q|_{X \times \{0\}})
\end{align*}
La deuxième égalité résulte du lemme \ref{lemma-silly-silly}.
La troisième égalité vient du fait que $c_p(Q)$ est une classe bivariante.
La quatrième égalité résulte du lemme \ref{lemma-gysin-at-infty}.
La cinquième égalité vient du fait que $c_p(Q)$ est une classe bivariante.
La dernière égalité vient de ce que $(W_0 \to W) \circ (W \to X)$
est l'identité de $X$ si l'on identifie $W_0$ à $X$ comme nous l'avons
fait ci-dessus. La même suite d'égalités permet de
démontrer la propriété pour $P'_p(Q)$.
\end{proof}

\begin{lemma}
\label{lemma-base-change-localized-chern-pre}
Dans le lemme \ref{lemma-localized-chern-pre}, soit $X' \to X$ un morphisme
localement de type fini. Notons
$Z'$, $b' : W' \to \mathbf{P}^1_{X'}$ et $T' \subset W'_\infty$
les changements de base de $Z$, $b : W \to \mathbf{P}^1_X$ et $T \subset W_\infty$.
Posons $Q' = (W' \to W)^*Q$. Alors la classe
$P'_p(Q')$, resp.\ $c'_p(Q')$ dans $A^p(Z' \to X')$, construite comme dans le
lemme \ref{lemma-localized-chern-pre} à l'aide de $b'$, $Q'$ et $T'$,
est la restriction (remarque \ref{remark-restriction-bivariant})
de la classe $P'_p(Q)$, resp.\ $c'_p(Q)$ dans $A^p(Z \to X)$.
\end{lemma}

\begin{proof}
Rappelons que la construction est la suivante :
$$
P'_p(Q) = (E \to Z)_* \circ P'_p(Q|_E) \circ C,\text{ resp. }
c'_p(Q) = (E \to Z)_* \circ c'_p(Q|_E) \circ C
$$
Le lemme résulte donc de la propriété de changement de base correspondante
pour $C$ (lemme \ref{lemma-base-change-gysin-at-infty})
et du fait que la même propriété de changement de base vaut pour les classes
construites dans le lemme \ref{lemma-silly} (nous omettons un petit détail).
\end{proof}

\begin{lemma}
\label{lemma-localized-chern-pre-independent}
Dans le lemme \ref{lemma-localized-chern-pre}, la classe bivariante
$P'_p(Q)$, resp.\ $c'_p(Q)$,
est indépendante du choix du sous-schéma fermé $T$.
De plus, étant donné un morphisme propre $g : W' \to W$ qui est un
isomorphisme au-dessus de $\mathbf{A}^1_X$, en posant $Q' = g^*Q$,
nous avons $P'_p(Q) = P'_p(Q')$, resp.\ $c'_p(Q) = c'_p(Q')$.
\end{lemma}

\begin{proof}
L'indépendance à l'égard de $T$ résulte immédiatement du
lemme \ref{lemma-silly-independent}.

\medskip\noindent
Soit $g : W' \to W$ un morphisme propre qui est un isomorphisme au-dessus de
$\mathbf{A}^1_X$. Remarquons que $T' = g^{-1}(T) \subset W'_\infty$
est un sous-schéma fermé satisfaisant à (A2), de sorte que l'opérateur
$P'_p(Q')$, resp.\ $c'_p(Q')$ dans $A^p(Z \to X)$,
correspondant à $b' = b \circ g : W' \to \mathbf{P}^1_X$
et à $Q'$, est défini. Notons $E' \subset W'_\infty$
l'image inverse de $Z$ dans $W'_\infty$. Rappelons que
$$
c'_p(Q') = (E' \to Z)_* \circ c'_p(Q'|_{E'}) \circ C'
$$
avec $C' \in A^0(W'_\infty \to X)$ et
$c'_p(Q'|_{E'}) \in A^p(E' \to W'_\infty)$.
D'après le lemme \ref{lemma-gysin-at-infty-independent}, nous avons
$g_{\infty, *} \circ C' = C$. Remarquons que $E'$ est aussi
l'image inverse de $E$ dans $W'_\infty$ par $g_\infty$.
De plus, puisque $Q' = g^*Q$, nous constatons que $c'_p(Q'|_{E'})$ est simplement la
restriction de $c'_p(Q|_E)$ aux schémas au-dessus de $W'_\infty$ ; voir la
remarque \ref{remark-restriction-bivariant}. Nous obtenons donc
\begin{align*}
c'_p(Q')
& = 
(E' \to Z)_* \circ c'_p(Q'|_{E'}) \circ C' \\
& =
(E \to Z)_* \circ (E' \to E)_* \circ c'_p(Q|_E) \circ C' \\
& =
(E \to Z)_* \circ c'_p(Q|_E) \circ g_{\infty, *} \circ C' \\
& =
(E \to Z)_* \circ c'_p(Q|_E) \circ C \\
& =
c'_p(Q)
\end{align*}
Dans la troisième égalité, nous avons utilisé que $c'_p(Q|_E)$
commute à l'image directe propre puisqu'il s'agit d'une
classe bivariante. L'égalité $P'_p(Q) = P'_p(Q')$
se démontre exactement de la même manière.
\end{proof}

\begin{lemma}
\label{lemma-homomorphism}
Dans le lemme \ref{lemma-localized-chern-pre}, supposons que $Q|_T$ soit isomorphe
à un $\mathcal{O}_T$-module localement libre de rang fini, de rang $< p$.
Notons $C \in A^0(W_\infty \to X)$ la classe du
lemme \ref{lemma-gysin-at-infty}. Alors
$$
C \circ c_p(Q|_{X \times \{0\}}) =
C \circ (Z \to X)_* \circ c'_p(Q) = c_p(Q|_{W_\infty}) \circ C
$$
\end{lemma}

\begin{proof}
La première égalité vaut car $c_p(Q|_{X \times \{0\}}) =
(Z \to X)_* \circ c'_p(Q)$ d'après le lemme \ref{lemma-localized-chern-pre}.
Nous pouvons démontrer la seconde égalité classe de cycle par classe de cycle
(voir le lemme \ref{lemma-bivariant-zero}). Comme la construction des
classes bivariantes du lemme est compatible au changement de base,
nous pouvons supposer disposer d'un $\alpha \in \CH_k(X)$ et devoir montrer que
$C \cap (Z \to X)_*(c'_p(Q) \cap \alpha) =
c_p(Q|_{W_\infty}) \cap C \cap \alpha$. Remarquons que
\begin{align*}
C \cap (Z \to X)_*(c'_p(Q) \cap \alpha)
& =
C \cap (Z \to X)_* (E \to Z)_*(c'_p(Q|_E) \cap C \cap \alpha) \\
& =
C \cap (W_\infty \to X)_*(E \to W_\infty)_*(c'_p(Q|_E) \cap C \cap \alpha) \\
& =
C \cap (W_\infty \to X)_*(E \to W_\infty)_*(c'_p(Q|_E) \cap i_\infty^*\beta) \\
& =
C \cap (W_\infty \to X)_*(c_p(Q|_{W_\infty}) \cap i_\infty^*\beta) \\
& =
C \cap (W_\infty \to X)_*i_\infty^*(c_p(Q) \cap \beta) \\
& =
i_\infty^*(c_p(Q) \cap \beta) \\
& =
c_p(Q|_{W_\infty}) \cap i_\infty^*\beta \\
& =
c_p(Q|_{W_\infty}) \cap C \cap \alpha
\end{align*}
comme souhaité. Pour la première égalité, nous avons utilisé
$c'_p(Q) = (E \to Z)_* \circ c'_p(Q|_E) \circ C$, où $E \subset W_\infty$
est l'image inverse de $Z$ et $c'_p(Q|_E)$ la classe construite
dans le lemme \ref{lemma-silly}. La deuxième égalité exprime simplement
que $E \to Z \to X$ est égal à $E \to W_\infty \to X$.
Pour la troisième égalité, choisissons $\beta \in \CH_{k + 1}(W)$ dont la restriction à
$b^{-1}(\mathbf{A}^1_X)$ est l'image inverse plate de $\alpha$, de sorte que
$C \cap \alpha = i_\infty^*\beta$ par construction. La quatrième égalité est le
lemme \ref{lemma-silly-silly}. La cinquième vient du fait que
$c_p(Q)$ est une classe bivariante et commute donc à $i_\infty^*$.
La sixième égalité est le lemme \ref{lemma-homomorphism-pre}.
La septième utilise de nouveau le fait que $c_p(Q)$ est une classe bivariante.
La dernière vaut puisque $C \cap \alpha = i_\infty^*\beta$.
\end{proof}

\begin{lemma}
\label{lemma-homomorphism-commute}
Dans le lemme \ref{lemma-localized-chern-pre}, soit $Y \to X$ un morphisme
localement de type fini et soit $c \in A^*(Y \to X)$ une classe bivariante.
Alors
$$
P'_p(Q) \circ c = c \circ P'_p(Q)
\quad\text{resp.}\quad
c'_p(Q) \circ c = c \circ c'_p(Q)
$$
dans $A^*(Y \times_X Z \to X)$.
\end{lemma}

\begin{proof}
Soit $E \subset W_\infty$ l'image inverse de $Z$.
Rappelons que $P'_p(Q) = (E \to Z)_* \circ P'_p(Q|_E) \circ C$,
resp.\ $c'_p(Q) = (E \to Z)_* \circ c'_p(Q|_E) \circ C$,
où $C$ est comme dans le lemme \ref{lemma-gysin-at-infty} et
$P'_p(Q|_E)$, resp.\ $c'_p(Q|_E)$ sont comme dans le
lemme \ref{lemma-silly}.
D'après le lemme \ref{lemma-gysin-at-infty-commutes},
nous voyons que $C$ commute à $c$,
et le lemme \ref{lemma-silly-commutes} montre que
$P'_p(Q|_E)$, resp.\ $c'_p(Q|_E)$ commute à $c$.
Comme $c$ est une classe bivariante, elle commute par définition à l'image
directe propre par $E \to Z$. Cela achève la démonstration.
\end{proof}

\begin{lemma}
\label{lemma-localized-chern-pre-compose}
Dans le lemme \ref{lemma-localized-chern-pre}, supposons que $Q|_T$ soit nul. Dans
$A^*(Z \to X)$, nous avons
\begin{align*}
P'_1(Q) & = c'_1(Q), \\
P'_2(Q) & = c'_1(Q)^2 - 2c'_2(Q), \\
P'_3(Q) & = c'_1(Q)^3 - 3c'_1(Q)c'_2(Q) + 3c'_3(Q), \\
P'_4(Q) & = c'_1(Q)^4 - 4c'_1(Q)^2c'_2(Q) +
4c'_1(Q)c'_3(Q) + 2c'_2(Q)^2 - 4c'_4(Q),
\end{align*}
et ainsi de suite, la multiplication étant celle de la remarque \ref{remark-ring-loc-classes}.
\end{lemma}

\begin{proof}
L'énoncé a un sens car le faisceau nul est de rang $< 1$ et
les classes $c'_p(Q)$ sont donc définies pour tout $p \geq 1$.
Dans la démonstration du lemme \ref{lemma-localized-chern-pre}, nous avons construit
les classes $P'_p(Q)$ et $c'_p(Q)$ en utilisant la classe bivariante
$C \in A^0(W_\infty \to X)$ du lemme \ref{lemma-gysin-at-infty}
et les classes bivariantes
$P'_p(Q|_E)$ et $c'_p(Q|_E)$ du lemme \ref{lemma-silly} pour la restriction
$Q|_E$ de $Q$ à l'image inverse $E$ de $Z$ dans $W_\infty$.
Remarquons que le lemme \ref{lemma-silly-compose} donne la relation voulue
entre $P'_p(Q|_E)$ et $c'_p(Q|_E)$. Rappelons que
$$
P'_p(Q) = (E \to Z)_* \circ P'_p(Q|_E) \circ C
\quad\text{et}\quad
c'_p(Q) = (E \to Z)_* \circ c'_p(Q|_E) \circ C
$$
Pour achever la démonstration, il suffit de montrer que les multiplications définies
dans la remarque \ref{remark-ring-loc-classes} sur les classes $a_p = c'_p(Q)$
et sur les classes $b_p = c'_p(Q|_E)$ concordent :
$$
a_{p_1}a_{p_2} \ldots a_{p_r} =
(E \to Z)_* \circ b_{p_1}b_{p_2} \ldots b_{p_r} \circ C
$$
Nous omettons certains détails. Si $r = 1$, c'est vrai.
Pour $r > 1$, remarquons que, d'après la remarque \ref{remark-res-push}, la multiplication de la
remarque \ref{remark-ring-loc-classes} s'effectue
en insérant $(Z \to X)_*$, resp.\ $(E \to W_\infty)_*$ entre
les facteurs du produit
$a_{p_1}a_{p_2} \ldots a_{p_r}$, resp.\ $b_{p_1}b_{p_2} \ldots b_{p_r}$,
et en prenant les compositions comme classes bivariantes.
Or, d'après le lemme \ref{lemma-silly}, nous avons
$$
(E \to W_\infty)_* \circ b_{p_i} = c_{p_i}(Q|_{W_\infty})
$$
et, d'après le lemme \ref{lemma-homomorphism}, nous avons
$$
C \circ (Z \to X)_* \circ a_{p_i} = c_{p_i}(Q|_{W_\infty}) \circ C
$$
pour $i = 2, \ldots, r$. Un calcul
montre que le membre de gauche et le membre de droite de l'égalité voulue
se simplifient tous deux en
$$
(E \to Z)_* \circ c'_{p_1}(Q|_E) \circ
c_{p_2}(Q|_{W_\infty}) \circ \ldots \circ
c_{p_r}(Q|_{W_\infty}) \circ C
$$
ce qui achève la démonstration.
\end{proof}

\begin{lemma}
\label{lemma-localized-chern-pre-sum-c}
Dans le lemme \ref{lemma-localized-chern-pre}, supposons que $Q|_T$ soit isomorphe
à un $\mathcal{O}_T$-module localement libre de rang fini, de rang $< p$.
Supposons donné un autre objet parfait $Q' \in D(\mathcal{O}_W)$
dont les classes de Chern sont définies et tel que $Q'|_T$ soit isomorphe à un
$\mathcal{O}_T$-module localement libre de rang fini, de rang $< p'$, placé
en degré cohomologique $0$. Avec la notation de la
remarque \ref{remark-ring-loc-classes}, posons
$$
c^{(p)}(Q) = 1 + c_1(Q|_{X \times \{0\}}) + \ldots +
c_{p - 1}(Q|_{X \times \{0\}}) +
c'_{p}(Q) + c'_{p + 1}(Q) + \ldots
$$
dans $A^{(p)}(Z \to X)$, où $c'_i(Q)$, pour $i \geq p$, est comme dans le
lemme \ref{lemma-localized-chern-pre}. De même pour $c^{(p')}(Q')$ et
$c^{(p + p')}(Q \oplus Q')$.
Alors $c^{(p + p')}(Q \oplus Q') = c^{(p)}(Q)c^{(p')}(Q')$
dans $A^{(p + p')}(Z \to X)$.
\end{lemma}

\begin{proof}
Rappelons que l'image de $c'_i(Q)$ dans $A^p(X)$ est égale à
$c_i(Q|_{X \times \{0\}})$ pour $i \geq p$, et de même pour
$Q'$ et $Q \oplus Q'$ ; voir le lemme \ref{lemma-localized-chern-pre}.
L'égalité en degrés $< p + p'$ résulte donc de
l'additivité du lemme \ref{lemma-additivity-on-perfect}.

\medskip\noindent
Prenons $n \geq p + p'$.
Comme dans la démonstration du lemme \ref{lemma-localized-chern-pre},
notons $E \subset W_\infty$ l'image inverse de $Z$.
Remarquons que nous avons l'égalité
$$
c^{(p + p')}(Q|_E \oplus Q'|_E) =
c^{(p)}(Q|_E)c^{(p')}(Q'|_E)
$$
dans $A^{(p + p')}(E \to W_\infty)$ d'après le lemme \ref{lemma-silly-sum-c}.
Puisque, par construction,
$$
c'_p(Q \oplus Q') = (E \to Z)_* \circ c'_p(Q|_E \oplus Q'|_E) \circ C
$$
nous concluons qu'il suffit de montrer, pour tout $i + j = n$, que
$$
(E \to Z)_* \circ c^{(p)}_i(Q|_E)c^{(p')}_j(Q'|_E) \circ C
=
c^{(p)}_i(Q)c^{(p')}_j(Q')
$$
dans $A^n(Z \to X)$, où la multiplication est celle de la
remarque \ref{remark-ring-loc-classes} dans les deux membres. Il y a
trois cas, selon que $i \geq p$, $j \geq p'$, ou les deux.

\medskip\noindent
Supposons $i \geq p$ et $j \geq p'$. Dans ce cas, les produits sont
définis en insérant $(E \to W_\infty)_*$, resp.\ $(Z \to X)_*$ entre
les deux facteurs et en prenant les compositions comme classes bivariantes ; voir la
remarque \ref{remark-res-push}.
Autrement dit, nous devons montrer
$$
(E \to Z)_* \circ c'_i(Q|_E) \circ
(E \to W_\infty)_* \circ c'_j(Q'|_E) \circ C =
c'_i(Q) \circ (Z \to X)_* \circ c'_j(Q')
$$
D'après le lemme \ref{lemma-silly}, le membre de gauche est égal à
$$
(E \to Z)_* \circ c'_i(Q|_E) \circ c_j(Q'|_{W_\infty}) \circ C
$$
Puisque $c'_i(Q) = (E \to Z)_* \circ c'_i(Q|_E) \circ C$,
le membre de droite est égal à
$$
(E \to Z)_* \circ c'_i(Q|_E) \circ C \circ (Z \to X)_* \circ c'_j(Q')
$$
ce qui est immédiatement égal à l'expression ci-dessus
d'après le lemme \ref{lemma-homomorphism}.

\medskip\noindent
Supposons $i \geq p$ et $j < p$. En développant les produits
dans ce cas, nous devons montrer
$$
(E \to Z)_* \circ c'_i(Q|_E) \circ c_j(Q'|_{W_\infty}) \circ C =
c'_i(Q) \circ c_j(Q'|_{X \times \{0\}})
$$
En utilisant de nouveau $c'_i(Q) = (E \to Z)_* \circ c'_i(Q|_E) \circ C$,
nous voyons qu'il suffit de montrer $c_j(Q'|_{W_\infty}) \circ C =
C \circ c_j(Q'|_{X \times \{0\}})$, ce qui fait partie du
lemme \ref{lemma-homomorphism}.

\medskip\noindent
Supposons $i < p$ et $j \geq p'$. En développant les produits
dans ce cas, nous devons montrer
$$
(E \to Z)_* \circ c_i(Q|_E) \circ c'_j(Q'|_E) \circ C =
c_i(Q|_{Z \times \{0\}}) \circ c'_j(Q')
$$
Cependant, puisque $c'_j(Q|_E)$ et $c'_j(Q')$ sont des
classes bivariantes, elles commutent à l'intersection avec les classes de Chern
(lemme \ref{lemma-cap-commutative-chern}). Il suffit donc de démontrer
$$
(E \to Z)_* \circ c'_j(Q'|_E) \circ c_i(Q|_{W_\infty}) \circ C =
c'_j(Q') \circ c_i(Q|_{X \times \{0\}})
$$
ce qui nous ramène au cas traité au paragraphe précédent.
\end{proof}

\begin{lemma}
\label{lemma-localized-chern-pre-sum-P}
Dans le lemme \ref{lemma-localized-chern-pre}, supposons que $Q|_T$ soit nul.
Supposons donné un autre objet parfait $Q' \in D(\mathcal{O}_W)$
dont les classes de Chern sont définies et dont la restriction $Q'|_T$ est nulle.
Dans ce cas, les classes
$P'_p(Q), P'_p(Q'), P'_p(Q \oplus Q') \in A^p(Z \to X)$
construites dans le lemme \ref{lemma-localized-chern-pre}
satisfont $P'_p(Q \oplus Q') = P'_p(Q) + P'_p(Q')$.
\end{lemma}

\begin{proof}
Cela résulte immédiatement de la construction de ces
classes et du lemme \ref{lemma-silly-sum-P}.
\end{proof}









 
\section{Classes de Chern localisées}
\label{section-localized-chern}

\noindent
Aperçu de la construction. Soit $F$ un corps,
soit $X$ une variété sur $F$, soit $E$ un objet parfait de
$D(\mathcal{O}_X)$, et soit $Z \subset X$ un sous-schéma fermé
tel que $E|_{X \setminus Z} = 0$. Nous voulons alors construire des éléments
$$
c_p(Z \to X, E) \in A^p(Z \to X)
$$
Nous le ferons en construisant un diagramme
$$
\xymatrix{
W \ar[d]_f \ar[r]_q & X \\
\mathbf{P}^1_F
}
$$
et un objet parfait $Q$ de $D(\mathcal{O}_W)$ tels que
\begin{enumerate}
\item $f$ est plat, et $f$, $q$ sont propres ; pour $t \in \mathbf{P}^1_F$,
notons $W_t$ la fibre de $f$,
$q_t : W_t \to X$ la restriction de $q$, et $Q_t = Q|_{W_t}$ ;
\item $q_t : W_t \to X$ est un isomorphisme et $Q_t = q_t^*E$
pour $t \in \mathbf{A}^1_F$ ;
\item $q_\infty : W_\infty \to X$ est un isomorphisme au-dessus de $X \setminus Z$ ;
\item si $T \subset W_\infty$ est l'adhérence de
$q_\infty^{-1}(X \setminus Z)$, alors $Q_\infty|_T$ est nul.
\end{enumerate}
L'idée est de considérer ceci comme une famille $\{(W_t, Q_t)\}$
paramétrée par $t \in \mathbf{P}^1$.
Pour $t \not = \infty$, nous voyons que $c_p(Q_t)$ est simplement $c_p(E)$
sur les groupes de Chow de $Q_t = X$. Mais pour $t = \infty$, nous voyons que
$c_p(Q_\infty)$ envoie les classes sur $Q_\infty$ sur des classes à support dans
$E = q_\infty^{-1}(Z)$ puisque $Q_\infty|_T = 0$.
Nous considérons $E$ comme le lieu exceptionnel de $q_\infty : W_\infty \to X$.
Puisque tout $\alpha \in \CH_*(X)$ donne naissance à une « famille »
de cycles $\alpha_t \in \CH_*(W_t)$, il est naturel de définir
$c_p(Z \to X, E) \cap \alpha$ comme
l'image directe $(E \to Z)_*(c_p(Q_\infty) \cap \alpha_\infty)$.

\medskip\noindent
Pour réaliser cette construction, il y a deux ingrédients principaux : (1) la construction de
$W$ et de $Q$ est une sorte de construction algébrique du graphe de Macpherson ; elle
est effectuée dans Compléments sur la platitude, section \ref{flat-section-blowup-complexes-III}.
(2) La construction de la classe elle-même à partir de $W$ et $Q$ est faite dans la
section \ref{section-preparation-localized-chern-II}, en s'appuyant sur les
sections \ref{section-gysin-at-infty} et
\ref{section-preparation-localized-chern}.

\begin{situation}
\label{situation-loc-chern}
Soit $(S, \delta)$ comme dans la situation \ref{situation-setup}. Soit $X$ un schéma
localement de type fini sur $S$. Soit $i : Z \to X$ une immersion fermée.
Soit $E \in D(\mathcal{O}_X)$ un objet. Soit $p \geq 0$. Supposons que
\begin{enumerate}
\item $E$ soit un objet parfait de $D(\mathcal{O}_X)$ ;
\item la restriction $E|_{X \setminus Z}$ soit nulle, resp.\ isomorphe à un
$\mathcal{O}_{X \setminus Z}$-module localement libre de rang fini, de rang $< p$,
placé en degré cohomologique $0$ ; et
\item au moins une\footnote{Veuillez ignorer cette condition technique en
première lecture ; voir la discussion de la remarque \ref{remark-loc-chern}.}
des assertions suivantes soit vraie :
(a) $X$ est quasi-compact ;
(b) les composantes irréductibles de $X$ sont quasi-compactes ;
(c) il existe un complexe localement borné de
$\mathcal{O}_X$-modules localement libres de rang fini représentant $E$ ; ou
(d) il existe un morphisme $X \to X'$ de schémas localement de type fini
sur $S$ tel que $E$ soit l'image inverse d'un objet parfait sur $X'$ et que
les composantes irréductibles de $X'$ soient quasi-compactes.
\end{enumerate}
\end{situation}

\begin{lemma}
\label{lemma-independent-loc-chern}
Dans la situation \ref{situation-loc-chern}, il existe une classe bivariante canonique
$$
P_p(Z \to X, E) \in A^p(Z \to X),
\quad\text{resp.}\quad
c_p(Z \to X, E) \in A^p(Z \to X)
$$
ayant la propriété suivante :
\begin{equation}
\label{equation-defining-property-localized-classes}
i_* \circ P_p(Z \to X, E) = P_p(E),
\quad\text{resp.}\quad
i_* \circ c_p(Z \to X, E) = c_p(E)
\end{equation}
comme classes bivariantes sur $X$ (où $\circ$ est comme dans le
lemme \ref{lemma-push-proper-bivariant}).
\end{lemma}

\begin{proof}
La construction de ces classes bivariantes est la suivante. Soient
$$
b : W \longrightarrow \mathbf{P}^1_X
\quad\text{et}\quad
T \longrightarrow W_\infty
\quad\text{et}\quad
Q
$$
l'éclatement, l'objet parfait $Q$ de $D(\mathcal{O}_W)$
et l'immersion fermée construits dans
Compléments sur la platitude, section \ref{flat-section-blowup-complexes-III}
et le lemme \ref{flat-lemma-graph-construction}.
Soit $T' \subset T$ le sous-schéma ouvert et fermé tel que
$Q|_{T'}$ soit nul, resp.\ isomorphe à un
$\mathcal{O}_{T'}$-module localement libre de rang fini, de rang $< p$,
placé en degré cohomologique $0$. D'après la condition (2) de la
situation \ref{situation-loc-chern}, les morphismes
$$
T' \to T \to W_\infty \to X
$$
sont tous des isomorphismes de schémas au-dessus du sous-schéma ouvert $X \setminus Z$ de $X$.
Nous vérifions ci-dessous que les classes de Chern de $Q$ sont définies.
Rappelant que $Q|_{X \times \{0\}} \cong E$ par construction, nous concluons
que la classe bivariante construite dans le lemme \ref{lemma-localized-chern-pre}
à partir de $W, b, Q, T'$ nous donne les classes
$$
P_p(Z \to X, E) = P'_p(Q) \in A^p(Z \to X)
$$
et
$$
c_p(Z \to X, E) = c'_p(Q) \in A^p(Z \to X)
$$
satisfaisant à (\ref{equation-defining-property-localized-classes}).

\medskip\noindent
Dans ce paragraphe, nous démontrons que les classes de Chern de $Q$ sont définies
(définition \ref{definition-defined-on-perfect}) ; nous suggérons au lecteur de sauter
ce passage. Si l'hypothèse (3)(a) ou (3)(b) de la situation \ref{situation-loc-chern}
est satisfaite, c'est-à-dire si les composantes irréductibles de $X$ sont quasi-compactes, alors il en
va de même pour $W$ (car $W \to X$ est propre). Nous concluons donc que
les classes de Chern de tout objet parfait de $D(\mathcal{O}_W)$ sont définies d'après le
lemme \ref{lemma-chern-classes-defined}. Si (3)(c) est satisfaite, c'est-à-dire si
$E$ peut être représenté par un complexe localement borné de modules localement
libres de rang fini, alors l'objet $Q$ peut être représenté par un complexe localement borné
de $\mathcal{O}_W$-modules localement libres de rang fini d'après la partie (5) de
Compléments sur la platitude, lemme \ref{flat-lemma-graph-construction}. Les
classes de Chern de $Q$ sont donc définies. Enfin, supposons (3)(d) satisfaite, c'est-à-dire
qu'il existe un morphisme $X \to X'$ de schémas localement de type fini
sur $S$ tel que $E$ soit l'image inverse d'un objet parfait $E'$ sur $X'$ et que
les composantes irréductibles de $X'$ soient quasi-compactes.
Soient $b' : W' \to \mathbf{P}^1_{X'}$ et $Q' \in D(\mathcal{O}_{W'})$
le morphisme et l'objet parfait construits comme
dans Compléments sur la platitude, section \ref{flat-section-blowup-complexes-III},
à partir du triplet
$(\mathbf{P}^1_{X'}, (\mathbf{P}^1_{X'})_\infty, L(p')^*E')$.
La discussion ci-dessus montre que les classes de Chern de $Q'$
sont définies. Puisque $b$ et $b'$ ont été construits par une application de
Compléments sur la platitude, lemme \ref{flat-lemma-complex-and-divisor-blowup-pre},
il résulte de
Compléments sur la platitude, lemme \ref{flat-lemma-complex-and-divisor-blowup-base-change},
qu'il existe un morphisme $W \to W'$ tel que
$Q = L(W \to W')^*Q'$. Le
lemme \ref{lemma-chern-classes-defined}
montre alors que les classes de Chern de $Q$ sont définies.
\end{proof}

\begin{definition}
\label{definition-localized-chern}
Avec $(S, \delta)$, $X$, $E \in D(\mathcal{O}_X)$ et $i : Z \to X$ comme dans la
situation \ref{situation-loc-chern}.
\begin{enumerate}
\item Si la restriction $E|_{X \setminus Z}$ est nulle, alors, pour tout
$p \geq 0$, nous définissons
$$
P_p(Z \to X, E) \in A^p(Z \to X)
$$
par la construction du lemme \ref{lemma-independent-loc-chern},
et nous définissons le {\it caractère de Chern localisé} par la formule
$$
ch(Z \to X, E) =
\sum\nolimits_{p = 0, 1, 2, \ldots} \frac{P_p(Z \to X, E)}{p!}
\quad\text{dans}\quad \prod\nolimits_{p \geq 0} A^p(Z \to X) \otimes \mathbf{Q}
$$
\item Si la restriction $E|_{X \setminus Z}$ est isomorphe à un
$\mathcal{O}_{X \setminus Z}$-module localement libre de rang fini, de rang $< p$,
placé en degré cohomologique $0$, nous définissons alors la
{\it $p$-ième classe de Chern localisée} $c_p(Z \to X, E)$ par la construction
du lemme \ref{lemma-independent-loc-chern}.
\end{enumerate}
\end{definition}

\noindent
Dans la situation de la définition, supposons que $E|_{X \setminus Z}$ soit nul.
Pour être précis, nous avons alors l'égalité
$$
i_* \circ ch(Z \to X, E) = ch(E)
$$
dans $A^*(X) \otimes \mathbf{Q}$, car nous avons démontré ci-dessus
l'égalité (\ref{equation-defining-property-localized-classes}).

\medskip\noindent
Voici une vérification de cohérence importante.

\begin{lemma}
\label{lemma-base-change-loc-chern}
Dans la situation \ref{situation-loc-chern},
soit $f : X' \to X$ un morphisme de schémas localement de type fini.
Notons $E' = f^*E$ et $Z' = f^{-1}(Z)$. Alors la classe bivariante
de la définition \ref{definition-localized-chern}
$$
P_p(Z' \to X', E') \in A^p(Z' \to X'),
\quad\text{resp.}\quad
c_p(Z' \to X', E') \in A^p(Z' \to X')
$$
construite comme dans le lemme \ref{lemma-independent-loc-chern}
à partir de $X', Z', E'$ est la restriction
(remarque \ref{remark-restriction-bivariant}) de la
classe bivariante $P_p(Z \to X, E) \in A^p(Z \to X)$,
resp.\ $c_p(Z \to X, E) \in A^p(Z \to X)$.
\end{lemma}

\begin{proof}
Notons $p : \mathbf{P}^1_X \to X$ et $p' : \mathbf{P}^1_{X'} \to X'$
les morphismes structuraux.
Rappelons que $b : W \to \mathbf{P}^1_X$ et $b' : W' \to \mathbf{P}^1_{X'}$
sont les morphismes construits à partir des triplets
$(\mathbf{P}^1_X, (\mathbf{P}^1_X)\infty, p^*E)$ et
$(\mathbf{P}^1_{X'}, (\mathbf{P}^1_{X'})\infty, (p')^*E')$
dans Compléments sur la platitude, lemme \ref{flat-lemma-complex-and-divisor-blowup-pre}.
En outre, $Q = L\eta_{\mathcal{I}_\infty}p^*E$ et
$Q = L\eta_{\mathcal{I}'_\infty}(p')^*E'$
où
$\mathcal{I}_\infty \subset \mathcal{O}_W$ est le
faisceau d'idéaux de $W_\infty$ et
$\mathcal{I}'_\infty \subset \mathcal{O}_{W'}$ est le
faisceau d'idéaux de $W'_\infty$.
Ensuite, $h : \mathbf{P}^1_{X'} \to \mathbf{P}^1_X$ est un morphisme de
schémas tel que l'image inverse du diviseur de Cartier effectif
$(\mathbf{P}^1_X)_\infty$ soit le diviseur de Cartier effectif
$(\mathbf{P}^1_{X'})_\infty$ et que $h^*p^*E = (p')^*E'$.
D'après Compléments sur la platitude, lemme
\ref{flat-lemma-complex-and-divisor-blowup-base-change},
nous obtenons un diagramme commutatif
$$
\xymatrix{
W' \ar[rd]_{b'} \ar[r]_-g &
\mathbf{P}^1_{X'} \times_{\mathbf{P}^1_X} W \ar[d]_r \ar[r]_-q &
W \ar[d]^b \\
&
\mathbf{P}^1_{X'} \ar[r] &
\mathbf{P}^1_X
}
$$
tel que $W'$ soit le « transformé strict » de $\mathbf{P}^1_{X'}$
par rapport à $b$ et que $Q' = (q \circ g)^*Q$.
Rappelons maintenant que $P_p(Z \to X, E) = P'_p(Q)$,
resp.\ $c_p(Z \to X, E) = c'_p(Q)$, où $P'_p(Q)$, resp.\ $c'_p(Q)$,
sont construits dans le lemme \ref{lemma-localized-chern-pre}
à l'aide de $b, Q, T'$, où $T'$ est un sous-schéma fermé $T' \subset W_\infty$
possédant les deux propriétés suivantes :
(a) $T'$ contient tous les points de $W_\infty$ situés au-dessus de $X \setminus Z$,
et (b) $Q|_{T'}$ est nul, resp.\ isomorphe à un module localement libre de rang fini,
de rang $< p$ et placé en degré $0$.
Dans la construction du lemme \ref{lemma-localized-chern-pre},
nous avons choisi un sous-schéma fermé particulier $T'$ possédant les propriétés (a) et (b),
mais le choix précis de $T'$ est sans importance ; voir le
lemme \ref{lemma-localized-chern-pre-independent}.

\medskip\noindent
Ensuite, d'après le lemme \ref{lemma-base-change-localized-chern-pre},
la restriction de la classe bivariante $P_p(Z \to X, E) = P'_p(Q)$,
resp.\ $c_p(Z \to X, E) = c_p(Q')$,
à $X'$ correspond à la classe $P'_p(q^*Q)$, resp.\ $c'_p(q^*Q)$,
construite comme dans le lemme \ref{lemma-localized-chern-pre} à l'aide de
$r : \mathbf{P}^1_{X'} \times_{\mathbf{P}^1_X} W \to \mathbf{P}^1_{X'}$,
du complexe $q^*Q$ et de l'image inverse $q^{-1}(T')$.

\medskip\noindent
Maintenant, d'après la seconde assertion du
lemme \ref{lemma-localized-chern-pre-independent},
nous avons $P'_p(Q') = P'_p(q^*Q)$, resp.\ $c'_p(q^*Q) = c'_p(Q')$.
Puisque $P_p(Z' \to X', E') = P'_p(Q')$, resp.\ $c_p(Z' \to X', E') = c'_p(Q')$,
nous concluons que le lemme est vrai.
\end{proof}

\begin{remark}
\label{remark-loc-chern}
Dans la situation \ref{situation-loc-chern}, il aurait été plus naturel
de remplacer l'hypothèse (3) par l'hypothèse : « les classes de Chern
de $E$ sont définies ». En fait, en combinant les
lemmes \ref{lemma-independent-loc-chern} et
\ref{lemma-base-change-loc-chern}
avec le lemme \ref{lemma-envelope-bivariant},
il est facile d'étendre la définition à ce cas légèrement plus général.
Si cela nous est un jour nécessaire, nous le ferons ici.
\end{remark}

\begin{lemma}
\label{lemma-loc-chern-after-pushforward}
Dans la situation \ref{situation-loc-chern}, nous avons
$$
P_p(Z \to X, E) \cap i_*\alpha = P_p(E|_Z) \cap \alpha,
\quad\text{resp.}\quad
c_p(Z \to X, E) \cap i_*\alpha = c_p(E|_Z) \cap \alpha
$$
dans $\CH_*(Z)$ pour tout $\alpha \in \CH_*(Z)$.
\end{lemma}

\begin{proof}
Nous démontrons seulement la seconde égalité et omettons la démonstration de la première.
Puisque $c_p(Z \to X, E)$ est une classe bivariante et que le changement de
base de $Z \to X$ par $Z \to X$ est $\text{id} : Z \to Z$, nous avons
$c_p(Z \to X, E) \cap i_*\alpha = c_p(Z \to X, E) \cap \alpha$.
D'après le lemme \ref{lemma-base-change-loc-chern}, la restriction de
$c_p(Z \to X, E)$ à $Z$ (!) est la classe de Chern localisée pour
$\text{id} : Z \to Z$ et $E|_Z$. Le résultat résulte donc de
(\ref{equation-defining-property-localized-classes}) avec $X = Z$.
\end{proof}

\begin{lemma}
\label{lemma-loc-chern-disjoint}
Dans la situation \ref{situation-loc-chern},
si $\alpha \in \CH_k(X)$ a un support disjoint de $Z$, alors
$P_p(Z \to X, E) \cap \alpha = 0$, resp.\ $c_p(Z \to X, E) \cap \alpha = 0$.
\end{lemma}

\begin{proof}
Cela résulte immédiatement de la construction des classes de Chern localisées.
Cela résulte aussi du fait que nous pouvons calculer
$c_p(Z \to X, E) \cap \alpha$ en restreignant d'abord $c_p(Z \to X, E)$ au
support de $\alpha$, puis en utilisant
le lemme \ref{lemma-base-change-loc-chern}
pour voir que cette restriction est nulle.
\end{proof}

\begin{lemma}
\label{lemma-loc-chern-shrink-Z}
Dans la situation \ref{situation-loc-chern},
supposons que $Z \subset Z' \subset X$, où $Z'$ est un sous-schéma fermé de $X$.
Alors
$P_p(Z' \to X, E) = (Z \to Z')_* \circ P_p(Z \to X, E)$,
resp.\ $c_p(Z' \to X, E) = (Z \to Z')_* \circ c_p(Z \to X, E)$
(avec $\circ$ comme dans le lemme \ref{lemma-push-proper-bivariant}).
\end{lemma}

\begin{proof}
La construction de $P_p(Z' \to X, E)$,
resp.\ $c_p(Z' \to X, E)$ dans le lemme \ref{lemma-independent-loc-chern}
utilise exactement le même morphisme
$b : W \to \mathbf{P}^1_X$ et le même objet parfait $Q$ de $D(\mathcal{O}_W)$.
Nous pouvons alors utiliser le lemme \ref{lemma-silly-shrink} pour conclure.
Certains détails sont omis.
\end{proof}

\begin{lemma}
\label{lemma-loc-chern-agree}
Dans le lemme \ref{lemma-silly}, supposons que $E_2$ soit la restriction d'un objet parfait
$E \in D(\mathcal{O}_X)$ dont la restriction à $X_1$ est nulle,
resp.\ est isomorphe à un $\mathcal{O}_{X_1}$-module localement libre de rang fini
de rang $< p$ placé en degré cohomologique $0$. Alors la classe
$P'_p(E_2)$, resp.\ $c'_p(E_2)$ du lemme \ref{lemma-silly} coïncide avec
$P_p(X_2 \to X, E)$, resp.\ $c_p(X_2 \to X, E)$ de la
définition \ref{definition-localized-chern}, pourvu que $E$ satisfasse
l'hypothèse (3) de la situation \ref{situation-loc-chern}.
\end{lemma}

\begin{proof}
Les hypothèses sur $E$ impliquent qu'il existe un ouvert $U \subset X$
contenant $X_1$ tel que $E|_U$ soit nul, resp.\ soit isomorphe à un
$\mathcal{O}_U$-module localement libre de rang $< p$. Voir Compléments d'algèbre, lemme
\ref{more-algebra-lemma-lift-perfect-from-residue-field}.
Soit $Z \subset X$ le complémentaire de $U$ dans $X$, muni de
la structure de sous-schéma fermé induite réduite. Alors
$P_p(X_2 \to X, E) = (Z \to X_2)_* \circ P_p(Z \to X, E)$,
resp.\ $c_p(X_2 \to X, E) = (Z \to X_2)_* \circ c_p(Z \to X, E)$
par le lemme \ref{lemma-loc-chern-shrink-Z}.
Nous pouvons maintenant démontrer que $P_p(X_2 \to X, E)$, resp.\ $c_p(X_2 \to X, E)$
satisfait la caractérisation de $P'_p(E_2)$, resp.\ $c'_p(E_2)$
donnée dans le lemme \ref{lemma-silly}. En effet, par la relation
$P_p(X_2 \to X, E) = (Z \to X_2)_* \circ P_p(Z \to X, E)$,
resp.\ $c_p(X_2 \to X, E) = (Z \to X_2)_* \circ c_p(Z \to X, E)$
que nous venons de démontrer et le fait que $X_1 \cap Z = \emptyset$,
la composée $P_p(X_2 \to X, E) \circ i_{1, *}$,
resp.\ $c_p(X_2 \to X, E) \circ i_{1, *}$ est nulle
par le lemme \ref{lemma-loc-chern-disjoint}.
D'autre part,
$P_p(X_2 \to X, E) \circ i_{2, *} = P_p(E_2)$,
resp.\ $c_p(X_2 \to X, E) \circ i_{2, *} = c_p(E_2)$
par le lemme \ref{lemma-loc-chern-after-pushforward}.
\end{proof}





\section{Deux lemmes techniques}
\label{section-tools-loc-chern}

\noindent
Dans cette section, nous développons quelques outils supplémentaires afin de pouvoir travailler
plus commodément avec les classes de Chern localisées. Le lemme suivant
est une version plus précise d'un fait que nous avons déjà rencontré dans
les démonstrations des lemmes \ref{lemma-localized-chern-pre-compose} et
\ref{lemma-localized-chern-pre-sum-c}.

\begin{lemma}
\label{lemma-homomorphism-final}
Soit $(S, \delta)$ comme dans la situation \ref{situation-setup}. Soit $X$
localement de type fini sur $S$. Soit $b : W \longrightarrow \mathbf{P}^1_X$
un morphisme propre de schémas. Soit $n \geq 1$. Pour $i = 1, \ldots, n$,
soit $Z_i \subset X$ un sous-schéma fermé, soit $Q_i \in D(\mathcal{O}_W)$
un objet parfait, soit $p_i \geq 0$ un entier et soient
$T_i \subset W_\infty$, $i = 1, \ldots, n$, des sous-schémas fermés.
Notons $W_i = b^{-1}(\mathbf{P}^1_{Z_i})$. Supposons que
\begin{enumerate}
\item pour $i = 1, \ldots, n$, l'hypothèse du
lemme \ref{lemma-localized-chern-pre} soit satisfaite pour
$b, Z_i, Q_i, T_i, p_i$,
\item $Q_i|_{W \setminus W_i}$ soit nul, resp.\ soit isomorphe à un module
localement libre de rang $< p_i$ placé en degré cohomologique $0$,
\item $Q_i$ sur $W$ satisfasse
l'hypothèse (3) de la situation \ref{situation-loc-chern}.
\end{enumerate}
Alors $P'_{p_n}(Q_n) \circ \ldots \circ P'_{p_1}(Q_1)$ est égal à
$$
(W_{n, \infty} \cap \ldots \cap W_{1, \infty} \to
Z_n \cap \ldots \cap Z_1)_* \circ
P'_{p_n}(Q_n|_{W_{n, \infty}}) \circ \ldots \circ P'_{p_1}(Q_1|_{W_{1, \infty}})
\circ C
$$
dans $A^{p_n + \ldots + p_1}(Z_n \cap \ldots \cap Z_1 \to X)$,
resp.\ $c'_{p_n}(Q_n) \circ \ldots \circ c'_{p_1}(Q_1)$ est égal à
$$
(W_{n, \infty} \cap \ldots \cap W_{1, \infty} \to
Z_n \cap \ldots \cap Z_1)_* \circ
c'_{p_n}(Q_n|_{W_{n, \infty}}) \circ \ldots \circ c'_{p_1}(Q_1|_{W_{1, \infty}})
\circ C
$$
dans $A^{p_n + \ldots + p_1}(Z_n \cap \ldots \cap Z_1 \to X)$.
\end{lemma}

\begin{proof}
Démontrons l'assertion sur les classes de Chern par récurrence sur $n$ ;
l'assertion sur $P_p(-)$ se démontre exactement de la même manière.
Le cas $n = 1$ est la construction de $c'_{p_1}(Q_1)$, car
$W_{1, \infty}$ est l'image inverse de $Z_1$ dans $W_\infty$.
Pour $n > 1$, nous avons par récurrence que
$c'_{p_n}(Q_n) \circ \ldots \circ c'_{p_1}(Q_1)$ est égal à
$$
c'_{p_n}(Q_n) \circ
(W_{n - 1, \infty} \cap \ldots \cap W_{1, \infty} \to
Z_{n - 1} \cap \ldots \cap Z_1)_* \circ
c'_{p_{n - 1}}(Q_{n - 1}|_{W_{n - 1}, \infty}) \circ \ldots \circ
c'_{p_1}(Q_1|_{W_{1, \infty}})
\circ C
$$
Par le lemme \ref{lemma-base-change-localized-chern-pre}, la restriction de
$c'_{p_n}(Q_n)$ à $Z_{n - 1} \cap \ldots \cap Z_1$ est calculée par
le sous-ensemble fermé $Z_n \cap \ldots \cap Z_1$, le morphisme
$b' : W_{n - 1} \cap \ldots \cap W_1 \to
\mathbf{P}^1_{Z_{n - 1} \cap \ldots \cap Z_1}$
et la restriction de $Q_n$ à $W_{n - 1} \cap \ldots \cap W_1$.
Observons que $(b')^{-1}(Z_n) = W_n \cap \ldots \cap W_1$
et que $(W_n \cap \ldots \cap W_1)_\infty =
W_{n, \infty} \cap \ldots \cap W_{1, \infty}$.
Notons $C_{n - 1} \in A^0(W_{n - 1, \infty} \cap \ldots \cap W_{1, \infty} \to
Z_{n - 1} \cap \ldots \cap Z_1)$ la classe du lemme \ref{lemma-gysin-at-infty}.
Nous concluons que la restriction de $c'_{p_n}(Q_n)$ à
$Z_{n - 1} \cap \ldots \cap Z_1$ est
\begin{align*}
&
(W_{n, \infty} \cap \ldots \cap W_{1, \infty} \to Z_n \cap \ldots \cap Z_1)_*
\circ
c'_{p_n}(Q_n|_{(W_n \cap \ldots \cap W_1)_\infty})
\circ
C_{n - 1} \\
& =
(W_{n, \infty} \cap \ldots \cap W_{1, \infty} \to Z_n \cap \ldots \cap Z_1)_*
\circ
c'_{p_n}(Q_n|_{W_{n, \infty}})
\circ
C_{n - 1}
\end{align*}
où l'égalité résulte du lemme \ref{lemma-base-change-silly}
(nous omettons d'écrire la restriction à droite). Ainsi, l'expression ci-dessus devient
\begin{align*}
(W_{n, \infty} \cap \ldots \cap W_{1, \infty} \to
Z_n \cap \ldots \cap Z_1)_* \circ
c'_{p_n}(Q_n|_{W_n, \infty}) \circ \\
C_{n - 1} \circ
(W_{n - 1, \infty} \cap \ldots \cap W_{1, \infty} \to
Z_{n - 1} \cap \ldots \cap Z_1)_* \\
\circ
c'_{p_{n - 1}}(Q_{n - 1}|_{W_{n - 1}, \infty}) \circ \ldots \circ
c'_{p_1}(Q_1|_{W_{1, \infty}})
\circ C
\end{align*}
Par le lemme \ref{lemma-homomorphism-pre},
nous savons que la composée
$C_{n - 1} \circ (W_{n - 1, \infty} \cap \ldots \cap W_{1, \infty} \to
Z_{n - 1} \cap \ldots \cap Z_1)_*$
est l'identité sur les éléments de l'image de l'homomorphisme de Gysin
$$
(W_{n - 1, \infty} \cap \ldots \cap W_{1, \infty} \to
W_{n - 1} \cap \ldots \cap W_1)^*
$$
Il suffit donc de montrer que tout élément de l'image de
$c'_{p_{n - 1}}(Q_{n - 1}|_{W_{n - 1}, \infty}) \circ \ldots \circ
c'_{p_1}(Q_1|_{W_{1, \infty}}) \circ C$
est dans l'image de l'homomorphisme de Gysin. Nous pouvons écrire
$$
c'_{p_i}(Q_i|_{W_{i, \infty}}) = \text{restriction de } c_{p_i}(W_i \to W, Q_i)
\text{ à } W_{i, \infty}
$$
par le lemme \ref{lemma-loc-chern-agree} et les hypothèses (2) et (3) sur $Q_i$
dans l'énoncé du lemme. Ainsi, si $\beta \in \CH_{k + 1}(W)$
se restreint à l'image inverse plate de $\alpha$ sur $b^{-1}(\mathbf{A}^1_X)$,
alors
\begin{align*}
& c'_{p_{n - 1}}(Q_{n - 1}|_{W_{n - 1}, \infty}) \cap \ldots \cap
c'_{p_1}(Q_1|_{W_{1, \infty}})
\cap C \cap \alpha \\
& =
c'_{p_{n - 1}}(Q_{n - 1}|_{W_{n - 1}, \infty}) \cap \ldots \cap
c'_{p_1}(Q_1|_{W_{1, \infty}})
\cap i_\infty^* \beta \\
& =
c_{p_{n - 1}}(W_{n - 1} \to W, Q_{n - 1}) \cap \ldots \cap
c_{p_{n - 1}}(W_1 \to W, Q_1) \cap i_\infty^* \beta \\
& =
(W_{n - 1, \infty} \cap \ldots \cap W_{1, \infty} \to
W_{n - 1} \cap \ldots \cap W_1)^*
\left(c_{p_{n - 1}}(W_{n - 1} \to W, Q_{n - 1}) \cap \ldots \cap
c_{p_1}(W_1 \to W, Q_1) \cap \beta\right)
\end{align*}
comme désiré. En effet, pour la dernière égalité, nous utilisons le fait que
$c_{p_i}(W_i \to W, Q_i)$ est une classe bivariante et donc
commute à $i_\infty^*$ par définition.
\end{proof}

\noindent
Le lemme suivant nous donne une flexibilité considérable
si nous voulons calculer les classes de Chern localisées d'un complexe.

\begin{lemma}
\label{lemma-independent-loc-chern-bQ}
Supposons que $(S, \delta), X, Z, b : W \to \mathbf{P}^1_X, Q, T, p$
satisfassent les hypothèses du lemme \ref{lemma-localized-chern-pre}.
Soit $F \in D(\mathcal{O}_X)$ un objet parfait tel que
\begin{enumerate}
\item la restriction de $Q$ à $b^{-1}(\mathbf{A}^1_X)$ soit
isomorphe à l'image inverse de $F$,
\item $F|_{X \setminus Z}$ soit nul, resp.\ soit isomorphe à un
$\mathcal{O}_{X \setminus Z}$-module localement libre de rang $< p$
placé en degré cohomologique $0$, et
\item $Q$ sur $W$ et $F$ sur $X$ satisfassent l'hypothèse (3) de la
situation \ref{situation-loc-chern}.
\end{enumerate}
Alors la classe $P'_p(Q)$, resp.\ $c'_p(Q)$ dans $A^p(Z \to X)$ construite
dans le lemme \ref{lemma-localized-chern-pre}
est égale à $P_p(Z \to X, F)$, resp.\ $c_p(Z \to X, F)$
de la définition \ref{definition-localized-chern}.
\end{lemma}

\begin{proof}
Les hypothèses sont préservées par changement de base par un morphisme
$X' \to X$ localement de type fini. Il suffit donc de montrer que
$P_p(Z \to X, F) \cap \alpha = P'_p(Q) \cap \alpha$,
resp.\ $c_p(Z \to X, F) \cap \alpha = c'_p(Q) \cap \alpha$
pour tout $\alpha \in \CH_k(X)$. Choisissons $\beta \in \CH_{k + 1}(W)$
dont la restriction à $b^{-1}(\mathbf{A}^1_X)$ est égale à
l'image inverse plate de $\alpha$, comme dans la construction de
$C$ dans le lemme \ref{lemma-gysin-at-infty}.
Notons $W' = b^{-1}(Z)$ et notons
$E = W'_\infty \subset W_\infty$ l'image inverse de $Z$
par $W_\infty \to X$.
Le lemme résulte de
la suite d'égalités suivante (le cas de $P_p$ est analogue)
\begin{align*}
c'_p(Q) \cap \alpha
& =
(E \to Z)_*(c'_p(Q|_E) \cap i_\infty^*\beta) \\
& =
(E \to Z)_*(c_p(E \to W_\infty, Q|_{W_\infty}) \cap i_\infty^*\beta) \\
& =
(W'_\infty \to Z)_*(c_p(W' \to W, Q) \cap i_\infty^*\beta) \\
& =
(W'_\infty \to Z)_*((i'_\infty)^*(c_p(W' \to W, Q) \cap \beta)) \\
& =
(W'_\infty \to Z)_*((i'_\infty)^*(c_p(Z' \to X, F) \cap \beta)) \\
& =
(W'_0 \to Z)_*((i'_0)^*(c_p(Z' \to X, F) \cap \beta)) \\
& =
(W'_0 \to Z)_*(c_p(Z' \to X, F) \cap i_0^*\beta)) \\
& =
c_p(Z \to X, F) \cap \alpha
\end{align*}
La première égalité est la construction de $c'_p(Q)$ dans le
lemme \ref{lemma-localized-chern-pre}.
La deuxième est le lemme \ref{lemma-loc-chern-agree}.
Le changement de base de $W' \to W$ par $W_\infty \to W$ est le
morphisme $E = W'_\infty \to W_\infty$. La troisième égalité résulte donc
du lemme \ref{lemma-base-change-loc-chern}. La quatrième
égalité, dans laquelle $i'_\infty : W'_\infty \to W'$ est le
morphisme d'inclusion, résulte du fait que $c_p(W' \to W, Q)$
est une classe bivariante. Pour la cinquième égalité, observons que
$c_p(W' \to W, Q)$ et $c_p(Z' \to X, F)$
se restreignent à la même classe bivariante dans
$A^p((b')^{-1} \to b^{-1}(\mathbf{A}^1_X))$ par
l'hypothèse (1) du lemme, qui dit que $Q$ et $F$ se restreignent
au même objet de $D(\mathcal{O}_{b^{-1}(\mathbf{A}^1_X)})$ ;
utilisons le lemme \ref{lemma-base-change-loc-chern}.
Puisque $(i'_\infty)^*$ annule les cycles à support dans $W'_\infty$
(voir la remarque \ref{remark-gysin-on-cycles}), nous concluons que la cinquième égalité
est vraie. La sixième égalité est satisfaite parce que $W'_\infty$ et $W'_0$
sont les images inverses des diviseurs de Cartier effectifs rationnellement équivalents
$D_0, D_\infty$ dans $\mathbf{P}^1_Z$ et donc $i_\infty^*\beta$ et
$i_0^*\beta$ s'envoient sur la même classe de cycle sur $W'$; plus précisément, tous deux
représentent la classe
$c_1(\mathcal{O}_{\mathbf{P}^1_Z}(1)) \cap c_p(Z \to X, F_) \cap \beta$ par le
lemme \ref{lemma-support-cap-effective-Cartier}.
La septième égalité est satisfaite parce que $c_p(Z \to X, F)$ est
une classe bivariante. Par construction, $W'_0 = Z$ et $i_0^*\beta = \alpha$,
ce qui explique pourquoi la dernière égalité est satisfaite.
\end{proof}







\section{Propriétés des classes de Chern localisées}
\label{section-properties-loc-chern}

\noindent
Les principaux résultats de cette section sont l'additivité et la multiplicativité
des classes de Chern localisées.

\begin{lemma}
\label{lemma-loc-chern-character}
Dans la situation \ref{situation-loc-chern}, supposons que $E|_{X \setminus Z}$ soit nul.
Alors
\begin{align*}
P_1(Z \to X, E) & = c_1(Z \to X, E), \\
P_2(Z \to X, E) & = c_1(Z \to X, E)^2 - 2c_2(Z \to X, E), \\
P_3(Z \to X, E) & = c_1(Z \to X, E)^3 - 3c_1(Z \to X, E)c_2(Z \to X, E)
+ 3c_3(Z \to X, E),
\end{align*}
et ainsi de suite, où les produits sont pris dans l'algèbre $A^{(1)}(Z \to X)$
de la remarque \ref{remark-ring-loc-classes}.
\end{lemma}

\begin{proof}
L'énoncé a un sens parce que le faisceau nul est de rang $< 1$ et
donc les classes $c_p(Z \to X, E)$ sont définies pour tout $p \geq 1$.
Le résultat lui-même découle immédiatement du plus général
lemme \ref{lemma-localized-chern-pre-compose}, puisque les classes de Chern localisées
ont été définies au moyen du procédé du
lemme \ref{lemma-localized-chern-pre}
dans la section \ref{section-localized-chern}.
\end{proof}

\begin{lemma}
\label{lemma-loc-chern-classes-commute}
Dans la situation \ref{situation-loc-chern},
soit $Y \to X$ localement de type fini et soit $c \in A^*(Y \to X)$.
Alors
$$
P_p(Z \to X, E) \circ c = c \circ P_p(Z \to X, E),
$$
respectivement
$$
c_p(Z \to X, E) \circ c = c \circ c_p(Z \to X, E)
$$
dans $A^*(Y \times_X Z \to X)$.
\end{lemma}

\begin{proof}
Cela résulte du lemme \ref{lemma-homomorphism-commute}.
Plus précisément, soient
$$
b : W \to \mathbf{P}^1_X
\quad\text{et}\quad
Q
\quad\text{et}\quad
T' \subset T \subset W_\infty
$$
comme dans la démonstration du lemme \ref{lemma-independent-loc-chern}.
Par définition, $c_p(Z \to X, E) = c'_p(Q)$
comme opérations bivariantes,
où le membre de droite est la classe bivariante construite dans le
lemme \ref{lemma-localized-chern-pre} à l'aide de $W, b, Q, T'$.
Par le lemme \ref{lemma-homomorphism-commute}, nous avons
$P'_p(Q) \circ c = c \circ P'_p(Q)$, resp.\ $c'_p(Q) \circ c = c \circ c'_p(Q)$
dans $A^*(Y \times_X Z \to X)$, et nous concluons.
\end{proof}

\begin{remark}
\label{remark-loc-chern-classes}
Dans la situation \ref{situation-loc-chern}, il est commode de définir
$$
c^{(p)}(Z \to X, E) = 1 + c_1(E) + \ldots + c_{p - 1}(E) +
c_p(Z \to X, E) + c_{p + 1}(Z \to X, E) + \ldots
$$
comme un élément de l'algèbre $A^{(p)}(Z \to X)$ considérée dans la
remarque \ref{remark-ring-loc-classes}.
\end{remark}

\begin{lemma}
\label{lemma-additivity-loc-chern-c}
Soit $(S, \delta)$ comme dans la situation \ref{situation-setup}.
Soit $X$ localement de type fini sur $S$. Soit $Z \to X$
une immersion fermée. Soit
$$
E_1 \to E_2 \to E_3 \to E_1[1]
$$
un triangle distingué d'objets parfaits de $D(\mathcal{O}_X)$.
Supposons que
\begin{enumerate}
\item les restrictions $E_1|_{X \setminus Z}$ et $E_3|_{X \setminus Z}$
soient isomorphes à des $\mathcal{O}_{X \setminus Z}$-modules localement libres de rang fini
de rangs $< p_1$ et $< p_3$ placés en degré $0$, et
\item au moins l'une des conditions suivantes soit satisfaite :
(a) $X$ est quasi-compact,
(b) $X$ a des composantes irréductibles quasi-compactes,
(c) $E_3 \to E_1[1]$ peut être représenté par un morphisme de complexes
localement bornés de $\mathcal{O}_X$-modules localement libres de rang fini, ou
(d) il existe une enveloppe $f : Y \to X$ telle que $Lf^*E_3 \to Lf^*E_1[1]$
puisse être représenté par un morphisme de complexes localement bornés de
$\mathcal{O}_Y$-modules localement libres de rang fini.
\end{enumerate}
Avec la notation de la remarque \ref{remark-loc-chern-classes}, nous avons
$$
c^{(p_1 + p_3)}(Z \to X, E_2) = c^{(p_1)}(Z \to X, E_1)c^{(p_3)}(Z \to X, E_3)
$$
dans $A^{(p_1 + p_3)}(Z \to X)$.
\end{lemma}

\begin{proof}
Observons que les hypothèses impliquent que $E_2|_{X \setminus Z}$ est nul,
resp.\ est isomorphe à un $\mathcal{O}_{X \setminus Z}$-module localement libre de rang fini
de rang $< p_1 + p_3$. L'énoncé a donc un sens.

\medskip\noindent
Soit $f : Y \to X$ une enveloppe. En développant les membres de gauche et de droite
de la formule dans l'énoncé du lemme, nous voyons que nous devons démontrer
certaines égalités de classes dans $A^*(X)$ et dans $A^*(Z \to X)$. Grâce à
l'unicité du lemme \ref{lemma-envelope-bivariant}, il suffit de démontrer les
relations correspondantes dans $A^*(Y)$ et $A^*(Z \to Y)$. Puisque, de plus,
la construction des classes en jeu est compatible au changement de base
(lemme \ref{lemma-base-change-loc-chern}), nous pouvons remplacer $X$ par $Y$
et le triangle distingué par son image inverse.

\medskip\noindent
Dans la démonstration du lemme \ref{lemma-additivity-on-perfect}, nous avons
vu que les conditions (2)(a), (2)(b) et (2)(c) impliquent la condition
(2)(d). Combiné à la discussion du paragraphe précédent, cela nous
ramène au cas examiné dans le paragraphe suivant.

\medskip\noindent
Soit $\varphi^\bullet : \mathcal{E}_3^\bullet[-1] \to \mathcal{E}_1^\bullet$
un morphisme de complexes localement bornés de modules localement libres de rang fini
sur $\mathcal{O}_X$, représentant le morphisme $E_3[-1] \to E_1$
dans la catégorie dérivée. Considérons le schéma
$X' = \mathbf{A}^1 \times X$, de projection
$g : X' \to X$. Posons $Z' = g^{-1}(Z) = \mathbf{A}^1 \times Z$.
Notons $t$ la coordonnée sur $\mathbf{A}^1$. Considérons le cône
$\mathcal{C}^\bullet$ du morphisme de complexes
$$
t g^*\varphi^\bullet :
g^*\mathcal{E}_3^\bullet[-1]
\longrightarrow
g^*\mathcal{E}_1^\bullet
$$
sur $X'$. Nous obtenons un triangle distingué
$$
g^*\mathcal{E}_1^\bullet \to \mathcal{C}^\bullet \to
g^*\mathcal{E}_3^\bullet \to g^*\mathcal{E}_1^\bullet[1]
$$
où les trois premiers termes forment une suite exacte courte scindée terme à terme
de complexes. Manifestement, $\mathcal{C}^\bullet$ est un
complexe borné de $\mathcal{O}_{X'}$-modules localement libres de rang fini
dont la restriction à $X' \setminus Z'$ est isomorphe à un
module localement libre de rang fini
sur $\mathcal{O}_{X' \setminus Z'}$, de rang $< p_1 + p_3$,
placé en degré $0$. Nous avons donc les classes de Chern localisées
$$
c_p(Z' \to X', \mathcal{C}^\bullet) \in A^p(Z' \to X')
$$
pour $p \geq p_1 + p_3$. Pour tout $\alpha \in \CH_k(X)$, considérons
$$
c_p(Z' \to X', \mathcal{C}^\bullet) \cap g^*\alpha
\in \CH_{k + 1 - p}(\mathbf{A}^1 \times X)
$$
Si nous nous restreignons à $t = 0$, alors le morphisme $t g^*\varphi^\bullet$
se restreint à zéro et $\mathcal{C}^\bullet|_{t = 0}$
est la somme directe de $\mathcal{E}_1^\bullet$ et de $\mathcal{E}_3^\bullet$.
Par compatibilité des classes de Chern localisées au changement de base
(lemme \ref{lemma-base-change-loc-chern}), nous concluons que
$$
i_0^* \circ c^{(p_1 + p_3)}(Z' \to X', \mathcal{C}^\bullet) \circ g^* =
c^{(p_1 + p_2)}(Z \to X, E_1 \oplus E_3)
$$
dans $A^{(p_1 + p_3)}(Z \to X)$. D'autre part, si nous nous restreignons à $t = 1$,
alors le morphisme $t g^*\varphi^\bullet$
se restreint à $\varphi$ et $\mathcal{C}^\bullet|_{t = 1}$
est un complexe borné de modules localement libres de rang fini représentant $E_2$.
Nous concluons que
$$
i_1^* \circ c^{(p_1 + p_3)}(Z' \to X', \mathcal{C}^\bullet) \circ g^* =
c^{(p_1 + p_2)}(Z \to X, E_2)
$$
dans $A^{(p_1 + p_3)}(Z \to X)$. Puisque $i_0^* = i_1^*$ par définition de
l'équivalence rationnelle (plus précisément, cela résulte des formules du
lemme \ref{lemma-linebundle-formulae}), nous concluons que
$$
c^{(p_1 + p_2)}(Z \to X, E_2) = c^{(p_1 + p_2)}(Z \to X, E_1 \oplus E_3)
$$
Cela nous ramène au cas examiné dans le paragraphe suivant.

\medskip\noindent
Supposons que $E_2 = E_1 \oplus E_3$ et que les triplets $(X, Z, E_i)$
soient comme dans la situation \ref{situation-loc-chern}.
Pour $i = 1, 3$, soient
$$
b_i : W_i \to \mathbf{P}^1_X
\quad\text{et}\quad
Q_i
\quad\text{et}\quad
T'_i \subset T_i \subset W_{i, \infty}
$$
comme dans la démonstration du lemme \ref{lemma-independent-loc-chern}.
Par définition,
$$
c_p(Z \to X, E_i) = c'_p(Q_i)
$$
où le membre de droite est la classe bivariante construite dans le
lemme \ref{lemma-localized-chern-pre} à l'aide de $W_i, b_i, Q_i, T'_i$.
Posons $W = W_1 \times_{b_1, \mathbf{P}^1_X, b_2} W_2$ et considérons
le diagramme cartésien
$$
\xymatrix{
W \ar[d]_{g_1} \ar[rd]^b \ar[r]_{g_3} & W_3 \ar[d]^{b_3} \\
W_1 \ar[r]^{b_1} & \mathbf{P}^1_X
}
$$
Bien sûr, $b^{-1}(\mathbf{A}^1)$ s'envoie isomorphiquement sur $\mathbf{A}^1_X$.
Observons que $T' = g_1^{-1}(T'_1) \cap g_2^{-1}(T'_2)$ contient encore
tous les points de $W_\infty$ situés au-dessus de $X \setminus Z$.
Par le lemme \ref{lemma-localized-chern-pre-independent}, nous pouvons utiliser
$W$, $b$, $g_i^*\mathcal{Q}_i$ et
$T'$ pour construire $c_p(Z \to X, E_i)$ pour $i = 1, 3$.
De plus, par l'indépendance plus forte donnée dans le
lemme \ref{lemma-independent-loc-chern-bQ}, nous pouvons utiliser
$W$, $b$, $g_1^*Q_1 \oplus g_3^*Q_3$ et $T'$
pour calculer les classes $c_p(Z \to X, E_2)$.
L'égalité désirée résulte donc du
lemme \ref{lemma-localized-chern-pre-sum-c}.
\end{proof}

\begin{lemma}
\label{lemma-additivity-loc-chern-P}
Soit $(S, \delta)$ comme dans la situation \ref{situation-setup}.
Soit $X$ localement de type fini sur $S$. Soit $Z \to X$
une immersion fermée. Soit
$$
E_1 \to E_2 \to E_3 \to E_1[1]
$$
un triangle distingué d'objets parfaits de $D(\mathcal{O}_X)$.
Supposons que
\begin{enumerate}
\item les restrictions $E_1|_{X \setminus Z}$ et $E_3|_{X \setminus Z}$
soient nulles, et
\item au moins l'une des conditions suivantes soit satisfaite :
(a) $X$ est quasi-compact,
(b) $X$ a des composantes irréductibles quasi-compactes,
(c) $E_3 \to E_1[1]$ peut être représenté par un morphisme de complexes
localement bornés de $\mathcal{O}_X$-modules localement libres de rang fini, ou
(d) il existe une enveloppe $f : Y \to X$ telle que $Lf^*E_3 \to Lf^*E_1[1]$
puisse être représenté par un morphisme de complexes localement bornés de
$\mathcal{O}_Y$-modules localement libres de rang fini.
\end{enumerate}
Alors nous avons
$$
P_p(Z \to X, E_2) = P_p(Z \to X, E_1) + P_p(Z \to X, E_3)
$$
pour tout $p \in \mathbf{Z}$ et, par conséquent,
$ch(Z \to X, E_2) = ch(Z \to X, E_1) + ch(Z \to X, E_3)$.
\end{lemma}

\begin{proof}
La démonstration est exactement la même que celle du
lemme \ref{lemma-additivity-loc-chern-c},
si ce n'est qu'elle utilise le
lemme \ref{lemma-localized-chern-pre-sum-P}
à la toute fin. Pour $p > 0$, nous pouvons déduire ce lemme
du lemme \ref{lemma-additivity-loc-chern-c} avec $p_1 = p_3 = 1$
et de la relation entre $P_p(Z \to X, E)$ et $c_p(Z \to X, E)$ donnée dans le
lemme \ref{lemma-loc-chern-character}. Le cas $p = 0$ peut se démontrer
directement (il n'est intéressant que si $X$ a une composante connexe
entièrement contenue dans $Z$).
\end{proof}

\begin{lemma}
\label{lemma-loc-chern-tensor-product}
Dans la situation \ref{situation-setup}, soit $X$ localement de type fini sur $S$.
Soient $Z_i \subset X$, $i = 1, 2$, des sous-schémas fermés. Soient $F_i$, $i = 1, 2$,
des objets parfaits de $D(\mathcal{O}_X)$. Supposons, pour $i = 1, 2$, que
$F_i|_{X \setminus Z_i}$ soit nul\footnote{Il existe vraisemblablement une
variante de ce lemme où nous supposons seulement que $F_i|_{X \setminus Z_i}$
est isomorphe à un $\mathcal{O}_{X \setminus Z_i}$-module localement libre de rang fini
de rang $< p_i$.} et que $F_i$ sur $X$ satisfasse l'hypothèse
(3) de la situation \ref{situation-loc-chern}. Notons
$r_i = P_0(Z_i \to X, F_i) \in A^0(Z_i \to X)$.
Alors nous avons
$$
c_1(Z_1 \cap Z_2 \to X, F_1 \otimes_{\mathcal{O}_X}^\mathbf{L} F_2) =
r_1 c_1(Z_2 \to X, F_2) + r_2 c_1(Z_1 \to X, F_1)
$$
dans $A^1(Z_1 \cap Z_2 \to X)$ et
\begin{align*}
c_2(Z_1 \cap Z_2 \to X, F_1 \otimes_{\mathcal{O}_X}^\mathbf{L} F_2)
& =
r_1 c_2(Z_2 \to X, F_2) +
r_2 c_2(Z_1 \to X, F_1) + \\
& {r_1 \choose 2} c_1(Z_2 \to X, F_2)^2 + \\
& (r_1r_2 - 1) c_1(Z_2 \to X, F_2)c_1(Z_1 \to X, F_1) + \\
& {r_2 \choose 2} c_1(Z_1 \to X, F_1)^2
\end{align*}
dans $A^2(Z_1 \cap Z_2 \to X)$, et ainsi de suite pour les classes de Chern supérieures.
De même, nous avons
$$
ch(Z_1 \cap Z_2 \to X, F_1 \otimes_{\mathcal{O}_X}^\mathbf{L} F_2) =
ch(Z_1 \to X, F_1) ch(Z_2 \to X, F_2)
$$
dans $\prod_{p \geq 0} A^p(Z_1 \cap Z_2 \to X) \otimes \mathbf{Q}$.
Plus précisément, nous avons
$$
P_p(Z_1 \cap Z_2 \to X, F_1 \otimes_{\mathcal{O}_X}^\mathbf{L} F_2) =
\sum\nolimits_{p_1 + p_2 = p}
{p \choose p_1} P_{p_1}(Z_1 \to X, F_1) P_{p_2}(Z_2 \to X, F_2)
$$
dans $A^p(Z_1 \cap Z_2 \to X)$.
\end{lemma}

\begin{proof}
Choisissons des morphismes propres $b_i : W_i \to \mathbf{P}^1_X$ et
$Q_i \in D(\mathcal{O}_{W_i})$, ainsi que des sous-schémas fermés
$T_i \subset W_{i, \infty}$, comme dans la construction des
classes de Chern localisées pour $F_i$, ou plus généralement comme dans le
lemme \ref{lemma-independent-loc-chern-bQ}. Choisissons un diagramme
commutatif
$$
\xymatrix{
W \ar[d]^{g_1} \ar[rd]^b \ar[r]_{g_2} & W_2 \ar[d]^{b_2} \\
W_1 \ar[r]^{b_1} & \mathbf{P}^1_X
}
$$
où tous les morphismes sont propres et sont des isomorphismes au-dessus de
$\mathbf{A}^1_X$. Par exemple, nous pouvons prendre pour $W$ l'adhérence
du graphe de l'isomorphisme entre
$b_1^{-1}(\mathbf{A}^1_X)$ et $b_2^{-1}(\mathbf{A}^1_X)$.
Par le lemme \ref{lemma-independent-loc-chern-bQ}, nous pouvons utiliser
$W$, $b = b_i \circ g_i$, $Lg_i^*Q_i$ et
$g_i^{-1}(T_i)$ pour construire les classes de Chern localisées
$c_p(Z_i \to X, F_i)$. Nous nous ramenons ainsi à la situation décrite
dans le paragraphe suivant.

\medskip\noindent
Supposons que nous ayons
\begin{enumerate}
\item un morphisme propre $b : W \to \mathbf{P}^1_X$ qui est un isomorphisme
au-dessus de $\mathbf{A}^1_X$,
\item $E_i \subset W_\infty$ est l'image inverse de $Z_i$,
\item des objets parfaits $Q_i \in D(\mathcal{O}_W)$ dont les classes de Chern
sont définies, tels que
\begin{enumerate}
\item la restriction de $Q_i$ à $b^{-1}(\mathbf{A}^1_X)$ soit
l'image inverse de $F_i$, et
\item il existe un sous-schéma fermé $T_i \subset W_\infty$ contenant
tous les points de $W_\infty$ situés au-dessus de $X \setminus Z_i$ tel que
$Q_i|_{T_i}$ soit nul.
\end{enumerate}
\end{enumerate}
Par le lemme \ref{lemma-independent-loc-chern-bQ}, nous avons
$$
c_p(Z_i \to X, F_i) = c'_p(Q_i) =
(E_i \to Z_i)_* \circ c'_p(Q_i|_{E_i}) \circ C
$$
et
$$
P_p(Z_i \to X, F_i) = P'_p(Q_i) =
(E_i \to Z_i)_* \circ P'_p(Q_i|_{E_i}) \circ C
$$
pour $i = 1, 2$. Ensuite, nous observons que
$Q = Q_1 \otimes_{\mathcal{O}_W}^\mathbf{L} Q_2$
satisfait (3)(a) et (3)(b) pour $F_1 \otimes_{\mathcal{O}_X}^\mathbf{L} F_2$
et $T_1 \cup T_2$. Nous voyons donc que
$$
c_p(Z_1 \cap Z_2 \to X, F_1 \otimes_{\mathcal{O}_X}^\mathbf{L} F_2) =
(E_1 \cap E_2 \to Z_1 \cap Z_2)_* \circ
c'_p(Q|_{E_1 \cap E_2}) \circ C
$$
et
$$
P_p(Z_1 \cap Z_2 \to X, F_1 \otimes_{\mathcal{O}_X}^\mathbf{L} F_2) =
(E_1 \cap E_2 \to Z_1 \cap Z_2)_* \circ
P'_p(Q|_{E_1 \cap E_2}) \circ C
$$
par le même lemme. Par le lemme \ref{lemma-silly-tensor-product},
les classes $c'_p(Q|_{E_1 \cap E_2})$ et $P'_p(Q|_{E_1 \cap E_2})$
peuvent être développées de la manière appropriée en fonction des classes
$c'_p(Q_i|_{E_i})$ et $P'_p(Q_i|_{E_i})$. Enfin, le
lemme \ref{lemma-homomorphism-final}
nous dit que les polynômes en $c'_p(Q_i|_{E_i})$ et $P'_p(Q_i|_{E_i})$
coïncident avec les polynômes correspondants en
$c'_p(Q_i)$ et $P'_p(Q_i)$, comme désiré.
\end{proof}






\section{Éclatement à l'infini}
\label{section-blowup-Z-first}

\noindent
Soit $X$ un schéma. Soit $Z \subset X$ un sous-schéma fermé découpé
par un faisceau quasi-cohérent d'idéaux de type fini. Notons $X' \to X$
l'éclatement de centre $Z$. Soit $b : W \to \mathbf{P}^1_X$
l'éclatement de centre $\infty(Z)$. Notons $E \subset W$ le diviseur
exceptionnel. Il existe un diagramme commutatif
$$
\xymatrix{
X' \ar[r] \ar[d] & W \ar[d]^b \\
X \ar[r]^\infty & \mathbf{P}^1_X
}
$$
dont les flèches horizontales sont des immersions fermées
(Diviseurs, lemme \ref{divisors-lemma-strict-transform}). Notons $E \subset W$
le diviseur exceptionnel et $W_\infty \subset W$ l'image inverse
de $(\mathbf{P}^1_X)_\infty$. Alors les assertions suivantes sont satisfaites :
\begin{enumerate}
\item $b$ est un isomorphisme au-dessus de
$\mathbf{A}^1_X \cup \mathbf{P}^1_{X \setminus Z}$,
\item $X'$ est un diviseur de Cartier effectif sur $W$,
\item $X' \cap E$ est le diviseur exceptionnel de $X' \to X$,
\item $W_\infty = X' + E$ comme diviseurs de Cartier effectifs sur $W$,
\item $E = \underline{\text{Proj}}_Z(\mathcal{C}_{Z/X, *}[S])$, où $S$
est une variable placée en degré $1$,
\item $X' \cap E = \underline{\text{Proj}}_Z(\mathcal{C}_{Z/X, *})$,
\item
\label{item-cone-is-open}
$E \setminus X' = E \setminus (X' \cap E) =
\underline{\Spec}_Z(\mathcal{C}_{Z/X, *}) = C_ZX$,
\item
\label{item-find-Z-in-blowup}
il existe une immersion fermée $\mathbf{P}^1_Z \to W$ dont la
composée avec $b$ est le morphisme d'inclusion
$\mathbf{P}^1_Z \to \mathbf{P}^1_X$ et dont le changement de base par $\infty$
est la composée $Z \to C_ZX \to E \to W_\infty$, où la première
flèche est le sommet du cône.
\end{enumerate}
Rappelons que $\mathcal{C}_{Z/X, *}$ est l'algèbre conormale de $Z$ dans $X$ ;
voir Diviseurs, définition \ref{divisors-definition-conormal-sheaf}, et
que $C_ZX$ est le cône normal de $Z$ dans $X$ ; voir
Diviseurs, définition \ref{divisors-definition-normal-cone}.

\medskip\noindent
Donnons maintenant la démonstration des assertions numérotées ci-dessus. Nous recommandons
vivement au lecteur d'étudier quelques exemples plutôt que de lire les
démonstrations.

\medskip\noindent
L'assertion (1) résulte de l'assertion correspondante de Diviseurs, lemme
\ref{divisors-lemma-blowing-up-gives-effective-Cartier-divisor}.
Observons que $E \subset W$ est un diviseur de Cartier effectif par le même lemme.

\medskip\noindent
Observons que $W_\infty$ est un diviseur de Cartier effectif par
Diviseurs, lemme \ref{divisors-lemma-blow-up-pullback-effective-Cartier}.
Puisque $E \subset W_\infty$, nous pouvons écrire $W_\infty = D + E$ pour un certain
diviseur de Cartier effectif $D$; voir
Diviseurs, lemme \ref{divisors-lemma-difference-effective-Cartier-divisors}.
Nous verrons ci-dessous que $D = X'$, ce qui démontrera (2) et (4).

\medskip\noindent
Puisque $X'$ est le transformé strict de l'immersion fermée
$\infty : X \to \mathbf{P}^1_X$ (voir ci-dessus), il s'ensuit que le diviseur
exceptionnel de $X' \to X$ est égal à l'intersection $X' \cap E$
(par exemple parce que tous deux sont découpés par l'image inverse du
faisceau d'idéaux de $Z$ sur $X'$). Cela démontre (3).

\medskip\noindent
L'intersection de $\infty(Z)$ avec $\mathbf{P}^1_Z$ est le diviseur de Cartier
effectif $(\mathbf{P}^1_Z)_\infty$ ; le transformé strict
de $\mathbf{P}^1_Z$ par l'éclatement $b$ s'envoie donc isomorphiquement sur
$\mathbf{P}^1_Z$ (voir Diviseurs, lemmes \ref{divisors-lemma-strict-transform}
et \ref{divisors-lemma-blow-up-effective-Cartier-divisor}).
Cela nous donne le morphisme $\mathbf{P}^1_Z \to W$ mentionné en (8).
C'est une immersion fermée puisque $b$ est séparé ; voir
Schémas, lemme \ref{schemes-lemma-section-immersion}.

\medskip\noindent
Supposons que $\Spec(A) \subset X$ soit un ouvert affine et que $Z \cap \Spec(A)$
corresponde à l'idéal de type fini $I \subset A$.
Un voisinage affine de $\infty(Z \cap \Spec(A))$ est un
espace affine sur $A$ de coordonnée $s = T_0/T_1$. Notons
$J = (I, s) \subset A[s]$ l'idéal engendré par $I$ et $s$.
Soit $B = A[s] \oplus J \oplus J^2 \oplus \ldots$ l'algèbre de Rees
de $(A[s], J)$. Observons que
$$
J^n =
I^n \oplus sI^{n - 1} \oplus s^2I^{n - 2} \ldots \oplus s^nA
\oplus s^{n + 1}A \oplus \ldots
$$
comme $A$-sous-module de $A[s]$ pour tout $n \geq 0$. Considérons le sous-schéma ouvert
$$
\text{Proj}(B) = \text{Proj}(A[s] \oplus J \oplus J^2 \oplus \ldots)
\subset W
$$
Enfin, notons $S$ l'élément $s \in J$ considéré comme un élément de degré $1$
de $B$.

\medskip\noindent
Puisque la formation de $\text{Proj}$ commute au changement de base
(Constructions, lemme \ref{constructions-lemma-base-change-map-proj}),
nous voyons que
$$
E = \text{Proj}(B \otimes_{A[s]} A/I) =
\text{Proj}((A/I \oplus I/I^2 \oplus I^2/I^3 \oplus \ldots)[S])
$$
La vérification de l'égalité $B \otimes_{A[s]} A/I = \bigoplus J^n/J^{n + 1}$
telle qu'elle a été donnée
résulte immédiatement de notre description des puissances $J^n$ ci-dessus.
Cela démontre (5), car l'algèbre conormale de $Z \cap \Spec(A)$
dans $\Spec(A)$ correspond à la $A$-algèbre graduée
$A/I \oplus I/I^2 \oplus I^2/I^3 \oplus \ldots$ par
Diviseurs, lemme \ref{divisors-lemma-affine-conormal-sheaf}.

\medskip\noindent
Rappelons que $\text{Proj}(B)$ est recouvert par les ouverts affines
$D_+(S)$ et $D_+(f^{(1)})$, pour $f \in I$, qui sont
les spectres des algèbres d'éclatement affines $A[s][\frac{J}{s}]$
et $A[s][\frac{J}{f}]$ ; voir
Diviseurs, lemme \ref{divisors-lemma-blowing-up-affine}, et
Algèbre, définition \ref{algebra-definition-blow-up}.
Nous allons décrire chacun de ces ouverts affines, ce qui achèvera la
démonstration.

\medskip\noindent
L'ouvert $D_+(S)$, c'est-à-dire le spectre de $A[s][\frac{J}{s}]$.
Il résulte de la description ci-dessus des puissances de $J$
que
$$
A[s][\textstyle{\frac{J}{s}}] = \sum s^{-n}I^n[s] \subset A[s, s^{-1}]
$$
L'élément $s$ est un non-diviseur de zéro dans cet anneau et définit aussi bien le diviseur
exceptionnel $E$ que $W_\infty$. Ainsi, $D \cap D_+(S) = \emptyset$.
Enfin, le quotient de $A[s][\frac{J}{s}]$ par $s$ est l'algèbre conormale
$$
A/I \oplus I/I^2 \oplus I^2/I^3 \oplus \ldots
$$
Cela démontre (7).

\medskip\noindent
L'ouvert $D_+(f^{(1)})$, c'est-à-dire le spectre de $A[s][\frac{J}{f}]$.
Il résulte de la description ci-dessus des puissances de $J$ que
$$
A[s][\textstyle{\frac{J}{f}}] =
A[\textstyle{\frac{I}{f}}][\textstyle{\frac{s}{f}}]
$$
où $\frac{s}{f}$ est une variable. L'élément $f$ est un non-diviseur de zéro
dans cet anneau, et son schéma des zéros définit le diviseur exceptionnel $E$.
Puisque $s$ définit $W_\infty$ et que $s = f \cdot \frac{s}{f}$,
nous concluons que $\frac{s}{f}$ définit
le diviseur $D$ construit ci-dessus. Nous voyons alors que
$$
D \cap D_+(f^{(1)}) = \Spec(A[\textstyle{\frac{I}{f}}])
$$
qui est l'ouvert correspondant de l'éclatement $X'$ au-dessus de $\Spec(A)$.
Plus précisément, le morphisme surjectif d'$A[s]$-algèbres graduées
$B \to A \oplus I \oplus I^2 \oplus \ldots$
vers l'algèbre de Rees de $(A, I)$ correspond à l'immersion
fermée $X' \to W$ au-dessus de $\Spec(A[s])$.
Cela démontre $D = X'$, comme désiré.

\medskip\noindent
Démontrons (6). Observons que le schéma des zéros de $\frac{s}{f}$
dans le paragraphe précédent est la restriction du schéma des zéros de $S$
à l'ouvert affine $D_+(f^{(1)})$. Nous voyons donc que $S = 0$ définit
$X' \cap E$ sur $E$. Ainsi, (6) résulte de (5).

\medskip\noindent
Enfin, nous devons démontrer la dernière partie de (8). C'est clair
car le morphisme $\mathbf{P}^1_Z \to W$ est décrit localement sur des ouverts affines
par la surjection
$$
B \to B \otimes_{A[s]} A/I =
(A/I \oplus I/I^2 \oplus I^2/I^3 \oplus \ldots)[S] \to
A/I[S]
$$
et l'identification $\text{Proj}(A/I[S]) = \Spec(A/I)$.
Certains détails sont omis.






\section{Homomorphismes de Gysin en codimension supérieure}
\label{section-gysin-higher-codimension}

\noindent
Soit $(S, \delta)$ comme dans la situation \ref{situation-setup}. Soit $X$ un schéma
localement de type fini sur $S$. Dans cette section, nous allons considérer des
triplets
$$
(Z \to X, \mathcal{N}, \sigma : \mathcal{N}^\vee \to \mathcal{C}_{Z/X})
$$
formés d'une immersion fermée $Z \to X$, d'un
$\mathcal{O}_Z$-module localement libre $\mathcal{N}$ et d'une surjection
$\sigma : \mathcal{N}^\vee \to \mathcal{C}_{Z/X}$ du dual
de $\mathcal{N}$ vers le faisceau conormal de $Z$ dans $X$ ; voir
Morphismes, section \ref{morphisms-section-conormal-sheaf}.
Nous dirons que
$\mathcal{N}$ est un {\it faisceau normal virtuel de $Z$ dans $X$}.

\begin{lemma}
\label{lemma-pullback-virtual-normal-sheaf}
Soit $(S, \delta)$ comme dans la situation \ref{situation-setup}. Soit
$$
\xymatrix{
Z' \ar[r] \ar[d]_g & X' \ar[d]^f \\
Z \ar[r] & X
}
$$
un diagramme cartésien de schémas localement de type fini sur $S$
dont les flèches horizontales sont des immersions fermées.
Si $\mathcal{N}$ est un faisceau normal virtuel pour $Z$ dans $X$, alors
$\mathcal{N}' = g^*\mathcal{N}$ est un faisceau normal virtuel pour
$Z'$ dans $X'$.
\end{lemma}

\begin{proof}
Cela résulte de la surjectivité de l'application
$g^*\mathcal{C}_{Z/X} \to \mathcal{C}_{Z'/X'}$ démontrée dans
Morphismes, lemme \ref{morphisms-lemma-conormal-functorial-flat}.
\end{proof}

\noindent
Soit $(S, \delta)$ comme dans la situation \ref{situation-setup}. Soit $X$ un schéma
localement de type fini sur $S$. Soit $\mathcal{N}$ un fibré normal virtuel
pour une immersion fermée $Z \to X$. Dans cette situation, posons
$$
p : N = \underline{\Spec}_Z(\text{Sym}(\mathcal{N}^\vee)) \longrightarrow Z
$$
le fibré vectoriel sur $Z$
dont les sections correspondent aux sections de $\mathcal{N}$.
Dans cette situation, nous avons des immersions fermées canoniques
$$
C_ZX \longrightarrow N_ZX \longrightarrow N
$$
La première immersion fermée est donnée par Diviseurs, équation
(\ref{divisors-equation-normal-cone-in-normal-bundle}),
et la seconde immersion fermée correspond à la surjection
$\text{Sym}(\mathcal{N}^\vee) \to \text{Sym}(\mathcal{C}_{Z/X})$
induite par $\sigma$.
Soit
$$
b : W \longrightarrow \mathbf{P}^1_X
$$
l'éclatement en $\infty(Z)$ construit dans la
section \ref{section-blowup-Z-first}. D'après le
lemme \ref{lemma-gysin-at-infty},
nous avons une classe bivariante canonique dans
$$
C \in A^0(W_\infty \to X)
$$
Considérons l'immersion ouverte $j : C_ZX \to W_\infty$ de
(\ref{item-cone-is-open}) et l'immersion fermée
$i : C_ZX \to N$ construite ci-dessus. D'après le lemme \ref{lemma-vectorbundle},
pour tout $\alpha \in \CH_k(X)$, il existe un unique
$\beta \in \CH_*(Z)$ tel que
$$
i_*j^*(C \cap \alpha) = p^*\beta
$$
Posons $c(Z \to X, \mathcal{N}) \cap \alpha = \beta$.

\begin{lemma}
\label{lemma-construction-gysin}
La construction ci-dessus définit une classe bivariante\footnote{La
notation $A^*(Z \to X)^\wedge$ est examinée dans la
remarque \ref{remark-completion-bivariant}.
Si $X$ est quasi-compact, alors $A^*(Z \to X)^\wedge = A^*(Z \to X)$.}
$$
c(Z \to X, \mathcal{N}) \in A^*(Z \to X)^\wedge
$$
et, de plus, la construction est compatible au changement de base
comme dans le lemme \ref{lemma-pullback-virtual-normal-sheaf}.
Si $\mathcal{N}$ est de rang constant $r$, alors
$c(Z \to X, \mathcal{N}) \in A^r(Z \to X)$.
\end{lemma}

\begin{proof}
Comme $i_* \circ j^* \circ C$ et $p^*$ sont tous deux des classes bivariantes
(voir les lemmes \ref{lemma-flat-pullback-bivariant} et
\ref{lemma-push-proper-bivariant}), nous pouvons utiliser l'égalité
$$
i_* \circ j^* \circ C  = p^* \circ c(Z \to X, \mathcal{N})
$$
(interprétée convenablement) pour définir $c(Z \to X, \mathcal{N})$
comme classe bivariante. Cela fonctionne parce que $p^*$ est toujours
bijectif sur les groupes de Chow d'après le lemme \ref{lemma-vectorbundle}.

\medskip\noindent
Soient $X' \to X$, $Z' \to X'$ et $\mathcal{N}'$ comme dans le
lemme \ref{lemma-pullback-virtual-normal-sheaf}. Écrivons
$c = c(Z \to X, \mathcal{N})$ et $c' = c(Z' \to X', \mathcal{N}')$.
La seconde assertion du lemme signifie que $c'$ est la restriction de $c$
comme dans la remarque \ref{remark-restriction-bivariant}. Puisque nous affirmons que cela
vaut pour tout $X'/X$ localement de type fini, un argument formel
montre qu'il suffit de vérifier que $c' \cap \alpha' = c \cap \alpha'$
pour $\alpha' \in \CH_k(X')$.
Pour le voir, remarquons que nous avons un diagramme commutatif
$$
\xymatrix{
C_{Z'}X' \ar[d] \ar[r] &
W'_\infty \ar[d] \ar[r] &
W' \ar[d] \ar[r] &
\mathbf{P}^1_{X'} \ar[d] \\
C_ZX \ar[r] &
W_\infty \ar[r] &
W \ar[r] &
\mathbf{P}^1_X
}
$$
qui induit des immersions fermées :
$$
W' \to W \times_{\mathbf{P}^1_X} \mathbf{P}^1_{X'},\quad
W'_\infty \to W_\infty \times_X X',\quad
C_{Z'}X' \to C_ZX \times_Z Z'
$$
Pour obtenir $c \cap \alpha'$, nous utilisons la classe $C \cap \alpha'$
définie au moyen du morphisme
$W \times_{\mathbf{P}^1_X} \mathbf{P}^1_{X'} \to \mathbf{P}^1_{X'}$
dans le lemme \ref{lemma-gysin-at-infty}.
Pour obtenir $c' \cap \alpha'$, d'autre part, nous utilisons la classe
$C' \cap \alpha'$ définie au moyen du morphisme $W' \to \mathbf{P}^1_{X'}$.
D'après le lemme \ref{lemma-gysin-at-infty-independent}, l'image directe de
$C' \cap \alpha'$ par l'immersion fermée
$W'_\infty \to (W \times_{\mathbf{P}^1_X} \mathbf{P}^1_{X'})_\infty$
est égale à $C \cap \alpha'$. Il en va donc de même pour les images inverses
sur les ouverts
$$
C_{Z'}X' \subset W'_\infty,\quad
C_ZX \times_Z Z' \subset (W \times_{\mathbf{P}^1_X} \mathbf{P}^1_{X'})_\infty
$$
d'après le lemme \ref{lemma-flat-pullback-proper-pushforward}.
Puisque nous avons un diagramme commutatif
$$
\xymatrix{
C_{Z'} X' \ar[d] \ar[r] & N' \ar@{=}[d] \\
C_ZX \times_Z Z' \ar[r] & N \times_Z Z'
}
$$
les images directes de ces classes donnent la même classe sur $N'$, ce qui
prouve que nous obtenons le même élément $c \cap \alpha' = c' \cap \alpha'$
dans $\CH_*(Z')$.
\end{proof}

\begin{lemma}
\label{lemma-gysin-decompose}
Soit $(S, \delta)$ comme dans la situation \ref{situation-setup}. Soit $X$ un schéma
localement de type fini sur $S$. Soit $\mathcal{N}$ un faisceau normal virtuel
pour un sous-schéma fermé $Z$ de $X$. Supposons que nous ayons une suite
exacte courte $0 \to \mathcal{N}' \to \mathcal{N} \to \mathcal{E} \to 0$
de $\mathcal{O}_Z$-modules localement libres de rang fini telle que la surjection donnée
$\sigma : \mathcal{N}^\vee \to \mathcal{C}_{Z/X}$ se factorise par une application
$\sigma' : (\mathcal{N}')^\vee \to \mathcal{C}_{Z/X}$.
Alors
$$
c(Z \to X, \mathcal{N}) = c_{top}(\mathcal{E}) \circ c(Z \to X, \mathcal{N}')
$$
comme classes bivariantes.
\end{lemma}

\begin{proof}
Notons $N' \to N$ l'immersion fermée de fibrés vectoriels correspondant
à la surjection $\mathcal{N}^\vee \to (\mathcal{N}')^\vee$. Nous avons alors
des immersions fermées
$$
C_ZX \to N' \to N
$$
Ainsi, la relation voulue entre les classes bivariantes résulte
immédiatement du lemme \ref{lemma-easy-virtual-class}.
\end{proof}

\begin{lemma}
\label{lemma-gysin-excess}
Soit $(S, \delta)$ comme dans la situation \ref{situation-setup}. Considérons
un diagramme cartésien
$$
\xymatrix{
Z' \ar[r] \ar[d]_g & X' \ar[d]^f \\
Z \ar[r] & X
}
$$
de schémas localement de type fini sur $S$, dont les flèches horizontales
sont des immersions fermées. Soient $\mathcal{N}$, resp.\ $\mathcal{N}'$,
un faisceau normal virtuel pour $Z \subset X$, resp.\ $Z' \to X'$.
Supposons donnée une suite exacte courte
$0 \to \mathcal{N}' \to g^*\mathcal{N} \to \mathcal{E} \to 0$
de modules localement libres de rang fini sur $Z'$, telle que le diagramme
$$
\xymatrix{
g^*\mathcal{N}^\vee \ar[r] \ar[d] &
(\mathcal{N}')^\vee \ar[d] \\
g^*\mathcal{C}_{Z/X} \ar[r] &
\mathcal{C}_{Z'/X'}
}
$$
soit commutatif. Alors nous avons
$$
res(c(Z \to X, \mathcal{N})) =
c_{top}(\mathcal{E}) \circ c(Z' \to X', \mathcal{N}')
$$
dans $A^*(Z' \to X')^\wedge$.
\end{lemma}

\begin{proof}
D'après le lemme \ref{lemma-construction-gysin}, nous avons
$res(c(Z \to X, \mathcal{N})) = c(Z' \to X', g^*\mathcal{N})$,
et l'égalité résulte du lemme \ref{lemma-gysin-decompose}.
\end{proof}

\noindent
Soit $(S, \delta)$ comme dans la situation \ref{situation-setup}. Soit $X$ un schéma
localement de type fini sur $S$. Soit $\mathcal{N}$ un faisceau normal virtuel
pour un sous-schéma fermé $Z$ de $X$. Soit $Y \to X$ un morphisme
localement de type fini. Supposons que $Z \times_X Y \to Y$ soit une
immersion fermée régulière, voir
Diviseurs, section \ref{divisors-section-regular-immersions}.
Dans ce cas, le faisceau conormal $\mathcal{C}_{Z \times_X Y/Y}$ est un
$\mathcal{O}_{Z \times_X Y}$-module localement libre de rang fini, et nous obtenons une suite
exacte courte
$$
0 \to \mathcal{E}^\vee \to
\mathcal{N}^\vee|_{Z \times_X Y} \to \mathcal{C}_{Z \times_X Y/Y} \to 0
$$
Le quotient $\mathcal{N}|_{Y \times_X Z} \to \mathcal{E}$ est appelé le
{\it faisceau normal d'excès} de la situation.

\begin{lemma}
\label{lemma-gysin-fundamental}
Dans la situation décrite juste ci-dessus, supposons que $\dim_\delta(Y) = n$
et que $\mathcal{C}_{Y \times_X Z/Z}$ soit de rang constant $r$.
Alors
$$
c(Z \to X, \mathcal{N}) \cap [Y]_n =
c_{top}(\mathcal{E}) \cap [Z \times_X Y]_{n - r}
$$
dans $\CH_*(Z \times_X Y)$.
\end{lemma}

\begin{proof}
La classe bivariante $c_{top}(\mathcal{E}) \in A^*(Z \times_X Y)$ a été
définie dans la remarque \ref{remark-top-chern-class}.
D'après le lemme \ref{lemma-construction-gysin}, nous pouvons remplacer $X$ par $Y$.
Ainsi, nous pouvons supposer que $Z \to X$ est une immersion fermée régulière
de codimension $r$, que $\dim_\delta(X) = n$, et nous devons
montrer que $c(Z \to X, \mathcal{N}) \cap [X]_n =
c_{top}(\mathcal{E}) \cap [Z]_{n - r}$ dans $\CH_*(Z)$.
D'après le lemme \ref{lemma-gysin-decompose}, nous pouvons même supposer que
$\mathcal{N}^\vee \to \mathcal{C}_{Z/X}$ est un isomorphisme.
Autrement dit, nous devons montrer que
$c(Z \to X, \mathcal{C}_{Z/X}^\vee) \cap [X]_n = [Z]_{n - r}$ dans $\CH_*(Z)$.

\medskip\noindent
Reprenons les étapes de la définition de
$c(Z \to X, \mathcal{C}_{Z/X}^\vee) \cap [X]_n$. Soit
$b : W \to \mathbf{P}^1_X$
l'éclatement de $\infty(Z)$. Nous devons d'abord calculer
$C \cap [X]_n$, où $C \in A^0(W_\infty \to X)$ est
la classe du lemme \ref{lemma-gysin-at-infty}.
Pour ce faire, remarquons que $[W]_{n + 1}$
est un cycle sur $W$ dont la restriction à $\mathbf{A}^1_X$ est
égale à l'image inverse plate de $[X]_n$. Ainsi, $C \cap [X]_n$
est égal à $i_\infty^*[W]_{n + 1}$. Comme $W_\infty$ est un
diviseur de Cartier effectif sur $W$, nous avons
$i_\infty^*[W]_{n + 1} = [W_\infty]_n$, voir le lemme \ref{lemma-easy-gysin}.
La restriction de cette classe à l'ouvert $C_ZX \subset W_\infty$
est bien sûr simplement $[C_ZX]_n$. Comme $Z \subset X$ est régulièrement
immergé, nous avons
$$
\mathcal{C}_{Z/X, *} = \text{Sym}(\mathcal{C}_{Z/X})
$$
comme $\mathcal{O}_Z$-algèbres graduées, voir
Diviseurs, lemme \ref{divisors-lemma-quasi-regular-immersion}.
Ainsi, $p : N = C_ZX \to Z$ est le morphisme structural du
fibré vectoriel associé au module localement libre de rang fini
$\mathcal{C}_{Z/X}$ de rang $r$. Il est alors clair que
$p^*[Z]_{n - r} = [C_ZX]_n$, et la démonstration est achevée.
\end{proof}

\begin{lemma}
\label{lemma-gysin-easy}
Soit $(S, \delta)$ comme dans la situation \ref{situation-setup}. Soit $X$ un schéma
localement de type fini sur $S$. Soit $\mathcal{N}$ un faisceau normal virtuel
pour un sous-schéma fermé $Z$ de $X$. Soit $Y \to X$ un morphisme
localement de type fini. Étant donnés des entiers $r$, $n$, supposons que
\begin{enumerate}
\item $\mathcal{N}$ soit localement libre de rang $r$,
\item toute composante irréductible de $Y$ soit de $\delta$-dimension $n$,
\item $\dim_\delta(Z \times_X Y) \leq n - r$, et
\item pour $\xi \in Z \times_X Y$ tel que $\delta(\xi) = n - r$,
l'anneau local $\mathcal{O}_{Y, \xi}$ soit de Cohen-Macaulay.
\end{enumerate}
Alors $c(Z \to X, \mathcal{N}) \cap [Y]_n = [Z \times_X Y]_{n - r}$
dans $\CH_{n - r}(Z \times_X Y)$.
\end{lemma}

\begin{proof}
L'énoncé a un sens puisque $Z \times_X Y$ est un sous-schéma fermé de $Y$.
Comme $\mathcal{N}$ est de rang $r$, nous savons que
$c(Z \to X, \mathcal{N}) \cap [Y]_n$ appartient à $\CH_{n - r}(Z \times_X Y)$.
Puisque $\dim_\delta(Z \cap Y) \leq n - r$, le groupe de Chow
$\CH_{n - r}(Z \times_X Y)$ est librement engendré par les
classes de cycles des composantes irréductibles $W \subset Z \times_X Y$
de $\delta$-dimension $n - r$. Soit $\xi \in W$ le point générique.
D'après l'hypothèse (2), nous voyons que $\dim(\mathcal{O}_{Y, \xi}) = r$.
D'autre part, comme $\mathcal{N}$ est de rang $r$ et que
$\mathcal{N}^\vee \to \mathcal{C}_{Z/X}$ est surjectif, nous voyons que
le faisceau d'idéaux de $Z$ est localement défini par $r$ équations.
Ainsi, le faisceau quasi-cohérent d'idéaux $\mathcal{I} \subset \mathcal{O}_Y$
de $Z \times_X Y$ dans $Y$ est localement engendré par $r$ éléments.
Comme $\mathcal{O}_{Y, \xi}$ est de Cohen-Macaulay de dimension $r$
et que $\mathcal{I}_\xi$ est un idéal de définition (puisque $\xi$ est
un point générique de $Z \times_X Y$), il s'ensuit que $\mathcal{I}_\xi$
est engendré par une suite régulière
(Algèbre, lemme \ref{algebra-lemma-reformulate-CM}).
D'après Diviseurs, lemme \ref{divisors-lemma-Noetherian-scheme-regular-ideal},
nous voyons que $\mathcal{I}$ est engendré par une suite régulière sur
un voisinage ouvert $V \subset Y$ de $\xi$. D'après notre description de
$\CH_{n - r}(Z \times_X Y)$, il suffit de montrer que
$c(Z \to X, \mathcal{N}) \cap [V]_n = [Z \times_X V]_{n - r}$
dans $\CH_{n - r}(Z \times_X V)$. Cela résulte du
lemme \ref{lemma-gysin-fundamental},
car le faisceau normal d'excès est $0$ sur $V$.
\end{proof}

\begin{lemma}
\label{lemma-gysin-agrees}
Soit $(S, \delta)$ comme dans la situation \ref{situation-setup}. Soit $X$ un schéma
localement de type fini sur $S$. Soit $(\mathcal{L}, s, i : D \to X)$
un triplet comme dans la définition \ref{definition-gysin-homomorphism}.
L'homomorphisme de Gysin $i^*$, considéré comme un élément de $A^1(D \to X)$
(voir le lemme \ref{lemma-gysin-bivariant}), est égal à la classe bivariante
$c(D \to X, \mathcal{N}) \in A^1(D \to X)$
construite à l'aide de $\mathcal{N} = i^*\mathcal{L}$,
considéré comme un faisceau normal virtuel pour $D$ dans $X$.
\end{lemma}

\begin{proof}
Nous utiliserons le critère du lemme \ref{lemma-bivariant-zero}.
Ainsi, nous pouvons supposer que $X$ est un schéma intègre et
nous devons montrer que $i^*[X]$ est égal à $c \cap [X]$.
Soit $n = \dim_\delta(X)$. Comme d'habitude, il y a deux cas.

\medskip\noindent
Si $X = D$, nous voyons que les deux classes sont représentées par
$c_1(\mathcal{N}) \cap [X]_n$. Voir le lemme \ref{lemma-gysin-fundamental}
et la définition \ref{definition-gysin-homomorphism}.

\medskip\noindent
Si $D \not = X$, alors $D \to X$ est un diviseur de Cartier effectif
et, en particulier, une immersion fermée régulière de codimension $1$.
De nouveau d'après le lemme \ref{lemma-gysin-fundamental}, nous concluons que
$c(D \to X, \mathcal{N}) \cap [X]_n = [D]_{n - 1}$. Par définition, il en va
de même pour l'homomorphisme de Gysin, et nous concluons
une nouvelle fois.
\end{proof}

\begin{lemma}
\label{lemma-gysin-commutes}
Soit $(S, \delta)$ comme dans la situation \ref{situation-setup}. Soit $X$ un schéma
localement de type fini sur $S$. Soit $Z \subset X$ un sous-schéma fermé
muni d'un faisceau normal virtuel $\mathcal{N}$. Soit $Y \to X$ localement de
type fini et soit $c \in A^*(Y \to X)$. Alors $c$ et $c(Z \to X, \mathcal{N})$
commutent (remarque \ref{remark-bivariant-commute}).
\end{lemma}

\begin{proof}
Pour le vérifier, nous pouvons utiliser le lemme \ref{lemma-bivariant-zero}.
Ainsi, nous pouvons supposer que $X$ est un schéma intègre, et nous devons montrer
$c \cap c(Z \to X, \mathcal{N}) \cap [X] =
c(Z \to X, \mathcal{N}) \cap c \cap [X]$ dans $\CH_*(Z \times_X Y)$.

\medskip\noindent
Si $Z = X$, alors $c(Z \to X, \mathcal{N}) = c_{top}(\mathcal{N})$ d'après le
lemme \ref{lemma-gysin-fundamental}, et cette classe commute
avec la classe bivariante $c$, voir le lemme \ref{lemma-cap-commutative-chern}.

\medskip\noindent
Supposons que $Z$ ne soit pas égal à $X$. D'après le lemme \ref{lemma-bivariant-zero},
il suffit même de prouver le résultat après avoir éclaté $X$ (le long d'un idéal non nul).
Éclatons $X$ le long du faisceau d'idéaux de $Z$. Cela nous ramène au cas
où $Z$ est un diviseur de Cartier effectif, voir
Diviseurs, lemme
\ref{divisors-lemma-blowing-up-gives-effective-Cartier-divisor},

\medskip\noindent
Si $Z$ est un diviseur de Cartier effectif, alors nous avons
$$
c(Z \to X, \mathcal{N}) =
c_{top}(\mathcal{E}) \circ i^*
$$
où $i^* \in A^1(Z \to X)$ est l'homomorphisme de Gysin
associé à $i : Z \to X$ (lemme \ref{lemma-gysin-bivariant})
et $\mathcal{E}$ est le dual du noyau de
$\mathcal{N}^\vee \to \mathcal{C}_{Z/X}$, voir les
lemmes \ref{lemma-gysin-decompose} et \ref{lemma-gysin-agrees}.
Nous concluons alors parce que les classes de Chern sont dans le centre de
l'anneau bivariant (au sens fort formulé dans le
lemme \ref{lemma-cap-commutative-chern}) et que $c$ commute
avec l'homomorphisme de Gysin $i^*$ par définition des classes bivariantes.
\end{proof}

\noindent
Soit $(S, \delta)$ comme dans la situation \ref{situation-setup}. Soit $X$ un
schéma intègre localement de type fini sur $S$, de $\delta$-dimension $n$.
Soient $Z \subset Y \subset X$ des sous-schémas fermés qui sont tous deux des
diviseurs de Cartier effectifs dans $X$. Notons $o : Y \to C_Y X$ la section nulle du
cône normal en droites de $Y$ dans $X$. Comme $C_YX$ est un fibré en droites sur $Y$,
nous obtenons une classe bivariante $o^* \in A^1(Y \to C_YX)$, voir le
lemme \ref{lemma-gysin-bivariant}.

\begin{lemma}
\label{lemma-relation-normal-cones}
Avec les notations ci-dessus, nous avons
$$
o^*[C_ZX]_n = [C_Z Y]_{n - 1}
$$
dans $\CH_{n - 1}(Y \times_{o, C_Y X} C_ZX)$.
\end{lemma}

\begin{proof}
Notons $W \to \mathbf{P}^1_X$ l'éclatement de $\infty(Z)$ comme dans la
section \ref{section-blowup-Z-first}.
De même, notons $W' \to \mathbf{P}^1_X$ l'éclatement de $\infty(Y)$.
Comme $\infty(Z) \subset \infty(Y)$, nous obtenons une inclusion opposée
des faisceaux d'idéaux, donc une application des algèbres graduées
définissant ces éclatements. Cela produit un morphisme rationnel de $W$
vers $W'$ qui admet en fait un représentant canonique
$$
W \supset U \longrightarrow W'
$$
Voir Constructions, lemme \ref{constructions-lemma-morphism-relative-proj}.
Un calcul local (omis) montre que $U$ contient au moins tous les points
de $W$ qui ne sont pas au-dessus de $\infty$, ainsi que le sous-schéma ouvert $C_Z X$ de la fibre
spéciale. Quitte à restreindre $U$, nous pouvons supposer que $U_\infty = C_Z X$ et que
$\mathbf{A}^1_X \subset U$. Un autre calcul local (omis)
montre que le morphisme $U_\infty \to W'_\infty$
induit le morphisme canonique $C_Z X \to C_Y X \subset W'_\infty$
de cônes normaux induit par l'inclusion des faisceaux d'idéaux
provenant de $Z \subset Y$. Notons $W'' \subset W$ la transformée stricte de
$\mathbf{P}^1_Y \subset \mathbf{P}^1_X$ dans $W$. Alors $W''$ est l'éclatement
de $\mathbf{P}^1_Y$ en $\infty(Z)$ d'après
Diviseurs, lemme \ref{divisors-lemma-strict-transform},
et donc $(W'' \cap U)_\infty = C_ZY$.

\medskip\noindent
Considérons le diviseur de Cartier effectif $i : \mathbf{P}^1_Y \to W'$
de (\ref{item-find-Z-in-blowup}) et sa classe bivariante associée
$i^* \in A^1(\mathbf{P}^1_Y \to W')$ donnée par le lemme \ref{lemma-gysin-bivariant}.
Notons de même $(i'_\infty)^* \in A^1(W'_\infty \to W')$
l'application de Gysin à l'infini. Remarquons que la restriction de $i'_\infty$
(remarque \ref{remark-restriction-bivariant}) à $U$ est la restriction de
$i_\infty^* \in A^1(W_\infty \to W)$ à $U$. D'une part, nous avons
$$
(i'_\infty)^* i^* [U]_{n + 1} =
i_\infty^* i^* [U]_{n + 1} =
i_\infty^* [(W'' \cap U)_\infty]_{n + 1} =
[C_ZY]_n
$$
parce que $i_\infty^*$ annule toutes les classes à support au-dessus de $\infty$, parce que
$i^*[U]$ et $[W'']$ coïncident comme cycles sur $\mathbf{A}^1$, et parce que
$C_ZY$ est la fibre de $W'' \cap U$ au-dessus de $\infty$.
D'autre part, nous avons
$$
(i'_\infty)^* i^* [U]_{n + 1} =
i^* i_\infty^*[U]_{n + 1} =
i^* [U_\infty] =
o^*[C_YX]_n
$$
parce que $(i'_\infty)^*$ et $i^*$ commutent
(lemme \ref{lemma-gysin-commutes-gysin})
et parce que la fibre de $i : \mathbf{P}^1_Y \to W'$ au-dessus de $\infty$
se factorise par $o : Y \to C_YX$ et l'immersion ouverte $C_YX \to W'_\infty$.
Le lemme en résulte.
\end{proof}

\begin{lemma}
\label{lemma-gysin-composition}
Soit $(S, \delta)$ comme dans la situation \ref{situation-setup}.
Soient $Z \subset Y \subset X$ des sous-schémas fermés d'un schéma localement
de type fini sur $S$.
Soit $\mathcal{N}$ un faisceau normal virtuel pour $Z \subset X$.
Soit $\mathcal{N}'$ un faisceau normal virtuel pour $Z \subset Y$.
Soit $\mathcal{N}''$ un faisceau normal virtuel pour $Y \subset X$.
Supposons qu'il existe un diagramme commutatif
$$
\xymatrix{
(\mathcal{N}'')^\vee|_Z \ar[r] \ar[d] &
\mathcal{N}^\vee \ar[r] \ar[d] &
(\mathcal{N}')^\vee \ar[d] \\
\mathcal{C}_{Y/X}|_Z \ar[r] &
\mathcal{C}_{Z/X} \ar[r] &
\mathcal{C}_{Z/Y}
}
$$
où la suite du bas provient de Compléments sur les morphismes, lemme
\ref{more-morphisms-lemma-transitivity-conormal}, et où la
suite du haut est une suite exacte courte. Alors
$$
c(Z \to X, \mathcal{N}) =
c(Z \to Y, \mathcal{N}') \circ c(Y \to X, \mathcal{N}'')
$$
dans $A^*(Z \to X)^\wedge$.
\end{lemma}

\begin{proof}
Remarquons que les hypothèses restent satisfaites après tout changement de base
par un morphisme $X' \to X$ localement de type fini (la suite
exacte courte de faisceaux normaux virtuels est localement scindée et
reste donc exacte après tout changement de base). Ainsi, pour vérifier
l'égalité des classes bivariantes, nous pouvons utiliser le lemme \ref{lemma-bivariant-zero}.
Nous pouvons donc supposer que $X$ est un schéma intègre, et nous devons montrer
$c(Z \to X, \mathcal{N}) \cap [X] =
c(Z \to Y, \mathcal{N}') \cap c(Y \to X, \mathcal{N}'') \cap [X]$.

\medskip\noindent
Si $Y = X$, alors nous avons
\begin{align*}
c(Z \to Y, \mathcal{N}') \cap c(Y \to X, \mathcal{N}'') \cap [X]
& =
c(Z \to Y, \mathcal{N}') \cap c_{top}(\mathcal{N}'') \cap [Y] \\
& =
c_{top}(\mathcal{N}''|_Z) \cap c(Z \to Y, \mathcal{N}') \cap [Y] \\
& =
c(Z \to X, \mathcal{N}) \cap [X] 
\end{align*}
La première égalité résulte du lemme \ref{lemma-gysin-decompose}.
La deuxième vient de ce que les classes de Chern commutent avec les classes bivariantes
(lemme \ref{lemma-cap-commutative-chern}).
La troisième égalité résulte du lemme \ref{lemma-gysin-decompose}.

\medskip\noindent
Supposons $Y \not = X$. D'après le lemme \ref{lemma-bivariant-zero},
il suffit même de prouver le résultat après avoir éclaté $X$ en un idéal non nul.
Éclatons $X$ le long du produit du faisceau d'idéaux de $Y$ et du faisceau
d'idéaux de $Z$. Cela nous ramène au cas où $Y$ et $Z$ sont tous deux des
diviseurs de Cartier effectifs sur $X$, voir
Diviseurs, lemmes
\ref{divisors-lemma-blowing-up-gives-effective-Cartier-divisor} et
\ref{divisors-lemma-blowing-up-two-ideals}.

\medskip\noindent
Notons $\mathcal{N}'' \to \mathcal{E}$ la surjection de
$\mathcal{O}_Z$-modules localement libres de rang fini telle que
$0 \to \mathcal{E}^\vee \to (\mathcal{N}'')^\vee \to \mathcal{C}_{Y/X} \to 0$
soit une suite exacte courte. Alors $\mathcal{N} \to \mathcal{E}|_Z$
est également une surjection. Notons $\mathcal{N}_1$ le noyau localement libre de rang fini
de cette application et remarquons que $\mathcal{N}^\vee \to \mathcal{C}_{Z/X}$
se factorise par $\mathcal{N}_1$.
D'après le lemme \ref{lemma-gysin-decompose}, nous avons
$$
c(Y \to X, \mathcal{N}'') = c_{top}(\mathcal{E}) \circ
c(Y \to X, \mathcal{C}_{Y/X}^\vee)
$$
et
$$
c(Z \to X, \mathcal{N}) = c_{top}(\mathcal{E}|_Z) \circ
c(Z \to X, \mathcal{N}_1)
$$
Comme les classes de Chern de fibrés commutent avec les classes bivariantes
(lemme \ref{lemma-cap-commutative-chern}),
il suffit de prouver
$$
c(Z \to X, \mathcal{N}_1) =
c(Z \to Y, \mathcal{N}') \circ c(Y \to X, \mathcal{C}_{Y/X}^\vee)
$$
dans $A^*(Z \to X)$. Pour cela, nous pouvons supposer que $\mathcal{N}'' = \mathcal{C}_{Y/X}$.
Cela nous ramène au cas examiné dans l'alinéa suivant.

\medskip\noindent
Dans cet alinéa, $Z$ et $Y$ sont des diviseurs de Cartier effectifs sur $X$,
qui est intègre de dimension $n$, et nous avons $\mathcal{N}'' = \mathcal{C}_{Y/X}$.
Dans ce cas, $c(Y \to X, \mathcal{C}_{Y/X}^\vee) \cap [X] = [Y]_{n - 1}$ d'après le
lemme \ref{lemma-gysin-fundamental}. Nous devons donc prouver que
$c(Z \to X, \mathcal{N}) \cap [X] = c(Z \to Y, \mathcal{N}') \cap [Y]_{n - 1}$.
Notons $N$ et $N'$ les fibrés vectoriels sur $Z$ associés à
$\mathcal{N}$ et $\mathcal{N}'$. Considérons le diagramme commutatif
$$
\xymatrix{
N' \ar[r]_i &
N \ar[r] &
(C_Y X) \times_Y Z \\
C_Z Y \ar[r] \ar[u] &
C_Z X \ar[u]
}
$$
de cônes et de fibrés vectoriels sur $Z$. Remarquons que $N'$ est un
diviseur de Cartier effectif relatif dans $N$ sur $Z$ et que
$$
\xymatrix{
N' \ar[d] \ar[r]_i & N \ar[d] \\
Z \ar[r]^-o & (C_Y X) \times_Y Z
}
$$
est cartésien, où $o$ est la section nulle du fibré en droites
$C_Y X$ sur $Y$. D'après le
lemme \ref{lemma-relation-normal-cones}, nous avons $o^*[C_ZX]_n = [C_Z Y]_{n - 1}$
dans
$$
\CH_{n - 1}(Y \times_{o, C_Y X} C_ZX) =
\CH_{n - 1}(Z \times_{o, (C_Y X) \times_Y Z} C_ZX)
$$
Par la propriété cartésienne de
ce carré, il s'ensuit que
$$
i^*[C_ZX]_n = [C_Z Y]_{n - 1}
$$
dans $\CH_{n - 1}(N')$. Remarquons maintenant que
$\gamma = c(Z \to X, \mathcal{N}) \cap [X]$ et
$\gamma' = c(Z \to Y, \mathcal{N}') \cap [Y]_{n - 1}$
sont caractérisés par $p^*\gamma = [C_Z X]_n$ dans $\CH_n(N)$
et par $(p')^*\gamma' = [C_Z Y]_{n - 1}$ dans $\CH_{n - 1}(N')$.
La démonstration est donc achevée, puisque $i^* \circ p^* = (p')^*$ d'après le
lemme \ref{lemma-relative-effective-cartier}.
\end{proof}

\begin{remark}[Variante pour les immersions]
\label{remark-gysin-for-immersion}
Soit $(S, \delta)$ comme dans la situation \ref{situation-setup}.
Soit $X$ un schéma localement de type fini sur $S$.
Soit $i : Z \to X$ une immersion de schémas.
Dans cette situation,
\begin{enumerate}
\item le faisceau conormal $\mathcal{C}_{Z/X}$
de $Z$ dans $X$ est défini
(Morphismes, définition \ref{morphisms-definition-conormal-sheaf}),
\item nous appelons un couple formé d'un $\mathcal{O}_Z$-module localement libre de rang fini
$\mathcal{N}$ et d'une surjection $\sigma : \mathcal{N}^\vee \to \mathcal{C}_{Z/X}$
un fibré normal virtuel pour l'immersion $Z \to X$,
\item choisissons un sous-schéma ouvert $U \subset X$ tel que $Z \to X$
se factorise par une immersion fermée $Z \to U$ et posons
$c(Z \to X, \mathcal{N}) = c(Z \to U, \mathcal{N}) \circ (U \to X)^*$.
\end{enumerate}
La classe bivariante $c(Z \to X, \mathcal{N})$ ne dépend pas du choix
du sous-schéma ouvert $U$. Tous les lemmes ont des analogues immédiats
pour cette construction légèrement plus générale. Nous omettons les détails.
\end{remark}







\section{Calcul de quelques classes}
\label{section-calculate}

\noindent
Pour aller plus loin, nous devons calculer les valeurs de certaines des
classes construites ci-dessus.

\begin{lemma}
\label{lemma-compute-koszul}
Soit $(S, \delta)$ comme dans la situation \ref{situation-setup}. Soit $X$
un schéma localement de type fini sur $S$. Soit $\mathcal{E}$
un $\mathcal{O}_X$-module localement libre de rang $r$.
Alors
$$
\prod\nolimits_{n = 0, \ldots, r} c(\wedge^n \mathcal{E})^{(-1)^n} =
1 - (r - 1)! c_r(\mathcal{E}) + \ldots
$$
\end{lemma}

\begin{proof}
Par le principe de scindage, nous pouvons ramener ceci à un calcul dans
l'anneau de polynômes aux racines de Chern $x_1, \ldots, x_r$ de $\mathcal{E}$. Voir la
section \ref{section-splitting-principle}. Remarquons que
$$
c(\wedge^n \mathcal{E}) =
\prod\nolimits_{1 \leq i_1 < \ldots < i_n \leq r}
(1 + x_{i_1} + \ldots + x_{i_n})
$$
Ainsi, le logarithme du membre de gauche de l'égalité du lemme est
$$
-
\sum\nolimits_{p \geq 1}
\sum\nolimits_{n = 0}^r
\sum\nolimits_{1 \leq i_1 < \ldots < i_n \leq r}
\frac{(-1)^{p + n}}{p}(x_{i_1} + \ldots + x_{i_n})^p
$$
Remarquons bien le signe moins devant. Cependant, nous avons
$$
\sum\nolimits_{p \geq 0}
\sum\nolimits_{n = 0}^r
\sum\nolimits_{1 \leq i_1 < \ldots < i_n \leq r}
\frac{(-1)^{p + n}}{p!}(x_{i_1} + \ldots + x_{i_n})^p
=
\prod (1 - e^{-x_i})
$$
Nous voyons donc que le premier terme non nul de notre classe de Chern
est de degré $r$ et égal à la valeur annoncée.
\end{proof}

\begin{lemma}
\label{lemma-compute-section}
Soit $(S, \delta)$ comme dans la situation \ref{situation-setup}. Soit $X$
un schéma localement de type fini sur $S$. Soit $\mathcal{C}$
un $\mathcal{O}_X$-module localement libre de rang $r$. Considérons les
morphismes
$$
X = \underline{\text{Proj}}_X(\mathcal{O}_X[T])
\xrightarrow{i}
E = \underline{\text{Proj}}_X(\text{Sym}^*(\mathcal{C})[T])
\xrightarrow{\pi}
X
$$
Alors $c_t(i_*\mathcal{O}_X) = 0$ pour $t = 1, \ldots, r - 1$ et, dans
$A^0(C \to E)$, nous avons
$$
p^* \circ \pi_* \circ c_r(i_*\mathcal{O}_X) = (-1)^{r - 1}(r - 1)! j^*
$$
où
$j : C \to E$ et $p : C \to X$ sont l'inclusion et le morphisme
structural du fibré vectoriel
$C = \underline{\Spec}(\text{Sym}^*(\mathcal{C}))$.
\end{lemma}

\begin{proof}
L'application canonique $\pi^*\mathcal{C} \to \mathcal{O}_E(1)$ s'annule
exactement le long de $i(X)$. Ainsi, le complexe de Koszul de l'application
$$
\pi^*\mathcal{C} \otimes \mathcal{O}_E(-1) \to \mathcal{O}_E
$$
est une résolution de $i_*\mathcal{O}_X$. En particulier, nous voyons que
$i_*\mathcal{O}_X$ est un objet parfait de $D(\mathcal{O}_E)$
dont les classes de Chern sont définies. L'annulation de $c_t(i_*\mathcal{O}_X)$
pour $t = 1, \ldots, t - 1$ résulte du lemme \ref{lemma-compute-koszul}.
Ce lemme donne également
$$
c_r(i_*\mathcal{O}_X) = - (r - 1)!
c_r(\pi^*\mathcal{C} \otimes \mathcal{O}_E(-1))
$$
D'autre part, d'après le lemme \ref{lemma-chern-classes-dual}, nous avons
$$
c_r(\pi^*\mathcal{C} \otimes \mathcal{O}_E(-1)) =
(-1)^r c_r(\pi^*\mathcal{C}^\vee \otimes \mathcal{O}_E(1))
$$
et $\pi^*\mathcal{C}^\vee \otimes \mathcal{O}_E(1)$ possède une section $s$
qui s'annule exactement le long de $i(X)$.

\medskip\noindent
Après avoir remplacé $X$ par un schéma localement de type fini sur $X$,
il suffit de prouver que les deux membres de l'égalité ont le
même effet sur un élément $\alpha \in \CH_*(E)$. Comme $C \to X$
est un fibré vectoriel, toute classe de cycle sur $C$ est de la forme $p^*\beta$
pour un certain $\beta \in \CH_*(X)$ (lemme \ref{lemma-vectorbundle}).
Ainsi, d'après le lemme \ref{lemma-restrict-to-open},
nous pouvons écrire $\alpha = \pi^*\beta + \gamma$, où $\gamma$
est à support dans $E \setminus C$. En utilisant les égalités ci-dessus,
il suffit de montrer que
$$
p^*(\pi_*(c_r(\pi^*\mathcal{C}^\vee \otimes \mathcal{O}_E(1)) \cap [W])) =
j^*[W]
$$
lorsque $W \subset E$ est un sous-schéma fermé intègre qui
soit (a) disjoint de $C$, soit (b) de la forme $W = \pi^{-1}Y$
pour un certain sous-schéma fermé intègre $Y \subset X$.
En utilisant la section $s$ et le lemme \ref{lemma-top-chern-class}, nous trouvons
dans le cas (a) $c_r(\pi^*\mathcal{C}^\vee \otimes \mathcal{O}_E(1)) \cap [W] = 0$
et, dans le cas (b),
$c_r(\pi^*\mathcal{C}^\vee \otimes \mathcal{O}_E(1)) \cap [W] = [i(Y)]$.
Le résultat s'en déduit aisément ; nous omettons les détails.
\end{proof}

\begin{lemma}
\label{lemma-agreement-with-loc-chern}
Soit $(S, \delta)$ comme dans la situation \ref{situation-setup}. Soit $i : Z \to X$
une immersion fermée régulière de codimension $r$
entre des schémas localement de type fini sur $S$.
Soit $\mathcal{N} = \mathcal{C}_{Z/X}^\vee$ le faisceau normal. Si $X$
est quasi-compact (ou possède des composantes irréductibles quasi-compactes), alors
$c_t(Z \to X, i_*\mathcal{O}_Z) = 0$ pour $t = 1, \ldots, r - 1$ et
$$
c_r(Z \to X, i_*\mathcal{O}_Z) = (-1)^{r - 1} (r - 1)! c(Z \to X, \mathcal{N})
\quad\text{dans}\quad
A^r(Z \to X)
$$
où $c_t(Z \to X, i_*\mathcal{O}_Z)$
est la classe de Chern localisée
de la définition \ref{definition-localized-chern}.
\end{lemma}

\begin{proof}
Pour tout $x \in Z$, nous pouvons choisir un voisinage ouvert affine
$\Spec(A) \subset X$ tel que $Z \cap \Spec(A) = V(f_1, \ldots, f_r)$,
où $f_1, \ldots, f_r \in A$ est une suite régulière.
Voir Diviseurs, définition \ref{divisors-definition-regular-immersion} et
lemme \ref{divisors-lemma-Noetherian-scheme-regular-ideal}.
Nous voyons alors que le complexe de Koszul de $f_1, \ldots, f_r$ est
une résolution de $A/(f_1, \ldots, f_r)$, par exemple d'après
Compléments d'algèbre, lemme \ref{more-algebra-lemma-regular-koszul-regular}.
Ainsi, $A/(f_1, \ldots, f_r)$ est parfait comme $A$-module.
Il s'ensuit que $F = i_*\mathcal{O}_Z$ est un objet parfait de
$D(\mathcal{O}_X)$ dont la restriction à $X \setminus Z$ est nulle.
L'hypothèse que $X$ est quasi-compact (ou possède des composantes
irréductibles quasi-compactes) signifie que les classes de Chern localisées
$c_t(Z \to X, i_*\mathcal{O}_Z)$ sont définies, voir
la situation \ref{situation-loc-chern} et la
définition \ref{definition-localized-chern}. En définitive,
nous concluons que l'énoncé a un sens.

\medskip\noindent
Notons $b : W \to \mathbf{P}^1_X$ l'éclatement en $\infty(Z)$
comme dans la section \ref{section-blowup-Z-first}. D'après (\ref{item-find-Z-in-blowup}),
nous avons une immersion fermée
$$
i' : \mathbf{P}^1_Z \longrightarrow W
$$
Nous affirmons que $Q = i'_*\mathcal{O}_{\mathbf{P}^1_Z}$
est un objet parfait de
$D(\mathcal{O}_W)$ et que $F$ et $Q$ satisfont aux hypothèses du
lemme \ref{lemma-independent-loc-chern-bQ}.

\medskip\noindent
Supposons l'affirmation démontrée. La conclusion du lemme \ref{lemma-independent-loc-chern-bQ}
est que nous avons
$$
c_p(Z \to X, F) = c'_p(Q) = (E \to Z)_* \circ c'_p(Q|_E) \circ C
$$
pour tout $p \geq 1$. Remarquons que $Q|_E$ est égal à l'image directe du
faisceau structural de $Z$ par le morphisme $Z \to E$ qui est le
changement de base de $i'$ par $\infty$.
Ainsi, l'annulation de $c_t(Z \to X, F)$ pour $1 \leq t \leq r - 1$
résulte du lemme \ref{lemma-compute-section} appliqué à $E \to Z$.
Comme $\mathcal{C}_{Z/X} = \mathcal{N}^\vee$
est localement libre, la classe bivariante $c(Z \to X, \mathcal{N})$
est caractérisée par la relation
$$
j^* \circ C = p^* \circ c(Z \to X, \mathcal{N})
$$
où $j : C_ZX \to W_\infty$ et $p : C_ZX \to Z$ sont les applications données.
(Rappelons que $C \in A^0(W_\infty \to X)$ est la classe du
lemme \ref{lemma-gysin-at-infty}.)
Ainsi, l'égalité affichée dans l'énoncé du lemme
résulte de l'égalité correspondante du lemme \ref{lemma-compute-section}.

\medskip\noindent
Démonstration de l'affirmation. Soient $A$ et $f_1, \ldots, f_r$ comme ci-dessus.
Considérons l'ouvert affine $\Spec(A[s]) \subset \mathbf{P}^1_X$
comme dans la section \ref{section-blowup-Z-first}. Rappelons que $s = 0$
définit $(\mathbf{P}^1_X)_\infty$ sur cet ouvert. Ainsi, au-dessus de
$\Spec(A[s])$, nous effectuons l'éclatement de l'idéal engendré par
la suite régulière $s, f_1, \ldots, f_r$. D'après Compléments d'algèbre, lemme
\ref{more-algebra-lemma-blowup-regular-sequence}, les $r + 1$
cartes affines sont des intersections complètes globales sur $A[s]$.
La carte correspondant à l'algèbre affine d'éclatement
$$
A[s][f_1/s, \ldots, f_r/s] = A[s, y_1, \ldots, y_r]/(sy_i - f_i)
$$
contient $i'(Z \cap \Spec(A))$ comme sous-schéma fermé défini par
$y_1, \ldots, y_r$. Comme $y_1, \ldots, y_r, sy_1 - f_1, \ldots, sy_r - f_r$
est une suite régulière dans l'anneau de polynômes $A[s, y_1, \ldots, y_r]$,
nous trouvons que $i'$ est une immersion régulière. Nous omettons quelques détails.
Comme ci-dessus, nous concluons que $Q = i'_*\mathcal{O}_{\mathbf{P}^1_Z}$
est un objet parfait de $D(\mathcal{O}_W)$. Toutes les
autres hypothèses sur $F$ et $Q$ du lemme \ref{lemma-independent-loc-chern-bQ}
(et du lemme \ref{lemma-localized-chern-pre}) se vérifient immédiatement.
\end{proof}

\begin{lemma}
\label{lemma-actual-computation}
Dans la situation du lemme \ref{lemma-agreement-with-loc-chern},
posons $\dim_\delta(X) = n$. Alors nous avons
\begin{enumerate}
\item $c_t(Z \to X, i_*\mathcal{O}_Z) \cap [X]_n = 0$ pour
$t = 1, \ldots, r - 1$,
\item $c_r(Z \to X, i_*\mathcal{O}_Z) \cap [X]_n =
(-1)^{r - 1}(r - 1)![Z]_{n - r}$,
\item $ch_t(Z \to X, i_*\mathcal{O}_Z) \cap [X]_n = 0$ pour
$t = 0, \ldots, r - 1$, et
\item $ch_r(Z \to X, i_*\mathcal{O}_Z) \cap [X]_n = [Z]_{n - r}$.
\end{enumerate}
\end{lemma}

\begin{proof}
Les parties (1) et (2) résultent immédiatement du
lemme \ref{lemma-agreement-with-loc-chern}
combiné au lemme \ref{lemma-gysin-fundamental}.
Nous en déduisons ensuite les parties (3) et (4) à l'aide de la relation
entre $ch_p = (1/p!)P_p$ et $c_p$ donnée dans le
lemme \ref{lemma-loc-chern-character}. (En effet,
$(-1)^{r - 1}(r - 1)!ch_r = c_r$ pourvu que
$c_1 = c_2 = \ldots = c_{r - 1} = 0$.)
\end{proof}







\section{Un opérateur d'Adams}
\label{section-adams}

\noindent
Nous effectuons le travail minimal nécessaire pour définir le deuxième opérateur d'Adams.
Soit $X$ un schéma. Rappelons que $\textit{Vect}(X)$ désigne la
catégorie des $\mathcal{O}_X$-modules localement libres de rang fini.
Rappelons en outre que nous avons construit le groupe $K$ d'indice zéro
$K_0(\textit{Vect}(X))$ associé à cette catégorie dans
Catégories dérivées de schémas, section \ref{perfect-section-K-groups}.
Enfin, $K_0(\textit{Vect}(X))$ est un anneau, voir
Catégories dérivées de schémas, remarque \ref{perfect-remark-K-ring}.

\begin{lemma}
\label{lemma-second-adams-operator}
Soit $X$ un schéma. Il existe un homomorphisme d'anneaux
$$
\psi^2 :
K_0(\textit{Vect}(X))
\longrightarrow
K_0(\textit{Vect}(X))
$$
qui envoie $[\mathcal{L}]$ sur $[\mathcal{L}^{\otimes 2}]$
lorsque $\mathcal{L}$ est inversible et qui est compatible aux images inverses.
\end{lemma}

\begin{proof}
Soit $X$ un schéma.
Soit $\mathcal{E}$ un $\mathcal{O}_X$-module localement libre de rang fini.
Nous allons considérer l'élément
$$
\psi^2(\mathcal{E}) = [\text{Sym}^2(\mathcal{E})] - [\wedge^2(\mathcal{E})]
$$
de $K_0(\textit{Vect}(X))$.

\medskip\noindent
Soit $X$ un schéma et considérons une suite exacte courte
$$
0 \to \mathcal{E} \to \mathcal{F} \to \mathcal{G} \to 0
$$
de $\mathcal{O}_X$-modules localement libres de rang fini. Considérons-la
comme une filtration de $\mathcal{F}$ à $2$ crans. La filtration
induite sur $\text{Sym}^2(\mathcal{F})$ possède $3$ crans, de
gradués associés $\text{Sym}^2(\mathcal{E})$, $\mathcal{E} \otimes \mathcal{F}$
et $\text{Sym}^2(\mathcal{G})$. Ainsi
$$
[\text{Sym}^2(\mathcal{F})] =
[\text{Sym}^2(\mathcal{E})] +
[\mathcal{E} \otimes \mathcal{F}] +
[\text{Sym}^2(\mathcal{G})]
$$
Exactement de la même manière, on montre que
$$
[\wedge^2(\mathcal{F})] =
[\wedge^2(\mathcal{E})] +
[\mathcal{E} \otimes \mathcal{F}] +
[\wedge^2(\mathcal{G})]
$$
Nous voyons donc que
$\psi^2(\mathcal{F}) = \psi^2(\mathcal{E}) + \psi^2(\mathcal{G})$.
Nous concluons que nous obtenons une application additive bien définie
$\psi^2 : K_0(\textit{Vect}(X)) \to K_0(\textit{Vect}(X))$.

\medskip\noindent
Il est clair que cette application commute aux images inverses.

\medskip\noindent
Il reste à montrer que $\psi^2$ est un homomorphisme d'anneaux.
Soit $X$ un schéma et soient $\mathcal{E}$ et $\mathcal{F}$ des
$\mathcal{O}_X$-modules localement libres de rang fini.
Remarquons qu'il existe une suite exacte courte
$$
0 \to \wedge^2(\mathcal{E}) \otimes \wedge^2(\mathcal{F}) \to
\text{Sym}^2(\mathcal{E} \otimes \mathcal{F}) \to
\text{Sym}^2(\mathcal{E}) \otimes \text{Sym}^2(\mathcal{F}) \to 0
$$
où la première application envoie $(e \wedge e') \otimes (f \wedge f')$ sur
$(e \otimes f)(e' \otimes f') - (e' \otimes f)(e \otimes f')$ et où
la seconde application envoie $(e \otimes f) (e' \otimes f')$ sur $ee' \otimes ff'$.
De même, il existe une suite exacte courte
$$
0 \to \text{Sym}^2(\mathcal{E}) \otimes \wedge^2(\mathcal{F}) \to
\wedge^2(\mathcal{E} \otimes \mathcal{F}) \to
\wedge^2(\mathcal{E}) \otimes \text{Sym}^2(\mathcal{F}) \to 0
$$
où la première application envoie $e e' \otimes f \wedge f'$ sur
$(e \otimes f) \wedge (e' \otimes f') + (e' \otimes f) \wedge (e \otimes f')$
et où la seconde application envoie
$(e \otimes f) \wedge (e' \otimes f')$ sur
$(e \wedge e') \otimes (f f')$.
Comme ci-dessus, cela prouve que l'application $\psi^2$ est multiplicative.
Puisqu'il est clair que $\psi^2(1) = 1$, cela achève la démonstration.
\end{proof}

\begin{remark}
\label{remark-adams-derived}
Soit $X$ un schéma tel que $2$ soit inversible sur $X$.
Alors l'opérateur d'Adams $\psi^2$ peut être défini sur le groupe $K$
$K_0(X) = K_0(D_{perf}(\mathcal{O}_X))$
(Catégories dérivées de schémas, définition \ref{perfect-definition-K-group})
d'une manière immédiate.
En effet, étant donné un complexe parfait $L$ sur $X$, nous obtenons une action
du groupe $\{\pm 1\}$ sur $L \otimes^\mathbf{L} L$ par permutation
des facteurs. Nous pouvons alors poser
$$
\psi^2(L) = [(L \otimes^\mathbf{L} L)^+] -
[(L \otimes^\mathbf{L} L)^-]
$$
où $(-)^+$ désigne la prise des invariants et $(-)^-$ celle des
anti-invariants (convenablement définis).
En utilisant l'exactitude de la prise des invariants et des anti-invariants, on peut
raisonner comme dans la démonstration du lemme \ref{lemma-second-adams-operator}
pour montrer que cette définition est bonne.
Lorsque $2$ n'est pas inversible sur $X$, la situation est beaucoup plus
compliquée et une autre approche doit être employée.
\end{remark}

\begin{lemma}
\label{lemma-minus-adams-operator}
Soit $X$ un schéma. Il existe un homomorphisme d'anneaux
$\psi^{-1} : K_0(\textit{Vect}(X)) \to K_0(\textit{Vect}(X))$
qui envoie $[\mathcal{E}]$ sur $[\mathcal{E}^\vee]$
lorsque $\mathcal{E}$ est localement libre de rang fini
et qui est compatible aux images inverses.
\end{lemma}

\begin{proof}
Il suffit de vérifier que la prise des duaux est compatible aux
suites exactes courtes et aux images inverses. C'est clair.
\end{proof}

\begin{remark}
\label{remark-chern-classes-K}
Soit $(S, \delta)$ comme dans la situation \ref{situation-setup}.
Soit $X$ localement de type fini sur $S$. La classe de Chern
définit une application canonique
$$
c : K_0(\textit{Vect}(X)) \longrightarrow \prod\nolimits_{i \geq 0} A^i(X)
$$
qui envoie un générateur $[\mathcal{E}]$ du membre de gauche sur
$c(\mathcal{E}) = 1 + c_1(\mathcal{E}) + c_2(\mathcal{E}) + \ldots$
et se prolonge multiplicativement. Ainsi, $-[\mathcal{E}]$ est envoyé sur
l'inverse formel $c(\mathcal{E})^{-1}$, ce qui explique le
produit infini du membre de droite. Cette application est bien définie d'après le
lemme \ref{lemma-additivity-chern-classes}.
\end{remark}

\begin{remark}
\label{remark-chern-character-K}
Soit $(S, \delta)$ comme dans la situation \ref{situation-setup}.
Soit $X$ localement de type fini sur $S$. Le caractère de Chern
définit un homomorphisme d'anneaux canonique
$$
ch : K_0(\textit{Vect}(X)) \longrightarrow
\prod\nolimits_{i \geq 0} A^i(X) \otimes \mathbf{Q}
$$
qui envoie un générateur $[\mathcal{E}]$ du membre de gauche sur
$ch(\mathcal{E})$ et se prolonge additivement. Cette application est bien définie
d'après le lemme \ref{lemma-chern-character-additive} et est un homomorphisme d'anneaux d'après le
lemme \ref{lemma-chern-character-multiplicative}.
\end{remark}

\begin{lemma}
\label{lemma-adams-and-chern}
Soit $(S, \delta)$ comme dans la situation \ref{situation-setup}.
Soit $X$ localement de type fini sur $S$. Si $\psi^2$ est
comme dans le lemme \ref{lemma-second-adams-operator}, et si $c$ et $ch$ sont comme dans les
remarques \ref{remark-chern-classes-K} et \ref{remark-chern-character-K},
alors nous avons $c_i(\psi^2(\alpha)) = 2^i c_i(\alpha)$ et
$ch_i(\psi^2(\alpha)) = 2^i ch_i(\alpha)$
pour tout $\alpha \in K_0(\textit{Vect}(X))$.
\end{lemma}

\begin{proof}
Remarquons que l'application $\prod_{i \geq 0} A^i(X) \to \prod_{i \geq 0} A^i(X)$
qui multiplie par $2^i$ sur $A^i(X)$ est un homomorphisme d'anneaux. Ainsi, puisque $\psi^2$
est également un homomorphisme d'anneaux, il suffit de prouver les formules pour des générateurs additifs
de $K_0(\textit{Vect}(X))$. Nous pouvons donc supposer que $\alpha = [\mathcal{E}]$
pour un certain $\mathcal{O}_X$-module localement libre de rang fini $\mathcal{E}$.
Par la construction des classes de Chern de $\mathcal{E}$, nous nous
ramenons immédiatement au cas où $\mathcal{E}$ est de rang constant $r$, voir la
remarque \ref{remark-extend-to-finite-locally-free}.
Dans ce cas, nous pouvons choisir un morphisme projectif lisse $p : P \to X$
tel que la restriction $A^*(X) \to A^*(P)$ soit injective
et que $p^*\mathcal{E}$ possède une filtration finie dont les
gradués associés sont des $\mathcal{O}_P$-modules inversibles $\mathcal{L}_j$, voir le
lemme \ref{lemma-splitting-principle}. Alors
$[p^*\mathcal{E}] = \sum [\mathcal{L}_j]$ et donc
$\psi^2([p^\mathcal{E}]) = \sum [\mathcal{L}_j^{\otimes 2}]$
par définition de $\psi^2$. En posant $x_j  = c_1(\mathcal{L}_j)$,
nous avons
$$
c(\alpha) = \prod (1 + x_j)
\quad\text{et}\quad
c(\psi^2(\alpha)) = \prod (1 + 2 x_j)
$$
dans $\prod A^i(P)$, et nous avons
$$
ch(\alpha) = \sum \exp(x_j)
\quad\text{et}\quad
ch(\psi^2(\alpha)) = \sum \exp(2 x_j)
$$
dans $\prod A^i(P)$. Le résultat voulu résulte de ces formules.
\end{proof}

\begin{remark}
\label{remark-perf-Z-cohomology-K}
Soit $X$ un schéma localement noethérien.
Soit $Z \subset X$ un sous-schéma fermé. Considérons la sous-catégorie triangulée
strictement pleine et saturée
$$
D_{Z, perf}(\mathcal{O}_X) \subset D(\mathcal{O}_X)
$$
formée des complexes parfaits de $\mathcal{O}_X$-modules
dont les faisceaux de cohomologie ont leur support ensembliste dans $Z$.
Notons $\textit{Coh}_Z(X) \subset \textit{Coh}(X)$
la sous-catégorie de Serre des $\mathcal{O}_X$-modules cohérents dont le support
ensembliste est contenu dans $Z$. Remarquons que, pour
$E \in D_{Z, perf}(\mathcal{O}_X)$, localement pour la topologie de Zariski sur $X$,
seul un nombre fini des faisceaux de cohomologie $H^i(E)$ sont non nuls
(et ils ont tous leur support ensembliste dans $Z$).
Nous pouvons donc définir
$$
K_0(D_{Z, perf}(\mathcal{O}_X))
\longrightarrow
K_0(\textit{Coh}_Z(X)) = K'_0(Z)
$$
(égalité d'après le lemme \ref{lemma-K-coherent-supported-on-closed}) par la règle
$$
E \longmapsto
[\bigoplus\nolimits_{i \in \mathbf{Z}} H^{2i}(E)] -
[\bigoplus\nolimits_{i \in \mathbf{Z}} H^{2i + 1}(E)]
$$
Cela fonctionne parce que, pour tout triangle distingué de
$D_{Z, perf}(\mathcal{O}_X)$, nous avons une suite exacte longue de
faisceaux de cohomologie.
\end{remark}

\begin{remark}
\label{remark-perf-Z-regular}
Soient $X$, $Z$, $D_{Z, perf}(\mathcal{O}_X)$ comme dans la
remarque \ref{remark-perf-Z-cohomology-K}. Supposons que $X$ soit régulier.
Il existe alors une application canonique
$$
K_0(\textit{Coh}(Z)) \longrightarrow K_0(D_{Z, perf}(\mathcal{O}_X))
$$
définie comme suit. Pour tout $\mathcal{O}_Z$-module cohérent
$\mathcal{F}$, notons $\mathcal{F}[0]$ l'objet de $D(\mathcal{O}_X)$
qui a $\mathcal{F}$ en degré $0$ et est nul dans les autres degrés.
Alors $\mathcal{F}[0]$ est un complexe parfait sur $X$ d'après
Catégories dérivées de schémas, lemme \ref{perfect-lemma-perfect-on-regular}.
Ainsi, $\mathcal{F}[0]$ est un objet de $D_{Z, perf}(\mathcal{O}_X)$.
D'autre part, étant donnée une suite exacte courte
$0 \to \mathcal{F} \to \mathcal{F}' \to \mathcal{F}'' \to 0$ de
$\mathcal{O}_Z$-modules cohérents, nous obtenons un triangle distingué
$\mathcal{F}[0] \to \mathcal{F}'[0] \to \mathcal{F}''[0] \to \mathcal{F}[1]$,
voir Catégories dérivées, section \ref{derived-section-canonical-delta-functor}.
Cela montre que nous obtenons une application
$K_0(\textit{Coh}(Z)) \to K_0(D_{Z, perf}(\mathcal{O}_X))$
en envoyant $[\mathcal{F}]$ sur $[\mathcal{F}[0]]$,
avec nos excuses pour cette affreuse notation.
\end{remark}

\begin{lemma}
\label{lemma-perf-Z-regular}
Soit $X$ un schéma noethérien régulier.
Soit $Z \subset X$ un sous-schéma fermé. Les applications construites
dans les remarques \ref{remark-perf-Z-cohomology-K} et
\ref{remark-perf-Z-regular} sont inverses l'une de l'autre, et nous obtenons
$K'_0(Z) = K_0(D_{Z, perf}(\mathcal{O}_X))$.
\end{lemma}

\begin{proof}
Il est clair que la composée
$$
K_0(\textit{Coh}(Z)) \longrightarrow
K_0(D_{Z, perf}(\mathcal{O}_X)) \longrightarrow
K_0(\textit{Coh}(Z))
$$
est l'identité. Il suffit donc de montrer que la première flèche est
surjective. Soit $E$ un objet de $D_{Z, perf}(\mathcal{O}_X)$.
Rappelons que $D_{perf}(\mathcal{O}_X) = D^b_{\textit{Coh}}(\mathcal{O}_X)$
d'après Catégories dérivées de schémas, lemme
\ref{perfect-lemma-perfect-on-regular}.
Ainsi, les cohomologies $H^i(E)$ sont cohérentes, peuvent être considérées
comme des objets de $D_{Z, perf}(\mathcal{O}_X)$ et seul un nombre fini
d'entre elles sont non nulles. En utilisant les triangles distingués des troncations canoniques, le
lecteur voit que
$$
[E] = \sum (-1)^i[H^i(E)[0]]
$$
dans $K_0(D_{Z, perf}(\mathcal{O}_X))$. Il suffit alors de
montrer que $[\mathcal{F}[0]]$ est dans l'image de l'application
pour tout $\mathcal{O}_X$-module cohérent dont le support ensembliste
est contenu dans $Z$. Puisque nous pouvons trouver une filtration finie de
$\mathcal{F}$ dont les sous-quotients sont des $\mathcal{O}_Z$-modules,
la démonstration est achevée.
\end{proof}

\begin{remark}
\label{remark-localized-chern-classes-K}
Soit $(S, \delta)$ comme dans la situation \ref{situation-setup}.
Soit $X$ localement de type fini sur $S$.
Soit $Z \subset X$ un sous-schéma fermé et soit
$D_{Z, perf}(\mathcal{O}_X)$ comme dans la
remarque \ref{remark-perf-Z-cohomology-K}.
Si $X$ est quasi-compact (ou, plus généralement, si les composantes
irréductibles de $X$ sont quasi-compactes), alors
les classes de Chern localisées définissent une application canonique
$$
c(Z \to X, -) : K_0(D_{Z, perf}(\mathcal{O}_X)) \longrightarrow
A^0(X) \times \prod\nolimits_{i \geq 1} A^i(Z \to X)
$$
qui envoie un générateur $[E]$ du membre de gauche sur
$$
c(Z \to X, E) = 1 + c_1(Z \to X, E) + c_2(Z \to X, E) + \ldots
$$
et se prolonge multiplicativement (le produit du membre de
droite étant celui de la remarque \ref{remark-ring-loc-classes}).
La condition de quasi-compacité sur $X$ garantit que les
classes de Chern localisées sont définies (situation \ref{situation-loc-chern} et
définition \ref{definition-localized-chern})
et que ces classes de Chern localisées transforment les triangles distingués en
produits correspondants dans les anneaux de Chow bivariants
(lemme \ref{lemma-additivity-loc-chern-c}).
\end{remark}

\begin{remark}
\label{remark-localized-chern-character-K}
Soit $(S, \delta)$ comme dans la situation \ref{situation-setup}.
Soit $X$ localement de type fini sur $S$.
Soit $Z \subset X$ un sous-schéma fermé et soit
$D_{Z, perf}(\mathcal{O}_X)$ comme dans la
remarque \ref{remark-perf-Z-cohomology-K}.
Si les composantes irréductibles de $X$ sont quasi-compactes, alors
le caractère de Chern localisé définit une application canonique additive
et multiplicative
$$
ch(Z \to X, -) : K_0(D_{Z, perf}(\mathcal{O}_X)) \longrightarrow
\prod\nolimits_{i \geq 0} A^i(Z \to X) \otimes \mathbf{Q}
$$
qui envoie un générateur $[E]$ du membre de gauche sur
$ch(Z \to X, E)$ et se prolonge additivement.
En effet, la condition sur les composantes irréductibles de $X$ garantit que le
caractère de Chern localisé est défini (situation \ref{situation-loc-chern} et
définition \ref{definition-localized-chern})
et que ces caractères de Chern localisés
transforment les triangles distingués en les
sommes correspondantes dans les anneaux de Chow bivariants
(lemme \ref{lemma-additivity-loc-chern-P}).
La multiplication sur
$K_0(D_{Z, perf}(X))$ est définie à l'aide du produit tensoriel dérivé
(Catégories dérivées de schémas, remarque \ref{perfect-remark-perf-Z}) ;
ainsi, $ch(Z \to X, \alpha \beta) = ch(Z \to X, \alpha) ch(Z \to X, \beta)$ d'après le
lemme \ref{lemma-loc-chern-tensor-product}.
Si $X$ est quasi-compact, alors l'image de l'application $ch(Z \to X, -)$ est
contenue dans $A^*(Z \to X) \otimes \mathbf{Q}$ ; nous omettons
les détails.
\end{remark}

\begin{remark}
\label{remark-chern-classes-agree}
Soit $(S, \delta)$ comme dans la situation \ref{situation-setup}.
Soit $X$ localement de type fini sur $S$ et supposons que $X$
soit quasi-compact (ou, plus généralement, que les composantes irréductibles
de $X$ soient quasi-compactes). Avec $Z = X$ et les notations des
remarques \ref{remark-localized-chern-classes-K} et
\ref{remark-localized-chern-character-K},
nous avons $D_{Z, perf}(\mathcal{O}_X) = D_{perf}(\mathcal{O}_X)$,
et nous voyons que
$$
K_0(D_{Z, perf}(\mathcal{O}_X)) = K_0(D_{perf}(\mathcal{O}_X)) = K_0(X)
$$
voir
Catégories dérivées de schémas, définition \ref{perfect-definition-K-group}.
Nous obtenons donc
$$
c : K_0(X) \to \prod A^i(X)
\quad\text{et}\quad
ch : K_0(X) \to \prod A^i(X) \otimes \mathbf{Q}
$$
comme cas particulier des remarques \ref{remark-localized-chern-classes-K} et
\ref{remark-localized-chern-character-K}. Bien sûr, nous aurions aussi pu
utiliser directement la définition \ref{definition-defined-on-perfect} et les
lemmes \ref{lemma-additivity-on-perfect} et
\ref{lemma-chern-classes-perfect-tensor-product} pour construire ces applications
(car on voit immédiatement que cela donne les mêmes classes).
Rappelons qu'il existe une application canonique $K_0(\textit{Vect}(X)) \to K_0(X)$
qui envoie un module localement libre de rang fini sur lui-même, considéré
comme un complexe parfait (placé en degré $0$), voir
Catégories dérivées de schémas, section \ref{perfect-section-K-groups}.
Alors le diagramme
$$
\xymatrix{
K_0((\textit{Vect}(X)) \ar[rd]_c \ar[rr] & &
K_0(D_{perf}(\mathcal{O}_X)) = K_0(X) \ar[ld]^c \\
& \prod A^i(X)
}
$$
est commutatif, où la flèche sud-est est celle construite dans la
remarque \ref{remark-chern-classes-K}. De même, le diagramme
$$
\xymatrix{
K_0((\textit{Vect}(X)) \ar[rd]_{ch} \ar[rr] & &
K_0(D_{perf}(\mathcal{O}_X)) = K_0(X) \ar[ld]^{ch} \\
& \prod A^i(X) \otimes \mathbf{Q}
}
$$
est commutatif, où la flèche sud-est est celle construite dans la
remarque \ref{remark-chern-character-K}.
\end{remark}











\section{Retour sur les groupes de Chow et les groupes K}
\label{section-chow-and-K-II}

\noindent
Cette section fait suite à la section \ref{section-chow-and-K}.
Soit $(S, \delta)$ comme dans la situation \ref{situation-setup}.
Soit $X$ localement de type fini sur $S$. Le groupe K
$K'_0(X) = K_0(\textit{Coh}(X))$ des faisceaux cohérents sur $X$
possède une filtration croissante canonique
$$
F_kK'_0(X) =
\Im\Big(K_0(\textit{Coh}_{\leq k}(X)) \to K_0(\textit{Coh}(X)\Big)
$$
Celle-ci est appelée filtration par la dimension des supports. Remarquons que
$$
\text{gr}_k K'_0(X) \subset K'_0(X)/F_{k - 1}K'_0(X) =
K_0(\textit{Coh}(X)/\textit{Coh}_{\leq k - 1}(X))
$$
où l'égalité est valable
d'après Homologie, lemme \ref{homology-lemma-serre-subcategory-K-groups}.
La discussion de la remarque \ref{remark-good-cases-K-A} montre
qu'il existe des applications canoniques
$$
\CH_k(X) \longrightarrow \text{gr}_k K'_0(X)
$$
définies en envoyant la classe d'un sous-schéma fermé intègre
$Z \subset X$ de $\delta$-dimension $k$ sur la classe de
$[\mathcal{O}_Z]$ dans le membre de droite.

\begin{proposition}
\label{proposition-K-tensor-Q}
Soit $(S, \delta)$ comme dans la situation \ref{situation-setup}. Supposons donnée une
immersion fermée $X \to Y$ de schémas localement de type fini sur $S$,
avec $Y$ régulier et quasi-compact. Alors la composée
$$
K'_0(X) \to
K_0(D_{X, perf}(\mathcal{O}_Y)) \to
A^*(X \to Y) \otimes \mathbf{Q} \to
\CH_*(X) \otimes \mathbf{Q}
$$
de l'application $\mathcal{F} \mapsto \mathcal{F}[0]$ de la
remarque \ref{remark-perf-Z-regular}, de l'application $ch(X \to Y, -)$ de la
remarque \ref{remark-localized-chern-character-K} et
de l'application $c \mapsto c \cap [Y]$ induit un isomorphisme
$$
K'_0(X) \otimes \mathbf{Q}
\longrightarrow
\CH_*(X) \otimes \mathbf{Q}
$$
qui dépend du choix de $Y$. De plus, l'application canonique
$$
\CH_k(X) \otimes \mathbf{Q}
\longrightarrow
\text{gr}_k K'_0(X) \otimes \mathbf{Q}
$$
(voir ci-dessus) est un isomorphisme d'espaces vectoriels sur $\mathbf{Q}$ pour tout
$k \in \mathbf{Z}$.
\end{proposition}

\begin{proof}
Comme $Y$ est régulier, la construction de la
remarque \ref{remark-perf-Z-regular} s'applique.
Comme $Y$ est quasi-compact, la construction de la
remarque \ref{remark-localized-chern-character-K} s'applique.
Le schéma $Y$ est localement équidimensionnel
(lemme \ref{lemma-locally-equidimensional}) ; ainsi,
le « cycle fondamental » $[Y]$ est défini
comme un élément de $\CH_*(Y)$, voir la remarque \ref{remark-fundamental-class}.
En combinant ceci avec l'application $\CH_k(X) \to \text{gr}_kK'_0(X)$
construite ci-dessus, nous voyons qu'il suffit de prouver :
\begin{enumerate}
\item si $\mathcal{F}$ est un $\mathcal{O}_X$-module cohérent
dont le support est de $\delta$-dimension $\leq k$, alors
la composée ci-dessus envoie $[\mathcal{F}]$ dans
$\bigoplus_{k' \leq k} \CH_{k'}(X) \otimes \mathbf{Q}$ ;
\item si $Z \subset X$ est un sous-schéma fermé intègre
de $\delta$-dimension $k$, alors la composée ci-dessus
envoie $[\mathcal{O}_Z]$ sur un élément dont la composante de degré $k$
est la classe de $[Z]$ dans $\CH_k(X) \otimes \mathbf{Q}$.
\end{enumerate}
En effet, si cela vaut, nos applications induisent des applications
$\text{gr}_kK'_0(X) \otimes \mathbf{Q} \to CH_k(X) \otimes \mathbf{Q}$
qui sont inverses des applications canoniques
$\CH_k(X) \otimes \mathbf{Q} \to \text{gr}_k K'_0(X) \otimes \mathbf{Q}$
données avant la proposition.

\medskip\noindent
Étant donné un $\mathcal{O}_X$-module cohérent $\mathcal{F}$,
la composée ci-dessus envoie $[\mathcal{F}]$ sur
$$
ch(X \to Y, \mathcal{F}[0]) \cap [Y] \in \CH_*(X) \otimes \mathbf{Q}
$$
Si $\mathcal{F}$ a son support ensembliste dans un sous-schéma fermé
$Z \subset X$, alors nous avons
$$
ch(X \to Y, \mathcal{F}[0]) = (Z \to X)_* \circ ch(Z \to Y, \mathcal{F}[0])
$$
d'après le lemme \ref{lemma-loc-chern-shrink-Z}. Nous concluons que, dans ce
cas, nous aboutissons dans l'image de $\CH_*(Z) \to \CH_*(X)$. Nous obtenons donc
la condition (1).

\medskip\noindent
Soit $Z \subset X$ un sous-schéma fermé intègre de $\delta$-dimension $k$.
La composée ci-dessus envoie $[\mathcal{O}_Z]$ sur l'élément
$$
ch(X \to Y, \mathcal{O}_Z[0]) \cap [Y] =
(Z \to X)_* ch(Z \to Y, \mathcal{O}_Z[0]) \cap [Y]
$$
par le même argument que ci-dessus.
Il suffit donc de prouver que la composante de degré $k$ de
$ch(Z \to Y, \mathcal{O}_Z[0]) \cap [Y] \in
\CH_*(Z) \otimes \mathbf{Q}$ est $[Z]$.
Comme $\CH_k(Z) = \mathbf{Z}$, pour
le prouver, nous pouvons remplacer $Y$ par un voisinage ouvert du
point générique $\xi$ de $Z$. Comme l'idéal maximal de l'anneau
local régulier $\mathcal{O}_{X, \xi}$ est engendré par une
suite régulière (Algèbre, lemme \ref{algebra-lemma-regular-ring-CM}),
nous pouvons supposer que l'idéal de $Z$ est engendré par une suite régulière, voir
Diviseurs, lemme \ref{divisors-lemma-Noetherian-scheme-regular-ideal}.
Nous déduisons donc le résultat du lemme \ref{lemma-actual-computation}.
\end{proof}







\section{Produits d'intersection rationnels sur les schémas réguliers}
\label{section-intersection-regular}

\noindent
Nous montrerons que $\CH_*(X) \otimes \mathbf{Q}$ possède un produit
d'intersection si $X$ est noethérien, régulier, de dimension finie et à
diagonale affine. La construction repose sur le résultat suivant
(qui est un corollaire de la proposition de la section précédente).

\begin{lemma}
\label{lemma-K-tensor-Q}
Soit $(S, \delta)$ comme dans la situation \ref{situation-setup}.
Soit $X$ un schéma régulier quasi-compact de type fini sur $S$, à
diagonale affine, et tel que $\delta_{X/S} : X \to \mathbf{Z}$ soit bornée.
Alors la composée
$$
K_0(\textit{Vect}(X)) \otimes \mathbf{Q}
\longrightarrow
A^*(X) \otimes \mathbf{Q}
\longrightarrow
\CH_*(X) \otimes \mathbf{Q}
$$
de l'application $ch$ de la remarque \ref{remark-chern-character-K} et
de l'application $c \mapsto c \cap [X]$ est un isomorphisme.
\end{lemma}

\begin{proof}
Nous avons $K'_0(X) = K_0(X) = K_0(\textit{Vect}(X))$ d'après
Catégories dérivées de schémas, lemmes \ref{perfect-lemma-Kprime-K},
\ref{perfect-lemma-regular-resolution-property} et
\ref{perfect-lemma-K-is-old-K}.
D'après la remarque \ref{remark-chern-classes-agree},
la composée donnée coïncide avec l'application de la
proposition \ref{proposition-K-tensor-Q} pour $X = Y$.
Le résultat découle donc de la proposition.
\end{proof}

\noindent
Soient $X, S, \delta$ comme dans le lemme \ref{lemma-K-tensor-Q}.
Pour simplifier, travaillons avec des cycles d'une codimension donnée, voir la
section \ref{section-cycles-codimension}.
Soit $[X]$ le cycle fondamental de $X$, voir la
remarque \ref{remark-fundamental-class}.
Choisissons $\alpha \in CH^i(X)$ et
$\beta \in CH^j(X)$. D'après le lemme, nous pouvons trouver un unique
$\alpha' \in K_0(\textit{Vect}(X)) \otimes \mathbf{Q}$ tel que
$ch(\alpha') \cap [X]  = \alpha$.
Bien sûr, cela signifie que $ch_{i'}(\alpha') \cap [X] = 0$
si $i' \not = i$ et que $ch_i(\alpha') \cap [X] = \alpha$.
D'après le lemme \ref{lemma-adams-and-chern}, nous voyons que
$\alpha'' = 2^{-i}\psi^2(\alpha')$ est une autre solution.
Par unicité, nous obtenons $\alpha'' = \alpha'$, et nous concluons
que $ch_{i'}(\alpha) = 0$ dans $A^{i'}(X) \otimes \mathbf{Q}$
pour $i' \not = i$. Nous pouvons alors définir
$$
\alpha \cdot \beta = ch(\alpha') \cap \beta =
ch_i(\alpha') \cap \beta
$$
dans $\CH^{i + j}(X) \otimes \mathbf{Q}$ grâce à la propriété de $\alpha'$
observée ci-dessus. C'est un accouplement symétrique : en effet, si nous choisissons
$\beta' \in K_0(\textit{Vect}(X)) \otimes \mathbf{Q}$ relevant
$\beta$, alors nous obtenons
$$
\alpha \cdot \beta = ch(\alpha') \cap \beta =
ch(\alpha') \cap ch(\beta') \cap [X]
$$
et nous savons que les classes de Chern commutent.
Le produit d'intersection est associatif pour la même raison :
$$
(\alpha \cdot \beta) \cdot \gamma =
ch(\alpha') \cap ch(\beta') \cap ch(\gamma') \cap [X]
$$
car nous savons que la composition des classes bivariantes est associative.
Une meilleure formulation est peut-être la suivante : il existe
un unique produit d'intersection commutatif et associatif sur
$\CH^*(X) \otimes \mathbf{Q}$, compatible à la graduation, tel que
l'isomorphisme
$K_0(\textit{Vect}(X)) \otimes \mathbf{Q} \to \CH^*(X) \otimes \mathbf{Q}$
soit un isomorphisme d'anneaux.






\section{Applications de Gysin pour les morphismes d'intersection complète locale}
\label{section-koszul}

\noindent
Avant de lire cette section, nous suggérons au lecteur de se renseigner sur les
immersions régulières
(Diviseurs, section \ref{divisors-section-regular-immersions}) et les
morphismes d'intersection complète locale
(Compléments sur les morphismes, section \ref{more-morphisms-section-lci}).

\medskip\noindent
Soit $(S, \delta)$ comme dans la situation \ref{situation-setup}.
Soit $i : X \to Y$ une
immersion régulière\footnote{Voir
Diviseurs, définition \ref{divisors-definition-regular-immersion}.
Remarquons que les immersions régulières sont la même chose que les
immersions Koszul-régulières ou quasi-régulières
pour les schémas localement noethériens, voir
Diviseurs, lemme \ref{divisors-lemma-regular-immersion-noetherian}.
Nous l'utiliserons sans autre mention dans cette section.}
de schémas localement de type fini sur $S$.
En particulier, le faisceau conormal $\mathcal{C}_{X/Y}$ est localement libre de rang fini
(voir Diviseurs, lemme \ref{divisors-lemma-quasi-regular-immersion}). Ainsi, le
faisceau normal
$$
\mathcal{N}_{X/Y} = \SheafHom_{\mathcal{O}_X}(\mathcal{C}_{X/Y}, \mathcal{O}_X)
$$
est lui aussi localement libre de rang fini, et nous avons une surjection
$\mathcal{N}_{X/Y}^\vee \to \mathcal{C}_{X/Y}$ (car un isomorphisme
est aussi une surjection).
La construction de la section \ref{section-gysin-higher-codimension}
nous donne une classe bivariante canonique
$$
i^! = c(X \to Y, \mathcal{N}_{X/Y}) \in A^*(X \to Y)^\wedge
$$
Nous avons besoin de deux lemmes sur cette notion.

\begin{lemma}
\label{lemma-composition-regular-immersion}
Soit $(S, \delta)$ comme dans la situation \ref{situation-setup}.
Soient $i : X \to Y$ et $j : Y \to Z$ des immersions régulières
de schémas localement de type fini sur $S$. Alors
$j \circ i$ est une immersion régulière et
$(j \circ i)^! = i^! \circ j^!$.
\end{lemma}

\begin{proof}
La première assertion résulte de
Diviseurs, lemme \ref{divisors-lemma-composition-regular-immersion}.
D'après Diviseurs, lemme \ref{divisors-lemma-transitivity-conormal-quasi-regular},
il existe une suite exacte courte
$$
0 \to
i^*(\mathcal{C}_{Y/Z}) \to
\mathcal{C}_{X/Z} \to
\mathcal{C}_{X/Y} \to 0
$$
Le résultat découle donc du lemme plus général \ref{lemma-gysin-composition}.
\end{proof}

\begin{lemma}
\label{lemma-section-smooth}
Soit $(S, \delta)$ comme dans la situation \ref{situation-setup}.
Soit $p : P \to X$ un morphisme lisse de schémas localement de type fini
sur $S$, et soit $s : X \to P$ une section. Alors $s$ est une
immersion régulière et $1 = s^! \circ p^*$ dans $A^*(X)^\wedge$,
où $p^* \in A^*(P \to X)^\wedge$ est la classe bivariante
du lemme \ref{lemma-flat-pullback-bivariant}.
\end{lemma}

\begin{proof}
La première assertion est Diviseurs, lemme
\ref{divisors-lemma-section-smooth-regular-immersion}.
Il suffit de montrer que $s^! \cap p^*[Z] = [Z]$ dans
$\CH_*(X)$ pour tout sous-schéma fermé intègre $Z \subset X$,
car les hypothèses sont préservées par changement de base par $X' \to X$
localement de type fini. Après avoir remplacé $P$ par un voisinage ouvert
de $s(Z)$, nous pouvons supposer que $P \to X$ est lisse de dimension relative constante $r$.
Posons $\dim_\delta(Z) = n$. Alors toute composante irréductible de $p^{-1}(Z)$
est de dimension $r + n$, et $p^*[Z]$ est donné par $[p^{-1}(Z)]_{n + r}$.
Remarquons que $s(X) \cap p^{-1}(Z) = s(Z)$ schématiquement. Ainsi, d'après la
même référence que ci-dessus, $s(X) \cap p^{-1}(Z)$ est un sous-schéma fermé
régulièrement immergé dans $\overline{p}^{-1}(Z)$, de codimension $r$.
Nous concluons par le lemme \ref{lemma-gysin-fundamental}.
\end{proof}

\noindent
Soit $(S, \delta)$ comme dans la situation \ref{situation-setup}. Considérons un
diagramme commutatif
$$
\xymatrix{
X \ar[rd]_f \ar[rr]_i & & P \ar[ld]^g \\
& Y
}
$$
de schémas localement de type fini sur $S$, tel que $g$ soit lisse
et $i$ une immersion régulière. En combinant la classe bivariante
$i^!$ examinée ci-dessus à la classe bivariante $g^* \in A^*(P \to Y)^\wedge$
du lemme \ref{lemma-flat-pullback-bivariant}, nous obtenons
$$
f^! = i^! \circ g^* \in A^*(X \to Y)
$$
Remarquons que le morphisme $f$ est un morphisme d'intersection complète locale, voir
Compléments sur les morphismes, définition \ref{more-morphisms-definition-lci}.
Réciproquement, si $f : X \to Y$ est un morphisme d'intersection complète locale
de schémas localement noethériens et $f = g \circ i$ avec $g$ lisse, alors
$i$ est une immersion régulière. Nous affirmons que notre construction de $f^!$
ne dépend que du morphisme $f$ et non du choix de la factorisation
$f = g \circ i$.

\begin{lemma}
\label{lemma-lci-gysin-well-defined}
Soit $(S, \delta)$ comme dans la situation \ref{situation-setup}.
Soit $f : X \to Y$ un morphisme d'intersection complète locale
de schémas localement de type fini sur $S$.
La classe bivariante $f^!$ est indépendante du choix de
la factorisation $f = g \circ i$ avec $g$ lisse (pourvu
qu'une telle factorisation existe).
\end{lemma}

\begin{proof}
Étant donnée une seconde factorisation $f = g' \circ i'$, nous pouvons
considérer le morphisme lisse $g'' : P \times_Y P' \to Y$, une
immersion $i'' : X \to P \times_Y P'$ et la factorisation
$f = g'' \circ i''$. Nous pouvons donc supposer que nous avons un diagramme
$$
\xymatrix{
& P' \ar[d]^p \ar[rd]^{g'} \\
X \ar[r]^i \ar[ru]^{i'} & P \ar[r]^g & Y
}
$$
où $p$ est un morphisme lisse. Alors $(g')^* = p^* \circ g^*$
(lemme \ref{lemma-compose-flat-pullback}), et il suffit donc
de montrer que $i^! = (i')^! \circ p^*$
dans $A^*(X \to P)$. Considérons le diagramme commutatif
$$
\xymatrix{
& X \times_P P' \ar[d]^{\overline{p}} \ar[r]_j & P' \ar[d]^p \\
X \ar[ru]^s \ar[r]^1 & X \ar[r]^i & P
}
$$
où $s =(1, i')$. Alors $s$ et $j$ sont des immersions régulières
(d'après Diviseurs, lemme \ref{divisors-lemma-section-smooth-regular-immersion},
et Diviseurs, lemme \ref{divisors-lemma-flat-base-change-regular-immersion}),
et $i' = j \circ s$. D'après le lemme \ref{lemma-composition-regular-immersion},
nous avons $(i')^! = s^! \circ j^!$.
Comme le carré est cartésien, la classe bivariante $j^!$
est la restriction (remarque \ref{remark-restriction-bivariant})
de $i^!$ à $P'$, voir le lemme \ref{lemma-construction-gysin}.
Comme les classes bivariantes commutent aux images inverses plates,
nous trouvons $j^! \circ p^* = \overline{p}^* \circ i^!$.
Il suffit donc de montrer que $s^! \circ \overline{p}^* = \text{id}$,
ce qui est fait dans le lemme \ref{lemma-section-smooth}.
\end{proof}

\begin{definition}
\label{definition-lci-gysin}
Soit $(S, \delta)$ comme dans la situation \ref{situation-setup}.
Soit $f : X \to Y$ un morphisme d'intersection complète locale
de schémas localement de type fini sur $S$. Nous disons que
{\it l'application de Gysin pour $f$ existe} si nous pouvons écrire
$f = g \circ i$ avec $g$ lisse et $i$ une immersion.
Dans ce cas, nous définissons
{\it l'application de Gysin} $f^! = i^! \circ g^* \in A^*(X \to Y)$ comme ci-dessus.
\end{definition}

\noindent
Il résulte de la définition que, pour une immersion régulière,
ceci coïncide avec la construction antérieure et que, pour un morphisme
lisse, ceci coïncide avec l'image inverse plate. En fait, cette coïncidence
vaut pour tous les morphismes syntomiques.

\begin{lemma}
\label{lemma-lci-gysin-flat}
Soit $(S, \delta)$ comme dans la situation \ref{situation-setup}.
Soit $f : X \to Y$ un morphisme d'intersection complète locale
de schémas localement de type fini sur $S$. Si l'application de Gysin
existe pour $f$ et si $f$ est plat, alors $f^!$ est égal à la
classe bivariante du lemme \ref{lemma-flat-pullback-bivariant}.
\end{lemma}

\begin{proof}
Choisissons une factorisation $f = g \circ i$, avec $i : X \to P$
une immersion et $g : P \to Y$ lisse. Remarquons que, pour
tout morphisme $Y' \to Y$ localement de type fini,
les changements de base $f'$, $g'$, $i'$ satisfont aux mêmes
hypothèses (voir Morphismes, lemmes \ref{morphisms-lemma-base-change-smooth}
et \ref{morphisms-lemma-base-change-syntomic}, et
Compléments sur les morphismes, lemme \ref{more-morphisms-lemma-flat-lci}).
Nous nous ramenons donc à prouver que $f^*[Y] = i^!(g^*[Y])$ lorsque $Y$
est intègre, voir le lemme \ref{lemma-bivariant-zero}. Posons $n = \dim_\delta(Y)$.
Après avoir décomposé $X$ et $P$ en composantes connexes, nous
pouvons supposer que $f$ est plat de dimension relative $r$ et
que $g$ est lisse de dimension relative $t$.
Alors $f^*[Y] = [X]_{n + s}$ et $g^*[Y] = [P]_{n + t}$.
D'autre part, $i$ est une immersion régulière de codimension $t - s$.
Ainsi, $i^![P]_{n + t} = [X]_{n + s}$ (lemme \ref{lemma-gysin-fundamental}),
et la démonstration est achevée.
\end{proof}

\begin{lemma}
\label{lemma-lci-gysin-composition}
Soit $(S, \delta)$ comme dans la situation \ref{situation-setup}.
Soient $f : X \to Y$ et $g : Y \to Z$ des morphismes d'intersection complète locale
de schémas localement de type fini sur $S$. Supposons que l'application de Gysin
existe pour $g \circ f$ et $g$. Alors l'application de Gysin existe pour $f$
et $(g \circ f)^! = f^! \circ g^!$.
\end{lemma}

\begin{proof}
Remarquons que $g \circ f$ est un morphisme d'intersection complète locale
d'après Compléments sur les morphismes, lemme \ref{more-morphisms-lemma-composition-lci},
et que l'énoncé du lemme a donc un sens.
Si $X \to P$ est une immersion de $X$ dans un schéma $P$ lisse sur $Z$,
alors $X \to P \times_Z Y$ est une immersion de $X$ dans un schéma lisse
sur $Y$. Cela prouve la première assertion du lemme.
Soit $Y \to P'$ une immersion de $Y$ dans un schéma $P'$ lisse sur $Z$.
Considérons le diagramme commutatif
$$
\xymatrix{
X \ar[r] \ar[d] &
P \times_Z Y \ar[r]_a \ar[ld]^p &
P \times_Z P' \ar[ld]^q \\
Y \ar[r]_b \ar[d] &
P' \ar[ld] \\
Z
}
$$
Ici, les flèches horizontales sont des immersions régulières, les flèches sud-ouest
sont lisses, et le carré est cartésien. Par conséquent,
$a^! \circ q^* = p^* \circ b^!$, puisque les classes bivariantes commutent
aux images inverses plates. En combinant ce fait avec les
lemmes \ref{lemma-composition-regular-immersion} et
\ref{lemma-compose-flat-pullback},
le lecteur constate que l'énoncé du lemme est vrai.
Nous omettons un petit détail.
\end{proof}

\begin{lemma}
\label{lemma-lci-gysin-commutes}
Soit $(S, \delta)$ comme dans la situation \ref{situation-setup}.
Considérons un diagramme commutatif
$$
\xymatrix{
X'' \ar[d] \ar[r] &
X' \ar[d] \ar[r] &
X \ar[d]^f \\
Y'' \ar[r] &
Y' \ar[r] &
Y
}
$$
de schémas localement de type fini sur $S$, dont les deux carrés sont cartésiens.
Supposons que $f : X \to Y$ soit un morphisme d'intersection complète locale
tel que l'application de Gysin existe pour $f$. Soit $c \in A^*(Y'' \to Y')$. Notons
$res(f^!) \in A^*(X' \to Y')$ la restriction de $f^!$ à $Y'$
(remarque \ref{remark-restriction-bivariant}). Alors $c$ et $res(f^!)$ commutent
(remarque \ref{remark-bivariant-commute}).
\end{lemma}

\begin{proof}
Choisissons une factorisation $f = g \circ i$ avec $g$ lisse et $i$ une immersion.
Comme $f^! = i^! \circ g^!$, il suffit de prouver le lemme pour $g^!$
(qui est donné par l'image inverse plate) et pour $i^!$. Le résultat pour l'image inverse plate
fait partie de la définition d'une classe bivariante. Le cas de $i^!$ découle
immédiatement du lemme \ref{lemma-gysin-commutes}.
\end{proof}

\begin{lemma}
\label{lemma-lci-gysin-easy}
Soit $(S, \delta)$ comme dans la situation \ref{situation-setup}.
Considérons un diagramme cartésien
$$
\xymatrix{
X' \ar[d]_{f'} \ar[r] &
X \ar[d]^f \\
Y' \ar[r] &
Y
}
$$
de schémas localement de type fini sur $S$. Supposons que
\begin{enumerate}
\item $f$ soit un morphisme d'intersection complète locale et que
l'application de Gysin existe pour $f$,
\item $X$, $X'$, $Y$, $Y'$ vérifient les conditions équivalentes du
lemme \ref{lemma-locally-equidimensional},
\item pour $x' \in X'$, d'images $x$, $y'$ et $y$
dans $X$, $Y'$ et $Y$, on ait $n_{x'} - n_{y'} = n_x - n_y$,
où $n_{x'}$, $n_x$, $n_{y'}$ et $n_y$ sont comme dans le lemme, et
\item pour tout point générique $\xi \in X'$, l'anneau local
$\mathcal{O}_{Y', f'(\xi)}$ soit de Cohen-Macaulay.
\end{enumerate}
Alors $f^![Y'] = [X']$, où $[Y']$ et $[X']$ sont comme dans la
remarque \ref{remark-fundamental-class}.
\end{lemma}

\begin{proof}
Rappelons que $n_{x'}$ est la valeur commune de $\delta(\xi)$
lorsque $\xi$ parcourt les points génériques des composantes irréductibles
passant par $x'$. De plus, les fonctions
$x' \mapsto n_{x'}$, $x \mapsto n_x$, $y' \mapsto n_{y'}$ et
$y \mapsto n_y$ sont localement constantes. Soient $X'_n$, $X_n$, $Y'_n$
et $Y_n$ les sous-schémas ouverts et fermés de $X'$, $X$, $Y'$ et
$Y$ où la fonction vaut $n$. Rappelons que
$[X'] = \sum [X'_n]_n$ et $[Y'] = \sum [Y'_n]_n$.
Cela étant dit, il est clair que, pour démontrer le lemme, nous
pouvons remplacer $X'$ par l'une de ses composantes connexes
et $X$, $Y'$, $Y$ par la composante connexe dans laquelle
elle s'envoie. Nous savons alors que $X'$, $X$, $Y'$ et
$Y$ sont $\delta$-équidimensionnels, au sens où
chaque composante irréductible a la même $\delta$-dimension.
Notons $n'$, $n$, $m'$ et $m$ cette valeur commune
pour $X'$, $X$, $Y'$ et $Y$. La dernière hypothèse
signifie que $n' - m' = n - m$.

\medskip\noindent
Choisissons une factorisation $f = g \circ i$, où $i : X \to P$
est une immersion et $g : P \to Y$ est lisse. Comme $X$ est connexe,
nous voyons que la dimension relative de $P \to Y$ aux points de $i(X)$
est constante. Ainsi, après avoir remplacé $P$ par un voisinage ouvert
de $i(X)$, nous pouvons supposer que $P \to Y$ est de dimension relative constante
et que $i : X \to P$ est une immersion fermée.
Notons $g' : Y' \times_Y P \to Y'$ le changement de base de $g$ et
$i' : X' \to Y' \times_Y P$ le changement de base de $i$.
Il est clair que $g^*[Y] = [P]$ et $(g')^*[Y'] = [Y' \times_Y P]$.
Enfin, si $\xi' \in X'$ est un point générique, alors
$\mathcal{O}_{Y' \times_Y P, i'(\xi)}$ est de Cohen-Macaulay.
En effet, l'homomorphisme d'anneaux locaux
$\mathcal{O}_{Y', f'(\xi)} \to \mathcal{O}_{Y' \times_Y P, i'(\xi)}$
est plat à fibre régulière
(voir Algèbre, section \ref{algebra-section-smooth-overview}),
un anneau local régulier est de Cohen-Macaulay
(Algèbre, lemme \ref{algebra-lemma-regular-ring-CM}),
$\mathcal{O}_{Y', f'(\xi)}$ est de Cohen-Macaulay par l'hypothèse
(4), et le résultat voulu découle de
Algèbre, lemme \ref{algebra-lemma-CM-goes-up}.
Nous nous ramenons donc au cas traité dans le paragraphe suivant.

\medskip\noindent
Supposons que $f$ soit une immersion fermée régulière et que $X'$, $X$, $Y'$ et
$Y$ soient $\delta$-équidimensionnels, de $\delta$-dimensions
$n'$, $n$, $m'$ et $m$, avec $m' - n' = m - n$.
Dans ce cas, le résultat découle immédiatement du
lemme \ref{lemma-gysin-easy}.
\end{proof}

\begin{remark}
\label{remark-gysin-chern-classes}
Soit $(S, \delta)$ comme dans la situation \ref{situation-setup}.
Soit $f : X \to Y$ un morphisme d'intersection complète locale
de schémas localement de type fini sur $S$. Supposons que l'application
de Gysin existe pour $f$. Alors
$f^! \circ c_i(\mathcal{E}) = c_i(f^*\mathcal{E}) \circ f^!$,
et de même pour le caractère de Chern, voir le
lemme \ref{lemma-lci-gysin-commutes}.
Si $X$ et $Y$ vérifient les conditions équivalentes du
lemme \ref{lemma-locally-equidimensional} et si $Y$ est de Cohen-Macaulay
(par exemple), alors $f^![Y] = [X]$ par le lemme \ref{lemma-lci-gysin-easy}.
Dans ce cas, nous obtenons aussi
$f^!(c_i(\mathcal{E}) \cap [Y]) = c_i(f^*\mathcal{E}) \cap [X]$,
et de même pour le caractère de Chern.
\end{remark}

\begin{lemma}
\label{lemma-compare-gysin-base-change}
Soit $(S, \delta)$ comme dans la situation \ref{situation-setup}.
Considérons un carré cartésien
$$
\xymatrix{
X' \ar[d]_{f'} \ar[r]_{g'} &
X \ar[d]^f \\
Y' \ar[r]^g &
Y
}
$$
de schémas localement de type fini sur $S$. Supposons que
\begin{enumerate}
\item $f$ et $f'$ soient tous deux des morphismes d'intersection complète locale, et
\item l'application de Gysin existe pour $f$.
\end{enumerate}
Alors $\mathcal{C} = \Ker(H^{-1}((g')^*\NL_{X/Y}) \to H^{-1}(\NL_{X'/Y'}))$
est un $\mathcal{O}_{X'}$-module localement libre de rang fini, l'application de Gysin
existe pour $f'$, et nous avons
$$
res(f^!) = c_{top}(\mathcal{C}^\vee) \circ (f')^!
$$
dans $A^*(X' \to Y')$.
\end{lemma}

\begin{proof}
Le fait que $\mathcal{C}$ soit localement libre de rang fini découle immédiatement de
Compléments d'algèbre, lemme \ref{more-algebra-lemma-base-change-lci-bis}.
Choisissons une factorisation $f = h \circ i$ avec $h : P \to Y$ lisse et $i$
une immersion. Nous pouvons alors factoriser $f' = h' \circ i'$, où $h' : P' \to Y'$
et $i' : X' \to P'$ sont les changements de base correspondants. On a le diagramme
$$
\xymatrix{
X' \ar[r] \ar[d] &
P' \ar[r] \ar[d] &
Y' \ar[d] \\
X \ar[r]^i &
P \ar[r]^h &
Y
}
$$
En particulier, nous voyons que l'application de Gysin existe pour $f'$. Par
Compléments sur les morphismes, lemme \ref{more-morphisms-lemma-get-NL},
nous avons
$$
\NL_{X/Y} = \left( \mathcal{C}_{X/P} \to i^*\Omega_{P/Y} \right)
$$
où $\mathcal{C}_{X/P}$ est le faisceau conormal de l'immersion $i$.
De même pour la version avec primes. Nous avons
$(g')^*i^*\Omega_{P/Y} = (i')^*\Omega_{P'/Y'}$, car
$\Omega_{P/Y}$ devient $\Omega_{P'/Y'}$ par image inverse, d'après
Morphismes, lemme \ref{morphisms-lemma-base-change-differentials}.
Rappelons aussi que $(g')^*\mathcal{C}_{X/P} \to \mathcal{C}_{X'/P'}$
est surjectif, voir
Morphismes, lemme \ref{morphisms-lemma-conormal-functorial-flat}.
Nous en déduisons que le faisceau $\mathcal{C}$ est canoniquement
isomorphe au noyau de l'application
$(g')^*\mathcal{C}_{X/P} \to \mathcal{C}_{X'/P'}$
de modules localement libres de rang fini. Rappelons que $i^!$ est défini
à l'aide de $\mathcal{N} = \mathcal{C}_{Z/X}^\vee$, et de même
pour $(i')^!$. Ainsi, nous avons
$$
res(i^!) = c_{top}(\mathcal{C}^\vee) \circ (i')^!
$$
dans $A^*(X' \to P')$, par application du lemme \ref{lemma-gysin-excess}.
Enfin, puisque $f^! = i^! \circ h^*$,
$(f')^! = (i')^! \circ (h')^*$ et $(h')^* = res(h^*)$, nous concluons.
\end{proof}

\begin{lemma}[Formule de l'éclatement]
\label{lemma-blow-up-formula}
Soit $(S, \delta)$ comme dans la situation \ref{situation-setup}.
Soit $i : Z \to X$ une immersion fermée régulière de schémas
localement de type fini sur $S$. Soit $b : X' \to X$
l'éclatement de centre $Z$. On a le diagramme
$$
\xymatrix{
E \ar[r]_j \ar[d]_\pi & X' \ar[d]^b \\
Z \ar[r]^i & X
}
$$
Supposons que l'application de Gysin existe pour $b$. Alors
$$
res(b^!) = c_{top}(\mathcal{F}^\vee) \circ \pi^*
$$
dans $A^*(E \to Z)$, où $\mathcal{F}$ est le noyau de l'application canonique
$\pi^*\mathcal{C}_{Z/X} \to \mathcal{C}_{E/X'}$.
\end{lemma}

\begin{proof}
Observons que le morphisme $b$ est un morphisme d'intersection complète locale
par Compléments d'algèbre, lemme \ref{more-algebra-lemma-blowup-regular-sequence},
et que l'énoncé a donc un sens. Comme $Z \to X$ est une
immersion régulière (et donc a fortiori quasi-régulière), nous voyons que $\mathcal{C}_{Z/X}$
est localement libre de rang fini et que l'application
$\text{Sym}^*(\mathcal{C}_{Z/X}) \to \mathcal{C}_{Z/X, *}$
est un isomorphisme, voir
Diviseurs, lemme \ref{divisors-lemma-quasi-regular-immersion}.
Comme $E = \text{Proj}(\mathcal{C}_{Z/X, *})$, nous concluons
que $E = \mathbf{P}(\mathcal{C}_{Z/X})$
est un fibré en espaces projectifs sur $Z$.
Ainsi $E \to Z$ est lisse et, en particulier, un morphisme d'intersection complète
locale. Le lemme \ref{lemma-compare-gysin-base-change}
s'applique donc, et nous voyons que
$$
res(b^!) = c_{top}(\mathcal{C}^\vee) \circ \pi^!
$$
avec $\mathcal{C}$ comme dans l'énoncé de ce lemme.
Bien entendu, $\pi^* = \pi^!$ par le lemme \ref{lemma-lci-gysin-flat}.
Il reste à montrer que $\mathcal{F}$ est égal au
noyau $\mathcal{C}$ de l'application
$H^{-1}(j^*\NL_{X'/X}) \to H^{-1}(\NL_{E/Z})$.

\medskip\noindent
Comme $E \to Z$ est lisse, nous avons $H^{-1}(\NL_{E/Z}) = 0$, voir
Compléments sur les morphismes, lemme \ref{more-morphisms-lemma-NL-smooth}.
Il suffit donc de montrer que $\mathcal{F}$ s'identifie
à $H^{-1}(j^*\NL_{X'/X})$. Par Compléments sur les morphismes, Lemmes
\ref{more-morphisms-lemma-get-triangle-NL} et
\ref{more-morphisms-lemma-NL-immersion}, nous avons une suite exacte
$$
0 \to H^{-1}(j^*\NL_{X'/X}) \to H^{-1}(\NL_{E/X}) \to
\mathcal{C}_{E/X'} \to \ldots
$$
Les mêmes lemmes appliqués à $E \to Z \to X$ donnent un isomorphisme
$\pi^*\mathcal{C}_{Z/X} = H^{-1}(\pi^*\NL_{Z/X}) \to H^{-1}(\NL_{E/X})$.
Nous pouvons donc conclure.
\end{proof}

\begin{lemma}
\label{lemma-lci-gysin-product-regular}
Soit $(S, \delta)$ comme dans la situation \ref{situation-setup}.
Soit $f : X \to Y$ un morphisme entre des schémas localement de type
fini sur $S$, tels que $X$ et $Y$ soient quasi-compacts,
réguliers, à diagonale affine et de dimension finie.
Alors $f$ est un morphisme d'intersection complète locale.
Supposons de plus que l'application de Gysin existe pour $f$. Alors
$$
f^!(\alpha \cdot \beta) = f^!\alpha \cdot f^!\beta
$$
dans $\CH^*(X) \otimes \mathbf{Q}$, où le produit d'intersection
est celui de la section \ref{section-intersection-regular}.
\end{lemma}

\begin{proof}
Le premier énoncé découle de
Compléments sur les morphismes, Lemme
\ref{more-morphisms-lemma-morphism-regular-schemes-is-lci}.
Observons que $f^![Y] = [X]$, voir le lemme \ref{lemma-lci-gysin-easy}.
Écrivons $\alpha = ch(\alpha') \cap [Y]$ et $\beta = ch(\beta') \cap [Y]$,
où $\alpha', \beta' \in K_0(\textit{Vect}(X)) \otimes \mathbf{Q}$,
comme dans la section \ref{section-intersection-regular}.
En posant $c = ch(\alpha')$ et $c' = ch(\beta')$, nous obtenons
$\alpha \cdot \beta = c \cap c' \cap [Y]$ par construction.
D'après le lemme \ref{lemma-lci-gysin-commutes}, nous savons que $f^!$
commute avec $c$ et $c'$. Par conséquent,
\begin{align*}
f^!(\alpha \cdot \beta)
& =
f^!(c \cap c' \cap [Y]) \\
& =
c \cap c' \cap f^![Y] \\
& =
c \cap c' \cap [X] \\
& =
(c \cap [X]) \cdot (c' \cap [X]) \\
& =
(c \cap f^![Y]) \cdot (c' \cap f^![Y]) \\
& =
f^!(\alpha) \cdot f^!(\beta)
\end{align*}
comme voulu.
\end{proof}

\begin{lemma}
\label{lemma-projection-formula-regular}
Soit $(S, \delta)$ comme dans la situation \ref{situation-setup}.
Soit $f : X \to Y$ un morphisme entre des schémas localement de type
fini sur $S$, tels que $X$ et $Y$ soient quasi-compacts,
réguliers, à diagonale affine et de dimension finie.
Alors $f$ est un morphisme d'intersection complète locale.
Supposons de plus que l'application de Gysin existe pour $f$
et que $f$ soit propre. Alors
$$
f_*(\alpha \cdot f^!\beta) = f_*\alpha \cdot \beta
$$
dans $\CH^*(Y) \otimes \mathbf{Q}$, où le produit d'intersection
est celui de la section \ref{section-intersection-regular}.
\end{lemma}

\begin{proof}
Le premier énoncé découle de
Compléments sur les morphismes, Lemme
\ref{more-morphisms-lemma-morphism-regular-schemes-is-lci}.
Observons que $f^![Y] = [X]$, voir le lemme \ref{lemma-lci-gysin-easy}.
Écrivons $\alpha = ch(\alpha') \cap [X]$ et $\beta = ch(\beta') \cap [Y]$,
avec $\alpha' \in K_0(\textit{Vect}(X)) \otimes \mathbf{Q}$ et
$\beta' \in K_0(\textit{Vect}(Y)) \otimes \mathbf{Q}$,
comme dans la section \ref{section-intersection-regular}.
Posons $c = ch(\alpha')$ et $c' = ch(\beta')$. Nous avons
\begin{align*}
f_*(\alpha \cdot f^!\beta)
& =
f_*(c \cap f^!(c' \cap [Y]_e)) \\
& =
f_*(c \cap c' \cap f^![Y]_e) \\
& =
f_*(c \cap c' \cap [X]_d) \\
& =
f_*(c' \cap c \cap [X]_d) \\
& =
c' \cap f_*(c \cap [X]_d) \\
& =
\beta \cdot f_*(\alpha)
\end{align*}
La première égalité résulte de la construction du produit d'intersection.
D'après le lemme \ref{lemma-lci-gysin-commutes}, nous savons que $f^!$
commute avec $c'$. Le fait que les classes de Chern soient dans le centre
de l'anneau bivariant justifie d'intervertir l'ordre dans lequel on intersecte
$[X]$ avec $c$ et $c'$. On peut commuter $c'$ avec $f_*$, car $c'$
est une classe bivariante. La dernière égalité résulte à nouveau de la construction
du produit d'intersection.
\end{proof}













\section{Applications de Gysin pour les diagonales}
\label{section-gysin-for-diagonal}

\noindent
Soit $(S, \delta)$ comme dans la situation \ref{situation-setup}. Soit $f : X \to Y$
un morphisme lisse de schémas localement de type fini sur $S$. Alors le
morphisme diagonal $\Delta : X \longrightarrow X \times_Y X$
est une immersion régulière, voir
Compléments sur les morphismes, lemme \ref{more-morphisms-lemma-smooth-diagonal-perfect}.
Nous avons donc l'application de Gysin
$$
\Delta^! \in A^*(X \to X \times_Y X)^\wedge
$$
construite à la section \ref{section-koszul}. Si $X \to Y$ est de
dimension relative constante $d$, alors $\Delta^! \in A^d(X \to X \times_Y X)$.

\begin{lemma}
\label{lemma-diagonal-identity}
Dans la situation ci-dessus, nous avons $\Delta^! \circ \text{pr}_i^! = 1$ dans $A^0(X)$.
\end{lemma}

\begin{proof}
Observons que les projections $\text{pr}_i : X \times_Y X \to X$ sont
lisses, de sorte que nous disposons aussi d'applications de Gysin pour ces projections.
Le lemme a donc un sens et est un cas particulier du
lemme \ref{lemma-lci-gysin-composition}.
\end{proof}

\begin{proposition}
\label{proposition-compute-bivariant}
\begin{reference}
\cite[Proposition 17.4.2]{F}
\end{reference}
Soit $(S, \delta)$ comme dans la situation \ref{situation-setup}.
Soient $f : X \to Y$ et $g : Y \to Z$ des morphismes de schémas localement
de type fini sur $S$. Si $g$ est lisse de dimension relative $d$, alors
$A^p(X \to Y) = A^{p - d}(X \to Z)$.
\end{proposition}

\begin{proof}
Nous utiliserons le fait que les morphismes lisses sont des morphismes d'intersection complète
locale dont les applications de Gysin existent (voir la section \ref{section-koszul}).
En particulier, nous avons $g^! \in A^{-d}(Y \to Z)$. Nous pouvons alors envoyer
$c \in A^p(X \to Y)$ sur $c \circ g^! \in A^{p - d}(X \to Z)$.

\medskip\noindent
Réciproquement, soit $c' \in A^{p - d}(X \to Z)$. Notons $res(c')$ la restriction
(remarque \ref{remark-restriction-bivariant}) de $c'$ par le morphisme $Y \to Z$.
Puisque le diagramme
$$
\xymatrix{
X \times_Z Y \ar[r]_{\text{pr}_2} \ar[d]_{\text{pr}_1} & Y \ar[d]^g \\
X \ar[r]^f & Z
}
$$
est cartésien, nous obtenons $res(c') \in A^{p - d}(X \times_Z Y \to Y)$.
Soit $\Delta : Y \to Y \times_Z Y$ la diagonale, et notons
$res(\Delta^!)$ la restriction de $\Delta^!$
à $X \times_Z Y$ par le morphisme $X \times_Z Y \to Y \times_Z Y$.
Puisque le diagramme
$$
\xymatrix{
X \ar[r] \ar[d] & X \times_Z Y \ar[d] \\
Y \ar[r]^-\Delta & Y \times_Z Y
}
$$
est cartésien, nous voyons que $res(\Delta^!) \in A^d(X \to X \times_Z Y)$.
En composant ces deux restrictions, nous obtenons
$$
res(\Delta^!) \circ res(c') \in A^p(X \to Y)
$$
Nous avons ainsi construit des applications $A^p(X \to Y) \to A^{p - d}(X \to Z)$
et $A^{p - d}(X \to Z) \to A^p(X \to Y)$. Pour achever la démonstration, nous
allons montrer que ces applications sont réciproques l'une de l'autre.

\medskip\noindent
Commençons avec $c \in A^p(X \to Y)$. Considérons le diagramme
$$
\xymatrix{
X \ar[d] \ar[r] & Y \ar[d] \\
X \times_Z Y \ar[r] \ar[d]^{\text{pr}_1} &
Y \times_Z Y \ar[r]_{p_2} \ar[d]^{p_1} &
Y \ar[d]^g \\
X \ar[r]^f &
Y \ar[r]^g &
Z
}
$$
dont les carrés sont cartésiens. Les deux carrés inférieurs de ce diagramme
montrent que $res(c \circ g^!) = res(c) \cap p_2^!$, où, dans cette formule,
$res(c)$ désigne la restriction de $c$ par $p_1$. En considérant le carré supérieur
du diagramme et en utilisant le lemme \ref{lemma-lci-gysin-commutes},
nous obtenons $c \circ \Delta^! = res(\Delta^!) \circ res(c)$.
Nous calculons
\begin{align*}
res(\Delta^!) \circ res(c \circ g^!)
& =
res(\Delta^!) \circ res(c) \circ p_2^! \\
& =
c \circ \Delta^! \circ p_2^! \\
& =
c
\end{align*}
La dernière égalité découle du lemme \ref{lemma-diagonal-identity}.

\medskip\noindent
Réciproquement, partons de $c' \in A^{p - d}(X \to Z)$. En considérant
le rectangle inférieur du diagramme ci-dessus, nous obtenons
$res(c') \circ g^! = \text{pr}_1^! \circ c'$.
Nous calculons
\begin{align*}
res(\Delta^!) \circ res(c') \circ g^!
& =
res(\Delta^!) \circ \text{pr}_1^! \circ c' \\
& =
c'
\end{align*}
La dernière égalité résulte du fait que les deux carrés de gauche du
diagramme montrent que
$\text{id} = res(\Delta^! \circ p_1^!) = res(\Delta^!) \circ \text{pr}_1^!$.
Ceci achève la démonstration.
\end{proof}









\section{Produit extérieur}
\label{section-exterior-product}

\noindent
Soit $k$ un corps. Dans cette section, nous travaillons sur $S = \Spec(k)$ avec
$\delta : S \to \mathbf{Z}$ définie en envoyant l'unique point sur $0$, voir
l'Exemple \ref{example-field}.

\medskip\noindent
Considérons un carré cartésien
$$
\xymatrix{
X \times_k Y \ar[r] \ar[d] & Y \ar[d] \\
X \ar[r] & \Spec(k) = S
}
$$
de schémas localement de type fini sur $k$. Il existe alors une application canonique
$$
\times :
\CH_n(X) \otimes_{\mathbf{Z}} \CH_m(Y)
\longrightarrow
\CH_{n + m}(X \times_k Y)
$$
qui est uniquement déterminée par la règle suivante :
étant donnés des sous-schémas fermés intègres $X' \subset X$
et $Y' \subset Y$ de dimensions $n$ et $m$, nous avons
$$
[X'] \times [Y'] = [X' \times_k Y']_{n + m}
$$
dans $\CH_{n + m}(X \times_k Y)$.

\begin{lemma}
\label{lemma-exterior-product-well-defined}
L'application
$\times : \CH_n(X) \otimes_{\mathbf{Z}} \CH_m(Y) \to \CH_{n + m}(X \times_k Y)$
est bien définie.
\end{lemma}

\begin{proof}
Une première remarque est que, si $\alpha = \sum n_i[X_i]$
et $\beta = \sum m_j[Y_j]$, avec $X_i \subset X$ et $Y_j \subset Y$
des familles localement finies de sous-schémas fermés intègres
de dimensions $n$ et $m$, alors
les $X_i \times_k Y_j$ forment une famille localement finie
de sous-schémas fermés de $X \times_k Y$, de
dimensions $n + m$, et nous pouvons effectivement considérer
$$
\alpha \times \beta = \sum n_i m_j [X_i \times_k Y_j]_{n + m}
$$
comme un $(n + m)$-cycle sur $X \times_k Y$. Nous obtenons ainsi une
application additive
$\times : Z_n(X) \otimes_{\mathbf{Z}} Z_m(Y) \to Z_{n + m}(X \times_k Y)$.
Il s'agit de montrer que
ce procédé est compatible avec l'équivalence rationnelle.

\medskip\noindent
Soit $i : X' \to X$ le morphisme d'inclusion d'un
sous-schéma fermé intègre de dimension $n$.
Alors l'image inverse plate par le morphisme $p' : X' \to \Spec(k)$ est un élément
$(p')^* \in A^{-n}(X' \to \Spec(k))$ d'après le
lemme \ref{lemma-flat-pullback-bivariant},
et donc $c' = i_* \circ (p')^* \in A^{-n}(X \to \Spec(k))$ d'après le
lemme \ref{lemma-push-proper-bivariant}.
Cela produit des applications
$$
c' \cap - : \CH_m(Y) \longrightarrow \CH_{m + n}(X \times_k Y)
$$
dont le lecteur vérifie aisément qu'elles envoient $[Y']$ sur $[X' \times_k Y']_{n + m}$
pour tout sous-schéma fermé intègre $Y' \subset Y$ de dimension
$m$. Par conséquent, la construction
$([X'], [Y']) \mapsto [X' \times_k Y']_{n + m}$
se factorise par l'équivalence rationnelle en la seconde variable, c'est-à-dire
donne une application bien définie
$Z_n(X) \otimes_{\mathbf{Z}} \CH_m(Y) \to \CH_{n + m}(X \times_k Y)$.
Par symétrie, il en est de même pour l'autre variable, et nous concluons.
\end{proof}

\begin{lemma}
\label{lemma-chow-cohomology-towards-point}
Soit $k$ un corps. Soit $X$ un schéma localement de type fini sur $k$.
Nous avons alors une identification canonique
$$
A^p(X \to \Spec(k)) = \CH_{-p}(X)
$$
pour tout $p \in \mathbf{Z}$.
\end{lemma}

\begin{proof}
Considérons l'élément $[\Spec(k)] \in \CH_0(\Spec(k))$. Nous obtenons une application
$A^p(X \to \Spec(k)) \to \CH_{-p}(X)$ qui envoie $c$ sur $c \cap [\Spec(k)]$.

\medskip\noindent
Réciproquement, supposons que nous ayons $\alpha \in \CH_{-p}(X)$.
Nous pouvons alors définir $c_\alpha \in A^p(X \to \Spec(k))$ comme
suit : étant donnés $X' \to \Spec(k)$ et $\alpha' \in \CH_n(X')$,
nous posons
$$
c_\alpha \cap \alpha' = \alpha \times \alpha'
$$
dans $\CH_{n - p}(X \times_k X')$. Pour montrer qu'il s'agit d'une classe
bivariante, écrivons $\alpha = \sum n_i[X_i]$ comme dans la
définition \ref{definition-cycles}. Considérons la composée
$$
\coprod X_i \xrightarrow{g} X \to \Spec(k)
$$
et notons $f : \coprod X_i \to \Spec(k)$ cette composée.
Alors $g$ est propre et $f$ est plat de dimension relative $-p$.
L'image inverse par $f$ est une classe bivariante
$f^* \in A^p(\coprod X_i \to \Spec(k))$ d'après le
lemme \ref{lemma-flat-pullback-bivariant}.
Notons $\nu \in A^0(\coprod X_i)$ la classe bivariante
qui multiplie un cycle par $n_i$ sur la $i$-ième composante.
Ainsi, $\nu \circ f^* \in A^p(\coprod X_i \to X)$.
Enfin, nous avons une classe bivariante
$$
g_* \circ \nu \circ f^*
$$
par le lemme \ref{lemma-push-proper-bivariant}. Le lecteur vérifie aisément
que $c_\alpha$ est égale à cette classe et qu'elle est donc
elle-même une classe bivariante.

\medskip\noindent
Pour achever la démonstration, nous devons montrer que les deux constructions
sont réciproques l'une de l'autre. Comme $c_\alpha \cap [\Spec(k)] = \alpha$,
c'est clair dans un sens. Pour l'autre sens, soit
$c \in A^p(X \to \Spec(k))$ et posons $\alpha = c \cap [\Spec(k)]$.
Il suffit de démontrer que
$$
c \cap [X'] = c_\alpha \cap [X']
$$
lorsque $X'$ est un schéma intègre localement de type fini sur $\Spec(k)$,
voir le lemme \ref{lemma-bivariant-zero}. Or $p' : X' \to \Spec(k)$ est alors
plat de dimension relative $\dim(X')$, et donc
$[X'] = (p')^*[\Spec(k)]$. Ainsi, le fait que les classes bivariantes
$c$ et $c_\alpha$ coïncident sur $[\Spec(k)]$ implique qu'elles
coïncident après intersection avec $[X']$, ce qui achève la démonstration.
\end{proof}

\begin{lemma}
\label{lemma-chow-cohomology-towards-point-commutes}
Soit $k$ un corps. Soit $X$ un schéma localement de type fini sur $k$.
Soit $c \in A^p(X \to \Spec(k))$. Soit $Y \to Z$ un morphisme de schémas
localement de type fini sur $k$. Soit $c' \in A^q(Y \to Z)$. Alors
$c \circ c' = c' \circ c$ dans $A^{p + q}(X \times_k Y \to Z)$.
\end{lemma}

\begin{proof}
Dans la démonstration du lemme \ref{lemma-chow-cohomology-towards-point},
nous avons vu que $c$ est donnée par une combinaison d'images
directes propres, de multiplications par des entiers sur les composantes
connexes et d'images inverses plates. Comme $c'$ commute avec chacune de
ces opérations par définition des classes bivariantes, nous concluons.
Certains détails sont omis.
\end{proof}

\begin{remark}
\label{remark-commuting-exterior}
La conséquence des lemmes \ref{lemma-chow-cohomology-towards-point}
et \ref{lemma-chow-cohomology-towards-point-commutes} est la suivante.
Soit $k$ un corps. Soit $X$ un schéma localement de type fini sur $k$.
Soit $\alpha \in \CH_*(X)$. Soit $Y \to Z$ un morphisme de schémas
localement de type fini sur $k$. Soit $c' \in A^q(Y \to Z)$. Alors
$$
\alpha \times (c' \cap \beta) = c' \cap (\alpha \times \beta)
$$
dans $\CH_*(X \times_k Y)$ pour tout $\beta \in \CH_*(Z)$. En effet, cela
s'obtient en prenant pour $c = c_\alpha \in A^*(X \to \Spec(k))$ la classe bivariante
correspondant à $\alpha$ ; voir la démonstration du
lemme \ref{lemma-chow-cohomology-towards-point}.
\end{remark}

\begin{lemma}
\label{lemma-exterior-product-associative}
Le produit extérieur est associatif. Plus précisément, soit $k$ un corps,
soient $X, Y, Z$ des schémas localement de type fini sur $k$, et soient
$\alpha \in \CH_*(X)$, $\beta \in \CH_*(Y)$, $\gamma \in \CH_*(Z)$.
Alors $(\alpha \times \beta) \times \gamma =
\alpha \times (\beta \times \gamma)$ dans $\CH_*(X \times_k Y \times_k Z)$.
\end{lemma}

\begin{proof}
Omis. Indication : associativité du produit fibré des schémas.
\end{proof}







\section{Produits d'intersection}
\label{section-intersection-product}

\noindent
Soit $k$ un corps. Dans cette section, nous travaillons sur $S = \Spec(k)$ avec
$\delta : S \to \mathbf{Z}$ définie en envoyant l'unique point sur $0$ ; voir
l'exemple \ref{example-field}.

\medskip\noindent
Soit $X$ un schéma lisse sur $k$. La classe bivariante $\Delta^!$
de la section \ref{section-gysin-for-diagonal} nous permet de définir une sorte de
produit d'intersection sur les groupes de Chow des schémas localement de type fini sur $X$.
Supposons en effet que $Y \to X$ et $Z \to X$ soient des morphismes de schémas
localement de type fini. Observons alors que
$$
Y \times_X Z =  (Y \times_k Z) \times_{X \times_k X, \Delta} X
$$
Nous pouvons donc considérer la suite d'applications suivante
$$
\CH_n(Y) \otimes_\mathbf{Z} \CH_m(Z)
\xrightarrow{\times}
\CH_{n + m}(Y \times_k Z)
\xrightarrow{\Delta^!}
\CH_{n + m - *}(Y \times_X Z)
$$
Ici, la première flèche est le produit extérieur construit dans la
section \ref{section-exterior-product}, et la seconde est l'homomorphisme de Gysin
pour la diagonale étudié dans la
section \ref{section-gysin-for-diagonal}. Si $X$ est équidimensionnel
de dimension $d$, nous aboutissons à $\CH_{n + m - d}(Y \times_X Z)$ ;
en général, nous pouvons décomposer selon les parties situées au-dessus des sous-schémas
ouverts et fermés de $X$ sur lesquels $X$ est de dimension donnée.
Pour $\alpha \in \CH_*(Y)$ et $\beta \in \CH_*(Z)$, nous noterons
$$
\alpha \cdot \beta = \Delta^!(\alpha \times \beta)
\in \CH_*(Y \times_X Z)
$$
Dans le cas particulier où $X = Y = Z$, nous obtenons une multiplication
$$
\CH_*(X) \times \CH_*(X) \to \CH_*(X),\quad
(\alpha, \beta) \mapsto \alpha \cdot \beta
$$
appelée {\it produit d'intersection}. Nous observons que
ce produit est manifestement symétrique. Son associativité résulte du
lemme suivant.

\begin{lemma}
\label{lemma-associative}
Le produit défini ci-dessus est associatif. Plus précisément, soit $k$ un corps,
soit $X$ lisse sur $k$,
soient $Y, Z, W$ des schémas localement de type fini sur $X$, et soient
$\alpha \in \CH_*(Y)$, $\beta \in \CH_*(Z)$, $\gamma \in \CH_*(W)$.
Alors $(\alpha \cdot \beta) \cdot \gamma =
\alpha \cdot (\beta \cdot \gamma)$ dans $\CH_*(Y \times_X Z \times_X W)$.
\end{lemma}

\begin{proof}
Par le lemme \ref{lemma-exterior-product-associative}, nous avons
$(\alpha \times \beta) \times \gamma =
\alpha \times (\beta \times \gamma)$ dans $\CH_*(Y \times_k Z \times_k W)$.
Considérons les immersions fermées
$$
\Delta_{12} : X \times_k X \longrightarrow X \times_k X \times_k X,
\quad (x, x') \mapsto (x, x, x')
$$
et
$$
\Delta_{23} : X \times_k X \longrightarrow X \times_k X \times_k X,
\quad (x, x') \mapsto (x, x', x')
$$
Notons $\Delta_{12}^!$ et $\Delta_{23}^!$ les classes bivariantes
correspondantes ; observons que $\Delta_{12}^!$ est la restriction
(remarque \ref{remark-restriction-bivariant}) de $\Delta^!$
à $X \times_k X \times_k X$ par l'application $\text{pr}_{12}$ et que
$\Delta_{23}^!$ est la restriction de $\Delta^!$
à $X \times_k X \times_k X$ par l'application $\text{pr}_{23}$.
Ainsi, la restriction de $\Delta_{12}^!$ par $\Delta_{23}$ est manifestement
$\Delta^!$, et la restriction de $\Delta_{23}^!$ par $\Delta_{12}$ est
également $\Delta^!$. Par le lemme \ref{lemma-gysin-commutes}, nous avons donc
$$
\Delta^! \circ \Delta_{12}^! =
\Delta^! \circ \Delta_{23}^! 
$$
Nous pouvons maintenant démontrer le lemme par la suite d'égalités suivante :
\begin{align*}
(\alpha \cdot \beta) \cdot \gamma
& =
\Delta^!(\Delta^!(\alpha \times \beta) \times \gamma) \\
& =
\Delta^!(\Delta_{12}^!((\alpha \times \beta) \times \gamma)) \\
& =
\Delta^!(\Delta_{23}^!((\alpha \times \beta) \times \gamma)) \\
& =
\Delta^!(\Delta_{23}^!(\alpha \times (\beta \times \gamma)) \\
& =
\Delta^!(\alpha \times \Delta^!(\beta \times \gamma)) \\
& =
\alpha \cdot (\beta \cdot \gamma)
\end{align*}
Toutes les égalités découlent clairement de ce qui précède, sauf peut-être
la deuxième et l'avant-dernière. L'égalité
$\Delta_{23}^!(\alpha \times (\beta \times \gamma)) =
\alpha \times \Delta^!(\beta \times \gamma)$ résulte de la
remarque \ref{remark-commuting-exterior}. Il en va de même pour la deuxième
égalité.
\end{proof}

\begin{lemma}
\label{lemma-identify-chow-for-smooth}
Soit $k$ un corps. Soit $X$ un schéma lisse sur $k$, équidimensionnel
de dimension $d$. L'application
$$
A^p(X) \longrightarrow \CH_{d - p}(X),\quad
c \longmapsto c \cap [X]_d
$$
est un isomorphisme. Par cet isomorphisme, la composition des classes
bivariantes devient le produit d'intersection défini ci-dessus.
\end{lemma}

\begin{proof}
Notons $g : X \to \Spec(k)$ le morphisme structural.
L'application est la composée des isomorphismes
$$
A^p(X) \to A^{p - d}(X \to \Spec(k)) \to \CH_{d - p}(X)
$$
Le premier est l'isomorphisme $c \mapsto c \circ g^*$ de la
proposition \ref{proposition-compute-bivariant},
et le second est l'isomorphisme $c \mapsto c \cap [\Spec(k)]$ du
lemme \ref{lemma-chow-cohomology-towards-point}.
La démonstration du lemme \ref{lemma-chow-cohomology-towards-point}
montre que l'inverse de la seconde flèche envoie $\alpha \in \CH_{d - p}(X)$
sur la classe bivariante $c_\alpha$ qui envoie $\beta \in \CH_*(Y)$,
pour $Y$ localement de type fini sur $k$,
sur $\alpha \times \beta$ dans $\CH_*(X \times_k Y)$. La démonstration de la
proposition \ref{proposition-compute-bivariant} montre à son tour que l'inverse
de la première flèche envoie $c_\alpha$ sur la classe bivariante
qui envoie $\beta \in \CH_*(Y)$, pour $Y \to X$ localement de type fini,
sur $\Delta^!(\alpha \times \beta) = \alpha \cdot \beta$.
Le dernier énoncé du lemme en résulte.
\end{proof}

\begin{lemma}
\label{lemma-lci-gysin-product}
Soit $k$ un corps. Soit $f : X \to Y$ un morphisme de schémas lisses
sur $k$. Alors l'homomorphisme de Gysin existe pour $f$ et
$f^!(\alpha \cdot \beta) = f^!\alpha \cdot f^!\beta$.
\end{lemma}

\begin{proof}
Observons que $X \to X \times_k Y$ est une immersion de $X$ dans un schéma
lisse sur $Y$. L'homomorphisme de Gysin existe donc pour $f$
(définition \ref{definition-lci-gysin}).
Pour démontrer la formule, nous pouvons décomposer $X$ et $Y$ en leurs
composantes connexes ; nous pouvons donc supposer $X$ lisse sur $k$
et équidimensionnel de dimension $d$, et $Y$ lisse sur $k$
et équidimensionnel de dimension $e$. Observons que
$f^![Y]_e = [X]_d$ (voir par exemple le lemme \ref{lemma-lci-gysin-easy}).
Écrivons $\alpha = c \cap [Y]_e$ et $\beta = c' \cap [Y]_e$,
de sorte que $\alpha \cdot \beta = c \cap c' \cap [Y]_e$ ;
voir le lemme \ref{lemma-identify-chow-for-smooth}.
Par le lemme \ref{lemma-lci-gysin-commutes}, nous savons que $f^!$
commute à la fois avec $c$ et $c'$. Par conséquent,
\begin{align*}
f^!(\alpha \cdot \beta)
& =
f^!(c \cap c' \cap [Y]_e) \\
& =
c \cap c' \cap f^![Y]_e \\
& =
c \cap c' \cap [X]_d \\
& =
(c \cap [X]_d) \cdot (c' \cap [X]_d) \\
& =
(c \cap f^![Y]_e) \cdot (c' \cap f^![Y]_e) \\
& =
f^!(\alpha) \cdot f^!(\beta)
\end{align*}
comme souhaité, où nous avons utilisé le Lemme \ref{lemma-identify-chow-for-smooth}
pour $X$ également.

\medskip\noindent
On peut donner une autre démonstration en prouvant que
$(f \times f)^!(\alpha \times \beta) = f^!\alpha \times f^!\beta$
et en utilisant le lemme \ref{lemma-lci-gysin-composition}.
\end{proof}

\begin{lemma}
\label{lemma-projection-formula}
Soit $k$ un corps. Soit $f : X \to Y$ un morphisme propre de schémas lisses
sur $k$. Alors l'homomorphisme de Gysin existe pour $f$ et
$f_*(\alpha \cdot f^!\beta) = f_*\alpha \cdot \beta$.
\end{lemma}

\begin{proof}
Observons que $X \to X \times_k Y$ est une immersion de $X$ dans un schéma
lisse sur $Y$. L'homomorphisme de Gysin existe donc pour $f$
(définition \ref{definition-lci-gysin}).
Pour démontrer la formule, nous pouvons décomposer $X$ et $Y$ en leurs
composantes connexes ; nous pouvons donc supposer $X$ lisse sur $k$
et équidimensionnel de dimension $d$, et $Y$ lisse sur $k$
et équidimensionnel de dimension $e$. Observons que
$f^![Y]_e = [X]_d$ (voir par exemple le lemme \ref{lemma-lci-gysin-easy}).
Écrivons $\alpha = c \cap [X]_d$ et $\beta = c' \cap [Y]_e$ ;
voir le lemme \ref{lemma-identify-chow-for-smooth}. Nous avons
\begin{align*}
f_*(\alpha \cdot f^!\beta)
& =
f_*(c \cap f^!(c' \cap [Y]_e)) \\
& =
f_*(c \cap c' \cap f^![Y]_e) \\
& =
f_*(c \cap c' \cap [X]_d) \\
& =
f_*(c' \cap c \cap [X]_d) \\
& =
c' \cap f_*(c \cap [X]_d) \\
& =
\beta \cdot f_*(\alpha)
\end{align*}
La première égalité résulte du lemme \ref{lemma-identify-chow-for-smooth}
appliqué à $X$. Par le lemme \ref{lemma-lci-gysin-commutes}, $f^!$
commute avec $c'$. La commutativité du produit
d'intersection justifie d'inverser l'ordre des intersections de $[X]_d$ avec $c$ et $c'$
(par le lemme). On peut commuter $c'$ avec $f_*$ puisque $c'$
est une classe bivariante. La dernière égalité résulte encore du lemme.
\end{proof}

\begin{lemma}
\label{lemma-intersect-properly}
Soit $k$ un corps. Soit $X$ un schéma intègre lisse sur $k$.
Soient $Y, Z \subset X$ des sous-schémas fermés intègres. Posons
$d = \dim(Y) + \dim(Z) - \dim(X)$. Supposons que
\begin{enumerate}
\item $\dim(Y \cap Z) \leq d$, et
\item $\mathcal{O}_{Y, \xi}$ et $\mathcal{O}_{Z, \xi}$
soient de Cohen-Macaulay pour tout $\xi \in Y \cap Z$ tel que
$\delta(\xi) = d$.
\end{enumerate}
Alors $[Y] \cdot [Z] = [Y \cap Z]_d$ dans $\CH_d(X)$.
\end{lemma}

\begin{proof}
Rappelons que $[Y] \cdot [Z] = \Delta^!([Y \times Z])$, où
$\Delta^! = c(\Delta : X \to X \times X, \mathcal{T}_{X/k})$
est un homomorphisme de Gysin en codimension supérieure
(section \ref{section-gysin-higher-codimension}), avec
$\mathcal{T}_{X/k} = \SheafHom(\Omega_{X/k}, \mathcal{O}_X)$
localement libre de rang $\dim(X)$. Nous avons l'égalité de schémas
$$
Y \cap Z = X \times_{\Delta, (X \times X)} (Y \times Z)
$$
et $\dim(Y \times Z) = \dim(Y) + \dim(Z)$ ; les conditions 
(1), (2) et (3) du lemme \ref{lemma-gysin-easy} sont donc satisfaites.
Enfin, si $\xi \in Y \cap Z$, nous avons un homomorphisme local
plat
$$
\mathcal{O}_{Y, \xi} \longrightarrow
\mathcal{O}_{Y \times Z, \xi}
$$
dont la « fibre » est $\mathcal{O}_{Z, \xi}$. Il s'ensuit que si
$\mathcal{O}_{Y, \xi}$ et $\mathcal{O}_{Z, \xi}$
sont tous deux de Cohen-Macaulay, alors $\mathcal{O}_{Y \times Z, \xi}$ l'est aussi ; voir
Algèbre, lemme \ref{algebra-lemma-CM-goes-up}.
Nous voyons ainsi que toutes les hypothèses du
lemme \ref{lemma-gysin-easy} sont satisfaites, et nous concluons.
\end{proof}

\begin{lemma}
\label{lemma-intersect-regularly-embedded}
Soit $k$ un corps. Soit $X$ un schéma lisse sur $k$. Soit $i : Y \to X$
une immersion fermée régulière. Soit $\alpha \in \CH_*(X)$. Si $Y$ est
équidimensionnel de dimension $e$, alors
$\alpha \cdot [Y]_e = i_*(i^!(\alpha))$ dans $\CH_*(X)$.
\end{lemma}

\begin{proof}
Après avoir décomposé $X$ en composantes connexes, nous pouvons supposer, et supposons, $X$
équidimensionnel de dimension $d$. Écrivons $\alpha = c \cap [X]_n$
avec $x \in A^*(X)$ ; voir le lemme \ref{lemma-identify-chow-for-smooth}. Alors
$$
i_*(i^!(\alpha)) = i_*(i^!(c \cap [X]_n)) =
i_*(c \cap i^![X]_n) = i_*(c \cap [Y]_e) =
c \cap i_*[Y]_e = \alpha \cdot [Y]_e
$$
La première égalité résulte du choix de $c$, puis la deuxième du
lemme \ref{lemma-lci-gysin-commutes}. La troisième vient de ce que
$i^![X]_d = [Y]_e$ dans $\CH_*(Y)$ (lemme \ref{lemma-lci-gysin-easy}).
La quatrième tient à ce que les classes bivariantes commutent à l'image directe propre.
La dernière égalité résulte du lemme \ref{lemma-identify-chow-for-smooth}.
\end{proof}

\begin{lemma}
\label{lemma-intersection-regular-smooth}
Soit $k$ un corps. Soit $X$ un schéma lisse sur $k$, qui est
quasi-compact et dont la diagonale est affine. Alors le produit
d'intersection sur $\CH^*(X)$ construit dans cette section coïncide,
après tensorisation par $\mathbf{Q}$, avec le produit d'intersection
construit dans la section \ref{section-intersection-regular}.
\end{lemma}

\begin{proof}
Soient $\alpha \in \CH^i(X)$ et $\beta \in \CH^j(X)$. Écrivons
$\alpha = ch(\alpha') \cap [X]$ et $\beta = ch(\beta') \cap [X]$
$\alpha', \beta' \in K_0(\textit{Vect}(X)) \otimes \mathbf{Q}$
comme dans la section \ref{section-intersection-regular}.
Posons $c = ch(\alpha')$ et $c' = ch(\beta')$.
Le produit d'intersection de la section \ref{section-intersection-regular}
donne alors $c \cap c' \cap [X]$. C'est égal à $\alpha \cdot \beta$
par le lemme \ref{lemma-identify-chow-for-smooth} (ou plutôt par sa
généralisation selon laquelle $A^i(X) \to \CH^i(X)$, $c \mapsto c \cap [X]$
est un isomorphisme pour tout schéma lisse $X$ sur $k$).
\end{proof}








\section{Produit extérieur sur les anneaux de Dedekind}
\label{section-exterior-product-dim-1}

\noindent
Soit $S$ un schéma localement noethérien qui possède un recouvrement ouvert
par des spectres d'anneaux de Dedekind. Posons $\delta(s) = 0$ pour $s \in S$ fermé
et $\delta(s) = 1$ sinon. Alors $(S, \delta)$ est un cas particulier de notre
situation générale \ref{situation-setup} ; voir
l'exemple \ref{example-domain-dimension-1}.
Observons que $S$ est normal
(Algèbre, lemme \ref{algebra-lemma-characterize-Dedekind})
et est donc une réunion disjointe de schémas intègres normaux
(Propriétés, lemme \ref{properties-lemma-normal-locally-Noetherian}).
Ainsi, tous les arguments ci-dessous se ramènent au cas où $S$ est
irréductible. En revanche, nous permettons que $S$ soit non séparé (ainsi, $S$
pourrait être la droite affine avec $0$ dédoublé, par exemple).

\medskip\noindent
Considérons un carré cartésien
$$
\xymatrix{
X \times_S Y \ar[r] \ar[d] & Y \ar[d] \\
X \ar[r] & S
}
$$
de schémas localement de type fini sur $S$. Nous affirmons qu'il existe une application canonique
$$
\times :
\CH_n(X) \otimes_{\mathbf{Z}} \CH_m(Y)
\longrightarrow
\CH_{n + m - 1}(X \times_S Y)
$$
qui est uniquement déterminée par la règle suivante :
étant donnés des sous-schémas fermés intègres $X' \subset X$
et $Y' \subset Y$ de $\delta$-dimensions $n$ et $m$, nous posons
\begin{enumerate}
\item $[X'] \times [Y'] = [X' \times_S Y']_{n + m - 1}$ si
$X'$ ou $Y'$ domine une composante irréductible de $S$,
\item $[X'] \times [Y'] = 0$ si ni $X'$ ni $Y'$ ne domine de
composante irréductible de $S$.
\end{enumerate}

\begin{lemma}
\label{lemma-exterior-product-well-defined-dim-1}
L'application $\times : \CH_n(X) \otimes_{\mathbf{Z}} \CH_m(Y) \to
\CH_{n + m - 1}(X \times_S Y)$ est bien définie.
\end{lemma}

\begin{proof}
Considérons les cycles de dimensions $n$ et $m$, $\alpha = \sum_{i \in I} n_i[X_i]$
et $\beta = \sum_{j \in J} m_j[Y_j]$, où $X_i \subset X$ et $Y_j \subset Y$
sont des familles localement finies de sous-schémas fermés intègres de
$\delta$-dimensions $n$ et $m$. Soit $K \subset I \times J$ l'ensemble
des couples $(i, j) \in I \times J$ tels que $X_i$ ou $Y_j$ domine
une composante irréductible de $S$.
Alors $\{X_i \times_S Y_j\}_{(i, j) \in K}$ est une famille localement finie
de sous-schémas fermés de $X \times_S Y$, de
$\delta$-dimension $n + m - 1$. Nous pouvons donc bien considérer
$$
\alpha \times \beta =
\sum\nolimits_{(i, j) \in K} n_i m_j [X_i \times_S Y_j]_{n + m - 1}
$$
comme un $(n + m - 1)$-cycle sur $X \times_S Y$. Nous obtenons ainsi une
application additive
$\times : Z_n(X) \otimes_{\mathbf{Z}} Z_m(Y) \to Z_{n + m}(X \times_S Y)$.
Le problème est de montrer que
ce procédé est compatible avec l'équivalence rationnelle.

\medskip\noindent
Soit $i : X' \to X$ le morphisme d'inclusion d'un sous-schéma fermé intègre
de $\delta$-dimension $n$ qui domine une composante irréductible
de $S$. Alors $p' : X' \to S$ est plat de dimension relative $n - 1$ ; voir
Compléments d'algèbre, lemme \ref{more-algebra-lemma-dedekind-torsion-free-flat}.
L'image inverse plate par $p'$ est donc un élément
$(p')^* \in A^{-n + 1}(X' \to S)$ par le
lemme \ref{lemma-flat-pullback-bivariant},
et donc $c' = i_* \circ (p')^* \in A^{-n + 1}(X \to S)$ par le
lemme \ref{lemma-push-proper-bivariant}.
Cela donne des applications
$$
c' \cap - : \CH_m(Y) \longrightarrow \CH_{m + n - 1}(X \times_S Y)
$$
qui envoient $[Y']$ sur $[X' \times_S Y']_{n + m - 1}$ pour tout
sous-schéma fermé intègre $Y' \subset Y$ de $\delta$-dimension $m$.

\medskip\noindent
Soit $i : X' \to X$ le morphisme d'inclusion d'un sous-schéma fermé intègre
de $\delta$-dimension $n$ tel que la composée $X' \to X \to S$ 
se factorise par un point fermé $s \in S$. Comme $s$ est un point fermé
du spectre d'un anneau de Dedekind, nous voyons que $s$ est un diviseur de
Cartier effectif sur $S$ dont le fibré normal est trivial. Notons
$c \in A^1(s \to S)$ l'homomorphisme de Gysin ; voir le
lemme \ref{lemma-gysin-bivariant}. Le morphisme $p' : X' \to s$
est plat de dimension relative $n$. L'image inverse plate par $p'$
est donc un élément $(p')^* \in A^{-n}(X' \to S)$ par le
lemme \ref{lemma-flat-pullback-bivariant}.
Ainsi,
$$
c' = i_* \circ (p')^* \circ c \in A^{-n}(X \to S)
$$
par le lemme \ref{lemma-push-proper-bivariant}. Cela donne des applications
$$
c' \cap - : \CH_m(Y) \longrightarrow \CH_{m + n - 1}(X \times_S Y)
$$
qui, pour tout sous-schéma fermé intègre $Y' \subset Y$
de $\delta$-dimension $m$,
envoient $[Y']$ sur $[X' \times_S Y']_{n + m - 1}$ si $Y'$ domine
une composante irréductible de $S$, et sur $0$ sinon.

\medskip\noindent
Des deux paragraphes précédents, nous concluons que
la construction $([X'], [Y']) \mapsto [X' \times_S Y']_{n + m - 1}$
se factorise par l'équivalence rationnelle en la seconde variable, c'est-à-dire
qu'elle donne une application bien définie
$Z_n(X) \otimes_{\mathbf{Z}} \CH_m(Y) \to \CH_{n + m - 1}(X \times_S Y)$.
Par symétrie, il en va de même pour l'autre variable, et nous concluons.
\end{proof}

\begin{lemma}
\label{lemma-chow-cohomology-towards-base-dim-1}
Soit $(S, \delta)$ comme ci-dessus. Soit $X$ un schéma localement de type fini
sur $S$. Nous avons alors une identification canonique
$$
A^p(X \to S) = \CH_{1 - p}(X)
$$
pour tout $p \in \mathbf{Z}$.
\end{lemma}

\begin{proof}
Considérons l'élément $[S]_1 \in \CH_1(S)$. Nous obtenons une application
$A^p(X \to S) \to \CH_{1 - p}(X)$ en envoyant $c$ sur $c \cap [S]_1$.

\medskip\noindent
Réciproquement, supposons donné $\alpha \in \CH_{1 - p}(X)$.
Nous pouvons alors définir $c_\alpha \in A^p(X \to S)$ de la manière
suivante : étant donnés $X' \to S$ et $\alpha' \in \CH_n(X')$,
nous posons
$$
c_\alpha \cap \alpha' = \alpha \times \alpha'
$$
dans $\CH_{n - p}(X \times_S X')$. Pour montrer qu'il s'agit d'une classe
bivariante, écrivons $\alpha = \sum_{i \in I} n_i[X_i]$ comme dans la
définition \ref{definition-cycles}. En particulier, le morphisme
$$
g : \coprod\nolimits_{i \in I} X_i \longrightarrow X
$$
est propre. Fixons $i \in I$. Si $X_i$ domine une composante irréductible
de $S$, alors le morphisme structural $p_i : X_i \to S$ est plat et nous avons
$\xi_i = p_i^* \in A^p(X_i \to S)$. D'autre part, si $p_i$ se factorise
en $p'_i : X_i \to s_i$ suivi de l'inclusion $s_i \to S$
d'un point fermé, alors nous avons
$\xi_i = (p'_i)^* \circ c_i \in A^p(X_i \to S)$
où $c_i \in A^1(s_i \to S)$ est l'homomorphisme de Gysin et
$(p'_i)^*$ est l'image inverse plate. Observons que
$$
A^p(\coprod\nolimits_{i \in I} X_i \to S) =
\prod\nolimits_{i \in I} A^p(X_i \to S)
$$
Nous avons donc
$$
\xi = \sum n_i \xi_i \in A^p(\coprod\nolimits_{i \in I} X_i \to S)
$$
Enfin, puisque $g$ est propre, nous avons une classe bivariante
$$
g_* \circ \xi \in A^p(X \to S)
$$
par le lemme \ref{lemma-push-proper-bivariant}. Le lecteur vérifie aisément
que $c_\alpha$ est égale à cette classe
(comparer avec la démonstration du
lemme \ref{lemma-exterior-product-well-defined-dim-1})
et est donc elle-même une classe bivariante.

\medskip\noindent
Pour achever la démonstration, nous devons montrer que les deux constructions
sont réciproques l'une de l'autre. Comme $c_\alpha \cap [S]_1 = \alpha$,
c'est clair dans un sens. Pour l'autre sens, soit
$c \in A^p(X \to S)$ et posons $\alpha = c \cap [S]_1$.
Il suffit de démontrer que
$$
c \cap [X'] = c_\alpha \cap [X']
$$
lorsque $X'$ est un schéma intègre localement de type fini sur $S$ ;
voir le lemme \ref{lemma-bivariant-zero}. Or, soit $p' : X' \to S$
est plat de dimension relative $\dim_\delta(X') - 1$, et donc
$[X'] = (p')^*[S]_1$, soit $X' \to S$ se factorise en $X' \to s \to S$,
et donc $[X'] = (p')^*(s \to S)^*[S]_1$. Ainsi, le fait que les
classes bivariantes $c$ et $c_\alpha$ coïncident sur $[S]_1$
implique qu'elles coïncident après intersection avec $[X']$ (puisque les classes bivariantes
commutent aux images inverses plates et aux homomorphismes de Gysin), ce qui achève la démonstration.
\end{proof}

\begin{lemma}
\label{lemma-chow-cohomology-towards-base-dim-1-commutes}
Soit $(S, \delta)$ comme ci-dessus. Soit $X$ un schéma localement de type fini
sur $S$. Soit $c \in A^p(X \to S)$. Soit $Y \to Z$ un morphisme de schémas
localement de type fini sur $S$. Soit $c' \in A^q(Y \to Z)$. Alors
$c \circ c' = c' \circ c$ dans $A^{p + q}(X \times_S Y \to X \times_S Z)$.
\end{lemma}

\begin{proof}
Dans la démonstration du lemme \ref{lemma-chow-cohomology-towards-base-dim-1},
nous avons vu que $c$ est donnée par une combinaison d'images
directes propres, de multiplications par des entiers sur les composantes
connexes, d'images inverses plates et d'homomorphismes de Gysin. Comme $c'$ commute avec chacune de
ces opérations par définition des classes bivariantes, nous concluons.
Certains détails sont omis.
\end{proof}

\begin{remark}
\label{remark-commuting-exterior-dim-1}
La conséquence des lemmes \ref{lemma-chow-cohomology-towards-base-dim-1}
et \ref{lemma-chow-cohomology-towards-base-dim-1-commutes} est la suivante.
Soit $(S, \delta)$ comme ci-dessus. Soit $X$ un schéma localement de type fini
sur $S$.
Soit $\alpha \in \CH_*(X)$. Soit $Y \to Z$ un morphisme de schémas
localement de type fini sur $S$. Soit $c' \in A^q(Y \to Z)$. Alors
$$
\alpha \times (c' \cap \beta) = c' \cap (\alpha \times \beta)
$$
dans $\CH_*(X \times_S Y)$ pour tout $\beta \in \CH_*(Z)$. En effet, cela
s'obtient en prenant pour $c = c_\alpha \in A^*(X \to S)$ la classe bivariante
correspondant à $\alpha$ ; voir la démonstration du
lemme \ref{lemma-chow-cohomology-towards-base-dim-1}.
\end{remark}

\begin{lemma}
\label{lemma-exterior-product-associative-dim-1}
Le produit extérieur est associatif. Plus précisément, soit $(S, \delta)$
comme ci-dessus, soient $X, Y, Z$ des schémas localement de type fini sur $S$, et soient
$\alpha \in \CH_*(X)$, $\beta \in \CH_*(Y)$, $\gamma \in \CH_*(Z)$.
Alors $(\alpha \times \beta) \times \gamma =
\alpha \times (\beta \times \gamma)$ dans $\CH_*(X \times_S Y \times_S Z)$.
\end{lemma}

\begin{proof}
Omis. Indication : associativité du produit fibré des schémas.
\end{proof}















\section{Produits d'intersection sur les anneaux de Dedekind}
\label{section-intersection-product-dim-1}

\noindent
Soit $S$ un schéma localement noethérien qui admet un recouvrement ouvert
par des spectres d'anneaux de Dedekind. Posons $\delta(s) = 0$ pour $s \in S$ fermé
et $\delta(s) = 1$ sinon. Alors $(S, \delta)$ est un cas particulier de notre
Situation générale \ref{situation-setup} ; voir
l'Exemple \ref{example-domain-dimension-1} et la discussion de la
section \ref{section-exterior-product-dim-1}.

\medskip\noindent
Soit $X$ un schéma lisse sur $S$. La classe bivariante $\Delta^!$
de la section \ref{section-gysin-for-diagonal} nous permet de définir une sorte de
produit d'intersection sur les groupes de Chow des schémas localement de type fini sur $X$.
Plus précisément, supposons que $Y \to X$ et $Z \to X$ soient des morphismes de schémas
qui sont localement de type fini. Observons alors que
$$
Y \times_X Z =  (Y \times_S Z) \times_{X \times_S X, \Delta} X
$$
Nous pouvons donc considérer la suite d'applications suivante
$$
\CH_n(Y) \otimes_\mathbf{Z} \CH_m(Y)
\xrightarrow{\times}
\CH_{n + m - 1}(Y \times_S Z)
\xrightarrow{\Delta^!}
\CH_{n + m - *}(Y \times_X Z)
$$
Ici, la première flèche est le produit extérieur construit à la
section \ref{section-exterior-product-dim-1} et la seconde flèche
est l'homomorphisme de Gysin pour la diagonale étudié à la
section \ref{section-gysin-for-diagonal}. Si $X$ est équidimensionnel
de dimension $d$, alors $X \to S$ est lisse de dimension relative $d - 1$
et nous aboutissons donc à $\CH_{n + m - d}(Y \times_X Z)$.
En général, nous pouvons décomposer selon les parties situées au-dessus des sous-schémas
ouverts et fermés de $X$ où $X$ a une dimension donnée.
Étant donnés $\alpha \in \CH_*(Y)$ et $\beta \in \CH_*(Z)$, nous noterons
$$
\alpha \cdot \beta = \Delta^!(\alpha \times \beta)
\in \CH_*(Y \times_X Z)
$$
Dans le cas particulier où $X = Y = Z$, nous obtenons une multiplication
$$
\CH_*(X) \times \CH_*(X) \to \CH_*(X),\quad
(\alpha, \beta) \mapsto \alpha \cdot \beta
$$
appelée le {\it produit d'intersection}. Nous observons que
ce produit est clairement symétrique. L'associativité résulte du
lemme suivant.

\begin{lemma}
\label{lemma-associative-dim-1}
Le produit défini ci-dessus est associatif. Plus précisément, avec
$(S, \delta)$ comme ci-dessus, soit $X$ lisse sur $S$,
soient $Y, Z, W$ des schémas localement de type fini sur $X$, et soient
$\alpha \in \CH_*(Y)$, $\beta \in \CH_*(Z)$, $\gamma \in \CH_*(W)$.
Alors $(\alpha \cdot \beta) \cdot \gamma =
\alpha \cdot (\beta \cdot \gamma)$ dans $\CH_*(Y \times_X Z \times_X W)$.
\end{lemma}

\begin{proof}
Par le Lemme \ref{lemma-exterior-product-associative-dim-1}, nous avons
$(\alpha \times \beta) \times \gamma =
\alpha \times (\beta \times \gamma)$ dans $\CH_*(Y \times_S Z \times_S W)$.
Considérons les immersions fermées
$$
\Delta_{12} : X \times_S X \longrightarrow X \times_S X \times_S X,
\quad (x, x') \mapsto (x, x, x')
$$
et
$$
\Delta_{23} : X \times_S X \longrightarrow X \times_S X \times_S X,
\quad (x, x') \mapsto (x, x', x')
$$
Notons $\Delta_{12}^!$ et $\Delta_{23}^!$ les classes bivariantes
correspondantes ; observons que $\Delta_{12}^!$ est la restriction
(Remarque \ref{remark-restriction-bivariant}) de $\Delta^!$
à $X \times_S X \times_S X$ par l'application $\text{pr}_{12}$ et que
$\Delta_{23}^!$ est la restriction de $\Delta^!$
à $X \times_S X \times_S X$ par l'application $\text{pr}_{23}$.
Ainsi, la restriction de $\Delta_{12}^!$ par $\Delta_{23}$
est clairement $\Delta^!$, et la restriction de $\Delta_{23}^!$ par $\Delta_{12}$ est
également $\Delta^!$. Ainsi, par le Lemme \ref{lemma-gysin-commutes}, nous avons
$$
\Delta^! \circ \Delta_{12}^! =
\Delta^! \circ \Delta_{23}^! 
$$
Nous pouvons maintenant démontrer le lemme par la suite d'égalités suivante :
\begin{align*}
(\alpha \cdot \beta) \cdot \gamma
& =
\Delta^!(\Delta^!(\alpha \times \beta) \times \gamma) \\
& =
\Delta^!(\Delta_{12}^!((\alpha \times \beta) \times \gamma)) \\
& =
\Delta^!(\Delta_{23}^!((\alpha \times \beta) \times \gamma)) \\
& =
\Delta^!(\Delta_{23}^!(\alpha \times (\beta \times \gamma)) \\
& =
\Delta^!(\alpha \times \Delta^!(\beta \times \gamma)) \\
& =
\alpha \cdot (\beta \cdot \gamma)
\end{align*}
Toutes les égalités ressortent clairement de ce qui précède, sauf peut-être
la deuxième et l'avant-dernière. L'égalité
$\Delta_{23}^!(\alpha \times (\beta \times \gamma)) =
\alpha \times \Delta^!(\beta \times \gamma)$ résulte de la
Remarque \ref{remark-commuting-exterior}. Il en va de même pour la deuxième
égalité.
\end{proof}

\begin{lemma}
\label{lemma-identify-chow-for-smooth-dim-1}
Soit $(S, \delta)$ comme ci-dessus. Soit $X$ un schéma lisse sur $S$,
équidimensionnel de dimension $d$. L'application
$$
A^p(X) \longrightarrow \CH_{d - p}(X),\quad
c \longmapsto c \cap [X]_d
$$
est un isomorphisme. Par cet isomorphisme, la composition des classes
bivariantes devient le produit d'intersection défini ci-dessus.
\end{lemma}

\begin{proof}
Notons $g : X \to S$ le morphisme structural.
L'application est la composée des isomorphismes
$$
A^p(X) \to A^{p - d + 1}(X \to S) \to \CH_{d - p}(X)
$$
Le premier est l'isomorphisme $c \mapsto c \circ g^*$ de la
Proposition \ref{proposition-compute-bivariant},
et le second est l'isomorphisme $c \mapsto c \cap [S]_1$ du
Lemme \ref{lemma-chow-cohomology-towards-base-dim-1}.
La démonstration du Lemme \ref{lemma-chow-cohomology-towards-base-dim-1}
montre que l'inverse de la seconde flèche envoie $\alpha \in \CH_{d - p}(X)$
sur la classe bivariante $c_\alpha$ qui envoie $\beta \in \CH_*(Y)$,
pour $Y$ localement de type fini sur $k$,
sur $\alpha \times \beta$ dans $\CH_*(X \times_k Y)$. La démonstration de la
Proposition \ref{proposition-compute-bivariant} montre à son tour que l'inverse
de la première flèche envoie $c_\alpha$ sur la classe bivariante
qui envoie $\beta \in \CH_*(Y)$, pour $Y \to X$ localement de type fini,
sur $\Delta^!(\alpha \times \beta) = \alpha \cdot \beta$.
Le dernier résultat du lemme en découle.
\end{proof}










\section{Classes de Todd}
\label{section-todd-classes}

\noindent
Une dernière classe associée à un fibré vectoriel $\mathcal{E}$
de rang $r$ est sa {\it classe de Todd} $Todd(\mathcal{E})$.
En termes des racines de Chern $x_1, \ldots, x_r$, elle est
définie par
$$
Todd(\mathcal{E})
=
\prod\nolimits_{i = 1}^r
\frac{x_i}{1 - e^{-x_i}}
$$
En termes des classes de Chern $c_i = c_i(\mathcal{E})$,
nous avons
$$
Todd(\mathcal{E})
=
1
+
\frac{1}{2}c_1
+
\frac{1}{12}(c_1^2 + c_2)
+
\frac{1}{24}c_1c_2
+
\frac{1}{720}(-c_1^4 + 4c_1^2c_2 + 3c_2^2 + c_1c_3 - c_4)
+
\ldots
$$
Nous avons fait les remarques appropriées sur les dénominateurs
dans la section précédente. Pour toute suite exacte
donnée
$$
0
\to
{\mathcal E}_1
\to
{\mathcal E}
\to
{\mathcal E}_2
\to
0
$$
nous avons
$$
Todd({\mathcal E}) = Todd({\mathcal E}_1) Todd({\mathcal E}_2).
$$







\section{Grothendieck--Riemann--Roch}
\label{section-grr}

\noindent
Soit $(S, \delta)$ comme dans la
Situation \ref{situation-setup}.
Soient $X, Y$ localement de type fini sur $S$.
Soit $\mathcal{E}$ un faisceau localement libre de type fini
${\mathcal E}$ sur $X$, de rang $r$.
Soit $f : X \to Y$ un morphisme propre et lisse.
Supposons que les $R^if_*\mathcal{E}$ soient des faisceaux localement libres
sur $Y$, de rang fini.
Le théorème de Grothendieck--Riemann--Roch affirme dans ce
cas que
$$
f_*(Todd(T_{X/Y}) ch(\mathcal{E}))
=
\sum (-1)^i ch(R^if_*\mathcal{E})
$$
Ici,
$$
T_{X/Y} = \SheafHom_{\mathcal{O}_X}(\Omega_{X/Y}, \mathcal{O}_X)
$$
est le fibré tangent relatif de $X$ sur $Y$. Si $Y = \Spec(k)$,
où $k$ est un corps, nous pouvons reformuler ceci sous la forme
$$
\chi(X, \mathcal{E}) = \deg(Todd(T_{X/k}) ch(\mathcal{E}))
$$
Le théorème est plus général et devient plus facile à démontrer
lorsqu'il est formulé avec la bonne généralité. Nous y reviendrons
ailleurs (insérer ici une référence future).







\section{Changement de schéma de base}
\label{section-change-base}

\noindent
Dans cette section, nous expliquons comment comparer les théories pour différents
schémas de base.

\begin{situation}
\label{situation-setup-base-change}
Ici, $(S, \delta)$ et $(S', \delta')$ sont comme dans la
Situation \ref{situation-setup}. En outre, $g : S' \to S$
est un morphisme plat de schémas et $c \in \mathbf{Z}$
est un entier tel que : pour tout $s \in S$ et tout $s' \in S'$
qui est un point générique d'une composante irréductible de $g^{-1}(\{s\})$, nous avons
$\delta'(s') = \delta(s) + c$.
\end{situation}

\noindent
Nous verrons que, pour un schéma $X$ localement de type fini sur $S$,
il existe une application bien définie $\CH_k(X) \to \CH_{k + c}(X \times_S S')$
de groupes de Chow qui commute (dans l'ensemble) avec les opérations
que nous avons définies dans ce chapitre.

\begin{lemma}
\label{lemma-dimension-base-change}
Dans la Situation \ref{situation-setup-base-change}, soit $X \to S$ localement
de type fini. Notons $X' \to S'$ le changement de base par $S' \to S$.
Si $X$ est intègre avec $\dim_\delta(X) = k$, alors
toute composante irréductible $Z'$ de $X'$ vérifie $\dim_{\delta'}(Z') = k + c$.
\end{lemma}

\begin{proof}
La projection $X' \to X$ est plate comme changement de base du morphisme plat
$S' \to S$ (Morphismes, Lemme \ref{morphisms-lemma-base-change-flat}).
Ainsi, tout point générique $x'$ d'une composante irréductible
de $X'$ s'envoie sur le point générique $x \in X$ (car les générisations
se relèvent le long de $X' \to X$ par
Morphismes, Lemme \ref{morphisms-lemma-generalizations-lift-flat}).
Soit $s \in S$ l'image de $x$.
Rappelons que le schéma $S'_s = S' \times_S s$
a le même espace topologique sous-jacent que $g^{-1}(\{s\})$
(Schémas, Lemme \ref{schemes-lemma-fibre-topological}).
Nous pouvons voir $x'$ comme un point du schéma $S'_s \times_s x$, lequel
est muni d'un monomorphisme $S'_s \times_s x \to S' \times_S X$.
Bien entendu, $x'$ est aussi un point générique d'une composante irréductible
de $S'_s \times_s x$.
En utilisant la platitude de $\Spec(\kappa(x)) \to \Spec(\kappa(s)) = s$
et en raisonnant comme ci-dessus, nous voyons que $x'$ s'envoie sur un point générique $s'$
d'une composante irréductible de $g^{-1}(\{s\})$. Ainsi,
$\delta'(s') = \delta(s) + c$ par hypothèse.
Nous avons $\dim_x(X_s) = \dim_{x'}(X_{s'})$ par
Morphismes, Lemme \ref{morphisms-lemma-dimension-fibre-after-base-change}.
Puisque $x$ est un point générique d'une composante irréductible $X_s$
(c'est un schéma irréductible, mais nous n'en avons pas besoin) et que
$x'$ est un point générique d'une composante irréductible de $X'_{s'}$, nous concluons
que $\text{trdeg}_{\kappa(s)}(\kappa(x)) = 
\text{trdeg}_{\kappa(s')}(\kappa(x'))$
par Morphismes, Lemme \ref{morphisms-lemma-dimension-fibre-at-a-point}.
Alors
$$
\delta_{X'/S'}(x') = \delta(s') + \text{trdeg}_{\kappa(s')}(\kappa(x')) =
\delta(s) + c + \text{trdeg}_{\kappa(s)}(\kappa(x)) = \delta_{X/S}(x) + c
$$
Ceci prouve ce que nous voulions par la Définition \ref{definition-delta-dimension}.
\end{proof}

\noindent
Dans la Situation \ref{situation-setup-base-change}, soit $X \to S$ localement
de type fini. Notons $X' \to S'$ le changement de base par $g : S' \to S$.
Il existe un unique homomorphisme
$$
g^* : Z_k(X) \longrightarrow Z_{k + c}(X')
$$
qui, étant donné un sous-schéma fermé intègre $Z \subset X$ de
$\delta$-dimension $k$, envoie $[Z]$ sur $[Z \times_S S']_{k + c}$.
Ceci a un sens par le Lemme \ref{lemma-dimension-base-change}.

\begin{lemma}
\label{lemma-pullback-coherent-base-change}
Dans la Situation \ref{situation-setup-base-change}, soit $X \to S$
localement de type fini et soit $X' \to S$ le changement de base par $S' \to S$.
\begin{enumerate}
\item Soit $Z \subset X$ un sous-schéma fermé tel que
$\dim_\delta(Z) \leq k$ et de changement de base $Z' \subset X'$. Alors nous avons
$\dim_{\delta'}(Z')) \leq k + c$
et $[Z']_{k + c} = g^*[Z]_k$ dans $Z_{k + c}(X')$.
\item Soit $\mathcal{F}$ un faisceau cohérent sur $X$ tel que
$\dim_\delta(\text{Supp}(\mathcal{F})) \leq k$, et soit
$\mathcal{F}'$ son changement de base sur $X'$.
Alors nous avons $\dim_\delta(\text{Supp}(\mathcal{F}')) \leq k + c$
et $g^*[\mathcal{F}]_k = [\mathcal{F}']_{k + c}$
dans $Z_{k + c}(X')$.
\end{enumerate}
\end{lemma}

\begin{proof}
La démonstration est exactement la même que celle du
Lemme \ref{lemma-pullback-coherent},
et nous suggérons au lecteur de la passer.

\medskip\noindent
Les assertions sur les dimensions résultent du
Lemme \ref{lemma-dimension-base-change}.
La partie (1) résulte de la partie (2) par le Lemme \ref{lemma-cycle-closed-coherent}
et par le fait que le changement de base du module cohérent $\mathcal{O}_Z$
est $\mathcal{O}_{Z'}$.

\medskip\noindent
Démonstration de (2). Puisque $X$, $X'$ sont localement noethériens, nous pouvons appliquer
Cohomologie des schémas, Lemme \ref{coherent-lemma-coherent-Noetherian}, pour voir
que $\mathcal{F}$ est de type fini, donc $\mathcal{F}'$ est
de type fini (Modules, Lemme \ref{modules-lemma-pullback-finite-type}),
et donc $\mathcal{F}'$ est cohérent
(de nouveau Cohomologie des schémas, Lemme \ref{coherent-lemma-coherent-Noetherian}).
Ainsi, le lemme a un sens. Soit $W \subset X$ un sous-schéma fermé intègre
de $\delta$-dimension $k$, et soit $W' \subset X'$ un sous-schéma
fermé intègre de $\delta'$-dimension $k + c$ s'envoyant dans $W$
par $X' \to X$. Nous devons montrer que le coefficient $n$ de
$[W']$ dans $g^*[\mathcal{F}]_k$ coïncide avec le coefficient
$m$ de $[W']$ dans $[\mathcal{F}']_{k + c}$. Soient $\xi \in W$ et
$\xi' \in W'$ les points génériques. Posons
$A = \mathcal{O}_{X, \xi}$, $B = \mathcal{O}_{X', \xi'}$
et $M = \mathcal{F}_\xi$ comme $A$-module. (Notons que
$M$ est de longueur finie par nos hypothèses de dimension, mais qu'en fait nous
n'avons pas besoin de le vérifier. Voir le
Lemme \ref{lemma-length-finite}.)
Nous avons $\mathcal{F}'_{\xi'} = B \otimes_A M$.
Ainsi, nous voyons que
$$
n = \text{longueur}_B(B \otimes_A M)
\quad
\text{et}
\quad
m = \text{longueur}_A(M) \text{longueur}_B(B/\mathfrak m_AB)
$$
L'égalité résulte donc de
Algèbre, Lemme \ref{algebra-lemma-pullback-module}.
\end{proof}

\begin{lemma}
\label{lemma-pullback-base-change}
Dans la Situation \ref{situation-setup-base-change}, soit $X \to S$ localement
de type fini et soit $X' \to S'$ le changement de base par $S' \to S$.
L'application $g^* : Z_k(X) \to Z_{k + c}(X')$ ci-dessus passe au quotient par l'équivalence
rationnelle et donne une application
$$
g^* : \CH_k(X) \longrightarrow \CH_{k + c}(X')
$$
de groupes de Chow.
\end{lemma}

\begin{proof}
Supposons que $\alpha \in Z_k(X)$ soit un $k$-cycle rationnellement équivalent
à zéro. Par le Lemme \ref{lemma-rational-equivalence-family},
il existe une famille localement finie de sous-schémas fermés intègres
$W_i \subset X \times \mathbf{P}^1$ de $\delta$-dimension $k$
qui ne sont pas contenus dans les diviseurs
$(X \times \mathbf{P}^1)_0$ ou $(X \times \mathbf{P}^1)_\infty$
de $X \times \mathbf{P}^1$, telle que
$\alpha = \sum ([(W_i)_0]_k - [(W_i)_\infty]_k)$.
Il suffit donc de prouver que, pour $W \subset X \times \mathbf{P}^1$
fermé intègre de $\delta$-dimension $k$ et non contenu dans les diviseurs
$(X \times \mathbf{P}^1)_0$ ou $(X \times \mathbf{P}^1)_\infty$
de $X \times \mathbf{P}^1$, nous avons
\begin{enumerate}
\item le changement de base $W' \subset X' \times \mathbf{P}^1$ satisfait aux
hypothèses du Lemme \ref{lemma-closed-subscheme-cross-p1}, avec
$k$ remplacé par $k + c$, et
\item $g^*[W_0]_k = [(W')_0]_{k + c}$ et
$g^*[W_\infty]_k = [(W')_\infty]_{k + c}$.
\end{enumerate}
La partie (2) résulte immédiatement du
Lemme \ref{lemma-pullback-coherent-base-change} et du fait que
$(W')_0$ est le changement de base de $W_0$ (par associativité des produits fibrés).
Pour la partie (1), l'assertion sur les dimensions résulte d'abord
du Lemme \ref{lemma-dimension-base-change}.
Soit ensuite $w' \in (W')_0$, d'image $w \in W_0$
et $z \in \mathbf{P}^1_S$. Notons $t \in \mathcal{O}_{\mathbf{P}^1_S, z}$
l'équation usuelle de $0 : S \to \mathbf{P}^1_S$.
Puisque $\mathcal{O}_{W, w} \to \mathcal{O}_{W', w'}$ est plat
et puisque $t$ est un non-diviseur de zéro sur $\mathcal{O}_{W, w}$
(car $W$ est intègre et $W \not = W_0$), nous voyons que
$t$ est aussi un non-diviseur de zéro dans $\mathcal{O}_{W', w'}$. Ainsi,
$W'$ n'a aucun point associé situé sur $W'_0$.
\end{proof}

\begin{lemma}
\label{lemma-pullback-base-change-pullback}
Dans la Situation \ref{situation-setup-base-change}, soient $Y \to X \to S$ localement
de type fini et soit $Y' \to X' \to S'$ le changement de base par $S' \to S$.
Supposons $f : Y \to X$ plat de dimension relative $r$. Alors $f' : Y' \to X'$
est plat de dimension relative $r$ et les diagrammes
$$
\vcenter{
\xymatrix{
Z_{k + r}(Y) \ar[r]_{g^*} & Z_{k + c + r}(Y') \\
Z_k(X) \ar[r]^{g^*} \ar[u]^{f^*} & Z_{k + c}(X') \ar[u]_{(f')^*}
}
}
\quad\text{et}\quad
\vcenter{
\xymatrix{
\CH_{k + r}(Y) \ar[r]_{g^*} & \CH_{k + c + r}(Y') \\
\CH_k(X) \ar[r]^{g^*} \ar[u]^{f^*} & \CH_{k + c}(X') \ar[u]_{(f')^*}
}
}
$$
de groupes de cycles et de Chow commutent.
\end{lemma}

\begin{proof}
Il suffit de montrer que le premier diagramme commute. Pour le voir, soit
$Z \subset X$ un sous-schéma fermé intègre de $\delta$-dimension $k$
et notons $Z' \subset X'$ son changement de base. Par construction,
nous avons $g^*[Z] = [Z']_{k + c}$. Par le Lemme \ref{lemma-pullback-coherent},
nous avons $(f')^*g^*[Z] = [Z' \times_{X'} Y']_{k + c + r}$.
Réciproquement, nous avons $f^*[Z] = [Z \times_X Y]_{k + r}$ par la
Définition \ref{definition-flat-pullback}. Par le
Lemme \ref{lemma-pullback-coherent-base-change},
nous avons $g^*f^*[Z] = [(Z \times_X Y)']_{k + r + c}$.
Puisque $(Z \times_X Y)' = Z' \times_{X'} Y'$ par
associativité du produit fibré, nous concluons.
\end{proof}

\begin{lemma}
\label{lemma-pullback-base-change-pushforward}
Dans la Situation \ref{situation-setup-base-change}, soient $Y \to X \to S$ localement
de type fini et soit $Y' \to X' \to S'$ le changement de base par $S' \to S$.
Supposons $f : Y \to X$ propre. Alors $f' : Y' \to X'$ est propre et les diagrammes
$$
\vcenter{
\xymatrix{
Z_k(Y) \ar[r]_{g^*} \ar[d]_{f_*} & Z_{k + c}(Y') \ar[d]^{f'_*} \\
Z_k(X) \ar[r]^{g^*} & Z_{k + c}(X')
}
}
\quad\text{et}\quad
\vcenter{
\xymatrix{
\CH_k(Y) \ar[r]_{g^*} \ar[d]_{f_*} & \CH_{k + c}(Y') \ar[d]^{f'_*} \\
\CH_k(X) \ar[r]^{g^*} & \CH_{k + c}(X')
}
}
$$
de groupes de cycles et de Chow commutent.
\end{lemma}

\begin{proof}
Il suffit de montrer que le premier diagramme commute. Pour le voir, soit
$Z \subset Y$ un sous-schéma fermé intègre de $\delta$-dimension $k$
et notons $Z' \subset X'$ son changement de base. Par construction,
nous avons $g^*[Z] = [Z']_{k + c}$. Par le Lemme \ref{lemma-cycle-push-sheaf},
nous avons $(f')_*g^*[Z] = [f'_*\mathcal{O}_{Z'}]_{k + c}$.
Par le même lemme, nous avons $f_*[Z] = [f_*\mathcal{O}_Z]_k$. Par le
Lemme \ref{lemma-pullback-coherent-base-change},
nous avons $g^*f_*[Z] = [(X' \to X)^*f_*\mathcal{O}_Z]_{k + r}$.
Il suffit donc de montrer que
$$
(X' \to X)^*f_*\mathcal{O}_Z \cong f'_*\mathcal{O}_{Z'}
$$
comme modules cohérents sur $X'$. Puisque $X' \to X$ est plat et que
$\mathcal{O}_{Z'} = (Y' \to Y)^*\mathcal{O}_Z$, ceci
résulte du changement de base plat, voir
Cohomologie des schémas, Lemme \ref{coherent-lemma-flat-base-change-cohomology}.
\end{proof}

\begin{lemma}
\label{lemma-pullback-base-change-c1}
Dans la Situation \ref{situation-setup-base-change}, soit $X \to S$ localement
de type fini et soit $X' \to S'$ le changement de base par $S' \to S$.
Soit $\mathcal{L}$ un $\mathcal{O}_X$-module inversible de
changement de base $\mathcal{L}'$ sur $X'$. Alors le
diagramme
$$
\xymatrix{
\CH_k(X) \ar[r]_{g^*} \ar[d]_{c_1(\mathcal{L}) \cap -} &
\CH_{k + c}(X') \ar[d]^{c_1(\mathcal{L}') \cap -} \\
\CH_{k - 1}(X) \ar[r]^{g^*} & \CH_{k + c - 1}(X')
}
$$
de groupes de Chow commute.
\end{lemma}

\begin{proof}
Soit $p : L \to X$ le fibré en droites associé à $\mathcal{L}$,
de section nulle $o : X \to L$. Pour $\alpha \in CH_k(X)$, nous
savons que $\beta = c_1(\mathcal{L}) \cap \alpha$
est l'unique élément de $\CH_{k - 1}(X)$ tel que
$o_*\alpha = - p^*\beta$, voir les Lemmes \ref{lemma-linebundle} et
\ref{lemma-linebundle-formulae}.
La même caractérisation vaut après image inverse. Le lemme résulte donc des
Lemmes \ref{lemma-pullback-base-change-pullback} et
\ref{lemma-pullback-base-change-pushforward}.
\end{proof}

\begin{lemma}
\label{lemma-pullback-base-change-chern-classes}
Dans la Situation \ref{situation-setup-base-change}, soit $X \to S$ localement
de type fini et soit $X' \to S'$ le changement de base par $S' \to S$.
Soit $\mathcal{E}$ un $\mathcal{O}_X$-module localement libre de type fini de
rang $r$, de changement de base $\mathcal{E}'$ sur $X'$. Alors le
diagramme
$$
\xymatrix{
\CH_k(X) \ar[r]_{g^*} \ar[d]_{c_i(\mathcal{E}) \cap -} &
\CH_{k + c}(X') \ar[d]^{c_i(\mathcal{E}') \cap -} \\
\CH_{k - i}(X) \ar[r]^{g^*} & \CH_{k + c - i}(X')
}
$$
de groupes de Chow commute pour tout $i$.
\end{lemma}

\begin{proof}
Posons $P = \mathbf{P}(\mathcal{E})$. Le changement de base $P'$ de $P$
est égal à $\mathbf{P}(\mathcal{E}')$. Puisque nous savons déjà que
l'image inverse plate et le cup-produit avec $c_1$ d'un module inversible
commutent au changement de base (Lemmes \ref{lemma-pullback-base-change-pullback} et
\ref{lemma-pullback-base-change-c1})
le lemme résulte de la caractérisation du cap-produit
avec $c_i(\mathcal{E})$ donnée dans le Lemme \ref{lemma-determine-intersections}.
\end{proof}

\begin{lemma}
\label{lemma-compose-base-change}
Soient $(S, \delta)$, $(S', \delta')$, $(S'', \delta'')$ comme dans la
Situation \ref{situation-setup}. Soient $g : S' \to S$ et $g' : S'' \to S'$
des morphismes plats de schémas, et soient $c, c' \in \mathbf{Z}$
des entiers tels que $S, \delta, S', \delta', g, c$ et
$S', \delta', S'', g', c'$ soient comme dans la
Situation \ref{situation-setup-base-change}.
Soit $X \to S$ localement de type fini et notons $X' \to S'$
et $X'' \to S''$ les changements de base par $S' \to S$ et $S'' \to S$.
Alors
\begin{enumerate}
\item $S, \delta, S'', \delta'', g \circ g', c + c'$ est comme dans la
Situation \ref{situation-setup-base-change},
\item les applications $g^* : Z_k(X) \to Z_{k + c}(X')$ et
$(g')^* : Z_{k + c}(X') \to Z_{k + c + c'}(X'')$ se
composent pour donner l'application $(g \circ g')^* : Z_k(X) \to Z_{k + c + c'}(X'')$, et
\item les applications $g^* : \CH_k(X) \to \CH_{k + c}(X')$ et
$(g')^* : \CH_{k + c}(X') \to \CH_{k + c + c'}(X'')$ du
Lemme \ref{lemma-pullback-base-change}
se composent pour donner l'application $(g \circ g')^* : \CH_k(X) \to \CH_{k + c + c'}(X'')$
du Lemme \ref{lemma-pullback-base-change}.
\end{enumerate}
\end{lemma}

\begin{proof}
Soit $s \in S$, et soit $s'' \in S''$ un point générique d'une composante
irréductible de $(g \circ g')^{-1}(\{s\})$. Posons $s' = g'(s'')$.
Clairement, $s''$ est un point générique d'une composante irréductible de
$(g')^{-1}(\{s'\})$. De plus, puisque $g'$ est plat et que les générisations
se relèvent donc le long de $g'$ (Morphismes, Lemme \ref{morphisms-lemma-base-change-flat}),
nous voyons que $s'$ est aussi un point générique d'une composante irréductible
de $g^{-1}(\{s\})$. Ainsi, par hypothèse, $\delta'(s') = \delta(s) + c$
et $\delta''(s'') = \delta'(s') + c'$. Nous concluons que
$\delta''(s'') = \delta(s) + c + c'$ et que la première partie de cette
assertion est vraie.

\medskip\noindent
Pour la deuxième partie, soit $Z \subset X$ un sous-schéma fermé intègre
de $\delta$-dimension $k$. Notons $Z' \subset X'$ et $Z'' \subset X''$
les changements de base. Par définition, nous avons $g^*[Z] = [Z']_{k + c}$.
Par le Lemme \ref{lemma-pullback-coherent-base-change}, nous avons
$(g')^*[Z']_{k + c} = [Z'']_{k + c + c'}$. Ceci prouve la dernière assertion.
\end{proof}

\begin{lemma}
\label{lemma-chow-limit}
Dans la Situation \ref{situation-setup-base-change}, supposons $c = 0$
et supposons que $S' = \lim_{i \in I} S_i$ soit une limite projective filtrante
de schémas $S_i$ affines sur $S$ tels que
\begin{enumerate}
\item avec $\delta_i$ égal à $S_i \to S \xrightarrow{\delta} \mathbf{Z}$,
le couple $(S_i, \delta_i)$ est comme dans la Situation \ref{situation-setup},
\item $S_i, \delta_i, S, \delta, S \to S_i, c = 0$ est comme dans la
Situation \ref{situation-setup-base-change},
\item $S_i, \delta_i, S_{i'}, \delta_{i'}, S_i \to S_{i'}, c = 0$ 
pour $i \geq i'$ est comme dans la Situation \ref{situation-setup-base-change}.
\end{enumerate}
Alors, pour un schéma quasi-compact $X$ de type fini sur $S$,
de changements de base $X'$ et $X_i$ par $S' \to S$ et $S_i \to S$, nous avons
$Z_k(X') = \colim Z_k(X_i)$ et $\CH_k(X') = \colim \CH_k(X_i)$.
\end{lemma}

\begin{proof}
Par le résultat du Lemme \ref{lemma-compose-base-change}, nous obtenons
un système de groupes de cycles $Z_k(X_i)$ et un système de
groupes de Chow $\CH_k(X_i)$ ainsi que des applications
$\colim Z_k(X_i) \to Z_k(X')$ et $\colim \CH_i(X_i) \to \CH_k(X')$.
Nous pouvons remplacer $S$ par un ouvert quasi-compact par lequel $X \to S$
se factorise ; nous pouvons donc supposer, et nous supposons, que tous les schémas intervenant dans
cette démonstration sont noethériens (et donc quasi-compacts et quasi-séparés).

\medskip\noindent
Montrons que l'application $\colim Z_k(X_i) \to Z_k(X')$ est surjective.
Soit en effet $Z' \subset X'$
un sous-schéma fermé intègre de $\delta'$-dimension $k$. Par
Limites, Lemme \ref{limits-lemma-descend-finite-presentation},
nous pouvons trouver un $i$ et un morphisme $Z_i \to X_i$ de présentation finie
dont le changement de base est $Z'$. En augmentant $i$, nous pouvons supposer que $Z_i$
est un sous-schéma fermé de $X_i$, voir
Limites, Lemme \ref{limits-lemma-descend-closed-immersion-finite-presentation}.
Alors $Z' \to X_i$ se factorise par $Z_i$ et nous pouvons remplacer $Z_i$
par l'image schématique de $Z' \to X_i$. Nous voyons ainsi
que nous pouvons supposer $Z_i$ fermé intègre dans $X_i$.
Par le Lemme \ref{lemma-dimension-base-change}, nous concluons que
$\dim_{\delta_i}(Z_i) = \dim_{\delta'}(Z') = k$.
Ainsi, $Z_k(X_i) \to Z_k(X')$ envoie $[Z_i]$ sur $[Z']$ et
nous concluons à la surjectivité.

\medskip\noindent
Montrons que l'application $\colim Z_k(X_i) \to Z_k(X')$ est injective.
Soit $\alpha_i = \sum n_j[Z_j] \in Z_k(X_i)$ un cycle dont l'image
dans $Z_k(X')$ est nulle. Nous pouvons supposer, et supposons, $Z_j \not = Z_{j'}$ si
$j \not = j'$ et $n_j \not = 0$ pour tout $j$.
Notons $Z'_j \subset X'$ le changement de base de $Z_j$.
Par le Lemme \ref{lemma-dimension-base-change}, chaque composante irréductible
de $Z'_j$ a pour $\delta'$-dimension $k$. De plus, comme $Z_j$ est irréductible
et $Z'_j \to Z_j$ est plat (comme changement de base de $S' \to S$),
nous voyons que $Z'_j \to Z_j$ est dominant.
Ainsi, si $Z'_j$ est non vide, une composante irréductible, disons
$Z'$, de $Z'_j$ domine $Z_j$. Il s'ensuit que $Z'$
ne peut être une composante irréductible de $Z'_{j'}$ pour $j' \not = j$.
Ainsi, si $Z'_j$ est non vide, nous voyons que $(S' \to S_i)^*\alpha_i
= \sum [Z'_j]_r$ est non nul (car le coefficient de $Z'$ serait non nul).
Nous voyons donc que $Z'_j = \emptyset$ pour tout $j$. Or, cela signifie
que le changement de base de $Z_j$ par une application de transition $S_{i'} \to S_i$
est vide par Limites, Lemme \ref{limits-lemma-limit-nonempty}.
Ainsi, $\alpha_i$ s'annule dans la colimite comme voulu.

\medskip\noindent
La surjectivité de $\colim Z_k(X_i) \to Z_k(X')$ implique que
$\colim \CH_k(X_i) \to \CH_k(X')$ est surjective.
Pour terminer la démonstration, montrons que cette application est injective.
Soit $\alpha_i \in \CH_k(X_i)$
un cycle dont l'image $\alpha' \in \CH_k(X')$ est nulle.
Il existe alors des sous-schémas fermés intègres
$W_l' \subset X'$, $l = 1, \ldots, r$, de $\delta"$-dimension $k + 1$
et des fonctions rationnelles non nulles $f'_l$ sur $W'_l$
tels que $\alpha' = \sum_{l = 1, \ldots, r} \text{div}_{W'_l}(f'_l)$.
En raisonnant comme ci-dessus, nous pouvons trouver un $i$ et des sous-schémas fermés intègres
$W_{i, l} \subset X_i$ de $\delta_i$-dimension $k + 1$
dont le changement de base est $W'_l$.
En augmentant $i$, nous pouvons supposer que nous avons des fonctions rationnelles
$f_{i, l}$ sur $W_{i, l}$. En effet, nous pouvons considérer $f'_l$ comme une
section du faisceau structural sur un ouvert non vide $U'_l \subset W'_l$ ;
nous pouvons descendre ces ouverts par Limites, Lemme \ref{limits-lemma-descend-opens},
et, en augmentant $i$, descendre $f'_l$ par
Limites, Lemme \ref{limits-lemma-descend-section}.
Nous affirmons que
$$
\alpha_i = \sum\nolimits_{l = 1, \ldots, r} \text{div}_{W_{i, l}}(f_{i, l})
$$
après avoir éventuellement augmenté $i$.

\medskip\noindent
Pour prouver l'affirmation, soit $Z'_{l, j} \subset W'_l$ une famille
finie de sous-schémas fermés intègres de $\delta'$-dimension $k$
tels que $f'_l$ soit une fonction régulière inversible en dehors de
$\bigcup_j Y'_{l, j}$. En augmentant $i$ (par les arguments ci-dessus),
nous pouvons supposer qu'il existe des sous-schémas fermés intègres $Z_{i, l, j} \subset W_i$
de $\delta_i$-dimension $k$ tels que $f_{i, l}$ soit une
fonction régulière inversible en dehors de $\bigcup_j Z_{i, l, j}$.
Nous pouvons alors écrire
$$
\text{div}_{W'_l}(f'_l) = \sum n_{l, j} [Z'_{l, j}]
$$
et
$$
\text{div}_{W_{i, l}}(f_{i, l}) = \sum n_{i, l, j} [Z_{i, l, j}]
$$
Pour prouver l'affirmation, il suffit de montrer que $n_{l, i} = n_{i, l, j}$.
En effet, ceci impliquera que $\beta_i =
\alpha_i - \sum\nolimits_{l = 1, \ldots, r} \text{div}_{W_{i, l}}(f_{i, l})$
est un cycle sur $X_i$ dont l'image inverse sur $X'$ est nulle comme cycle !
Il s'ensuit que l'image inverse de $\beta_i$ est nulle comme cycle sur $X_{i'}$
pour un certain $i' \geq i$, par un argument facile que nous omettons.

\medskip\noindent
Pour prouver l'égalité $n_{l, i} = n_{i, l, j}$, choisissons un
point générique $\xi' \in Z'_{l, j}$ et notons
$\xi \in Z_{i, l, j}$ son image, qui est également un point générique.
Alors, l'homomorphisme d'anneaux locaux
$$
\mathcal{O}_{W_{i, l}, \xi}
\longrightarrow
\mathcal{O}_{W'_l, \xi'}
$$
est plat car $W'_l \to W_{i, l}$ est le changement de base du morphisme
plat $S' \to S_i$. Nous avons aussi
$\mathfrak m_\xi \mathcal{O}_{W'_l, \xi'} = \mathfrak m_{\xi'}$
car l'image inverse de $Z_{i, l, j}$ est $Z'_{l, j}$ ! Ainsi, l'égalité de
$$
n_{l, j} = \text{ord}_{Z'_{l, j}}(f'_l) =
\text{ord}_{\mathcal{O}_{W'_l, \xi'}}(f'_l)
\quad\text{et}\quad
n_{i, l, j} = \text{ord}_{Z_{i, l, j}}(f_{i, l}) =
\text{ord}_{\mathcal{O}_{W_{i, l}, \xi}}(f_{i, l})
$$
résulte de Algèbre, Lemme \ref{algebra-lemma-pullback-module},
et de la construction de $\text{ord}$ dans
Algèbre, section \ref{algebra-section-orders-of-vanishing}.
\end{proof}












\section{Annexe A : autre approche du lemme clé}
\label{section-appendix-A}

\noindent
Dans cette annexe, nous définissons d'abord les déterminants $\det_\kappa(M)$
des modules de longueur finie $M$ sur des anneaux locaux $(R, \mathfrak m, \kappa)$,
voir la Sous-section \ref{subsection-determinants-finite-length}.
Le déterminant $\det_\kappa(M)$ est de dimension $1$ comme $\kappa$-espace vectoriel.
Nous l'utilisons dans la Sous-section \ref{subsection-periodic-complexes-determinants}
pour définir le déterminant $\det_\kappa(M, \varphi, \psi) \in \kappa^*$
d'un complexe $(2, 1)$-périodique exact $(M, \varphi, \psi)$
avec $M$ de longueur finie. Dans la Sous-section \ref{subsection-symbols},
nous utilisons ces déterminants pour construire un symbole modéré
$d_R(a, b) = \det_\kappa(R/ab, a, b)$ pour un couple de non-diviseurs de zéro
$a, b \in R$ lorsque $R$ est noethérien de dimension $1$.
Bien qu'il ne fasse aucun doute que
$$
d_R(a, b) = \partial_R(a, b)
$$
où $\partial_R$ est comme dans la section \ref{section-tame-symbol},
nous n'avons pas (encore) ajouté la vérification. L'avantage du
symbole modéré construit dans cette annexe est qu'il s'étend
(par exemple) aux couples d'endomorphismes injectifs $\varphi, \psi$
d'un $R$-module de type fini $M$ de dimension $1$ tels que
$\varphi(\psi(M)) = \psi(\varphi(M))$. Dans la
Sous-section \ref{subsection-length-determinant},
nous relions quotients de Herbrand et déterminants.
Une version facile à énoncer du résultat principal
(Proposition \ref{proposition-length-determinant-periodic-complex})
est la formule
$$
-e_R(M, \varphi, \psi) =
\text{ord}_R(\det\nolimits_K(M_K, \varphi, \psi))
$$
lorsque $(M, \varphi, \psi)$ est un complexe $(2, 1)$-périodique
dont le quotient de Herbrand $e_R$ (Définition \ref{definition-periodic-length})
est défini
sur un anneau local noethérien intègre de dimension $1$, noté $R$, de corps des fractions $K$.
Nous utilisons cette proposition pour donner une autre démonstration du lemme clé
(Lemme \ref{lemma-milnor-gersten-low-degree})
pour le symbole modéré construit dans cette annexe, voir le
Lemme \ref{lemma-secondary-ramification}.


\subsection{Déterminants de modules de longueur finie}
\label{subsection-determinants-finite-length}

\noindent
Le contenu de cette section est apparenté à celui de
l'article \cite{determinant} et à celui de la
thèse \cite{Joe}.

\medskip\noindent
Soit $(R, \mathfrak m, \kappa)$ un anneau local. Soit
$\varphi : M \to M$ un endomorphisme $R$-linéaire d'un
$R$-module de longueur finie $M$. Dans Compléments sur l'algèbre, section
\ref{more-algebra-section-determinants-finite-length},
nous avons déjà défini le déterminant $\det_\kappa(\varphi)$
(ainsi que la trace et le polynôme caractéristique)
de $\varphi$ relativement à $\kappa$. Dans cette section, nous allons
construire un espace vectoriel canonique de dimension $1$ sur $\kappa$
$\det_\kappa(M)$ tel que
$\det_\kappa(\varphi : M \to M) : \det_\kappa(M) \to \det_\kappa(M)$
est égal à la multiplication par $\det_\kappa(\varphi)$.
Si $M$ est annulé par $\mathfrak m$, alors $M$ peut être considéré
comme un $\kappa$-espace vectoriel de dimension finie et nous avons alors
$\det_\kappa(M) = \wedge^n_\kappa(M)$, où $n = \dim_\kappa(M)$.
Notre construction généralisera ceci à tous les modules de longueur finie
sur $R$ et, si $R$ contient son corps résiduel, le déterminant
$\det_\kappa(M)$ sera donné par le déterminant usuel en un sens
approprié, voir la Remarque \ref{remark-explain-determinant}.

\begin{definition}
\label{definition-determinant}
Soit $R$ un anneau local d'idéal maximal $\mathfrak m$ et de
corps résiduel $\kappa$. Soit $M$ un $R$-module de longueur finie.
Posons $l = \text{longueur}_R(M)$.
\begin{enumerate}
\item Étant donnés des éléments $x_1, \ldots, x_r \in M$, nous notons
$\langle x_1, \ldots, x_r \rangle = Rx_1 + \ldots + Rx_r$ le
sous-$R$-module de $M$ engendré par $x_1, \ldots, x_r$.
\item Nous dirons qu'un $l$-uplet d'éléments
$(e_1, \ldots, e_l)$ de $M$ est {\it admissible} si
$\mathfrak m e_i \subset \langle e_1, \ldots, e_{i - 1} \rangle$
pour $i = 1, \ldots, l$.
\item Un {\it symbole} $[e_1, \ldots, e_l]$ signifiera que
$(e_1, \ldots, e_l)$ est un $l$-uplet admissible.
\item Une {\it relation admissible} entre symboles est l'une des suivantes :
\begin{enumerate}
\item si $(e_1, \ldots, e_l)$ est une suite admissible et si,
pour un certain $1 \leq a \leq l$, nous avons
$e_a \in \langle e_1, \ldots, e_{a - 1}\rangle$, alors
$[e_1, \ldots, e_l] = 0$,
\item si $(e_1, \ldots, e_l)$ est une suite admissible et si,
pour un certain $1 \leq a \leq l$, nous avons $e_a = \lambda e'_a + x$
avec $\lambda \in R^*$, et
$x \in \langle e_1, \ldots, e_{a - 1}\rangle$, alors
$$
[e_1, \ldots, e_l] =
\overline{\lambda} [e_1, \ldots, e_{a - 1}, e'_a, e_{a + 1}, \ldots, e_l]
$$
où $\overline{\lambda} \in \kappa^*$ est l'image de $\lambda$ dans
le corps résiduel, et
\item si $(e_1, \ldots, e_l)$ est une suite admissible et si
$\mathfrak m e_a \subset \langle e_1, \ldots, e_{a - 2}\rangle$, alors
$$
[e_1, \ldots, e_l] =
- [e_1, \ldots, e_{a - 2}, e_a, e_{a - 1}, e_{a + 1}, \ldots, e_l].
$$
\end{enumerate}
\item
Nous définissons le {\it déterminant du $R$-module de longueur finie $M$} par
$$
\det\nolimits_\kappa(M) =
\left\{
\frac{\kappa\text{-espace vectoriel engendré par les symboles}}
{\kappa\text{-combinaisons linéaires de relations admissibles}}
\right\}
$$
\end{enumerate}
\end{definition}

\noindent
Soulignons que l'on a toujours $l = \text{longueur}_R(M)$. Soulignons aussi qu'il
n'en résulte pas que le symbole $[e_1, \ldots, e_l]$ soit
additif en ses entrées (ce ne sera typiquement pas le cas).
Avant de pouvoir montrer que le déterminant $\det_\kappa(M)$ est effectivement
de dimension $1$, nous devons montrer qu'il est de dimension au plus $1$.

\begin{lemma}
\label{lemma-dimension-at-most-one}
Avec les notations ci-dessus, nous avons $\dim_\kappa(\det_\kappa(M)) \leq 1$.
\end{lemma}

\begin{proof}
Fixons une suite admissible $(f_1, \ldots, f_l)$ de $M$ telle que
$$
\text{longueur}_R(\langle f_1, \ldots, f_i\rangle) = i
$$
pour $i = 1, \ldots, l$. Une telle suite admissible existe précisément parce que
$M$ est de longueur $l$. Nous allons montrer que tout élément de
$\det_\kappa(M)$ est un multiple par un élément de $\kappa$ du symbole
$[f_1, \ldots, f_l]$. Ceci prouvera le lemme.

\medskip\noindent
Soit $(e_1, \ldots, e_l)$ une suite admissible de $M$.
Il suffit de montrer que $[e_1, \ldots, e_l]$ est un multiple
de $[f_1, \ldots, f_l]$. Supposons d'abord que
$\langle e_1, \ldots, e_l\rangle \not = M$. Il existe alors
un $i \in [1, \ldots, l]$ tel que
$e_i \in \langle e_1, \ldots, e_{i - 1}\rangle$. Il résulte immédiatement
de la première relation admissible que
$[e_1, \ldots, e_n] = 0$ dans $\det_\kappa(M)$.
Nous pouvons donc supposer que $\langle e_1, \ldots, e_l\rangle = M$.
En particulier, il existe un plus petit indice $i \in \{1, \ldots, l\}$
tel que $f_1 \in \langle e_1, \ldots, e_i\rangle$. Ceci signifie
que $e_i = \lambda f_1 + x$, avec
$x \in \langle e_1, \ldots, e_{i - 1}\rangle$ et $\lambda \in R^*$.
Par la deuxième relation admissible, ceci signifie que
$[e_1, \ldots, e_l] =
\overline{\lambda}[e_1, \ldots, e_{i - 1}, f_1, e_{i + 1}, \ldots, e_l]$.
Notons que $\mathfrak m f_1 = 0$. Ainsi, en appliquant la troisième
relation admissible $i - 1$ fois, nous voyons que
$$
[e_1, \ldots, e_l] =
(-1)^{i - 1}\overline{\lambda}
[f_1, e_1, \ldots, e_{i - 1}, e_{i + 1}, \ldots, e_l].
$$
Notons que nous avons aussi
$ \langle f_1, e_1, \ldots, e_{i - 1}, e_{i + 1}, \ldots, e_l\rangle = M$.
Supposons par récurrence que nous ayons prouvé que notre symbole
initial est égal au produit d'un scalaire par
$$
[f_1, \ldots, f_j, e_{j + 1}, \ldots, e_l]
$$
pour une certaine suite admissible $(f_1, \ldots, f_j, e_{j + 1}, \ldots, e_l)$
dont les éléments engendrent $M$, c'est-à-dire \ avec
$\langle f_1, \ldots, f_j, e_{j + 1}, \ldots, e_l\rangle = M$.
Nous trouvons alors le plus petit $i$ tel que
$f_{j + 1} \in \langle f_1, \ldots, f_j, e_{j + 1}, \ldots, e_i\rangle$
et nous procédons comme ci-dessus pour voir que
$$
[f_1, \ldots, f_j, e_{j + 1}, \ldots, e_l]
=
(\text{scalaire}) [f_1, \ldots, f_j, f_{j + 1}, e_{j + 1},
\ldots, \hat{e_i}, \ldots, e_l]
$$
En continuant ainsi, nous obtenons le résultat désiré.
\end{proof}

\noindent
Avant de montrer que $\det_\kappa(M)$ est toujours de dimension $1$,
montrons qu'il coïncide avec la puissance extérieure maximale usuelle dans
le cas où le module est un espace vectoriel sur $\kappa$.

\begin{lemma}
\label{lemma-compare-det}
Soit $R$ un anneau local d'idéal maximal $\mathfrak m$ et de
corps résiduel $\kappa$. Soit $M$ un $R$-module de longueur finie
annulé par $\mathfrak m$. Soit $l = \dim_\kappa(M)$.
Alors l'application
$$
\det\nolimits_\kappa(M) \longrightarrow \wedge^l_\kappa(M),
\quad
[e_1, \ldots, e_l] \longmapsto e_1 \wedge \ldots \wedge e_l
$$
est un isomorphisme.
\end{lemma}

\begin{proof}
Il est clair que la règle décrite dans le lemme donne une application $\kappa$-linéaire
puisque toutes les relations admissibles sont satisfaites par les
symboles usuels $e_1 \wedge \ldots \wedge e_l$. Cette application est aussi clairement
surjective. Puisque, par le Lemme \ref{lemma-dimension-at-most-one}, le membre de gauche
est de dimension au plus un,
nous voyons que l'application est un isomorphisme.
\end{proof}

\begin{lemma}
\label{lemma-determinant-dimension-one}
Soit $R$ un anneau local d'idéal maximal $\mathfrak m$ et de
corps résiduel $\kappa$. Soit $M$ un $R$-module de longueur finie.
Le déterminant $\det_\kappa(M)$ défini ci-dessus est un $\kappa$-espace vectoriel
de dimension $1$. Il est engendré par le symbole
$[f_1, \ldots, f_l]$ pour toute suite admissible telle
que $\langle f_1, \ldots f_l \rangle = M$.
\end{lemma}

\begin{proof}
Nous savons que $\det_\kappa(M)$ est de dimension au plus $1$, et en fait qu'il
est engendré par $[f_1, \ldots, f_l]$, par le
Lemme \ref{lemma-dimension-at-most-one} et sa démonstration.
Nous allons montrer par récurrence sur $l = \text{longueur}(M)$
qu'il est non nul. Pour $l = 1$, ceci résulte du Lemme \ref{lemma-compare-det}.
Choisissons un élément non nul $f \in M$
tel que $\mathfrak m f = 0$. Posons $\overline{M} = M /\langle f \rangle$,
et notons l'application quotient $x \mapsto \overline{x}$.
Nous allons définir une application surjective
$$
\psi : \det\nolimits_k(M) \to \det\nolimits_\kappa(\overline{M})
$$
qui prouvera le lemme puisque, par récurrence, le déterminant de
$\overline{M}$ est non nul.

\medskip\noindent
Nous définissons $\psi$ sur les symboles comme suit.
Soit $(e_1, \ldots, e_l)$ une suite admissible.
Si $f \not \in \langle e_1, \ldots, e_l \rangle$, alors
nous posons simplement $\psi([e_1, \ldots, e_l]) = 0$.
Si $f \in \langle e_1, \ldots, e_l \rangle$, alors nous choisissons
un $i$ minimal tel que $f \in \langle e_1, \ldots, e_i \rangle$.
Nous pouvons écrire $e_i = \lambda f + x$ pour une unité $\lambda \in R$
et $x \in \langle e_1, \ldots, e_{i - 1} \rangle$.
Dans ce cas, nous posons
$$
\psi([e_1, \ldots, e_l]) =
(-1)^i
\overline{\lambda}[\overline{e}_1, \ldots,
\overline{e}_{i - 1},
\overline{e}_{i + 1}, \ldots, \overline{e}_l].
$$
Notons que nous avons bien
$(\overline{e}_1, \ldots,
\overline{e}_{i - 1},
\overline{e}_{i + 1}, \ldots, \overline{e}_l)$
est une suite admissible dans $\overline{M}$, de sorte que ceci a un sens.
Montrons qu'étendre cette règle $\kappa$-linéairement aux
combinaisons linéaires de symboles conduit bien à une application sur les
déterminants. Pour cela, nous devons montrer que les relations
admissibles sont envoyées sur zéro.

\medskip\noindent
Relations de type (a). Supposons que $(e_1, \ldots, e_l)$ soit une
suite admissible et que, pour un certain $1 \leq a \leq l$, nous ayons
$e_a \in \langle e_1, \ldots, e_{a - 1}\rangle$.
Supposons que $f \in \langle e_1, \ldots, e_i\rangle$, avec $i$ minimal.
Alors $i \not = a$ et
$\overline{e}_a \in \langle \overline{e}_1, \ldots,
\hat{\overline{e}_i}, \ldots, \overline{e}_{a - 1}\rangle$ si $i < a$
ou
$\overline{e}_a \in \langle \overline{e}_1, \ldots,
\overline{e}_{a - 1}\rangle$ si $i > a$.
Ainsi, la même relation admissible pour $\det_\kappa(\overline{M})$ force
le symbole $[\overline{e}_1, \ldots,
\overline{e}_{i - 1},
\overline{e}_{i + 1}, \ldots, \overline{e}_l]$
à être nul comme voulu.

\medskip\noindent
Relations de type (b). Supposons que $(e_1, \ldots, e_l)$ soit une
suite admissible et que, pour un certain $1 \leq a \leq l$, nous ayons
$e_a = \lambda e'_a + x$, avec $\lambda \in R^*$, et
$x \in \langle e_1, \ldots, e_{a - 1}\rangle$.
Supposons que $f \in \langle e_1, \ldots, e_i\rangle$, avec $i$ minimal.
Écrivons $e_i = \mu f + y$, avec $y \in \langle e_1, \ldots, e_{i - 1}\rangle$.
Si $i < a$, alors l'égalité désirée est
$$
(-1)^i
\overline{\lambda}
[\overline{e}_1,
\ldots,
\overline{e}_{i - 1},
\overline{e}_{i + 1},
\ldots,
\overline{e}_l]
=
(-1)^i
\overline{\lambda}
[\overline{e}_1,
\ldots,
\overline{e}_{i - 1},
\overline{e}_{i + 1},
\ldots,
\overline{e}_{a - 1},
\overline{e}'_a,
\overline{e}_{a + 1},
\ldots,
\overline{e}_l]
$$
ce qui résulte de $\overline{e}_a = \lambda \overline{e}'_a + \overline{x}$
et de la relation admissible correspondante pour $\det_\kappa(\overline{M})$.
Si $i > a$, alors l'égalité désirée est
$$
(-1)^i
\overline{\lambda}
[\overline{e}_1,
\ldots,
\overline{e}_{i - 1},
\overline{e}_{i + 1},
\ldots,
\overline{e}_l]
=
(-1)^i
\overline{\lambda}
[\overline{e}_1,
\ldots,
\overline{e}_{a - 1},
\overline{e}'_a,
\overline{e}_{a + 1},
\ldots,
\overline{e}_{i - 1},
\overline{e}_{i + 1},
\ldots,
\overline{e}_l]
$$
ce qui résulte de $\overline{e}_a = \lambda \overline{e}'_a + \overline{x}$
et de la relation admissible correspondante pour $\det_\kappa(\overline{M})$.
Le cas intéressant est celui où $i = a$. Dans ce cas, nous avons
$e_a = \lambda e'_a + x = \mu f + y$. Donc aussi
$e'_a = \lambda^{-1}(\mu f + y - x)$. Ainsi, nous voyons que
$$
\psi([e_1, \ldots, e_l])
= (-1)^i \overline{\mu}
[\overline{e}_1, \ldots,
\overline{e}_{i - 1},
\overline{e}_{i + 1}, \ldots, \overline{e}_l]
=
\psi(
\overline{\lambda}
[e_1, \ldots, e_{a - 1}, e'_a, e_{a + 1}, \ldots, e_l]
)
$$
comme voulu.

\medskip\noindent
Relations de type (c). Supposons que $(e_1, \ldots, e_l)$
soit une suite admissible et que
$\mathfrak m e_a \subset \langle e_1, \ldots, e_{a - 2}\rangle$.
Supposons que $f \in \langle e_1, \ldots, e_i\rangle$, avec $i$ minimal.
Écrivons $e_i = \lambda f + x$, avec $x \in \langle e_1, \ldots, e_{i - 1}\rangle$.
Nous distinguons $4$ cas :

\medskip\noindent
Cas 1 : $i < a - 1$. L'égalité désirée est
\begin{align*}
& (-1)^i
\overline{\lambda}
[\overline{e}_1,
\ldots,
\overline{e}_{i - 1},
\overline{e}_{i + 1},
\ldots,
\overline{e}_l] \\
& =
(-1)^{i + 1}
\overline{\lambda}
[\overline{e}_1,
\ldots,
\overline{e}_{i - 1},
\overline{e}_{i + 1},
\ldots,
\overline{e}_{a - 2},
\overline{e}_a,
\overline{e}_{a - 1},
\overline{e}_{a + 1},
\ldots,
\overline{e}_l]
\end{align*}
ce qui résulte de la relation admissible de type (c) pour
$\det_\kappa(\overline{M})$.

\medskip\noindent
Cas 2 : $i > a$. L'égalité désirée est
\begin{align*}
& (-1)^i
\overline{\lambda}
[\overline{e}_1,
\ldots,
\overline{e}_{i - 1},
\overline{e}_{i + 1},
\ldots,
\overline{e}_l] \\
& =
(-1)^{i + 1}
\overline{\lambda}
[\overline{e}_1,
\ldots,
\overline{e}_{a - 2},
\overline{e}_a,
\overline{e}_{a - 1},
\overline{e}_{a + 1},
\ldots,
\overline{e}_{i - 1},
\overline{e}_{i + 1},
\ldots,
\overline{e}_l]
\end{align*}
ce qui résulte de la relation admissible de type (c) pour
$\det_\kappa(\overline{M})$.

\medskip\noindent
Cas 3 : $i = a$. Nous écrivons $e_a = \lambda f + \mu e_{a - 1} + y$,
avec $y \in \langle e_1, \ldots, e_{a - 2}\rangle$. Alors
$$
\psi([e_1, \ldots, e_l]) =
(-1)^a
\overline{\lambda}
[\overline{e}_1,
\ldots,
\overline{e}_{a - 1},
\overline{e}_{a + 1},
\ldots,
\overline{e}_l]
$$
par définition. Si $\overline{\mu}$ est non nul, alors nous avons
$e_{a - 1} = - \mu^{-1} \lambda f + \mu^{-1}e_a - \mu^{-1} y$
et nous obtenons
$$
\psi(-[e_1, \ldots, e_{a - 2}, e_a, e_{a - 1}, e_{a + 1}, \ldots, e_l]) =
(-1)^a
\overline{\mu^{-1}\lambda}
[\overline{e}_1,
\ldots,
\overline{e}_{a - 2},
\overline{e}_a,
\overline{e}_{a + 1},
\ldots,
\overline{e}_l]
$$
par définition. Puisque, dans $\overline{M}$, nous avons
$\overline{e}_a = \mu \overline{e}_{a - 1} + \overline{y}$, nous voyons
que les deux résultats sont égaux par la relation (a) pour $\det_\kappa(\overline{M})$.
Si, en revanche, $\overline{\mu}$ est nul, alors nous pouvons écrire
$e_a = \lambda f + y$, avec $y \in \langle e_1, \ldots, e_{a - 2}\rangle$,
et nous avons
$$
\psi(-[e_1, \ldots, e_{a - 2}, e_a, e_{a - 1}, e_{a + 1}, \ldots, e_l]) =
(-1)^a
\overline{\lambda}
[\overline{e}_1,
\ldots,
\overline{e}_{a - 1},
\overline{e}_{a + 1},
\ldots,
\overline{e}_l]
$$
ce qui est égal à $\psi([e_1, \ldots, e_l])$.

\medskip\noindent
Cas 4 : $i = a - 1$. Ici, nous avons
$$
\psi([e_1, \ldots, e_l]) =
(-1)^{a - 1}
\overline{\lambda}
[\overline{e}_1,
\ldots,
\overline{e}_{a - 2},
\overline{e}_a,
\ldots,
\overline{e}_l]
$$
par définition. Si $f \not \in \langle e_1, \ldots, e_{a - 2}, e_a \rangle$,
alors
$$
\psi(-[e_1, \ldots, e_{a - 2}, e_a, e_{a - 1}, e_{a + 1}, \ldots, e_l]) =
(-1)^{a + 1}\overline{\lambda}
[\overline{e}_1,
\ldots,
\overline{e}_{a - 2},
\overline{e}_a,
\ldots,
\overline{e}_l]
$$
Puisque $(-1)^{a - 1} = (-1)^{a + 1}$, les deux expressions sont les mêmes.
Enfin, supposons $f \in \langle e_1, \ldots, e_{a - 2}, e_a \rangle$.
Dans ce cas, nous voyons que $e_{a - 1} = \lambda f + x$, avec
$x \in \langle e_1, \ldots, e_{a - 2}\rangle$, et
$e_a = \mu f + y$, avec $y \in \langle e_1, \ldots, e_{a - 2}\rangle$,
pour des unités $\lambda, \mu \in R$.
Nous concluons que nous avons à la fois
$e_a \in \langle e_1, \ldots, e_{a - 1} \rangle$ et
$e_{a - 1} \in \langle e_1, \ldots, e_{a - 2}, e_a\rangle$.
Dans ce cas, une relation de type (a) s'applique à la fois à
$[e_1, \ldots, e_l]$ et
$[e_1, \ldots, e_{a - 2}, e_a, e_{a - 1}, e_{a + 1}, \ldots, e_l]$
et la compatibilité de $\psi$ avec celles-ci montrée ci-dessus fait voir que les deux
$$
\psi([e_1, \ldots, e_l])
\quad\text{et}\quad
\psi([e_1, \ldots, e_{a - 2}, e_a, e_{a - 1}, e_{a + 1}, \ldots, e_l])
$$
sont nuls, comme voulu.

\medskip\noindent
À ce stade, nous avons montré que $\psi$ est bien définie, et il ne reste
qu'à montrer qu'elle est surjective. Pour le voir, soit
$(\overline{f}_2, \ldots, \overline{f}_l)$ une suite admissible
dans $\overline{M}$. Nous pouvons choisir des relèvements $f_2, \ldots, f_l \in M$, et
alors $(f, f_2, \ldots, f_l)$ est une suite admissible dans $M$.
Puisque $\psi([f, f_2, \ldots, f_l]) = [f_2, \ldots, f_l]$, nous avons gagné.
\end{proof}

\noindent
Soit $R$ un anneau local d'idéal maximal $\mathfrak m$ et de
corps résiduel $\kappa$. Notons que si $\varphi : M \to N$ est un
isomorphisme de $R$-modules de longueur finie, alors nous obtenons un
isomorphisme
$$
\det\nolimits_\kappa(\varphi) :
\det\nolimits_\kappa(M)
\to
\det\nolimits_\kappa(N)
$$
simplement par la règle
$$
\det\nolimits_\kappa(\varphi)([e_1, \ldots, e_l])
=
[\varphi(e_1), \ldots, \varphi(e_l)]
$$
pour tout symbole $[e_1, \ldots, e_l]$ de $M$.
Ainsi, nous voyons que $\det\nolimits_\kappa$ est un foncteur
\begin{equation}
\label{equation-functor}
\left\{
\begin{matrix}
\text{les }R\text{-modules de longueur finie}\\
\text{avec isomorphismes}
\end{matrix}
\right\}
\longrightarrow
\left\{
\begin{matrix}
1\text{ est la dimension des }\kappa\text{-espaces vectoriels}\\
\text{avec isomorphismes}
\end{matrix}
\right\}
\end{equation}
Ceci est typique d'un « foncteur déterminant »
(voir \cite{Knudsen}), tout comme la propriété d'additivité
suivante.

\begin{lemma}
\label{lemma-det-exact-sequences}
Soit $(R, \mathfrak m, \kappa)$ un anneau local.
Pour toute suite exacte courte
$$
0 \to K \to L \to M \to 0
$$
de $R$-modules de longueur finie, il existe un isomorphisme canonique
$$
\gamma_{K \to L \to M} :
\det\nolimits_\kappa(K) \otimes_\kappa \det\nolimits_\kappa(M)
\longrightarrow
\det\nolimits_\kappa(L)
$$
défini sur les symboles non nuls par la règle
$$
[e_1, \ldots, e_k]
\otimes
[\overline{f}_1, \ldots, \overline{f}_m]
\longrightarrow
[e_1, \ldots, e_k, f_1, \ldots, f_m]
$$
ayant les propriétés suivantes :
\begin{enumerate}
\item Pour tout isomorphisme de suites exactes courtes, c'est-à-dire pour
tout diagramme commutatif
$$
\xymatrix{
0 \ar[r] &
K \ar[r] \ar[d]^u &
L \ar[r] \ar[d]^v &
M \ar[r] \ar[d]^w &
0 \\
0 \ar[r] &
K' \ar[r] &
L' \ar[r] &
M' \ar[r] &
0
}
$$
dont les lignes sont exactes courtes et où $u, v, w$ sont des isomorphismes, nous avons
$$
\gamma_{K' \to L' \to M'} \circ
(\det\nolimits_\kappa(u) \otimes \det\nolimits_\kappa(w))
=
\det\nolimits_\kappa(v) \circ
\gamma_{K \to L \to M},
$$
\item pour tout carré commutatif de $R$-modules de longueur finie
dont les lignes et les colonnes sont exactes
$$
\xymatrix{
& 0 \ar[d] & 0 \ar[d] & 0 \ar[d] & \\
0 \ar[r] & A \ar[r] \ar[d] & B \ar[r] \ar[d] & C \ar[r] \ar[d] & 0 \\
0 \ar[r] & D \ar[r] \ar[d] & E \ar[r] \ar[d] & F \ar[r] \ar[d] & 0 \\
0 \ar[r] & G \ar[r] \ar[d] & H \ar[r] \ar[d] & I \ar[r] \ar[d] & 0 \\
& 0  & 0  & 0  &
}
$$
le diagramme suivant est commutatif
$$
\xymatrix{
\det\nolimits_\kappa(A) \otimes
\det\nolimits_\kappa(C) \otimes
\det\nolimits_\kappa(G) \otimes
\det\nolimits_\kappa(I)
\ar[dd]_{\epsilon}
\ar[rrr]_-{\gamma_{A \to B \to C} \otimes \gamma_{G \to H \to I}}
& & &
\det\nolimits_\kappa(B) \otimes
\det\nolimits_\kappa(H)
\ar[d]^{\gamma_{B \to E \to H}}
\\
& & & \det\nolimits_\kappa(E)
\\
\det\nolimits_\kappa(A) \otimes
\det\nolimits_\kappa(G) \otimes
\det\nolimits_\kappa(C) \otimes
\det\nolimits_\kappa(I)
\ar[rrr]^-{\gamma_{A \to D \to G} \otimes \gamma_{C \to F \to I}}
& & &
\det\nolimits_\kappa(D) \otimes
\det\nolimits_\kappa(F)
\ar[u]_{\gamma_{D \to E \to F}}
}
$$
où $\epsilon$ est la permutation des facteurs du produit tensoriel
multipliée par $(-1)^{cg}$, avec $c = \text{longueur}_R(C)$ et $g = \text{longueur}_R(G)$,
et
\item l'application $\gamma_{K \to L \to M}$ coïncide avec l'isomorphisme usuel
si $0 \to K \to L \to M \to 0$ est en fait une suite exacte courte
de $\kappa$-espaces vectoriels.
\end{enumerate}
\end{lemma}

\begin{proof}
L'intérêt de prendre des symboles non nuls dans la description explicite
de l'application $\gamma_{K \to L \to M}$ tient simplement à ce que, si $(e_1, \ldots, e_l)$
est une suite admissible dans $K$, et si
$(\overline{f}_1, \ldots, \overline{f}_m)$ est une suite admissible dans
$M$, il n'est pas garanti que $(e_1, \ldots, e_l, f_1, \ldots, f_m)$
soit une suite admissible dans $L$ (où, bien sûr, $f_i \in L$ désigne
un relèvement de $\overline{f}_i$). Toutefois, si le symbole
$[e_1, \ldots, e_l]$ est non nul dans $\det_\kappa(K)$, alors
nécessairement $K = \langle e_1, \ldots, e_k\rangle$ (voir la
démonstration du Lemme \ref{lemma-dimension-at-most-one}), et
dans ce cas $(e_1, \ldots, e_k, f_1, \ldots, f_m)$ est bien
une suite admissible.
De plus, par les relations admissibles de type (b) pour $\det_\kappa(L)$,
nous voyons que la valeur de $[e_1, \ldots, e_k, f_1, \ldots, f_m]$ dans
$\det_\kappa(L)$ est indépendante du choix des relèvements
$f_i$ dans ce cas également. Compte tenu de cette remarque, il est clair
qu'une relation admissible entre $e_1, \ldots, e_k$ dans $K$
se traduit par une relation admissible entre
$e_1, \ldots, e_k, f_1, \ldots, f_m$ dans $L$, et
de même pour une relation admissible entre les
$\overline{f}_1, \ldots, \overline{f}_m$.
Ainsi, $\gamma$ définit une application linéaire d'espaces vectoriels comme affirmé dans le lemme.

\medskip\noindent
Par le Lemme \ref{lemma-determinant-dimension-one}, nous savons que
$\det_\kappa(L)$ est engendré par n'importe quel
symbole $[x_1, \ldots, x_{k + m}]$ tel que
$(x_1, \ldots, x_{k + m})$ soit une suite admissible
avec $L = \langle x_1, \ldots, x_{k + m}\rangle$. Il est donc
clair que l'application $\gamma_{K \to L \to M}$ est surjective et
qu'elle est donc un isomorphisme.

\medskip\noindent
La propriété (1) vaut car
\begin{eqnarray*}
& & \det\nolimits_\kappa(v)([e_1, \ldots, e_k, f_1, \ldots, f_m]) \\
& = &
[v(e_1), \ldots, v(e_k), v(f_1), \ldots, v(f_m)] \\
& = &
\gamma_{K' \to L' \to M'}([u(e_1), \ldots, u(e_k)]
\otimes [w(f_1), \ldots, w(f_m)]).
\end{eqnarray*}
La propriété (2) signifie qu'étant donnés un symbole
$[\alpha_1, \ldots, \alpha_a]$ engendrant $\det_\kappa(A)$,
un symbole $[\gamma_1, \ldots, \gamma_c]$ engendrant $\det_\kappa(C)$,
un symbole $[\zeta_1, \ldots, \zeta_g]$ engendrant $\det_\kappa(G)$, et
un symbole $[\iota_1, \ldots, \iota_i]$ engendrant $\det_\kappa(I)$,
nous avons
\begin{eqnarray*}
& & [\alpha_1, \ldots, \alpha_a, \tilde\gamma_1, \ldots, \tilde\gamma_c,
\tilde\zeta_1, \ldots, \tilde\zeta_g, \tilde\iota_1, \ldots, \tilde\iota_i] \\
& = &
(-1)^{cg} [\alpha_1, \ldots, \alpha_a, \tilde\zeta_1, \ldots, \tilde\zeta_g,
\tilde\gamma_1, \ldots, \tilde\gamma_c, \tilde\iota_1, \ldots, \tilde\iota_i]
\end{eqnarray*}
(pour des relèvements convenables $\tilde{x}$ dans $E$) dans $\det_\kappa(E)$.
Ceci vaut car nous pouvons utiliser les relations admissibles de type (c)
$cg$ fois dans l'ordre suivant : faire passer
$\tilde\zeta_1$ devant les éléments
$\tilde\gamma_c, \ldots, \tilde\gamma_1$
(ce qui est permis puisque $\mathfrak m\tilde\zeta_1 \subset A$),
puis faire passer $\tilde\zeta_2$ devant les éléments
$\tilde\gamma_c, \ldots, \tilde\gamma_1$
(ce qui est permis puisque $\mathfrak m\tilde\zeta_2 \subset A + R\tilde\zeta_1$),
et ainsi de suite.

\medskip\noindent
La partie (3) du lemme est évidente.
Ceci termine la démonstration.
\end{proof}

\noindent
Nous pouvons utiliser les applications $\gamma$ du lemme pour définir des applications
$\gamma$ plus générales comme suit. Supposons que $(R, \mathfrak m, \kappa)$ soit un
anneau local. Soit $M$ un $R$-module de longueur finie et supposons donnée
une filtration finie (voir
Homologie, Définition \ref{homology-definition-filtered})
$$
0 = F^m \subset F^{m - 1} \subset \ldots \subset F^{n + 1} \subset F^n = M
$$
alors il existe un isomorphisme canonique bien défini
$$
\gamma_{(M, F)} :
\det\nolimits_\kappa(F^{m - 1}/F^m) \otimes_\kappa \ldots \otimes_k 
\det\nolimits_\kappa(F^n/F^{n + 1})
\longrightarrow
\det\nolimits_\kappa(M)
$$
Pour le construire, nous utilisons les isomorphismes du Lemme \ref{lemma-det-exact-sequences}
provenant des suites exactes courtes
$0 \to F^{i - 1}/F^i \to M/F^i \to M/F^{i - 1} \to 0$.
La partie (2) du Lemme \ref{lemma-det-exact-sequences}, avec $G = 0$, montre
que nous obtenons le même isomorphisme si nous utilisons les suites exactes courtes
$0 \to F^i \to F^{i - 1} \to F^{i - 1}/F^i \to 0$.

\medskip\noindent
Voici un autre résultat typique sur les foncteurs déterminant.
Il n'est pas difficile à montrer. La partie délicate consiste habituellement à établir
l'existence d'un foncteur déterminant.

\begin{lemma}
\label{lemma-uniqueness-det}
Soit $(R, \mathfrak m, \kappa)$ un anneau local quelconque.
Le foncteur
$$
\det\nolimits_\kappa :
\left\{
\begin{matrix}
\text{les }R\text{-modules de longueur finie} \\
\text{avec isomorphismes}
\end{matrix}
\right\}
\longrightarrow
\left\{
\begin{matrix}
1\text{ est la dimension des }\kappa\text{-espaces vectoriels} \\
\text{avec isomorphismes}
\end{matrix}
\right\}
$$
muni des applications $\gamma_{K \to L \to M}$ est caractérisé par
les propriétés suivantes
\begin{enumerate}
\item sa restriction à la sous-catégorie des modules annulés
par $\mathfrak m$ est isomorphe au foncteur déterminant usuel
(voir le Lemme \ref{lemma-compare-det}), et
\item les points (1), (2) et (3) du Lemme \ref{lemma-det-exact-sequences}
sont vérifiés.
\end{enumerate}
\end{lemma}

\begin{proof}
Omis.
\end{proof}

\begin{lemma}
\label{lemma-determinant-quotient-ring}
Soit $(R', \mathfrak m') \to (R, \mathfrak m)$ un homomorphisme d'anneaux
locaux qui induit un isomorphisme sur les corps résiduels $\kappa$.
Alors, pour tout $R$-module de longueur finie, la restriction $M_{R'}$
est un $R'$-module de longueur finie et il existe un isomorphisme canonique
$$
\det\nolimits_{R, \kappa}(M)
\longrightarrow
\det\nolimits_{R', \kappa}(M_{R'})
$$
Cet isomorphisme est fonctoriel en $M$ et compatible avec les
isomorphismes $\gamma_{K \to L \to M}$ du Lemme \ref{lemma-det-exact-sequences}
définis pour $\det_{R, \kappa}$ et $\det_{R', \kappa}$.
\end{lemma}

\begin{proof}
Si la longueur de $M$ comme $R$-module est $l$, alors la longueur
de $M$ comme $R'$-module (c'est-à-dire de $M_{R'}$) est aussi $l$, voir
Algèbre, Lemme \ref{algebra-lemma-pushdown-module}.
Notons qu'une suite admissible $x_1, \ldots, x_l$ de $M$
sur $R$ est une suite admissible de $M$ sur $R'$, puisque $\mathfrak m'$
s'envoie dans $\mathfrak m$.
L'isomorphisme est obtenu en envoyant le symbole
$[x_1, \ldots, x_l] \in \det\nolimits_{R, \kappa}(M)$
sur le symbole correspondant
$[x_1, \ldots, x_l] \in \det\nolimits_{R', \kappa}(M)$.
Il est immédiat de vérifier que ceci est fonctoriel pour les
isomorphismes et compatible avec les isomorphismes
$\gamma$ du Lemme \ref{lemma-det-exact-sequences}.
\end{proof}

\begin{remark}
\label{remark-explain-determinant}
Soit $(R, \mathfrak m, \kappa)$ un anneau local et supposons soit que
la caractéristique de $\kappa$ est nulle, soit qu'elle vaut $p$ et que $p R = 0$.
Soient $M_1, \ldots, M_n$ des $R$-modules de longueur finie.
Nous montrerons ci-dessous qu'il existe un
idéal $I \subset \mathfrak m$ annulant $M_i$ pour $i = 1, \ldots, n$
et une section $\sigma : \kappa \to R/I$ de la surjection canonique
$R/I \to \kappa$. La restriction $M_{i, \kappa}$ de $M_i$ par $\sigma$
est un $\kappa$-espace vectoriel de dimension $l_i = \text{longueur}_R(M_i)$ et,
en utilisant le Lemme \ref{lemma-determinant-quotient-ring}, nous voyons que
$$
\det\nolimits_\kappa(M_i) = \wedge_\kappa^{l_i}(M_{i, \kappa})
$$
Ces isomorphismes sont compatibles avec les isomorphismes
$\gamma_{K \to M \to L}$ du Lemme \ref{lemma-det-exact-sequences}
pour les suites exactes courtes de $R$-modules de longueur finie annulés
par $I$. Nous en concluons que vérifier une propriété de
$\det_\kappa$ se réduit souvent à vérifier les propriétés correspondantes
du déterminant usuel sur la catégorie des espaces vectoriels de dimension
finie.

\medskip\noindent
Pour $I$, nous pouvons prendre l'annulateur
(Algèbre, Définition \ref{algebra-definition-annihilator})
du module $M = \bigoplus M_i$. Dans ce cas, nous voyons que
$R/I \subset \text{End}_R(M)$ est donc de longueur finie.
Ainsi, $R/I$ est un anneau local artinien de corps résiduel $\kappa$.
Puisqu'un anneau local artinien est complet, nous voyons que $R/I$
admet un anneau de coefficients par le théorème de structure de Cohen
(Algèbre, Théorème \ref{algebra-theorem-cohen-structure-theorem}),
qui est un corps par notre hypothèse sur $R$.
\end{remark}

\noindent
Voici un cas où nous pouvons calculer le déterminant d'une application linéaire.
En fait, il n'y a rien de mystérieux dans aucun cas ; voir
l'Exemple \ref{example-determinant-map} pour un exemple quelconque.

\begin{lemma}
\label{lemma-times-u-determinant}
Soit $R$ un anneau local de corps résiduel $\kappa$.
Soit $u \in R^*$ une unité.
Soit $M$ un module de longueur finie sur $R$.
Notons $u_M : M \to M$ l'application de multiplication par $u$.
Alors
$$
\det\nolimits_\kappa(u_M) :
\det\nolimits_\kappa(M)
\longrightarrow
\det\nolimits_\kappa(M)
$$
est la multiplication par $\overline{u}^l$, où $l = \text{longueur}_R(M)$
et $\overline{u} \in \kappa^*$ est l'image de $u$.
\end{lemma}

\begin{proof}
Notons $f_M \in \kappa^*$ l'élément tel que
$\det\nolimits_\kappa(u_M) = f_M \text{id}_{\det\nolimits_\kappa(M)}$.
Supposons que $0 \to K \to L \to M \to 0$ soit une suite exacte
courte de $R$-modules de type fini. Nous voyons alors que
$u_k$, $u_L$, $u_M$ donnent un isomorphisme de suites exactes courtes.
Ainsi, par le Lemme \ref{lemma-det-exact-sequences} (1), nous concluons que
$f_K f_M = f_L$.
Ceci signifie que, par récurrence sur la longueur, il suffit de prouver le
lemme dans le cas de longueur $1$, où il est trivial.
\end{proof}

\begin{example}
\label{example-determinant-map}
Considérons l'anneau local $R = \mathbf{Z}_p$.
Posons $M = \mathbf{Z}_p/(p^2) \oplus \mathbf{Z}_p/(p^3)$.
Soit $u : M \to M$ l'application donnée par la matrice
$$
u =
\left(
\begin{matrix}
a & b \\
pc & d
\end{matrix}
\right)
$$
où $a, b, c, d \in \mathbf{Z}_p$, et $a, d \in \mathbf{Z}_p^*$.
Dans ce cas, $\det_\kappa(u)$ est égal à la multiplication par
$a^2d^3 \bmod p \in \mathbf{F}_p^*$. Ceci se voit facilement
en considérant l'effet de $u$ sur le symbole
$[p^2e, pe, pf, e, f]$, où $e = (0 , 1) \in M$ et
$f = (1, 0) \in M$.
\end{example}





\subsection{Complexes périodiques et déterminants}
\label{subsection-periodic-complexes-determinants}

\noindent
Soit $R$ un anneau local de corps résiduel $\kappa$.
Soit $(M, \varphi, \psi)$ un complexe $(2, 1)$-périodique sur $R$.
Supposons que $M$ soit de longueur finie et que $(M, \varphi, \psi)$ soit
exact. Nous allons utiliser la construction du déterminant pour définir
un invariant de cette situation. Voir la
Sous-section \ref{subsection-determinants-finite-length}.
Abrégeons
$K_\varphi = \Ker(\varphi)$,
$I_\varphi = \Im(\varphi)$,
$K_\psi = \Ker(\psi)$, et
$I_\psi = \Im(\psi)$.
Les suites exactes courtes
$$
0 \to K_\varphi \to M \to I_\varphi \to 0, \quad
0 \to K_\psi \to M \to I_\psi \to 0
$$
donnent des isomorphismes
$$
\gamma_\varphi :
\det\nolimits_\kappa(K_\varphi)
\otimes
\det\nolimits_\kappa(I_\varphi)
\longrightarrow
\det\nolimits_\kappa(M), \quad
\gamma_\psi :
\det\nolimits_\kappa(K_\psi)
\otimes
\det\nolimits_\kappa(I_\psi)
\longrightarrow
\det\nolimits_\kappa(M),
$$
voir le lemme \ref{lemma-det-exact-sequences}.
D'autre part, l'exactitude du complexe donne les égalités
$K_\varphi = I_\psi$, et $K_\psi = I_\varphi$
et donc un isomorphisme
$$
\sigma :
\det\nolimits_\kappa(K_\varphi)
\otimes
\det\nolimits_\kappa(I_\varphi)
\longrightarrow
\det\nolimits_\kappa(K_\psi)
\otimes
\det\nolimits_\kappa(I_\psi)
$$
en permutant les facteurs. Avec ces notations, nous pouvons définir notre invariant.

\begin{definition}
\label{definition-periodic-determinant}
Soit $R$ un anneau local de corps résiduel $\kappa$.
Soit $(M, \varphi, \psi)$ un complexe $(2, 1)$-périodique sur $R$.
Supposons que $M$ soit de longueur finie et que $(M, \varphi, \psi)$ soit
exact. Le {\it déterminant de $(M, \varphi, \psi)$} est
l'élément
$$
\det\nolimits_\kappa(M, \varphi, \psi) \in \kappa^*
$$
tel que la composée
$$
\det\nolimits_\kappa(M)
\xrightarrow{\gamma_\psi \circ \sigma \circ \gamma_\varphi^{-1}}
\det\nolimits_\kappa(M)
$$
soit la multiplication par
$(-1)^{\text{longueur}_R(I_\varphi)\text{longueur}_R(I_\psi)}
\det\nolimits_\kappa(M, \varphi, \psi)$.
\end{definition}

\begin{remark}
\label{remark-more-elementary}
Voici une description plus concrète du déterminant
introduit ci-dessus. Soit $R$ un anneau local de corps résiduel $\kappa$.
Soit $(M, \varphi, \psi)$ un complexe $(2, 1)$-périodique sur $R$.
Supposons que $M$ soit de longueur finie et que $(M, \varphi, \psi)$ soit
exact. Abrégeons $I_\varphi = \Im(\varphi)$,
$I_\psi = \Im(\psi)$ comme ci-dessus.
Supposons que $\text{longueur}_R(I_\varphi) = a$ et que
$\text{longueur}_R(I_\psi) = b$, de sorte que $a + b = \text{longueur}_R(M)$
par exactitude. Choisissons des suites admissibles
$x_1, \ldots, x_a \in I_\varphi$ et $y_1, \ldots, y_b \in I_\psi$
telles que le symbole $[x_1, \ldots, x_a]$ engendre $\det_\kappa(I_\varphi)$
et que le symbole $[x_1, \ldots, x_b]$ engendre $\det_\kappa(I_\psi)$.
Choisissons $\tilde x_i \in M$ tels que $\varphi(\tilde x_i) = x_i$.
Choisissons $\tilde y_j \in M$ tels que $\psi(\tilde y_j) = y_j$.
Alors $\det_\kappa(M, \varphi, \psi)$ est caractérisé
par l'égalité
$$
[x_1, \ldots, x_a, \tilde y_1, \ldots, \tilde y_b]
=
(-1)^{ab} \det\nolimits_\kappa(M, \varphi, \psi)
[y_1, \ldots, y_b, \tilde x_1, \ldots, \tilde x_a]
$$
dans $\det_\kappa(M)$. Cela explique également le signe.
\end{remark}

\begin{lemma}
\label{lemma-periodic-determinant-shift}
Soit $R$ un anneau local de corps résiduel $\kappa$.
Soit $(M, \varphi, \psi)$ un complexe $(2, 1)$-périodique sur $R$.
Supposons que $M$ soit de longueur finie et que $(M, \varphi, \psi)$ soit
exact. Alors
$$
\det\nolimits_\kappa(M, \varphi, \psi)
\det\nolimits_\kappa(M, \psi, \varphi)
= 1.
$$
\end{lemma}

\begin{proof}
La démonstration est omise.
\end{proof}

\begin{lemma}
\label{lemma-periodic-determinant-sign}
Soit $R$ un anneau local de corps résiduel $\kappa$.
Soit $(M, \varphi, \varphi)$ un complexe $(2, 1)$-périodique sur $R$.
Supposons que $M$ soit de longueur finie et que $(M, \varphi, \varphi)$ soit
exact. Alors $\text{longueur}_R(M) = 2 \text{longueur}_R(\Im(\varphi))$
et
$$
\det\nolimits_\kappa(M, \varphi, \varphi)
=
(-1)^{\text{longueur}_R(\Im(\varphi))}
=
(-1)^{\frac{1}{2}\text{longueur}_R(M)}
$$
\end{lemma}

\begin{proof}
Cela résulte directement de la règle des signes dans les définitions.
\end{proof}

\begin{lemma}
\label{lemma-periodic-determinant-easy-case}
Soit $R$ un anneau local de corps résiduel $\kappa$.
Soit $M$ un $R$-module de longueur finie.
\begin{enumerate}
\item si $\varphi : M \to M$ est un isomorphisme, alors
$\det_\kappa(M, \varphi, 0) = \det_\kappa(\varphi)$.
\item si $\psi : M \to M$ est un isomorphisme, alors
$\det_\kappa(M, 0, \psi) = \det_\kappa(\psi)^{-1}$.
\end{enumerate}
\end{lemma}

\begin{proof}
Démontrons (1). Posons $\psi = 0$. Avec les notations
précédant la définition \ref{definition-periodic-determinant}, nous pouvons alors identifier
$K_\varphi = I_\psi = 0$, $I_\varphi = K_\psi = M$.
Sous ces identifications, l'application
$$
\gamma_\varphi :
\kappa \otimes \det\nolimits_\kappa(M)
=
\det\nolimits_\kappa(K_\varphi)
\otimes
\det\nolimits_\kappa(I_\varphi)
\longrightarrow
\det\nolimits_\kappa(M)
$$
s'identifie à $\det_\kappa(\varphi^{-1})$. D'autre part,
l'application $\gamma_\psi$ s'identifie à l'application identité. Par conséquent,
$\gamma_\psi \circ \sigma \circ \gamma_\varphi^{-1}$ est égale
à $\det_\kappa(\varphi)$ dans ce cas. D'où le résultat.
Nous omettons la démonstration de (2).
\end{proof}

\begin{lemma}
\label{lemma-periodic-determinant}
Soit $R$ un anneau local de corps résiduel $\kappa$.
Supposons donnée une suite exacte courte de
complexes $(2, 1)$-périodiques
$$
0 \to (M_1, \varphi_1, \psi_1)
\to (M_2, \varphi_2, \psi_2)
\to (M_3, \varphi_3, \psi_3)
\to 0
$$
où tous les $M_i$ sont de longueur finie, et chaque $(M_1, \varphi_1, \psi_1)$ est exact.
Alors
$$
\det\nolimits_\kappa(M_2, \varphi_2, \psi_2) =
\det\nolimits_\kappa(M_1, \varphi_1, \psi_1)
\det\nolimits_\kappa(M_3, \varphi_3, \psi_3).
$$
. dans $\kappa^*$.
\end{lemma}

\begin{proof}
Abrégeons
$I_{\varphi, i} = \Im(\varphi_i)$,
$K_{\varphi, i} = \Ker(\varphi_i)$,
$I_{\psi, i} = \Im(\psi_i)$, et
$K_{\psi, i} = \Ker(\psi_i)$.
Remarquons que nous avons un carré commutatif
$$
\xymatrix{
& 0 \ar[d] & 0 \ar[d] & 0 \ar[d] & \\
0 \ar[r] &
K_{\varphi, 1} \ar[r] \ar[d] &
K_{\varphi, 2} \ar[r] \ar[d] &
K_{\varphi, 3} \ar[r] \ar[d] &
0 \\
0 \ar[r] &
M_1 \ar[r] \ar[d] &
M_2 \ar[r] \ar[d] &
M_3 \ar[r] \ar[d] &
0 \\
0 \ar[r] &
I_{\varphi, 1} \ar[r] \ar[d] &
I_{\varphi, 2} \ar[r] \ar[d] &
I_{\varphi, 3} \ar[r] \ar[d] &
0 \\
& 0  & 0  & 0  &
}
$$
de $R$-modules de longueur finie, à lignes et colonnes exactes.
La ligne supérieure est exacte puisqu'elle s'identifie à la
suite $I_{\psi, 1} \to I_{\psi, 2} \to I_{\psi, 3} \to 0$
des images, et de même pour la ligne inférieure. Il existe un diagramme analogue
faisant intervenir les modules $I_{\psi, i}$ et $K_{\psi, i}$.
Par définition, $\det_\kappa(M_2, \varphi_2, \psi_2)$
correspond, au signe près, à la composée des applications verticales de gauche
dans le diagramme suivant
$$
\xymatrix{
\det_\kappa(M_1) \otimes
\det_\kappa(M_3) \ar[r]^\gamma
\ar[d]^{\gamma^{-1} \otimes \gamma^{-1}} &
\det_\kappa(M_2)
\ar[d]^{\gamma^{-1}} \\
\det\nolimits_\kappa(K_{\varphi, 1})
\otimes
\det\nolimits_\kappa(I_{\varphi, 1})
\otimes
\det\nolimits_\kappa(K_{\varphi, 3})
\otimes
\det\nolimits_\kappa(I_{\varphi, 3})
\ar[d]^{\sigma \otimes \sigma}
\ar[r]^-{\gamma \otimes \gamma} &
\det\nolimits_\kappa(K_{\varphi, 2})
\otimes
\det\nolimits_\kappa(I_{\varphi, 2})
\ar[d]^\sigma
\\
\det\nolimits_\kappa(K_{\psi, 1})
\otimes
\det\nolimits_\kappa(I_{\psi, 1})
\otimes
\det\nolimits_\kappa(K_{\psi, 3})
\otimes
\det\nolimits_\kappa(I_{\psi, 3})
\ar[d]^{\gamma \otimes \gamma}
\ar[r]^-{\gamma \otimes \gamma}
&
\det\nolimits_\kappa(K_{\psi, 2})
\otimes
\det\nolimits_\kappa(I_{\psi, 2})
\ar[d]^\gamma \\
\det_\kappa(M_1)
\otimes
\det_\kappa(M_3) \ar[r]^\gamma
&
\det_\kappa(M_2)
}
$$
Les carrés supérieur et inférieur sont commutatifs au signe près
en appliquant le lemme \ref{lemma-det-exact-sequences} (2).
Le carré du milieu est trivialement
commutatif (nous ne faisons que permuter les facteurs). Nous voyons donc
que
$\det\nolimits_\kappa(M_2, \varphi_2, \psi_2) =
\epsilon \det\nolimits_\kappa(M_1, \varphi_1, \psi_1)
\det\nolimits_\kappa(M_3, \varphi_3, \psi_3)
$
pour un certain signe $\epsilon$. Ce signe se calcule :
le rectangle extérieur du diagramme ci-dessus commute au signe près
\begin{eqnarray*}
\epsilon & = &
(-1)^{\text{longueur}(I_{\varphi, 1})\text{longueur}(K_{\varphi, 3})
+ \text{longueur}(I_{\psi, 1})\text{longueur}(K_{\psi, 3})} \\
& = &
(-1)^{\text{longueur}(I_{\varphi, 1})\text{longueur}(I_{\psi, 3})
+ \text{longueur}(I_{\psi, 1})\text{longueur}(I_{\varphi, 3})}
\end{eqnarray*}
(démonstration omise). Il en résulte aisément que les signes
conviennent également.
\end{proof}

\begin{example}
\label{example-dual-numbers}
Soit $k$ un corps.
Considérons l'anneau $R = k[T]/(T^2)$ des nombres duaux sur $k$.
Notons $t$ la classe de $T$ dans $R$.
Soient $M = R$, $\varphi = ut$ et $\psi = vt$, où $u, v \in k^*$.
Dans ce cas, $\det_k(M)$ a pour générateur $e = [t, 1]$.
Nous identifions $I_\varphi = K_\varphi = I_\psi = K_\psi = (t)$.
Alors $\gamma_\varphi(t \otimes t) = u^{-1}[t, 1]$
(puisque $u^{-1} \in M$ est un relèvement de $t \in I_\varphi$)
et $\gamma_\psi(t \otimes t) = v^{-1}[t, 1]$ (pour la même raison).
Nous voyons donc que $\det_k(M, \varphi, \psi) = -u/v \in k^*$.
\end{example}

\begin{example}
\label{example-Zp}
Soient $R = \mathbf{Z}_p$ et $M = \mathbf{Z}_p/(p^l)$.
Soient $\varphi = p^b u$ et $\varphi = p^a v$, où $a, b \geq 0$,
$a + b = l$ et $u, v \in \mathbf{Z}_p^*$.
Alors un calcul analogue à celui de l'exemple \ref{example-dual-numbers}
montre que
\begin{eqnarray*}
\det\nolimits_{\mathbf{F}_p}(\mathbf{Z}_p/(p^l), p^bu, p^av) & = &
(-1)^{ab}u^a/v^b \bmod p \\
& = &
(-1)^{\text{ord}_p(\alpha)\text{ord}_p(\beta)}
\frac{\alpha^{\text{ord}_p(\beta)}}{\beta^{\text{ord}_p(\alpha)}} \bmod p
\end{eqnarray*}
où $\alpha = p^bu, \beta = p^av \in \mathbf{Z}_p$.
Voir le lemme \ref{lemma-symbol-is-usual-tame-symbol}
pour un cas plus général (et une démonstration).
\end{example}

\begin{example}
\label{example-generic-vector-space}
Soit $R = k$ un corps.
Soit $M = k^{\oplus a} \oplus k^{\oplus b}$, de dimension $l = a + b$.
Soient $\varphi$ et $\psi$ les matrices diagonales suivantes
$$
\varphi = \text{diag}(u_1, \ldots, u_a, 0, \ldots, 0),
\quad
\psi = \text{diag}(0, \ldots, 0, v_1, \ldots, v_b)
$$
où $u_i, v_j \in k^*$. Dans ce cas, nous avons
$$
\det\nolimits_k(M, \varphi, \psi)
=
\frac{u_1 \ldots u_a}{v_1 \ldots v_b}.
$$
On peut le voir par un calcul direct, ou en calculant le cas $l = 1$
et en utilisant l'additivité du lemme \ref{lemma-periodic-determinant}.
\end{example}

\begin{example}
\label{example-special-vector-space}
Soit $R = k$ un corps.
Soit $M = k^{\oplus a} \oplus k^{\oplus a}$, de dimension $l = 2a$.
Soient $\varphi$ et $\psi$ les matrices par blocs suivantes
$$
\varphi =
\left(
\begin{matrix}
0 & U \\
0 & 0
\end{matrix}
\right),
\quad
\psi =
\left(
\begin{matrix}
0 & V \\
0 & 0
\end{matrix}
\right),
$$
où $U, V \in \text{Mat}(a \times a, k)$ sont inversibles.
Dans ce cas, nous avons
$$
\det\nolimits_k(M, \varphi, \psi)
=
(-1)^a\frac{\det(U)}{\det(V)}.
$$
On peut le voir par un calcul direct.
Le cas $a = 1$ est analogue au calcul de
l'exemple \ref{example-dual-numbers}.
\end{example}

\begin{example}
\label{example-a-la-oort}
Soit $R = k$ un corps.
Soit $M = k^{\oplus 4}$.
Soit
$$
\varphi =
\left(
\begin{matrix}
  0 &   0 &   0 &   0 \\
u_1 &   0 &   0 &   0 \\
  0 &   0 &   0 &   0 \\
  0 &   0 & u_2 &   0
\end{matrix}
\right)
\quad
\varphi =
\left(
\begin{matrix}
  0 &   0 &   0 &   0 \\
  0 &   0 & v_2 &   0 \\
  0 &   0 &   0 &   0 \\
v_1 &   0 &   0 &   0
\end{matrix}
\right)
\quad
$$
où $u_1, u_2, v_1, v_2 \in k^*$.
Alors nous avons
$$
\det\nolimits_k(M, \varphi, \psi) = -\frac{u_1u_2}{v_1v_2}.
$$
\end{example}

\noindent
Nous en venons maintenant à l'analogue du fait que le déterminant
d'une composée d'endomorphismes linéaires est le produit des
déterminants. Pour éviter des formules très longues, nous
écrivons $I_\varphi = \Im(\varphi)$ et
$K_\varphi = \Ker(\varphi)$
pour toute application de $R$-modules $\varphi : M \to M$.
Nous notons également $\varphi\psi = \varphi \circ \psi$
pour une paire de morphismes $\varphi, \psi : M \to M$.

\begin{lemma}
\label{lemma-multiplicativity-determinant}
Soit $R$ un anneau local de corps résiduel $\kappa$.
Soit $M$ un $R$-module de longueur finie.
Soient $\alpha, \beta, \gamma$ des endomorphismes de $M$.
Supposons que
\begin{enumerate}
\item $I_\alpha = K_{\beta\gamma}$, et de même pour toute permutation
de $\alpha, \beta, \gamma$,
\item $K_\alpha = I_{\beta\gamma}$, et de même pour toute permutation
de $\alpha, \beta, \gamma$.
\end{enumerate}
Alors
\begin{enumerate}
\item le triplet $(M, \alpha, \beta\gamma)$
est un complexe $(2, 1)$-périodique exact ;
\item le triplet $(I_\gamma, \alpha, \beta)$
est un complexe $(2, 1)$-périodique exact ;
\item le triplet $(M/K_\beta, \alpha, \gamma)$
est un complexe $(2, 1)$-périodique exact ;
\item nous avons
$$
\det\nolimits_\kappa(M, \alpha, \beta\gamma)
=
\det\nolimits_\kappa(I_\gamma, \alpha, \beta)
\det\nolimits_\kappa(M/K_\beta, \alpha, \gamma).
$$
\end{enumerate}
\end{lemma}

\begin{proof}
Il est clair que les hypothèses impliquent la partie (1) du lemme.

\medskip\noindent
Pour voir la partie (1), remarquons que les hypothèses impliquent que
$I_{\gamma\alpha} = I_{\alpha\gamma}$, et de même pour les noyaux
et pour toute autre paire de morphismes.
De plus, nous voyons que
$I_{\gamma\beta} =I_{\beta\gamma} = K_\alpha \subset I_\gamma$, et
de même pour toute autre paire. En particulier, nous obtenons une suite exacte courte
$$
0 \to I_{\beta\gamma} \to I_\gamma \xrightarrow{\alpha} I_{\alpha\gamma} \to 0
$$
et, de même, nous obtenons une suite exacte courte
$$
0 \to I_{\alpha\gamma} \to I_\gamma \xrightarrow{\beta} I_{\beta\gamma} \to 0.
$$
Cela démontre que $(I_\gamma, \alpha, \beta)$ est un complexe $(2, 1)$-périodique
exact. La partie (2) du lemme est donc vraie.

\medskip\noindent
Pour voir que $\alpha$ et $\gamma$ induisent des endomorphismes bien définis
de $M/K_\beta$, nous devons vérifier que $\alpha(K_\beta) \subset K_\beta$
et $\gamma(K_\beta) \subset K_\beta$. Cela est vrai car
$\alpha(K_\beta) = \alpha(I_{\gamma\alpha}) = I_{\alpha\gamma\alpha}
\subset I_{\alpha\gamma} = K_\beta$, et de même dans l'autre cas.
Le noyau de l'application $\alpha : M/K_\beta \to M/K_\beta$ est
$K_{\beta\alpha}/K_\beta = I_\gamma/K_\beta$. De même,
le noyau de $\gamma : M/K_\beta \to M/K_\beta$ est égal à
$I_\alpha/K_\beta$. Nous concluons donc que (3) est vraie.

\medskip\noindent
Introduisons $r = \text{longueur}_R(K_\alpha)$,
$s = \text{longueur}_R(K_\beta)$ et $t = \text{longueur}_R(K_\gamma)$.
D'après les suites exactes ci-dessus et nos hypothèses, nous avons
$\text{longueur}_R(I_\alpha) = s + t$, $\text{longueur}_R(I_\beta) = r + t$,
$\text{longueur}_R(I_\gamma) = r + s$, et
$\text{longueur}(M) = r + s + t$.
Choisissons
\begin{enumerate}
\item une suite admissible $x_1, \ldots, x_r \in K_\alpha$
engendrant $K_\alpha$,
\item une suite admissible $y_1, \ldots, y_s \in K_\beta$
engendrant $K_\beta$,
\item une suite admissible $z_1, \ldots, z_t \in K_\gamma$
engendrant $K_\gamma$,
\item des éléments $\tilde x_i \in M$ tels que $\beta\gamma\tilde x_i = x_i$,
\item des éléments $\tilde y_i \in M$ tels que $\alpha\gamma\tilde y_i = y_i$,
\item des éléments $\tilde z_i \in M$ tels que $\beta\alpha\tilde z_i = z_i$.
\end{enumerate}
Avec ces choix, la suite
$y_1, \ldots, y_s, \alpha\tilde z_1, \ldots, \alpha\tilde z_t$
est une suite admissible dans $I_\alpha$ qui l'engendre.
Ainsi, d'après la remarque \ref{remark-more-elementary}, le déterminant
$D = \det_\kappa(M, \alpha, \beta\gamma)$ est
l'unique élément de $\kappa^*$ tel que
\begin{align*}
[y_1, \ldots, y_s,
\alpha\tilde z_1, \ldots, \alpha\tilde z_s,
\tilde x_1, \ldots, \tilde x_r] \\
= (-1)^{r(s + t)} D
[x_1, \ldots, x_r,
\gamma\tilde y_1, \ldots, \gamma\tilde y_s,
\tilde z_1, \ldots, \tilde z_t]
\end{align*}
D'après la même remarque, nous voyons que
$D_1 = \det_\kappa(M/K_\beta, \alpha, \gamma)$
est caractérisé par
$$
[y_1, \ldots, y_s,
\alpha\tilde z_1, \ldots, \alpha\tilde z_t,
\tilde x_1, \ldots, \tilde x_r]
=
(-1)^{rt} D_1
[y_1, \ldots, y_s,
\gamma\tilde x_1, \ldots, \gamma\tilde x_r,
\tilde z_1, \ldots, \tilde z_t]
$$
D'après la même remarque, nous voyons que
$D_2 = \det_\kappa(I_\gamma, \alpha, \beta)$ est caractérisé par
$$
[y_1, \ldots, y_s,
\gamma\tilde x_1, \ldots, \gamma\tilde x_r,
\tilde z_1, \ldots, \tilde z_t]
=
(-1)^{rs} D_2
[x_1, \ldots, x_r,
\gamma\tilde y_1, \ldots, \gamma\tilde y_s,
\tilde z_1, \ldots, \tilde z_t]
$$
En combinant les formules ci-dessus, nous voyons que $D = D_1 D_2$,
comme voulu.
\end{proof}


\begin{lemma}
\label{lemma-tricky}
Soit $R$ un anneau local de corps résiduel $\kappa$.
Soit $\alpha : (M, \varphi, \psi) \to (M', \varphi', \psi')$
un morphisme de complexes $(2, 1)$-périodiques sur $R$.
Supposons que
\begin{enumerate}
\item $M$ et $M'$ soient de longueur finie ;
\item $(M, \varphi, \psi)$ et $(M', \varphi', \psi')$ soient exacts ;
\item les applications $\varphi$ et $\psi$ induisent l'application nulle sur
$K = \Ker(\alpha)$ ; et
\item les applications $\varphi$ et $\psi$ induisent l'application nulle sur
$Q = \Coker(\alpha)$.
\end{enumerate}
Notons $N = \alpha(M) \subset M'$. Nous obtenons deux suites exactes courtes
de complexes $(2, 1)$-périodiques
$$
\begin{matrix}
0 \to (N, \varphi', \psi') \to (M', \varphi', \psi') \to (Q, 0, 0) \to 0 \\
0 \to (K, 0, 0) \to (M, \varphi, \psi) \to (N, \varphi', \psi') \to 0
\end{matrix}
$$
qui induisent deux isomorphismes $\alpha_i : Q \to K$, $i = 0, 1$. Alors
$$
\det\nolimits_\kappa(M, \varphi, \psi)
=
\det\nolimits_\kappa(\alpha_0^{-1} \circ \alpha_1)
\det\nolimits_\kappa(M', \varphi', \psi')
$$
En particulier, si $\alpha_0 = \alpha_1$, alors
$\det\nolimits_\kappa(M, \varphi, \psi) =
\det\nolimits_\kappa(M', \varphi', \psi')$.
\end{lemma}

\begin{proof}
Il existe (au moins) deux façons de démontrer ce lemme. L'une consiste à produire un
énorme diagramme commutatif en utilisant les propriétés des déterminants.
L'autre consiste à employer la caractérisation des déterminants au moyen
de suites admissibles d'éléments. C'est la seconde approche que nous
utiliserons.

\medskip\noindent
Expliquons d'abord précisément ce que sont les applications $\alpha_i$.
À savoir, $\alpha_0$ est la composée
$$
\alpha_0 : Q = H^0(Q, 0, 0) \to H^1(N, \varphi', \psi') \to H^2(K, 0, 0) = K
$$
et $\alpha_1$ est la composée
$$
\alpha_1 : Q = H^1(Q, 0, 0) \to H^2(N, \varphi', \psi') \to H^3(K, 0, 0) = K
$$
provenant des morphismes de bord des suites exactes courtes de complexes
affichées dans le lemme. Le fait que les
complexes $(M, \varphi, \psi)$ et $(M', \varphi', \psi')$ soient exacts
implique que ces applications sont des isomorphismes.

\medskip\noindent
Nous utiliserons les notations $I_\varphi = \Im(\varphi)$,
$K_\varphi = \Ker(\varphi)$, et de même pour les autres applications.
L'exactitude pour $M$ et $M'$
signifie que $K_\varphi = I_\psi$, ainsi que trois égalités analogues.
Introduisons $k = \text{longueur}_R(K)$, $a = \text{longueur}_R(I_\varphi)$
et $b = \text{longueur}_R(I_\psi)$. Nous voyons alors que $\text{longueur}_R(M) = a + b$,
ainsi que $\text{longueur}_R(N) = a + b - k$, $\text{longueur}_R(Q) = k$
et $\text{longueur}_R(M') = a + b$. Les suites exactes ci-dessous montreront
qu'également $\text{longueur}_R(I_{\varphi'}) = a$ et
$\text{longueur}_R(I_{\psi'}) = b$.

\medskip\noindent
L'hypothèse $K \subset K_\varphi = I_\psi$ signifie que
$\varphi$ se factorise par $N$ pour donner une suite exacte
$$
0 \to \alpha(I_\psi) \to N \xrightarrow{\varphi\alpha^{-1}} I_\psi \to 0.
$$
Ici, $\varphi\alpha^{-1}(x') = y$ signifie que $x' = \alpha(x)$ et $y = \varphi(x)$.
De même, nous avons
$$
0 \to \alpha(I_\varphi) \to N \xrightarrow{\psi\alpha^{-1}} I_\varphi \to 0.
$$
L'hypothèse selon laquelle $\psi'$ induit l'application nulle sur
$Q$ signifie que $I_{\psi'} = K_{\varphi'} \subset N$.
Cela signifie que le quotient $\varphi'(N) \subset I_{\varphi'}$
s'identifie à $Q$. Remarquons que $\varphi'(N) = \alpha(I_\varphi)$.
Nous en concluons donc qu'il existe un isomorphisme
$$
\varphi' : Q \to I_{\varphi'}/\alpha(I_\varphi)
$$
décrit simplement par
$\varphi'(x' \bmod N) = \varphi'(x') \bmod \alpha(I_\varphi)$.
Exactement de la même manière, nous obtenons
$$
\psi' : Q \to I_{\psi'}/\alpha(I_\psi)
$$
Enfin, remarquons que $\alpha_0$ est la composée
$$
\xymatrix{
Q \ar[r]^-{\varphi'} &
I_{\varphi'}/\alpha(I_\varphi)
\ar[rrr]^-{\psi\alpha^{-1}|_{I_{\varphi'}/\alpha(I_\varphi)}} & & &
K
}
$$
et, de même,
$\alpha_1 = \varphi\alpha^{-1}|_{I_{\psi'}/\alpha(I_\psi)} \circ \psi'$.

\medskip\noindent
Pour abréger les formules ci-dessous, nous écrirons désormais $\alpha x$
au lieu de $\alpha(x)$. Il ne devrait en résulter aucune confusion puisque
toutes les applications sont désignées par des lettres grecques et les éléments par des lettres romaines.
Nous allons choisir
\begin{enumerate}
\item une suite admissible $z_1, \ldots, z_k \in K$
engendrant $K$,
\item des éléments $z'_i \in M$ tels que $\varphi z'_i = z_i$,
\item des éléments $z''_i \in M$ tels que $\psi z''_i = z_i$,
\item des éléments $x_{k + 1}, \ldots, x_a \in I_\varphi$ tels
que $z_1, \ldots, z_k, x_{k + 1}, \ldots, x_a$ soit une suite admissible
engendrant $I_\varphi$,
\item des éléments $\tilde x_i \in M$ tels que $\varphi \tilde x_i = x_i$,
\item des éléments $y_{k + 1}, \ldots, y_b \in I_\psi$ tels que
$z_1, \ldots, z_k, y_{k + 1}, \ldots, y_b$ soit une suite admissible
engendrant $I_\psi$,
\item des éléments $\tilde y_i \in M$ tels que $\psi \tilde y_i = y_i$, et
\item des éléments $w_1, \ldots, w_k \in M'$ tels que
$w_1 \bmod N, \ldots, w_k \bmod N$ forment une suite admissible
dans $Q$ qui engendre $Q$.
\end{enumerate}
D'après la remarque \ref{remark-more-elementary}, l'élément
$D = \det_\kappa(M, \varphi, \psi) \in \kappa^*$ est
caractérisé par
\begin{eqnarray*}
& &
[z_1, \ldots, z_k,
x_{k + 1}, \ldots, x_a,
z''_1, \ldots, z''_k,
\tilde y_{k + 1}, \ldots, \tilde y_b] \\
& = &
(-1)^{ab} D
[z_1, \ldots, z_k,
y_{k + 1}, \ldots, y_b,
z'_1, \ldots, z'_k,
\tilde x_{k + 1}, \ldots, \tilde x_a]
\end{eqnarray*}
Remarquons que, d'après la discussion ci-dessus,
$\alpha x_{k + 1}, \ldots, \alpha x_a, \varphi w_1, \ldots, \varphi w_k$
est une suite admissible engendrant $I_{\varphi'}$, et que
$\alpha y_{k + 1}, \ldots, \alpha y_b, \psi w_1, \ldots, \psi w_k$
est une suite admissible engendrant $I_{\psi'}$.
Ainsi, d'après la remarque \ref{remark-more-elementary}, l'élément
$D' = \det_\kappa(M', \varphi', \psi') \in \kappa^*$ est
caractérisé par
\begin{eqnarray*}
& &
[\alpha x_{k + 1}, \ldots, \alpha x_a,
\varphi' w_1, \ldots, \varphi' w_k,
\alpha \tilde y_{k + 1}, \ldots, \alpha \tilde y_b,
w_1, \ldots, w_k]
\\
& = &
(-1)^{ab} D'
[\alpha y_{k + 1}, \ldots, \alpha y_b,
\psi' w_1, \ldots, \psi' w_k,
\alpha \tilde x_{k + 1}, \ldots, \alpha \tilde x_a,
w_1, \ldots, w_k]
\end{eqnarray*}
Remarquons que, dans la première formule affichée, resp.\ la seconde,
les $k$ premières, resp.\ dernières, entrées des symboles des deux côtés
sont les mêmes. Ces formules sont donc réellement équivalentes aux
égalités
\begin{eqnarray*}
& &
[\alpha x_{k + 1}, \ldots, \alpha x_a,
\alpha z''_1, \ldots, \alpha z''_k,
\alpha \tilde y_{k + 1}, \ldots, \alpha \tilde y_b] \\
& = &
(-1)^{ab} D
[\alpha y_{k + 1}, \ldots, \alpha y_b,
\alpha z'_1, \ldots, \alpha z'_k,
\alpha \tilde x_{k + 1}, \ldots, \alpha \tilde x_a]
\end{eqnarray*}
et
\begin{eqnarray*}
& &
[\alpha x_{k + 1}, \ldots, \alpha x_a,
\varphi' w_1, \ldots, \varphi' w_k,
\alpha \tilde y_{k + 1}, \ldots, \alpha \tilde y_b]
\\
& = &
(-1)^{ab} D'
[\alpha y_{k + 1}, \ldots, \alpha y_b,
\psi' w_1, \ldots, \psi' w_k,
\alpha \tilde x_{k + 1}, \ldots, \alpha \tilde x_a]
\end{eqnarray*}
dans $\det_\kappa(N)$. Remarquons que
$\varphi' w_1, \ldots, \varphi' w_k$
et
$\alpha z''_1, \ldots, z''_k$
sont des suites admissibles engendrant le module
$I_{\varphi'}/\alpha(I_\varphi)$. Écrivons
$$
[\varphi' w_1, \ldots, \varphi' w_k]
= \lambda_0 [\alpha z''_1, \ldots, \alpha z''_k]
$$
dans $\det_\kappa(I_{\varphi'}/\alpha(I_\varphi))$
pour un certain $\lambda_0 \in \kappa^*$. De même,
écrivons
$$
[\psi' w_1, \ldots, \psi' w_k]
= \lambda_1 [\alpha z'_1, \ldots, \alpha z'_k]
$$
dans $\det_\kappa(I_{\psi'}/\alpha(I_\psi))$
pour un certain $\lambda_1 \in \kappa^*$. D'une part,
il est clair que
$$
\alpha_i([w_1, \ldots, w_k]) = \lambda_i[z_1, \ldots, z_k]
$$
pour $i = 0, 1$, d'après notre description de $\alpha_i$ ci-dessus,
ce qui signifie que
$$
\det\nolimits_\kappa(\alpha_0^{-1} \circ \alpha_1)
=
\lambda_1/\lambda_0
$$
et,
d'autre part, il est clair que
\begin{eqnarray*}
& &
\lambda_0 [\alpha x_{k + 1}, \ldots, \alpha x_a,
\alpha z''_1, \ldots, \alpha z''_k,
\alpha \tilde y_{k + 1}, \ldots, \alpha \tilde y_b] \\
& = &
[\alpha x_{k + 1}, \ldots, \alpha x_a,
\varphi' w_1, \ldots, \varphi' w_k,
\alpha \tilde y_{k + 1}, \ldots, \alpha \tilde y_b]
\end{eqnarray*}
et
\begin{eqnarray*}
& &
\lambda_1[\alpha y_{k + 1}, \ldots, \alpha y_b,
\alpha z'_1, \ldots, \alpha z'_k,
\alpha \tilde x_{k + 1}, \ldots, \alpha \tilde x_a] \\
& = &
[\alpha y_{k + 1}, \ldots, \alpha y_b,
\psi' w_1, \ldots, \psi' w_k,
\alpha \tilde x_{k + 1}, \ldots, \alpha \tilde x_a]
\end{eqnarray*}
ce qui implique $\lambda_0 D = \lambda_1 D'$. Le lemme en résulte.
\end{proof}








\subsection{Symboles}
\label{subsection-symbols}

\noindent
Le bon degré de généralité pour cette construction est peut-être la
situation du lemme suivant.

\begin{lemma}
\label{lemma-pre-symbol}
Soit $A$ un anneau local noethérien.
Soit $M$ un $A$-module fini de dimension $1$.
Supposons que $\varphi, \psi : M \to M$ soient deux applications injectives de
$A$-modules, et supposons que $\varphi(\psi(M)) = \psi(\varphi(M))$,
par exemple si $\varphi$ et $\psi$ commutent.
Alors $\text{longueur}_A(M/\varphi\psi M) < \infty$
et $(M/\varphi\psi M, \varphi, \psi)$ est un complexe
$(2, 1)$-périodique exact.
\end{lemma}

\begin{proof}
Soit $\mathfrak q$ un idéal premier minimal du support de $M$.
Alors $M_{\mathfrak q}$ est un $A_{\mathfrak q}$-module de longueur finie,
voir Algèbre, lemme \ref{algebra-lemma-support-point}.
Ainsi, $\varphi$ et $\psi$
induisent des isomorphismes $M_{\mathfrak q} \to M_{\mathfrak q}$.
Le support de $M/\varphi\psi M$ est donc $\{\mathfrak m_A\}$
et celui-ci est ainsi de longueur finie (voir le lemme cité ci-dessus).
Enfin, le noyau de $\varphi$ sur $M/\varphi\psi M$
est manifestement $\psi M/\varphi\psi M$ ; le noyau
de $\varphi$ est donc l'image de $\psi$ sur $M/\varphi\psi M$.
De même dans l'autre sens, puisque $M/\varphi\psi M = M/\psi\varphi M$
par hypothèse.
\end{proof}

\begin{lemma}
\label{lemma-symbol-defined}
Soit $A$ un anneau local noethérien. Soient $a, b \in A$.
\begin{enumerate}
\item si $M$ est un $A$-module fini de dimension $1$
tel que $a, b$ soient des non-diviseurs de zéro dans $M$, alors
$\text{longueur}_A(M/abM) < \infty$ et
$(M/abM, a, b)$ est un complexe exact $(2, 1)$-périodique ;
\item si $a, b$ sont des non-diviseurs de zéro et $\dim(A) = 1$,
alors $\text{longueur}_A(A/(ab)) < \infty$ et
$(A/(ab), a, b)$ est un complexe exact $(2, 1)$-périodique.
\end{enumerate}
En particulier, dans ces cas,
$\det_\kappa(M/abM, a, b) \in \kappa^*$,
resp.\ $\det_\kappa(A/(ab), a, b) \in \kappa^*$
sont définis.
\end{lemma}

\begin{proof}
Cela résulte du lemme \ref{lemma-pre-symbol}.
\end{proof}

\begin{definition}
\label{definition-symbol-M}
Soit $A$ un anneau local noethérien de corps résiduel $\kappa$.
Soient $a, b \in A$.
Soit $M$ un $A$-module fini de dimension $1$
tel que $a, b$ soient des non-diviseurs de zéro dans $M$.
Nous définissons le {\it symbole associé à $M, a, b$}
comme l'élément
$$
d_M(a, b) =
\det\nolimits_\kappa(M/abM, a, b) \in \kappa^*
$$
\end{definition}

\begin{lemma}
\label{lemma-multiplicativity-symbol}
Soit $A$ un anneau local noethérien.
Soient $a, b, c \in A$. Soit $M$ un $A$-module fini
tel que $\dim(\text{Supp}(M)) = 1$. Supposons que $a, b, c$ soient des non-diviseurs de zéro dans $M$.
Alors
$$
d_M(a, bc) = d_M(a, b) d_M(a, c)
$$
et $d_M(a, b)d_M(b, a) = 1$.
\end{lemma}

\begin{proof}
La première assertion résulte du lemme \ref{lemma-multiplicativity-determinant}
appliqué à $M/abcM$ et aux endomorphismes $\alpha, \beta, \gamma$ donnés par
la multiplication par $a, b, c$.
La seconde provient du lemme \ref{lemma-periodic-determinant-shift}.
\end{proof}

\begin{definition}
\label{definition-tame-symbol}
Soit $A$ un anneau intègre local noethérien de dimension $1$
et de corps résiduel $\kappa$.
Soit $K$ le corps des fractions de $A$.
Nous définissons le {\it symbole modéré} de $A$ comme l'application
$$
K^* \times K^* \longrightarrow \kappa^*,
\quad
(x, y) \longmapsto d_A(x, y)
$$
où $d_A(x, y)$ est étendu à $K^* \times K^*$ par la multiplicativité du
lemme \ref{lemma-multiplicativity-symbol}.
\end{definition}

\noindent
Il est clair que nous pouvons étendre plus généralement $d_M(-, -)$ à
certains anneaux de fractions de $A$ (même si $A$ n'est pas intègre).

\begin{lemma}
\label{lemma-symbol-when-equal}
Soient $A$ un anneau local noethérien et $M$ un $A$-module fini de
dimension $1$. Soit $a \in A$ un non-diviseur de zéro dans $M$.
Alors $d_M(a, a) = (-1)^{\text{longueur}_A(M/aM)}$.
\end{lemma}

\begin{proof}
C'est une conséquence immédiate du lemme \ref{lemma-periodic-determinant-sign}.
\end{proof}

\begin{lemma}
\label{lemma-symbol-when-one-is-a-unit}
Soit $A$ un anneau local noethérien.
Soit $M$ un $A$-module fini de dimension $1$.
Soient $b \in A$ un non-diviseur de zéro dans $M$ et $u \in A^*$.
Alors
$$
d_M(u, b) = u^{\text{longueur}_A(M/bM)} \bmod \mathfrak m_A.
$$
En particulier, si $M = A$, alors
$d_A(u, b) = u^{\text{ord}_A(b)} \bmod \mathfrak m_A$.
\end{lemma}

\begin{proof}
Remarquons que, dans ce cas, $M/ubM = M/bM$, sur lequel la multiplication
par $b$ est nulle. Ainsi, $d_M(u, b) = \det_\kappa(u|_{M/bM})$
d'après le lemme \ref{lemma-periodic-determinant-easy-case}. Le lemme
résulte alors du lemme \ref{lemma-times-u-determinant}.
\end{proof}

\begin{lemma}
\label{lemma-symbol-short-exact-sequence}
Soit $A$ un anneau local noethérien.
Soient $a, b \in A$.
Soit
$$
0 \to M \to M' \to M'' \to 0
$$
une suite exacte courte de $A$-modules de dimension $1$
tels que $a, b$ soient des non-diviseurs de zéro dans
les trois $A$-modules.
Alors
$$
d_{M'}(a, b) = d_M(a, b) d_{M''}(a, b)
$$
dans $\kappa^*$.
\end{lemma}

\begin{proof}
Il est facile de voir que cela donne une suite exacte courte
de complexes $(2, 1)$-périodiques exacts
$$
0 \to
(M/abM, a, b) \to
(M'/abM', a, b) \to
(M''/abM'', a, b) \to 0
$$
Le lemme résulte donc du lemme \ref{lemma-periodic-determinant}.
\end{proof}

\begin{lemma}
\label{lemma-symbol-compare-modules}
Soit $A$ un anneau local noethérien.
Soit $\alpha : M \to M'$ un homomorphisme de
$A$-modules finis de dimension $1$.
Soient $a, b \in A$. Supposons que
\begin{enumerate}
\item $a$ et $b$ soient des non-diviseurs de zéro dans $M$ et $M'$ ; et
\item $\dim(\Ker(\alpha)), \dim(\Coker(\alpha)) \leq 0$.
\end{enumerate}
Alors $d_M(a, b) = d_{M'}(a, b)$.
\end{lemma}

\begin{proof}
Si $a \in A^*$, l'égalité résulte de
l'égalité $\text{longueur}(M/bM) = \text{longueur}(M'/bM')$
et du lemme \ref{lemma-symbol-when-one-is-a-unit}.
De même, si $b$ est une unité, le lemme est encore vrai
(par la symétrie du lemme \ref{lemma-multiplicativity-symbol}).
Nous pouvons donc supposer que $a, b \in \mathfrak m_A$.
Cela implique en particulier que $\mathfrak m$ n'est pas
un idéal premier associé à $M$, et donc que $\alpha : M \to M'$
est injectif. Nous pouvons ainsi considérer $M$ comme un sous-module de $M'$.
Par hypothèse, $M'/M$ est un $A$-module fini de support
$\{\mathfrak m_A\}$, et il est donc de longueur finie.
Remarquons que, pour tout troisième module $M''$ tel que $M \subset M'' \subset M'$,
les applications $M \to M''$ et $M'' \to M'$ satisfont elles aussi les hypothèses du lemme.
Cela nous ramène, par récurrence sur la longueur de $M'/M$,
au cas où $\text{longueur}_A(M'/M) = 1$.
Enfin, dans ce cas, considérons l'application
$$
\overline{\alpha} : M/abM \longrightarrow M'/abM'.
$$
Par construction, le conoyau $Q$ de $\overline{\alpha}$ est de
longueur $1$. Puisque $a, b \in \mathfrak m_A$, ils agissent trivialement sur
$Q$. Il s'ensuit aussi que le noyau $K$ de $\overline{\alpha}$ est de
longueur $1$, et donc que $a$ et $b$ agissent également trivialement sur $K$.
Nous pouvons ainsi appliquer le lemme \ref{lemma-tricky}. Il suffit donc de voir
que les deux applications $\alpha_i : Q \to K$ sont les mêmes.
En fait, elles sont toutes deux égales à l'application
$q = x' \bmod \Im(\overline{\alpha}) \mapsto abx' \in K$.
Nous omettons la vérification.
\end{proof}

\begin{lemma}
\label{lemma-compute-symbol-M}
Soit $A$ un anneau local noethérien.
Soit $M$ un $A$-module fini tel que $\dim(\text{Supp}(M)) = 1$.
Soient $a, b \in A$ des non-diviseurs de zéro dans $M$.
Soient $\mathfrak q_1, \ldots, \mathfrak q_t$ les idéaux premiers minimaux
du support de $M$. Alors
$$
d_M(a, b)
=
\prod\nolimits_{i = 1, \ldots, t}
d_{A/\mathfrak q_i}(a, b)^{
\text{longueur}_{A_{\mathfrak q_i}}(M_{\mathfrak q_i})}
$$
comme éléments de $\kappa^*$.
\end{lemma}

\begin{proof}
Choisissons une filtration par des $A$-sous-modules
$$
0 = M_0 \subset M_1 \subset \ldots \subset M_n = M
$$
telle que chaque quotient $M_j/M_{j - 1}$ soit isomorphe
à $A/\mathfrak p_j$ pour un certain idéal premier $\mathfrak p_j$
de $A$. Voir Algèbre, lemme \ref{algebra-lemma-filter-Noetherian-module}.
Pour chaque $j$, nous avons soit $\mathfrak p_j = \mathfrak q_i$
pour un certain $i$, soit $\mathfrak p_j = \mathfrak m_A$. De plus,
pour $i$ fixé, le nombre de $j$ tels que
$\mathfrak p_j = \mathfrak q_i$ est égal à
$\text{longueur}_{A_{\mathfrak q_i}}(M_{\mathfrak q_i})$, d'après
Algèbre, lemme \ref{algebra-lemma-filter-minimal-primes-in-support}.
Ainsi, $d_{M_j}(a, b)$ est défini pour chaque $j$, et
$$
d_{M_j}(a, b)
=
\left\{
\begin{matrix}
d_{M_{j - 1}}(a, b) d_{A/\mathfrak q_i}(a, b) &
\text{si} & \mathfrak p_j = \mathfrak q_i \\
d_{M_{j - 1}}(a, b) & \text{si} & \mathfrak p_j = \mathfrak m_A
\end{matrix}
\right.
$$
d'après le lemme \ref{lemma-symbol-short-exact-sequence} dans le premier cas
et le lemme \ref{lemma-symbol-compare-modules} dans le second. D'où le lemme.
\end{proof}

\begin{lemma}
\label{lemma-symbol-is-usual-tame-symbol}
Soit $A$ un anneau de valuation discrète de corps des fractions $K$.
Pour tous $x, y \in K$ non nuls, nous avons
$$
d_A(x, y)
=
(-1)^{\text{ord}_A(x)\text{ord}_A(y)}
\frac{x^{\text{ord}_A(y)}}{y^{\text{ord}_A(x)}} \bmod \mathfrak m_A,
$$
autrement dit, le symbole est égal au symbole modéré usuel.
\end{lemma}

\begin{proof}
Par multiplicativité, il suffit de le démontrer lorsque $x, y \in A$.
Soit $t \in A$ une uniformisante.
Écrivons $x = t^bu$ et $y = t^bv$ pour certains $a, b \geq 0$
et $u, v \in A^*$. Posons $l = a + b$. Alors
$t^{l - 1}, \ldots, t^b$ est une suite admissible dans
$(x)/(xy)$ et $t^{l - 1}, \ldots, t^a$ est une suite admissible
dans $(y)/(xy)$. Ainsi, d'après la remarque \ref{remark-more-elementary},
nous voyons que $d_A(x, y)$ est caractérisé par l'équation
$$
[t^{l - 1}, \ldots, t^b, v^{-1}t^{b - 1}, \ldots, v^{-1}]
=
(-1)^{ab} d_A(x, y)
[t^{l - 1}, \ldots, t^a, u^{-1}t^{a - 1}, \ldots, u^{-1}].
$$
Ainsi, d'après les relations d'admissibilité pour les
symboles $[x_1, \ldots, x_l]$, nous voyons que
$$
d_A(x, y) = (-1)^{ab} u^a/v^b \bmod \mathfrak m_A
$$
comme voulu.
\end{proof}

\begin{lemma}
\label{lemma-symbol-is-steinberg-prepare}
Soit $A$ un anneau local noethérien.
Soient $a, b \in A$.
Soit $M$ un $A$-module fini de dimension $1$ tel que, dans celui-ci,
chacun de $a$, $b$ et $b - a$ soit un non-diviseur de zéro.
Alors
$$
d_M(a, b - a)d_M(b, b) = d_M(b, b - a)d_M(a, b)
$$
dans $\kappa^*$.
\end{lemma}

\begin{proof}
D'après le lemme \ref{lemma-compute-symbol-M}, il suffit de montrer la relation lorsque
$M = A/\mathfrak q$ pour un certain idéal premier $\mathfrak q \subset A$ tel que
$\dim(A/\mathfrak q) = 1$.

\medskip\noindent
Dans le cas $M = A/\mathfrak q$, nous pouvons remplacer $A$ par $A/\mathfrak q$ et
$a, b$ par leurs images dans $A/\mathfrak q$. Nous pouvons donc supposer que $A = M$
et que $A$ est un anneau intègre local noethérien de dimension $1$. La raison en est
que les corps résiduels $\kappa$ de $A$ et de $A/\mathfrak q$ sont
les mêmes, et que, pour tout $A/\mathfrak q$-module $M$, les déterminants
pris sur $A$ ou sur $A/\mathfrak q$ s'identifient canoniquement.
Voir le lemme \ref{lemma-determinant-quotient-ring}.

\medskip\noindent
Il suffit de montrer la relation lorsque
$a, b$ appartiennent tous deux à l'idéal maximal. En effet, le cas où l'un
d'eux, ou les deux, est une unité résulte des lemmes \ref{lemma-symbol-when-one-is-a-unit}
et \ref{lemma-symbol-when-equal}.

\medskip\noindent
Choisissons une extension $A \subset A'$ et des factorisations
$a = ta'$, $b = tb'$ comme dans le
lemme \ref{lemma-Noetherian-domain-dim-1-two-elements}.
Remarquons qu'aussi $b - a = t(b' - a')$ et que
$A' = (a', b') = (a', b' - a') = (b' - a', b')$.
Ici et dans la suite, nous considérons $A'$ comme un $A$-module, et
$a, b, a', b', t$ comme des endomorphismes du $A$-module $A'$.
Nous emploierons les notations $d^A_{A'}(a', b')$, etc., pour désigner
$$
d^A_{A'}(a', b')
=
\det\nolimits_\kappa(A'/a'b'A', a', b')
$$
qui est défini par le lemme \ref{lemma-pre-symbol}. L'indice supérieur ${}^A$
sert à le distinguer du symbole déjà défini
$d_{A'}(a', b')$, qui est différent (par exemple parce qu'il prend ses valeurs
dans le corps résiduel de $A'$, lequel peut différer de $\kappa$).
D'après le lemme \ref{lemma-symbol-compare-modules}, nous voyons que
$d_A(a, b) = d^A_{A'}(a, b)$,
et de même pour les autres combinaisons.
En utilisant ceci et la multiplicativité, nous voyons qu'il suffit de démontrer
$$
d^A_{A'}(a', b' - a') d^A_{A'}(b', b')
=
d^A_{A'}(b', b' - a') d^A_{A'}(a', b')
$$
Maintenant, puisque $(a', b') = A'$, et de même pour les autres idéaux, nous avons
$$
\begin{matrix}
A'/(a'(b' - a')) & \cong & A'/(a') \oplus A'/(b' - a') \\
A'/(b'(b' - a')) & \cong & A'/(b') \oplus A'/(b' - a') \\
A'/(a'b') & \cong & A'/(a') \oplus A'/(b')
\end{matrix}
$$
De plus, remarquons que la multiplication par $b' - a'$ sur
$A/(a')$ est égale à la multiplication par $b'$, et que
la multiplication par $b' - a'$ sur $A/(b')$ est égale à la multiplication par $-a'$.
En utilisant les lemmes
\ref{lemma-periodic-determinant-easy-case} et
\ref{lemma-periodic-determinant}
nous concluons que
$$
\begin{matrix}
d^A_{A'}(a', b' - a') & = &
\det\nolimits_\kappa(b'|_{A'/(a')})^{-1}
\det\nolimits_\kappa(a'|_{A'/(b' - a')}) \\
d^A_{A'}(b', b' - a') & = &
\det\nolimits_\kappa(-a'|_{A'/(b')})^{-1}
\det\nolimits_\kappa(b'|_{A'/(b' - a')}) \\
d^A_{A'}(a', b') & = &
\det\nolimits_\kappa(b'|_{A'/(a')})^{-1}
\det\nolimits_\kappa(a'|_{A'/(b')})
\end{matrix}
$$
Nous concluons donc que
$$
(-1)^{\text{longueur}_A(A'/(b'))}
d^A_{A'}(a', b' - a')
=
d^A_{A'}(b', b' - a') d^A_{A'}(a', b')
$$
le signe provenant du $-a'$ dans la deuxième égalité ci-dessus.
D'autre part, d'après le lemme \ref{lemma-periodic-determinant-sign}, nous avons
$d^A_{A'}(b', b') = (-1)^{\text{longueur}_A(A'/(b'))}$, et le lemme
est démontré.
\end{proof}

\noindent
Le symbole modéré est un symbole de Steinberg.

\begin{lemma}
\label{lemma-symbol-is-steinberg}
Soit $A$ un anneau intègre local noethérien de dimension $1$
et de corps des fractions $K$. Pour $x \in K \setminus \{0, 1\}$,
nous avons
$$
d_A(x, 1 -x) = 1
$$
\end{lemma}

\begin{proof}
Écrivons $x = a/b$ avec $a, b \in A$.
Puisque $1 - x = (b - a)/b$, l'hypothèse implique
que $b - a \not = 0$ aussi. Nous calculons donc
$$
d_A(x, 1 - x)
=
d_A(a, b - a)d_A(a, b)^{-1}d_A(b, b - a)^{-1}d_A(b, b)
$$
Nous devons ainsi montrer que
$d_A(a, b - a) d_A(b, b) = d_A(b, b - a) d_A(a, b)$.
C'est le lemme \ref{lemma-symbol-is-steinberg-prepare}.
\end{proof}










\subsection{Longueurs et déterminants}
\label{subsection-length-determinant}

\noindent
Dans cette section, nous utilisons le déterminant pour comparer des réseaux.
Le lemme clé est le suivant.

\begin{lemma}
\label{lemma-key-lemma}
Soit $R$ un anneau local noethérien.
Soit $\mathfrak q \subset R$ un idéal premier tel que $\dim(R/\mathfrak q) = 1$.
Soit $\varphi : M \to N$ un homomorphisme de $R$-modules finis.
Supposons qu'il existe $x_1, \ldots, x_l \in M$ et $y_1, \ldots, y_l \in M$
ayant les propriétés suivantes :
\begin{enumerate}
\item $M = \langle x_1, \ldots, x_l\rangle$ ;
\item $\langle x_1, \ldots, x_i\rangle / \langle x_1, \ldots, x_{i - 1}\rangle
\cong R/\mathfrak q$ pour $i = 1, \ldots, l$ ;
\item $N = \langle y_1, \ldots, y_l\rangle$ ; et
\item $\langle y_1, \ldots, y_i\rangle / \langle y_1, \ldots, y_{i - 1}\rangle
\cong R/\mathfrak q$ pour $i = 1, \ldots, l$.
\end{enumerate}
Alors $\varphi$ est injectif si et seulement si $\varphi_{\mathfrak q}$ est un
isomorphisme, et dans ce cas nous avons
$$
\text{longueur}_R(\Coker(\varphi)) = \text{ord}_{R/\mathfrak q}(f)
$$
où $f \in \kappa(\mathfrak q)$ est l'élément tel que
$$
[\varphi(x_1), \ldots, \varphi(x_l)] = f [y_1, \ldots, y_l]
$$
dans $\det_{\kappa(\mathfrak q)}(N_{\mathfrak q})$.
\end{lemma}

\begin{proof}
Remarquons d'abord que le lemme est vrai dans le cas $l = 1$.
En effet, dans ce cas, $x_1$ est une base de $M$ sur $R/\mathfrak q$,
$y_1$ est une base de $N$ sur $R/\mathfrak q$, et nous avons
$\varphi(x_1) = fy_1$ pour un certain $f \in R$. Ainsi, $\varphi$ est injectif
si et seulement si $f \not \in \mathfrak q$. De plus,
$\Coker(\varphi) = R/(f, \mathfrak q)$, et donc le lemme
résulte de la définition de $\text{ord}_{R/q}(f)$
(voir Algèbre, définition \ref{algebra-definition-ord}).

\medskip\noindent
En fait, supposons plus généralement que $\varphi(x_i) = f_iy_i$ pour certains
$f_i \in R$, $f_i \not \in \mathfrak q$. Alors les applications induites
$$
\langle x_1, \ldots, x_i\rangle / \langle x_1, \ldots, x_{i - 1}\rangle
\longrightarrow
\langle y_1, \ldots, y_i\rangle / \langle y_1, \ldots, y_{i - 1}\rangle
$$
sont toutes injectives et ont des conoyaux isomorphes à
$R/(f_i, \mathfrak q)$. Nous voyons donc que
$$
\text{longueur}_R(\Coker(\varphi)) = \sum \text{ord}_{R/\mathfrak q}(f_i).
$$
D'autre part, il est clair que
$$
[\varphi(x_1), \ldots, \varphi(x_l)] = f_1 \ldots f_l [y_1, \ldots, y_l]
$$
dans ce cas, d'après la relation d'admissibilité (b) pour les symboles.
Nous voyons donc que le résultat est encore vrai dans ce cas.

\medskip\noindent
Nous démontrons le cas général par récurrence sur $l$. Supposons $l > 1$.
Soit $i \in \{1, \ldots, l\}$ minimal tel que
$\varphi(x_1) \in \langle y_1, \ldots, y_i\rangle$.
Nous allons raisonner par récurrence sur $i$.
Si $i = 1$, nous obtenons un diagramme commutatif
$$
\xymatrix{
0 \ar[r] &
\langle x_1 \rangle \ar[r] \ar[d] &
\langle x_1, \ldots, x_l \rangle \ar[r] \ar[d] &
\langle x_1, \ldots, x_l \rangle / \langle x_1 \rangle \ar[r] \ar[d] &
0 \\
0 \ar[r] &
\langle y_1 \rangle \ar[r] &
\langle y_1, \ldots, y_l \rangle \ar[r] &
\langle y_1, \ldots, y_l \rangle / \langle y_1 \rangle \ar[r] &
0
}
$$
et le lemme résulte du lemme du serpent et de la récurrence sur $l$.
Supposons maintenant que $i > 1$.
Écrivons $\varphi(x_1) = a_1 y_1 + \ldots + a_{i - 1} y_{i - 1} + a y_i$
avec $a_j, a \in R$ et $a \not \in \mathfrak q$ (sinon
$i$ ne serait pas minimal). Posons
$$
x'_j =
\left\{
\begin{matrix}
x_j & \text{si} & j = 1 \\
ax_j & \text{si} & j \geq 2
\end{matrix}
\right.
\quad\text{et}\quad
y'_j =
\left\{
\begin{matrix}
y_j & \text{si} & j < i \\
ay_j & \text{si} & j \geq i
\end{matrix}
\right.
$$
Soient $M' = \langle x'_1, \ldots, x'_l \rangle$ et
$N' = \langle y'_1, \ldots, y'_l \rangle$.
Puisque $\varphi(x'_1) = a_1 y'_1 + \ldots + a_{i - 1} y'_{i - 1} + y'_i$
par construction, et que pour $j > 1$ nous avons
$\varphi(x'_j) = a\varphi(x_i) \in \langle y'_1, \ldots, y'_l\rangle$,
nous obtenons un diagramme commutatif de $R$-modules et d'applications
$$
\xymatrix{
M' \ar[d] \ar[r]_{\varphi'} & N' \ar[d] \\
M \ar[r]^\varphi & N
}
$$
D'après le résultat du deuxième paragraphe de la démonstration, nous savons
que $\text{longueur}_R(M/M') = (l - 1)\text{ord}_{R/\mathfrak q}(a)$
et, de même,
$\text{longueur}_R(M/M') = (l - i + 1)\text{ord}_{R/\mathfrak q}(a)$.
Une chasse au diagramme montre alors que
$$
\text{longueur}_R(\Coker(\varphi')) =
\text{longueur}_R(\Coker(\varphi)) + i\ \text{ord}_{R/\mathfrak q}(a).
$$
D'autre part, il est clair qu'en écrivant
$$
[\varphi(x_1), \ldots, \varphi(x_l)] = f [y_1, \ldots, y_l],
\quad
[\varphi'(x'_1), \ldots, \varphi(x'_l)] = f' [y'_1, \ldots, y'_l]
$$
nous avons $f' = a^if$. Il suffit donc de démontrer le lemme dans le
cas où $\varphi(x_1) = a_1y_1 + \ldots a_{i - 1}y_{i - 1} + y_i$,
c'est-à-dire dans le cas où $a = 1$. Rappelons ensuite que
$$
[y_1, \ldots, y_l] = [y_1, \ldots, y_{i - 1},
a_1y_1 + \ldots a_{i - 1}y_{i - 1} + y_i, y_{i + 1}, \ldots, y_l]
$$
d'après les relations d'admissibilité pour les symboles. La suite
$y_1, \ldots, y_{i - 1},
a_1y_1 + \ldots + a_{i - 1}y_{i - 1} + y_i, y_{i + 1}, \ldots, y_l$
satisfait encore les conditions (3) et (4) du lemme.
Ainsi, nous pouvons en fait
supposer que $\varphi(x_1) = y_i$. Dans ce cas, remarquons que
$\mathfrak q x_1 = 0$, ce qui implique aussi $\mathfrak q y_i = 0$.
Nous avons
$$
[y_1, \ldots, y_l] =
- [y_1, \ldots, y_{i - 2}, y_i, y_{i - 1}, y_{i + 1}, \ldots, y_l]
$$
d'après la troisième des relations d'admissibilité qui définissent
$\det_{\kappa(\mathfrak q)}(N_{\mathfrak q})$. Nous pouvons donc
remplacer $y_1, \ldots, y_l$ par
la suite
$y'_1, \ldots, y'_l =
y_1, \ldots, y_{i - 2}, y_i, y_{i - 1}, y_{i + 1}, \ldots, y_l$
(qui satisfait elle aussi les conditions (3) et (4) du lemme).
Cela diminue manifestement l'invariant $i$ de $1$, et nous concluons par récurrence
sur $i$.
\end{proof}

\noindent
Pour utiliser le lemme précédent, montrons qu'il existe souvent des suites d'éléments
ayant les propriétés requises.

\begin{lemma}
\label{lemma-good-sequence-exists}
Soit $R$ un anneau local noethérien.
Soit $\mathfrak q \subset R$ un idéal premier.
Soit $M$ un $R$-module fini tel que
$\mathfrak q$ soit l'un des idéaux premiers minimaux du support de $M$.
Alors il existe $x_1, \ldots, x_l \in M$ tels que
\begin{enumerate}
\item le support de $M / \langle x_1, \ldots, x_l\rangle$ ne contienne pas
$\mathfrak q$ ; et
\item $\langle x_1, \ldots, x_i\rangle / \langle x_1, \ldots, x_{i - 1}\rangle
\cong R/\mathfrak q$ pour $i = 1, \ldots, l$.
\end{enumerate}
De plus, dans ce cas, $l = \text{longueur}_{R_\mathfrak q}(M_\mathfrak q)$.
\end{lemma}

\begin{proof}
La condition selon laquelle $\mathfrak q$ est un idéal premier minimal du support
de $M$ implique que $l = \text{longueur}_{R_\mathfrak q}(M_\mathfrak q)$
est fini (voir Algèbre, lemme \ref{algebra-lemma-support-point}).
Nous pouvons donc trouver $y_1, \ldots, y_l \in M_{\mathfrak q}$
tels que
$\langle y_1, \ldots, y_i\rangle / \langle y_1, \ldots, y_{i - 1}\rangle
\cong \kappa(\mathfrak q)$ pour $i = 1, \ldots, l$.
Nous pouvons trouver $f_i \in R$, $f_i \not \in \mathfrak q$, tels que
$f_i y_i$ soit l'image d'un certain élément $z_i \in M$.
De plus, puisque $R$ est noethérien, nous pouvons écrire
$\mathfrak q = (g_1, \ldots, g_t)$ pour certains $g_j \in R$.
Par hypothèse, $g_j y_i \in \langle y_1, \ldots, y_{i - 1} \rangle$
dans le module $M_{\mathfrak q}$.
Grâce à notre choix de $z_i$, nous pouvons trouver d'autres éléments
$f_{ji} \in R$, $f_{ij} \not \in \mathfrak q$, tels que
$f_{ij} g_j z_i \in \langle z_1, \ldots, z_{i - 1} \rangle$
(égalité dans le module $M$).
Le lemme résulte du choix
$$
x_1 = f_{11}f_{12}\ldots f_{1t}z_1,
\quad
x_2 = f_{11}f_{12}\ldots f_{1t}f_{21}f_{22}\ldots f_{2t}z_2,
$$
et ainsi de suite. En effet, puisque tous les éléments $f_i, f_{ij}$ sont inversibles
dans $R_{\mathfrak q}$, nous avons encore
$R_{\mathfrak q}x_1 + \ldots + R_{\mathfrak q}x_i /
R_{\mathfrak q}x_1 + \ldots + R_{\mathfrak q}x_{i - 1}
\cong \kappa(\mathfrak q)$ pour $i = 1, \ldots, l$.
Par construction, $\mathfrak q x_i \in \langle x_1, \ldots, x_{i - 1}\rangle$.
Ainsi, $\langle x_1, \ldots, x_i\rangle / \langle x_1, \ldots, x_{i - 1}\rangle$
est un $R$-module engendré par un élément et annulé par $\mathfrak q$,
tel que sa localisation en $\mathfrak q$ donne un espace vectoriel de dimension $q$
sur $\kappa(\mathfrak q)$.
Il est donc isomorphe à $R/\mathfrak q$.
\end{proof}

\noindent
Voici le résultat principal de cette section.
Nous verrons ci-dessous les différentes
conséquences de cette proposition.
Le lecteur est invité à démontrer d'abord lui-même le plus facile
lemme \ref{lemma-application-herbrand-quotient}.

\begin{proposition}
\label{proposition-length-determinant-periodic-complex}
Soit $R$ un anneau local noethérien de corps résiduel $\kappa$.
Supposons que $(M, \varphi, \psi)$ soit un complexe $(2, 1)$-périodique
sur $R$. Supposons que
\begin{enumerate}
\item $M$ soit un $R$-module fini ;
\item les modules de cohomologie de $(M, \varphi, \psi)$ soient de longueur finie ; et
\item $\dim(\text{Supp}(M)) = 1$.
\end{enumerate}
Soient $\mathfrak q_i$, $i = 1, \ldots, t$, les idéaux
premiers minimaux du support de $M$. Alors nous avons\footnote{
Évidemment, nous pourrions supprimer le signe moins en redéfinissant
$\det_\kappa(M, \varphi, \psi)$ comme l'inverse de sa
valeur actuelle ; voir la définition \ref{definition-periodic-determinant}.}
$$
- e_R(M, \varphi, \psi) =
\sum\nolimits_{i = 1, \ldots, t}
\text{ord}_{R/\mathfrak q_i}\left(
\det\nolimits_{\kappa(\mathfrak q_i)}
(M_{\mathfrak q_i}, \varphi_{\mathfrak q_i}, \psi_{\mathfrak q_i})
\right)
$$
\end{proposition}

\begin{proof}
Nous nous ramenons d'abord au cas $t = 1$ de la manière suivante.
Remarquons que
$\text{Supp}(M) = \{\mathfrak m, \mathfrak q_1, \ldots, \mathfrak q_t\}$,
où $\mathfrak m \subset R$ est l'idéal maximal.
Notons $M_i$ l'image de $M \to M_{\mathfrak q_i}$,
de sorte que $\text{Supp}(M_i) = \{\mathfrak m, \mathfrak q_i\}$.
L'application $\varphi$ (resp.\ $\psi$) induit une application de $R$-modules
$\varphi_i : M_i \to M_i$ (resp.\ $\psi_i : M_i \to M_i$).
Nous obtenons ainsi un morphisme de complexes $(2, 1)$-périodiques
$$
(M, \varphi, \psi) \longrightarrow
\bigoplus\nolimits_{i = 1, \ldots, t} (M_i, \varphi_i, \psi_i).
$$
Le noyau et le conoyau de cette application ont leur support contenu dans
$\{\mathfrak m\}$. Ainsi, d'après le lemme \ref{lemma-compare-periodic-lengths},
nous avons
$$
e_R(M, \varphi, \psi) =
\sum\nolimits_{i = 1, \ldots, t}
e_R(M_i, \varphi_i, \psi_i)
$$
D'autre part, nous avons clairement $M_{\mathfrak q_i} = M_{i, \mathfrak q_i}$,
et donc les termes du membre de droite de la formule du
lemme sont égaux aux expressions
$$
\text{ord}_{R/\mathfrak q_i}\left(
\det\nolimits_{\kappa(\mathfrak q_i)}
(M_{i, \mathfrak q_i}, \varphi_{i, \mathfrak q_i}, \psi_{i, \mathfrak q_i})
\right)
$$
Autrement dit, si nous pouvons démontrer le lemme pour chacun des modules
$M_i$, alors le lemme est vrai. Cela nous ramène au cas $t = 1$.

\medskip\noindent
Supposons donné un complexe $(2, 1)$-périodique $(M, \varphi, \psi)$
sur un anneau local noethérien, où $M$ est un $R$-module fini,
$\text{Supp}(M) = \{\mathfrak m, \mathfrak q\}$ et les modules de cohomologie
sont de longueur finie. Dans ce cas, la démonstration
résulte du lemme \ref{lemma-key-lemma} et d'une comptabilité soigneuse.
Notons
$K_\varphi = \Ker(\varphi)$,
$I_\varphi = \Im(\varphi)$,
$K_\psi = \Ker(\psi)$, et
$I_\psi = \Im(\psi)$.
Puisque $R$ est noethérien, ce sont tous des $R$-modules finis.
Posons
$$
a = \text{longueur}_{R_{\mathfrak q}}(I_{\varphi, \mathfrak q})
= \text{longueur}_{R_{\mathfrak q}}(K_{\psi, \mathfrak q}),
\quad
b = \text{longueur}_{R_{\mathfrak q}}(I_{\psi, \mathfrak q})
= \text{longueur}_{R_{\mathfrak q}}(K_{\varphi, \mathfrak q}).
$$
Ces égalités viennent de ce que le complexe devient exact après localisation en
$\mathfrak q$. Remarquons que $l = \text{longueur}_{R_{\mathfrak q}}(M_{\mathfrak q})$
est égal à $l = a + b$.

\medskip\noindent
Nous allons utiliser le lemme \ref{lemma-good-sequence-exists}
pour choisir des suites d'éléments dans des $R$-modules finis
$N$ dont le support est contenu dans $\{\mathfrak m, \mathfrak q\}$.
Dans ce cas, $N_{\mathfrak q}$ est de longueur finie, disons $n \in \mathbf{N}$.
Appelons « bonne suite » une suite $w_1, \ldots, w_n \in N$
ayant les propriétés (1) et (2) du lemme \ref{lemma-good-sequence-exists}.
Remarquons que le quotient
$N/\langle w_1, \ldots, w_n \rangle$ de $N$ par le sous-module engendré par
une bonne suite a son support (contenu) dans $\{\mathfrak m\}$
et est donc de longueur finie (Algèbre, lemme \ref{algebra-lemma-support-point}).
De plus, le symbole
$[w_1, \ldots, w_n] \in \det_{\kappa(\mathfrak q)}(N_{\mathfrak q})$
est un générateur ; voir le lemme \ref{lemma-determinant-dimension-one}.

\medskip\noindent
Cela étant dit, choisissons de bonnes suites
$$
\begin{matrix}
x_1, \ldots, x_b & \text{dans} & K_\varphi, &
t_1, \ldots, t_a & \text{dans} & K_\psi, \\
y_1, \ldots, y_a & \text{dans} & I_\varphi \cap \langle t_1, \ldots t_a\rangle, &
s_1, \ldots, s_b & \text{dans} & I_\psi \cap \langle x_1, \ldots, x_b\rangle.
\end{matrix}
$$
Nous allons modifier légèrement nos choix comme suit.
Choisissons des relèvements $\tilde y_i \in M$ de $y_i \in I_\varphi$
et $\tilde s_i \in M$ de $s_i \in I_\psi$. Il se peut
que $\mathfrak q \tilde y_1 \subset \langle x_1, \ldots, x_b\rangle$ ne soit pas vérifié,
et il se peut que
$\mathfrak q \tilde s_1 \subset \langle t_1, \ldots, t_a\rangle$ ne le soit pas non plus.
Cependant, en utilisant le fait que $\mathfrak q$ est de type fini (comme dans la démonstration
du lemme \ref{lemma-good-sequence-exists}), nous pouvons trouver
$d \in R$, $d \not \in \mathfrak q$, tel que
$\mathfrak q d\tilde y_1 \subset \langle x_1, \ldots, x_b\rangle$
et
$\mathfrak q d\tilde s_1 \subset \langle t_1, \ldots, t_a\rangle$.
Ainsi, après avoir remplacé $y_i$ par $dy_i$,
$\tilde y_i$ par $d\tilde y_i$, $s_i$ par $ds_i$ et $\tilde s_i$
par $d\tilde s_i$, nous voyons que nous pouvons aussi supposer que
$x_1, \ldots, x_b, \tilde y_1, \ldots, \tilde y_b$
et $t_1, \ldots, t_a, \tilde s_1, \ldots, \tilde s_b$
sont de bonnes suites dans $M$.

\medskip\noindent
Enfin, choisissons une bonne suite
$z_1, \ldots, z_l$ dans le $R$-module fini
$$
\langle
x_1, \ldots, x_b, \tilde y_1, \ldots, \tilde y_a
\rangle
\cap
\langle
t_1, \ldots, t_a, \tilde s_1, \ldots, \tilde s_b
\rangle.
$$
Remarquons que c'est aussi une bonne suite dans $M$.

\medskip\noindent
Puisque $I_{\varphi, \mathfrak q} = K_{\psi, \mathfrak q}$,
il existe un unique élément $h \in \kappa(\mathfrak q)$ tel que
$[y_1, \ldots, y_a] = h [t_1, \ldots, t_a]$
dans $\det_{\kappa(\mathfrak q)}(K_{\psi, \mathfrak q})$.
De même, comme $I_{\psi, \mathfrak q} = K_{\varphi, \mathfrak q}$,
il existe un unique élément $h \in \kappa(\mathfrak q)$ tel que
$[s_1, \ldots, s_b] = g [x_1, \ldots, x_b]$
dans $\det_{\kappa(\mathfrak q)}(K_{\varphi, \mathfrak q})$.
Nous pouvons également procéder ainsi avec les trois bonnes suites que nous avons
dans $M$. Au total, nous obtenons les identités suivantes :
\begin{align*}
[y_1, \ldots, y_a]
& =
h [t_1, \ldots, t_a] \\
[s_1, \ldots, s_b]
& =
g [x_1, \ldots, x_b] \\
[z_1, \ldots, z_l]
& =
f_\varphi [x_1, \ldots, x_b, \tilde y_1, \ldots, \tilde y_a] \\
[z_1, \ldots, z_l]
& =
f_\psi [t_1, \ldots, t_a, \tilde s_1, \ldots, \tilde s_b]
\end{align*}
pour certains $g, h, f_\varphi, f_\psi \in \kappa(\mathfrak q)$.

\medskip\noindent
Toutes ces
notations étant en place, calculons $\det_{\kappa(\mathfrak q)}(M, \varphi, \psi)$.
Considérons l'élément $[z_1, \ldots, z_l]$.
Sous l'application $\gamma_\psi \circ \sigma \circ \gamma_\varphi^{-1}$
de la définition \ref{definition-periodic-determinant}, nous avons
\begin{eqnarray*}
[z_1, \ldots, z_l] & = &
f_\varphi [x_1, \ldots, x_b, \tilde y_1, \ldots, \tilde y_a] \\
& \mapsto & f_\varphi [x_1, \ldots, x_b] \otimes [y_1, \ldots, y_a] \\
& \mapsto &
f_\varphi h/g [t_1, \ldots, t_a] \otimes [s_1, \ldots, s_b] \\
& \mapsto &
f_\varphi h/g [t_1, \ldots, t_a, \tilde s_1, \ldots, \tilde s_b] \\
& = &
f_\varphi h/f_\psi g [z_1, \ldots, z_l]
\end{eqnarray*}
Cela signifie que
$\det_{\kappa(\mathfrak q)}
(M_{\mathfrak q}, \varphi_{\mathfrak q}, \psi_{\mathfrak q})$
est égal à $f_\varphi h/f_\psi g$ au signe près.

\medskip\noindent
Abrégeons les quantités suivantes :
\begin{eqnarray*}
k_\varphi & = & \text{longueur}_R(K_\varphi/\langle x_1, \ldots, x_b\rangle) \\
k_\psi    & = & \text{longueur}_R(K_\psi/\langle t_1, \ldots, t_a\rangle) \\
i_\varphi & = & \text{longueur}_R(I_\varphi/\langle y_1, \ldots, y_a\rangle) \\
i_\psi    & = & \text{longueur}_R(I_\psi/\langle s_1, \ldots, s_a\rangle) \\
m_\varphi & = & \text{longueur}_R(M/
\langle x_1, \ldots, x_b, \tilde y_1, \ldots, \tilde y_a\rangle) \\
m_\psi    & = & \text{longueur}_R(M/
\langle t_1, \ldots, t_a, \tilde s_1, \ldots, \tilde s_b\rangle) \\
\delta_\varphi & = & \text{longueur}_R(
\langle x_1, \ldots, x_b, \tilde y_1, \ldots, \tilde y_a\rangle
\langle z_1, \ldots, z_l\rangle) \\
\delta_\psi & = & \text{longueur}_R(
\langle t_1, \ldots, t_a, \tilde s_1, \ldots, \tilde s_b\rangle
\langle z_1, \ldots, z_l\rangle)
\end{eqnarray*}
En utilisant les suites exactes $0 \to K_\varphi \to M \to I_\varphi \to 0$,
nous obtenons $m_\varphi = k_\varphi + i_\varphi$. De même, nous avons
$m_\psi = k_\psi + i_\psi$. Nous avons
$\delta_\varphi + m_\varphi = \delta_\psi + m_\psi$, puisque cette quantité
est égale à la colongueur de $\langle z_1, \ldots, z_l \rangle$
dans $M$. Enfin, nous avons
$$
\delta_\varphi = \text{ord}_{R/\mathfrak q}(f_\varphi),
\quad
\delta_\psi = \text{ord}_{R/\mathfrak q}(f_\psi)
$$
d'après notre première application du lemme clé \ref{lemma-key-lemma}.

\medskip\noindent
Calculons ensuite la multiplicité du complexe périodique :
\begin{eqnarray*}
e_R(M, \varphi, \psi) & = &
\text{longueur}_R(K_\varphi/I_\psi) - \text{longueur}_R(K_\psi/I_\varphi) \\
& = &
\text{longueur}_R(
\langle x_1, \ldots, x_b\rangle/
\langle s_1, \ldots, s_b\rangle)
+ k_\varphi - i_\psi \\
& & -
\text{longueur}_R(
\langle t_1, \ldots, t_a\rangle/
\langle y_1, \ldots, y_a\rangle)
- k_\psi + i_\varphi \\
& = &
\text{ord}_{R/\mathfrak q}(g/h) + k_\varphi - i_\psi - k_\psi + i_\varphi \\
& = &
\text{ord}_{R/\mathfrak q}(g/h) + m_\varphi - m_\psi \\
& = &
\text{ord}_{R/\mathfrak q}(g/h) + \delta_\psi - \delta_\varphi \\
& = &
\text{ord}_{R/\mathfrak q}(f_\psi g/f_\varphi h)
\end{eqnarray*}
où nous avons utilisé deux fois le lemme clé \ref{lemma-key-lemma} dans la troisième égalité.
D'après notre calcul de $\det_{\kappa(\mathfrak q)}
(M_{\mathfrak q}, \varphi_{\mathfrak q}, \psi_{\mathfrak q})$,
cela démontre la proposition.
\end{proof}

\noindent
Dans la plupart des applications, le lemme suivant suffit.

\begin{lemma}
\label{lemma-application-herbrand-quotient}
Soit $R$ un anneau local noethérien d'idéal maximal $\mathfrak m$.
Soient $M$ un $R$-module fini et $\psi : M \to M$ une
application de $R$-modules. Supposons que
\begin{enumerate}
\item $\Ker(\psi)$ et $\Coker(\psi)$ soient de longueur finie ; et
\item $\dim(\text{Supp}(M)) \leq 1$.
\end{enumerate}
Écrivons
$\text{Supp}(M) = \{\mathfrak m, \mathfrak q_1, \ldots, \mathfrak q_t\}$
et notons $f_i \in \kappa(\mathfrak q_i)^*$ l'élément tel que
$\det_{\kappa(\mathfrak q_i)}(\psi_{\mathfrak q_i}) :
\det_{\kappa(\mathfrak q_i)}(M_{\mathfrak q_i})
\to \det_{\kappa(\mathfrak q_i)}(M_{\mathfrak q_i})$
soit la multiplication par $f_i$. Alors
nous avons
$$
\text{longueur}_R(\Coker(\psi))
-
\text{longueur}_R(\Ker(\psi))
=
\sum\nolimits_{i = 1, \ldots, t}
\text{ord}_{R/\mathfrak q_i}(f_i).
$$
\end{lemma}

\begin{proof}
Rappelons que $H^0(M, 0, \psi) = \Coker(\psi)$ et
$H^1(M, 0, \psi) = \Ker(\psi)$ ; voir les remarques précédant
la définition \ref{definition-periodic-length}.
Le lemme résulte de la combinaison de la
proposition \ref{proposition-length-determinant-periodic-complex} avec le
lemme \ref{lemma-periodic-determinant-easy-case}.

\medskip\noindent
Autre démonstration. Ramenons-nous au cas
$\text{Supp}(M) = \{\mathfrak m, \mathfrak q\}$
comme dans la démonstration de la
proposition \ref{proposition-length-determinant-periodic-complex}.
Combinons alors directement les
lemmes \ref{lemma-key-lemma} et
\ref{lemma-good-sequence-exists}
pour démontrer ce cas particulier de la
proposition \ref{proposition-length-determinant-periodic-complex}.
La comptabilité est bien moindre dans ce cas, et le lecteur est
invité à la mener. Les détails sont omis.
\end{proof}




\subsection{Application au lemme clé}
\label{subsection-application-tame-symbol}

\noindent
Dans cette section, nous appliquons les résultats ci-dessus pour montrer
l'analogue du lemme clé (lemme \ref{lemma-milnor-gersten-low-degree})
avec le symbole modéré $d_A$ construit ci-dessus. Voir la
remarque \ref{remark-gersten-complex-milnor}
pour la relation avec la $K$-théorie de Milnor.

\begin{lemma}[Lemme clé]
\label{lemma-secondary-ramification}
\begin{reference}
Lorsque $A$ est un anneau excellent, c'est \cite[Proposition 1]{Kato-Milnor-K}.
\end{reference}
Soit $A$ un anneau intègre local noethérien de dimension $2$ et de corps des fractions $K$.
Soient $f, g \in K^*$.
Soient $\mathfrak q_1, \ldots, \mathfrak q_t$ les idéaux premiers de hauteur
$1$, notés $\mathfrak q$, de $A$, tels que $f$ ou $g$ n'appartienne pas à
$A^*_{\mathfrak q}$.
Alors nous avons
$$
\sum\nolimits_{i = 1, \ldots, t}
\text{ord}_{A/\mathfrak q_i}(d_{A_{\mathfrak q_i}}(f, g))
=
0
$$
Nous pouvons aussi écrire ceci sous la forme
$$
\sum\nolimits_{\text{hauteur}(\mathfrak q) = 1}
\text{ord}_{A/\mathfrak q}(d_{A_{\mathfrak q}}(f, g))
=
0
$$
puisque, pour tout idéal premier $\mathfrak q$ de hauteur un
de $A$ tel que $f, g \in A^*_{\mathfrak q}$,
nous avons $d_{A_{\mathfrak q}}(f, g) = 1$ d'après le
lemme \ref{lemma-symbol-when-one-is-a-unit}.
\end{lemma}

\begin{proof}
Puisque les symboles modérés $d_{A_{\mathfrak q}}(f, g)$ sont additifs
(lemme \ref{lemma-multiplicativity-symbol}) et que les fonctions d'ordre
$\text{ord}_{A/\mathfrak q}$
sont additives (Algèbre, lemme \ref{algebra-lemma-ord-additive}),
il suffit de démontrer la formule lorsque $f = a \in A$ et
$g = b \in A$. Dans ce cas, nous voyons qu'il faut montrer
$$
\sum\nolimits_{\text{hauteur}(\mathfrak q) = 1}
\text{ord}_{A/\mathfrak q}(\det\nolimits_\kappa(A_{\mathfrak q}/(ab), a, b))
= 0
$$
D'après la proposition \ref{proposition-length-determinant-periodic-complex},
cela équivaut à montrer que
$$
e_A(A/(ab), a, b) = 0.
$$
Puisque le complexe
$A/(ab) \xrightarrow{a} A/(ab) \xrightarrow{b} A/(ab) \xrightarrow{a} A/(ab)$
est exact, nous avons terminé.
\end{proof}





\section{Annexe B : Approches alternatives}
\label{section-appendix-chow}

\noindent
Dans cette annexe, nous tentons d'abord de relier brièvement les résultats
du texte principal à la $K$-théorie des faisceaux cohérents. En
particulier, nous décrivons comment le produit par $c_1$ d'un
module inversible est lié au produit tensoriel par ce module inversible ; voir le
lemme \ref{lemma-coherent-sheaf-cap-c1}.
Ces considérations sont manifestement très intéressantes et
méritent une présentation beaucoup plus détaillée et développée.


\subsection{Équivalence rationnelle et groupes de K-théorie}
\label{subsection-rational-equivalence-K-groups}

\noindent
Cette section prolonge la section \ref{section-chow-and-K}.
Le lemme suivant est motivé par
Homologie, lemme \ref{homology-lemma-serre-subcategory-K-groups}.

\begin{lemma}
\label{lemma-maps-between-coherent-sheaves}
Soit $(S, \delta)$ comme dans la situation \ref{situation-setup}.
Soit $X$ un schéma localement de type fini sur $S$.
Soit $\mathcal{F}$ un faisceau cohérent sur $X$.
Soit
$$
\xymatrix{
\ldots \ar[r] &
\mathcal{F} \ar[r]^\varphi &
\mathcal{F} \ar[r]^\psi &
\mathcal{F} \ar[r]^\varphi &
\mathcal{F} \ar[r] & \ldots
}
$$
un complexe comme dans Homologie, équation (\ref{homology-equation-cyclic-complex}).
Supposons que
\begin{enumerate}
\item $\dim_\delta(\text{Supp}(\mathcal{F})) \leq k + 1$ ;
\item $\dim_\delta(\text{Supp}(H^i(\mathcal{F}, \varphi, \psi))) \leq k$
pour $i = 0, 1$.
\end{enumerate}
Alors nous avons
$$
[H^0(\mathcal{F}, \varphi, \psi)]_k
\sim_{rat}
[H^1(\mathcal{F}, \varphi, \psi)]_k
$$
comme $k$-cycles sur $X$.
\end{lemma}

\begin{proof}
Soit $\{W_j\}_{j \in J}$ la famille des composantes irréductibles
de $\text{Supp}(\mathcal{F})$
de $\delta$-dimension $k + 1$. Remarquons que $\{W_j\}$
est une famille localement finie de fermés de
$X$, d'après le lemme \ref{lemma-length-finite}.
Pour chaque $j$, soit $\xi_j \in W_j$ le point générique.
Posons
$$
f_j = \det\nolimits_{\kappa(\xi_j)}
(\mathcal{F}_{\xi_j}, \varphi_{\xi_j}, \psi_{\xi_j})
\in
R(W_j)^*.
$$
Voir la définition \ref{definition-periodic-determinant} pour les notations.
Nous affirmons que
$$
- [H^0(\mathcal{F}, \varphi, \psi)]_k + [H^1(\mathcal{F}, \varphi, \psi)]_k
=
\sum (W_j \to X)_*\text{div}(f_j)
$$
Si nous démontrons ceci, le lemme en résulte.

\medskip\noindent
Soit $Z \subset X$ un sous-schéma fermé intègre de $\delta$-dimension $k$.
Pour démontrer l'égalité ci-dessus, il suffit de montrer que le coefficient $n$
de $[Z]$ dans
$
[H^0(\mathcal{F}, \varphi, \psi)]_k - [H^1(\mathcal{F}, \varphi, \psi)]_k
$
est égal au coefficient $m$ de $[Z]$ dans
$
\sum (W_j \to X)_*\text{div}(f_j)
$.
Soit $\xi \in Z$ le point générique.
Considérons l'anneau local $A = \mathcal{O}_{X, \xi}$.
Posons $M = \mathcal{F}_\xi$, considéré comme un $A$-module.
Notons $\varphi, \psi : M \to M$ les applications induites par $\varphi, \psi$ sur
le germe.
Par notre choix de $\xi \in Z$, nous avons $\delta(\xi) = k$,
et donc $\dim(\text{Supp}(M)) = 1$.
Enfin, les sous-schémas fermés intègres
$W_j$ passant par $\xi$ correspondent aux idéaux premiers minimaux
$\mathfrak q_i$ de $\text{Supp}(M)$.
Dans chaque cas, l'élément $f_j \in R(W_j)^*$ correspond à
l'élément $\det_{\kappa(\mathfrak q_i)}(M_{\mathfrak q_i}, \varphi, \psi)$
de $\kappa(\mathfrak q_i)^*$. Nous voyons donc que
$$
n = - e_A(M, \varphi, \psi)
$$
et
$$
m =
\sum
\text{ord}_{A/\mathfrak q_i}
(\det\nolimits_{\kappa(\mathfrak q_i)}(M_{\mathfrak q_i}, \varphi, \psi))
$$
Le résultat découle ainsi de la
proposition \ref{proposition-length-determinant-periodic-complex}.
\end{proof}

\begin{lemma}
\label{lemma-cycles-rational-equivalence-K-group}
Soit $(S, \delta)$ comme dans la situation \ref{situation-setup}.
Soit $X$ un schéma localement de type fini sur $S$.
L'application
$$
\CH_k(X) \longrightarrow
K_0(\textit{Coh}_{\leq k + 1}(X)/\textit{Coh}_{\leq k - 1}(X))
$$
du lemme \ref{lemma-from-chow-to-K} induit une bijection de
$\CH_k(X)$ sur l'image $B_k(X)$ de l'application
$$
K_0(\textit{Coh}_{\leq k}(X)/\textit{Coh}_{\leq k - 1}(X))
\longrightarrow
K_0(\textit{Coh}_{\leq k + 1}(X)/\textit{Coh}_{\leq k - 1}(X)).
$$
\end{lemma}

\begin{proof}
D'après le lemme \ref{lemma-cycles-k-group}, nous avons
$Z_k(X) = K_0(\textit{Coh}_{\leq k}(X)/\textit{Coh}_{\leq k - 1}(X))$,
de manière compatible avec l'application du lemme \ref{lemma-from-chow-to-K}.
Supposons donc donné un élément $[A] - [B]$ de
$K_0(\textit{Coh}_{\leq k}(X)/\textit{Coh}_{\leq k - 1}(X))$
qui s'envoie sur zéro dans $B_k(X)$, c'est-à-dire sur zéro dans
$K_0(\textit{Coh}_{\leq k + 1}(X)/\textit{Coh}_{\leq k - 1}(X))$.
Nous devons montrer que $[A] - [B]$ correspond à un cycle
rationnellement équivalent à zéro sur $X$.
Supposons que $[A] = [\mathcal{A}]$ et $[B] = [\mathcal{B}]$
pour certains faisceaux cohérents $\mathcal{A}, \mathcal{B}$ sur
$X$, à support de $\delta$-dimension $\leq k$.
L'hypothèse selon laquelle $[A] - [B]$ s'envoie sur zéro dans le groupe
$K_0(\textit{Coh}_{\leq k + 1}(X)/\textit{Coh}_{\leq k - 1}(X))$
signifie qu'il existe des faisceaux cohérents
$\mathcal{A}', \mathcal{B}'$ sur $X$, à support de
$\delta$-dimension $\leq k - 1$, tels que
$[\mathcal{A} \oplus \mathcal{A}'] - [\mathcal{B} \oplus \mathcal{B}']$
soit nul dans $K_0(\textit{Coh}_{k + 1}(X))$ (utiliser la partie (1) de
Homologie, lemme \ref{homology-lemma-serre-subcategory-K-groups}).
D'après la partie (2) de
Homologie, lemme \ref{homology-lemma-serre-subcategory-K-groups},
cela signifie qu'il existe un complexe $(2, 1)$-périodique
$(\mathcal{F}, \varphi, \psi)$ dans la catégorie $\textit{Coh}_{\leq k + 1}(X)$
tel que
$\mathcal{A} \oplus \mathcal{A}' = H^0(\mathcal{F}, \varphi, \psi)$
et $\mathcal{B} \oplus \mathcal{B}' = H^1(\mathcal{F}, \varphi, \psi)$.
D'après le lemme \ref{lemma-maps-between-coherent-sheaves},
cela implique que
$$
[\mathcal{A} \oplus \mathcal{A}']_k
\sim_{rat}
[\mathcal{B} \oplus \mathcal{B}']_k
$$
Cela montre que $[A] - [B]$ s'envoie sur un cycle rationnellement
équivalent à zéro par l'application
$$
K_0(\textit{Coh}_{\leq k}(X)/\textit{Coh}_{\leq k - 1}(X))
\longrightarrow Z_k(X)
$$
du lemme \ref{lemma-cycles-k-group}. C'est ce que nous
devions démontrer, et la démonstration est achevée.
\end{proof}














\subsection{Diviseurs de Cartier et groupes de K-théorie}
\label{subsection-cartier-coherent}

\noindent
Dans cette section, nous décrivons comment l'intersection avec la
première classe de Chern d'un faisceau inversible $\mathcal{L}$
correspond au produit tensoriel par $\mathcal{L} - \mathcal{O}$
dans les groupes de $K$-théorie.

\begin{lemma}
\label{lemma-no-embedded-points-modules}
Soit $A$ un anneau local noethérien.
Soit $M$ un $A$-module fini.
Soient $a, b \in A$.
Supposons que
\begin{enumerate}
\item $\dim(A) = 1$ ;
\item $a$ et $b$ soient tous deux des non-diviseurs de zéro dans $A$ ;
\item $A$ n'ait pas d'idéaux premiers immergés ;
\item $M$ n'ait pas d'idéaux premiers associés immergés ;
\item $\text{Supp}(M) = \Spec(A)$.
\end{enumerate}
Soit $I = \{x \in A \mid x(a/b) \in A\}$.
Soient $\mathfrak q_1, \ldots, \mathfrak q_t$ les idéaux premiers minimaux
de $A$. Alors $(a/b)IM \subset M$ et
$$
\text{longueur}_A(M/(a/b)IM)
-
\text{longueur}_A(M/IM)
=
\sum\nolimits_i
\text{longueur}_{A_{\mathfrak q_i}}(M_{\mathfrak q_i})
\text{ord}_{A/\mathfrak q_i}(a/b)
$$
\end{lemma}

\begin{proof}
Puisque $M$ n'a pas d'idéaux premiers associés immergés et que
le support de $M$ est $\Spec(A)$, nous voyons que
$\text{Ass}(M) = \{\mathfrak q_1, \ldots, \mathfrak q_t\}$.
Ainsi, $a$ et $b$ sont des non-diviseurs de zéro dans $M$. Remarquons que
\begin{align*}
& \text{longueur}_A(M/(a/b)IM) \\
& = \text{longueur}_A(bM/aIM) \\
& = \text{longueur}_A(M/aIM)
-
\text{longueur}_A(M/bM) \\
& = \text{longueur}_A(M/aM) + \text{longueur}_A(aM/aIM) - \text{longueur}_A(M/bM) \\
& = \text{longueur}_A(M/aM) + \text{longueur}_A(M/IM) - \text{longueur}_A(M/bM)
\end{align*}
car l'application injective $b : M \to bM$ envoie $(a/b)IM$ sur $aIM$,
et l'application injective $a : M \to aM$ envoie $IM$ sur $aIM$.
Le membre de gauche de l'équation du lemme est donc
égal à
$$
\text{longueur}_A(M/aM) - \text{longueur}_A(M/bM).
$$
En appliquant la deuxième formule du
lemme \ref{lemma-additivity-divisors-restricted}
respectivement avec $x = a, b$,
et en utilisant Algèbre, définition \ref{algebra-definition-ord},
des fonctions $\text{ord}$, nous obtenons le résultat.
\end{proof}

\begin{lemma}
\label{lemma-no-embedded-points}
Soit $(S, \delta)$ comme dans la situation \ref{situation-setup}.
Soit $X$ un schéma localement de type fini sur $S$.
Soit $\mathcal{L}$ un $\mathcal{O}_X$-module inversible.
Soit $\mathcal{F}$ un $\mathcal{O}_X$-module cohérent.
Soit $s \in \Gamma(X, \mathcal{K}_X(\mathcal{L}))$ une
section méromorphe de $\mathcal{L}$.
Supposons que
\begin{enumerate}
\item $\dim_\delta(X) \leq k + 1$ ;
\item $X$ n'ait pas de points immergés ;
\item $\mathcal{F}$ n'ait pas de points associés immergés ;
\item le support de $\mathcal{F}$ soit $X$ ; et
\item la section $s$ soit méromorphe régulière.
\end{enumerate}
Dans cette situation, soit $\mathcal{I} \subset \mathcal{O}_X$
l'idéal des dénominateurs de $s$ ; voir
Diviseurs,
définition \ref{divisors-definition-regular-meromorphic-ideal-denominators}.
Nous avons alors ce qui suit :
\begin{enumerate}
\item il existe des suites exactes courtes
$$
\begin{matrix}
0 &
\to &
\mathcal{I}\mathcal{F} &
\xrightarrow{1} &
\mathcal{F} &
\to &
\mathcal{Q}_1 &
\to &
0 \\
0 &
\to &
\mathcal{I}\mathcal{F} &
\xrightarrow{s} &
\mathcal{F} \otimes_{\mathcal{O}_X} \mathcal{L} &
\to &
\mathcal{Q}_2 &
\to &
0
\end{matrix}
$$
\item les faisceaux cohérents $\mathcal{Q}_1$, $\mathcal{Q}_2$
sont à support de $\delta$-dimension $\leq k$ ;
\item la section $s$ se restreint en une section méromorphe régulière
$s_i$ sur toute composante irréductible $X_i$ de
$X$ de $\delta$-dimension $k + 1$ ; et
\item en écrivant $[\mathcal{F}]_{k + 1} = \sum m_i[X_i]$, nous avons
$$
[\mathcal{Q}_2]_k - [\mathcal{Q}_1]_k
=
\sum m_i(X_i \to X)_*\text{div}_{\mathcal{L}|_{X_i}}(s_i)
$$
dans $Z_k(X)$ ; en particulier,
$$
[\mathcal{Q}_2]_k - [\mathcal{Q}_1]_k
=
c_1(\mathcal{L}) \cap [\mathcal{F}]_{k + 1}
$$
dans $\CH_k(X)$.
\end{enumerate}
\end{lemma}

\begin{proof}
Rappelons, d'après Diviseurs, lemme \ref{divisors-lemma-make-maps-regular-section},
l'existence d'applications injectives
$1 : \mathcal{I}\mathcal{F} \to \mathcal{F}$ et
$s : \mathcal{I}\mathcal{F} \to \mathcal{F} \otimes_{\mathcal{O}_X}\mathcal{L}$
dont les conoyaux sont supportés par un fermé partout non dense $T$.
Notons $\mathcal{Q}_i$ leurs conoyaux comme dans le lemme.
Nous concluons que $\dim_\delta(\text{Supp}(\mathcal{Q}_i)) \leq k$.
D'après Diviseurs, lemmes \ref{divisors-lemma-pullback-meromorphic-sections-defined}
et \ref{divisors-lemma-meromorphic-sections-pullback}, les images réciproques $s_i$
sont définies et sont des sections méromorphes régulières de $\mathcal{L}|_{X_i}$.
L'égalité des cycles dans (4) implique l'égalité des classes de cycles
dans (4). Il reste donc seulement à montrer que
$$
[\mathcal{Q}_2]_k - [\mathcal{Q}_1]_k
=
\sum m_i(X_i \to X)_*\text{div}_{\mathcal{L}|_{X_i}}(s_i)
$$
est vraie dans $Z_k(X)$. Pour le voir, soit $Z \subset X$ un sous-schéma fermé intègre
de $\delta$-dimension $k$. Soit $\xi \in Z$ le point générique.
Soient $A = \mathcal{O}_{X, \xi}$ et $M = \mathcal{F}_\xi$.
De plus, choisissons un générateur $s_\xi \in \mathcal{L}_\xi$.
Nous pouvons alors écrire $s = (a/b) s_\xi$, où $a, b \in A$ sont des
non-diviseurs de zéro. Dans ce cas,
$I = \mathcal{I}_\xi = \{x \in A \mid x(a/b) \in A\}$.
Dans ce cas, le coefficient de $[Z]$ dans le membre de gauche est
$$
\text{longueur}_A(M/(a/b)IM) - \text{longueur}_A(M/IM)
$$
et le coefficient de $[Z]$ dans le membre de droite
est
$$
\sum
\text{longueur}_{A_{\mathfrak q_i}}(M_{\mathfrak q_i})
\text{ord}_{A/\mathfrak q_i}(a/b)
$$
où $\mathfrak q_1, \ldots, \mathfrak q_t$ sont les idéaux premiers minimaux
de l'anneau local de dimension $1$ noté $A$. Le résultat
découle donc du lemme \ref{lemma-no-embedded-points-modules}.
\end{proof}

\begin{lemma}
\label{lemma-coherent-sheaf-cap-c1}
Soit $(S, \delta)$ comme dans la situation \ref{situation-setup}.
Soit $X$ un schéma localement de type fini sur $S$.
Soit $\mathcal{L}$ un $\mathcal{O}_X$-module inversible.
Soit $\mathcal{F}$ un $\mathcal{O}_X$-module cohérent.
Supposons $\dim_\delta(\text{Supp}(\mathcal{F})) \leq k + 1$.
Alors l'élément
$$
[\mathcal{F} \otimes_{\mathcal{O}_X} \mathcal{L}]
-
[\mathcal{F}]
\in
K_0(\textit{Coh}_{\leq k + 1}(X)/\textit{Coh}_{\leq k - 1}(X))
$$
appartient au sous-groupe $B_k(X)$ du
lemme \ref{lemma-cycles-rational-equivalence-K-group} et s'envoie sur
l'élément $c_1(\mathcal{L}) \cap [\mathcal{F}]_{k + 1}$ par
l'application $B_k(X) \to \CH_k(X)$.
\end{lemma}

\begin{proof}
Soit
$$
0 \to \mathcal{K} \to \mathcal{F} \to \mathcal{F}' \to 0
$$
la suite exacte courte construite dans
Diviseurs, lemme \ref{divisors-lemma-remove-embedded-points}.
Cela signifie en particulier que $\mathcal{F}'$ n'a pas de points
associés immergés.
Puisque le support de $\mathcal{K}$ est partout non dense dans le
support de $\mathcal{F}$, nous voyons que
$\dim_\delta(\text{Supp}(\mathcal{K})) \leq k$. Nous pouvons
appliquer de nouveau
Diviseurs, lemme \ref{divisors-lemma-remove-embedded-points},
en partant de $\mathcal{K}$, pour obtenir une suite exacte courte
$$
0 \to \mathcal{K}'' \to \mathcal{K} \to \mathcal{K}' \to 0
$$
où maintenant $\dim_\delta(\text{Supp}(\mathcal{K}'')) < k$
et où $\mathcal{K}'$ n'a pas de points associés immergés.
Supposons que nous puissions démontrer le lemme pour les faisceaux cohérents
$\mathcal{F}'$ et $\mathcal{K}'$. Alors, d'après
les égalités
$$
[\mathcal{F}]_{k + 1}
=
[\mathcal{F}']_{k + 1}
+ [\mathcal{K}']_{k + 1}
+ [\mathcal{K}'']_{k + 1}
$$
(utiliser le lemme \ref{lemma-additivity-sheaf-cycle}),
$$
[\mathcal{F} \otimes_{\mathcal{O}_X} \mathcal{L}]
-
[\mathcal{F}]
=
[\mathcal{F}' \otimes_{\mathcal{O}_X} \mathcal{L}]
-
[\mathcal{F}']
+
[\mathcal{K}' \otimes_{\mathcal{O}_X} \mathcal{L}]
-
[\mathcal{K}']
+
[\mathcal{K}'' \otimes_{\mathcal{O}_X} \mathcal{L}]
-
[\mathcal{K}'']
$$
(utiliser le fait que $\otimes \mathcal{L}$ est exact)
et l'annulation triviale de $[\mathcal{K}'']_{k + 1}$ et de
$[\mathcal{K}'' \otimes_{\mathcal{O}_X} \mathcal{L}]
- [\mathcal{K}'']$ dans
$K_0(\textit{Coh}_{\leq k + 1}(X)/\textit{Coh}_{\leq k - 1}(X))$,
le résultat est vrai
pour $\mathcal{F}$. Cela signifie que nous pouvons supposer que
le faisceau $\mathcal{F}$ n'a pas de points associés immergés.

\medskip\noindent
Supposons que $X$ et $\mathcal{F}$ soient comme dans le lemme, et supposons de plus
que $\mathcal{F}$ n'a pas de points associés immergés. Considérons le
faisceau d'idéaux $\mathcal{I} \subset \mathcal{O}_X$, le sous-schéma fermé
correspondant $i : Z \to X$ et le $\mathcal{O}_Z$-module cohérent
$\mathcal{G}$ construits dans
Diviseurs, lemme \ref{divisors-lemma-no-embedded-points-endos}.
Rappelons que $Z$ est un schéma localement noethérien sans points immergés,
que $\mathcal{G}$ est un faisceau cohérent sans points
associés immergés, que $\text{Supp}(\mathcal{G}) = Z$,
et que $i_*\mathcal{G} = \mathcal{F}$.
De plus, posons $\mathcal{N} = \mathcal{L}|_Z$.

\medskip\noindent
D'après Diviseurs, lemme \ref{divisors-lemma-regular-meromorphic-section-exists},
le faisceau inversible $\mathcal{N}$ possède une section méromorphe régulière $s$
sur $Z$. Notons $\mathcal{J} \subset \mathcal{O}_Z$ le faisceau
des dénominateurs de $s$. D'après le lemme \ref{lemma-no-embedded-points},
il existe des suites exactes courtes
$$
\begin{matrix}
0 &
\to &
\mathcal{J}\mathcal{G} &
\xrightarrow{1} &
\mathcal{G} &
\to &
\mathcal{Q}_1 &
\to &
0 \\
0 &
\to &
\mathcal{J}\mathcal{G} &
\xrightarrow{s} &
\mathcal{G} \otimes_{\mathcal{O}_Z} \mathcal{N} &
\to &
\mathcal{Q}_2 &
\to &
0
\end{matrix}
$$
telles que $\dim_\delta(\text{Supp}(\mathcal{Q}_i)) \leq k$ et
telles que le cycle
$
[\mathcal{Q}_2]_k - [\mathcal{Q}_1]_k
$
soit un représentant de $c_1(\mathcal{N}) \cap [\mathcal{G}]_{k + 1}$.
Nous voyons (en utilisant le fait que
$i_*(\mathcal{G} \otimes \mathcal{N}) = \mathcal{F} \otimes \mathcal{L}$
par la formule de projection ; voir
Cohomologie, lemme \ref{cohomology-lemma-projection-formula})
que
$$
[\mathcal{F} \otimes_{\mathcal{O}_X} \mathcal{L}]
-
[\mathcal{F}]
=
[i_*\mathcal{Q}_2] - [i_*\mathcal{Q}_1]
$$
dans $K_0(\textit{Coh}_{\leq k + 1}(X)/\textit{Coh}_{\leq k - 1}(X))$.
Cela montre déjà que
$[\mathcal{F} \otimes_{\mathcal{O}_X} \mathcal{L}] - [\mathcal{F}]$
est un élément de $B_k(X)$. De plus, nous avons
\begin{eqnarray*}
[i_*\mathcal{Q}_2]_k - [i_*\mathcal{Q}_1]_k
& = &
i_*\left( [\mathcal{Q}_2]_k - [\mathcal{Q}_1]_k \right) \\
& = &
i_*\left(c_1(\mathcal{N}) \cap [\mathcal{G}]_{k + 1} \right) \\
& = &
c_1(\mathcal{L}) \cap i_*[\mathcal{G}]_{k + 1} \\
& = &
c_1(\mathcal{L}) \cap [\mathcal{F}]_{k + 1}
\end{eqnarray*}
d'après ce qui précède et les lemmes \ref{lemma-pushforward-cap-c1}
et \ref{lemma-cycle-push-sheaf}. Et cela coïncide avec
l'image de l'élément par $B_k(X) \to \CH_k(X)$, par définition.
Le lemme est donc démontré.
\end{proof}


































\begin{multicols}{2}[\section{Autres chapitres}]
\noindent
Préliminaires
\begin{enumerate}
\item \hyperref[introduction-section-phantom]{Introduction}
\item \hyperref[conventions-section-phantom]{Conventions}
\item \hyperref[sets-section-phantom]{Théorie des ensembles}
\item \hyperref[categories-section-phantom]{Catégories}
\item \hyperref[topology-section-phantom]{Topologie}
\item \hyperref[sheaves-section-phantom]{Faisceaux sur les espaces}
\item \hyperref[sites-section-phantom]{Sites et faisceaux}
\item \hyperref[stacks-section-phantom]{Champs}
\item \hyperref[fields-section-phantom]{Corps}
\item \hyperref[algebra-section-phantom]{Algèbre commutative}
\item \hyperref[brauer-section-phantom]{Groupes de Brauer}
\item \hyperref[homology-section-phantom]{Algèbre homologique}
\item \hyperref[derived-section-phantom]{Catégories dérivées}
\item \hyperref[simplicial-section-phantom]{Méthodes simpliciales}
\item \hyperref[more-algebra-section-phantom]{Compléments d'algèbre}
\item \hyperref[smoothing-section-phantom]{Lissage des morphismes d'anneaux}
\item \hyperref[modules-section-phantom]{Faisceaux de modules}
\item \hyperref[sites-modules-section-phantom]{Modules sur les sites}
\item \hyperref[injectives-section-phantom]{Objets injectifs}
\item \hyperref[cohomology-section-phantom]{Cohomologie des faisceaux}
\item \hyperref[sites-cohomology-section-phantom]{Cohomologie sur les sites}
\item \hyperref[dga-section-phantom]{Algèbre différentielle graduée}
\item \hyperref[dpa-section-phantom]{Algèbre à puissances divisées}
\item \hyperref[sdga-section-phantom]{Faisceaux différentiels gradués}
\item \hyperref[hypercovering-section-phantom]{Hyperrecouvrements}
\end{enumerate}
Schémas
\begin{enumerate}
\setcounter{enumi}{25}
\item \hyperref[schemes-section-phantom]{Schémas}
\item \hyperref[constructions-section-phantom]{Constructions de schémas}
\item \hyperref[properties-section-phantom]{Propriétés des schémas}
\item \hyperref[morphisms-section-phantom]{Morphismes de schémas}
\item \hyperref[coherent-section-phantom]{Cohomologie des schémas}
\item \hyperref[divisors-section-phantom]{Diviseurs}
\item \hyperref[limits-section-phantom]{Limites de schémas}
\item \hyperref[varieties-section-phantom]{Variétés}
\item \hyperref[topologies-section-phantom]{Topologies sur les schémas}
\item \hyperref[descent-section-phantom]{Descente}
\item \hyperref[perfect-section-phantom]{Catégories dérivées de schémas}
\item \hyperref[more-morphisms-section-phantom]{Compléments sur les morphismes}
\item \hyperref[flat-section-phantom]{Compléments sur la platitude}
\item \hyperref[groupoids-section-phantom]{Schémas en groupoïdes}
\item \hyperref[more-groupoids-section-phantom]{Compléments sur les schémas en groupoïdes}
\item \hyperref[etale-section-phantom]{Morphismes étales de schémas}
\end{enumerate}
Thèmes de théorie des schémas
\begin{enumerate}
\setcounter{enumi}{41}
\item \hyperref[chow-section-phantom]{Homologie de Chow}
\item \hyperref[intersection-section-phantom]{Théorie de l'intersection}
\item \hyperref[pic-section-phantom]{Schémas de Picard des courbes}
\item \hyperref[weil-section-phantom]{Théories cohomologiques de Weil}
\item \hyperref[adequate-section-phantom]{Modules adéquats}
\item \hyperref[dualizing-section-phantom]{Complexes dualisants}
\item \hyperref[duality-section-phantom]{Dualité pour les schémas}
\item \hyperref[discriminant-section-phantom]{Discriminants et différentes}
\item \hyperref[derham-section-phantom]{Cohomologie de de Rham}
\item \hyperref[local-cohomology-section-phantom]{Cohomologie locale}
\item \hyperref[algebraization-section-phantom]{Géométrie algébrique et formelle}
\item \hyperref[curves-section-phantom]{Courbes algébriques}
\item \hyperref[resolve-section-phantom]{Résolution des surfaces}
\item \hyperref[models-section-phantom]{Réduction semi-stable}
\item \hyperref[functors-section-phantom]{Foncteurs et morphismes}
\item \hyperref[equiv-section-phantom]{Catégories dérivées de variétés}
\item \hyperref[pione-section-phantom]{Groupes fondamentaux de schémas}
\item \hyperref[etale-cohomology-section-phantom]{Cohomologie étale}
\item \hyperref[crystalline-section-phantom]{Cohomologie cristalline}
\item \hyperref[proetale-section-phantom]{Cohomologie pro-étale}
\item \hyperref[relative-cycles-section-phantom]{Cycles relatifs}
\item \hyperref[more-etale-section-phantom]{Compléments de cohomologie étale}
\item \hyperref[trace-section-phantom]{La formule des traces}
\end{enumerate}
Espaces algébriques
\begin{enumerate}
\setcounter{enumi}{64}
\item \hyperref[spaces-section-phantom]{Espaces algébriques}
\item \hyperref[spaces-properties-section-phantom]{Propriétés des espaces algébriques}
\item \hyperref[spaces-morphisms-section-phantom]{Morphismes d'espaces algébriques}
\item \hyperref[decent-spaces-section-phantom]{Espaces algébriques décents}
\item \hyperref[spaces-cohomology-section-phantom]{Cohomologie des espaces algébriques}
\item \hyperref[spaces-limits-section-phantom]{Limites d'espaces algébriques}
\item \hyperref[spaces-divisors-section-phantom]{Diviseurs sur les espaces algébriques}
\item \hyperref[spaces-over-fields-section-phantom]{Espaces algébriques sur des corps}
\item \hyperref[spaces-topologies-section-phantom]{Topologies sur les espaces algébriques}
\item \hyperref[spaces-descent-section-phantom]{Descente et espaces algébriques}
\item \hyperref[spaces-perfect-section-phantom]{Catégories dérivées d'espaces}
\item \hyperref[spaces-more-morphisms-section-phantom]{Compléments sur les morphismes d'espaces}
\item \hyperref[spaces-flat-section-phantom]{Platitude sur les espaces algébriques}
\item \hyperref[spaces-groupoids-section-phantom]{Groupoïdes dans les espaces algébriques}
\item \hyperref[spaces-more-groupoids-section-phantom]{Compléments sur les groupoïdes dans les espaces}
\item \hyperref[bootstrap-section-phantom]{Amorçage}
\item \hyperref[spaces-pushouts-section-phantom]{Sommes amalgamées d'espaces algébriques}
\end{enumerate}
Thèmes de géométrie
\begin{enumerate}
\setcounter{enumi}{81}
\item \hyperref[spaces-chow-section-phantom]{Groupes de Chow des espaces}
\item \hyperref[groupoids-quotients-section-phantom]{Quotients de groupoïdes}
\item \hyperref[spaces-more-cohomology-section-phantom]{Compléments sur la cohomologie des espaces}
\item \hyperref[spaces-simplicial-section-phantom]{Espaces simpliciaux}
\item \hyperref[spaces-duality-section-phantom]{Dualité pour les espaces}
\item \hyperref[formal-spaces-section-phantom]{Espaces algébriques formels}
\item \hyperref[restricted-section-phantom]{Algébrisation des espaces formels}
\item \hyperref[spaces-resolve-section-phantom]{Résolution des surfaces revisitée}
\end{enumerate}
Théorie des déformations
\begin{enumerate}
\setcounter{enumi}{89}
\item \hyperref[formal-defos-section-phantom]{Théorie formelle des déformations}
\item \hyperref[defos-section-phantom]{Théorie des déformations}
\item \hyperref[cotangent-section-phantom]{Le complexe cotangent}
\item \hyperref[examples-defos-section-phantom]{Problèmes de déformation}
\end{enumerate}
Champs algébriques
\begin{enumerate}
\setcounter{enumi}{93}
\item \hyperref[algebraic-section-phantom]{Champs algébriques}
\item \hyperref[examples-stacks-section-phantom]{Exemples de champs}
\item \hyperref[stacks-sheaves-section-phantom]{Faisceaux sur les champs algébriques}
\item \hyperref[criteria-section-phantom]{Critères de représentabilité}
\item \hyperref[artin-section-phantom]{Axiomes d'Artin}
\item \hyperref[quot-section-phantom]{Espaces de Quot et de Hilbert}
\item \hyperref[stacks-properties-section-phantom]{Propriétés des champs algébriques}
\item \hyperref[stacks-morphisms-section-phantom]{Morphismes de champs algébriques}
\item \hyperref[stacks-limits-section-phantom]{Limites de champs algébriques}
\item \hyperref[stacks-cohomology-section-phantom]{Cohomologie des champs algébriques}
\item \hyperref[stacks-perfect-section-phantom]{Catégories dérivées de champs}
\item \hyperref[stacks-introduction-section-phantom]{Introduction aux champs algébriques}
\item \hyperref[stacks-more-morphisms-section-phantom]{Compléments sur les morphismes de champs}
\item \hyperref[stacks-geometry-section-phantom]{La géométrie des champs}
\end{enumerate}
Thèmes de théorie des espaces de modules
\begin{enumerate}
\setcounter{enumi}{107}
\item \hyperref[moduli-section-phantom]{Champs de modules}
\item \hyperref[moduli-curves-section-phantom]{Espaces de modules de courbes}
\end{enumerate}
Divers
\begin{enumerate}
\setcounter{enumi}{109}
\item \hyperref[examples-section-phantom]{Exemples}
\item \hyperref[exercises-section-phantom]{Exercices}
\item \hyperref[guide-section-phantom]{Guide bibliographique}
\item \hyperref[desirables-section-phantom]{Desiderata}
\item \hyperref[coding-section-phantom]{Style de programmation}
\item \hyperref[obsolete-section-phantom]{Obsolète}
\item \hyperref[fdl-section-phantom]{Licence de documentation libre GNU}
\item \hyperref[index-section-phantom]{Index généré automatiquement}
\end{enumerate}
\end{multicols}


\bibliography{my}
\bibliographystyle{amsalpha}

\end{document}
